% Options for packages loaded elsewhere
\PassOptionsToPackage{unicode}{hyperref}
\PassOptionsToPackage{hyphens}{url}
\PassOptionsToPackage{dvipsnames,svgnames,x11names}{xcolor}
%
\documentclass[
  10pt,
]{article}
\usepackage{amsmath,amssymb}
\usepackage{iftex}
\ifPDFTeX
  \usepackage[T1]{fontenc}
  \usepackage[utf8]{inputenc}
  \usepackage{textcomp} % provide euro and other symbols
\else % if luatex or xetex
  \usepackage{unicode-math} % this also loads fontspec
  \defaultfontfeatures{Scale=MatchLowercase}
  \defaultfontfeatures[\rmfamily]{Ligatures=TeX,Scale=1}
\fi
\usepackage{lmodern}
\ifPDFTeX\else
  % xetex/luatex font selection
  \setmainfont[]{DejaVu Serif}
  \setmonofont[]{DejaVu Sans Mono}
  \setmathfont[]{Latin Modern Math}
\fi
% Use upquote if available, for straight quotes in verbatim environments
\IfFileExists{upquote.sty}{\usepackage{upquote}}{}
\IfFileExists{microtype.sty}{% use microtype if available
  \usepackage[]{microtype}
  \UseMicrotypeSet[protrusion]{basicmath} % disable protrusion for tt fonts
}{}
\makeatletter
\@ifundefined{KOMAClassName}{% if non-KOMA class
  \IfFileExists{parskip.sty}{%
    \usepackage{parskip}
  }{% else
    \setlength{\parindent}{0pt}
    \setlength{\parskip}{6pt plus 2pt minus 1pt}}
}{% if KOMA class
  \KOMAoptions{parskip=half}}
\makeatother
\usepackage{xcolor}
\usepackage[margin=25mm]{geometry}
\usepackage{color}
\usepackage{fancyvrb}
\newcommand{\VerbBar}{|}
\newcommand{\VERB}{\Verb[commandchars=\\\{\}]}
\DefineVerbatimEnvironment{Highlighting}{Verbatim}{commandchars=\\\{\}}
% Add ',fontsize=\small' for more characters per line
\newenvironment{Shaded}{}{}
\newcommand{\AlertTok}[1]{\textcolor[rgb]{1.00,0.00,0.00}{\textbf{#1}}}
\newcommand{\AnnotationTok}[1]{\textcolor[rgb]{0.38,0.63,0.69}{\textbf{\textit{#1}}}}
\newcommand{\AttributeTok}[1]{\textcolor[rgb]{0.49,0.56,0.16}{#1}}
\newcommand{\BaseNTok}[1]{\textcolor[rgb]{0.25,0.63,0.44}{#1}}
\newcommand{\BuiltInTok}[1]{\textcolor[rgb]{0.00,0.50,0.00}{#1}}
\newcommand{\CharTok}[1]{\textcolor[rgb]{0.25,0.44,0.63}{#1}}
\newcommand{\CommentTok}[1]{\textcolor[rgb]{0.38,0.63,0.69}{\textit{#1}}}
\newcommand{\CommentVarTok}[1]{\textcolor[rgb]{0.38,0.63,0.69}{\textbf{\textit{#1}}}}
\newcommand{\ConstantTok}[1]{\textcolor[rgb]{0.53,0.00,0.00}{#1}}
\newcommand{\ControlFlowTok}[1]{\textcolor[rgb]{0.00,0.44,0.13}{\textbf{#1}}}
\newcommand{\DataTypeTok}[1]{\textcolor[rgb]{0.56,0.13,0.00}{#1}}
\newcommand{\DecValTok}[1]{\textcolor[rgb]{0.25,0.63,0.44}{#1}}
\newcommand{\DocumentationTok}[1]{\textcolor[rgb]{0.73,0.13,0.13}{\textit{#1}}}
\newcommand{\ErrorTok}[1]{\textcolor[rgb]{1.00,0.00,0.00}{\textbf{#1}}}
\newcommand{\ExtensionTok}[1]{#1}
\newcommand{\FloatTok}[1]{\textcolor[rgb]{0.25,0.63,0.44}{#1}}
\newcommand{\FunctionTok}[1]{\textcolor[rgb]{0.02,0.16,0.49}{#1}}
\newcommand{\ImportTok}[1]{\textcolor[rgb]{0.00,0.50,0.00}{\textbf{#1}}}
\newcommand{\InformationTok}[1]{\textcolor[rgb]{0.38,0.63,0.69}{\textbf{\textit{#1}}}}
\newcommand{\KeywordTok}[1]{\textcolor[rgb]{0.00,0.44,0.13}{\textbf{#1}}}
\newcommand{\NormalTok}[1]{#1}
\newcommand{\OperatorTok}[1]{\textcolor[rgb]{0.40,0.40,0.40}{#1}}
\newcommand{\OtherTok}[1]{\textcolor[rgb]{0.00,0.44,0.13}{#1}}
\newcommand{\PreprocessorTok}[1]{\textcolor[rgb]{0.74,0.48,0.00}{#1}}
\newcommand{\RegionMarkerTok}[1]{#1}
\newcommand{\SpecialCharTok}[1]{\textcolor[rgb]{0.25,0.44,0.63}{#1}}
\newcommand{\SpecialStringTok}[1]{\textcolor[rgb]{0.73,0.40,0.53}{#1}}
\newcommand{\StringTok}[1]{\textcolor[rgb]{0.25,0.44,0.63}{#1}}
\newcommand{\VariableTok}[1]{\textcolor[rgb]{0.10,0.09,0.49}{#1}}
\newcommand{\VerbatimStringTok}[1]{\textcolor[rgb]{0.25,0.44,0.63}{#1}}
\newcommand{\WarningTok}[1]{\textcolor[rgb]{0.38,0.63,0.69}{\textbf{\textit{#1}}}}
\usepackage{longtable,booktabs,array}
\usepackage{calc} % for calculating minipage widths
% Correct order of tables after \paragraph or \subparagraph
\usepackage{etoolbox}
\makeatletter
\patchcmd\longtable{\par}{\if@noskipsec\mbox{}\fi\par}{}{}
\makeatother
% Allow footnotes in longtable head/foot
\IfFileExists{footnotehyper.sty}{\usepackage{footnotehyper}}{\usepackage{footnote}}
\makesavenoteenv{longtable}
\setlength{\emergencystretch}{3em} % prevent overfull lines
\providecommand{\tightlist}{%
  \setlength{\itemsep}{0pt}\setlength{\parskip}{0pt}}
\setcounter{secnumdepth}{5}
\usepackage{amsmath,amssymb,amsthm,mathtools}
\usepackage{longtable,booktabs,array}
\usepackage{xurl}
\allowdisplaybreaks
\setlength{\emergencystretch}{4em}
\setcounter{tocdepth}{2}
\setlength{\parskip}{0.35em}
\providecommand{\tightlist}{\setlength{\itemsep}{0pt}\setlength{\parskip}{0pt}}
\newcommand{\Sym}{\operatorname{Sym}}
\newcommand{\GL}{\operatorname{GL}}
\newcommand{\SL}{\operatorname{SL}}
\newcommand{\Spec}{\operatorname{Spec}}
\newcommand{\Hom}{\operatorname{Hom}}
\ifLuaTeX
  \usepackage{selnolig}  % disable illegal ligatures
\fi
\usepackage{bookmark}
\IfFileExists{xurl.sty}{\usepackage{xurl}}{} % add URL line breaks if available
\urlstyle{same}
\hypersetup{
  pdftitle={GCT\_PLETH: Plethysm, theta correspondence, and Grothendieck operations},
  pdfauthor={GCT\_PLETH\_20261121\_run01 --- GPT-6 Astra Pro and delegated proof audits},
  colorlinks=true,
  linkcolor={Maroon},
  filecolor={Maroon},
  citecolor={Blue},
  urlcolor={Blue},
  pdfcreator={LaTeX via pandoc}}

\title{GCT\_PLETH: Plethysm, theta correspondence, and Grothendieck
operations}
\usepackage{etoolbox}
\makeatletter
\providecommand{\subtitle}[1]{% add subtitle to \maketitle
  \apptocmd{\@title}{\par {\large #1 \par}}{}{}
}
\makeatother
\subtitle{Checked proofs, explicit positive formulas, and a continuing
research ledger}
\author{GCT\_PLETH\_20261121\_run01 --- GPT-6 Astra Pro and delegated
proof audits}
\date{21 November 2026 --- working checkpoint}

\begin{document}
\maketitle

{
\hypersetup{linkcolor=}
\setcounter{tocdepth}{3}
\tableofcontents
}
\section{State of this research
record}\label{state-of-this-research-record}

This is a continuing research record for the supplied prompt. The
checked findings below are stated with their exact scope. The final part
separates the remaining mathematical and input gates. A checkpoint is a
dated snapshot of active work; it is not a declaration that the open
questions have been settled.

The main constructive result is a terminating positive lattice-point
procedure for every symmetric plethysm coefficient (G-009), every Schur
plethysm coefficient (G-010), and every image multiplicity of the
specified Hadamard--Howe map (G-011). It constructs the finite data from
flag minors, balanced torus invariants, finite-group averages, and
elimination. The proof gives an explicit invariant algebra or coherent
module and a bijection from its standard monomials to a finite disjoint
union of bounded polytope lattice sets. No previously computed plethysm
coefficient is used as an oracle.

The construction contains large preprocessing steps: a Hilbert basis, a
finite-group degree bound, and a Gröbner basis. Their size has no
polynomial bound in the varying parameters in this work. Mulmuley's GCT
VII, Problem 1.1, asks for a positive formula in the complexity class
\#P; the missing polynomial description and verification bounds are
recorded explicitly. The construction proved here gives finite positive
formulas in the literal lattice-count sense and leaves that complexity
requirement open.

Several formulas avoid the general elimination step. G-012 gives four
fixed polytopes for every symmetric outer-degree-three coefficient.
T-037 now gives a complete positive formula for every outer-degree-four
Schur label and every output partition, using at most 108 fixed
polytopes with at most five variables. Its proof evaluates the entire
four-point ideal model, including unequal opposite gaps, and its
evaluator uses a fixed number of integer arithmetic operations.
T-012--T-013 give finite explicit Littlewood--Richardson polytopes
through stable width six, and T-018 evaluates an entire width-seven
sector.

T-032 and G-015 construct an exact projective-frame ideal model in every
outer degree, with an explicit balanced-chart inverse and proof of every
transverse order. In outer degree five, T-038 and A-029 supply the
actual equivariant quadratic resolution, cubic algebra, Rees
presentation and linear powers. G-017 evaluates the quotient-bundle
module at every point order by an explicit matrix factorization and an
equivariant free multiplication map. G-019 evaluates all symmetric
quotient powers and all point orders together: its fixed positive
formula has at most sixteen active lattice tags, and its complete
vanishing set consists of 44 explicitly listed parameter-label
exceptions. G-018, T-040 and T-041 give further exterior,
shifted-quotient and hook families. These formulas include every outer
Schur label of five in their stated geometric parameter families.

For the actual canonical map, G-013 transports prime-adjacent square
preimages through an explicit block specialization and multiplies them,
with exact numerical semigroup certificates for ranks 2 through 100.
G-014 supplies explicit preimages for every target highest-weight
component whose tail has size at most the coordinate degree.
R-017--R-020 prove second-row injectivity and settle the entire
first-row boundary band of common polarization. R-021 supplies
determinant preimages at every even-rank quadratic boundary.
R-022--R-023 transfer these through an actual annihilator injection and
exact rectangular prolongation. R-024 proves recursive componentwise
surjectivity from the established degree-five base and the square-bound
bases, including positive exact image ranks below the full-map range.
The remaining critical annihilator has the smaller first- and second-row
support stated in R-023.

The archive audit now inventories 123 individual assertions: 80 proved,
25 refuted, and 18 exact gates. All thirteen supplied target leaves have
been compared with the original Montpellier scans. A-021 proves the
exact lattice-ring statement over every characteristic-zero discrete
valuation ring. A-023 gives an explicit characteristic-zero local-ring
counterexample even with base K0 equal to the integers. A-022 proves the
hyperbolic characteristic-class formula over every scheme with 2
invertible using an actual injective splitting construction. The
remaining historical and mathematical inputs are retained assertion by
assertion.

\subsection{Conventions and checking}\label{conventions-and-checking}

All representation multiplicities are in characteristic zero unless an
integral or modular ring is explicitly specified. The notation
\(a_{m,n}^{\lambda}\) means the multiplicity in
\(\operatorname{Sym}^m(\operatorname{Sym}^n V)\). The canonical map uses
normalized tensor averaging, and each integral comparison distinguishes
symmetric and divided powers by its actual map.

An identifier labels a finding or explicit reformulation produced in
this run; it does not claim historical priority for the underlying
coefficient family. Known results are attributed with exact locators.
Proof readings, independent re-proofs, finite computations, and
background uses are separate events in
\texttt{\small CHECK\_\allowbreak{}EVENTS.tsv}. The model is GPT-6 Astra
Pro; the runtime does not expose the reasoning-effort setting, so it is
recorded as unexposed. No Lean proof was executed, and none of the
prompt's named Lean source files was present in the supplied ZIP.

The package includes executable exact-arithmetic computations, complete
finite tables, proof-review reports, a question register, and the timed
ledger. The manuscript is assembled from the current proof reports; the
source register records the original literature and Grothendieck
locators. The entire checked mathematical argument is contained in the
present TeX source, with its finite certificates also available as
machine-readable files.

\part{Checked theta and positive formulas}

\section{Theta, balanced weights, and the finite
projection}\label{theta-balanced-weights-and-the-finite-projection}

Run:
\texttt{\small GCT\_\allowbreak{}PLETH-\allowbreak{}20261121-\allowbreak{}theta}.
Author/checker: GPT-6 Astra Pro, effort not exposed by the runtime.
Mathematical statements below are unconditional statements over
characteristic zero. They are exact realizations and explicit
consequences of established duality, with no claim of historical
novelty.

\subsection{Notation}\label{notation}

Let \texttt{\small V=\allowbreak{}C\textasciicircum{}r}, let
\texttt{\small m,\allowbreak{}n>=\allowbreak{}1}, and let
\texttt{\small lambda} be a partition of \texttt{\small mn}. Set

\[
a_{m,n}^{\lambda}=\dim\operatorname{Hom}_{GL(V)}
(S_\lambda V,\operatorname{Sym}^m(\operatorname{Sym}^n V))
\]

in the stable range
\texttt{\small r>=\allowbreak{}length(\allowbreak{}lambda)\allowbreak{}}.
More generally,
\texttt{\small p\_\allowbreak{}\{mu,\allowbreak{}nu\}\textasciicircum{}lambda}
is the coefficient of \texttt{\small s\_\allowbreak{}lambda} in
\texttt{\small s\_\allowbreak{}mu[\allowbreak{}s\_\allowbreak{}nu]\allowbreak{}}.
Specht modules are written
\texttt{\small [\allowbreak{}mu]\allowbreak{}}.

\subsection{T-001. The actual compact dual pair and its vacuum
character}\label{t-001.-the-actual-compact-dual-pair-and-its-vacuum-character}

Equip \texttt{\small V=\allowbreak{}C\textasciicircum{}r} and
\texttt{\small E=\allowbreak{}C\textasciicircum{}m} with positive
definite Hermitian forms. Put \texttt{\small W=\allowbreak{}V tensor E},
considered as a real vector space, with symplectic form
\texttt{\small omega=\allowbreak{}Im h\_\allowbreak{}W}. The images of
\texttt{\small U(\allowbreak{}V)\allowbreak{}} and
\texttt{\small U(\allowbreak{}E)\allowbreak{}} in
\texttt{\small Sp(\allowbreak{}W\_\allowbreak{}R)\allowbreak{}} are
mutual centralizers.

\textbf{Proof.} A real-linear map commuting with
\texttt{\small U(\allowbreak{}V)\allowbreak{}} commutes with scalar
multiplication by \texttt{\small i}, hence is complex-linear. The
complex commutant of the natural
\texttt{\small U(\allowbreak{}V)\allowbreak{}} action on
\texttt{\small V tensor E} is
\texttt{\small 1\_\allowbreak{}V tensor End\_\allowbreak{}C(\allowbreak{}E)\allowbreak{}}:
this follows either by the matrix units generated by the natural
representation or by Schur's lemma on the irreducible natural
\texttt{\small U(\allowbreak{}V)\allowbreak{}} module. A complex-linear
map preserving \texttt{\small omega} preserves
\texttt{\small h\_\allowbreak{}W}, because its real part is recovered
from \texttt{\small omega} and the complex structure. Thus the
symplectic centralizer is exactly
\texttt{\small 1\_\allowbreak{}V tensor U(\allowbreak{}E)\allowbreak{}}.
Interchanging the two factors proves the reverse centralizer identity.
This proof also covers \texttt{\small r=\allowbreak{}1} or
\texttt{\small m=\allowbreak{}1}.

Choose the positive-energy oscillator convention in which the polynomial
Fock module on a Hermitian space \texttt{\small W} is

\[
\mathcal F_{\mathrm{alg}}(W)=\operatorname{Sym}(W)\otimes\chi_0,
\qquad \chi_0^2=\det_W
\]

on the metaplectic preimage of
\texttt{\small U(\allowbreak{}W)\allowbreak{}}. Explicitly this preimage
is the double cover of \texttt{\small U(\allowbreak{}W)\allowbreak{}}
with elements
\texttt{\small (\allowbreak{}u,\allowbreak{}zeta)\allowbreak{}}
satisfying
\texttt{\small zeta\textasciicircum{}2=\allowbreak{}det(\allowbreak{}u)\allowbreak{}},
and the action on the homogeneous polynomial summand is
\texttt{\small zeta Sym\textasciicircum{}d(\allowbreak{}u)\allowbreak{}}.
Its infinitesimal action is
\texttt{\small d Sym(\allowbreak{}A)\allowbreak{}+\allowbreak{}(\allowbreak{}tr\_\allowbreak{}W A)\allowbreak{}/\allowbreak{}2};
on a diagonal rotation in the \texttt{\small j}th coordinate the
eigenvalue is
\texttt{\small k\_\allowbreak{}j+\allowbreak{}1/\allowbreak{}2}, the
usual number-operator realization of the oscillator. This specifies the
convention and the character completely. The contragredient convention
dualizes the polynomial space and inverts the displayed character.

For the homomorphism
\texttt{\small U(\allowbreak{}V)\allowbreak{} x U(\allowbreak{}E)\allowbreak{} -\allowbreak{}> U(\allowbreak{}W)\allowbreak{}}
the vacuum character has square

\[
(\det_V)^m(\det_E)^r.
\]

Consequently twisting the restricted oscillator by
\texttt{\small chi\_\allowbreak{}0\textasciicircum{}\{-\allowbreak{}1\}}
gives exactly the polynomial Cauchy module. The twist is part of the
dictionary; it is not discarded silently. On permutation matrices in
\texttt{\small U(\allowbreak{}E)\allowbreak{}}, its square is
\texttt{\small sgn\textasciicircum{}r}.

\subsection{T-002. Balanced-weight Howe
realization}\label{t-002.-balanced-weight-howe-realization}

Let

\[
M_{\lambda,n}=(S_\lambda E)_{(n,\ldots,n)}.
\]

Here \texttt{\small S\_\allowbreak{}m} acts by literal permutation
matrices on \texttt{\small E}, not by an unspecified choice of
representatives in \texttt{\small SU(\allowbreak{}m)\allowbreak{}}.
There is an isomorphism of
\texttt{\small GL(\allowbreak{}V)\allowbreak{} x S\_\allowbreak{}m}
modules

\[
(\operatorname{Sym}^n V)^{\otimes m}
\simeq \bigoplus_{\lambda\vdash mn} S_\lambda V\otimes M_{\lambda,n}.
\tag{T2.1}
\]

In particular,

\[
\boxed{a_{m,n}^{\lambda}=\dim M_{\lambda,n}^{S_m}}.
\tag{T2.2}
\]

\textbf{Proof of the precise realization.} Polynomial Howe/Cauchy
duality is

\[
\operatorname{Sym}(V\otimes E)
\simeq\bigoplus_\lambda S_\lambda V\otimes S_\lambda E.
\]

Take the weight
\texttt{\small (\allowbreak{}n,\allowbreak{}...,\allowbreak{}n)\allowbreak{}}
for the diagonal torus acting on \texttt{\small E}. On the left, a
polynomial has degree \texttt{\small n} in each of the \texttt{\small m}
column spaces \texttt{\small V tensor e\_\allowbreak{}i}; this weight
space is
\texttt{\small (\allowbreak{}Sym\textasciicircum{}n V)\allowbreak{}\textasciicircum{}\{tensor m\}}.
A permutation matrix permutes those column spaces and gives exactly the
ordinary tensor-factor action. The right side gives (T2.1). Taking
\texttt{\small S\_\allowbreak{}m} invariants gives (T2.2), because
characteristic-zero symmetrization identifies invariant tensors with
\texttt{\small Sym\textasciicircum{}m(\allowbreak{}Sym\textasciicircum{}n V)\allowbreak{}}.

The exact literature locator in the supplied source is Landsberg,
\emph{Geometric Complexity Theory: an introduction for geometers},
subsection ``A duality theorem for weight zero spaces and plethysms,''
theorem \texttt{\small gaythm}, lines 766--798; it attributes the
statement to Gay (1976). The displayed corollary \texttt{\small gaycor}
(ii) writes \texttt{\small S\_\allowbreak{}d} where this construction
requires the permutation group on
\texttt{\small dim(\allowbreak{}E)\allowbreak{}=\allowbreak{}m} letters.
The proof above fixes the group and the action rather than relying on
that typographical line.

\subsection{T-003. General Schur plethysm, both oscillator and
wreath-product
versions}\label{t-003.-general-schur-plethysm-both-oscillator-and-wreath-product-versions}

For \texttt{\small mu} a partition of \texttt{\small m}, the
balanced-weight identity refines to

\[
\boxed{p_{\mu,(n)}^\lambda
=\dim\operatorname{Hom}_{S_m}([\mu],M_{\lambda,n}).}
\tag{T3.1}
\]

Indeed, Schur--Weyl duality on
\texttt{\small (\allowbreak{}Sym\textasciicircum{}n V)\allowbreak{}\textasciicircum{}\{tensor m\}}
gives
\texttt{\small sum\_\allowbreak{}mu S\_\allowbreak{}mu(\allowbreak{}Sym\textasciicircum{}n V)\allowbreak{} tensor [\allowbreak{}mu]\allowbreak{}};
comparing it with (T2.1) proves the formula.

For arbitrary \texttt{\small nu} of size \texttt{\small n}, write
\texttt{\small W\_\allowbreak{}nu=\allowbreak{}S\_\allowbreak{}nu V}.
Its invariant Hermitian form gives a homomorphism
\texttt{\small rho\_\allowbreak{}nu:\allowbreak{}U(\allowbreak{}V)\allowbreak{}-\allowbreak{}>U(\allowbreak{}W\_\allowbreak{}nu)\allowbreak{}}.
The compact dual pair
\texttt{\small U(\allowbreak{}W\_\allowbreak{}nu)\allowbreak{},\allowbreak{}U(\allowbreak{}E)\allowbreak{}}
on the polynomial Fock space
\texttt{\small Sym(\allowbreak{}W\_\allowbreak{}nu tensor E)\allowbreak{}}
has Howe component

\[
\operatorname{Hom}_{U(E)}(S_\mu E,
\operatorname{Sym}(W_\nu\otimes E))=S_\mu(W_\nu).
\tag{T3.2}
\]

Here \texttt{\small dim E=\allowbreak{}m} suffices, and the vacuum twist
is handled as in T-001. Restrict the first member along
\texttt{\small rho\_\allowbreak{}nu}; then the multiplicity of
\texttt{\small S\_\allowbreak{}lambda V} in (T3.2) is precisely
\texttt{\small p\_\allowbreak{}\{mu,\allowbreak{}nu\}\textasciicircum{}lambda}.
This is an actual compact-oscillator realization. The restriction along
the explicitly specified polynomial representation
\texttt{\small rho\_\allowbreak{}nu} is the remaining branching
computation.

For a finite description, identify the \texttt{\small mn} tensor
positions with an \texttt{\small m} by \texttt{\small n} grid and put
\texttt{\small H=\allowbreak{}S\_\allowbreak{}n wr S\_\allowbreak{}m}.
Let

\[
U_{\mu,\nu}=[\nu]^{\otimes m}\otimes\operatorname{Inf}_{S_m}^H[\mu],
\]

where block permutations permute the tensor factors in
\texttt{\small [\allowbreak{}nu]\allowbreak{}\textasciicircum{}\{tensor m\}}.
There is a natural isomorphism

\[
S_\mu(S_\nu V)\simeq
\operatorname{Hom}_{H}(U_{\mu,\nu},V^{\otimes mn}),
\qquad
p_{\mu,\nu}^\lambda=
\dim\operatorname{Hom}_H(U_{\mu,\nu},[\lambda]|_H).
\tag{T3.3}
\]

\textbf{Proof.} Apply \texttt{\small S\_\allowbreak{}n}-Schur--Weyl
duality separately to the \texttt{\small m} blocks. Taking
\texttt{\small Hom\_\allowbreak{}\{S\_\allowbreak{}n\textasciicircum{}m\}(\allowbreak{}[\allowbreak{}nu]\allowbreak{}\textasciicircum{}\{tensor m\},\allowbreak{}-\allowbreak{})\allowbreak{}}
gives
\texttt{\small (\allowbreak{}S\_\allowbreak{}nu V)\allowbreak{}\textasciicircum{}\{tensor m\}}
and preserves the block-permutation action. Taking
\texttt{\small Hom\_\allowbreak{}\{S\_\allowbreak{}m\}(\allowbreak{}[\allowbreak{}mu]\allowbreak{},\allowbreak{}-\allowbreak{})\allowbreak{}}
gives
\texttt{\small S\_\allowbreak{}mu(\allowbreak{}S\_\allowbreak{}nu V)\allowbreak{}}.
This nested Hom space is precisely the Hom space for the semidirect
product \texttt{\small H}. Full Schur--Weyl duality for
\texttt{\small V\textasciicircum{}\{tensor mn\}} then gives the second
equality. Every map used is functorial in \texttt{\small V}.

\subsection{T-004. The exact positive polytope count produced by
interlacing}\label{t-004.-the-exact-positive-polytope-count-produced-by-interlacing}

For
\texttt{\small length(\allowbreak{}lambda)\allowbreak{}<=\allowbreak{}m},
pad \texttt{\small lambda} by zeros to length \texttt{\small m}; for
greater length the balanced space and symmetric plethysm coefficient
vanish. Define
\texttt{\small GT(\allowbreak{}lambda;\allowbreak{}n\textasciicircum{}m)\allowbreak{}}
by integer arrays
\texttt{\small x\_\allowbreak{}i\textasciicircum{}(\allowbreak{}k)\allowbreak{}},
\texttt{\small 1<=\allowbreak{}i<=\allowbreak{}k<=\allowbreak{}m}, with

\[
x_i^{(m)}=\lambda_i,\qquad
x_i^{(k+1)}\ge x_i^{(k)}\ge x_{i+1}^{(k+1)},\qquad
\sum_{i=1}^k x_i^{(k)}=kn\quad(1\le k\le m).
\tag{T4.1}
\]

The inequalities define a bounded rational polytope. The usual
\texttt{\small U(\allowbreak{}m)\allowbreak{}} to
\texttt{\small U(\allowbreak{}m-\allowbreak{}1)\allowbreak{}}
interlacing rule, iterated, gives

\[
\#(GT(\lambda;n^m)\cap\mathbb Z^{m(m+1)/2})
=\dim M_{\lambda,n}
=K_{\lambda,(n^m)}
=\boxed{\sum_{\mu\vdash m}(\dim[\mu])p_{\mu,(n)}^\lambda}.
\tag{T4.2}
\]

\textbf{Proof of the specified specialization.} Iterated branching
associates a one-dimensional torus weight line to every interlacing
pattern. Its \texttt{\small k}th torus weight is
\texttt{\small sum\_\allowbreak{}i x\_\allowbreak{}i\textasciicircum{}(\allowbreak{}k)\allowbreak{}-\allowbreak{}sum\_\allowbreak{}i x\_\allowbreak{}i\textasciicircum{}(\allowbreak{}k-\allowbreak{}1)\allowbreak{}}.
The row-sum equations in (T4.1) say exactly that every such weight
equals \texttt{\small n}, so they select a basis of
\texttt{\small M\_\allowbreak{}\{lambda,\allowbreak{}n\}}. Decomposition
of this \texttt{\small S\_\allowbreak{}m} module into Specht modules,
with multiplicities supplied by (T3.1), yields the final equality. The
standard interlacing theorem itself is carried as background; this is
its explicit balanced-weight specialization.

Thus the interlacing count controls the displayed positive weighted sum
of outer Schur plethysms. The desired symmetric outer power is the
summand indexed by
\texttt{\small mu=\allowbreak{}(\allowbreak{}m)\allowbreak{}}, and T-002
identifies its projector exactly.

\subsection{T-005. Determinant twist, parity obstruction, and
transport}\label{t-005.-determinant-twist-parity-obstruction-and-transport}

On the determinant-normalized representation

\[
\widetilde M=S_\lambda E\otimes(\det E)^{-n},
\]

the balanced weight becomes torus weight zero. Its permutation action on
that zero-weight space is

\[
\widetilde M_0\simeq M_{\lambda,n}\otimes\operatorname{sgn}^{n}.
\tag{T5.1}
\]

Consequently

\[
a_{m,n}^\lambda
=\dim\operatorname{Hom}_{S_m}(\operatorname{sgn}^{n},\widetilde M_0).
\tag{T5.2}
\]

The determinant of a permutation matrix is its sign, which proves
(T5.1); tensoring by a one-dimensional character proves (T5.2). This
supplies the exact morphism for normalization, including its finite
character.

For \texttt{\small m>=\allowbreak{}2} and
\texttt{\small lambda=\allowbreak{}(\allowbreak{}n\textasciicircum{}m)\allowbreak{}},
\texttt{\small S\_\allowbreak{}lambda E=\allowbreak{}(\allowbreak{}det E)\allowbreak{}\textasciicircum{}n}.
Its balanced weight space has dimension one and the permutation action
is \texttt{\small sgn\textasciicircum{}n}. Hence

\[
\boxed{a_{m,n}^{(n^m)}=\begin{cases}1&n\text{ even},\\0&n\text{ odd}.\end{cases}}
\tag{T5.3}
\]

There is exactly one interlacing pattern in (T4.1), since each row is
forced to consist of \texttt{\small n}s. For odd \texttt{\small n}, the
polytope count is one and the desired coefficient is zero. This refutes
the specific proposal of using those unprojected interlacing counts for
the symmetric plethysm coefficient. It makes no claim about all possible
positive formulas. The same pattern contributes to outer Schur type
\texttt{\small (\allowbreak{}1\textasciicircum{}m)\allowbreak{}} in
(T4.2).

More generally, for \texttt{\small t>=\allowbreak{}0}, functorial
determinant twisting gives

\[
S_{\lambda+(t^m)}E\simeq S_\lambda E\otimes(\det E)^t.
\]

It follows that

\[
p_{\mu,(n+t)}^{\lambda+(t^m)}
=\begin{cases}
p_{\mu,(n)}^\lambda&t\text{ even},\\
p_{\mu',(n)}^\lambda&t\text{ odd},
\end{cases}
\tag{T5.4}
\]

where \texttt{\small mu'} is the conjugate partition. In particular,

\[
\boxed{a_{m,n+2t}^{\lambda+(2t)^m}=a_{m,n}^{\lambda}}.
\tag{T5.5}
\]

\textbf{Proof.} The balanced weight shifts from
\texttt{\small (\allowbreak{}n\textasciicircum{}m)\allowbreak{}} to
\texttt{\small (\allowbreak{}(\allowbreak{}n+\allowbreak{}t)\allowbreak{}\textasciicircum{}m)\allowbreak{}};
its \texttt{\small S\_\allowbreak{}m} representation tensors with
\texttt{\small sgn\textasciicircum{}t}. The identity
\texttt{\small [\allowbreak{}mu]\allowbreak{} tensor sgn=\allowbreak{}[\allowbreak{}mu']\allowbreak{}}
and T-003 give (T5.4). Equation (T5.5) is its even case. This is an
exact transport of every nonvanishing, and of every vanishing, on these
rays.

The rectangular parity statement already occurs in the supplied
Landsberg source, exercise \texttt{\small snssv} immediately after
theorem \texttt{\small gaythm}, with attribution to Howe, Proposition
4.3(a). T5.1--T5.5 record the explicit character dictionary needed here;
no novelty claim is made for determinant-shift identities.

\subsection{T-006. A complete positive formula for outer degree
two}\label{t-006.-a-complete-positive-formula-for-outer-degree-two}

For \texttt{\small n>=\allowbreak{}1},

\[
\boxed{\operatorname{Sym}^2(\operatorname{Sym}^n V)
=\bigoplus_{j=0}^{\lfloor n/2\rfloor}S_{(2n-2j,2j)}V.}
\tag{T6.1}
\]

The equality is understood in characteristic zero; summands of excessive
length vanish. Its coefficient is the lattice count

\[
a_{2,n}^\lambda
=\#\{j\in\mathbb Z:0\le j\le n/2,
\lambda=(2n-2j,2j)\}.
\tag{T6.2}
\]

\textbf{Short proof in the balanced-weight form.} Partitions of
\texttt{\small 2n} occurring in T-002 have length at most two, so write
\texttt{\small lambda=\allowbreak{}(\allowbreak{}2n-\allowbreak{}k,\allowbreak{}k)\allowbreak{}},
\texttt{\small 0<=\allowbreak{}k<=\allowbreak{}n}. For
\texttt{\small E=\allowbreak{}C\textasciicircum{}2},

\[
S_{(2n-k,k)}E=\operatorname{Sym}^{2n-2k}E\otimes(\det E)^k.
\]

Its balanced weight line is spanned by
\texttt{\small e\_\allowbreak{}1\textasciicircum{}\{n-\allowbreak{}k\}e\_\allowbreak{}2\textasciicircum{}\{n-\allowbreak{}k\}}
tensored with the determinant factor. The transposition fixes the
monomial and multiplies the determinant factor by
\texttt{\small (\allowbreak{}-\allowbreak{}1)\allowbreak{}\textasciicircum{}k}.
T-002 retains exactly even \texttt{\small k=\allowbreak{}2j}, each with
multiplicity one. This provides a positive formula after evaluating the
finite projector, rather than dropping it. The underlying decomposition
is classical; this proof is included only to exhibit the exact
Howe-to-polytope dictionary.

For the smallest sign obstruction,
\texttt{\small Sym\textasciicircum{}2(\allowbreak{}Sym\textasciicircum{}3 V)\allowbreak{}=\allowbreak{}S\_\allowbreak{}(\allowbreak{}6)\allowbreak{}V+\allowbreak{}S\_\allowbreak{}(\allowbreak{}4,\allowbreak{}2)\allowbreak{}V},
whereas the GT count for
\texttt{\small (\allowbreak{}3,\allowbreak{}3)\allowbreak{}} equals one.

\subsection{T-007. Inner degree two: established multiplicity-free
positive
formula}\label{t-007.-inner-degree-two-established-multiplicity-free-positive-formula}

The classical symmetric-matrix identity is

\[
\operatorname{Sym}^m(\operatorname{Sym}^2 V)
=\bigoplus_{\alpha\vdash m}S_{2\alpha}V.
\tag{T7.1}
\]

An exact locator is Weyman, \emph{Cohomology of Vector Bundles and
Syzygies}, Proposition 2.3.8; Macdonald, \emph{Symmetric Functions and
Hall Polynomials}, second edition, Chapter I, §8, Example 6. It is
independently quoted with these locators in Fischer--Ikenmeyer,
\emph{The computational complexity of plethysm coefficients}, footnote 4
immediately before Theorem 3.5,
\url{https://pure.mpg.de/rest/items/item_3265369_1/component/file_3265370/content?download=true}
. The original books were not opened in this subtask, so the check event
records the exact extent of verification.

Thus
\texttt{\small a\_\allowbreak{}(\allowbreak{}m,\allowbreak{}2)\allowbreak{}\textasciicircum{}lambda}
is the number of integer points \texttt{\small alpha} satisfying
\texttt{\small 2alpha=\allowbreak{}lambda},
\texttt{\small alpha\_\allowbreak{}1>=\allowbreak{}...>=\allowbreak{}alpha\_\allowbreak{}r>=\allowbreak{}0},
\texttt{\small sum alpha\_\allowbreak{}i=\allowbreak{}m}: zero or one.
Combined with T-006 this proves the Foulkes coefficient inequality for
the pair of outer degrees \texttt{\small 2,\allowbreak{}m}: every
constituent of
\texttt{\small Sym\textasciicircum{}2(\allowbreak{}Sym\textasciicircum{}m V)\allowbreak{}}
occurs once in
\texttt{\small Sym\textasciicircum{}m(\allowbreak{}Sym\textasciicircum{}2 V)\allowbreak{}}.
This is a coefficient statement; the actual Hadamard--Howe rank is
handled by the grid projector supplied separately to the root agent.

\subsection{T-008. Explicit additive support
semigroups}\label{t-008.-explicit-additive-support-semigroups}

For fixed \texttt{\small m}, define

\[
\Gamma_m=\{(n,\lambda):|\lambda|=mn,\ell(\lambda)\le m,
a_{m,n}^\lambda>0\}.
\]

This set is closed under addition
\texttt{\small (\allowbreak{}n,\allowbreak{}lambda)\allowbreak{}+\allowbreak{}(\allowbreak{}n',\allowbreak{}lambda')\allowbreak{}=\allowbreak{}(\allowbreak{}n+\allowbreak{}n',\allowbreak{}lambda+\allowbreak{}lambda')\allowbreak{}}.

\textbf{Proof in a domain, with all group actions specified.} Let
\texttt{\small E=\allowbreak{}F=\allowbreak{}C\textasciicircum{}m}, let
\texttt{\small U\_\allowbreak{}F} be a maximal upper unipotent subgroup,
and consider

\[
R=\operatorname{Sym}(E\otimes F)^{U_F}.
\]

It is a subring of a polynomial ring, hence a domain. Cauchy duality
gives
\texttt{\small R=\allowbreak{}sum\_\allowbreak{}lambda S\_\allowbreak{}lambda E tensor L\_\allowbreak{}lambda},
where \texttt{\small L\_\allowbreak{}lambda} is the one-dimensional
highest-weight line of \texttt{\small S\_\allowbreak{}lambda F}.
Multiplication respects the right-torus weight, sending the
\texttt{\small lambda} and \texttt{\small lambda'} summands to
\texttt{\small lambda+\allowbreak{}lambda'}. It also respects the left
diagonal-torus weights and the left permutation action. Nonzero vectors
in
\texttt{\small M\_\allowbreak{}(\allowbreak{}lambda,\allowbreak{}n)\allowbreak{}\textasciicircum{}\{S\_\allowbreak{}m\} tensor L\_\allowbreak{}lambda}
and
\texttt{\small M\_\allowbreak{}(\allowbreak{}lambda',\allowbreak{}n')\allowbreak{}\textasciicircum{}\{S\_\allowbreak{}m\} tensor L\_\allowbreak{}lambda'}
have nonzero product, which has balanced weight
\texttt{\small (\allowbreak{}n+\allowbreak{}n')\allowbreak{}\textasciicircum{}m},
right highest weight \texttt{\small lambda+\allowbreak{}lambda'}, and
remains permutation invariant. T-002 gives the claimed additive support.

For fixed \texttt{\small n}, the analogous support set

\[
\Gamma^{n}=\{(m,\lambda):a_{m,n}^\lambda>0\}
\]

is also additive: use the domain
\texttt{\small Sym(\allowbreak{}Sym\textasciicircum{}n V)\allowbreak{}}
and the product of nonzero highest-weight vectors. The preceding proof
specifies the less immediate fixed-outer-degree realization, which is
useful for transporting nonvanishing in the direction of inner degree.

These semigroup properties construct further positive instances from
positive seeds. T-005 gives their exact determinant translations.
Neither calculation evaluates every finite
\texttt{\small S\_\allowbreak{}m} projector.

\subsection{T-009. A precise obstruction to homogeneous Ehrhart
formulas}\label{t-009.-a-precise-obstruction-to-homogeneous-ehrhart-formulas}

Kahle--Michałek, \emph{Obstructions to combinatorial formulas for
plethysm}, Electronic Journal of Combinatorics 25(1) (2018), P1.41,
Theorem 5 and Remark 6, give

\[
a_{3,4s}^{(7s,5s)}=a_{4s,3}^{(7s,5s)}
=\phi(s)=\frac{s+r(s)}3,
\]

where
\texttt{\small r(\allowbreak{}0)\allowbreak{},\allowbreak{}...,\allowbreak{}r(\allowbreak{}5)\allowbreak{}=\allowbreak{}(\allowbreak{}3,\allowbreak{}-\allowbreak{}1,\allowbreak{}1,\allowbreak{}0,\allowbreak{}2,\allowbreak{}-\allowbreak{}2)\allowbreak{}}
periodically. This function is not the Ehrhart function of a rational
polytope:
\texttt{\small phi(\allowbreak{}1)\allowbreak{}=\allowbreak{}0} and
\texttt{\small -\allowbreak{}phi(\allowbreak{}-\allowbreak{}1)\allowbreak{}=\allowbreak{}1}
violate Ehrhart reciprocity. The theorem also excludes one-dimensional
inhomogeneous polytopes. Its exact scope leaves higher-dimensional
inhomogeneous constructions untreated. Source:
\url{https://www.combinatorics.org/ojs/index.php/eljc/article/download/v25i1p41/pdf/}
, printed page 4, Theorem 5 and Remark 6. No reproof of the known
theorem is claimed. The attached computation checks the displayed values
by independent partition counting for
\texttt{\small 1<=\allowbreak{}s<=\allowbreak{}40}.

\subsection{Exact projection and subsequent positive-count
constructions}\label{exact-projection-and-subsequent-positive-count-constructions}

For symmetric plethysm the exact finite projection is

\[
\operatorname{rank}\left(\frac1{m!}\sum_{\sigma\in S_m}
\rho_{\lambda,n}(\sigma)\right)
\quad\text{on }M_{\lambda,n}.
\]

T-004 computes the dimension of its ambient balanced-weight space, and
T-005 proves by an explicit determinant-sign calculation that this
dimension alone can differ from the projected dimension. Subsequent work
now evaluates the projected space by a terminating positive-count
construction: G-009 in
\texttt{\small work/\allowbreak{}root/\allowbreak{}invariant\_\allowbreak{}polytope\_\allowbreak{}report.md}
gives an explicit invariant-algebra presentation and standard-monomial
polytope sums for every symmetric plethysm coefficient. G-010 in
\texttt{\small work/\allowbreak{}root/\allowbreak{}general\_\allowbreak{}polytope\_\allowbreak{}extensions.md}
extends this to general Schur plethysm; G-011 gives the analogous
positive counts for the actual canonical-map image, kernel, and cokernel
modules. Their input consists of specified polynomials and finite group
actions, with no coefficient oracle. Their output size and running time
have no proved polynomial bound in the varying parameters.

The smaller direct formulas remain useful: T-006--T-007 evaluate
classical families, G-012 in
\texttt{\small work/\allowbreak{}root/\allowbreak{}rank\_\allowbreak{}three\_\allowbreak{}presentation.md}
gives four explicit polytopes for the complete outer-degree-three
family, T-012--T-013 give the stable width-five and width-six formulas,
and T-016--T-020 in
\texttt{\small work/\allowbreak{}theta/\allowbreak{}width7/\allowbreak{}width7\_\allowbreak{}report.md}
evaluate the three-parameter end-node sector of the width-seven
exterior-power factor. T-024--T-028 in
\texttt{\small work/\allowbreak{}theta/\allowbreak{}end\_\allowbreak{}node/\allowbreak{}}
now give a uniform small polytope and an arithmetic nonvanishing test
for every outer Schur functor with inner symmetric power and output
shape
\texttt{\small (\allowbreak{}alpha,\allowbreak{}beta\textasciicircum{}(\allowbreak{}m-\allowbreak{}2)\allowbreak{},\allowbreak{}gamma)\allowbreak{}},
together with an explicit toric degeneration for the symmetric-outer
algebra. T-009 excludes its precisely stated homogeneous Ehrhart model
and is consistent with these finite sums of graded fibers. The remaining
research problem in this direction is a small, efficient construction
outside the proved direct families, together with uniform inequalities
for the still-open canonical-map parameter families. The existence of a
terminating positive count is no longer left as an unevaluated gate.

\subsection{Source and checking notes}\label{source-and-checking-notes}

\begin{enumerate}
\def\labelenumi{\arabic{enumi}.}
\tightlist
\item
  Supplied Landsberg source was read at the theorem/proof and
  surrounding definitions; group actions and the two noted source typos
  were checked directly.
\item
  The compact-centralizer proof and explicit vacuum-character
  calculation in T-001 were read in full after drafting; the action
  convention is included to prevent a hidden metaplectic twist.
\item
  All derived proofs T-002--T-006 and T-008 were read in full; details
  of the balanced permutation action, determinant sign, and
  multiplication domain were checked.
\item
  T-007 is carried as an established identity with checked secondary
  exact locators; its original book proofs have not been read in this
  run.
\item
  T-009 is a cited theorem whose statement and reciprocity calculation
  were read in the primary article. Its complete inhomogeneous-polytope
  proof was read; no independent proof of that stronger clause is
  claimed.
\end{enumerate}

\section{Additional exact constructions and
limits}\label{additional-exact-constructions-and-limits}

Run
\texttt{\small GCT\_\allowbreak{}PLETH\_\allowbreak{}20261121\_\allowbreak{}run01};
model GPT-6 Astra Pro; effort setting not exposed. Findings below are
proved formulations for this investigation; no priority claim is made.
All vector spaces in the representation calculations have characteristic
zero unless an integral model is explicitly specified.

\subsection{G-001. A balanced weight space with no monomial
basis}\label{g-001.-a-balanced-weight-space-with-no-monomial-basis}

Let \texttt{\small E=\allowbreak{}C\textasciicircum{}5},
\texttt{\small lambda=\allowbreak{}(\allowbreak{}5,\allowbreak{}2,\allowbreak{}1,\allowbreak{}1,\allowbreak{}1)\allowbreak{}},
and let permutation matrices act on \[
M=(S_{(5,2,1,1,1)}E)_{(2,2,2,2,2)}.
\] Then \texttt{\small M} is isomorphic as an
\texttt{\small S\_\allowbreak{}5} representation to
\texttt{\small sgn tensor U\_\allowbreak{}5}, where
\texttt{\small U\_\allowbreak{}5=\allowbreak{}\{(\allowbreak{}z\_\allowbreak{}1,\allowbreak{}...,\allowbreak{}z\_\allowbreak{}5)\allowbreak{}:\allowbreak{}sum z\_\allowbreak{}i=\allowbreak{}0\}}.
There is no basis of \texttt{\small M} that
\texttt{\small S\_\allowbreak{}5} permutes up to nonzero scalar factors.
In particular there is no such basis indexed by its four balanced
Gelfand--Tsetlin patterns.

\textbf{Proof.} The multiplication map \[
\operatorname{Sym}^{4}E\otimes E\longrightarrow\operatorname{Sym}^{5}E
\] is surjective, and its kernel is
\texttt{\small S\_\allowbreak{}(\allowbreak{}4,\allowbreak{}1)\allowbreak{}E},
by the degree-one Pieri rule. On weight
\texttt{\small (\allowbreak{}1\textasciicircum{}5)\allowbreak{}} the
source has basis \[
v_i=\left(\prod_{j\ne i}e_j\right)\otimes e_i\quad(1\le i\le5).
\] Permutation matrices permute these five vectors, and multiplication
sends all of them to
\texttt{\small e\_\allowbreak{}1e\_\allowbreak{}2e\_\allowbreak{}3e\_\allowbreak{}4e\_\allowbreak{}5}.
The kernel on this weight is exactly \texttt{\small U\_\allowbreak{}5}.
Tensoring with \texttt{\small det E} shifts the highest weight from
\texttt{\small (\allowbreak{}4,\allowbreak{}1,\allowbreak{}0,\allowbreak{}0,\allowbreak{}0)\allowbreak{}}
to
\texttt{\small (\allowbreak{}5,\allowbreak{}2,\allowbreak{}1,\allowbreak{}1,\allowbreak{}1)\allowbreak{}}
and the balanced weight from
\texttt{\small (\allowbreak{}1\textasciicircum{}5)\allowbreak{}} to
\texttt{\small (\allowbreak{}2\textasciicircum{}5)\allowbreak{}}; its
restriction to permutation matrices is the sign character. This gives
the stated isomorphism, including the character twist.

A monomial basis of size four would give a homomorphism
\texttt{\small A\_\allowbreak{}5→\allowbreak{}S\_\allowbreak{}4}. The
group \texttt{\small A\_\allowbreak{}5} is simple: its conjugacy-class
sizes are
\texttt{\small 1,\allowbreak{}15,\allowbreak{}20,\allowbreak{}12,\allowbreak{}12},
and no sum containing 1 of a proper nonempty subcollection of the four
nonidentity classes is a proper divisor of 60. To verify these class
sizes, the double transpositions number 15 and have centralizer of order
4 in \texttt{\small A\_\allowbreak{}5}; the 3-cycles number 20 and have
centralizer of order 3; the 24 five-cycles split into two classes of
size 12 since their centralizer has order 5. Lagrange's theorem proves
simplicity from the listed sizes. The homomorphism to
\texttt{\small S\_\allowbreak{}4} therefore has trivial image, since 60
exceeds 24. Thus every element of \texttt{\small A\_\allowbreak{}5}
would act diagonally in the proposed basis. The resulting image is
abelian. A nonabelian simple group has trivial abelian image, so
\texttt{\small A\_\allowbreak{}5} would act trivially on
\texttt{\small M}. But \texttt{\small (\allowbreak{}123)\allowbreak{}}
acts nontrivially on
\texttt{\small e\_\allowbreak{}1-\allowbreak{}e\_\allowbreak{}2} in
\texttt{\small U\_\allowbreak{}5}. On \texttt{\small A\_\allowbreak{}5}
the sign twist is trivial. This is a contradiction.

The argument identifies a specific failure of an orbit-counting approach
to the balanced-weight normalizer action. It makes no assertion
excluding other lattice-point descriptions. The exact group-scheme map
is
\texttt{\small S\_\allowbreak{}5→\allowbreak{}GL(\allowbreak{}E)\allowbreak{}}
by permutation matrices, followed by
\texttt{\small S\_\allowbreak{}lambda}; the determinant character
records its passage through
\texttt{\small SL(\allowbreak{}E)\allowbreak{}}.

\subsection{G-002. Exact nonnegative determinant certificates for
plethysm and the canonical
map}\label{g-002.-exact-nonnegative-determinant-certificates-for-plethysm-and-the-canonical-map}

Let \texttt{\small N=\allowbreak{}mn}, let
\texttt{\small [\allowbreak{}lambda]\allowbreak{}\_\allowbreak{}Z} be
the integral Specht lattice in its standard polytabloid basis, and let
\texttt{\small rho} be its permutation action. Let
\texttt{\small H=\allowbreak{}S\_\allowbreak{}n wr S\_\allowbreak{}m} be
the row wreath subgroup of \texttt{\small S\_\allowbreak{}N}, and
\texttt{\small K=\allowbreak{}S\_\allowbreak{}m wr S\_\allowbreak{}n}
the column wreath subgroup, with cells indexed by the same
\texttt{\small m by n} grid. Set \[
A_H=\sum_{h\in H}\rho(h),\qquad A_K=\sum_{k\in K}\rho(k),\qquad B=A_KA_H.
\tag{G2.1}
\] These are finite integer matrices. Their ranks over \texttt{\small Q}
are respectively \[
\operatorname{rank} A_H=a_{m,n}^{\lambda},\qquad
\operatorname{rank} B=r_{m,n}^{\lambda},
\tag{G2.2}
\] where \texttt{\small r} is the multiplicity of
\texttt{\small S\_\allowbreak{}lambda} in the image of the normalized
Hadamard--Howe map. For either
\texttt{\small C=\allowbreak{}A\_\allowbreak{}H} or
\texttt{\small C=\allowbreak{}B}, \[
\det(I+tCC^{\mathsf T})=
\sum_{j\ge0}t^j
\sum_{\substack{|I|=|J|=j}}\det(C_{I,J})^2.
\tag{G2.3}
\] Its degree is the corresponding rank. This gives exact finite sums of
squares as certificates for every multiplicity and every map rank.

\textbf{Proof.} Averaging
\texttt{\small P\_\allowbreak{}H=\allowbreak{}A\_\allowbreak{}H/\allowbreak{}|H|}
is an idempotent with image
\texttt{\small [\allowbreak{}lambda]\allowbreak{}\textasciicircum{}H},
because left multiplication permutes the terms of the average, and the
average fixes every invariant vector. Schur--Weyl duality followed by
row symmetrization identifies this invariant space with the multiplicity
space of \texttt{\small S\_\allowbreak{}lambda} in
\texttt{\small Sym\textasciicircum{}m(\allowbreak{}Sym\textasciicircum{}n V)\allowbreak{}}.
The normalized grid map on that multiplicity space is
\texttt{\small P\_\allowbreak{}K|\_\allowbreak{}[\allowbreak{}lambda]\allowbreak{}\textasciicircum{}H};
hence \texttt{\small P\_\allowbreak{}KP\_\allowbreak{}H} has the same
image, and multiplication by the positive nonzero integers
\texttt{\small |K||H|} does not change its rank. This proves (G2.2) with
the specified map, rather than an arbitrary equivariant map having the
same source and target.

The coefficient of \texttt{\small t\textasciicircum{}j} in
\texttt{\small det(\allowbreak{}I+\allowbreak{}tD)\allowbreak{}} is the
sum of all principal \texttt{\small j by j} minors of \texttt{\small D},
by multilinearity of determinant. Apply this to
\texttt{\small D=\allowbreak{}CC\textasciicircum{}T} and then apply
Cauchy--Binet to each principal minor. The result is (G2.3). A real
matrix of rank \texttt{\small s} has a nonzero \texttt{\small s by s}
minor and every larger minor vanishes. The coefficient at
\texttt{\small t\textasciicircum{}s} is consequently positive and all
coefficients above it are zero. This proves the degree assertion without
relying on a choice of invariant inner product for the Specht basis.

The finite sum in (G2.3) includes determinants that must themselves be
evaluated. The formula therefore records the exact remaining
nonzero-minor test; no efficient lattice-point rule or complexity-class
containment is inferred from it. In scheme language the rank conditions
are the determinantal closed subschemes of the affine space of matrices,
cut out by these minors. The displayed integral matrices give actual
integral points of those schemes.

\subsection{G-003. The integral Chow pullback, divided powers, and the
image
algebra}\label{g-003.-the-integral-chow-pullback-divided-powers-and-the-image-algebra}

Let \texttt{\small E=\allowbreak{}Z\textasciicircum{}r}, and let \[
\mu:E^{\times n}\longrightarrow\operatorname{Sym}^n E,
\qquad(v_1,\ldots,v_n)\longmapsto v_1\cdots v_n
\tag{G3.1}
\] be the multiplication morphism of affine schemes over
\texttt{\small Z}. Its degree-\texttt{\small d} coordinate pullback is
\[
\mu_d^*:
\operatorname{Sym}^d\bigl((\operatorname{Sym}^nE)^\vee\bigr)
\longrightarrow
\left((\operatorname{Sym}^d E^\vee)^{\otimes n}\right)^{S_n}
=\Gamma^n(\operatorname{Sym}^dE^\vee).
\tag{G3.2}
\] Here
\texttt{\small (\allowbreak{}Sym\textasciicircum{}n E)\allowbreak{}\textasciicircum{}vee=\allowbreak{}Gamma\textasciicircum{}n(\allowbreak{}E\textasciicircum{}vee)\allowbreak{}}
and the right equality is the divided-power identification for a finite
free module. With the rational norm isomorphisms
\texttt{\small Sym\textasciicircum{}s(\allowbreak{}F)\allowbreak{}→\allowbreak{}(\allowbreak{}F\textasciicircum{}tensor s)\allowbreak{}\textasciicircum{}\{S\_\allowbreak{}s\}},
normalized by \texttt{\small 1/\allowbreak{}s!}, equation (G3.2) becomes
the canonical grid map \[
h_{d,n}:\operatorname{Sym}^d(\operatorname{Sym}^n E_{\mathbb Q}^\vee)
\longrightarrow\operatorname{Sym}^n(\operatorname{Sym}^d E_{\mathbb Q}^\vee).
\tag{G3.3}
\]

\textbf{Proof.} A homogeneous degree-\texttt{\small d} function on
\texttt{\small Sym\textasciicircum{}n E} pulls back to a polynomial of
multidegree
\texttt{\small (\allowbreak{}d,\allowbreak{}...,\allowbreak{}d)\allowbreak{}}
in the \texttt{\small n} input vectors. The multidegree component of the
polynomial ring is
\texttt{\small (\allowbreak{}Sym\textasciicircum{}d E\textasciicircum{}vee)\allowbreak{}\textasciicircum{}tensor n}.
The map (G3.1) is unchanged by permuting input vectors, so its pullback
belongs to its \texttt{\small S\_\allowbreak{}n} invariants. For a free
module, the orbit sums of tensor basis words give both a basis of
symmetric tensors and the defining basis of divided powers; this proves
the integral equality in (G3.2). The dual of coinvariant symmetric
powers is the invariant tensor module, yielding the other divided-power
identification.

Over \texttt{\small Q}, the inverse of the quotient from invariant
tensors to coinvariant symmetric tensors is normalized averaging. Apply
that explicit map in both the \texttt{\small d}-row and
\texttt{\small n}-column tensor arrangements. Coordinate pullback
multiplies entries along the transposed arrangement; the resulting
invariant tensor is precisely column averaging of the row-invariant
tensor. This is the grid definition of (G3.3), with no unrecorded
integral identification of symmetric and divided powers.

Set
\texttt{\small A\_\allowbreak{}Z=\allowbreak{}im(\allowbreak{}mu\textasciicircum{}*)\allowbreak{}}
inside the polynomial ring
\texttt{\small Z[\allowbreak{}E\textasciicircum{}times n]\allowbreak{}},
with the grading just specified. Its kernel is an ideal in the
polynomial coordinate ring of \texttt{\small Sym\textasciicircum{}n E},
and \[
\mathcal C_{n,r}=\operatorname{Spec}A_{\mathbb Z}
\] is the scheme-theoretic image of (G3.1). This is EGA I, Proposition
9.5.1 and Corollary 9.5.2, printed p.176, with Definition 9.5.3 on
p.177. The image is a subring of a torsion-free abelian group, hence is
torsion-free, hence is flat over \texttt{\small Z}: for a principal
ideal domain, a module is flat precisely when it is torsion-free. One
proof of the latter statement writes a torsion-free module as the
filtered union of its finitely generated free submodules and uses
preservation of flatness under filtered colimits. Each graded part is
finite free as a subgroup of a finite free multidegree component. This
proves flatness of this specific integral model. It does not identify
every special fiber with the scheme-theoretic image of the corresponding
fiber map; the precise specialization map is
\texttt{\small A\_\allowbreak{}Z tensor k→\allowbreak{}im(\allowbreak{}mu\_\allowbreak{}k\textasciicircum{}*)\allowbreak{}},
and injectivity of that map remains a separate saturation question.

Over \texttt{\small Q}, the degree-\texttt{\small d} Schur
multiplicities of \texttt{\small A\_\allowbreak{}Q} are the ranks
\texttt{\small r\_\allowbreak{}(\allowbreak{}d,\allowbreak{}n)\allowbreak{}\textasciicircum{}lambda}
of G-002, with the usual dual convention for coordinate functions. The
rank equation thus has an explicit morphism to the orbit-closure
coordinate ring. The generic span of \texttt{\small n} independent
linear factors identifies the dense product orbit with that of
\texttt{\small x\_\allowbreak{}1...x\_\allowbreak{}n} when
\texttt{\small r>=\allowbreak{}n}, so this is the affine Chow cone in
that range. The construction of (G3.1) is valid for every
\texttt{\small r}.

The projective Hilbert and Quot functors with fixed Hilbert polynomial
are represented by the schemes of FGA, Bourbaki exposé 221, Theorems 3.1
and 3.2, pp.259--260; the Hilbert specialization is after Proposition
3.8, p.265. They parametrize flat families of projective closed
subschemes and coherent quotient sheaves with that Hilbert polynomial.
The finite matrices in (G3.2) remain the explicitly identified maps
whose minors determine the graded ranks; the representing property alone
supplies no new nonzero minor.

\subsection{G-004. Purity and the invariant projector on a finite étale
example}\label{g-004.-purity-and-the-invariant-projector-on-a-finite-uxe9tale-example}

Let
\texttt{\small X=\allowbreak{}Spec F\_\allowbreak{}(\allowbreak{}q\textasciicircum{}2)\allowbreak{}}
over \texttt{\small F\_\allowbreak{}q}, let \texttt{\small ell} be prime
to \texttt{\small q}, and let \texttt{\small sigma} generate its
two-element Galois group. The degree-zero étale cohomology of
\texttt{\small X} after base change to an algebraic closure is
\texttt{\small Q\_\allowbreak{}ell e\_\allowbreak{}0 directsum Q\_\allowbreak{}ell e\_\allowbreak{}1};
Frobenius interchanges the two basis vectors. The two idempotents \[
e_+=\tfrac12(1+\sigma),\qquad e_-=\tfrac12(1-\sigma)
\tag{G4.1}
\] have one-dimensional images with Frobenius eigenvalues
\texttt{\small 1} and \texttt{\small -\allowbreak{}1} respectively. Both
are pure of weight zero, since both eigenvalues have complex absolute
value one, but their Galois-invariant dimensions are respectively one
and zero.

\textbf{Proof.} The étale cohomology of a two-point geometric scheme in
degree zero is its functions, and higher cohomology vanishes. The
Frobenius action permutes the two embeddings of the degree-two finite
field extension, so its eigenvectors are
\texttt{\small e\_\allowbreak{}0+\allowbreak{}e\_\allowbreak{}1} and
\texttt{\small e\_\allowbreak{}0-\allowbreak{}e\_\allowbreak{}1}.
Formula (G4.1), the two eigenvalues, and the invariant dimensions follow
by direct matrix multiplication. This works also for
\texttt{\small ell=\allowbreak{}2}, since the coefficient field is
\texttt{\small Q\_\allowbreak{}2} and \texttt{\small 1/\allowbreak{}2}
exists. In the category of rational Artin motives the same idempotents
are the two projectors on this finite étale motive.

This supplies an exact test of the proposed passage from a pure
cohomological object to a symmetric-group-invariant multiplicity: the
group projector must be carried along the realization. Purity alone
assigns the same weight to the two displayed lines while the invariant
projector has ranks one and zero. The example does not rule out a
geometric construction that includes that projector. The motivic
interpretation applies to this explicit finite étale object, without
assuming a general category of motives has been specified for archive
cote12/scan041.

\subsection{G-005. A positive polytope formulation for every two-row
output}\label{g-005.-a-positive-polytope-formulation-for-every-two-row-output}

This section carries the known marked-tree theorem of
Pak--Panova--Swanson, \emph{A combinatorial interpretation for certain
plethysm and Kronecker coefficients} (4 November 2025), Theorems 1.1/4.5
and 1.2/5.2, into explicit polytope language. The known combinatorial
interpretation is not claimed as a new theorem of this run.

For a nonempty outer partition \texttt{\small mu}, inner row
\texttt{\small (\allowbreak{}k)\allowbreak{}} with
\texttt{\small k>=\allowbreak{}1}, and output
\texttt{\small (\allowbreak{}N-\allowbreak{}r,\allowbreak{}r)\allowbreak{}}
where \texttt{\small N=\allowbreak{}k|mu|} and
\texttt{\small 0<=\allowbreak{}r<=\allowbreak{}N/\allowbreak{}2}, take
the finite set of unmarked GOH trees
\texttt{\small G(\allowbreak{}mu,\allowbreak{}k)\allowbreak{}} defined
in that work, Definitions 3.5 and 3.8. The empty outer partition has the
sole coefficient one at the empty output. The definition, including all
integer partition labels, is reproduced below so that the formula is
executable from the paper. For a tree \texttt{\small T} order its leaves
in the prescribed lexicographic depth-first order, and write their
nonnegative labels as
\texttt{\small a\_\allowbreak{}1,\allowbreak{}...,\allowbreak{}a\_\allowbreak{}t}.
Put \[
s(T)=\frac{N-\sum_i a_i}{2}.
\] Define the bounded rational polytope
\texttt{\small P\_\allowbreak{}T(\allowbreak{}r)\allowbreak{}} in
coordinates
\texttt{\small z\_\allowbreak{}1,\allowbreak{}...,\allowbreak{}z\_\allowbreak{}t}
by \[
z_1=0,\qquad z_t=r-s(T),\qquad z_i\le z_{i+1},
\tag{G5.1}
\] \[
z_{i+1}-z_i\le a_{i+1},\qquad
z_{i+1}+z_i\le a_1+\cdots+a_i
\quad(1\le i<t).
\tag{G5.2}
\] Then \[
p_{\mu,(k)}^{(N-r,r)}
=\sum_{T\in G(\mu,k)}\#(P_T(r)\cap\mathbb Z^t).
\tag{G5.3}
\] For the symmetric outer index
\texttt{\small (\allowbreak{}m)\allowbreak{}}, one may instead use the
KOH trees \texttt{\small T(\allowbreak{}n,\allowbreak{}m)\allowbreak{}}
for \texttt{\small N=\allowbreak{}mn}, and the same formula also equals
\[
g((mn-r,r),(n^m),(n^m)).
\tag{G5.4}
\]

\textbf{Proof of the asserted form.} The marks of Pak--Panova--Swanson,
Definitions 4.4 and 5.1, satisfy
\texttt{\small 0=\allowbreak{}z\_\allowbreak{}1<=\allowbreak{}...<=\allowbreak{}z\_\allowbreak{}t=\allowbreak{}r-\allowbreak{}s(\allowbreak{}T)\allowbreak{}}
and
\texttt{\small z\_\allowbreak{}(\allowbreak{}i+\allowbreak{}1)\allowbreak{}-\allowbreak{}z\_\allowbreak{}i<=\allowbreak{}min(\allowbreak{}sum\_\allowbreak{}(\allowbreak{}j<=\allowbreak{}i)\allowbreak{}a\_\allowbreak{}j-\allowbreak{}2z\_\allowbreak{}i,\allowbreak{}a\_\allowbreak{}(\allowbreak{}i+\allowbreak{}1)\allowbreak{})\allowbreak{}}.
Separating this minimum into its two linear inequalities gives exactly
(G5.2). Thus the identity map on mark vectors is a bijection between
their marked trees with fixed underlying tree and the integer points of
\texttt{\small P\_\allowbreak{}T(\allowbreak{}r)\allowbreak{}}. The
endpoint and monotonicity bound every coordinate between zero and
\texttt{\small r-\allowbreak{}s(\allowbreak{}T)\allowbreak{}}; a
negative endpoint gives the empty polytope. Summing the cited theorem
over underlying trees proves (G5.3) and (G5.4). The proof of the change
of language is complete and uses no subtraction of lattice-point counts.

For a general inner partition \texttt{\small nu}, evaluation in rank two
gives the following exact transfer. Partitions with three or more rows
have
\texttt{\small S\_\allowbreak{}nu(\allowbreak{}C\textasciicircum{}2)\allowbreak{}=\allowbreak{}0}.
For
\texttt{\small nu=\allowbreak{}(\allowbreak{}b+\allowbreak{}k,\allowbreak{}b)\allowbreak{}},
the isomorphism \[
S_{(b+k,b)}(\mathbb C^2)=\det^b\otimes\operatorname{Sym}^k\mathbb C^2
\] and homogeneity of \texttt{\small S\_\allowbreak{}mu} give the
determinant shift by \texttt{\small b|mu|} in each output row. This is
precisely Pak--Panova--Swanson, Lemma 2.1. The shifted output with
negative second row has coefficient zero. For
\texttt{\small k=\allowbreak{}0}, the inner representation is a line,
and its outer Schur functor is zero for
\texttt{\small ell(\allowbreak{}mu)\allowbreak{}>1} and is
\texttt{\small det\textasciicircum{}(\allowbreak{}b|mu|)\allowbreak{}}
for \texttt{\small mu=\allowbreak{}(\allowbreak{}|mu|)\allowbreak{}}.
This covers the rank-one edge case without invoking a zero-parameter
tree convention.

The corresponding vector-bundle statement is obtained from the rank-two
frame torsor and the above determinant-line isomorphism. Every polytope
computes the coefficient of the associated Schur functor in
\texttt{\small K\_\allowbreak{}0(\allowbreak{}BGL\_\allowbreak{}2)\allowbreak{}}
over a characteristic-zero field. The \texttt{\small SL\_\allowbreak{}2}
Clebsch--Gordan branching chain behind the marks is an exact chain of
interlacing inequalities; it supplies a successful instance of the
branching approach outside the stable-tail range.

\subsubsection{Explicit tree indexing}\label{explicit-tree-indexing}

For a partition \texttt{\small alpha}, write
\texttt{\small m\_\allowbreak{}j(\allowbreak{}alpha)\allowbreak{}} for
the multiplicity of its part \texttt{\small j}, \texttt{\small alpha'}
for its conjugate, and
\texttt{\small Q\_\allowbreak{}j(\allowbreak{}alpha)\allowbreak{}=\allowbreak{}sum\_\allowbreak{}(\allowbreak{}x<=\allowbreak{}j)\allowbreak{}alpha'\_\allowbreak{}x}.
A KOH vertex is labeled
\texttt{\small (\allowbreak{}alpha,\allowbreak{}a,\allowbreak{}b)\allowbreak{}}
with \texttt{\small a>=\allowbreak{}0},
\texttt{\small b>=\allowbreak{}1}, \texttt{\small alpha partition of b}.
A vertex with \texttt{\small b=\allowbreak{}1} is a leaf with scalar
label \texttt{\small a}. A vertex with \texttt{\small b>=\allowbreak{}2}
has one child for each distinct part \texttt{\small j} of
\texttt{\small alpha}, in increasing order, with parameters \[
a^{(j)}=(a+2)j-2Q_j(\alpha),\qquad
b^{(j)}=m_j(\alpha),
\] and its partition label ranges over the partitions of
\texttt{\small b\textasciicircum{}(\allowbreak{}j)\allowbreak{}}. A
negative
\texttt{\small a\textasciicircum{}(\allowbreak{}j)\allowbreak{}} gives
no tree. The root parameters of
\texttt{\small T(\allowbreak{}n,\allowbreak{}m)\allowbreak{}} are
\texttt{\small (\allowbreak{}a,\allowbreak{}b)\allowbreak{}=\allowbreak{}(\allowbreak{}n,\allowbreak{}m)\allowbreak{}}.

A GOH root of type
\texttt{\small (\allowbreak{}mu,\allowbreak{}k)\allowbreak{}}, for
\texttt{\small ell=\allowbreak{}ell(\allowbreak{}mu)\allowbreak{}} and
\texttt{\small n=\allowbreak{}|mu|}, carries a sequence of partitions
\texttt{\small nu\textasciicircum{}(\allowbreak{}0)\allowbreak{},\allowbreak{}...,\allowbreak{}nu\textasciicircum{}(\allowbreak{}ell)\allowbreak{}}
with \[
\nu^{(0)}=(1^n),\qquad
|\nu^{(i)}|=\sum_{j\ge i+1}\mu_j,
\] \[
P_j^i=Q_j(\nu^{(i+1)})-2Q_j(\nu^{(i)})+Q_j(\nu^{(i-1)})\ge0
\quad(1\le i<\ell,\ 1\le j\le n).
\] Let
\texttt{\small m=\allowbreak{}(\allowbreak{}nu\textasciicircum{}(\allowbreak{}1)\allowbreak{})\allowbreak{}'\_\allowbreak{}1},
with \texttt{\small m=\allowbreak{}0} for an empty partition, and retain
configurations with \texttt{\small m<=\allowbreak{}k}. For every pair
\texttt{\small (\allowbreak{}i,\allowbreak{}j)\allowbreak{}} with
\texttt{\small m\_\allowbreak{}j(\allowbreak{}nu\textasciicircum{}(\allowbreak{}i)\allowbreak{})\allowbreak{}>0},
attach a KOH tree of type
\texttt{\small (\allowbreak{}P\_\allowbreak{}j\textasciicircum{}i,\allowbreak{}m\_\allowbreak{}j(\allowbreak{}nu\textasciicircum{}(\allowbreak{}i)\allowbreak{})\allowbreak{})\allowbreak{}},
ordered lexicographically. For \texttt{\small m<k}, attach also a final
KOH tree of type
\texttt{\small (\allowbreak{}n,\allowbreak{}k-\allowbreak{}m)\allowbreak{}}.
This defines the complete finite index set in (G5.3). Finiteness and the
equal-center property \texttt{\small s(\allowbreak{}T)\allowbreak{}} are
the cited Proposition 3.4 and Proposition 3.9; their proofs are repeated
application of the KOH and GOH identities, Theorems 3.1 and 3.6.

The printed induction proof of Lemma 4.2 has a local indexing slip: its
last inequality writes
\texttt{\small |a|-\allowbreak{}2k\_\allowbreak{}t}, whereas
substitution into Lemma 4.1 gives
\texttt{\small |a|-\allowbreak{}2k\_\allowbreak{}(\allowbreak{}t-\allowbreak{}1)\allowbreak{}}.
The statement of Lemma 4.2 and both marking definitions have the correct
latter index. The polytope above uses those correct statements.

A second local correction concerns Lemma 4.1, printed p.8. Put
\(L=\min(r-2i,s-2j)\). Direct multiplication of finite geometric sums
gives \[
(1-q)(q^i+\cdots+q^{r-i})(q^j+\cdots+q^{s-j})
=q^{i+j}(1+\cdots+q^L)-q^{r+s-i-j+1}(1+\cdots+q^{-L}).
\] The printed negative leading exponent omits the final one. For
\(r=s=1,i=j=0\), the left side is \(1+q-q^2-q^3\), while the printed
expression is \(1-q^2\). The smallest exponent in the corrected negative
term is \(i+j+\max(r-2i,s-2j)+1>(r+s)/2\). Thus its actual coefficient
assertion below the midpoint follows, and the corrected Lemma 4.2
induction and the tree formulas retain their stated conclusions. The
independent reviewer inspected the PDF page visually and checked this
entire repair.

\section{Outer degree three and the canonical two-row
blocks}\label{outer-degree-three-and-the-canonical-two-row-blocks}

\subsection{G-006. Canonical-map isomorphisms on every two-row
block}\label{g-006.-canonical-map-isomorphisms-on-every-two-row-block}

For every \texttt{\small m,\allowbreak{}n>=\allowbreak{}1}, the
canonical map
\texttt{\small h\_\allowbreak{}(\allowbreak{}m,\allowbreak{}n)\allowbreak{}}
is an isomorphism on each isotypic multiplicity space indexed by a
partition with at most two rows. In particular \[
r_{m,n}^{\lambda}=a_{m,n}^{\lambda}=a_{n,m}^{\lambda}
\quad(\ell(\lambda)\le2).
\] This is the natural-transformation form of the classical canonical
Hermite reciprocity theorem, stated in the supplied Landsberg survey,
Theorem 7.21, §7.7. Its precise transfer to arbitrary ambient dimension
is as follows.

\textbf{Proof of the transfer.} A natural transformation between
homogeneous characteristic-zero polynomial functors is the direct sum,
over partitions \texttt{\small lambda}, of the identity on
\texttt{\small S\_\allowbreak{}lambda} tensored with a fixed linear map
between its multiplicity spaces. Evaluation on
\texttt{\small C\textasciicircum{}2} kills exactly the summands of
length greater than two. Canonical Hermite reciprocity says the
evaluated transformation is an isomorphism. Consequently every retained
multiplicity map is invertible, and those same maps occur after
evaluation in every larger dimension. This proves the displayed rank
equality, including the specified canonical map.

The geometric realization is the multiplication morphism of G-003 in
rank two: every binary form factors into linear forms over
\texttt{\small C}, so the Chow image is the whole affine space of binary
forms. Its coordinate pullback is injective. Both sides have dimension
\texttt{\small binom(\allowbreak{}m+\allowbreak{}n,\allowbreak{}m)\allowbreak{}},
giving the canonical isomorphism. This is the exact morphism relating
the factorization statement, the coordinate ring, and the canonical map;
the source theorem is known.

\subsection{G-007. Explicit one-dimensional lattice counts for every
outer Schur index of size
three}\label{g-007.-explicit-one-dimensional-lattice-counts-for-every-outer-schur-index-of-size-three}

The coefficients in this family are classical (Thrall and Chen). Their
explicit current primary presentation is Maas--Gariépy and Tétreault,
\emph{3-Plethysms of homogeneous and elementary symmetric functions},
arXiv:2209.14942v1, §4.2, pp.8--9; their tableau realization is Theorem
1, §7.3, p.20. Here the formula is put in the literal permutation-matrix
realization of T-002, retaining the two finite-character corrections.

Let
\texttt{\small lambda=\allowbreak{}(\allowbreak{}lambda\_\allowbreak{}1,\allowbreak{}lambda\_\allowbreak{}2,\allowbreak{}lambda\_\allowbreak{}3)\allowbreak{}}
be a partition of \texttt{\small 3n}, padded by zeros. Set \[
L=\max(\lambda_3,2n-\lambda_1),\qquad
U=\min(\lambda_2,2n-\lambda_2,n),\qquad K=U-L+1.
\] Define \texttt{\small d\_\allowbreak{}2=\allowbreak{}0} for even
\texttt{\small K} and
\texttt{\small d\_\allowbreak{}2=\allowbreak{}(\allowbreak{}-\allowbreak{}1)\allowbreak{}\textasciicircum{}L}
for odd \texttt{\small K}. Let \[
(b_1,b_2,b_3)=(\lambda_1+2,\lambda_2+1,\lambda_3)\pmod3.
\] Set \texttt{\small d\_\allowbreak{}3=\allowbreak{}0} for repeated
residues; for three distinct residues set
\texttt{\small d\_\allowbreak{}3} equal to the sign of the permutation
carrying the ordered residue list
\texttt{\small (\allowbreak{}2,\allowbreak{}1,\allowbreak{}0)\allowbreak{}}
to
\texttt{\small (\allowbreak{}b\_\allowbreak{}1,\allowbreak{}b\_\allowbreak{}2,\allowbreak{}b\_\allowbreak{}3)\allowbreak{}}.
Then \[
p_{(3),(n)}^{\lambda}=\frac{K+3d_2+2d_3}{6},\qquad
p_{(1,1,1),(n)}^{\lambda}=\frac{K-3d_2+2d_3}{6},\qquad
p_{(2,1),(n)}^{\lambda}=\frac{K-d_3}{3}.
\tag{G7.1}
\] Each quotient is the following positive lattice-point formula: for
its numerator \texttt{\small B} and denominator \texttt{\small D}, \[
\boxed{\#\{z\in\mathbb Z:1\le z,\ Dz\le B\}.}
\tag{G7.2}
\] The residue tests are finite case selection, and within each selected
chamber the displayed bounds are affine linear in
\texttt{\small lambda,\allowbreak{}n}. The intervals are allowed to be
empty. The statement imposes no homogeneous-dilation requirement.

\textbf{Proof of the normalizer dictionary.} The balanced space
\texttt{\small M=\allowbreak{}(\allowbreak{}S\_\allowbreak{}lambda C\textasciicircum{}3)\allowbreak{}\_\allowbreak{}(\allowbreak{}n,\allowbreak{}n,\allowbreak{}n)\allowbreak{}}
is a representation of \texttt{\small S\_\allowbreak{}3}. Restricting
first to \texttt{\small GL\_\allowbreak{}2 times GL\_\allowbreak{}1},
the interlacing highest weights contributing balanced degree are
\texttt{\small (\allowbreak{}2n-\allowbreak{}k,\allowbreak{}k)\allowbreak{}}
for exactly \texttt{\small L<=\allowbreak{}k<=\allowbreak{}U}. Each
contributes a single
\texttt{\small (\allowbreak{}n,\allowbreak{}n)\allowbreak{}} weight
line. The transposition of the first two basis vectors acts on this line
by
\texttt{\small (\allowbreak{}-\allowbreak{}1)\allowbreak{}\textasciicircum{}k},
as the explicit determinant calculation in T-006 shows. Thus its trace
is
\texttt{\small sum\_\allowbreak{}(\allowbreak{}k=\allowbreak{}L)\allowbreak{}\textasciicircum{}U(\allowbreak{}-\allowbreak{}1)\allowbreak{}\textasciicircum{}k=\allowbreak{}d\_\allowbreak{}2}
and the identity trace is \texttt{\small K}.

A three-cycle permutes the torus weight spaces of
\texttt{\small S\_\allowbreak{}lambda C\textasciicircum{}3}. Only a
weight constant on all three coordinates contributes to its trace; the
total degree \texttt{\small 3n} makes that weight
\texttt{\small (\allowbreak{}n,\allowbreak{}n,\allowbreak{}n)\allowbreak{}}.
Its trace on \texttt{\small M} therefore equals its trace on the whole
representation. The cycle has eigenvalues
\texttt{\small 1,\allowbreak{}zeta,\allowbreak{}zeta\textasciicircum{}2}
for a primitive cube root \texttt{\small zeta}. The Weyl alternant
formula evaluates its character as \[
\frac{\det(x_j^{\lambda_i+3-i})_{1\le i,j\le3}}
{\det(x_j^{3-i})_{1\le i,j\le3}},
\qquad(x_1,x_2,x_3)=(1,\zeta,\zeta^2).
\] The denominator is nonzero. Repeated exponent residues give two equal
rows and value zero; distinct residues give the stated row-permutation
sign. Thus the three-cycle trace is \texttt{\small d\_\allowbreak{}3}.

The three irreducible characters of \texttt{\small S\_\allowbreak{}3}
have values
\texttt{\small (\allowbreak{}1,\allowbreak{}1,\allowbreak{}1)\allowbreak{}},
\texttt{\small (\allowbreak{}1,\allowbreak{}-\allowbreak{}1,\allowbreak{}1)\allowbreak{}},
and
\texttt{\small (\allowbreak{}2,\allowbreak{}0,\allowbreak{}-\allowbreak{}1)\allowbreak{}}
on identity, transpositions and three-cycles, whose class sizes are
\texttt{\small 1,\allowbreak{}3,\allowbreak{}2}. Taking these character
inner products proves (G7.1), via T-003. Each is consequently a
nonnegative integer. For \texttt{\small B=\allowbreak{}Dc} with
\texttt{\small c} that nonnegative multiplicity, the integer solutions
of (G7.2) are exactly \texttt{\small 1,\allowbreak{}...,\allowbreak{}c};
the interval endpoints were computed by the explicit bounded trace
calculation, rather than by a general coefficient oracle. This proves
(G7.2). The interlacing interval is nonempty for every partition of
\texttt{\small 3n} of length at most three; this also follows from the
inequalities
\texttt{\small lambda\_\allowbreak{}1>=\allowbreak{}n>=\allowbreak{}lambda\_\allowbreak{}3}
and
\texttt{\small lambda\_\allowbreak{}1+\allowbreak{}lambda\_\allowbreak{}2+\allowbreak{}lambda\_\allowbreak{}3=\allowbreak{}3n}
in the definitions of \texttt{\small L,\allowbreak{}U}.

The geometric statement is the equality in the characteristic-zero
representation ring
\texttt{\small K\_\allowbreak{}0(\allowbreak{}BGL\_\allowbreak{}r)\allowbreak{}}
and its pullback along the frame torsor of a locally free
rank-\texttt{\small r} sheaf. The two traces are taken along the
explicit constant subgroup scheme
\texttt{\small S\_\allowbreak{}3→\allowbreak{}GL\_\allowbreak{}3} of
permutation matrices, and the determinant factor in the transposition
action is retained. Hence this evaluates, rather than omits, the
normalizer projector for the entire outer-degree-three family.

\section{Positive lattice-point formulas from explicit invariant
presentations}\label{positive-lattice-point-formulas-from-explicit-invariant-presentations}

\subsection{G-008. A repeated-odd-part nonvanishing family in unbounded
width}\label{g-008.-a-repeated-odd-part-nonvanishing-family-in-unbounded-width}

Put
\texttt{\small S=\allowbreak{}<4,\allowbreak{}6,\allowbreak{}7>=\allowbreak{}\{0,\allowbreak{}4,\allowbreak{}6,\allowbreak{}7,\allowbreak{}8\} union \{k:\allowbreak{}k>=\allowbreak{}10\}}.
For every \texttt{\small k in S} with \texttt{\small k>=\allowbreak{}4}
and every \texttt{\small j>=\allowbreak{}0}, \[
\boxed{c((k^{4j}))\ge1.}
\tag{G8.1}
\] The empty partition has coefficient one. More generally, let
\texttt{\small lambda} be a partition with no part one. Its even parts
have arbitrary multiplicities. For each odd part \texttt{\small i in S},
its multiplicity has the form
\texttt{\small 4q\_\allowbreak{}i+\allowbreak{}epsilon\_\allowbreak{}i},
with \texttt{\small q\_\allowbreak{}i>=\allowbreak{}0} and
\texttt{\small epsilon\_\allowbreak{}i in \{0,\allowbreak{}1\}}; for the
remaining odd parts \texttt{\small 3,\allowbreak{}5,\allowbreak{}9}, its
multiplicity is at most one. Then \[
\boxed{c(\lambda)\ge1.}
\tag{G8.2}
\] The finite plethysm coefficients in T-010's stability range have the
same positive value. The theorem concerns a specified sufficient family;
it does not assert an exhaustive support classification.

\textbf{Proof, including the exact finite seeds.} Let
\texttt{\small W=\allowbreak{}C\textasciicircum{}4}. The invariant
algebra
\texttt{\small Sym(\allowbreak{}Sym\textasciicircum{}4 W)\allowbreak{}\textasciicircum{}(\allowbreak{}SL\_\allowbreak{}4)\allowbreak{}}
is a subalgebra of a polynomial algebra and therefore a domain. The
exact seed multiplicities in degrees
\texttt{\small 4,\allowbreak{}6,\allowbreak{}7} are each one. The finite
certificates are supplied in
\texttt{\small QUATERNARY\_\allowbreak{}QUARTIC\_\allowbreak{}INVARIANTS.tsv}
and independently in
\texttt{\small QUARTIC\_\allowbreak{}SEEDS\_\allowbreak{}INDEPENDENT\_\allowbreak{}CHECK.tsv},
with the three complete class-sum files. Here is the precise
mathematical verification of those finite numbers. The torus character
is the coefficient of \texttt{\small t\textasciicircum{}k} in \[
\prod_{\substack{\alpha\in\mathbb Z_{\ge0}^4\\|\alpha|=4}}
(1-tx^\alpha)^{-1}.
\tag{G8.3}
\] There are 35 factors. With
\texttt{\small rho=\allowbreak{}(\allowbreak{}3,\allowbreak{}2,\allowbreak{}1,\allowbreak{}0)\allowbreak{}}
and \texttt{\small M\_\allowbreak{}k(\allowbreak{}w)\allowbreak{}} its
integer weight coefficients, Weyl's alternant identity gives the
invariant multiplicity \[
\sum_{\sigma\in S_4}\operatorname{sgn}(\sigma)
M_k((k,k,k,k)+\rho-\sigma\rho).
\tag{G8.4}
\] The supplied program expands the 35 geometric series by integer
dynamic programming, computes all 24 terms of (G8.4), and obtains one in
each of the three asserted degrees. Only coordinates at most
\texttt{\small k+\allowbreak{}3} can contribute; all factor weights are
nonnegative, so truncation beyond that box preserves all required
coefficients. The independent program instead expands the wreath-product
cycle index and sums symmetric-group characters by Murnaghan--Nakayama;
it obtains the same three values. This is an exact finite computation
certificate, not a proof-assistant claim.

Choose nonzero invariant seed vectors. Their products are nonzero, so
every degree in \texttt{\small S} contains an invariant. The displayed
description of \texttt{\small S} follows because
\texttt{\small 10=\allowbreak{}4+\allowbreak{}6},
\texttt{\small 11=\allowbreak{}4+\allowbreak{}7},
\texttt{\small 12=\allowbreak{}6+\allowbreak{}6},
\texttt{\small 13=\allowbreak{}6+\allowbreak{}7}, and adding four covers
every subsequent integer; direct checking settles
\texttt{\small 0,\allowbreak{}...,\allowbreak{}9}.

Under the symmetric-function involution \texttt{\small omega}, the even
inner degree four gives \[
\omega(h_k[h_4])=h_k[e_4].
\tag{G8.5}
\] Thus the coefficient of
\texttt{\small (\allowbreak{}k\textasciicircum{}4)\allowbreak{}} in
\texttt{\small Sym\textasciicircum{}k(\allowbreak{}Sym\textasciicircum{}4 W)\allowbreak{}}
equals the coefficient of
\texttt{\small (\allowbreak{}4\textasciicircum{}k)\allowbreak{}} in
\texttt{\small Sym\textasciicircum{}k(\allowbreak{}Lambda\textasciicircum{}4 C\textasciicircum{}k)\allowbreak{}}.
This is the determinant representation
\texttt{\small det\textasciicircum{}4}, so a nonzero highest vector of
weight \texttt{\small 4 omega\_\allowbreak{}k} exists in that symmetric
algebra. Its \texttt{\small j}th power is nonzero and has highest weight
\texttt{\small 4j omega\_\allowbreak{}k}. This is the constituent for
tail
\texttt{\small (\allowbreak{}k\textasciicircum{}(\allowbreak{}4j)\allowbreak{})\allowbreak{}}
in the exact module M\_k of T-010, proving (G8.1).

For (G8.2), work in rank
\texttt{\small r=\allowbreak{}lambda\_\allowbreak{}1}. Embed every seed
highest vector from rank \texttt{\small i} using the first
\texttt{\small i} coordinates into
\texttt{\small Sym\textasciicircum{}i(\allowbreak{}Lambda\textasciicircum{}4 C\textasciicircum{}r)\allowbreak{}}.
It remains highest for \texttt{\small GL\_\allowbreak{}r}: positive root
operators involving a basis coordinate beyond \texttt{\small i}
annihilate it, and the operators within the first \texttt{\small i}
coordinates already do so. Multiply the prescribed powers of all these
vectors inside the single polynomial algebra
\texttt{\small Sym(\allowbreak{}Lambda\textasciicircum{}4 C\textasciicircum{}r)\allowbreak{}}.
This product is nonzero. The additional pure highest vectors for the
arbitrary even parts live in their corresponding symmetric-algebra
factors; the part four contributes another power in the same domain. The
residual odd parts with multiplicities
\texttt{\small epsilon\_\allowbreak{}i=\allowbreak{}0,\allowbreak{}1}
select scalars or highest wedge vectors in the corresponding
exterior-algebra factors. Their tensor product is nonzero, highest, and
has total highest weight \texttt{\small lambda'}. T-010 proves (G8.2).
This proof records the factor in which multiplication occurs, so that no
unsupported domain assertion about the full tensor product of exterior
algebras is used.

\subsection{G-009. A terminating positive-polytope procedure for all
symmetric
plethysm}\label{g-009.-a-terminating-positive-polytope-procedure-for-all-symmetric-plethysm}

For every positive integer \texttt{\small m}, the following finite
procedure computes an integer matrix \texttt{\small D\_\allowbreak{}m}
with columns
\texttt{\small d\_\allowbreak{}j=\allowbreak{}(\allowbreak{}n\_\allowbreak{}j,\allowbreak{}lambda\textasciicircum{}(\allowbreak{}j)\allowbreak{})\allowbreak{}},
all \texttt{\small n\_\allowbreak{}j>=\allowbreak{}1}, and a finite
collection \texttt{\small C\_\allowbreak{}m} of pairs
\texttt{\small (\allowbreak{}u,\allowbreak{}S)\allowbreak{}} with
\texttt{\small u in Z\_\allowbreak{}\{≥\allowbreak{}0\}\textasciicircum{}q},
\texttt{\small S subset \{1,\allowbreak{}...,\allowbreak{}q\}}. The
output depends only on \texttt{\small m}. For all
\texttt{\small n>=\allowbreak{}0} and partitions \texttt{\small lambda}
of \texttt{\small mn}, of length at most \texttt{\small m}, it gives \[
\boxed{
a_{m,n}^{\lambda}=
\sum_{(u,S)\in C_m}
\#\left\{z\in\mathbb Z_{\ge0}^{S}:
D_{m,S}z=(n,\lambda)-D_m u\right\}.
}
\tag{G9.1}
\] Each summand is the integer-point count of a bounded rational
polytope. The construction uses flag minors, a finitely generated monoid
of balanced exponents, Reynolds averages of bounded ordinary degree, and
polynomial elimination. It makes no assumption about any unknown
plethysm coefficient, any canonical basis, or the sign of any character
sum. No polynomial bound in the varying parameter \texttt{\small m} for
its output size or running time is asserted.

\subsubsection{The finite construction}\label{the-finite-construction}

Let
\texttt{\small X=\allowbreak{}(\allowbreak{}x\_\allowbreak{}ij)\allowbreak{}}
be an \texttt{\small m by m} matrix of indeterminates over
\texttt{\small Q}. For every nonempty subset
\texttt{\small I subset \{1,\allowbreak{}...,\allowbreak{}m\}}, let
\texttt{\small p\_\allowbreak{}I(\allowbreak{}X)\allowbreak{}} be the
minor with row set \texttt{\small I} and first \texttt{\small |I|}
columns. Introduce a formal variable \texttt{\small z\_\allowbreak{}I}
for each such subset. Let \[
\mathcal H_m=\left\{b\in\mathbb Z_{\ge0}^{2^m-1}:
\sum_{I\ni1}b_I=\cdots=\sum_{I\ni m}b_I\right\}.
\tag{G9.2}
\] Compute its finite Hilbert basis
\texttt{\small b\textasciicircum{}(\allowbreak{}1)\allowbreak{},\allowbreak{}...,\allowbreak{}b\textasciicircum{}(\allowbreak{}s)\allowbreak{}},
and put
\texttt{\small q\_\allowbreak{}l=\allowbreak{}product\_\allowbreak{}I p\_\allowbreak{}I\textasciicircum{}\{b\_\allowbreak{}I\textasciicircum{}(\allowbreak{}l)\allowbreak{}\}}.
The common sum in (G9.2) is its positive integer
\texttt{\small n\_\allowbreak{}l}; its right-torus highest weight is \[
\lambda^{(l)}=\sum_I b_I^{(l)}\omega_{|I|},
\quad\omega_j=(1^j,0^{m-j}).
\tag{G9.3}
\]

Row permutation matrices send each formal
\texttt{\small z\_\allowbreak{}I} to
\texttt{\small epsilon(\allowbreak{}sigma,\allowbreak{}I)\allowbreak{} z\_\allowbreak{}(\allowbreak{}sigma I)\allowbreak{}},
where \texttt{\small epsilon} is the sign needed to sort the permuted
row indices. They consequently permute Hilbert-basis monomials up to
signs. Let
\texttt{\small Y\_\allowbreak{}1,\allowbreak{}...,\allowbreak{}Y\_\allowbreak{}s}
carry this signed permutation action and map
\texttt{\small Y\_\allowbreak{}l} to \texttt{\small q\_\allowbreak{}l}.
Form every Reynolds average \[
\mathcal R(Y^v)=\frac1{m!}\sum_{\sigma\in S_m}\sigma(Y^v),
\qquad1\le|v|\le m!.
\tag{G9.4}
\] Substitute
\texttt{\small Y\_\allowbreak{}l=\allowbreak{}q\_\allowbreak{}l(\allowbreak{}X)\allowbreak{}}
and discard zero polynomials. Denote the remaining explicit polynomials
by
\texttt{\small f\_\allowbreak{}1,\allowbreak{}...,\allowbreak{}f\_\allowbreak{}q};
duplicates may be retained. Their degrees
\texttt{\small (\allowbreak{}n\_\allowbreak{}j,\allowbreak{}lambda\textasciicircum{}(\allowbreak{}j)\allowbreak{})\allowbreak{}}
give the columns of \texttt{\small D\_\allowbreak{}m}.

Compute the elimination ideal \[
J=\bigl(Z_j-f_j(X):1\le j\le q\bigr)
\cap\mathbb Q[Z_1,\ldots,Z_q].
\tag{G9.5}
\] Use any fixed monomial order to compute a homogeneous Gröbner basis
of \texttt{\small J}, and let \texttt{\small J\_\allowbreak{}0} be the
monomial ideal of its leading terms. Let
\texttt{\small B\_\allowbreak{}j} be the largest exponent of
\texttt{\small Z\_\allowbreak{}j} in the finite minimal monomial
generating set of \texttt{\small J\_\allowbreak{}0}, with
\texttt{\small B\_\allowbreak{}j=\allowbreak{}0} when it never occurs.
Partition the nonnegative exponent lattice, independently in each
coordinate, into the singleton cells
\texttt{\small 0,\allowbreak{}...,\allowbreak{}B\_\allowbreak{}j-\allowbreak{}1}
and the final interval
\texttt{\small B\_\allowbreak{}j+\allowbreak{}Z\_\allowbreak{}\{≥\allowbreak{}0\}}.
A resulting cell is uniquely
\texttt{\small u+\allowbreak{}Z\_\allowbreak{}\{≥\allowbreak{}0\}\textasciicircum{}S},
with \texttt{\small u\_\allowbreak{}j=\allowbreak{}B\_\allowbreak{}j} on
its free coordinates and a selected smaller integer on the other
coordinates. Retain exactly those cells with
\texttt{\small Z\textasciicircum{}u notin J\_\allowbreak{}0}. These
retained pairs form \texttt{\small C\_\allowbreak{}m}.

All steps are finite and effective: Gordan's lemma and Hilbert-basis
algorithms apply to (G9.2), finite-group Noether generation applies to
(G9.4), and Buchberger elimination terminates for (G9.5). Exact cited
locators are supplied in the background register. The Gröbner-basis step
uses only the explicitly constructed polynomials
\texttt{\small f\_\allowbreak{}j(\allowbreak{}X)\allowbreak{}}.

\subsubsection{Proof of the invariant-ring
realization}\label{proof-of-the-invariant-ring-realization}

Let \texttt{\small E=\allowbreak{}F=\allowbreak{}Q\textasciicircum{}m},
with left and right general linear actions on
\texttt{\small P=\allowbreak{}Sym(\allowbreak{}E tensor F)\allowbreak{}}.
Let \texttt{\small U} be the upper unipotent group on the right, and let
\texttt{\small T\_\allowbreak{}0} be the determinant-one diagonal torus
on the left. Cauchy duality gives \[
R=P^U=\bigoplus_\lambda S_\lambda E\otimes L_\lambda,
\tag{G9.6}
\] where \texttt{\small L\_\allowbreak{}lambda} is the right
highest-weight line. The subalgebra generated by the flag minors is all
of \texttt{\small R}. To see this, the flag minors of size
\texttt{\small j} span its component
\texttt{\small Lambda\textasciicircum{}j E tensor L\_\allowbreak{}(\allowbreak{}omega\_\allowbreak{}j)\allowbreak{}}.
The product of components with multiplicities
\texttt{\small lambda\_\allowbreak{}j-\allowbreak{}lambda\_\allowbreak{}(\allowbreak{}j+\allowbreak{}1)\allowbreak{}}
has right weight \texttt{\small lambda}. It is nonzero because
\texttt{\small R} is a domain. Its image is a nonzero left
\texttt{\small GL(\allowbreak{}E)\allowbreak{}} submodule of the
irreducible right-weight component
\texttt{\small S\_\allowbreak{}lambda E tensor L\_\allowbreak{}lambda},
hence it equals that whole component. Every partition is a sum of these
fundamental weights. This proves the claimed generation, including the
concrete generators used above.

The torus \texttt{\small T\_\allowbreak{}0} fixes exactly the equal
left-weight components
\texttt{\small (\allowbreak{}n\textasciicircum{}m)\allowbreak{}}.
Therefore \[
R^{T_0}=\bigoplus_{n,\lambda}
M_{\lambda,n}\otimes L_\lambda.
\tag{G9.7}
\] The left permutation matrices normalize
\texttt{\small T\_\allowbreak{}0} and commute with the right action.
Taking their invariants gives the actual finitely generated algebra \[
A_m=(R^{T_0})^{S_m},\qquad
\dim(A_m)_{(n,\lambda)}=a_{m,n}^{\lambda}.
\tag{G9.8}
\] The equality of dimensions is T-002 with its literal permutation
action.

For completeness, the finite presentation above indeed generates this
exact algebra. The polynomial ring in formal flag-minor variables
surjects equivariantly onto \texttt{\small R}. Extracting a torus weight
is exact: a preimage of a homogeneous invariant decomposes into weight
components, and its zero-weight component still maps to that invariant.
The torus-invariant monomials are exactly (G9.2); their Hilbert basis
consequently gives the surjection
\texttt{\small Q[\allowbreak{}Y\_\allowbreak{}1,\allowbreak{}...,\allowbreak{}Y\_\allowbreak{}s]\allowbreak{}→\allowbreak{}R\textasciicircum{}\{T\_\allowbreak{}0\}}.
The basis is preserved by the signed permutation action because
indecomposability in the monoid is preserved by coordinate permutations.
Averaging over \texttt{\small S\_\allowbreak{}m} makes invariants exact
as well. Noether's degree bound says that the averages of monomials of
ordinary \texttt{\small Y}-degree at most \texttt{\small m!} generate
\texttt{\small Q[\allowbreak{}Y]\allowbreak{}\textasciicircum{}\{S\_\allowbreak{}m\}}.
Their images thus generate precisely \texttt{\small A\_\allowbreak{}m}.
Their explicit image kernel is (G9.5). Every image generator is
homogeneous for
\texttt{\small (\allowbreak{}n,\allowbreak{}lambda)\allowbreak{}}, and
\texttt{\small n\_\allowbreak{}j>0}, because every nonconstant monomial
in (G9.2) has positive common row sum.

\subsubsection{Proof of the positive
count}\label{proof-of-the-positive-count}

Residue classes of monomials outside \texttt{\small J\_\allowbreak{}0}
form a basis of
\texttt{\small Q[\allowbreak{}Z]\allowbreak{}/\allowbreak{}J}, by
Gröbner division. The ideal is multihomogeneous for
\texttt{\small D\_\allowbreak{}m}, so division preserves the grading and
this is a homogeneous basis of \texttt{\small A\_\allowbreak{}m}. For
any monomial generator \texttt{\small Z\textasciicircum{}a} of
\texttt{\small J\_\allowbreak{}0}, the test
\texttt{\small v\_\allowbreak{}j>=\allowbreak{}a\_\allowbreak{}j} is
constant throughout each cell of the above partition: a fixed coordinate
is constant, and a free coordinate satisfies
\texttt{\small v\_\allowbreak{}j>=\allowbreak{}B\_\allowbreak{}j>=\allowbreak{}a\_\allowbreak{}j}.
Consequently membership in \texttt{\small J\_\allowbreak{}0} is constant
on the entire cell. Retaining its lower corner outside
\texttt{\small J\_\allowbreak{}0} retains exactly the standard
monomials, with no overlaps. Its exponents have the unique form
\texttt{\small u+\allowbreak{}z},
\texttt{\small z in Z\_\allowbreak{}\{≥\allowbreak{}0\}\textasciicircum{}S},
and its degree is
\texttt{\small D\_\allowbreak{}m u+\allowbreak{}D\_\allowbreak{}(\allowbreak{}m,\allowbreak{}S)\allowbreak{}z}.
Taking degree
\texttt{\small (\allowbreak{}n,\allowbreak{}lambda)\allowbreak{}} proves
(G9.1).

Finally its first degree equation is \[
\sum_{j\in S}n_j z_j=n-\sum_j n_j u_j,
\] with every \texttt{\small n\_\allowbreak{}j>=\allowbreak{}1}. Every
nonnegative real solution has each coordinate at most \texttt{\small n}
when the right side is nonnegative; a negative right side gives an empty
set. Thus the real polyhedra being counted are bounded. This completes
the proof for all parameters, including vanishing coefficients.

\subsubsection{The fully computed rank-two
instance}\label{the-fully-computed-rank-two-instance}

For \texttt{\small m=\allowbreak{}2}, the flag algebra is
\texttt{\small Q[\allowbreak{}x\_\allowbreak{}1,\allowbreak{}x\_\allowbreak{}2,\allowbreak{}delta]\allowbreak{}};
the right highest weights are
\texttt{\small (\allowbreak{}1,\allowbreak{}0)\allowbreak{},\allowbreak{}(\allowbreak{}1,\allowbreak{}0)\allowbreak{},\allowbreak{}(\allowbreak{}1,\allowbreak{}1)\allowbreak{}}.
The balanced torus algebra is
\texttt{\small Q[\allowbreak{}u,\allowbreak{}v]\allowbreak{}}, with
\texttt{\small u=\allowbreak{}x\_\allowbreak{}1x\_\allowbreak{}2},
\texttt{\small v=\allowbreak{}delta}. The transposition fixes
\texttt{\small u} and negates \texttt{\small v}, so \[
A_2=\mathbb Q[u,v^2],\qquad
\deg u=(1;2,0),\quad\deg v^2=(2;2,2).
\] Its initial ideal is zero, and G-009 gives the single polytope \[
\{(a,b)\in\mathbb R_{\ge0}^2:
a+2b=n,\ (2a+2b,2b)=\lambda\}.
\] This recovers T-006 through the exact invariant-algebra morphism. It
is also an explicit check that the row-permutation determinant sign has
survived the construction.

\subsubsection{Geometric and complexity
content}\label{geometric-and-complexity-content}

The generators and relations define affine schemes over
\texttt{\small Q}; the maps are given by the flag minors, the torus
monomial morphism, and the finite-group Reynolds morphism. Clearing the
finitely many rational denominators gives models over a localization of
\texttt{\small Z}, with the characteristic-zero fiber stated above. The
initial ideal gives a multigraded flat Gröbner degeneration, and its
standard monomials are the lattice points used in (G9.1). The count is
the class multiplicity in the same representation
\texttt{\small K\_\allowbreak{}0} and same balanced normalizer module as
T-002, now evaluated by a finite presentation.

The process includes a Hilbert-basis computation in
\texttt{\small 2\textasciicircum{}m-\allowbreak{}1} variables, Reynolds
enumeration through ordinary degree \texttt{\small m!}, polynomial
elimination, and a finite cell list depending on an initial ideal. No
bound polynomial in the varying input size has been proved for any of
those complete precomputation data. Thus this supplies a terminating
positive-polytope construction in the literal counting sense. The
smaller formulas T-012--T-013 and G-007 avoid that precomputation in
their stated families. GCT questions asking for efficiently constructed
polytopes or uniform nonvanishing certificates retain those additional
requirements; neither a complexity-class separation nor their failure
follows from the construction.

\section{General Schur coefficients and canonical-map
ranks}\label{general-schur-coefficients-and-canonical-map-ranks}

\subsection{G-010. The positive-polytope construction for general Schur
plethysm}\label{g-010.-the-positive-polytope-construction-for-general-schur-plethysm}

Let \texttt{\small mu} be a partition of
\texttt{\small m>=\allowbreak{}1}, \texttt{\small nu} a partition of
\texttt{\small n>=\allowbreak{}1}, and
\texttt{\small N=\allowbreak{}mn}. Let
\texttt{\small H=\allowbreak{}S\_\allowbreak{}n wr S\_\allowbreak{}m subset S\_\allowbreak{}N}
act on its grid of \texttt{\small N} letters, and let \[
W=[\nu]^{\otimes m}\otimes\operatorname{Inf}_{S_m}^H[\mu].
\] Use integral standard polytabloid Specht matrices and their rational
duals to specify the \texttt{\small H} action on
\texttt{\small W\textasciicircum{}vee} exactly. There is a terminating
construction of a finite set of bounded rational polytopes whose
lattice-point counts sum to
\texttt{\small p\_\allowbreak{}(\allowbreak{}mu,\allowbreak{}nu)\allowbreak{}\textasciicircum{}lambda},
for every output partition \texttt{\small lambda} of \texttt{\small N}.
The data use no unknown multiplicities. They are constructed as follows
from G-009.

Take the flag algebra \texttt{\small R\_\allowbreak{}N} and its balanced
torus algebra
\texttt{\small B\_\allowbreak{}N=\allowbreak{}R\_\allowbreak{}N\textasciicircum{}(\allowbreak{}T\_\allowbreak{}0)\allowbreak{}}.
Introduce variables spanning \texttt{\small W\textasciicircum{}vee},
each of auxiliary degree one, and form \[
C_{N,H,W}=(B_N\otimes\operatorname{Sym}(W^\vee))^H.
\tag{G10.1}
\] Its grading is
\texttt{\small (\allowbreak{}b,\allowbreak{}lambda,\allowbreak{}k)\allowbreak{}}:
balanced left weight
\texttt{\small (\allowbreak{}b\textasciicircum{}N)\allowbreak{}}, right
highest weight \texttt{\small lambda}, and auxiliary degree
\texttt{\small k}. It has \[
\boxed{\dim(C_{N,H,W})_{(1,\lambda,1)}=p_{\mu,\nu}^{\lambda}.}
\tag{G10.2}
\]

To produce its finite presentation, take the same Hilbert-basis
variables \texttt{\small Y} used for \texttt{\small B\_\allowbreak{}N},
add the \texttt{\small W\textasciicircum{}vee} variables
\texttt{\small Z}, average every monomial of ordinary
\texttt{\small (\allowbreak{}Y,\allowbreak{}Z)\allowbreak{}} degree from
one through \texttt{\small |H|}, and substitute the flag-minor
expressions for \texttt{\small Y}. The group action on \texttt{\small Z}
is its explicit rational Specht action; it need not permute these
variables. Discard zero images and introduce presentation variables
\texttt{\small T\_\allowbreak{}j} for the remaining polynomials
\texttt{\small f\_\allowbreak{}j(\allowbreak{}X,\allowbreak{}Z)\allowbreak{}}.
Compute the graph ideal
\texttt{\small (\allowbreak{}T\_\allowbreak{}j-\allowbreak{}f\_\allowbreak{}j(\allowbreak{}X,\allowbreak{}Z)\allowbreak{})\allowbreak{}}
and eliminate \textbf{both} the matrix variables \texttt{\small X} and
the auxiliary variables \texttt{\small Z}, leaving its intersection with
\texttt{\small Q[\allowbreak{}T]\allowbreak{}}. Apply the threshold-cell
construction of G-009 to this exact presentation kernel. The resulting
matrix has columns
\texttt{\small (\allowbreak{}b\_\allowbreak{}j,\allowbreak{}lambda\textasciicircum{}(\allowbreak{}j)\allowbreak{},\allowbreak{}k\_\allowbreak{}j)\allowbreak{}}
with
\texttt{\small b\_\allowbreak{}j+\allowbreak{}k\_\allowbreak{}j>=\allowbreak{}1}.
For each retained cell the polytope is \[
\{z\in\mathbb R_{\ge0}^{S}:
D_Sz=(1,\lambda,1)-Du\}.
\tag{G10.3}
\] The sum of its lattice-point counts is the coefficient in (G10.2).

\textbf{Proof.} The multilinear weight space
\texttt{\small (\allowbreak{}S\_\allowbreak{}lambda Q\textasciicircum{}N)\allowbreak{}\_\allowbreak{}(\allowbreak{}1\textasciicircum{}N)\allowbreak{}}
is the Specht representation
\texttt{\small [\allowbreak{}lambda]\allowbreak{}} under literal
permutation matrices. This is the multilinear Schur--Weyl identification
used in T-003 and recorded in Landsberg, Corollary 3.8(1). Thus the
auxiliary-degree-one component of (G10.1) in balanced degree one is \[
([\lambda]\otimes W^\vee)^H
=\operatorname{Hom}_H(W,[\lambda]),
\] which is the exact wreath-product formula T-003. This proves (G10.2).

The map from the polynomial ring on \texttt{\small Y,\allowbreak{}Z} to
\texttt{\small B\_\allowbreak{}N tensor Sym(\allowbreak{}W\textasciicircum{}vee)\allowbreak{}}
is equivariant and surjective. Characteristic-zero averaging is exact,
and Noether's degree bound for the finite group \texttt{\small H}
supplies exactly the stated invariant generators. Its proof only
requires a linear action, so the Specht matrices on \texttt{\small Z}
cause no additional assumption. Each average is homogeneous in all three
gradings: \texttt{\small H} preserves the two gradings of every
Hilbert-basis orbit, and preserves ordinary degree in \texttt{\small Z}.
Thus elimination and the standard-monomial cell argument of G-009 apply.
The linear functional \texttt{\small b+\allowbreak{}k} is positive on
every nonconstant generator, and has value two on the target. This
bounds every coordinate of every polytope by two divided by that
generator's positive grading. It proves boundedness and completes the
positive count. Many pieces are consequently empty; retaining them does
not change the formula.

For arbitrary outer Schur type and a row inner index, a smaller
construction uses \texttt{\small R\_\allowbreak{}m},
\texttt{\small H=\allowbreak{}S\_\allowbreak{}m}, and
\texttt{\small W=\allowbreak{}[\allowbreak{}mu]\allowbreak{}}, with
target
\texttt{\small (\allowbreak{}n,\allowbreak{}lambda,\allowbreak{}1)\allowbreak{}}.
It gives the same coefficient by T-003 and has the positive grading
bound \texttt{\small n+\allowbreak{}1}. Thus the general construction
includes both exact oscillator descriptions already established.

The algebraic-geometric statement is the invariant affine scheme (G10.1)
with the indicated finite-group quotient. Its auxiliary-degree-one
component is the covariant module for the explicit representation
\texttt{\small W}; its specified graded component dimension is
transported by the given wreath-product Hom isomorphism. All operations,
equations, and rational actions are specified, and after clearing
denominators give models over a localization of \texttt{\small Z}. No
polynomial-size bound for the construction as
\texttt{\small m,\allowbreak{}n,\allowbreak{}mu,\allowbreak{}nu} vary is
proved.

\subsection{G-011. A positive-polytope construction for each
canonical-map image
multiplicity}\label{g-011.-a-positive-polytope-construction-for-each-canonical-map-image-multiplicity}

For fixed \texttt{\small m,\allowbreak{}n>=\allowbreak{}1}, put
\texttt{\small N=\allowbreak{}mn}, let
\texttt{\small G=\allowbreak{}S\_\allowbreak{}N}, and take the row and
column wreath subgroups \texttt{\small H,\allowbreak{}K} used in G-002.
There is a terminating finite-presentation procedure, independent of the
output partition \texttt{\small lambda}, giving a finite sum of bounded
rational-polytope counts for \[
r_{m,n}^{\lambda}=\operatorname{rank}(P_K|_{[\lambda]^H}).
\tag{G11.1}
\] The procedure constructs the image of the specified canonical map,
including its possibly vanishing scalar on a multiplicity-one component.

\subsubsection{Explicit module
construction}\label{explicit-module-construction}

Let
\texttt{\small B=\allowbreak{}R\_\allowbreak{}N\textasciicircum{}(\allowbreak{}T\_\allowbreak{}0)\allowbreak{}}
be the balanced flag algebra, and let
\texttt{\small C=\allowbreak{}B\textasciicircum{}G=\allowbreak{}A\_\allowbreak{}N},
with its explicit homogeneous invariant generators
\texttt{\small c\_\allowbreak{}1(\allowbreak{}X)\allowbreak{},\allowbreak{}...,\allowbreak{}c\_\allowbreak{}q(\allowbreak{}X)\allowbreak{}}
from G-009. Let
\texttt{\small q\_\allowbreak{}1(\allowbreak{}X)\allowbreak{},\allowbreak{}...,\allowbreak{}q\_\allowbreak{}s(\allowbreak{}X)\allowbreak{}}
be the Hilbert-basis generators of \texttt{\small B}. Put
\texttt{\small h=\allowbreak{}|G|=\allowbreak{}N!}, and let
\texttt{\small u} range over the finite box
\texttt{\small 0<=\allowbreak{}u\_\allowbreak{}i<h}. Compute the
polynomial vectors, here ordinary scalar polynomials, \[
g_u(X)=P_KP_H\bigl(q_1(X)^{u_1}\cdots q_s(X)^{u_s}\bigr),
\tag{G11.2}
\] and discard zero images. Each retained generator has its explicitly
known degree
\texttt{\small delta\_\allowbreak{}u=\allowbreak{}(\allowbreak{}b\_\allowbreak{}u,\allowbreak{}lambda\textasciicircum{}(\allowbreak{}u)\allowbreak{})\allowbreak{}}.
Define the \texttt{\small C}-module \[
M=P_K(B^H)\subseteq B^K.
\tag{G11.3}
\] Its degree
\texttt{\small (\allowbreak{}1,\allowbreak{}lambda)\allowbreak{}} is the
image in (G11.1). Equations (G11.2) generate the entire module (G11.3),
not merely the degree under consideration.

To compute a finite presentation over
\texttt{\small Q[\allowbreak{}Z\_\allowbreak{}1,\allowbreak{}...,\allowbreak{}Z\_\allowbreak{}q]\allowbreak{}},
set
\texttt{\small Z\_\allowbreak{}j=\allowbreak{}c\_\allowbreak{}j(\allowbreak{}X)\allowbreak{}}.
The kernel of \[
\bigoplus_u\mathbb Q[Z]e_u\longrightarrow\mathbb Q[X],
\qquad e_u\longmapsto g_u(X)
\tag{G11.4}
\] is computable by syzygies and elimination. Explicitly, in
\texttt{\small Q[\allowbreak{}Z,\allowbreak{}X]\allowbreak{}} compute
the syzygy module of the row \[
(g_{u_1},\ldots,g_{u_t},Z_1-c_1(X),\ldots,Z_q-c_q(X)).
\] Project onto its first \texttt{\small t} components and intersect
with \texttt{\small Q[\allowbreak{}Z]\allowbreak{}\textasciicircum{}t}
by a module elimination order. This yields exactly the kernel in
(G11.4). Use a homogeneous module Gröbner basis. Its leading submodule
has the form
\texttt{\small directsum\_\allowbreak{}u J\_\allowbreak{}u e\_\allowbreak{}u},
with each \texttt{\small J\_\allowbreak{}u} a monomial ideal. Decompose
the standard monomials for each \texttt{\small J\_\allowbreak{}u} into
the threshold cells of G-009; discard a summand with
\texttt{\small J\_\allowbreak{}u=\allowbreak{}(\allowbreak{}1)\allowbreak{}}.
For each retained pair
\texttt{\small (\allowbreak{}u;\allowbreak{}v,\allowbreak{}S)\allowbreak{}}
count \[
\{z\in\mathbb Z_{\ge0}^{S}:
D_Sz=(1,\lambda)-\delta_u-Dv\}.
\tag{G11.5}
\] The sum of these nonnegative counts is (G11.1).

\subsubsection{Proof of generation and the rank
identity}\label{proof-of-generation-and-the-rank-identity}

Every element \texttt{\small q\_\allowbreak{}i} satisfies the monic
equation \[
\prod_{g\in G}(t-gq_i)=0\quad\text{at }t=q_i,
\] whose coefficients are in \texttt{\small C}. Reducing one exponent at
a time shows that the monomials \texttt{\small q\textasciicircum{}u},
\texttt{\small 0<=\allowbreak{}u\_\allowbreak{}i<h}, generate
\texttt{\small B} as a module over \texttt{\small C}. An invariant
coefficient commutes with both averages:
\texttt{\small P\_\allowbreak{}H(\allowbreak{}cf)\allowbreak{}=\allowbreak{}cP\_\allowbreak{}H(\allowbreak{}f)\allowbreak{}}
and
\texttt{\small P\_\allowbreak{}K(\allowbreak{}cf)\allowbreak{}=\allowbreak{}cP\_\allowbreak{}K(\allowbreak{}f)\allowbreak{}}
for \texttt{\small c in C}. Averaging the finite generating expression
first by \texttt{\small H} and then by \texttt{\small K} consequently
proves that (G11.2) generates (G11.3).

All actions commute with the right highest-weight grading and preserve
the balanced left grading. In degree
\texttt{\small (\allowbreak{}1,\allowbreak{}lambda)\allowbreak{}},
\texttt{\small B} is \texttt{\small [\allowbreak{}lambda]\allowbreak{}}
by the multilinear identification used in G-010. Thus (G11.3) in that
degree is exactly
\texttt{\small P\_\allowbreak{}K(\allowbreak{}[\allowbreak{}lambda]\allowbreak{}\textasciicircum{}H)\allowbreak{}},
with no substitution of an arbitrary map. This proves the asserted rank
identity.

The syzygy-elimination description is exact because a vector
\texttt{\small (\allowbreak{}a\_\allowbreak{}u)\allowbreak{}} belongs to
(G11.4)'s kernel precisely when
\texttt{\small sum a\_\allowbreak{}u(\allowbreak{}Z)\allowbreak{} g\_\allowbreak{}u(\allowbreak{}X)\allowbreak{}}
belongs to the ideal
\texttt{\small (\allowbreak{}Z\_\allowbreak{}j-\allowbreak{}c\_\allowbreak{}j(\allowbreak{}X)\allowbreak{})\allowbreak{}}.
Adding the coefficients expressing that ideal membership produces a
syzygy of the displayed row, and projecting a syzygy gives such an ideal
membership. The intersection enforces that the coefficients contain no
\texttt{\small X} variables. These two implications prove the claimed
kernel equality. Buchberger's module algorithm and elimination terminate
over the polynomial ring.

Assign
\texttt{\small deg e\_\allowbreak{}u=\allowbreak{}delta\_\allowbreak{}u}.
The kernel is homogeneous because (G11.4) respects degrees. Standard
module monomials therefore give a homogeneous vector-space basis of
\texttt{\small M}. The threshold-cell argument proves (G11.5) exactly as
in G-009. Every coefficient variable \texttt{\small Z\_\allowbreak{}j}
has positive balanced degree, so the first equation in (G11.5) bounds
all nonnegative real coordinates; a negative remaining degree gives the
empty set. This proves both positivity and boundedness.

The same finite data give the kernel and cokernel blocks: use the
presentation of \texttt{\small B\textasciicircum{}H} with generators
\texttt{\small P\_\allowbreak{}H(\allowbreak{}q\textasciicircum{}u)\allowbreak{}},
the presentation of \texttt{\small B\textasciicircum{}K} with generators
\texttt{\small P\_\allowbreak{}K(\allowbreak{}q\textasciicircum{}u)\allowbreak{}},
and the explicit images (G11.2). Module kernels and quotients have
homogeneous finite presentations over
\texttt{\small Q[\allowbreak{}Z]\allowbreak{}}, so the same
standard-monomial construction counts their multiplicities. This is a
terminating determination for each specified pair
\texttt{\small (\allowbreak{}m,\allowbreak{}n)\allowbreak{}}. The
theorem does not give a uniform inequality proving all open parameter
families of canonical injectivity or surjectivity.

In scheme language, (G11.3) is a coherent module on
\texttt{\small Spec C}; (G11.4) is a finite presentation of its
associated coherent sheaf, and the two idempotents are the precise
averaging correspondences. Module Gröbner degeneration preserves the
graded Hilbert function. Thus the positive lattice counts are
transported through an explicit coherent-module presentation to the
canonical map, rather than through its source and target dimensions
alone.

\section{A complete small presentation of the projected flag
algebra}\label{a-complete-small-presentation-of-the-projected-flag-algebra}

\subsection{G-012. Four explicit polytopes for the full
outer-degree-three
family}\label{g-012.-four-explicit-polytopes-for-the-full-outer-degree-three-family}

The algebra \texttt{\small A\_\allowbreak{}3} of G-009 has the
multihomogeneous presentation \[
\boxed{A_3=\mathbb Q[u,s,V,Z,d,W]/(d^2+4s^3+27u^2V,\ W^2-VZ).}
\tag{G12.1}
\] The respective degrees
\texttt{\small (\allowbreak{}n;\allowbreak{}lambda)\allowbreak{}} are \[
\begin{array}{c|rrrrrr}
&u&s&V&Z&d&W\\\hline
n&1&2&4&2&3&3\\
\lambda&(3,0,0)&(4,2,0)&(6,6,0)&(2,2,2)&(6,3,0)&(4,4,1).
\end{array}
\tag{G12.2}
\] Thus
\texttt{\small a\_\allowbreak{}(\allowbreak{}3,\allowbreak{}n)\allowbreak{}\textasciicircum{}lambda}
is the sum over
\texttt{\small epsilon,\allowbreak{}delta in \{0,\allowbreak{}1\}} of
the number of nonnegative integral quadruples
\texttt{\small (\allowbreak{}a,\allowbreak{}b,\allowbreak{}c,\allowbreak{}e)\allowbreak{}}
satisfying \[
\begin{aligned}
a+2b+4c+2e&=n-3\epsilon-3\delta,\\
3a+4b+6c+2e&=\lambda_1-6\epsilon-4\delta,\\
2b+6c+2e&=\lambda_2-3\epsilon-4\delta,\\
2e&=\lambda_3-\delta.
\end{aligned}
\tag{G12.3}
\] These are four bounded rational polytopes for partitions of length at
most three. Every coefficient with more than three parts is zero, by the
balanced-weight Cauchy identification in T-001. The coefficients are the
classical Thrall--Chen family. This presentation proves their direct
realization as standard monomials in the precise normalizer-invariant
flag algebra; it does not claim discovery of the coefficient family.

\subsubsection{Construction and proof}\label{construction-and-proof}

For a \texttt{\small 3 by 3} matrix, write the first-column minors as
\texttt{\small x\_\allowbreak{}1,\allowbreak{}x\_\allowbreak{}2,\allowbreak{}x\_\allowbreak{}3},
the oriented two-column minors as
\texttt{\small y\_\allowbreak{}1=\allowbreak{}p\_\allowbreak{}23,\allowbreak{}y\_\allowbreak{}2=\allowbreak{}-\allowbreak{}p\_\allowbreak{}13,\allowbreak{}y\_\allowbreak{}3=\allowbreak{}p\_\allowbreak{}12},
and the determinant as \texttt{\small z}. Then \[
R_3=\mathbb Q[x_1,x_2,x_3,y_1,y_2,y_3,z]/(x_1y_1+x_2y_2+x_3y_3).
\tag{G12.4}
\] The relation expands a determinant with two identical columns. These
minors generate by G-009. The target has dimension six: on the open
matrix chart where the first two leading minors and the determinant are
nonzero, right-upper-unipotent Gaussian elimination leaves six
independent parameters. The displayed relation is irreducible by Gauss's
lemma, applied as a linear polynomial in \texttt{\small y} with
relatively prime coefficients \texttt{\small x\_\allowbreak{}i}. Its
quotient is a domain of dimension six. The prime kernel has the same
dimension and contains this principal prime, so they are equal. This
proves (G12.4).

Set
\texttt{\small t\_\allowbreak{}i=\allowbreak{}x\_\allowbreak{}i y\_\allowbreak{}i},
\texttt{\small u=\allowbreak{}x\_\allowbreak{}1x\_\allowbreak{}2x\_\allowbreak{}3},
and
\texttt{\small v=\allowbreak{}y\_\allowbreak{}1y\_\allowbreak{}2y\_\allowbreak{}3}.
A monomial in \texttt{\small x,\allowbreak{}y} has balanced weight
precisely when
\texttt{\small a\_\allowbreak{}i-\allowbreak{}b\_\allowbreak{}i} is
constant. A nonnegative difference factors out a power of
\texttt{\small u}; a nonpositive difference factors out a power of
\texttt{\small v}; the rest is a product of the
\texttt{\small t\_\allowbreak{}i}. The torus-invariant polynomial ring
is consequently generated by
\texttt{\small u,\allowbreak{}v,\allowbreak{}t\_\allowbreak{}1,\allowbreak{}t\_\allowbreak{}2,\allowbreak{}t\_\allowbreak{}3,\allowbreak{}z}
with sole relation
\texttt{\small uv=\allowbreak{}t\_\allowbreak{}1t\_\allowbreak{}2t\_\allowbreak{}3},
as also follows from the unique expressions containing \texttt{\small u}
or \texttt{\small v} but not both. Exact torus weight extraction carries
the flag relation to
\texttt{\small sum t\_\allowbreak{}i=\allowbreak{}0}. Therefore \[
B_3=\mathbb Q[u,v,t_1,t_2,t_3,z]/(t_1+t_2+t_3,\ uv-t_1t_2t_3).
\tag{G12.5}
\]

A permutation sends \texttt{\small x\_\allowbreak{}i} to
\texttt{\small x\_\allowbreak{}(\allowbreak{}sigma i)\allowbreak{}} and
\texttt{\small y\_\allowbreak{}i} to
\texttt{\small sgn(\allowbreak{}sigma)\allowbreak{}y\_\allowbreak{}(\allowbreak{}sigma i)\allowbreak{}}.
It fixes \texttt{\small u}, negates each of
\texttt{\small v,\allowbreak{}z} under an odd permutation, and sends
\texttt{\small t\_\allowbreak{}i} to
\texttt{\small sgn(\allowbreak{}sigma)\allowbreak{}t\_\allowbreak{}(\allowbreak{}sigma i)\allowbreak{}}.
On the plane \texttt{\small sum t\_\allowbreak{}i=\allowbreak{}0} put \[
s=t_1t_2+t_1t_3+t_2t_3,\quad f=t_1t_2t_3,\quad
d=(t_1-t_2)(t_1-t_3)(t_2-t_3).
\] The elements \texttt{\small s,\allowbreak{}d} are invariant under
this twisted action, \texttt{\small f} has sign character, and \[
d^2=-4s^3-27f^2.
\tag{G12.6}
\] Its invariant and sign-isotypic polynomial subspaces are respectively
\[
\mathbb Q[s,d],\qquad f\mathbb Q[s,d].
\tag{G12.7}
\] Indeed ordinary alternating-group invariants are
\texttt{\small Q[\allowbreak{}s,\allowbreak{}f]\allowbreak{} directsum d Q[\allowbreak{}s,\allowbreak{}f]\allowbreak{}}:
average under an ordinary transposition and divide the alternating part
by the Vandermonde. Under a twisted transposition,
\texttt{\small s,\allowbreak{}d} stay fixed and \texttt{\small f}
changes sign. Taking even and odd parts in \texttt{\small f} and using
(G12.6) proves (G12.7).

Adjoining the sign variables \texttt{\small v,\allowbreak{}z}, formula
(G12.7) shows that before imposing \texttt{\small uv=\allowbreak{}f} the
invariants are generated by \[
u,s,d,\quad V=v^2,\quad W=vz,\quad Z=z^2,\quad fv,\quad fz.
\] After that relation \texttt{\small fv=\allowbreak{}uV},
\texttt{\small fz=\allowbreak{}uW}, so
\texttt{\small u,\allowbreak{}s,\allowbreak{}d,\allowbreak{}V,\allowbreak{}W,\allowbreak{}Z}
generate \texttt{\small A\_\allowbreak{}3}; exact averaging carries
every quotient invariant to one of these images. Equations (G12.6) and
\texttt{\small W\textasciicircum{}2=\allowbreak{}VZ} give (G12.1)'s
relations. Hence its map to \texttt{\small A\_\allowbreak{}3} is
surjective.

To exclude further relations, \texttt{\small s,\allowbreak{}f} are
algebraically independent on the two-dimensional \texttt{\small t} plane
by the elementary symmetric-polynomial theorem. In the fraction field of
(G12.5), take \texttt{\small v,\allowbreak{}z} independent and put
\texttt{\small u=\allowbreak{}f/\allowbreak{}v}. Then
\texttt{\small u,\allowbreak{}s,\allowbreak{}V=\allowbreak{}v\textasciicircum{}2,\allowbreak{}Z=\allowbreak{}z\textasciicircum{}2}
are algebraically independent. In
\texttt{\small F=\allowbreak{}Q(\allowbreak{}u,\allowbreak{}s,\allowbreak{}V,\allowbreak{}Z)\allowbreak{}},
the element
\texttt{\small D=\allowbreak{}-\allowbreak{}4s\textasciicircum{}3-\allowbreak{}27u\textasciicircum{}2V}
is nonsquare because its pole order at
\texttt{\small s=\allowbreak{}infinity} is three. Both
\texttt{\small VZ} and
\texttt{\small D/\allowbreak{}(\allowbreak{}VZ)\allowbreak{}} are
nonsquare by their odd valuation at \texttt{\small Z=\allowbreak{}0}.
Their square classes are independent. Thus the actual
\texttt{\small d,\allowbreak{}W} generate a degree-four extension of
\texttt{\small F}, and
\texttt{\small 1,\allowbreak{}d,\allowbreak{}W,\allowbreak{}dW} are
linearly independent over \texttt{\small F}.

The quotient in (G12.1) is free over
\texttt{\small Q[\allowbreak{}u,\allowbreak{}s,\allowbreak{}V,\allowbreak{}Z]\allowbreak{}}
with basis
\texttt{\small 1,\allowbreak{}d,\allowbreak{}W,\allowbreak{}dW}, since
its two relations are monic quadratics in different variables. Their
images are independent over \texttt{\small F}, hence over the polynomial
subring. The presentation is therefore injective as well as surjective.

Its homogeneous basis is \[
u^a s^b V^c Z^e d^\epsilon W^\delta,
\qquad a,b,c,e\ge0,\quad\epsilon,\delta\in\{0,1\}.
\] Degrees give (G12.3), whose first equation bounds every nonnegative
real coordinate. Equation (G9.8) identifies this dimension with the
desired coefficient. This proves the four-polytope formula with the
finite sign action evaluated explicitly.

The geometric construction is the affine flag cone, followed by the
diagonal-torus and constant-\texttt{\small S\_\allowbreak{}3} quotients.
The monic relations supply the four-element free module and its positive
monomial description. This is the precise algebraic-geometric
realization of the outer-degree-three normalizer projection.

\section{Stable plethysm: explicit positive polytopes through six
columns}\label{stable-plethysm-explicit-positive-polytopes-through-six-columns}

Run
\texttt{\small GCT\_\allowbreak{}PLETH-\allowbreak{}20261121-\allowbreak{}theta};
GPT-6 Astra Pro; effort not surfaced. This supplements
\texttt{\small theta\_\allowbreak{}report.md}. The new work here makes
the finite indexing and the actual polytope constraints explicit. The
underlying stability theorem, the character reduction, the Littlewood
identities and the Littlewood--Richardson rule are cited results, not
claimed discoveries.

\subsection{T-010. Exact current stable family and its known four-column
polytope}\label{t-010.-exact-current-stable-family-and-its-known-four-column-polytope}

Write
\texttt{\small lambda[\allowbreak{}N]\allowbreak{}=\allowbreak{}(\allowbreak{}N-\allowbreak{}|lambda|,\allowbreak{}lambda\_\allowbreak{}1,\allowbreak{}lambda\_\allowbreak{}2,\allowbreak{}...)\allowbreak{}}
and

\[
c(\lambda)=\langle s_\lambda,H[H-1-h_1]\rangle,
\qquad H=\sum_{i\ge0}h_i.
\]

Sarkar, \emph{Stability of plethysm coefficients and modified polynomial
induction}, arXiv:2610.07929v2 (14 October 2026), Theorem D = Theorem
6.4, proves

\[
a_{n,m}^{\lambda[mn]}=c(\lambda)
\quad\bigl(m\ge\lambda_1,\ n\ge\lfloor(|\lambda|+\lambda_1)/2\rfloor,
\ mn-|\lambda|\ge\lambda_1\bigr).
\tag{T10.1}
\]

For
\texttt{\small lambda=\allowbreak{}(\allowbreak{}4\textasciicircum{}u4,\allowbreak{}3\textasciicircum{}u3,\allowbreak{}2\textasciicircum{}u2,\allowbreak{}1\textasciicircum{}u1)\allowbreak{}},
his Theorem E = Theorem 7.3 identifies
\texttt{\small c(\allowbreak{}lambda)\allowbreak{}} with the integer
points
\texttt{\small (\allowbreak{}t,\allowbreak{}epsilon\_\allowbreak{}1,\allowbreak{}...,\allowbreak{}epsilon\_\allowbreak{}4)\allowbreak{}}
satisfying

\[
0\le\epsilon_i\le1,\quad t+\sum_i\epsilon_i\ge3,\quad
\epsilon_1-\epsilon_2=u_1,\quad \epsilon_3-\epsilon_4=u_3,
\quad\epsilon_2-\epsilon_3\le u_2,\quad t+\epsilon_4\le u_4.
\tag{T10.2}
\]

These are actual finite coefficients in the range (T10.1). The complete
proofs of Theorems 6.4 and 7.3 were read and checked in this run. Source
bytes were downloaded from
\texttt{\small https:\allowbreak{}/\allowbreak{}/\allowbreak{}arxiv.org/\allowbreak{}src/\allowbreak{}2610.07929v2}
and are retained in \texttt{\small sarkar/\allowbreak{}main.tex}.

Sarkar's Proposition 7.1 gives the exact character reduction, for
\texttt{\small k>=\allowbreak{}lambda\_\allowbreak{}1} and
\texttt{\small V=\allowbreak{}C\textasciicircum{}k},

\[
c(\lambda)=[S_{\lambda'}V:M_k],\qquad
M_k=\bigotimes_{i=2}^k F_i(\Lambda^iV),
\quad F_i=\operatorname{Sym}\ (i\text{ even}),\quad
F_i=\Lambda\ (i\text{ odd}).
\tag{T10.3}
\]

Here the bracket denotes multiplicity, \texttt{\small lambda'} denotes
conjugation, and each full symmetric or exterior algebra is graded by
ordinary scalar weight. Only total scalar degree
\texttt{\small |lambda|} contributes. The source's Remark 7.5 notes the
LR approach for five columns; T-012 and T-013 give full finite indexing
and polytope constraints through six columns.

\subsubsection{The exact lambda-ring dictionary for
(T10.3)}\label{the-exact-lambda-ring-dictionary-for-t10.3}

The involution \texttt{\small omega} of the symmetric-function ring
satisfies
\texttt{\small omega(\allowbreak{}p\_\allowbreak{}r)\allowbreak{}=\allowbreak{}(\allowbreak{}-\allowbreak{}1)\allowbreak{}\textasciicircum{}\{r-\allowbreak{}1\}p\_\allowbreak{}r}
and
\texttt{\small omega(\allowbreak{}s\_\allowbreak{}lambda)\allowbreak{}=\allowbreak{}s\_\allowbreak{}lambda'}.
On a homogeneous input \texttt{\small h\_\allowbreak{}i}, expansion in
power sums gives

\[
\omega(p_r[h_i])=(-1)^{(r-1)i}p_r[e_i].
\]

Indeed, each term \texttt{\small p\_\allowbreak{}\{r alpha\}} has sign
\texttt{\small (\allowbreak{}-\allowbreak{}1)\allowbreak{}\textasciicircum{}\{ri-\allowbreak{}length(\allowbreak{}alpha)\allowbreak{}\}},
while the corresponding term of
\texttt{\small p\_\allowbreak{}r[\allowbreak{}e\_\allowbreak{}i]\allowbreak{}}
has sign
\texttt{\small (\allowbreak{}-\allowbreak{}1)\allowbreak{}\textasciicircum{}\{i-\allowbreak{}length(\allowbreak{}alpha)\allowbreak{}\}}.
Exponentiating
\texttt{\small sum\_\allowbreak{}r p\_\allowbreak{}r/\allowbreak{}r}
yields

\[
\omega(H[h_i])=H[e_i]\ (i\text{ even}),\qquad
\omega(H[h_i])=E[e_i]\ (i\text{ odd}),
\quad E=\sum_{j\ge0}e_j.
\tag{T10.4}
\]

The completed product
\texttt{\small H[\allowbreak{}H-\allowbreak{}1-\allowbreak{}h\_\allowbreak{}1]\allowbreak{}=\allowbreak{}product\_\allowbreak{}\{i>=\allowbreak{}2\}H[\allowbreak{}h\_\allowbreak{}i]\allowbreak{}}
is well-defined degree by degree. Specialization to \texttt{\small k}
variables sends \texttt{\small e\_\allowbreak{}i} to zero for
\texttt{\small i>k}, giving exactly (T10.3). On locally free sheaves
this is the dictionary between symmetric powers of exterior powers and
exterior powers of exterior powers in the representation Grothendieck
ring. All factors of any given total scalar degree are finite locally
free polynomial functors. The decomposition into characteristic-zero
Schur summands takes place after base change to a characteristic-zero
field; the polynomial operations themselves are defined over
\texttt{\small Z}.

\subsection{T-011. An explicit Littlewood--Richardson
polytope}\label{t-011.-an-explicit-littlewoodrichardson-polytope}

For partitions
\texttt{\small alpha,\allowbreak{}beta,\allowbreak{}gamma} of length at
most \texttt{\small r}, padded by zeros, let
\texttt{\small L\_\allowbreak{}r(\allowbreak{}alpha,\allowbreak{}beta;\allowbreak{}gamma)\allowbreak{}}
consist of real arrays
\texttt{\small x\_\allowbreak{}(\allowbreak{}i,\allowbreak{}j)\allowbreak{}},
\texttt{\small 1<=\allowbreak{}i,\allowbreak{}j<=\allowbreak{}r}, with

\[
x_{ij}\ge0,\qquad x_{ij}=0\text{ for }i>j,
\tag{T11.1}
\] \[
\sum_{i=1}^r x_{ij}=\gamma_j-\alpha_j,
\qquad \sum_{j=1}^r x_{ij}=\beta_i,
\tag{T11.2}
\] \[
\alpha_j+\sum_{i=1}^{h}x_{ij}
\le\alpha_{j-1}+\sum_{i=1}^{h-1}x_{i,j-1}
\quad(2\le j\le r,\ 1\le h\le r),
\tag{T11.3}
\] \[
\sum_{s=1}^{j-1}x_{i,s}\ge\sum_{s=1}^{j}x_{i+1,s}
\quad(1\le i<r,\ 1\le j\le r).
\tag{T11.4}
\]

Boundary partitions are required to have nonnegative entries and to be
weakly decreasing. An invalid boundary gives the empty polytope. Its
integer points count
\texttt{\small c\_\allowbreak{}\{alpha,\allowbreak{}beta\}\textasciicircum{}gamma}.

In the later gluing constructions, membership in
\texttt{\small L\_\allowbreak{}r} means precisely the displayed linear
constraints also for real boundary variables. At a lattice point those
boundary variables become integral partitions, so the same tableau
bijection applies.

\textbf{Proof of the coordinate dictionary.} From an integer array fill
row \texttt{\small j} of the skew diagram
\texttt{\small gamma/\allowbreak{}alpha} with
\texttt{\small x\_\allowbreak{}(\allowbreak{}1,\allowbreak{}j)\allowbreak{}}
ones, then
\texttt{\small x\_\allowbreak{}(\allowbreak{}2,\allowbreak{}j)\allowbreak{}}
twos, and so on. The row and content conditions are (T11.2), and the
filling is unique. In row \texttt{\small j}, the last box with entry at
most \texttt{\small h} is at column
\texttt{\small alpha\_\allowbreak{}j+\allowbreak{}sum\_\allowbreak{}(\allowbreak{}i<=\allowbreak{}h)\allowbreak{}x\_\allowbreak{}(\allowbreak{}i,\allowbreak{}j)\allowbreak{}}.
For strict increase down columns its upper neighbor must have entry at
most \texttt{\small h-\allowbreak{}1}; (T11.3) states exactly that this
last column lies within the removed upper-row prefix or within the
upper-row prefix of entries at most \texttt{\small h-\allowbreak{}1}.
Applying this for every \texttt{\small h} gives column strictness, and
column strictness gives these inequalities.

Read each row from right to left, in order from top row to bottom row.
For the ballot inequality comparing \texttt{\small i} and
\texttt{\small i+\allowbreak{}1}, the largest deficit during row
\texttt{\small j} occurs after all its \texttt{\small i+\allowbreak{}1}
entries have been read and before any of its \texttt{\small i} entries
have been read. The required inequality at that point is exactly
(T11.4). Thus these conditions are equivalent to the lattice-word
condition. The lattice-word condition also forces every entry in row
\texttt{\small j} to be at most \texttt{\small j}, so the triangular
zeros in (T11.1) impose no additional restriction. This is a bijection
with Littlewood--Richardson tableaux. The classical LR rule gives the
coefficient; see Macdonald, Chapter I, §9, or the
representation-theoretic LR rule in the attached Landsberg source, §4,
equations preceding \texttt{\small pieri}.

Each coordinate is bounded by the prescribed total content. Thus this is
a bounded rational polytope, including the empty case.

\subsection{T-012. A fully explicit positive formula for stable tails of
width
five}\label{t-012.-a-fully-explicit-positive-formula-for-stable-tails-of-width-five}

Let \texttt{\small lambda\_\allowbreak{}1<=\allowbreak{}5}, put
\texttt{\small mu=\allowbreak{}lambda'}, padded to length five, and
\texttt{\small N=\allowbreak{}|lambda|}. Define the following finite set
\texttt{\small B\_\allowbreak{}5} of pairs
\texttt{\small (\allowbreak{}j,\allowbreak{}beta)\allowbreak{}}. Each
listed \texttt{\small beta} has multiplicity one in
\texttt{\small Lambda\textasciicircum{}j(\allowbreak{}Lambda\textasciicircum{}3 C\textasciicircum{}5)\allowbreak{}}:

\begin{longtable}[]{@{}rl@{}}
\toprule\noalign{}
j & beta \\
\midrule\noalign{}
\endhead
\bottomrule\noalign{}
\endlastfoot
0 & (0,0,0,0,0) \\
1 & (1,1,1,0,0) \\
2 & (2,2,1,1,0) \\
3 & (3,3,1,1,1); (3,2,2,2,0) \\
4 & (4,3,2,2,1); (3,3,3,3,0) \\
5 & (5,3,3,2,2); (4,4,3,3,1) \\
6 & (6,3,3,3,3); (5,4,4,3,2) \\
7 & (6,4,4,4,3); (5,5,5,3,3) \\
8 & (6,5,5,4,4) \\
9 & (6,6,5,5,5) \\
10 & (6,6,6,6,6) \\
\end{longtable}

For
\texttt{\small (\allowbreak{}j,\allowbreak{}beta)\allowbreak{} in B\_\allowbreak{}5}
and \texttt{\small epsilon in \{0,\allowbreak{}1\}}, define
\texttt{\small Q\_\allowbreak{}(\allowbreak{}lambda,\allowbreak{}beta,\allowbreak{}epsilon)\allowbreak{}}
using real variables \texttt{\small P,\allowbreak{}Q,\allowbreak{}q}, a
weakly decreasing nonnegative five-tuple \texttt{\small nu}, and two
arrays \texttt{\small x,\allowbreak{}y} as follows:

\[
P\ge Q\ge0,\quad q\ge0,\qquad
\alpha=(P,P,Q,Q,0),\quad \eta=(q,q,q,q,0),\quad
\gamma=\mu-(\epsilon^5)\ge0,
\tag{T12.1}
\] \[
x\in L_5(\alpha,\beta;\nu),\qquad
y\in L_5(\nu,\eta;\gamma),
\tag{T12.2}
\] \[
2(P+Q)+3j+4q+5\epsilon=N.
\tag{T12.3}
\]

Every condition is linear with integer coefficients. The target boundary
condition in (T12.1) makes the piece empty when needed. Then

\[
\boxed{c(\lambda)=
\sum_{(j,\beta)\in B_5}\ \sum_{\epsilon=0}^{1}
\#\bigl(Q_{\lambda,\beta,\epsilon}\cap\mathbb Z^{58}\bigr).}
\tag{T12.4}
\]

The ambient dimension counts three scalar parameters, five intermediate
coordinates, and two five-by-five arrays; redundant triangular zeros can
reduce it. The formula is a finite sum of 32 bounded rational polytopes,
with tags retaining disjointness. It makes no homogeneous-dilation
assertion.

\textbf{Proof.} By (T10.3), the required module is

\[
M_5=\operatorname{Sym}(\Lambda^2V)\otimes\Lambda(\Lambda^3V)
\otimes\operatorname{Sym}(\Lambda^4V)\otimes\Lambda(\Lambda^5V).
\]

The classical Littlewood decomposition of
\texttt{\small Sym(\allowbreak{}Lambda\textasciicircum{}2 V)\allowbreak{}}
has all even-column partitions, each once; in rank five these are
\texttt{\small (\allowbreak{}P,\allowbreak{}P,\allowbreak{}Q,\allowbreak{}Q,\allowbreak{}0)\allowbreak{}}.
This follows as well by applying (T10.4) to the symmetric-matrix
identity T-007. The natural perfect wedge pairing gives
\texttt{\small Lambda\textasciicircum{}4 V=\allowbreak{}V\textasciicircum{}* tensor det V};
hence

\[
\operatorname{Sym}^q(\Lambda^4V)
=\operatorname{Sym}^q(V^*)\otimes\det^q
=S_{(q,q,q,q,0)}V.
\]

Finally
\texttt{\small Lambda(\allowbreak{}Lambda\textasciicircum{}5 V)\allowbreak{}=\allowbreak{}1+\allowbreak{}det V}.
The complete finite decomposition of the second factor is the displayed
\texttt{\small B\_\allowbreak{}5}; the exact finite-character
certificate is described under T-013 below and in
\texttt{\small few\_\allowbreak{}columns.py}.

The multiplicity of \texttt{\small S\_\allowbreak{}mu V} is therefore

\[
\sum_{(j,\beta),\epsilon,P\ge Q\ge0,q\ge0,\nu}
c_{\alpha,\beta}^{\nu}c_{\nu,\eta}^{\mu-(\epsilon^5)}.
\]

Every summand has scalar degree
\texttt{\small 2(\allowbreak{}P+\allowbreak{}Q)\allowbreak{}+\allowbreak{}3j+\allowbreak{}4q+\allowbreak{}5epsilon}.
Only degree \texttt{\small N} contributes. Apply the integer-point
bijection T-011 twice, retaining
\texttt{\small P,\allowbreak{}Q,\allowbreak{}q,\allowbreak{}nu} as
coordinates rather than summing them externally. This proves (T12.4).
Equation (T12.3) bounds the three scalar parameters; the two LR
row/content conditions bound
\texttt{\small nu,\allowbreak{}x,\allowbreak{}y} by \texttt{\small N}.
Thus every piece is bounded.

\subsection{T-013. A fully explicit positive formula for stable tails of
width
six}\label{t-013.-a-fully-explicit-positive-formula-for-stable-tails-of-width-six}

Let \texttt{\small lambda\_\allowbreak{}1<=\allowbreak{}6}, put
\texttt{\small mu=\allowbreak{}lambda'} padded to length six, and
\texttt{\small N=\allowbreak{}|lambda|}. Let
\texttt{\small B\_\allowbreak{}6} be the complete 156-row list in
\texttt{\small EXTERIOR\_\allowbreak{}CUBE\_\allowbreak{}RANK\_\allowbreak{}6.tsv}.
Its pair \texttt{\small (\allowbreak{}j,\allowbreak{}beta)\allowbreak{}}
denotes one copy of
\texttt{\small S\_\allowbreak{}beta C\textasciicircum{}6} in
\texttt{\small Lambda\textasciicircum{}j(\allowbreak{}Lambda\textasciicircum{}3 C\textasciicircum{}6)\allowbreak{}};
every listed multiplicity is one,
\texttt{\small 0<=\allowbreak{}j<=\allowbreak{}20}, and
\texttt{\small |beta|=\allowbreak{}3j}.

For
\texttt{\small (\allowbreak{}j,\allowbreak{}beta)\allowbreak{} in B\_\allowbreak{}6}
and \texttt{\small q in \{0,\allowbreak{}...,\allowbreak{}6\}}, set

\[
\theta(q)=(q^{6-q},(q-1)^q).
\tag{T13.1}
\]

For \texttt{\small q=\allowbreak{}0} the second block has length zero
and
\texttt{\small theta(\allowbreak{}0)\allowbreak{}=\allowbreak{}(\allowbreak{}0\textasciicircum{}6)\allowbreak{}}.
Define
\texttt{\small R\_\allowbreak{}(\allowbreak{}lambda,\allowbreak{}beta,\allowbreak{}q)\allowbreak{}}
by variables
\texttt{\small P,\allowbreak{}Q,\allowbreak{}R,\allowbreak{}a,\allowbreak{}b,\allowbreak{}c,\allowbreak{}t},
intermediate partitions \texttt{\small nu,\allowbreak{}xi} of length at
most six, and three six-by-six arrays
\texttt{\small x,\allowbreak{}y,\allowbreak{}z} satisfying

\[
P\ge Q\ge R\ge0,\quad a\ge b\ge c\ge0,\quad t\ge0,
\tag{T13.2}
\] \[
\alpha=(P,P,Q,Q,R,R),\qquad
\eta=(a+b,a+b,a+c,a+c,b+c,b+c),\qquad
\gamma=\mu-(t^6)\ge0,
\tag{T13.3}
\] \[
x\in L_6(\alpha,\beta;\nu),\qquad
y\in L_6(\nu,\eta;\xi),\qquad
z\in L_6(\xi,\theta(q);\gamma),
\tag{T13.4}
\] \[
2(P+Q+R)+3j+4(a+b+c)+5q+6t=N.
\tag{T13.5}
\]

Then

\[
\boxed{c(\lambda)=
\sum_{(j,\beta)\in B_6}\ \sum_{q=0}^{6}
\#\bigl(R_{\lambda,\beta,q}\cap\mathbb Z^{127}\bigr).}
\tag{T13.6}
\]

There are 1092 tagged bounded rational polytopes in this formula. There
are seven scalar variables, twelve intermediate-partition coordinates,
and 108 array coordinates, totaling 127. Empty pieces are allowed. Thus
the finite coefficient in (T10.1) has an explicit positive lattice-point
formula for every tail of width at most six.

\textbf{Proof of the representation formula.} The exact factorization is

\[
M_6=\operatorname{Sym}(\Lambda^2V)\otimes\Lambda(\Lambda^3V)
\otimes\operatorname{Sym}(\Lambda^4V)\otimes\Lambda(\Lambda^5V)
\otimes\operatorname{Sym}(\Lambda^6V).
\]

The first factor has partitions
\texttt{\small alpha=\allowbreak{}(\allowbreak{}P,\allowbreak{}P,\allowbreak{}Q,\allowbreak{}Q,\allowbreak{}R,\allowbreak{}R)\allowbreak{}},
each once. The perfect wedge pairing
\texttt{\small Lambda\textasciicircum{}4 V=\allowbreak{}Lambda\textasciicircum{}2 V\textasciicircum{}* tensor det V}
sends the summand
\texttt{\small (\allowbreak{}a,\allowbreak{}a,\allowbreak{}b,\allowbreak{}b,\allowbreak{}c,\allowbreak{}c)\allowbreak{}}
in symmetric degree
\texttt{\small d=\allowbreak{}a+\allowbreak{}b+\allowbreak{}c} to its
dual tensored with \texttt{\small det\textasciicircum{}d}. Its highest
weight is

\[
(d-c,d-c,d-b,d-b,d-a,d-a)
=(a+b,a+b,a+c,a+c,b+c,b+c)=\eta.
\]

This proves the asserted indexing of the third factor with no omitted
determinant shift. The pairing
\texttt{\small Lambda\textasciicircum{}5 V=\allowbreak{}V\textasciicircum{}* tensor det V}
gives

\[
\Lambda^q(\Lambda^5V)=\Lambda^q(V^*)\otimes\det^q
=S_{\theta(q)}V,\qquad 0\le q\le6.
\]

The last factor has summands \texttt{\small det\textasciicircum{}t},
\texttt{\small t>=\allowbreak{}0}. The second factor is the fixed table
\texttt{\small B\_\allowbreak{}6}. Iterating the LR rule three times
yields

\[
\sum_{(j,\beta),q,P,Q,R,a,b,c,t,\nu,\xi}
c_{\alpha,\beta}^{\nu}
c_{\nu,\eta}^{\xi}
c_{\xi,\theta(q)}^{\mu-(t^6)}.
\]

The scalar degrees of the five factors are respectively
\texttt{\small 2(\allowbreak{}P+\allowbreak{}Q+\allowbreak{}R)\allowbreak{},\allowbreak{}3j,\allowbreak{}4(\allowbreak{}a+\allowbreak{}b+\allowbreak{}c)\allowbreak{},\allowbreak{}5q,\allowbreak{}6t}.
Keeping degree \texttt{\small N} gives (T13.5). Gluing the three
integer-point descriptions from T-011 along the displayed intermediate
partitions gives (T13.6). Every scalar variable in (T13.2) is bounded by
(T13.5), every intermediate partition has size at most
\texttt{\small N}, and each nonnegative array entry is bounded by a row
or content sum at most \texttt{\small N}. This proves boundedness as
well as finiteness of the formula.

\subsubsection{\texorpdfstring{Complete finite certificate for
\texttt{\small B\_\allowbreak{}5} and
\texttt{\small B\_\allowbreak{}6}}{Complete finite certificate for  and }}\label{complete-finite-certificate-for-b_5-and-b_6}

For \texttt{\small r=\allowbreak{}5,\allowbreak{}6} the exact torus
character of the entire finite exterior algebra is

\[
C_r(t;x)=\prod_{I\subset\{1,\ldots,r\},\ |I|=3}
\left(1+t\prod_{i\in I}x_i\right).
\tag{T13.7}
\]

Let
\texttt{\small rho=\allowbreak{}(\allowbreak{}r-\allowbreak{}1,\allowbreak{}...,\allowbreak{}0)\allowbreak{}}
and write \texttt{\small M\_\allowbreak{}j(\allowbreak{}w)\allowbreak{}}
for the coefficient of
\texttt{\small t\textasciicircum{}j x\textasciicircum{}w} in this finite
product, with zero for unsupported weights. The Weyl alternant identity
gives the exact multiplicity

\[
b_{j,\beta}=\sum_{\sigma\in S_r}\operatorname{sgn}(\sigma)
M_j(\beta+\rho-\sigma\rho).
\tag{T13.8}
\]

The integer program \texttt{\small few\_\allowbreak{}columns.py} expands
(T13.7) without floating point arithmetic, evaluates (T13.8) on every
dominant weight in its support, and writes the complete tables. Every
irreducible highest weight is a dominant weight in that support, so the
candidate set is exhaustive. The resulting tables have respectively 16
and 156 entries, all multiplicity one. As a separate check, the program
reconstructs the character of each listed Schur module by recursively
enumerating interlacing restrictions and compares every monomial in
every exterior degree with (T13.7). Both complete comparisons pass.
Dimension totals are
\texttt{\small 2\textasciicircum{}10=\allowbreak{}1024} and
\texttt{\small 2\textasciicircum{}20=\allowbreak{}1048576}.

This is a finite exact computation certificate, together with its
mathematical extraction formula and an independent character
reconstruction. It is not labeled a proof-assistant certificate. The
table is fixed data independent of the user-varying partition
\texttt{\small lambda}. Once the table is supplied, formulas (T12.4) and
(T13.6) use only counts of integer points in explicitly stated bounded
polytopes.

\subsection{T-014. Stable nonvanishing in an explicit family of every
width}\label{t-014.-stable-nonvanishing-in-an-explicit-family-of-every-width}

Let \texttt{\small lambda} range over the partitions with no part equal
to one, with arbitrary multiplicities of even parts, and with
multiplicity at most one for each odd part. Then

\[
\boxed{c(\lambda)\ge1.}
\tag{T14.1}
\]

The empty partition has
\texttt{\small c(\allowbreak{}empty)\allowbreak{}=\allowbreak{}1}, the
constant term of the completed character, and is included separately
from the nonempty case below.

Consequently, for the exact parameter range in (T10.1),

\[
a_{n,m}^{\lambda[mn]}\ge1.
\tag{T14.2}
\]

\textbf{Proof for nonempty \texttt{\small lambda}.} Put
\texttt{\small k=\allowbreak{}lambda\_\allowbreak{}1} and
\texttt{\small u\_\allowbreak{}i=\allowbreak{}m\_\allowbreak{}i(\allowbreak{}lambda)\allowbreak{}}.
For even \texttt{\small i}, the pure power of the highest wedge vector
\texttt{\small (\allowbreak{}e\_\allowbreak{}1 wedge ... wedge e\_\allowbreak{}i)\allowbreak{}\textasciicircum{}\{u\_\allowbreak{}i\}}
is a nonzero highest vector of weight
\texttt{\small u\_\allowbreak{}i omega\_\allowbreak{}i} in
\texttt{\small Sym\textasciicircum{}\{u\_\allowbreak{}i\}(\allowbreak{}Lambda\textasciicircum{}i V)\allowbreak{}}.
For odd \texttt{\small i}, the permitted values
\texttt{\small u\_\allowbreak{}i=\allowbreak{}0,\allowbreak{}1} select
respectively the scalar vector or the highest wedge vector in
\texttt{\small Lambda\textasciicircum{}\{u\_\allowbreak{}i\}(\allowbreak{}Lambda\textasciicircum{}i V)\allowbreak{}}.
Their tensor product is nonzero and is killed by all positive-root
operators. Its highest weight is

\[
\sum_{i=2}^k u_i\omega_i=\lambda',
\]

because the \texttt{\small j}th coordinate equals
\texttt{\small sum\_\allowbreak{}(\allowbreak{}i>=\allowbreak{}j)\allowbreak{}u\_\allowbreak{}i},
the height of the \texttt{\small j}th column of \texttt{\small lambda}.
It lies in scalar degree
\texttt{\small sum i u\_\allowbreak{}i=\allowbreak{}|lambda|}. Complete
reducibility in characteristic zero gives a constituent
\texttt{\small S\_\allowbreak{}lambda' V} of
\texttt{\small M\_\allowbreak{}k}. Equation (T10.3) proves (T14.1), and
(T10.1) transports it to (T14.2).

This explicitly includes all partitions with even parts and, for
example,
\texttt{\small (\allowbreak{}6\textasciicircum{}u)\allowbreak{}},
\texttt{\small (\allowbreak{}4\textasciicircum{}a,\allowbreak{}3,\allowbreak{}2\textasciicircum{}b)\allowbreak{}},
and
\texttt{\small (\allowbreak{}8\textasciicircum{}a,\allowbreak{}7,\allowbreak{}6\textasciicircum{}b,\allowbreak{}5,\allowbreak{}4\textasciicircum{}c,\allowbreak{}3,\allowbreak{}2\textasciicircum{}d)\allowbreak{}}.
It is a sufficient support family, with no claim that it exhausts the
stable support.

\subsection{Subsequent positive counts and the smaller direct-formula
problem}\label{subsequent-positive-counts-and-the-smaller-direct-formula-problem}

The stable coefficients of tails of width at most six have the small
explicit polytope formulas T-012 and T-013; T-014 proves support in a
specified family of every width. Formula (T10.3) remains exact at
arbitrary width. Subsequent statements T-016--T-020 in
\texttt{\small work/\allowbreak{}theta/\allowbreak{}width7/\allowbreak{}width7\_\allowbreak{}report.md}
give an exact three-variable positive polytope for the constituents of
shape
\texttt{\small (\allowbreak{}A,\allowbreak{}B\textasciicircum{}5,\allowbreak{}C)\allowbreak{}}
in
\texttt{\small Sym\textasciicircum{}d(\allowbreak{}Lambda\textasciicircum{}4 C\textasciicircum{}7)\allowbreak{}},
as well as further stable nonvanishing families and a precise
boundary-valuation formula for other weights.

The general existence question has also advanced: G-009--G-010 in
\texttt{\small work/\allowbreak{}root/\allowbreak{}invariant\_\allowbreak{}polytope\_\allowbreak{}report.md}
and
\texttt{\small work/\allowbreak{}root/\allowbreak{}general\_\allowbreak{}polytope\_\allowbreak{}extensions.md}
give terminating finite-presentation and standard-monomial polytope
constructions for all symmetric and general Schur plethysm coefficients,
including the unstabilized coefficients. G-011 does the same for the
actual canonical-map image, kernel, and cokernel modules. G-012 supplies
four direct polytopes for the full outer-degree-three family.
T-024--T-028 in
\texttt{\small work/\allowbreak{}theta/\allowbreak{}end\_\allowbreak{}node/\allowbreak{}}
further give a uniform polynomial-size polytope and an arithmetic
support test for every outer Schur functor with inner symmetric power
and output shape
\texttt{\small (\allowbreak{}alpha,\allowbreak{}beta\textasciicircum{}(\allowbreak{}m-\allowbreak{}2)\allowbreak{},\allowbreak{}gamma)\allowbreak{}},
with a direct toric algebra degeneration for symmetric outer powers.
Thus the current narrower task is to replace large parameter-dependent
Reynolds enumeration and elimination by small, efficiently specified
presentations or polyhedral formulas outside these direct families. The
width-seven valuation outside T-018 is one concrete calculation for that
purpose; it is not a remaining obstruction to the existence of a
terminating positive count.

\section{Width seven: metric covariants, exact end-node coefficients,
and the extension
valuation}\label{width-seven-metric-covariants-exact-end-node-coefficients-and-the-extension-valuation}

This report continues T-010--T-014. It supplies a complete positive
formula for a three-parameter family inside the first remaining
symmetric-power factor, together with further stable nonvanishing
families. The constructions and proofs are unconditional. No global
novelty claim is made.

All representations are over \(\mathbf C\). Put \(V=\mathbf C^7\),
\(G=GL(V)\), \[
R=\operatorname{Sym}(\Lambda^3V)=\mathbf C[\Lambda^3V^*],
\qquad S=\operatorname{Sym}(\Lambda^4V).
\] The coordinate-ring action is \((gf)(\phi)=f(g^{-1}\phi)\). Thus
\(R_d\) has positive central degree \(3d\). Highest weights are
decreasing integral seven-tuples; a superscript on an entry means
repetition.

For integers \(A\ge B\ge C\ge0\), \(d\ge0\), \(A+5B+C=4d\), define \[
\mathcal P(A,B,C)=\left\{(k,\ell,j)\in\mathbf R_{\ge0}^3:
\begin{array}{rcl}
2k+4\ell+4j&=&A,\\
2k+2\ell+4j&=&B,\\
2\ell+4j&=&C.
\end{array}\right\}.
\] The exact answer proved below is \[
\boxed{[S_{(A,B^5,C)}V:S_d]
=\#(\mathcal P(A,B,C)\cap\mathbf Z^3).}
\tag{W7.1}
\] This bounded polytope is empty or a point. Its unique possible point
is \[
k=(B-C)/2,\qquad \ell=(A-B)/2,\qquad j=(B+C-A)/4.
\tag{W7.2}
\]

\subsection{T-016. Cubic metric and polynomial quartic
cometric}\label{t-016.-cubic-metric-and-polynomial-quartic-cometric}

\subsubsection{Precisely located
background}\label{precisely-located-background}

The following classical input is used.

\begin{enumerate}
\def\labelenumi{\arabic{enumi}.}
\item
  A generic alternating three-form on \(V\) has stabilizer
  \(H=\mu_3\times G_2\subset GL(V)\), with \(G_2\subset SL(V)\) acting
  through its irreducible seven-dimensional representation. Its orbit
  \(O\) is open. The finite central factor is retained. An exact primary
  locator is Dave Anderson, \emph{Chern class formulas for \(G_2\)
  Schubert loci}, Appendix A, Proposition A.1 and Corollary A.2, printed
  page 29:
  \url{https://people.math.osu.edu/anderson.2804/papers/g2-chern.pdf}
\item
  The complement of \(O\) is one irreducible hypersurface with
  degree-seven equation. The cubic metric below is nondegenerate exactly
  on \(O\), and has rank four on the dense codimension-one boundary
  orbit. Laurent Manivel, \emph{Prehomogeneous spaces and projective
  geometry}, Example 4.3.3, printed pages 59--60, and Example 4.4.2,
  printed pages 61--62, give the orbit representatives, rank data,
  boundary diagram, and degree-seven invariant. Proposition 3.1.4,
  printed page 44, identifies fundamental relative invariants with
  irreducible hypersurface equations. The metric also appears in §5.1.2,
  printed page 68:
  \url{https://manivel.perso.math.cnrs.fr/trento2012.pdf}
\end{enumerate}

Here ``metric rank'' always means the rank of the explicitly defined
seven-by-seven matrix; the source's two rank labels need not be adopted.
The classification and generic-stabilizer statements are cited
background, not reproved as new results.

Choose a volume vector \(\tau\in\det V\), with dual volume form
\(\Omega\). Define \[
B_\phi(u,v)\Omega=(\iota_u\phi)\wedge(\iota_v\phi)\wedge\phi.
\tag{W7.3}
\] The contracted forms have degree two, so this is symmetric in
\(u,v\), and its coefficients are cubic in \(\phi\). Intrinsically its
values lie in \(\operatorname{Sym}^2V^*\otimes\det V^*\). The actual
coefficient morphism is \[
F_B:\operatorname{Sym}^2V\otimes\det V\longrightarrow R_3,\qquad
F_B(uv\otimes\tau)(\phi)
=\langle(\iota_u\phi)\wedge(\iota_v\phi)\wedge\phi,\tau\rangle.
\tag{W7.4}
\] Equivariance follows by pulling back the forms and moving \(g\)
across the evaluation pairing. This fixes the determinant twist. For
\(\phi=e^{123}+e^{145}+e^{167}\), evaluation at \(u=v=e_1\) gives
\(6\Omega\), so the map is nonzero. Its source is irreducible, hence it
is injective. In particular \(b=F_B(e_1^2\otimes\tau)\ne0\) is a highest
vector of weight \((3,1^6)\).

Let \(\Delta\) be an irreducible equation of the boundary. Its degree is
seven and its character for the stated action is \[
\Delta\in R_7,\qquad g\Delta=(\det g)^3\Delta.
\tag{W7.5}
\] Indeed every character is a determinant power, and scalar matrices
act on \(R_7\) by their twenty-first power.

Writing \(B=(B_{ij})\), its coefficient transformation is \[
gB=(\det g)\,g^{\mathsf T}Bg.
\] Thus \(\det B\) has degree twenty-one and character \(\det^9\). Its
zero locus is exactly the boundary; unique factorization gives
\(\det B=c_0\Delta^3\), \(c_0\ne0\). Rescale \(\Delta\) over
\(\mathbf C\) so that \(c_0=1\).

Every entry of \(\operatorname{adj}B\) is divisible by \(\Delta^2\). To
prove this, localize the polynomial ring \(R\) at the height-one prime
\((\Delta)\), obtaining a discrete valuation ring. Modulo its maximal
ideal the metric has rank four. Invertible row and column operations
reduce its first four rows and columns to an invertible block; the
entries of the remaining three-by-three Schur-complement block are
divisible by \(\Delta\). Every six-by-six minor therefore has valuation
at least two. Since the minors are polynomials and \(R\) is a unique
factorization domain, this is divisibility by \(\Delta^2\) in \(R\).

Consequently \[
C=\operatorname{adj}B/\Delta^2
\tag{W7.6}
\] is a matrix of quartic polynomials and \[
BC=CB=\Delta I_7.
\tag{W7.7}
\] The adjugate transformation gives \[
gC=(\det g)^2g^{-1}Cg^{-\mathsf T}.
\] Its coefficient space is the image of the equivariant morphism \[
F_C:\operatorname{Sym}^2V^*\otimes\det^2V\longrightarrow R_4.
\tag{W7.8}
\] The matrix is invertible on \(O\), so this map is nonzero; its
irreducible source makes it injective. Choose a nonzero highest vector
\(c\) in its image. Its weight is \((2^6,0)\).

Neither \(b\) nor \(c\) is divisible by the irreducible \(\Delta\),
since their degrees are three and four. Hence \[
v_\Delta(b^kc^\ell\Delta^j)=j
\quad(k,\ell\ge0,\ j\in\mathbf Z).
\tag{W7.9}
\] This is the exact extension test.

The symbolic receipt METRIC\_PATH\_CHECK.txt uses
\(\phi(t)=e^{123}+e^{145}+t e^{167}+e^{246}+e^{357}\). It obtains metric
ranks four at \(t=0\) and seven at \(t=1\), determinant \(-4374t^3\),
and adjugate order at least two, attained by an entry. This finite check
supplements the global divisibility proof.

\subsection{\texorpdfstring{T-017. The required \(G_2\)-invariant
spaces}{T-017. The required G\_2-invariant spaces}}\label{t-017.-the-required-g_2-invariant-spaces}

Write \(\omega_1\) for the highest weight of the seven-dimensional
\(G_2\)-module. Then \[
\operatorname{Sym}^mV|_{G_2}
=\bigoplus_{r=0}^{\lfloor m/2\rfloor}V_{(m-2r)\omega_1}
\quad(m\ge0).
\tag{W7.10}
\] Here is a proof of precisely this assertion. The nondegenerate
invariant quadratic form gives the usual contraction
\(\operatorname{Sym}^mV\to\operatorname{Sym}^{m-2}V\), and the harmonic
decomposition
\(\operatorname{Sym}^mV=H_m\oplus q\operatorname{Sym}^{m-2}V\) in
characteristic zero. A highest vector \(v\) is isotropic because \(v^2\)
has nonzero weight \(2\omega_1\), so \(v^m\in H_m\). It generates
\(V_{m\omega_1}\). The Weyl formula for the six positive roots of
\(G_2\) gives \[
\dim V_{m\omega_1}
=\frac{(m+1)(m+2)(m+3)(m+4)(2m+5)}{120}
=\binom{m+6}{6}-\binom{m+4}{6}
=\dim H_m.
\] Binomial terms with top below six are zero. Thus
\(H_m=V_{m\omega_1}\), and induction proves (W7.10). Complete
reducibility, harmonic decomposition, and the Weyl formula are the
standard background theorems used here.

For \(a,b\ge0\), let \(L_{a,b}\) be the irreducible \(SL_7\)-module with
highest weight \(a\omega_1+b\omega_6\), now using the \(SL_7\)
fundamental weights. Then \[
\boxed{\dim L_{a,b}^{G_2}
=\begin{cases}1,&a,b\text{ both even},\\0,&\text{otherwise}.\end{cases}}
\tag{W7.11}
\] To prove it, use the exact determinant identity \[
\operatorname{Sym}^aV\otimes\operatorname{Sym}^bV^*
=\operatorname{Sym}^aV\otimes S_{(b^6)}V\otimes\det^{-b}V.
\] Pieri's horizontal-strip rule, applied to the six-row rectangle,
gives \[
\operatorname{Sym}^aV\otimes\operatorname{Sym}^bV^*
=\bigoplus_{i=0}^{\min(a,b)}L_{a-i,b-i}
\tag{W7.12}
\] after restriction to \(SL_7\). This includes \(a=0\) or \(b=0\). By
(W7.10) and Schur's lemma the left side has \(G_2\)-invariant dimension
\[
H(a,b)=\#\{r:0\le r\le\min(a,b),\ r\equiv a\equiv b\pmod2\}.
\] Set \(H(a,b)=0\) if either argument is negative. Subtracting (W7.12)
at arguments \(a-1,b-1\) gives \(\dim L_{a,b}^{G_2}=H(a,b)-H(a-1,b-1)\).
Different parities give two zero terms; both odd give equal terms; both
even give difference one. The boundary cases follow from the zero
convention. This proves (W7.11).

The independent rectangular-cases checker verified the Pieri
calculation, the difference including all boundary cases, and the
harmonic dimension argument.

\subsection{T-018. Exact polynomial sector and positive polytope
formula}\label{t-018.-exact-polynomial-sector-and-positive-polytope-formula}

For \(a,b\ge0\) and \(s\in\mathbf Z\), set \[
\gamma=(a+b+s,(b+s)^5,s).
\tag{W7.13}
\] These are exactly the dominant integral seven-tuples with equal
second through sixth entries. They label rational
\(GL_7\)-representations. The central \(\mu_3\) acts by the character
determined by \(|\gamma|=a+6b+7s\).

Evaluation at the basepoint of \(O=G/H\) gives the actual isomorphism \[
\operatorname{Hom}_G(S_\gamma V,\mathbf C[O])
\simeq((S_\gamma V)^*)^H.
\tag{W7.14}
\] An \(H\)-invariant functional \(\ell\) corresponds explicitly to \[
v\longmapsto f_{\ell,v},\qquad
f_{\ell,v}(gH)=\ell(g^{-1}v).
\] Conversely evaluate an equivariant map at the identity coset. These
operations are inverse. Regularity is the matrix-coefficient
construction on \(G\) and its descent to the homogeneous space
\(O=G/H\).

By (W7.11), this space has dimension one exactly when \[
a=2k,\quad b=2\ell,\quad s=k+3j
\quad(k,\ell\ge0,\ j\in\mathbf Z),
\tag{W7.15}
\] and otherwise it has dimension zero. The congruence on \(s\) is
exactly the \(\mu_3\)-invariance condition; no determinant character has
been dropped.

For (W7.15), the nonzero rational highest vector \(b^kc^\ell\Delta^j\)
has the required weight and degree \[
d=3k+4\ell+7j,
\tag{W7.16}
\] \[
\gamma=(3k+2\ell+3j,(k+2\ell+3j)^5,k+3j).
\tag{W7.17}
\] It is regular on \(O\) and spans its one-dimensional highest-vector
space. It extends to \(R\) exactly when \(j\ge0\), by (W7.9). Indeed
\(\mathbf C[O]=R[\Delta^{-1}]\), and a fraction with only this possible
denominator is polynomial exactly when its \(\Delta\)-valuation is
nonnegative.

It follows that the subalgebra of \(R^U\) spanned by the highest weights
with equal second through sixth entries is exactly \[
\boxed{R_{\mathrm{ends}}=\mathbf C[b,c,\Delta].}
\tag{W7.18}
\] The three generators are algebraically independent: their weights in
(W7.17) determine the three exponents uniquely, so distinct monomials
have distinct torus weights. Their nonzero products are independent.
Thus (W7.18) identifies the algebra, its actual generators, and its
positive monomial basis.

Now use the degreewise exterior-duality isomorphism \[
\Lambda^4V\simeq\Lambda^3V^*\otimes\det V,
\qquad
S_d\simeq R_d^*\otimes\det^dV.
\tag{W7.19}
\] A constituent of highest weight \(\gamma\) goes to highest weight \[
\eta_i=d-\gamma_{8-i}\quad(1\le i\le7).
\tag{W7.20}
\] This identifies multiplicities with their determinant shifts.
Equations (W7.16)--(W7.17) give \[
\eta=(2k+4\ell+4j,(2k+2\ell+4j)^5,2\ell+4j).
\tag{W7.21}
\] Conversely \((A,B^5,C)\) determines precisely the parameters in
(W7.2). If parity or the central congruence fails, (W7.14)--(W7.15) give
zero; if an extension parameter is negative, the polynomial constituent
is absent. If all three are nonnegative integers, (W7.17) is a
polynomial constituent with multiplicity one, and (W7.19) supplies the
target. This proves (W7.1), including all zero coefficients.

The slice \(\ell=0\) says \[
[S_{(A^6,C)}V:S_d]
=\begin{cases}1,&A\text{ even},\\0,&A\text{ odd},\end{cases}
\qquad 6A+C=4d,\quad A\ge C\ge0.
\] The degree condition forces \(C\equiv0\pmod4\) in the nonzero case.
The slice \(k=0\) is also governed by the same explicit polytope,
including the triangle condition \(A\le2B\).

\subsubsection{A precise obstruction to discarding the extension
valuation}\label{a-precise-obstruction-to-discarding-the-extension-valuation}

For \(k=3,\ell=2,j=-1\), (W7.17) gives \[
d=10,\qquad \gamma=(10,4^5,0).
\] The function \(b^3c^2/\Delta\) is regular on \(G/H\), has this
polynomial dominant weight, and has a genuine simple pole on the
boundary. Therefore \[
\dim((S_{(10,4^5,0)}V)^*)^H=1,\qquad
[S_{(10,4^5,0)}V:R_{10}]=0.
\tag{W7.22}
\] Every entry of \(\gamma\) lies between zero and ten and
\(|\gamma|=30\), so it also satisfies the elementary torus
weight-support bounds for \(R_{10}\). The failure of this transfer is
proved by the localization morphism and its valuation.

Duality gives the corresponding zero \([S_{(10,6^5,0)}V:S_{10}]=0\). The
possible polytope point is \((3,2,-1)\), displaying the failed
inequality.

The exact verifier check\_width7.py computes the character of \(R_d\) by
Newton's recurrence and extracts Schur multiplicities by the Weyl
alternant. Every end-node partition in degrees zero through
ten---seventy-five partitions, including (W7.22)---agrees with (W7.18).
END\_NODE\_CHECK.tsv is the complete receipt. This is a finite
verification supplementary to the infinite proof.

\subsection{T-019. Stable positive
families}\label{t-019.-stable-positive-families}

T-010 realizes the stable coefficient
\(c_\lambda=\langle s_\lambda,H[H-1-h_1]\rangle\), for width at most
seven, as the multiplicity of \(S_{\lambda'}V\) in scalar degree
\(|\lambda|\) of \[
\mathcal M_7=
\operatorname{Sym}(\Lambda^2V)\otimes
\Lambda(\Lambda^3V)\otimes
\operatorname{Sym}(\Lambda^4V)\otimes
\Lambda(\Lambda^5V)\otimes
\operatorname{Sym}(\Lambda^6V)\otimes
\Lambda(\Lambda^7V).
\] Taking constants in every factor other than
\(\operatorname{Sym}(\Lambda^4V)\), (W7.21) yields \[
\boxed{c_{(7^{\,2\ell+4j},\,6^{\,2k},\,1^{\,2\ell})}\ge1
\quad(k,\ell,j\ge0).}
\tag{W7.23}
\] The empty partition has coefficient one. This is a stable
multiplicity bound; (W7.1) computes exactly the contribution of the
specified symmetric-power factor.

Differentiating \(\Delta\) produces the nonzero equivariant map \[
\Lambda^3V^*\otimes\mathbf C\Delta\longrightarrow R_6.
\tag{W7.24}
\] It is nonzero by Euler's identity
\(7\Delta=\sum x_i\partial_i\Delta\); its irreducible source makes it
injective. Its highest weight is \((3^4,2^3)\). Via (W7.19) it gives a
nonzero highest vector \(g\in S_6\) of weight \((4^3,3^4)\). Hence \[
c_{(7^3,3)}\ge1,\qquad
c_{(7^{\,4j+3h},3^h)}\ge1\quad(j,h\ge0).
\tag{W7.25}
\] For the second assertion choose a determinant vector in \(S_7\), of
weight \((4^7)\), supplied from \(\Delta\) by (W7.19), and multiply its
\(j\)-th power by \(g^h\) inside the polynomial domain \(S\). Degreewise
duality is used to produce individual constituents; all multiplication
occurs in the stated target algebra.

A more flexible sufficient family is \[
\boxed{\begin{split}
\lambda={}&(7^{\,2\ell+4j+3h+\epsilon_7},
6^{\,u_6},5^{\,\epsilon_5},4^{\,u_4},
3^{\,h+\epsilon_3},2^{\,u_2},1^{\,2\ell}),\\
&c_\lambda\ge1,
\end{split}}
\tag{W7.26}
\] for \(\ell,j,h,u_2,u_4,u_6\ge0\) and
\(\epsilon_3,\epsilon_5,\epsilon_7\in\{0,1\}\).

For a complete proof, inside the single domain \(S\) multiply the
cometric-derived highest vector of weight \((4,2^6)\) to power \(\ell\),
the determinant vector to power \(j\), the differentiated vector \(g\)
to power \(h\), and the pure highest vector
\(e_1\wedge e_2\wedge e_3\wedge e_4\) to power \(u_4\). The result is
nonzero and highest. In the symmetric factors for exterior degrees two
and six, take pure highest vectors to powers \(u_2,u_6\). In each of the
exterior-algebra factors of exterior degree three, five, and seven, take
either its scalar or its degree-one highest vector according to
\(\epsilon_i\). The tensor is nonzero, highest of weight \(\lambda'\),
and has scalar degree \(|\lambda|\). This proves (W7.26). No power
greater than one of a single vector is used in an exterior algebra.

Every stable conclusion transports to the finite coefficient
\(a_{n,m}^{\lambda[mn]}\) in T-010's exact range \[
m\ge\lambda_1,\qquad
n\ge\left\lfloor\frac{|\lambda|+\lambda_1}{2}\right\rfloor,
\qquad mn-|\lambda|\ge\lambda_1.
\] These bounds are Sarkar, arXiv:2610.07929v2, Theorem 6.4. This report
supplies the stated constituents of his stable module.

\subsection{T-020. The remaining data: a finite-dimensional valuation
filtration}\label{t-020.-the-remaining-data-a-finite-dimensional-valuation-filtration}

For each dominant integral weight \(\gamma\), fix a highest vector
\(v_\gamma\) of the rational irreducible \(L_\gamma\) and define \[
E_\gamma=(L_\gamma^*)^H,\qquad
F_{\gamma,\ge q}=
\{\ell\in E_\gamma:
v_\Delta(gH\mapsto\ell(g^{-1}v_\gamma))\ge q\}.
\tag{W7.27}
\] The zero function has infinite valuation. This is a vector-space
filtration, since the valuation of a sum is at least the smaller summand
valuation. The exact equality is \[
[L_\gamma:R]=\dim F_{\gamma,\ge0}.
\tag{W7.28}
\] The central character determines its degree whenever it occurs.
Equation (W7.28) follows from (W7.14) and
\(R=\{f\in R[\Delta^{-1}]:v_\Delta(f)\ge0\}\). If the highest matrix
coefficient is in \(R\), its entire irreducible span is in \(R\),
because \(R\) is \(G\)-stable. Thus testing this coefficient proves the
assertion for the full equivariant map, not only one function.

Multiplication by \(\Delta^t\) identifies the open-orbit highest spaces
at \(\gamma\) and \(\gamma+(3t)^7\), and increases their valuations by
\(t\). Thus the determinant-shift coefficients are precisely the
dimensions of steps in this filtration. T-018 calculates the entire
filtration for the end-node sector. Elsewhere, computing only the
\(G_2\times\mu_3\)-invariant dimensions leaves the extension condition
unevaluated. Example (W7.22) proves that replacing the filtered
dimension by the full invariant dimension can change the coefficient.

There is a precise published limitation on copying the rank-six
free-generator formula to all weights. Michel Brion, \emph{Invariants
d'un sous-groupe unipotent maximal d'un groupe semi-simple}, Annales de
l'Institut Fourier 33(1) (1983), Theorem 3, printed pages 12--13, and
case \((A_6,\omega_3)\) of §III.1, printed page 14, prove that the full
algebra \(R^U\) is not polynomial:
\url{https://www.numdam.org/item/10.5802/aif.902.pdf}

The cited case gives a concrete contradiction. Its quotient dimension is
at least \(35-21=14\). There is one generator in degree one, one
additional generator in degree two, and further generators have degree
at least three, using
\(\operatorname{Sym}^2(\Lambda^3V)=S_{(2^3)}V\oplus S_{(2,1^4)}V\).
Polynomiality would force the sum of generator degrees to be at least
\(1+2+12\cdot3=39\). Brion's Corollary 2 to Theorem 2 bounds this sum by
\(35\) in the polynomial case. This excludes that particular
free-polynomial-algebra formula for all weights. It permits finite
presentations, standard monomial bases, and positive polytope sums.
Formula (W7.18) supplies a polynomial sector with its exact multiplicity
interpretation.

A concrete next orbit-based computation is the filtration (W7.27)
outside this sector. Rota and Stein's nine symbolic covariants, printed
page 847 of \emph{Symbolic method in invariant theory}, PNAS 83 (1986),
844--847, give explicit tensor-valued candidates:
\url{https://pmc.ncbi.nlm.nih.gov/articles/PMC322966/} Their tensor
operations require conversion to a presentation or valuation basis
before serving as a basis of all \(R^U\). No assertion that these nine
diagrams are nine algebraically independent generators of \(R^U\) is
used.

\subsection{Verification status and
artifacts}\label{verification-status-and-artifacts}

The complete arguments T-016--T-020 were reread with attention to
coordinate-ring duals, the finite stabilizer, scalar degrees, both
determinant twists, adjugate divisibility, all zero-boundary cases in
(W7.11), and the use of degreewise duality versus multiplication in the
target algebra. No gap was found in the stated conclusions.

The orbit classification and generic-stabilizer identification are cited
established inputs. Their statements and relevant source sections were
read; their classification proofs were not independently reproduced. The
complete extension and coefficient arguments used here are supplied
above. The finite checker uses exact integers and symbolic one-variable
matrix calculations; it is not a proof-assistant certificate.

Accompanying files:

\begin{itemize}
\tightlist
\item
  check\_width7.py: Newton/Weyl verification and cubic-metric path
  calculation.
\item
  END\_NODE\_CHECK.tsv: all seventy-five end-node partitions through
  symmetric degree ten.
\item
  METRIC\_PATH\_CHECK.txt: exact metric matrix, determinant, and
  adjugate orders.
\item
  The parent CHECK\_EVENTS.tsv records actual proof-reading and
  computation events.
\end{itemize}

One initial symbolic-check run included an unjustified expected
normalization scalar for \(\det B\). The computed scalar was different;
the check was corrected to test the claimed invariant statement, namely
nonzero determinant order exactly three. All final tests passed. No
plethysm coefficient prediction failed.

\section{An explicit toric presentation for the two-row sector of outer
degree
four}\label{an-explicit-toric-presentation-for-the-two-row-sector-of-outer-degree-four}

\subsection{Scope and prior result}\label{scope-and-prior-result}

This report gives a small presentation and a finite SAGBI basis for the
subalgebra of G-009's \(A_4\) with output partitions of length at most
two. It yields two fixed, four-variable polytopes for all coefficients
\(a_{4,n}^{(4n-k,k)}\). The binary-quartic covariant algebra and this
coefficient family are classical; the point here is the complete
explicit normalizer-algebra morphism, the finite toric degeneration, and
the directly verified lattice count. No discovery of the classical
presentation or coefficient family is claimed.

Primary historical source: J. H. Grace and A. Young, \emph{The Algebra
of Invariants} (1903), §89, printed pp.~91--94, proves generation of the
quartic covariants by five forms; §93, printed pp.~97--98, derives their
syzygy. The source text was read in this run:
\url{https://archive.org/download/algebraofinvaria00graciala/algebraofinvaria00graciala_djvu.txt}
A scan is available at:
\url{https://commons.wikimedia.org/wiki/File:The_algebra_of_invariants_(IA_algebraofinvaria00graciala).pdf}
The coordinate normalization below is checked directly and does not
require translating the historical symbolic normalization.

\subsection{T-021. Exact normalizer-algebra
morphism}\label{t-021.-exact-normalizer-algebra-morphism}

Let \(F=\mathbf Q^2\), with basis \(e_1,e_2\), and let \[
P=\mathbf Q[a_1,b_1,\ldots,a_4,b_4]
=\operatorname{Sym}(\mathbf Q^4\otimes F).
\] The determinant-one diagonal torus on the first factor fixes
precisely the equal row-degree pieces \[
B=\bigoplus_{n\ge0}P_{(n,n,n,n)}
=\bigoplus_{n\ge0}(\operatorname{Sym}^nF)^{\otimes4}.
\] Literal row permutations act by \(S_4\). The upper unipotent group of
\(GL(F)\) acts infinitesimally by \[
D_{\rm rows}=\sum_{r=1}^4 a_r\frac{\partial}{\partial b_r}.
\] The desired two-row part of the projected flag algebra is \[
A_4^{\le2}=(B^{S_4})^{U_2}.
\tag{Q4.1}
\] This is the actual subalgebra of \(A_4\) having right highest weights
\((\lambda_1,\lambda_2,0,0)\). A polynomial of that right-torus weight
uses only the first two columns. All positive-root operators involving
later columns annihilate it; the remaining operator is \(D_{\rm rows}\).
This proves the precise identification with the flag construction.

Put
\(Q=\mathbf Q[x_0,x_1,x_2,x_3,x_4]=\operatorname{Sym}(\operatorname{Sym}^4F)\),
where \(x_i=e_1^{4-i}e_2^i\). Define \[
\Phi(x_i)=\binom4i^{-1}
\sum_{\substack{I\subseteq\{1,2,3,4\}\\|I|=i}}
\prod_{r\in I}b_r\prod_{r\notin I}a_r.
\tag{Q4.2}
\] This is the pullback of the multiplication map on four linear forms,
using the perfect symmetric-power pairing. It is a \(GL(F)\)-equivariant
graded algebra morphism \(Q\to B^{S_4}\). In particular, \[
D_{\rm rows}\Phi(x_i)=i\Phi(x_{i-1}),
\] as is also verified directly: each monomial with \(i-1\) selected
\(b\)'s is obtained \(5-i\) times before dividing by \(\binom4i\).

Every complex binary quartic factors into four linear forms. Therefore
the multiplication map is surjective over \(\mathbf C\), so its pullback
\(\Phi\) is injective, also over \(\mathbf Q\). In degree \(n\), its
source and target both have dimension \[
\dim\operatorname{Sym}^n(\operatorname{Sym}^4F)
=\binom{n+4}{4}
=\dim\operatorname{Sym}^4(\operatorname{Sym}^nF).
\] Consequently \(\Phi\) is surjective degree by degree and is an
algebra isomorphism. This is the explicit binary canonical Foulkes map
\(h_{n,4}\) in each degree, and it is multiplicative across degrees.

Let \[
D=\sum_{i=1}^4 i x_{i-1}\frac{\partial}{\partial x_i}.
\] Since characteristic zero identifies invariance under the additive
one-parameter group with the kernel of its locally nilpotent derivation,
(Q4.1)--(Q4.2) give \[
\boxed{A_4^{\le2}\simeq\ker(D:Q\to Q).}
\tag{Q4.3}
\] This proves the morphism relating the binary-quartic calculation to
the normalizer projection, including its multiplication.

\subsection{T-022. Complete presentation, positive basis, and SAGBI
degeneration}\label{t-022.-complete-presentation-positive-basis-and-sagbi-degeneration}

In \(Q\), define \[
\begin{aligned}
f&=x_0,\\
H&=x_0x_2-x_1^2,\\
T&=x_0^2x_3-3x_0x_1x_2+2x_1^3,\\
I&=x_0x_4-4x_1x_3+3x_2^2,\\
J&=x_0x_2x_4+2x_1x_2x_3-x_0x_3^2-x_1^2x_4-x_2^3.
\end{aligned}
\tag{Q4.4}
\] Direct differentiation proves that each lies in \(\ker D\). Direct
expansion gives \[
T^2+4H^3-f^2HI+f^3J=0.
\tag{Q4.5}
\] The exact presentation is \[
\boxed{\ker D=\mathbf Q[f,H,I,J,T]/(T^2+4H^3-f^2HI+f^3J).}
\tag{Q4.6}
\]

Here is a full proof of generation and absence of further relations,
formulated to identify the possible denominator obstruction explicitly.
After inverting \(f=x_0\), translate by the group parameter \(-x_1/f\).
The resulting coordinates have first coordinate \(f\), second coordinate
zero, and the other three coordinates \[
y_2=H/f,\qquad y_3=T/f^2,\qquad
y_4=(f^2I-3H^2)/f^3.
\] The inverse translation recovers the original coordinates. Thus \[
Q[f^{-1}]
=\mathbf Q[f^{\pm1},x_1/f,y_2,y_3,y_4],
\] and \(D(x_1/f)=1\), while \(D\) annihilates the other four
generators. Hence \[
(\ker D)[f^{-1}]
=\mathbf Q[f^{\pm1},H,T,I].
\tag{Q4.7}
\]

Let \(C_0\) be the abstract right side of (Q4.6), and map it to \(Q\) by
(Q4.4). Its defining equation is monic in \(T\), so it is free over the
abstract polynomial ring \(\mathbf Q[f,H,I,J]\) with basis \(1,T\); in
particular multiplication by \(f\) is injective. After inverting \(f\),
the equation solves for \(J\), and (Q4.7) shows that the map becomes an
isomorphism to the displayed polynomial subalgebra. Thus \(C_0\to Q\) is
injective before localization as well.

Modulo \(f\), this map is \[
\mathbf Q[H,T,I,J]/(T^2+4H^3)
\longrightarrow \mathbf Q[x_1,x_2,x_3,x_4],
\] \[
H\mapsto-x_1^2,\quad T\mapsto2x_1^3,\quad
I\mapsto-4x_1x_3+3x_2^2,\quad
J\mapsto2x_1x_2x_3-x_1^2x_4-x_2^3.
\tag{Q4.8}
\] This map is injective. The first two images give exactly the cusp
algebra \(\mathbf Q[x_1^2,x_1^3]\): reduction to \(H^aT^\epsilon\),
\(\epsilon=0,1\), yields distinct powers of \(x_1\). Over
\(\mathbf Q(x_1,x_2)\), the last two images are affine linear
coordinates in \(x_3,x_4\), with nonzero triangular coefficients
\(-4x_1\) and \(-x_1^2\). They are algebraically independent over
\(\mathbf Q(x_1)\), and hence over the cusp algebra. This proves the
asserted injectivity.

It follows that \(C_0\cap fQ=fC_0\). Every \(g\in\ker D\) belongs to
\(C_0[f^{-1}]\) by (Q4.7). Choose a smallest \(r\ge0\) for which
\(f^rg\in C_0\). If \(r>0\), then \(f^rg\in C_0\cap fQ=fC_0\),
contradicting minimality. Thus \(r=0\), proving (Q4.6).

In particular the monomials \[
f^aH^bI^cJ^dT^\epsilon,\qquad
a,b,c,d\ge0,\quad\epsilon\in\{0,1\},
\tag{Q4.9}
\] form a basis.

Now put positive ambient weights \((1,2,1,2,2)\) on
\((x_0,x_1,x_2,x_3,x_4)\), and refine the weighted order by
lexicographic order. The leading terms of \((f,H,T,I,J)\) are \[
x_0,\quad -x_1^2,\quad 2x_1^3,\quad -4x_1x_3,\quad -x_1^2x_4.
\tag{Q4.10}
\] The leading monomial of (Q4.9) is \[
x_0^a x_1^{\,2b+c+2d+3\epsilon}x_3^cx_4^d.
\] All these monomials are distinct: recover \(a,c,d\) from their
coordinates, recover \(\epsilon\) by parity of the remaining \(x_1\)
exponent, then recover \(b\). Therefore every nonzero element of \(C_0\)
has leading monomial in the algebra generated by (Q4.10). Conversely the
leading generator monomials belong to the initial algebra. This proves,
without an unverified subduction assumption, that the five generators
are a SAGBI basis and \[
\boxed{\operatorname{in}(A_4^{\le2})
\simeq\mathbf Q[x_0,x_1^2,x_1^3,x_1x_3,x_1^2x_4].}
\tag{Q4.11}
\] This is a toric algebra; normality is not asserted. Its only toric
relation, with the coefficient normalization in (Q4.10), is
\(T^2+4H^3=0\). The unique-exponent argument proves that there are no
further relations. The ambient weight filtration, transported along
(Q4.3), supplies the toric degeneration of the specified normalizer
subalgebra.

The generator degrees \((n;\lambda_1,\lambda_2)\) are \[
\begin{array}{c|ccccc}
 &f&H&I&J&T\\ \hline
n&1&2&2&3&3\\
\lambda&(4,0)&(6,2)&(4,4)&(6,6)&(9,3).
\end{array}
\] Thus, for \(0\le k\le2n\), \[
\boxed{
a_{4,n}^{(4n-k,k)}
=\sum_{\epsilon=0}^{1}
\#\left\{(a,b,c,d)\in\mathbf Z_{\ge0}^{4}:
\begin{array}{rcl}
a+2b+2c+3d&=&n-3\epsilon,\\
2b+4c+6d&=&k-3\epsilon
\end{array}\right\}.
}
\tag{Q4.12}
\] These are two bounded rational polytopes of fixed dimension,
specified by a constant matrix and the two input integers. The first
equation bounds every coordinate. Unlike the general construction, no
parameter-dependent Reynolds enumeration, Hilbert-basis calculation, or
polynomial elimination is required to produce these polytopes. They
recover a known coefficient family through the explicit normalizer
algebra and its toric filtration.

\subsection{T-023. A precise obstruction for the ordinary lexicographic
candidate}\label{t-023.-a-precise-obstruction-for-the-ordinary-lexicographic-candidate}

The same five generators do not form a SAGBI basis for ordinary
lexicographic order \(x_0>x_1>x_2>x_3>x_4\). Their leading monomials are
\[
x_0,\quad x_0x_2,\quad x_0^2x_3,\quad x_0x_4,\quad x_0x_2x_4.
\] But \(fJ-HI\) belongs to their algebra and has leading term
\(-x_0^2x_3^2\). A nonnegative product of the listed leading monomials
with zero \(x_2,x_4\) exponents can use only \(x_0\) and \(x_0^2x_3\).
Obtaining \(x_3^2\) then forces \(x_0\) exponent at least four, rather
than two. Thus the required leading monomial is absent. This proves the
failure of that particular order; (Q4.10) gives an explicit order for
which the same generators succeed.

\subsection{Exact checks and next
sector}\label{exact-checks-and-next-sector}

The script check\_binary\_quartic.py checks every derivation identity,
the defining relation, the five weighted leading terms, and the lex
obstruction by exact symbolic arithmetic. It also compares (Q4.12) with
independent Gaussian-rectangle weight differences for all
\(0\le n\le100\), \(0\le k\le2n\): 10,201 coefficients. Every value
agrees. BINARY\_QUARTIC\_CHECK.tsv and SYMBOLIC\_CHECK.txt provide the
receipts. These are finite checks supplementary to the complete proof.

The full \(A_4\), and already its output-length-three subalgebra,
contain additional highest weights and are not evaluated by (Q4.3). The
displayed morphism is exact for the stated two-row sector. The next
concrete task is an explicit presentation or filtration for the
length-three sector, using the three-column flag variables rather than
extending the binary map without a proved morphism.

\section{A uniform positive formula and toric degeneration for equal
middle
rows}\label{a-uniform-positive-formula-and-toric-degeneration-for-equal-middle-rows}

Run:
\texttt{\small GCT\_\allowbreak{}PLETH-\allowbreak{}20261121-\allowbreak{}theta}.
Author/checker: GPT-6 Astra Pro; reasoning-effort setting is not
surfaced by the runtime. All coefficient statements are in
characteristic zero. Every algebra map and parity character used below
is specified. These are direct constructions and deductions; no claim of
historical novelty is made.

\subsection{Result and scope}\label{result-and-scope}

Write \[
a_{m,n}^{\lambda}=[S_\lambda V:\operatorname{Sym}^m(\operatorname{Sym}^nV)]
\] in a dimension at least the length of \(\lambda\). Fix \(m\ge3\),
\(n\ge0\), and a partition \[
\lambda=(\alpha,\beta^{m-2},\gamma)\vdash mn,
\qquad \alpha\ge\beta\ge\gamma\ge0.
\] Put \[
N=\binom m2,\quad \epsilon=\beta\bmod2\in\{0,1\},\quad
L=\max(0,\beta-n),\quad
R=\beta-\gamma-N\epsilon-mL.
\tag{E.1}
\] Equivalently, \(R=\min(\alpha-\beta,\beta-\gamma)-N\epsilon\). The
equality follows from the exact identity
\((\alpha-\beta)-(\beta-\gamma)=m(n-\beta)\), using \(|\lambda|=mn\);
the basis bijection proving the formula is given in (E.26). The result
proved below is \[
\boxed{
a_{m,n}^{(\alpha,\beta^{m-2},\gamma)}
=\#\left\{(c_2,\ldots,c_m)\in\mathbf Z_{\ge0}^{m-1}:
\sum_{i=2}^m i c_i=R\right\}.}
\tag{E.2}
\] The set is empty for \(R<0\). This is one bounded rational polytope,
described directly by \(m-1\) variables, nonnegativity, and one
equation. In particular, \[
\boxed{a_{m,n}^{\lambda}>0\quad\Longleftrightarrow\quad R=0\text{ or }R\ge2.}
\tag{E.3}
\] The same construction gives an explicit flat toric degeneration of
the entire corresponding subalgebra of the balanced normalizer invariant
algebra. It involves no Reynolds enumeration, Hilbert-basis computation,
or elimination precomputation. For \(m=3\), the shape condition includes
every possible output partition of \(s_3[s_n]\). For \(m=4\), it
includes every \((\alpha,\beta,\beta,\gamma)\), with arbitrary
determinant shift \(\gamma\).

\subsection{T-024. Exact incidence and normalizer
algebra}\label{t-024.-exact-incidence-and-normalizer-algebra}

\subsubsection{1. The flag-algebra
morphism}\label{the-flag-algebra-morphism}

Let \(E=F=\mathbf Q^m\). Use the polynomial representation \[
P=\operatorname{Sym}(E\otimes F)=\mathbf Q[X_{ij}:1\le i,j\le m].
\] The two general linear groups act on the two factors. Choose the
upper unipotent group \(U_F\) so that the minors of the first \(j\)
columns are highest-weight vectors of right weight
\(\omega_j=(1^j,0^{m-j})\). Polynomial Cauchy duality gives \[
P^{U_F}=\bigoplus_{\lambda:\ell(\lambda)\le m} S_\lambda E,
\tag{E.4}
\] where each summand is marked by its right highest weight \(\lambda\).
This is the same polynomial Howe algebra used in T-002 and G-009. Taking
the determinant-one row torus \(T_0\) keeps the equal row-weight spaces,
and taking literal permutation matrices \(S_m\) then gives \[
\dim ((P^{U_F})^{T_0\rtimes S_m})_{n,\lambda}
=a_{m,n}^{\lambda}.
\tag{E.5}
\] Indeed, before taking \(U_F\), the row weight \((n^m)\) is
\((\operatorname{Sym}^nF)^{\otimes m}\); ordinary row permutations are
ordinary tensor permutations. Their invariants are
\(\operatorname{Sym}^m(\operatorname{Sym}^nF)\). Taking right highest
vectors gives (E.5). Thus the finite group here is fixed by the actual
representation, including its determinant character.

Define \[
x_i=X_{i1},\qquad
y_i=(-1)^{i-1}\det X_{[m]\setminus\{i\},[m-1]},\qquad
z=\det X.
\tag{E.6}
\] The span of the \(x_i\) is \(E\), the span of the \(y_i\) is
\(\Lambda^{m-1}E\simeq E^*\otimes\det E\), and \(z\) spans \(\det E\).
Their right weights are respectively \(\omega_1,\omega_{m-1},\omega_m\).
The relevant right-weight submonoid consists exactly of \[
\lambda=(\alpha,\beta^{m-2},\gamma)
=(\alpha-\beta)\omega_1+(\beta-\gamma)\omega_{m-1}+\gamma\omega_m.
\tag{E.7}
\] Consequently the corresponding subalgebra of \(P^{U_F}\) is generated
by the \(x_i,y_i,z\). Here is a direct generation argument. The
subalgebra generated by these three left-stable spaces is left
\(GL(E)\)-stable. The product of suitable nonzero highest vectors from
the three spaces is a nonzero polynomial in every degree (E.7), since
\(P\) is a domain. It lies in the single irreducible left module
\(S_\lambda E\) occurring in (E.4). A nonzero left-stable subspace of
that module is the whole module. This proves generation in every marked
degree, including weights on the boundary where some coefficient in
(E.7) vanishes.

There is precisely one defining relation: \[
\boxed{
P^{U_F}_{\mathrm{ends}}
\simeq\mathbf Q[x_1,\ldots,x_m,y_1,\ldots,y_m,z]
\big/\left(\sum_{i=1}^m x_i y_i\right).}
\tag{E.8}
\] The relation is the determinant expansion with the first column
repeated. To prove that there are no others, work over an algebraic
closure and prescribe \(x\ne0\), \(y\ne0\), with \(y(x)=0\). Choose
columns \(x,c_2,\ldots,c_{m-1}\) forming a basis of \(\ker(y)\). Their
exterior product is a nonzero scalar multiple of the prescribed oriented
\((m-1)\)-vector associated to \(y\). Since \(m\ge3\), one of
\(c_2,\ldots,c_{m-1}\) can be rescaled while leaving \(x\) fixed, so
that the exterior product is exactly the prescribed one. The final
column can then be chosen with arbitrary determinant \(z\), because its
pairing with the nonzero cofactor vector is a nonzero linear functional.
Thus the matrix map (E.6) reaches the dense open set \(x\ne0,y\ne0\) of
the incidence hypersurface times \(\mathbf A^1_z\).

For completeness, the incidence polynomial is irreducible: regard it as
a polynomial linear in \(y_1\) over
\(\mathbf Q[x_1,\ldots,x_m,y_2,\ldots,y_m]\). The coefficient \(x_1\)
and constant term \(\sum_{i>1}x_i y_i\) are relatively prime. Gauss's
lemma proves irreducibility. The resulting hypersurface is a domain, so
density proves the asserted kernel. The map and all equations have
rational coefficients, giving (E.8) over \(\mathbf Q\).

\subsubsection{2. Exact torus quotient and the signs of permutation
matrices}\label{exact-torus-quotient-and-the-signs-of-permutation-matrices}

A monomial \(x^a y^b z^t\) has row weights \[
a_i+|b|-b_i+t\qquad (1\le i\le m).
\tag{E.9}
\] It is fixed by \(T_0\) precisely when all differences \(a_i-b_i\) are
equal. Define \[
u=\prod_i x_i,\qquad v=\prod_i y_i,\qquad t_i=x_i y_i.
\tag{E.10}
\] If the common difference is \(k\ge0\), the monomial is
\(u^k\prod_i t_i^{b_i}z^t\); if it is \(k<0\), it is
\(v^{-k}\prod_i t_i^{a_i}z^t\). These expressions also prove, by their
distinct exponent vectors in the original \(x,y,z\), that the only
relation before the incidence quotient is \(uv=\prod_i t_i\). The
incidence relation has \(T_0\)-weight zero, and exact weight
decomposition passes to its quotient. More explicitly, if a weight-zero
element equals \((\sum_i x_i y_i)h\), then only the weight-zero part of
\(h\) contributes, because the multiplier has weight zero and the
polynomial ring is a domain. Therefore \[
\boxed{
B_{\mathrm{ends}}:=(P^{U_F}_{\mathrm{ends}})^{T_0}
\simeq\mathbf Q[u,v,z,t_1,\ldots,t_m]
/(\sum_i t_i,\ uv-\prod_i t_i).}
\tag{E.11}
\]

For a literal permutation matrix \(\sigma\), the exterior-power
identification in (E.6) gives \[
\sigma x_i=x_{\sigma(i)},\qquad
\sigma y_i=\operatorname{sgn}(\sigma)y_{\sigma(i)},\qquad
\sigma z=\operatorname{sgn}(\sigma)z.
\tag{E.12}
\] Thus \[
\sigma u=u,\qquad \sigma v=\operatorname{sgn}(\sigma)^m v,
\qquad \sigma t_i=\operatorname{sgn}(\sigma)t_{\sigma(i)}.
\tag{E.13}
\] Both equations in (E.11) generate stable ideals under this action.
These signs are retained in every subsequent invariant calculation.

\subsubsection{3. Alternating-group presentation and diagonal
involution}\label{alternating-group-presentation-and-diagonal-involution}

Let \(\mathfrak A_m\) denote the alternating group. Its action on the
\(t_i\) is ordinary permutation, and it fixes \(u,v,z\). Write \[
e_i=e_i(t_1,\ldots,t_m),\qquad
D=\prod_{i<j}(t_i-t_j),\qquad N=\binom m2.
\tag{E.14}
\] Then \[
\mathbf Q[t_1,\ldots,t_m]^{\mathfrak A_m}
=\mathbf Q[e_1,\ldots,e_m]\oplus D\mathbf Q[e_1,\ldots,e_m].
\tag{E.15}
\] To verify (E.15), average an alternating-group invariant with its
image under one odd permutation. Its plus part is symmetric; its minus
part transforms by the sign character under all permutations. A sign
polynomial vanishes at each \(t_i=t_j\), hence is divisible by every
\(t_i-t_j\), and therefore by their product \(D\). Its quotient by \(D\)
is symmetric. The fundamental theorem of symmetric polynomials gives
(E.15), and uniqueness follows from the different characters. Its
multiplication is specified by \[
D^2=\mathcal D_m(e_1,\ldots,e_m),
\tag{E.16}
\] where \(\mathcal D_m\) is the discriminant of
\(T^m-e_1T^{m-1}+e_2T^{m-2}-\cdots+(-1)^m e_m\).

Taking invariants of the two relations in (E.11) is exact. One can see
the exactness without a structural assumption on the quotient: averaging
any expression \(e_1h+(uv-e_m)k\) over \(\mathfrak A_m\) averages
\(h,k\), since both multipliers are invariant. Thus set \(e_1=0\) and
substitute \(e_m=uv\) in (E.15)--(E.16). The result is \[
\boxed{
B_{\mathrm{ends}}^{\mathfrak A_m}
\simeq
\mathbf Q[u,v,z,e_2,\ldots,e_{m-1},D]
\big/(D^2-\mathcal D_m(0,e_2,\ldots,e_{m-1},uv)).}
\tag{E.17}
\] In particular, this ring is free, with basis \(1,D\), over the
polynomial ring in \(u,v,z,e_2,\ldots,e_{m-1}\). This follows from the
monic quadratic relation; no assertion about algebraic independence has
been left to a dimension count.

The remaining quotient \(S_m/\mathfrak A_m\) acts diagonally. On an odd
permutation, (E.13) gives \[
u\longmapsto u,\quad v\longmapsto(-1)^m v,\quad
z\longmapsto-z,\quad e_i\longmapsto(-1)^i e_i,
\quad D\longmapsto(-1)^{N+1}D.
\tag{E.18}
\] The last sign has two factors: the ordinary Vandermonde permutation
sign and the factor \((-1)^N\) from negating all \(t_i\). Hence \[
\boxed{A_{\mathrm{ends}}:=(B_{\mathrm{ends}})^{S_m}
=(B_{\mathrm{ends}}^{\mathfrak A_m})^{\langle\tau\rangle},}
\tag{E.19}
\] where \(\tau\) is exactly the diagonal involution (E.18). A
vector-space basis of this algebra consists of the monomials \[
u^a v^b z^t\prod_{i=2}^{m-1}e_i^{c_i}D^\varepsilon,
\quad a,b,t,c_i\ge0,\quad\varepsilon\in\{0,1\},
\tag{E.20}
\] whose sign exponent \[
mb+t+\sum_{i=2}^{m-1} i c_i+(N+1)\varepsilon
\tag{E.21}
\] is even. The basis follows directly from (E.17); every basis vector
is an eigenvector of \(\tau\).

\subsection{T-025. One bounded polytope and a complete nonvanishing
test}\label{t-025.-one-bounded-polytope-and-a-complete-nonvanishing-test}

Set \(K=\sum_{i=2}^{m-1}i c_i+N\varepsilon\). The basis monomial (E.20)
has balanced row degree and right highest weight \[
\begin{aligned}
n&=a+(m-1)b+t+K,\\
\lambda_1&=m(a+b)+t+2K,\\
\lambda_2=\cdots=\lambda_{m-1}&=mb+t+K,\\
\lambda_m&=t.
\end{aligned}
\tag{E.22}
\] These formulas can be checked on each generator: \(u\) has right
weight \(m\omega_1\) and balanced degree 1; \(v\) has right weight
\(m\omega_{m-1}\) and balanced degree \(m-1\); \(z\) has weight
\(\omega_m\) and degree 1; each \(t_i\) has weight
\(\omega_1+\omega_{m-1}=(2,1^{m-2},0)\) and balanced degree 1. Thus
\(e_i\) has \(i\) times that weight and degree \(i\), while \(D\) has
\(N\) times that weight and degree \(N\).

For a fixed partition \((\alpha,\beta^{m-2},\gamma)\) and degree \(n\),
(E.22) gives \[
t=\gamma,\qquad a-b=n-\beta,\qquad
mb+\sum_{i=2}^{m-1}i c_i=\beta-\gamma-N\varepsilon.
\tag{E.23}
\] The first-row equation is then automatic from \(|\lambda|=mn\).
Substitution into the sign condition (E.21) gives simply \[
\beta+\varepsilon\equiv0\pmod2.
\tag{E.24}
\] Therefore \(\varepsilon=\epsilon=\beta\bmod2\) is uniquely
determined. The conditions \(a,b\ge0\) are equivalent to
\(b\ge L=\max(0,\beta-n)\). Set \(b=L+c_m\). The remaining equation
becomes exactly \[
\sum_{i=2}^m i c_i
=\beta-\gamma-N\epsilon-mL=R.
\tag{E.25}
\] The map from a lattice point in (E.25) to its basis vector is
explicit and bijective: \[
\begin{aligned}
b&=L+c_m, &a&=n-\beta+L+c_m, &t&=\gamma, &\varepsilon&=\epsilon.
\end{aligned}
\tag{E.26}
\] This proves (E.2), with neither a subtraction of counts nor an
unknown rank. Boundedness follows from \(0\le c_i\le R/i\) whenever
\(R\ge0\).

Since the denominators include 2 and 3, their additive semigroup is
exactly \(\{0\}\cup\{2,3,4,\ldots\}\): even integers are sums of 2's,
and every odd integer at least 3 is 3 plus a sum of 2's. No sum of the
allowed positive integers is 1. This proves (E.3), including all cases
\(n=0\), \(\gamma=0\), \(\alpha=\beta\), or \(\beta=\gamma\).

\subsubsection{Encoding and the precise efficiency
conclusion}\label{encoding-and-the-precise-efficiency-conclusion}

The polytope is constructed by integer arithmetic on
\(m,n,\beta,\gamma\), and its full description has \(O(m)\) integer
entries of bit length \(O(\log(m(n+1)+1))\). There is no
parameter-dependent algebra presentation to compute beforehand. A
lattice witness has \(O(m\log(m(n+1)+1))\) bits and is verified by one
weighted sum. Thus this coefficient family has a direct
counting-verifier interpretation in \(\#P\) for ordinary explicit
partition encoding. In that encoding, a nonzero middle part gives at
least \(m-1\) explicitly listed parts; when \(\beta=0\), the coefficient
is 1 and can be handled separately. A representation of partitions that
compresses long strings of equal parts changes the size convention; the
displayed \(O(m)\) description is not asserted polynomial in the length
of that compressed encoding. The nonvanishing criterion (E.3) is an
integer-arithmetic polynomial-time test even when the repeated middle
part is compressed. No polynomial-time assertion for exact counting with
variable \(m\) is needed here.

\subsubsection{Concrete subfamilies and
checks}\label{concrete-subfamilies-and-checks}

For \(m=3\), the coefficient is the number of solutions \(2c_2+3c_3=R\),
with \(R=\beta-\gamma-3(\beta\bmod2)-3\max(0,\beta-n)\). Root
independently compared this formula and its support test with the
separate finite-character formula G-007 for all 261,425 partitions
occurring as test inputs in degrees \(n=1,\ldots,100\). Every comparison
passed; see
\texttt{\small work/\allowbreak{}root/\allowbreak{}EQUAL\_\allowbreak{}MIDDLE\_\allowbreak{}M3\_\allowbreak{}CHECK.tsv}.

The rectangular agent independently expanded the power-sum cycle index
of \(h_m[h_n]\) and used Murnaghan--Nakayama characters. All 151
admissible end-shape cases for \(m=4,n=0,\ldots,6\),
\(m=5,n=0,\ldots,4\), and \(m=6,n=0,\ldots,3\) passed. Code:
\texttt{\small work/\allowbreak{}rectangular/\allowbreak{}verify\_\allowbreak{}endsector\_\allowbreak{}formula.py};
receipt:
\texttt{\small work/\allowbreak{}rectangular/\allowbreak{}ENDSECTOR\_\allowbreak{}INDEPENDENT\_\allowbreak{}CHECK.tsv}.
Those calculations do not derive their comparison values from the
incidence algebra.

\subsection{T-026. Explicit uniform toric
degeneration}\label{t-026.-explicit-uniform-toric-degeneration}

This section gives a flat algebra family whose general fiber is (E.19),
and whose special fiber is a semigroup algebra. The construction is
direct for every \(m\); it does not require a computation of a
generating set for the full flag normalizer algebra.

Start from (E.17). Assign temporary nonnegative weights \(w\) as
follows.

\begin{itemize}
\tightlist
\item
  For even \(m\), set \(w(u)=w(v)=1\), \(w(D)=m-1\), and give all other
  generators weight zero.
\item
  For odd \(m\), set \(w(e_{m-1})=2\), \(w(D)=m\), and give all other
  generators weight zero.
\end{itemize}

The discriminant has weighted degree \(m(m-1)=2N\) for the coefficient
weights \(\deg(e_i)=i\). For even \(m\), the unique term with greatest
possible exponent of \(e_m\) is \[
q_m e_m^{m-1},\qquad q_m\in\mathbf Q^\times.
\tag{E.27}
\] Its coefficient is nonzero because the polynomial \(T^m+(-1)^m e_m\)
has no repeated root for \(e_m\ne0\) in characteristic zero. The
coefficient degree forces an exponent at most \(m-1\), and equality
leaves degree zero for all other factors, proving uniqueness. After
\(e_m=uv\), its temporary weight is \(2(m-1)\), equal to \(2w(D)\);
every other term has smaller temporary weight.

For odd \(m\), the unique term with greatest possible exponent of
\(e_{m-1}\) is \[
r_m e_{m-1}^{m},\qquad r_m\in\mathbf Q^\times.
\tag{E.28}
\] Again the coefficient degree bounds the exponent by \(m\), with
equality unique. Its coefficient is nonzero because
\(T^m+(-1)^{m-1}e_{m-1}T\) has one simple root at zero and \(m-1\)
distinct nonzero roots for \(e_{m-1}\ne0\). Its temporary weight is
\(2m=2w(D)\), greater than all other terms.

If strictly positive generator weights are desired, add the original
balanced degree to each temporary weight. The original relation is
homogeneous of balanced degree \(2N\), so this leaves all comparisons of
its terms unchanged.

Introduce a parameter \(s\), and for a monomial \(M\) in the
discriminant let \(w(M)\) be its temporary weight. Define \[
\mathscr C=
\mathbf Q[s,u,v,z,e_2,\ldots,e_{m-1},D]
\Big/\left(D^2-
\sum_M [M]\mathcal D_m\;s^{2w(D)-w(M)}M\right),
\tag{E.29}
\] where \(\mathcal D_m=\mathcal D_m(0,e_2,\ldots,e_{m-1},uv)\). All
powers of \(s\) are nonnegative by (E.27)--(E.28). The relation is monic
in \(D\), so \(\mathscr C\) is free with basis \(1,D\) over
\(\mathbf Q[s,u,v,z,e_2,\ldots,e_{m-1}]\), and in particular is flat
over \(\mathbf Q[s]\). The fiber at \(s=1\) is (E.17). Over \(s\ne0\),
the exact isomorphism from the original algebra tensored with
\(\mathbf Q[s,s^{-1}]\) to \(\mathscr C[s^{-1}]\) sends each original
generator \(X\) to \(s^{-w(X)}X\). Its inverse sends the new generator
\(X\) to \(s^{w(X)}X_{\rm original}\). Substitution into the defining
equations verifies both maps. At zero its relation is respectively \[
D^2-q_m(uv)^{m-1}\quad(m\text{ even}),\qquad
D^2-r_m e_{m-1}^{m}\quad(m\text{ odd}).
\tag{E.30}
\] These are prime toric binomials, including over \(\mathbf Q\) after a
rational rescaling. In either case the exponent \(q\) on the other
monomial is odd. For a relation \(D^2-cY^q\), put \(Y'=cY\) and
\(D'=c^{(q-1)/2}D\), giving \((D')^2=(Y')^q\). In the even case take
\(Y=uv\) and rescale \(u\) alone. The odd case is the cusp semigroup
ring \(\mathbf Q[t^2,t^m]\) tensored with a polynomial ring. The even
case embeds as \[
u\longmapsto a^2,\qquad v\longmapsto b^2,
\qquad D\longmapsto a^{m-1}b^{m-1}.
\tag{E.31}
\] Its kernel is exactly the displayed binomial: the monomials
\(u^i v^j D^\varepsilon\), \(\varepsilon\in\{0,1\}\), have distinct
image exponents because \(m-1\) is odd. These arguments prove the domain
and semigroup assertions directly. Normality is not asserted.

The involution (E.18) extends to \(\mathscr C\) by fixing \(s\). It
preserves every term of the relation, since every discriminant term has
even sign. In the free \(\mathbf Q[s]\)-basis of monomials (E.20), its
action is diagonal and independent of \(s\). Therefore \[
\mathscr A=\mathscr C^{\langle\tau\rangle}
\tag{E.32}
\] is free over \(\mathbf Q[s]\), with basis the same sign-even
monomials (E.20)--(E.21). Its fibers commute with invariants, as is also
immediate from this basis. Its fiber at 1 is \(A_{\mathrm{ends}}\). Its
fiber at 0 is the subalgebra of the semigroup algebra (E.30) spanned by
monomials of even sign. Since the defining binomial preserves the sign,
this is a semigroup algebra for the kernel of a homomorphism of the
exponent semigroup to \(\mathbf Z/2\). It is finitely generated
explicitly by the even-sign generators and the pairwise products of the
odd-sign generators: any even-sign monomial can be grouped into those
factors. Thus (E.32) is an explicit flat toric degeneration of the
desired end-sector normalizer algebra.

The family respects the balanced degree and the right highest-weight
grading throughout, because the defining relation and its separate terms
all have the same such degrees. Its basis is the same basis used in
(E.25), so the polytope formula counts the graded Hilbert function of
both general and special fibers. This proves the precise morphism
between the plethysm coefficients, the incidence quotient, and the toric
model.

\subsection{A fully written outer-degree-four
illustration}\label{a-fully-written-outer-degree-four-illustration}

In the length-three part \(m=4,\gamma=0\), let \(s_2=e_2\), \(q=e_3^2\),
and \(w=e_3D\). In (E.18), \(e_3,D\) are odd and \(u,v,e_2\) are even.
Hence \[
A_{4,\mathrm{ends},\gamma=0}
\simeq\mathbf Q[u,v,s_2,q,w]/(w^2-q\mathcal D(s_2,q,uv)),
\tag{E.33}
\] where the depressed-quartic discriminant is \[
\mathcal D(s_2,q,p)
=256p^3-128s_2^2p^2+144s_2qp-27q^2
+16s_2^4p-4s_2^3q.
\tag{E.34}
\] The ring is free with basis \(1,w\) over \(\mathbf Q[u,v,s_2,q]\).
The generator degrees \((n;\lambda)\) are \[
\begin{array}{c|c|c}
u&1&(4,0,0)\\
v&3&(4,4,4)\\
s_2&2&(4,2,2)\\
q&6&(12,6,6)\\
w&9&(18,9,9).
\end{array}
\tag{E.35}
\] Thus one also obtains two four-variable polytopes \[
n=a+3b+2c+6d+9\delta,\qquad
\beta=4b+2c+6d+9\delta,
\quad a,b,c,d\ge0,\quad\delta\in\{0,1\}.
\tag{E.36}
\] They are equivalent, by the explicit basis analysis above, to the
single three-variable polytope (E.25). Formula (E.33) supplies a small
presentation in addition to the count; all determinant shifts are
included in the uniform construction (E.17)--(E.25).

\subsection{Relation to the main
gates}\label{relation-to-the-main-gates}

G-009--G-011 already give terminating positive-polytope constructions
for all symmetric plethysm, general Schur plethysm, and canonical-map
image multiplicities. T-024--T-026 replace their parameter-dependent
invariant computation by one uniform explicit presentation, one bounded
polytope, and an arithmetic nonvanishing test on the family (E.7). They
also give a direct toric degeneration in every outer degree. The
arbitrary-output efficient formula remains outside this proved family.
Nothing here asserts that the dimension of a canonical-map image equals
a coefficient; the actual canonical images retain G-011's two-projector
morphism.

\subsection{Sources, provenance, and proof
checks}\label{sources-provenance-and-proof-checks}

The input representation identity is polynomial Cauchy/Howe duality with
its balanced permutation action, proved in T-002 and re-derived in
(E.4)--(E.5). An exact primary locator supplied with the task is J. M.
Landsberg, \emph{Geometric Complexity Theory: an introduction for
geometers}, subsection ``A duality theorem for weight zero spaces and
plethysms,'' theorem \texttt{\small gaythm}, source lines 766--798; its
underlying result is attributed there to Gay (1976). The trace-free
incidence presentation, the alternating-group calculation, the full
determinant sign, the bijection (E.26), and the flat family
(E.29)--(E.32) are all proved in this report. No unlocated literature
theorem is used to infer the formula. The elementary discriminant
coefficient arguments are included explicitly in (E.27)--(E.28).

For the classical relative-invariant step, R. P. Stanley,
\emph{Invariants of finite groups and their applications to
combinatorics}, Bull. Amer. Math. Soc. \textbf{1} (1979), 475--511,
Proposition 4.7 and its symmetric-group specialization on printed
pp.~488--489, identifies alternating polynomials as the Vandermonde
times symmetric polynomials. The primary author-hosted source was read
directly: \url{https://math.mit.edu/~rstan/pubs/pubfiles/38.pdf}.
Equation (E.15) gives the exact index-two decomposition needed here,
with the later determinant twist computed separately.

The proof is to be recorded as a whole-proof check only after reading
every step of this report. Independent computation receipts and
independent proof readers are listed separately in
\texttt{\small CHECK\_\allowbreak{}EVENTS.tsv}; the numerical
comparisons are evidence additional to the proof, not replacements for
it.

\section{Arbitrary outer Schur functors on the equal-middle-row
family}\label{arbitrary-outer-schur-functors-on-the-equal-middle-row-family}

This report extends T-024--T-026 by identifying the full balanced
permutation module. All statements below are unconditional over
characteristic zero, and every coefficient is accompanied by an explicit
algebra or lattice-point morphism. Author/checker: GPT-6 Astra Pro, run
\texttt{\small GCT\_\allowbreak{}PLETH-\allowbreak{}20261121-\allowbreak{}theta};
effort setting is not exposed.

\subsection{T-027. Exact reflection-representation model and positive
tableau
formula}\label{t-027.-exact-reflection-representation-model-and-positive-tableau-formula}

Let \(m\ge3\), \(\mu\vdash m\), \(n\ge0\), and \[
\lambda=(\alpha,\beta^{m-2},\gamma)\vdash mn,
\qquad \alpha\ge\beta\ge\gamma\ge0.
\] Write \[
p_{\mu,(n)}^\lambda=[S_\lambda V:S_\mu(\operatorname{Sym}^nV)],
\qquad
S=\min(\alpha-\beta,\beta-\gamma).
\tag{S.1}
\] Let \(\mathfrak h_m\) be the \((m-1)\)-dimensional standard
representation of \(S_m\), realized as the quotient of its permutation
representation by the trivial line. Set \[
\nu=\begin{cases}\mu,&\beta\text{ even},\\ \mu',&\beta\text{ odd}.\end{cases}
\tag{S.2}
\] For a standard Young tableau \(T\) with entries \(1,\ldots,m\),
define \[
\operatorname{Des}(T)=\{i: i+1\text{ lies in a lower row than }i\},
\qquad
\operatorname{comaj}(T)=\sum_{i\in\operatorname{Des}(T)}(m-i).
\tag{S.3}
\] Then the exact formulas are \[
\boxed{
(S_\lambda\mathbf Q^m)_{(n^m)}
\simeq\operatorname{sgn}^{\beta}\otimes
\operatorname{Sym}^{S}(\mathfrak h_m)
\quad\text{as }S_m\text{-modules},}
\tag{S.4}
\] and \[
\boxed{
p_{\mu,(n)}^\lambda
=\sum_{T\in\operatorname{SYT}(\nu)}
\#\left\{(c_2,\ldots,c_m)\in\mathbf Z_{\ge0}^{m-1}:
\sum_{i=2}^m i c_i=S-\operatorname{comaj}(T)\right\}.}
\tag{S.5}
\] Negative right sides contribute zero. The sum is finite and entirely
explicit. T-028 below realizes it as the lattice-point count of one
polytope with a polynomial-size description.

\subsubsection{Proof of the module isomorphism, including its character
twist}\label{proof-of-the-module-isomorphism-including-its-character-twist}

Use the exact balanced incidence algebra from T-024: \[
B_{\mathrm{ends}}
=\mathbf Q[u,v,z,t_1,\ldots,t_m]
/(\sum_i t_i,\ uv-\prod_i t_i),
\tag{S.6}
\] with the right weights and literal permutation action in
(E.12)--(E.13) of
\texttt{\small end\_\allowbreak{}node\_\allowbreak{}report.md}. Put \[
k=n-\beta,\quad k_+=\max(k,0),\quad L=\max(-k,0).
\] Before imposing \(\sum_i t_i=0\), every monomial has a unique normal
form with at most one of \(u,v\) present, using \(uv=\prod_i t_i\).
Distinct such normal monomials have distinct exponent vectors in the
original \(x_i,y_i,z\), as proved in T-024. In the component of degree
\((n,\lambda)\), the normal form is necessarily \[
u^{k_+}v^L z^\gamma f(t_1,\ldots,t_m),
\qquad \deg f=\beta-\gamma-mL.
\tag{S.7}
\] This follows directly by subtracting the degree and middle-row
equations (E.22), which gives the fixed difference \(a-b=k\);
eliminating every factor \(uv\) increases the \(t\)-degree by \(m\). The
equation \(\sum_i t_i=0\) does not change that difference. Its part in
this component is exactly the same factor in (S.7) times a multiple of
\(\sum_i t_i\). Thus multiplication in (S.7) gives a vector-space
isomorphism \[
\left(\mathbf Q[t_1,\ldots,t_m]/(\sum_i t_i)\right)_S
\longrightarrow (B_{\mathrm{ends}})_{n,\lambda}.
\tag{S.8}
\] Here the degree equals \(S\), since \[
(\alpha-\beta)-(\beta-\gamma)=m(n-\beta)=mk
\] implies \[
\beta-\gamma-m\max(0,\beta-n)
=\min(\alpha-\beta,\beta-\gamma).
\tag{S.9}
\] This also proves its nonnegativity under the stated partition
inequalities.

For an odd permutation, the factor \(v^Lz^\gamma\) contributes the sign
\((-1)^{mL+\gamma}\). The degree-\(S\) polynomial in the \(t_i\)
contributes an additional \((-1)^S\) under their signed permutation
action. Since \(S+mL+\gamma=\beta\), the total twist is exactly
\(\operatorname{sgn}^{\beta}\). The quotient by \(\sum_i t_i\) is the
symmetric algebra on \(\mathfrak h_m\). Finally the polynomial Cauchy
identification (E.4) identifies the right-hand side of (S.8), with all
actions retained, with \((S_\lambda\mathbf Q^m)_{(n^m)}\). This proves
(S.4).

The refined balanced Howe identity T-003 gives \[
p_{\mu,(n)}^\lambda
=\dim\operatorname{Hom}_{S_m}([\mu],(S_\lambda\mathbf Q^m)_{(n^m)}).
\tag{S.10}
\] The exact sign-twist identity
\([\mu]\otimes\operatorname{sgn}\simeq[\mu']\) identifies this with the
multiplicity of \([\nu]\) in \(\operatorname{Sym}^{S}\mathfrak h_m\).
One character verification of this identity is as follows: the
symmetric-function involution \(p_i\mapsto(-1)^{i-1}p_i\) sends
\(s_\mu\) to \(s_{\mu'}\); under the Frobenius formula it multiplies the
coefficient on a cycle type \(\rho\) by \((-1)^{m-\ell(\rho)}\), exactly
the permutation sign. Equality of the resulting characters gives the
module isomorphism in characteristic zero.

\subsubsection{Cycle traces and a full standardization
proof}\label{cycle-traces-and-a-full-standardization-proof}

For a permutation with cycle lengths \(\rho_1,\ldots,\rho_\ell\), the
permutation representation has characteristic determinant \[
\det(1-q\sigma)=\prod_{j=1}^\ell(1-q^{\rho_j}).
\] Removing the trivial line gives \[
\sum_{d\ge0}\operatorname{tr}(\sigma;\operatorname{Sym}^d\mathfrak h_m)q^d
=\frac{1-q}{\prod_j(1-q^{\rho_j})}.
\tag{S.11}
\] Take its inner product with the character \(\chi^\nu\). The Frobenius
character formula, evaluated at the alphabet \(1,q,q^2,\ldots\), gives
\[
\begin{aligned}
\sum_{d\ge0}[\,[\nu]:\operatorname{Sym}^d\mathfrak h_m\,]q^d
&=(1-q)\sum_{\rho\vdash m}\frac{\chi^\nu(\rho)}{z_\rho}
\prod_j\frac1{1-q^{\rho_j}}\\
&=(1-q)s_\nu(1,q,q^2,\ldots).
\end{aligned}
\tag{S.12}
\] The equality uses precisely \(p_i(1,q,q^2,\ldots)=(1-q^i)^{-1}\). It
is a formal power-series equality, so no analytic specialization is
needed.

We now derive the positive numerator without assuming a coinvariant-ring
decomposition. Semistandard tableaux of shape \(\nu\), with weak rows
and strict columns, enumerate \(s_\nu(1,q,q^2,\ldots)\), with weight
\(q^{\sum_c(Q(c)-1)}\). Standardize a semistandard tableau by ordering
cells by increasing entry and breaking ties from left to right. Equal
entries cannot occupy the same column; their left-to-right order runs
weakly upward in row number, by the column-strict and row-weak
conditions. The result is a standard tableau \(T\), together with
integers \[
0\le a_1\le\cdots\le a_m,
\qquad a_i<a_{i+1}\text{ for every }i\in\operatorname{Des}(T),
\tag{S.13}
\] where \(a_i+1\) is the original entry in the cell labelled \(i\).

Conversely, fill the cell labelled \(i\) in \(T\) by \(a_i+1\). The rows
are weak because \(T\) increases across rows. For two cells in one
column with labels \(p<q\), the row number increases from the cell of
\(p\) to the cell of \(q\), so some consecutive step \(i\to i+1\),
\(p\le i<q\), goes to a lower row. At that step (S.13) is strict, and
hence the column entries are strict. Within an equality block of the
sequence \(a_i\), the row numbers can never increase. A later cell in a
higher row must then be farther right: otherwise it would lie weakly
northwest of the earlier cell, contradicting the increasing row and
column labels of \(T\). Thus standardizing the filled tableau recovers
exactly \(T\). This proves the bijection.

For a fixed \(T\), write \[
b_1=a_1,\qquad b_{i+1}=a_{i+1}-a_i-\mathbf1_{i\in\operatorname{Des}(T)}\ge0.
\] Then \[
\sum_{i=1}^m a_i
=\operatorname{comaj}(T)+\sum_{j=1}^m(m-j+1)b_j.
\tag{S.14}
\] It follows that \[
s_\nu(1,q,q^2,\ldots)
=\frac{\sum_{T\in\operatorname{SYT}(\nu)}q^{\operatorname{comaj}(T)}}
{\prod_{i=1}^m(1-q^i)}.
\tag{S.15}
\] Substitution into (S.12) cancels exactly the factor \(1-q\), proving
\[
\sum_{d\ge0}[\,[\nu]:\operatorname{Sym}^d\mathfrak h_m\,]q^d
=\frac{\sum_{T\in\operatorname{SYT}(\nu)}q^{\operatorname{comaj}(T)}}
{\prod_{i=2}^m(1-q^i)}.
\tag{S.16}
\] Taking the coefficient of \(q^S\) proves (S.5). For \(\mu=(m)\), the
shape \(\nu\) is either a row with comajor index 0 or a column with
comajor index \(\binom m2\), recovering exactly T-025.

\subsection{T-028. A single polynomial-size polytope and a support
criterion}\label{t-028.-a-single-polynomial-size-polytope-and-a-support-criterion}

\subsubsection{The explicit polytope}\label{the-explicit-polytope}

The cells of \(\nu\) form a set \(C\) of size \(m\); write \(r(c)\) for
the row number of a cell. Introduce variables \[
X_{c,j}\quad(c\in C,\ 1\le j\le m),\qquad
\delta_j\quad(1\le j<m),\qquad
c_i\quad(2\le i\le m).
\] There are \(m^2+2m-2\) variables. Impose \[
0\le X_{c,j}\le1,\quad0\le\delta_j\le1,\quad c_i\ge0;
\qquad
\sum_jX_{c,j}=1\ (c\in C),\quad
\sum_cX_{c,j}=1\ (1\le j\le m).
\tag{S.17}
\] Define the following linear forms, which introduce no new variables:
\[
\ell_c=\sum_{j=1}^m jX_{c,j},\qquad
r_j=\sum_{c\in C}r(c)X_{c,j}.
\tag{S.18}
\] For each adjacent pair of cells in a row or column, with \(c'\)
immediately right of or below \(c\), require \[
\ell_c+1\le\ell_{c'}.
\tag{S.19}
\] For every \(1\le j<m\), require \[
1-m(1-\delta_j)\le r_{j+1}-r_j\le m\delta_j.
\tag{S.20}
\] Finally require \[
\sum_{i=2}^m i c_i+\sum_{j=1}^{m-1}(m-j)\delta_j=S.
\tag{S.21}
\]

At every integral point, (S.17) makes \(X\) a permutation assignment of
the labels \(1,\ldots,m\) to the cells of \(\nu\). Equation (S.19) is
exactly the standard-tableau condition. Since \(1\le r_j\le m\), (S.20)
uniquely forces \(\delta_j=1\) when \(r_{j+1}>r_j\), and \(\delta_j=0\)
otherwise. Thus its weighted sum in (S.21) is exactly
\(\operatorname{comaj}(T)\). The remaining variables count the
partitions in (S.5), without any multiplicity from the encoding.
Conversely every pair counted in (S.5) gives exactly one integral point.
Therefore \[
\boxed{p_{\mu,(n)}^\lambda=\#(\mathcal P_{\mu,\lambda}\cap\mathbf Z^{m^2+2m-2}),}
\tag{S.22}
\] where \(\mathcal P_{\mu,\lambda}\) is the explicitly displayed
polytope (S.17)--(S.21). It is bounded: the \(X,\delta\) variables lie
in \([0,1]\), and \(0\le c_i\le S/i\) follows from (S.21). Its number of
variables, constraints, and nonzero matrix entries is polynomial in
\(m\). Its coefficients have bit length \(O(\log(m(n+1)+1))\). This
gives a direct \(\#P\) verifier in ordinary explicit partition encoding,
with no algebra precomputation. As in T-025, a nonzero \(\beta\)
supplies at least \(m-1\) explicitly listed parts of \(\lambda\); for
\(\beta=0\), the answer is directly 1 for \(\mu=(m)\) and 0 for the
other outer shapes.

\subsubsection{The exact nonvanishing
criterion}\label{the-exact-nonvanishing-criterion}

Define \[
b(\nu)=\sum_{r\ge1}(r-1)\nu_r,
\qquad d(\nu)=\#\{r:\nu_r>\nu_{r+1}\},
\qquad Q=S-b(\nu).
\tag{S.23}
\] Here \(d(\nu)\) is the number of removable corners, equivalently the
number of distinct positive part sizes. Then \[
\boxed{
p_{\mu,(n)}^\lambda>0
\Longleftrightarrow
\left[Q=0\right]\ \text{or}\ \left[Q\ge2\right]\ \text{or}\
\left[Q=1\text{ and }d(\nu)\ge2\right].}
\tag{S.24}
\]

To prove this, the smallest exponent in the semistandard generating
function \(s_\nu(1,q,q^2,\ldots)\) is \(b(\nu)\), occurring once: in row
\(r\), column strictness forces every entry to be at least \(r\), and
filling that row entirely by \(r\) attains the bound. A semistandard
tableau of weight \(b(\nu)+1\) must increase exactly one cell of this
minimum tableau by 1. Row weakness permits this only at the rightmost
cell of a row; column strictness permits it only when no cell lies below
it. These are exactly the \(d(\nu)\) removable corners, and all such
increases are valid. Thus \[
s_\nu(1,q,q^2,\ldots)=q^{b(\nu)}(1+d(\nu)q+O(q^2)).
\] Multiplication by \(\prod_{i=1}^m(1-q^i)=1-q+O(q^2)\) in (S.15) gives
\[
\sum_Tq^{\operatorname{comaj}(T)}
=q^{b(\nu)}(1+(d(\nu)-1)q+O(q^2)).
\tag{S.25}
\] It follows that the minimum comajor index is \(b(\nu)\), attained
once, and the number of tableaux with comajor index \(b(\nu)+1\) is
\(d(\nu)-1\). No sum of parts \(2,\ldots,m\) is 1, whereas every integer
at least 2 is such a sum because \(m\ge3\). Formula (S.5) now proves
every case of (S.24): the unique minimum tableau gives all \(Q=0\) and
\(Q\ge2\), and at \(Q=1\) only the tableaux counted by \(d(\nu)-1\) can
contribute.

This is an integer-arithmetic polynomial-time support test. When
\(\beta\) is odd, the statistic can be computed directly on \(\mu\),
since \[
b(\mu')=\sum_i\binom{\mu_i}{2},\qquad d(\mu')=d(\mu).
\tag{S.26}
\] The first identity counts pairs of cells in the same row; the second
follows by exchanging the horizontal and vertical boundary corners of
the Ferrers diagram. Thus a long transpose need not be explicitly
expanded to test support. Compressed strings of repeated parts can also
be handled by arithmetic sums over their blocks for this support test;
the \(O(m^2)\) tableau polytope itself uses the ordinary explicit size
convention stated above.

\subsection{Exact input identities and
validation}\label{exact-input-identities-and-validation}

The only representation-theoretic inputs are polynomial Cauchy/Howe
duality, the refined balanced-weight identity T-003, and the Frobenius
character identity from Schur--Weyl duality. Their action on the
relevant module is specified in (S.4), (S.8), and (S.10). The exact
primary locator supplied with the task for the balanced-weight
realization is Landsberg, \emph{Geometric Complexity Theory: an
introduction for geometers}, subsection ``A duality theorem for weight
zero spaces and plethysms,'' theorem \texttt{\small gaythm}, source
lines 766--798. The Frobenius evaluation used in (S.12) is displayed
explicitly, and (S.13)--(S.15) give a complete proof of the needed
principal-specialization tableau identity. No generic coefficient
oracle, coinvariant decomposition, or unproved positivity assertion
enters the argument.

The cycle-trace averaging in (S.11)--(S.12) is the character-valued
Molien formula: R. P. Stanley, \emph{Invariants of finite groups and
their applications to combinatorics}, Bull. Amer. Math. Soc. \textbf{1}
(1979), 475--511, Theorem 2.1 and its complete proof on printed
pp.~479--480. The primary author-hosted text was read in this run:
\url{https://math.mit.edu/~rstan/pubs/pubfiles/38.pdf}. The tableau
numerator is a classical representation-theoretic object; the report
proves its precise comajor convention directly, rather than importing a
statistic with a different tableau convention.

The author and the rectangular agent have each read the complete T-027
and T-028 proofs; their actual events are recorded separately in the
respective \texttt{\small CHECK\_\allowbreak{}EVENTS.tsv} files. The
independent reader checked both directions of standardization, the
comajor convention, the exact module twist, all polytope constraints,
and the support argument. The same agent independently expanded the
outer Schur power-sum cycle index and used Murnaghan--Nakayama to test
all 231 coefficients, and their support, for every outer \(\mu\) and
every admissible end shape with \(m=3,n=0,\ldots,4\) and
\(m=4,n=0,\ldots,3\). Every comparison passed. Code:
\texttt{\small work/\allowbreak{}rectangular/\allowbreak{}verify\_\allowbreak{}endsector\_\allowbreak{}schur.py};
receipt:
\texttt{\small work/\allowbreak{}rectangular/\allowbreak{}ENDSECTOR\_\allowbreak{}SCHUR\_\allowbreak{}INDEPENDENT\_\allowbreak{}CHECK.tsv}.
These computations are additional checks, not replacements for the proof
readings.

\section{A face-dimension obstruction to one multigraded polynomial
normalization}\label{a-face-dimension-obstruction-to-one-multigraded-polynomial-normalization}

\subsection{T-029. Statement}\label{t-029.-statement}

Let \(m\ge4\), \(E=F=\mathbf Q^m\), and \[
A_m=\left(\operatorname{Sym}(E\otimes F)^{U_F}\right)^{T_0\rtimes S_m},
\tag{N.1}
\] with the exact literal row-permutation action and determinant-one row
torus used in T-002, T-024, and G-009. Grade this algebra by the full
right dominant weight, expressed in the independent fundamental weights
\[
\omega_j=(1^j,0^{m-j}),\qquad 1\le j\le m.
\] Then there is no polynomial subalgebra \[
\mathbf Q[f_1,\ldots,f_{\binom m2+1}]\subset A_m
\tag{N.2}
\] over which \(A_m\) is finite and whose generators \(f_i\) are all
homogeneous for this full \(\mathbf N^m\)-grading. Equivalently, \(A_m\)
has no homogeneous system of parameters for the full dominant-weight
grading. The conclusion concerns precisely the displayed grading and
finiteness morphism. G-009's finite standard-monomial polytope
construction, T-024--T-028's direct face formulas, and normalizations
using a coarser grading remain available with their stated proofs.

The original rank-four instance was proposed by the root agent; the
argument below proves the general rank statement. No claim of historical
novelty is made.

\subsection{Proof}\label{proof}

\subsubsection{1. Exact flag-face
dimensions}\label{exact-flag-face-dimensions}

Put \(r=m-1\). For \(J\subseteq\{1,\ldots,r\}\), let \(C_J\) be the
subalgebra of \(\operatorname{Sym}(E\otimes F)^{U_F}\) with dominant
degrees supported in \(J\cup\{m\}\). Thus the determinant degree is
retained in every face. As in the generation proof of T-024, \(C_J\) is
generated by all first-column flag minors of sizes \(j\in J\), together
with the full determinant \(z\): their products give a nonzero vector in
every permitted irreducible Cauchy summand, and left \(GL(E)\)-stability
supplies that entire summand.

For \(J\ne\varnothing\), write \(\operatorname{Fl}_J(E)\) for the
variety of nested subspaces with dimensions the elements of \(J\). A
dense open part of \(\operatorname{Spec}C_J\) consists of such a partial
flag, a nonzero exterior-volume vector for each selected subspace, and a
nonzero full determinant. Every such set of data is realized by an
adapted basis of \(E\): at each newly selected dimension, rescale a
newly introduced basis vector to fix that exterior volume; since the
previous selected dimensions are smaller, their volumes remain fixed.
The final basis vector fixes the full determinant. Conversely, the
nonzero decomposable exterior vectors recover the flag and its selected
volumes. Hence \[
\dim C_J=\dim\operatorname{Fl}_J(E)+|J|+1.
\tag{N.3}
\] This can also be read on the usual affine chart for the partial flag,
where the selected volume choices form a product of \(|J|+1\)
multiplicative groups. For \(J=\varnothing\), \(C_J=\mathbf Q[z]\) has
dimension 1.

We justify the torus dimension subtraction by an explicit localization.
Let \(f_J\) be the product of all flag minors of the selected sizes. It
is nonzero in the matrix polynomial ring. For each fixed size \(j\),
every row occurs in exactly \(\binom{m-1}{j-1}\) minors, so their
product has equal row weights and is fixed by \(T_0\). Thus \(f_J\) is
\(T_0\)-invariant. In \((C_J)_{f_J}\), every selected minor is a unit.

Let \(L_J\) be the lattice generated by their \(T_0\)-weights. It has
rank \(r\). Indeed, choose any \(j\in J\). For any two different indices
\(a,b\), choose a \(j\)-subset containing \(a\) and not \(b\). Replacing
\(a\) by \(b\) in that subset shows that the difference of two selected
minor weights is \(\varepsilon_a-\varepsilon_b\). These differences span
rank \(m-1\) on the determinant-one torus.

Choose a \(\mathbf Z\)-basis of \(L_J\), and realize each basis weight
by a Laurent monomial \(v_i\) in the now invertible minors. Every
homogeneous element of \((C_J)_{f_J}\) can be multiplied by a unique
Laurent monomial in the \(v_i\) to have weight zero. Distinct exponent
vectors in the \(v_i\) have distinct weights. Consequently
multiplication gives the exact algebra isomorphism \[
(C_J)_{f_J}
\simeq ((C_J)_{f_J})^{T_0}
[v_1^{\pm1},\ldots,v_r^{\pm1}].
\tag{N.4}
\] No assertion about generic closed orbits is needed for this
isomorphism. Since \(f_J\) is invariant,
\(((C_J)_{f_J})^{T_0}=((C_J)^{T_0})_{f_J}\). Taking dimensions in (N.4)
and using the dense nonempty principal open gives \[
\dim (C_J)^{T_0}=\dim C_J-r\qquad(J\ne\varnothing).
\tag{N.5}
\] The finite group \(S_m\) preserves each face. Its invariant inclusion
is integral: every element satisfies its monic orbit polynomial. Hence
passing to \(S_m\)-invariants preserves dimension. If \(A_J\) is the
corresponding face algebra of \(A_m\), then \[
d(J):=\dim A_J
=\begin{cases}
1,&J=\varnothing,\\
1+|J|+\dim\operatorname{Fl}_J(E)-r,&J\ne\varnothing.
\end{cases}
\tag{N.6}
\] In the empty case, the invariant algebra is \(\mathbf Q[z^2]\), since
odd row permutations negate \(z\).

The projective flag dimension has a particularly useful exact count. A
positive root of type \(A_r\) is an interval \([a,b]\subseteq[1,r]\). It
contributes to \(\dim\operatorname{Fl}_J(E)\) precisely when its
interval meets \(J\). To verify this directly, cut \(1,\ldots,m\) into
consecutive blocks at the positions in \(J\). The flag dimension is the
sum of products of the sizes of distinct blocks. This counts pairs of
indices in different blocks, or equivalently intervals crossing at least
one selected cut. Therefore \[
d(J)=1+|J|+
\sum_{I\text{ a nonempty interval in }[1,r]}\mathbf1_{I\cap J\ne\varnothing}
-r\mathbf1_{J\ne\varnothing}.
\tag{N.7}
\] In particular, \(d([1,r])=\binom m2+1\), the dimension used in (N.2).

\subsubsection{2. Restriction of a finite homogeneous normalization to
every
face}\label{restriction-of-a-finite-homogeneous-normalization-to-every-face}

Assume that (N.2) exists. For a homogeneous generator \(f_i\), define
its root support to be the set of indices in \(\{1,\ldots,r\}\) with
nonzero fundamental-weight coefficient in its degree; the determinant
coefficient is ignored in this support. Let \(c_S\) count generators
with exact root support \(S\subseteq[1,r]\).

For a face \(J\), retain the polynomial generators whose root supports
lie in \(J\). They remain algebraically independent in \(A_J\).
Moreover, \(A_J\) is finite over their polynomial subalgebra. To prove
this, choose finitely many homogeneous module generators for the assumed
finite extension (N.2). Such a choice exists by taking homogeneous
components of a finite generating set. In any expression of a
face-degree element in terms of these generators, take its homogeneous
component of that degree. Since the degree monoid is \(\mathbf N^m\), a
summand with a positive coordinate outside the face cannot contribute.
Thus only module generators and polynomial variables with degrees in
that face remain. These finitely many module generators prove finiteness
after restriction.

Equality of dimensions in this finite extension now forces \[
d(J)=\sum_{S\subseteq J}c_S\qquad\text{for every }J\subseteq[1,r].
\tag{N.8}
\] In particular all the numbers obtained by Boolean-lattice Möbius
inversion of \(d\) must be nonnegative integers.

\subsubsection{3. The forced negative generator
count}\label{the-forced-negative-generator-count}

For nonempty \(S\subseteq[1,r]\), invert (N.7). The constant 1
contributes zero. The term \(|J|\) contributes 1 for a singleton \(S\)
and zero otherwise. For a fixed interval \(I\), \[
\sum_{J\subseteq S}(-1)^{|S|-|J|}\mathbf1_{I\cap J\ne\varnothing}
=\begin{cases}(-1)^{|S|+1},&S\subseteq I,\\0,&S\not\subseteq I.\end{cases}
\tag{N.9}
\] Indeed, replace the indicator by
\(1-\mathbf1_{J\subseteq S\setminus I}\) and sum the resulting binomial
identity. The analogous transform of \(\mathbf1_{J\ne\varnothing}\) is
\((-1)^{|S|+1}\). Finally, the number of intervals containing \(S\) is
\(\min(S)(r-\max(S)+1)=\min(S)(m-\max(S))\). Hence (N.8) forces \[
\boxed{
c_S=\mathbf1_{|S|=1}
+(-1)^{|S|+1}\left(\min(S)(m-\max(S))-(m-1)\right).}
\tag{N.10}
\] Take \(S=\{1,2,m-1\}\), which has three distinct elements for every
\(m\ge4\). Then \[
c_S=1-(m-1)=2-m<0,
\] contradicting its definition as a number of generators. This proves
T-029.

\subsubsection{The rank-four instance}\label{the-rank-four-instance}

For \(m=4\), (N.6) gives \[
d(\varnothing)=1,\quad d(\{1\})=d(\{3\})=2,\quad d(\{2\})=3,
\quad d(\{i,j\})=5,\quad d(\{1,2,3\})=7.
\] The forced counts on proper supports are \[
c_\varnothing=1,\ c_{\{1\}}=c_{\{3\}}=1,\ c_{\{2\}}=2,
\quad c_{\{1,2\}}=c_{\{2,3\}}=1,\ c_{\{1,3\}}=2.
\] They already total 9, whereas the full algebra has dimension 7; the
missing full-support count is \(-2\). This is the original rank-four
obstruction in explicit numerical form.

\subsection{Provenance and scope}\label{provenance-and-scope}

The flag-minor and Cauchy construction is the same exact morphism used
in T-002, T-024, and G-009. The flag dimensions, invariant localization,
finite-face restriction, and Möbius contradiction are proved in full
above. This theorem excludes a specific global construction by one
polynomial subring homogeneous in every highest-weight coordinate. It
does not leave the existence of a terminating positive lattice count
unresolved: G-009--G-011 already provide such counts through their
specified algebra and module presentations, and the direct formulas
T-025 and T-028 provide polynomial-size descriptions on their respective
faces.

\section{The paired-row face in outer degree
four}\label{the-paired-row-face-in-outer-degree-four}

\subsection{T-030. Exact small presentation and all five outer Schur
coefficients}\label{t-030.-exact-small-presentation-and-all-five-outer-schur-coefficients}

Let \(\mu\vdash4\), \(n\ge0\), and \[
\lambda=(a,a,b,b)\vdash4n,\qquad a\ge b\ge0.
\] Then \(a+b=2n\), so \[
r=(a-b)/2=n-b\in\mathbf Z_{\ge0}.
\] Define \(P(q)=\#\{(u,v)\in\mathbf Z_{\ge0}^2:2u+3v=q\}\), with
\(P(q)=0\) for \(q<0\). The complete coefficient list is \[
\boxed{
\begin{array}{c|c}
\mu & p_{\mu,(n)}^{(a,a,b,b)}\\ \hline
(4)&P(r-3(b\bmod2))\\
(1,1,1,1)&P(r-3((b+1)\bmod2))\\
(2,2)&\lfloor(r+2)/3\rfloor\\
(3,1),(2,1,1)&0.
\end{array}}
\tag{P.1}
\] The two occurrences of \(b\) in the parity shifts can equivalently be
replaced by \(n\). The equality of those forms is proved below by an
explicit solution bijection. Every nonzero entry is a count in one
bounded two-variable rational polytope; the \((2,2)\) entry counts
\(u+3v=r-1\), \(u,v\ge0\). In particular its support is exactly
\(r\ge1\); each of the first two entries is positive exactly when its
argument is 0 or at least 2.

The symmetric-outer algebra on this face also has the following small
presentation: \[
\boxed{
A_4^{\{2\}}
\simeq\mathbf Q[e,D,Z,W]
/(27W^2+Z(D^2+4e^3)).}
\tag{P.2}
\] It is free with basis \(1,W\) over \(\mathbf Q[e,D,Z]\). This is a
different flag face from the equal-middle-row family T-024--T-028; its
exact relation to the same balanced Howe algebra is proved next. No
novelty claim is made for these classical small-outer-degree
coefficients.

\subsubsection{1. The Plücker morphism and exact torus
quotient}\label{the-pluxfccker-morphism-and-exact-torus-quotient}

Let \(X\) be a \(4\times4\) matrix of variables. Put \[
p_{ij}=X_{i1}X_{j2}-X_{i2}X_{j1}\quad(i<j),\qquad z=\det X,
\] and extend \(p_{ji}=-p_{ij}\). The flag subalgebra of right weights
\((a,a,b,b)\) is generated by the six \(p_{ij}\) and \(z\). Indeed these
are the left modules \(\Lambda^2\mathbf Q^4\) and \(\det\), and products
of their highest vectors give every permitted irreducible Cauchy
summand, exactly as in T-024 and T-029. The exact presentation is \[
C_{\{2\}}\simeq
\mathbf Q[p_{12},p_{13},p_{14},p_{23},p_{24},p_{34},z]
/(p_{12}p_{34}-p_{13}p_{24}+p_{14}p_{23}).
\tag{P.3}
\] The relation follows by expanding a decomposable two-vector. For the
converse, on the open set \(p_{12}\ne0\), the first two columns \[
c_1=(1,0,-p_{23}/p_{12},-p_{24}/p_{12})^\mathsf T,
\quad c_2=(0,p_{12},p_{13},p_{14})^\mathsf T
\] have precisely the prescribed minors, with \(p_{34}\) determined by
the displayed relation. Choose two additional columns with any desired
determinant, by first choosing a complement to
\(\langle c_1,c_2\rangle\) and then rescaling its last vector. Thus the
matrix morphism reaches this dense open set for arbitrary \(z\). The
quadratic is irreducible by Gauss's lemma, regarding it as linear in
\(p_{34}\), since its coefficient \(p_{12}\) is relatively prime to its
constant term. Density therefore proves that its ideal is the entire
kernel.

In a monomial \(\prod p_{ij}^{a_{ij}}z^t\), row \(i\) has degree
\(t+\sum_{j\ne i}a_{ij}\). Equality of the four row degrees is
equivalent to \[
a_{12}=a_{34},\qquad a_{13}=a_{24},\qquad a_{14}=a_{23}.
\tag{P.4}
\] For example, adding the first two row equations and subtracting the
last two yields \(2(a_{12}-a_{34})=0\); the other two equalities follow
similarly, and they also imply equal row sums. Hence the torus-invariant
polynomial ring before the Plücker quotient is freely generated by \[
t_1=p_{12}p_{34},\qquad
t_2=-p_{13}p_{24},\qquad
t_3=p_{14}p_{23},\qquad z.
\tag{P.5}
\] Their monomials have distinct exponent vectors in the original six
variables. The relation in (P.3) is torus-invariant, so the exact
weight-zero quotient gives \[
\boxed{B_{\{2\}}\simeq\mathbf Q[t_1,t_2,t_3,z]/(t_1+t_2+t_3).}
\tag{P.6}
\] Each \(t_i\) has balanced degree 1 and right weight
\(2\omega_2=(2,2,0,0)\), while \(z\) has degree 1 and weight
\(\omega_4\). Consequently the component \((n,(a,a,b,b))\) is exactly
\(z^b\) times the degree-\(r\) part in the \(t_i\).

\subsubsection{2. The permutation action and its exact
representation}\label{the-permutation-action-and-its-exact-representation}

The group \(S_4\) permutes the three perfect matchings
\(12|34,13|24,14|23\), giving a surjection \(\pi:S_4\to S_3\) with
kernel the Klein four group. Directly on the generators
\((12),(23),(34)\), the oriented products (P.5) transform by minus the
corresponding matching permutation. For instance, \[
(12):(t_1,t_2,t_3)\longmapsto(-t_1,-t_3,-t_2),
\qquad
(23):(t_1,t_2,t_3)\longmapsto(-t_2,-t_1,-t_3).
\] The action of \((34)\) on the triple is the same as that of \((12)\).
Thus for every \(\sigma\in S_4\), \[
\sigma t_i=\operatorname{sgn}(\sigma)t_{\pi(\sigma)(i)},
\qquad \sigma z=\operatorname{sgn}(\sigma)z,
\qquad\operatorname{sgn}(\pi(\sigma))=\operatorname{sgn}(\sigma).
\tag{P.7}
\] Let \(\mathfrak h_3\) be the ordinary two-dimensional standard
representation of \(S_3\). Equations (P.6)--(P.7) first give the actual
module \[
(S_{(a,a,b,b)}\mathbf Q^4)_{(n^4)}
\simeq\operatorname{Inf}_{S_3}^{S_4}
\left(\operatorname{sgn}^{b}\otimes
\operatorname{Sym}^{r}(\mathfrak h_3\otimes\operatorname{sgn})\right).
\tag{P.8}
\] There is an explicit rational intertwiner
\(J:\mathfrak h_3\otimes\operatorname{sgn}\to\mathfrak h_3\). Equip the
standard plane with its nondegenerate invariant dot product and any
nonzero alternating form \(\Omega\). Define \(J\) by
\(\langle Jv,w\rangle=\Omega(v,w)\). It is invertible, and the
determinant of the standard representation is the sign character, so
\(J\rho(g)=\operatorname{sgn}(g)\rho(g)J\). This is exactly the required
intertwining identity. Applying \(\operatorname{Sym}^rJ\) in (P.8) gives
\[
\boxed{
(S_{(a,a,b,b)}\mathbf Q^4)_{(n^4)}
\simeq\operatorname{Inf}_{S_3}^{S_4}
(\operatorname{sgn}^{b}\otimes\operatorname{Sym}^{r}\mathfrak h_3).}
\tag{P.9}
\] Alternatively, pulling the degree-\(r\) sign outside the symmetric
power in (P.8) gives the equally exact twist
\(\operatorname{sgn}^{r+b}=\operatorname{sgn}^{n}\). Both realizations
are related by the specified invertible intertwiner.

The irreducible \(S_4\)-modules factoring through \(\pi\) are the
trivial module \([(4)]\), the sign module \([(1^4)]\), and the
matching-plane module \([(2,2)]\). The last identification follows, for
instance, from its character \((2,0,2,-1,0)\) on the cycle types
\(1^4,2\,1^2,2^2,3\,1,4\), the character of \([(2,2)]\). The two
remaining Specht modules are the standard \([(3,1)]\) and its sign twist
\([(2,1,1)]\). The Klein four group acts transitively on the four
letters, so its fixed space in the four-dimensional permutation
representation consists only of constants; its fixed space in the
standard quotient is zero. Its sign is trivial, so the sign twist also
has no fixed vector. Consequently neither remaining Specht module can
occur in (P.9).

For \(S_3\), formula (S.16) of T-027 gives the multiplicity generating
functions in \(\operatorname{Sym}(\mathfrak h_3)\): \[
\frac1{(1-q^2)(1-q^3)},\qquad
\frac{q^3}{(1-q^2)(1-q^3)},\qquad
\frac{q+q^2}{(1-q^2)(1-q^3)}
=\frac q{(1-q)(1-q^3)}
\tag{P.10}
\] for the trivial, sign, and standard modules respectively. Apply
(P.9), then the refined balanced Howe identity (S.10). The first two
series give the parity shifts in (P.1), and the third gives
\(u+3v=r-1\), with \(\lfloor(r+2)/3\rfloor\) solutions. This proves
(P.1).

For a purely lattice verification of the replacement \(b\mapsto n\),
observe that \(n-b=r\). When \(r\) is even, the parities agree. When
\(r\) is odd, every solution of \(2u+3v=r\) has \(v\) odd and therefore
\(v\ge1\). The bijection \((u,v)\mapsto(u,v-1)\) proves \(P(r)=P(r-3)\),
including the empty small case. This is exactly the equality needed when
the two parity shifts differ.

\subsubsection{3. The invariant-algebra
presentation}\label{the-invariant-algebra-presentation}

In the trace-free triple (P.6), let \[
e=e_2(t_1,t_2,t_3),\quad f=e_3(t_1,t_2,t_3),
\quad D=(t_1-t_2)(t_1-t_3)(t_2-t_3).
\] The alternating subgroup of \(S_3\) fixes a free rank-two algebra
over \(\mathbf Q[e,f]\), generated by \(D\) with \[
D^2=-4e^3-27f^2.
\tag{P.11}
\] This is exactly the alternating/symmetric decomposition proved in
T-024, specialized to three variables. Under the odd signed permutation
action (P.7), \(e,D\) are fixed and \(f,z\) are negated. Eliminate
\(f^2\) using (P.11). The sign-fixed monomials form a free basis
\(1,zf\) over \(\mathbf Q[e,D,z^2]\). Setting \(Z=z^2\) and \(W=zf\)
gives precisely (P.2), including its unique monic quadratic relation in
\(W\) after division by 27.

The generator degrees \((n;\lambda)\) are \[
\begin{array}{c|c|c}
e&2&(4,4,0,0)\\
D&3&(6,6,0,0)\\
Z&2&(2,2,2,2)\\
W&4&(7,7,1,1).
\end{array}
\tag{P.12}
\] In a standard basis monomial \(e^uD^vZ^wW^\varepsilon\),
\(\varepsilon\in\{0,1\}\), the final part is \(b=2w+\varepsilon\), while
\(r=2u+3v+3\varepsilon\). This directly recovers the symmetric-outer
count \(P(r-3(b\bmod2))\), with an explicit basis bijection.

\subsection{Checks and relation to the global
question}\label{checks-and-relation-to-the-global-question}

The Plücker morphism, its dense chart, the exact torus quotient, the
signed matching action, the intertwiner \(J\), and the coefficient and
algebra presentations are all proved above. The only shared input is the
already specified Cauchy/Howe and Specht-module dictionary T-002--T-003.
This face admits a homogeneous polynomial base of dimension three,
consistent with its dimension in T-029; the global obstruction in T-029
arises from trying to choose one such base simultaneously across all
faces.

An independent calculation using the cycle-index expansion and the
Murnaghan--Nakayama rule passed all 140 coefficients for these
paired-row shapes and all five outer partitions through inner degree 6.
It is recorded in
\texttt{\small work/\allowbreak{}rectangular/\allowbreak{}OUTER4\_\allowbreak{}PAIRED\_\allowbreak{}AND\_\allowbreak{}PROPER\_\allowbreak{}CHECK.tsv},
with reproduction code
\texttt{\small work/\allowbreak{}rectangular/\allowbreak{}verify\_\allowbreak{}outer4\_\allowbreak{}faces.py}.
The independent agent also read and checked the complete proof above.

For precise historical attribution, T. Kuppel, \emph{Asymptotics of
plethysm}, Math. Z. \textbf{311}, article 84 (2025),
\href{https://link.springer.com/article/10.1007/s00209-025-03879-4}{Theorem
5, its proof, and Remark 2}, treats the paired-row exceptional family
and states that its occurring Specht modules factor through the matching
quotient \(S_4\to S_3\). Those passages were opened and read in this
run. The explicit signed-coordinate presentation (P.3)--(P.12) supplies
the algebra morphism and fixes all determinant shifts for the formulas
here.

\section{Direct positive formulas on every proper fundamental face of
outer degree
four}\label{direct-positive-formulas-on-every-proper-fundamental-face-of-outer-degree-four}

\subsection{T-031. Statement and explicit
formulas}\label{t-031.-statement-and-explicit-formulas}

Let \(\lambda=(\lambda_1,\lambda_2,\lambda_3,\lambda_4)\vdash4n\),
allowing trailing zero parts. Define the fixed binary-quartic count \[
\begin{aligned}
F_4(N,K)=\sum_{\varepsilon=0}^1\#\{(a,b,c,d)\in\mathbf Z_{\ge0}^4:\quad
&N=a+2b+2c+3d+3\varepsilon,\\
&K=2b+4c+6d+3\varepsilon\}.
\end{aligned}
\tag{F.1}
\] Set it equal to zero when \(N<0\), \(K<0\), or \(K>2N\). T-021--T-022
prove the exact algebra map and basis giving
\(F_4(N,K)=a_{4,N}^{(4N-K,K)}\) for valid indices. Each term in (F.1) is
one bounded four-variable polytope.

For the lower face \(\lambda_3=\lambda_4=t\), put \(N=n-t\),
\(K=\lambda_2-t\). Then \[
\boxed{
a_{4,n}^{\lambda}=
\begin{cases}
F_4(N,K),&t\text{ even},\\
F_4(N-3,K-6),&t\text{ odd}.
\end{cases}}
\tag{F.2}
\] The invalid-index convention includes every boundary case.

For the upper face \(\lambda_1=\lambda_2=H\), put \(N=H-n\),
\(K=H-\lambda_3\). Then \[
\boxed{
a_{4,n}^{\lambda}=
\begin{cases}
F_4(N,K),&H\text{ even},\\
F_4(N-3,K-6),&H\text{ odd}.
\end{cases}}
\tag{F.3}
\]

For the middle face \(\lambda_2=\lambda_3=\beta\), put \[
R=\min(\lambda_1-\beta,\beta-\lambda_4)-6(\beta\bmod2).
\] T-025 gives the single-polytope formula \[
\boxed{a_{4,n}^{\lambda}
=\#\{(c_2,c_3,c_4)\in\mathbf Z_{\ge0}^3:
2c_2+3c_3+4c_4=R\}.}
\tag{F.4}
\] Thus all proper faces of the three simple-root fundamental
coordinates of \(A_4\), with every determinant shift retained, have
direct bounded-polytope formulas. Indeed a proper face sets at least one
of \(\lambda_1-\lambda_2\), \(\lambda_2-\lambda_3\),
\(\lambda_3-\lambda_4\) equal to zero, exactly the three cases above.
The formulas agree on their intersections because the proofs identify
each with the same balanced Howe multiplicity. T-030 supplies another
explicit presentation on the paired-row intersection.

\subsection{1. Exact determinant shift}\label{exact-determinant-shift}

For any rank \(m\), a partition \(\eta\), and \(t\ge0\), the canonical
highest-weight pairing identifies \[
S_{\eta+(t^m)}E\simeq(\det E)^t\otimes S_\eta E.
\tag{F.5}
\] One construction is to append \(t\) full-height columns to the column
exterior-power realization of the Schur module. Those columns contribute
\((\Lambda^mE)^{\otimes t}\); the Cartan product with the original Schur
module is nonzero on its highest vector and has the displayed
irreducible highest weight, giving the isomorphism. On the balanced
weight \((n^m)\), this gives the exact \(S_m\)-module morphism \[
(S_{\eta+(t^m)}E)_{(n^m)}
\simeq\operatorname{sgn}^{t}\otimes
(S_\eta E)_{((n-t)^m)},
\tag{F.6}
\] because literal permutation matrices have determinant equal to their
sign.

Apply this with \(m=4\), \(\eta=(\lambda_1-t,\lambda_2-t,0,0)\), and
\(t=\lambda_3=\lambda_4\). Put \(A=\lambda_1-t\), \(B=\lambda_2-t\),
\(N=n-t\). The equality \(A+B=4N\) and the partition inequalities imply
\(N\ge0\) and \(0\le B\le2N\). By the refined balanced Howe identity
T-003, (F.6) identifies the desired coefficient with the multiplicity of
\(S_{(A,B)}\mathbf Q^2\) in \[
\operatorname{Sym}^4(\operatorname{Sym}^N\mathbf Q^2)
\quad(t\text{ even}),\qquad
\Lambda^4(\operatorname{Sym}^N\mathbf Q^2)
\quad(t\text{ odd}).
\tag{F.7}
\] The restriction to two variables is exact: the Schur coefficient for
a partition of length at most two is unchanged by evaluating the
character in two variables, while every partition of greater length
vanishes there.

\subsection{2. Binary alternants and the exact exterior-to-symmetric
map}\label{binary-alternants-and-the-exact-exterior-to-symmetric-map}

The following map works in every number \(m\) of binary slots, so we
state and prove it in that form. Let \(F=\mathbf Q^2\), and write
\(v_i=(a_i,b_i)\) for the variables in slot \(i\). Put \[
[ij]=a_i b_j-a_j b_i,\qquad
\Delta_m=\prod_{1\le i<j\le m}[ij],\qquad
M=\binom m2.
\tag{F.8}
\] The polynomial \(\Delta_m\) has degree \(m-1\) in every slot,
transforms by the sign under slot permutations, and spans the
\(GL(F)\)-character \((\det F)^M\).

Every alternating polynomial of equal degree \(N\) in all \(m\) slots is
divisible by \(\Delta_m\). To see this, fix a pair \(i,j\). On the dense
part of the bracket hypersurface where \(v_j=c v_i\), equal slot
homogeneity gives \[
f(\ldots,v_i,\ldots,c v_i,\ldots)
=c^N f(\ldots,v_i,\ldots,v_i,\ldots)=0,
\] since interchanging the equal slots negates an alternating
polynomial. The bracket \([ij]\) is an irreducible determinant
polynomial, so it divides \(f\). Distinct brackets are nonassociate
irreducibles in the polynomial unique factorization domain; therefore
their product divides \(f\). Its quotient has equal slot degree
\(N-m+1\), and is symmetric because numerator and denominator have the
same permutation sign. Conversely multiplying any symmetric polynomial
of that slot degree by \(\Delta_m\) produces an alternating polynomial
of degree \(N\) in each slot. If \(N<m-1\), divisibility leaves a
negative degree in each slot, so the alternating space is zero.

Multiplication by this explicit polynomial therefore gives the
\(GL(F)\)-equivariant isomorphism \[
\boxed{
\Lambda^m(\operatorname{Sym}^N F)
\simeq(\det F)^{\binom m2}\otimes
\operatorname{Sym}^m(\operatorname{Sym}^{N-m+1}F)
\quad(N\ge m-1),}
\tag{F.9}
\] and the left side is zero for \(N<m-1\). The identification of
alternating tensors with \(\Lambda^m\), and of symmetric tensors with
\(\operatorname{Sym}^m\), is the characteristic-zero symmetrizer
identification. The determinant factor in (F.9) records the actual
\(GL_2\) action of \(\Delta_m\), rather than only its degree.

For \(m=4\), (F.9) shifts the inner degree by 3 and both output parts by
6. Thus the odd case of (F.7) has multiplicity \[
[S_{(A-6,B-6)}F:\operatorname{Sym}^4(\operatorname{Sym}^{N-3}F)]
=F_4(N-3,B-6),
\tag{F.10}
\] with value zero when \(N<3\) or \(B<6\). If these inequalities hold,
\(A\ge B\) and \(A+B=4N\) automatically give a valid partition after the
shift. Equations (F.7)--(F.10) prove (F.2), including its boundary
conventions.

\subsection{3. Exact complement duality and the upper
face}\label{exact-complement-duality-and-the-upper-face}

For a rank-\(m\) space \(E\), a partition \(\lambda\) of length at most
\(m\), and an integer \(H\ge\lambda_1\), define its rectangular
complement \[
\widehat\lambda=(H-\lambda_m,\ldots,H-\lambda_1).
\tag{F.11}
\] There is a perfect equivariant pairing \[
S_\lambda E\otimes S_{\widehat\lambda}E\longrightarrow(\det E)^H.
\tag{F.12}
\] To construct it, place \(\lambda\) in an \(m\times H\) rectangle.
Pair each column exterior power \(\Lambda^j E\) with the complementary
\(\Lambda^{m-j}E\) by exterior multiplication into \(\det E\); reverse
the order of the complementary columns to recover the partition (F.11).
Applying the Schur/Cartan projections gives the equivariant pairing
(F.12). It is nonzero: a highest column tensor and the complementary
lowest column tensor wedge to a nonzero volume in every column. The two
irreducible modules are contragredient after the displayed determinant
twist, as their highest weights show, so Schur's lemma makes the induced
pairing perfect. Equivalently, \[
S_\lambda E\simeq(\det E)^H\otimes(S_{\widehat\lambda}E)^*.
\tag{F.13}
\] On balanced weights this gives the exact module identity \[
(S_\lambda E)_{(n^m)}
\simeq\operatorname{sgn}^{H}\otimes
\left((S_{\widehat\lambda}E)_{((H-n)^m)}\right)^*.
\tag{F.14}
\] Since Specht modules are self-dual and tensoring by sign conjugates
their partitions, T-003 yields the full coefficient equality \[
p_{\mu,(n)}^\lambda
=p_{\mu^{[H]},(H-n)}^{\widehat\lambda},
\qquad
\mu^{[H]}=\begin{cases}\mu,&H\text{ even},\\\mu',&H\text{ odd}.\end{cases}
\tag{F.15}
\] All quantities are polynomial representations:
\(H\ge\lambda_1\ge n\), and \(|\widehat\lambda|=m(H-n)\). Thus the
complement and the inner degree are nonnegative. Equation (F.14), not a
dimension-only substitution, explains the sign and duality in (F.15).

For the upper face of rank four, take \(H=\lambda_1=\lambda_2\). Then \[
\widehat\lambda=(H-\lambda_4,H-\lambda_3,0,0),
\qquad N=H-n,\quad K=H-\lambda_3.
\] Its first two parts sum to \(4N\). Formula (F.15) turns the original
symmetric outer coefficient into the symmetric binary coefficient when
\(H\) is even and the alternating binary coefficient when \(H\) is odd.
Apply (F.9) in the latter case. This proves (F.3) and all its edge
cases.

\subsection{Provenance and scope}\label{provenance-and-scope-1}

The binary-quartic presentation used in (F.1) is classical and is proved
with the exact normalizer-algebra morphism in T-021--T-022; the primary
Grace--Young locators and direct symbolic checks are recorded there.
Formula (F.9) is the classical binary Vandermonde factorization, proved
here to fix its slot grading, determinant power, and zero range.
Equations (F.5) and (F.12) are the actual Schur-module maps carrying the
formulas across determinant shifts and complementary partitions. The
result completes all proper fundamental faces of the symmetric-outer
\(A_4\) algebra. T-032--T-036 subsequently identify the full-face
extension through an exact flag chart and four-point vanishing model,
and T-037 completes the positive formula for every outer partition of
four.

An independent exact cycle-index/Murnaghan--Nakayama calculation passed
all 197 lower- and upper-face cases through inner degree 6. The receipt
\texttt{\small work/\allowbreak{}rectangular/\allowbreak{}OUTER4\_\allowbreak{}PAIRED\_\allowbreak{}AND\_\allowbreak{}PROPER\_\allowbreak{}CHECK.tsv}
also includes the 140 paired-row cases of T-030, for 337 total checks;
reproduction code is
\texttt{\small work/\allowbreak{}rectangular/\allowbreak{}verify\_\allowbreak{}outer4\_\allowbreak{}faces.py}.
The independent proof reading checked the divisibility map, every
zero-range convention, the complement pairing, and the literal
determinant signs. T-032--T-034 subsequently give direct chamber
formulas for every outer Schur label in \(\lambda_2\le n\) and
\(\lambda_3\ge n\).

\section{An exact stable flag chart for every outer
degree}\label{an-exact-stable-flag-chart-for-every-outer-degree}

\subsection{T-032. The algebra morphism and the balanced
module}\label{t-032.-the-algebra-morphism-and-the-balanced-module}

Work over \(\mathbf Q\), or extend scalars to any field of
characteristic zero. Let \(m\ge2\), let \(X=(X_{ik})\) be an
\(m\times m\) matrix of polynomial generators, and let \[
R_m=\operatorname{Sym}(\mathbf Q^m\otimes\mathbf Q^m)^{U},
\qquad B_m=R_m^{T_0},
\tag{C.1}
\] where \(U\) is the right maximal unipotent group fixing the
first-column flag minors and \(T_0\) is the determinant-one diagonal
group acting on the rows. The literal row permutation matrices give the
\(S_m\) action. The component of right dominant weight \(\lambda\) and
equal row weight \((n^m)\) is \[
(B_m)_{n,\lambda}=(S_\lambda\mathbf Q^m)_{(n^m)},
\qquad |\lambda|=mn.
\tag{C.2}
\] This is the same Cauchy/Howe identification as T-002--T-003. No
quotient by a determinant character is made in (C.2).

Put \(x_i=X_{i1}\), \(u=\prod_i x_i\), and \[
\mathfrak h_m=\{(a_1,\ldots,a_m):\sum_i a_i=0\}.
\] For \(j\ge2\), write \(p_{j,I}\) for the minor with row set \(I\),
\(|I|=j\), and the first \(j\) columns, in the natural increasing row
order. Define the polynomial \[
v_{j,I}=\left(\prod_{i\notin I}x_i\right)p_{j,I}.
\tag{C.3}
\] It has row weight \((1^m)\) and right weight
\((m-j+1,1^{j-1},0^{m-j})\). The following map is an algebra
isomorphism: \[
\boxed{
\mathbf Q[u,u^{-1}]\otimes
\operatorname{Sym}(\mathfrak h_m\otimes\mathbf Q^{m-1})^{U_{m-1}}
\ \xrightarrow{\ \Phi\ }\ B_m[u^{-1}].}
\tag{C.4}
\] The variable \(u\) maps to \(\prod_i x_i\); a first-\(\ell\)-column
flag minor on \(\mathfrak h_m\) maps to the coordinates of the
alternating tensor \(v_{\ell+1}\) under the explicit identification \[
\Lambda^\ell\mathfrak h_m\longrightarrow
\mathbf1\wedge\Lambda^\ell\mathfrak h_m\subset\Lambda^{\ell+1}\mathbf Q^m,
\qquad w\longmapsto\mathbf1\wedge w,
\quad\mathbf1=(1,\ldots,1).
\tag{C.5}
\] Thus (C.4) is \(S_m\)-equivariant, with \(u\) fixed and with the
ordinary standard action on \(\mathfrak h_m\).

For every partition \(\lambda\vdash mn\) of length at most \(m\), the
inverse chart embeds its balanced component into \[
u^{n-\lambda_2}S_{(\lambda_2,\ldots,\lambda_m)}\mathfrak h_m.
\tag{C.6}
\] In the entire chamber \(n\ge\lambda_2\), this embedding is onto.
Hence \[
\boxed{
(S_\lambda\mathbf Q^m)_{(n^m)}
\simeq S_{(\lambda_2,\ldots,\lambda_m)}\mathfrak h_m
\qquad(n\ge\lambda_2),}
\tag{C.7}
\] as actual \(S_m\)-modules. For any \(\mu\vdash m\), the refined Howe
identity therefore gives \[
\boxed{
p_{\mu,(n)}^\lambda
=\dim\operatorname{Hom}_{S_m}
\bigl([\mu],S_{(\lambda_2,\ldots,\lambda_m)}\mathfrak h_m\bigr)
\quad(n\ge\lambda_2).}
\tag{C.8}
\] Here the restriction of the Schur module on the right is along the
specified standard representation of \(S_m\), so all normalizer signs
remain part of the map.

\subsection{Proof of the chart}\label{proof-of-the-chart}

\subsubsection{1. Flag minors and centered
ratios}\label{flag-minors-and-centered-ratios}

The algebra \(R_m\) is generated by its first-column flag minors. One
way to see generation without imposing an unproved presentation is the
Cauchy decomposition \(R_m=\bigoplus_\lambda S_\lambda\mathbf Q^m\),
with the right highest-weight grading. Products of highest vectors in
the fundamental exterior-power summands give a nonzero highest vector in
every summand; the algebra generated by these summands is stable under
left \(GL_m\), so it contains every entire summand. The same argument
applies to the flag algebra on \(\mathfrak h_m\).

On \(u\ne0\), every \(x_i\) is invertible. For \(k\ge2\), define the
ratio vector and its centered projection by \[
w_k=(X_{1k}/x_1,\ldots,X_{mk}/x_m),\qquad
\bar w_k=w_k-\frac{1}{m}\left(\sum_i(w_k)_i\right)\mathbf1.
\tag{C.9}
\] Dividing the rows of each flag determinant by \(x_i\) gives \[
v_j=u\,\mathbf1\wedge w_2\wedge\cdots\wedge w_j
=\mathbf1\wedge
\bigl(u\bar w_2\wedge\bar w_3\wedge\cdots\wedge\bar w_j\bigr).
\tag{C.10}
\] In particular these are exactly, under (C.5), the flag minors of the
\((m-1)\)-column matrix \[
F=(u\bar w_2,\bar w_3,\ldots,\bar w_m)
\quad\text{with columns in }\mathfrak h_m.
\tag{C.11}
\] This proves that the proposed assignment respects every relation in
the smaller flag algebra. Its coordinates are the polynomials (C.3), so
it defines the displayed map even before localizing its smaller-flag
factor. Row permutation commutes with centering and fixes
\(\mathbf1,u\); this proves equivariance, including the orientation
signs in the exterior coordinates.

\subsubsection{2. Surjectivity after
localization}\label{surjectivity-after-localization}

Expand any element of \(R_m\) as a linear combination of monomials in
flag minors and then take its torus-weight-zero component. Such
monomials individually have equal row weights. Consider a balanced
monomial of right weight \(\lambda\) and row degree \(n\). Replace every
factor \(p_{j,I}\), \(j\ge2\), using \[
p_{j,I}=u^{-1}\left(\prod_{i\in I}x_i\right)v_{j,I}.
\tag{C.12}
\] The total number of flag factors of height at least two is
\(\lambda_2\): in a product of fundamental highest-weight degrees it is
\(\sum_{j\ge2}(\lambda_j-\lambda_{j+1})\). Every remaining \(x_i\)
exponent is \(n\), because the original monomial had row degree \(n\).
Therefore the result is precisely \[
u^{n-\lambda_2}\times\text{a monomial of smaller flag degree }
(\lambda_2,\ldots,\lambda_m).
\tag{C.13}
\] The same substitution applies after multiplying by any power of
\(u^{-1}\). Since \(u\) is invariant, taking the torus-weight-zero part
commutes with its localization. Equation (C.13) proves surjectivity in
(C.4) and the component assertion (C.6).

\subsubsection{3. Injectivity by an explicit dense
chart}\label{injectivity-by-an-explicit-dense-chart}

Prescribe a nonzero scalar \(u_0\) and any matrix
\(F=(f_1,\ldots,f_{m-1})\) with columns in \(\mathfrak h_m\). Choose
nonzero \(x_i\) with product \(u_0\), and set \[
X_{i1}=x_i,\qquad
X_{i2}=x_i(f_1)_i/u_0,\qquad
X_{ik}=x_i(f_{k-1})_i\quad(3\le k\le m).
\tag{C.14}
\] All ratio vectors are centered, and (C.11) recovers the prescribed
\(F\) exactly. Hence a Laurent polynomial in \(u\) with coefficients in
the smaller flag algebra that vanishes after \(\Phi\) vanishes on all
\((u_0,F)\). The smaller flag algebra is defined as its polynomial
subalgebra on such matrices, and nonzero scalars are Zariski dense. Thus
the Laurent polynomial was zero. This proves injectivity.

\subsubsection{4. The stable chamber and the exact extension
condition}\label{the-stable-chamber-and-the-exact-extension-condition}

The smaller flag component of degree \((\lambda_2,\ldots,\lambda_m)\) is
the irreducible Schur module displayed in (C.6). When
\(n-\lambda_2\ge0\), each element of this component maps to a polynomial
in the generators (C.3), multiplied by a nonnegative power of the
polynomial \(u\). Its right weight is \(\lambda\): the tail coordinates
are the smaller flag degree and the first coordinate is forced by total
size \(mn\). It therefore belongs to \((B_m)_{n,\lambda}\). This proves
surjectivity in (C.7); injectivity and equivariance follow from (C.4).

For completeness the same chart gives a precise, unconditional
description outside the chamber. Put \(k=\max(0,\lambda_2-n)\), and for
an element \(F\) of the tail Schur module use \(F(v(X))\) to denote its
polynomial image under (C.3) and (C.5). Then \[
\boxed{
(B_m)_{n,\lambda}\simeq
\{F\in S_{(\lambda_2,\ldots,\lambda_m)}\mathfrak h_m:
u^k\mid F(v(X))\text{ in }\mathbf Q[X_{ik}]\}.}
\tag{C.15}
\] The isomorphism sends \(F\) to \(u^{n-\lambda_2}F(v(X))\).
Polynomiality is exactly the indicated divisibility in a polynomial
ring. If the quotient is polynomial, its rational \(U\)-invariance and
balanced torus weight are equalities of rational functions and hence
polynomial identities; thus it belongs to \(B_m\). This proves both
directions. Divisibility does not depend on the chosen expression of
\(F\), since the smaller flag relations map to zero. Since \(u\) is
\(S_m\)-fixed, the subspace in (C.15) is \(S_m\)-stable.

\subsection{Scope, provenance, and the next finite-group
map}\label{scope-provenance-and-the-next-finite-group-map}

The chart supplies an explicit algebra morphism behind the classical
first-row stability phenomenon. The general stability literature
includes M. Brion, \emph{Stable properties of plethysm: on two
conjectures of Foulkes}, Manuscripta Math. \textbf{80} (1993), 347--371,
\href{https://doi.org/10.1007/BF03026558}{DOI 10.1007/BF03026558}. The
publisher abstract and bibliographic record were opened in this run; the
complete 1993 text was not accessible there. Consequently no unverified
exact theorem locator is being used to justify (C.4)--(C.15): the whole
algebra and component proof is given above. T. Kuppel, \emph{Asymptotics
of plethysm}, Math. Z. \textbf{311}, article 84 (2025),
\href{https://link.springer.com/article/10.1007/s00209-025-03879-4}{§2.3,
Theorem 3}, records Brion's §2.6, Corollary 1 as a general
coefficient-stability result; that quoted result concerns the stated
coefficient sequence, while (C.7) specifies the actual module and its
map.

The remaining map in the stable chamber is precisely restriction along
\(S_m\to GL(\mathfrak h_m)\). For \(m=4\), this factors through the
explicitly sign-twisted orthogonal representation of dimension three.
T-033--T-034 compute that restriction by classical \(GL_3\)-to-\(SO_3\)
branching and a five-class character calculation, yielding small
positive polytopes. T-035--T-037 then complete the extension outside the
chamber through four flag fibers and two quadratic generators; T-037
gives the full outer-four formula with at most108 fixed polytopes in at
most5 variables. For general \(m\), (C.15) states exactly what
polynomial extension imposes outside this chart chamber. G-009--G-011
already provide terminating positive counts in all degrees; the direct
uniform extension question after T-037 lies beyond outer degree four.

\section{Small positive chamber formulas in outer degree
four}\label{small-positive-chamber-formulas-in-outer-degree-four}

\subsection{T-033. The two exact branching
steps}\label{t-033.-the-two-exact-branching-steps}

Let \(W=\mathfrak h_4\otimes\operatorname{sgn}\), where
\(\mathfrak h_4\) is the three-dimensional standard representation of
\(S_4\). Its invariant dot product is nondegenerate, and
\(\det(W)=\operatorname{sgn}\cdot\operatorname{sgn}^3=1\). Thus its
actual representation map factors as \[
S_4\longrightarrow SO(W)\subset GL(W).
\tag{H.1}
\] Over \(\mathbf C\), write \(H_L\) for the irreducible \(SO_3\)-module
of angular momentum \(L\), with character \(\sum_{j=-L}^L z^j\) on
\(\operatorname{diag}(z,1,z^{-1})\).

For nonnegative integers \(c,b\), set \[
A=\max(c,b),\qquad s=\min(c,b).
\] The exact restriction of \(S_{(c+b,b)}W\) to \(SO(W)\) is the
following direct sum, with every pair \((K,L)\) providing one summand:
\[
\boxed{
S_{(c+b,b)}W\big|_{SO(W)}
\simeq
\bigoplus_{\substack{K=s,s-2,\ldots\\K>0}}
\ \bigoplus_{L=K}^{K+A}H_L
\ \oplus\
\begin{cases}
\displaystyle\bigoplus_{j=0}^{\lfloor A/2\rfloor}H_{A-2j},
&s\text{ even},\\[2mm]
0,&s\text{ odd}.
\end{cases}}
\tag{H.2}
\] This is the classical Elliott branching rule, proved below in the
precise convention used here.

For each \(\tau\vdash4\), define the finite tag sets \[
\begin{array}{c|c}
\tau&E_\tau\\ \hline
(4)&\{0,9\}\\
(1,1,1,1)&\{3,6\}\\
(2,2)&\{2,4,5,7\}\\
(3,1)&\{2,3,4,5,6,7\}\\
(2,1,1)&\{1,3,4,5,6,8\}.
\end{array}
\tag{H.3}
\] Then the second restriction in (H.1) has the exact positive
multiplicity formula \[
\boxed{
[\tau:H_L]
=\sum_{e\in E_\tau}
\#\{(x,y)\in\mathbf Z_{\ge0}^2:4x+6y=L-e\}.}
\tag{H.4}
\] An equation with negative right side contributes zero.

\subsubsection{1. Full character proof of the Elliott
rule}\label{full-character-proof-of-the-elliott-rule}

The three-variable Weyl determinant, specialized to \((z,1,z^{-1})\),
gives \[
\chi_{c,b}(z)
=z^{-(c+b)}
\frac{(1-z^{c+1})(1-z^{b+1})(1-z^{c+b+2})}
{(1-z)^2(1-z^2)}.
\tag{H.5}
\] For example this follows by dividing the alternant with exponents
\(c+b+2,b+1,0\) by the alternant with exponents \(2,1,0\), then
factoring their Vandermonde products after substitution. Formula (H.5)
is symmetric in \(c,b\), so one may put \(c=A,b=s\) in the ensuing
identity without changing the character.

Let \(P(z)\) be the sum of \(z^L\) over all the pairs \((K,L)\) in
(H.2), with repetitions retained. Finite geometric summation yields \[
P(z)=
\frac{z-z^{A+2}-z^{s+2}+z^{A+s+3}
+\delta(1-z)^2}{(1-z)(1-z^2)},
\quad
\delta=
\begin{cases}1,&A,s\text{ both even},\\0,&\text{otherwise}.\end{cases}
\tag{H.6}
\] Here the positive-\(K\) terms sum
\(\sum_{j=0}^{\lfloor(s-1)/2\rfloor}
z^{s-2j}(1-z^{A+1})/(1-z)\); when \(s\) is even, the final band
contributes \(\sum_{j=0}^{\lfloor A/2\rfloor}z^{A-2j}\). Adding these
two displayed finite geometric series gives (H.6), separately for the
parities of \(A,s\).

Now \[
P(z)-z^{-1}P(z^{-1})
=(1-z^{-1})\chi_{c,b}(z).
\tag{H.7}
\] For a direct verification, the \(\delta\) term in (H.6) equals
\(\delta/(1+z)\), whose contribution to the left side vanishes. The
remaining numerator of that side, with denominator \((1-z)(1-z^2)\), is
\[
-z^{A+2}-z^{s+2}+z^{A+s+3}
+z^{-A}+z^{-s}-z^{-A-s-1}.
\] Expanding the numerator of the right side using (H.5) gives exactly
the same six terms. Finally, \[
(1-z^{-1})\sum_{j=-L}^{L}z^j=z^L-z^{-L-1}.
\] Thus (H.7) says that (H.5) equals the sum of the characters \(H_L\)
in (H.2). Complete reducibility of the compact \(SO_3\) action, followed
by complexification, makes this equality of characters the asserted
module decomposition.

The primary locator for the classical rule is D. J. Rowe, \emph{Vector
coherent state theory: a powerful tool for solving algebraic problems in
physics}, submitted April 8, 2003,
\href{https://www.physics.utoronto.ca/site-archive/rowe/pdf/Erice.pdf}{§5,
Theorem 1 (Elliott), equations (57)--(58), printed page 7}. Those
displayed ranges were opened and read in this run. The full calculation
above specifies the symmetric \(A,s\) convention and also covers every
zero boundary.

\subsubsection{2. The harmonic trace and five exact character
averages}\label{the-harmonic-trace-and-five-exact-character-averages}

The polynomial quotient \[
H_L\simeq
\operatorname{Sym}^L W/
\bigl(Q\,\operatorname{Sym}^{L-2}W\bigr)
\tag{H.8}
\] is the stated irreducible \(SO(W)\)-module, where \(Q\) is its
invariant quadratic tensor and a negative symmetric power is zero. For a
direct justification, after complexification the double cover
\(SL_2\to SO_3\) identifies \(W\) with
\(\operatorname{Sym}^2\mathbf C^2\). One explicit realization is the
adjoint action on trace-zero \(2\times2\) matrices, whose determinant is
the invariant quadratic and whose kernel is \(\{\pm I\}\). A highest
vector in \(\operatorname{Sym}^L W\) has \(SL_2\)-weight \(2L\); its
\(L\)-th power is not divisible by the nondegenerate rank-three
quadratic \(Q\). It supplies the irreducible
\(\operatorname{Sym}^{2L}\mathbf C^2\) in the quotient. The quotient
dimension is \(\binom{L+2}{2}-\binom{L}{2}=2L+1\), exactly its
dimension. This proves (H.8) and the displayed character of \(H_L\).

For \(g\in S_4\), taking traces in (H.8) gives \[
\sum_{L\ge0}\operatorname{tr}(g\mid H_L)q^L
=\frac{1-q^2}{\det(1-qg\mid W)}.
\tag{H.9}
\] In the class order \(1^4,2\,1^2,2^2,3\,1,4\), the class sizes are
\(1,6,3,8,6\). The eigenvalues of the signed standard representation
\(W\) give the five trace series \[
\frac{1+q}{(1-q)^2},\quad
\frac1{1+q},\quad
\frac1{1+q},\quad
\frac{1+q}{1+q+q^2},\quad
\frac{1+q}{1+q^2}.
\tag{H.10}
\] The full character table in the same class order is \[
\begin{array}{c|rrrrr}
(4)&1&1&1&1&1\\
(1^4)&1&-1&1&1&-1\\
(2,2)&2&0&2&-1&0\\
(3,1)&3&1&-1&0&-1\\
(2,1,1)&3&-1&-1&0&1.
\end{array}
\tag{H.11}
\] These entries follow directly from the trivial and sign characters,
the four-letter permutation character minus its constant subspace, its
sign twist, and the three-matching permutation character minus its
constant subspace. Their squared dimensions sum to \(24\), and their
class inner products are the identity matrix, so these give all five
irreducible characters with the indicated usual Specht labels.

Multiply each row of (H.11) entrywise by (H.10), weight by its class
size, and divide by \(24\). A common-denominator calculation gives \[
\sum_{L\ge0}[\tau:H_L]q^L
=\frac{\sum_{e\in E_\tau}q^e}{(1-q^4)(1-q^6)}.
\tag{H.12}
\] Every numerator is listed explicitly in (H.3). Expanding the two
geometric denominators proves (H.4). This calculation is the
trace-average proof of the graded multiplicity version of Molien's
formula; see Stanley (1979), Theorem 2.1, printed pages 479--480,
already opened and checked in T-027. All five rational identities were
also verified exactly by the accompanying symbolic script; the finite
identity check supplements the derivation.

\subsection{T-034. The explicit chamber
polytopes}\label{t-034.-the-explicit-chamber-polytopes}

Let \(\mu\vdash4\), \(n\ge0\), and
\(\lambda=(\lambda_1,\lambda_2,\lambda_3,\lambda_4)\vdash4n\), with
\(\lambda_2\le n\). Set \[
c=\lambda_2-\lambda_3,\quad
b=\lambda_3-\lambda_4,\quad
d=\lambda_4,\quad A=\max(c,b),\quad s=\min(c,b),
\tag{H.13}
\] and put \(\tau=\mu\) when \(c+d\) is even and \(\tau=\mu'\) when it
is odd.

For every \(e\in E_\tau\), define the following bounded rational
polytopes, each with the three integer variables \((j,x,y)\): \[
\begin{aligned}
\mathcal P_e^+(A,s)
=\{(j,x,y)\in\mathbf R_{\ge0}^3:\quad&
2j\le s-1,\\
&s-2j\le4x+6y+e\le s-2j+A\},
\end{aligned}
\tag{H.14}
\] and \[
\mathcal P_e^0(A)=
\{(j,x,y)\in\mathbf R_{\ge0}^3:
2j+4x+6y=A-e\}.
\tag{H.15}
\] The first polytope is empty when \(s=0\). Negative right sides make
the second empty. Then \[
\boxed{
p_{\mu,(n)}^\lambda
=
\sum_{e\in E_\tau}
\#\bigl(\mathcal P_e^+(A,s)\cap\mathbf Z^3\bigr)
+\mathbf1_{s\ {\rm even}}
\sum_{e\in E_\tau}
\#\bigl(\mathcal P_e^0(A)\cap\mathbf Z^3\bigr).}
\tag{H.16}
\] This uses at most twelve fixed three-variable polytopes, without a
degree-dependent precomputation or a plethysm coefficient oracle. Their
tags are retained, so an integer point appearing in distinct polytopes
contributes once for each corresponding summand.

\subsubsection{3. Proof of the module transport and the lattice
bijection}\label{proof-of-the-module-transport-and-the-lattice-bijection}

The exact chart T-032 gives \[
(S_\lambda\mathbf Q^4)_{(n^4)}
\simeq S_{(\lambda_2,\lambda_3,\lambda_4)}\mathfrak h_4
\simeq(\det\mathfrak h_4)^d\otimes S_{(c+b,b)}\mathfrak h_4.
\tag{H.17}
\] The second map is the exterior-column determinant shift from (F.5).
Since \(\det\mathfrak h_4=\operatorname{sgn}\) and
\(\mathfrak h_4=W\otimes\operatorname{sgn}\), homogeneity of the Schur
functor gives the actual \(S_4\)-module isomorphism \[
(S_\lambda\mathbf Q^4)_{(n^4)}
\simeq \operatorname{sgn}^{c+d}\otimes S_{(c+b,b)}W.
\tag{H.18}
\] Indeed the latter Schur degree is \(c+2b\), whose parity is \(c\).
Tensoring a Specht module by sign conjugates its partition, so the
desired multiplicity is \([\tau:S_{(c+b,b)}W]\). Compose (H.2) and
(H.4).

For each positive \(K\) in (H.2), write \(K=s-2j\). Its index condition
is exactly \(j\ge0\), \(2j\le s-1\). For each harmonic tag \(e\), write
\(L=4x+6y+e\). The range \(K\le L\le K+A\) is then exactly (H.14). The
\(K=0\) band, present when \(s\) is even, has \(L=A-2j\); its equality
with \(4x+6y+e\) is (H.15). Every index in the two branching sums has
one and only one such tagged integer point, and the converse is
immediate from these equations. This proves (H.16).

Boundedness is explicit: the first polytope bounds \(j\) by \((s-1)/2\)
and \(4x+6y\) by \(A+s\); the second bounds all three nonnegative
coordinates by its right side. In particular their description and all
integer witnesses have bit length \(O(\log(n+2))\) for fixed outer
degree four.

\subsubsection{4. The other chamber and the exact extension problem
completed in
T-037}\label{the-other-chamber-and-the-exact-extension-problem-completed-in-t-037}

The complement map (F.12)--(F.15) covers every \(\lambda\) with
\(\lambda_3\ge n\). Explicitly take \[
H=\lambda_1,\qquad
\widehat\lambda=(H-\lambda_4,H-\lambda_3,H-\lambda_2,0),
\quad\widehat n=H-n,
\] and let \(\widehat\mu=\mu\) for even \(H\) and \(\widehat\mu=\mu'\)
for odd \(H\). Then \[
p_{\mu,(n)}^\lambda
=p_{\widehat\mu,(\widehat n)}^{\widehat\lambda}.
\tag{H.19}
\] The inequality \(\widehat\lambda_2\le\widehat n\) is exactly
\(\lambda_3\ge n\), so (H.16) applies to the right side. Formula (H.19)
is obtained from an actual perfect Schur-module pairing, including the
sign and contragredient maps; it also covers equality at the chamber
wall.

T-031 already gives direct formulas on every proper fundamental face for
the symmetric outer partition. At this stage of the reductions, the
symmetric-outer rank-four sector still requiring a direct chart
computation has \[
\lambda_1>\lambda_2>\lambda_3>\lambda_4,
\qquad \lambda_2>n>\lambda_3.
\tag{H.20}
\] Its balanced module is precisely the \(S_4\)-stable divisibility
subspace (C.15), with \(k=\lambda_2-n>0\), in the same rank-three flag
module used above. T-035 identifies this subspace with vanishing along
four flag fibers, and T-037 completes its positive decomposition for
every parameter. Thus the chamber formulas above are now part of the
full outer-four formula, with at most 108 fixed polytopes of at most
five variables for arbitrary outer Schur labels. General terminating
positive-polytope constructions remain available from G-009--G-011 in
every rank.

\subsection{Exact checks and
provenance}\label{exact-checks-and-provenance}

The companion script \url{check_chambers.py} verifies all five rational
identities (H.12), the geometric-sum and full Elliott character
identities for \(0\le c,b\le24\), their \(625\) Weyl dimensions, and
\(1280\) tagged lattice-count comparisons. Its exact receipt is
\url{CHAMBER_CHARACTER_CHECK.tsv}. These checks do not replace the
symbolic proofs above. Independent cycle-index plethysm checks are
recorded separately when completed.

The Elliott rule is classical and its precise primary locator is given
above. The symmetric-group character average is explicitly calculated
here. The purpose of the composition is the exact flag/normalizer
transport and fixed small positive polytope list in the stated chambers.
No historical novelty claim is made.

\section{Four fibers and a positive formula in the central
chamber}\label{four-fibers-and-a-positive-formula-in-the-central-chamber}

\subsection{T-035. The exact four-point vanishing
model}\label{t-035.-the-exact-four-point-vanishing-model}

Let \(\lambda\vdash4n\) have at most four parts, and put \[
c=\lambda_2-\lambda_3,\qquad
b=\lambda_3-\lambda_4,\qquad
d=\lambda_4,\qquad
k=\max(0,\lambda_2-n).
\tag{V.1}
\] Let \(W=[(2,1,1)]=\mathfrak h_4\otimes\operatorname{sgn}\), with its
invariant quadratic form, and let \[
\mathcal F=\{([q],[t])\in\mathbf P(W)\times\mathbf P(W):q\cdot t=0\},
\qquad \pi:\mathcal F\to\mathbf P(W).
\tag{V.2}
\] The two projective factors have the usual line-bundle notation
\(\mathcal O_{\mathcal F}(c,b)\). Use the invariant form to identify
\(W\) with its dual when taking polynomial sections. Let \(Z\) be the
four projective points \[
[\ell_i],\qquad
\ell_i=e_i-\tfrac14\mathbf1\in\mathfrak h_4,\quad 1\le i\le4.
\tag{V.3}
\] The sign twist does not change their projective action, so \(Z\) is
\(S_4\)-stable. Let \(\mathcal I_Z\) be their reduced ideal sheaf. Then
there is an actual \(S_4\)-module isomorphism \[
\boxed{
(S_\lambda\mathbf Q^4)_{(n^4)}
\simeq
\operatorname{sgn}^{c+d}\otimes
H^0\!\left(\mathcal F,\,
\pi^*(\mathcal I_Z^k)\otimes
\mathcal O_{\mathcal F}(c,b)\right).}
\tag{V.4}
\] The map is the polynomial chart of T-032 with its exact extension
condition. The proof below identifies its divisibility order with
vanishing order at each of the four fibers, including all higher orders.

The points in (V.3) have particularly simple equations. In the linear
coordinates \[
X=q_1+q_2,\qquad Y=q_1+q_3,\qquad Z_0=q_1+q_4
\tag{V.5}
\] on \(\mathfrak h_4\), they are the four projective sign choices
\([1:1:1]\), \([1:-1:-1]\), \([-1:1:-1]\), \([-1:-1:1]\). Their
homogeneous ideal is \[
J=(X^2-Y^2,\ Y^2-Z_0^2).
\tag{V.6}
\] It is a radical complete intersection of two quadrics. The generator
space \[
U=J_2
\tag{V.7}
\] is the Specht module \([(2,2)]\), and \(\det U=\operatorname{sgn}\).

\subsubsection{1. The explicit rank-four
chart}\label{the-explicit-rank-four-chart}

Write the first three columns of the original matrix as \(x_i,s_i,w_i\),
and let \(y_i\) be the oriented cofactor of its first three columns, so
\[
\sum_i x_i y_i=\sum_i s_i y_i=\sum_i w_i y_i=0.
\] As in T-032, \(u=\prod_i x_i\), and define the following polynomial
balanced coordinates: \[
q_i=u\left(\frac{s_i}{x_i}
-\frac14\sum_j\frac{s_j}{x_j}\right),\qquad
t_i=x_i y_i,\qquad z=\det X.
\tag{V.8}
\] Every denominator in each summand of \(q_i\) divides \(u\). These
satisfy \[
\sum_i q_i=\sum_i t_i=0,\qquad q\cdot t=0.
\tag{V.9}
\] The rank-three flag algebra in T-032 is \[
R=\mathbf Q[q\in\mathfrak h_4,\ t\in\mathfrak h_4\otimes\operatorname{sgn}]
/(q\cdot t),
\tag{V.10}
\] together with its determinant variable \(z\); its component of tail
degree is \(z^dR_{c,b}\). This incidence presentation follows directly
from the line-and-hyperplane flag description: a nonzero pair \(q,t\)
with \(q\cdot t=0\) is a line in its annihilating plane, and can be
realized by two independent columns of a rank-three matrix. The
incidence quadratic is irreducible, as a nondegenerate bilinear form of
rank three. Its dense open set therefore gives the entire kernel. This
also verifies the presentation directly, in addition to (C.4).

Here \(q\) transforms as \(\mathfrak h_4\), \(t\) as \(W\), and \(z\) by
sign. Retwist the action on each \(q\)-generator by sign. Both \(q\) and
\(t\) then transform as \(W\), the quadratic relation becomes invariant,
and a component of bidegree \((c,b)\) changes by the scalar character
\(\operatorname{sgn}^c\). This is an explicitly defined automorphism
action on the same algebra. Including \(z^d\) gives exactly the twist
\(\operatorname{sgn}^{c+d}\) in (V.4).

The complete incidence ring has \[
R_{c,b}=H^0(\mathcal F,\mathcal O_{\mathcal F}(c,b))
\quad(c,b\ge0).
\tag{V.11}
\] For instance, use the exact sequence of the incidence hypersurface in
\(\mathbf P^2\times\mathbf P^2\), whose equation has bidegree \((1,1)\).
The first cohomology of \(\mathcal O(c-1,b-1)\) on this product vanishes
for these indices, so its global section quotient is exactly (V.10).
Equivalently the same equality is the rank-three flag Cauchy component
of T-032.

\subsubsection{2. Equality of the two orders of
vanishing}\label{equality-of-the-two-orders-of-vanishing}

It is enough to treat \(i=4\), since row permutation acts transitively
on the four indices. At a general point of the divisor \(x_4=0\), the
other \(x_i\) and \(s_4\) are nonzero. Equations (V.8) give \[
q\in\mathbf Q^\times\ell_4,\qquad t_4=0,
\] so the image is the general point of the fiber
\(\mathcal F_4=\pi^{-1}([\ell_4])\).

We verify that the valuation of a pulled-back section along \(x_4=0\) is
exactly its order in the ideal of \(\mathcal F_4\). Put \[
x_4=\epsilon,\quad s_4=S\ne0,\quad r_j=s_j/x_j\ (j=1,2,3).
\] Use the affine chart \(q_4-q_1\ne0\) of the first projective factor.
The ideal of the point \([\ell_4]\) there is generated by \[
a_2=\frac{q_2-q_1}{q_4-q_1},\qquad
a_3=\frac{q_3-q_1}{q_4-q_1}.
\] Their exact pullbacks are \[
a_j=\epsilon\,\frac{r_j-r_1}{S-\epsilon r_1},
\qquad j=2,3.
\tag{V.12}
\] The map \(\pi\) is a projective-line bundle, so \(a_2,a_3\) are also
regular parameters transverse to \(\mathcal F_4\) at its generic point.

The two leading coefficients \((r_2-r_1)/S,(r_3-r_1)/S\) in (V.12) are
algebraically independent over the function field of the image fiber.
This can be checked inside the original matrix without a genericity
assumption about stabilizers. Restrict to the subfamily \(w_4=0\), keep
\(x_1,x_2,x_3,S\) nonzero, and vary \(w_1,w_2,w_3\). At \(\epsilon=0\),
the cofactors \(y_j\), \(j\le3\), expand along the fourth row
\((0,S,0)\); hence the resulting \(t_j\) depend on
\(S,x_1,x_2,x_3,w_1,w_2,w_3\) and not on \(r_1,r_2,r_3\). They range
over the entire plane \(t_1+t_2+t_3=0\): after dividing by the nonzero
factors \(x_i,S\), they are the pairwise differences of \(w_j/x_j\),
with the usual cyclic signs. They therefore cover a dense open of the
projective fiber independently of the freely chosen differences
\(r_2-r_1,r_3-r_1\). A fourth column makes the determinant \(z\) any
desired nonzero value on this same dense set.

Now take a nonzero section, trivialized on this chart and a nonzero
coordinate chart of the second factor. If its first nonzero transverse
term has total degree \(h\) in \(a_2,a_3\), that term lies in \[
\operatorname{gr}^h_{\mathcal F_4}\simeq
\mathbf Q(\mathcal F_4)[a_2,a_3]_h.
\] After (V.12), its leading coefficient of \(\epsilon^h\) is a nonzero
homogeneous polynomial in the two algebraically independent leading
coefficients. It cannot cancel. The coordinate trivializations have
valuation zero at a generic point of the divisor. Thus \[
\operatorname{ord}_{x_4=0}(F(q(X),t(X),z(X)))
=\operatorname{ord}_{\mathcal F_4}(F).
\tag{V.13}
\] Multiplication by \(z^d\) has valuation zero there, as just verified.
The same argument applies to each \(i\).

The divisor factors \(x_i\) are pairwise nonassociate primes in the
original polynomial ring. Therefore divisibility by \(u^k\) is
equivalent to order at least \(k\) along every \(x_i=0\). By (V.13),
this is exactly order at least \(k\) along all four fibers. The
normal-jet quotients of each smooth fiber are locally free over that
integral fiber. A section of such a quotient that vanishes at its
generic point vanishes everywhere, so the generic order condition is the
full sheaf-ideal-power condition. The fibers are disjoint, and \(\pi\)
is smooth, so their ideal powers are precisely the pullbacks of the four
point-ideal powers, or equivalently of \(\mathcal I_Z^k\). This proves
(V.4) using the exact divisibility identity (C.15).

\subsubsection{3. The four-point ideal and its equivariant
resolution}\label{the-four-point-ideal-and-its-equivariant-resolution}

No projective point of the scheme (V.6) has \(Z_0=0\), since its
equations would force \(X=Y=0\). On \(Z_0=1\) the equations give
\(X^2=Y^2=1\), four reduced points in characteristic zero. These are
exactly (V.3), so (V.6) is its saturated homogeneous ideal. The quadrics
have no common factor and form a regular sequence.

The \(S_4\) action on \(X,Y,Z_0\) is by signed permutation matrices. On
their squares it is the permutation action on the three perfect
matchings of four letters. Thus their trace-free difference space \(U\)
is the inflated standard representation of \(S_3\), namely \([(2,2)]\),
with determinant sign. This is also seen by taking the kernel of
evaluation \(\operatorname{Sym}^2W\to\mathbf Q^4\), which is surjective
on these four points.

For \(k\ge1\), the \(k\)-th power of this two-generator
complete-intersection ideal has the exact equivariant resolution \[
0\longrightarrow
\operatorname{Sym}^{k-1}U\otimes\det U\otimes
\mathcal O_{\mathbf P(W)}(-2k-2)
\longrightarrow
\operatorname{Sym}^{k}U\otimes
\mathcal O_{\mathbf P(W)}(-2k)
\longrightarrow
\mathcal I_Z^k\longrightarrow0.
\tag{V.14}
\] To verify exactness directly, choose generators \(f,g\) of \(U\), and
let \(e_j\) denote the formal basis element mapped to \(f^{k-j}g^j\).
The final map sends these \(k+1\) formal monomials to their products in
the polynomial ring. Its \(k\) relations are \(g e_j-f e_{j+1}\),
\(0\le j<k\). They are independent over the polynomial ring and generate
all relations: reduce a relation modulo \(f\), use that \(g\) is a
nonzerodivisor modulo \(f\), subtract the last adjacent relation, and
descend on the number of terms. This yields exactness. The invariant
form of that map is \[
h\otimes(f\wedge g)\longmapsto
(hg)\otimes f-(hf)\otimes g,
\] which proves the \(\det U\) twist and full equivariance.

\subsection{T-036. A positive module formula when the opposite gaps
agree}\label{t-036.-a-positive-module-formula-when-the-opposite-gaps-agree}

Let \(b,k,d\ge0\), and set \[
\lambda=(d+2b+2k,\ d+b+2k,\ d+b,\ d),\qquad n=d+b+k.
\tag{V.15}
\] Equivalently the two opposite root gaps satisfy
\(\lambda_1-\lambda_2=\lambda_3-\lambda_4=b\), and the middle gap is
\(2k\). Then \[
\boxed{
(S_\lambda\mathbf Q^4)_{(n^4)}
\simeq\operatorname{sgn}^{d}\otimes
\left(
\operatorname{Sym}^{k}U\otimes\operatorname{Sym}^{b}W
\ \oplus\
\operatorname{sgn}\otimes
\operatorname{Sym}^{k-1}U\otimes\operatorname{Sym}^{b-1}W
\right).}
\tag{V.16}
\] A negative symmetric power means the zero module. In particular this
formula includes every boundary \(b=0\) or \(k=0\), and also every
strictly interior central-chamber member \(b,k>0\) of this family.

\subsubsection{4. Proof of the positive boundary cohomology
term}\label{proof-of-the-positive-boundary-cohomology-term}

Let \(\mathcal Q\) be the rank-two quotient bundle in \[
0\to\mathcal O_{\mathbf P(W)}(-1)\to
W\otimes\mathcal O_{\mathbf P(W)}\to\mathcal Q\to0.
\tag{V.17}
\] The fiber of \(\pi\) over a line \(q\) consists of lines in its
annihilator. The invariant quadratic form identifies the dual of that
annihilator with \(W/q\). Hence \[
\pi_*\mathcal O_{\mathcal F}(0,b)=\operatorname{Sym}^b\mathcal Q,
\qquad
R^1\pi_*\mathcal O_{\mathcal F}(0,b)=0
\quad(b\ge0).
\tag{V.18}
\] This is the polynomial degree-\(b\) description on each
projective-line fiber, and it is compatible with the specified group
action.

Taking symmetric powers of (V.17), whose kernel has rank one, gives \[
0\to\mathcal O(-1)\otimes\operatorname{Sym}^{b-1}W
\to\mathcal O\otimes\operatorname{Sym}^{b}W
\to\operatorname{Sym}^{b}\mathcal Q\to0.
\tag{V.19}
\] The first map multiplies by the tautological vector. Locally the
original line inclusion splits, so tensoring (V.19) with
\(\mathcal I_Z^k(c)\) preserves its exactness.

In the family (V.15), \(c=2k\). The resolution (V.14), together with the
elementary cohomology of line bundles on \(\mathbf P^2\), gives \[
\begin{aligned}
H^0(\mathcal I_Z^k(2k))&=\operatorname{Sym}^k U,\\
H^0(\mathcal I_Z^k(2k-1))&=0,\\
H^1(\mathcal I_Z^k(2k))&=0,\\
H^1(\mathcal I_Z^k(2k-1))
&=\operatorname{Sym}^{k-1}U\otimes\det U.
\end{aligned}
\tag{V.20}
\] Indeed the two resolution twists are \((-2,0)\) in the first and
third lines and \((-3,-1)\) in the second and fourth. On
\(\mathbf P^2\), \(H^1(\mathcal O(a))=0\),
\(H^2(\mathcal O(-2))=H^2(\mathcal O(-1))=0\), and
\(H^2(\mathcal O(-3))\) is one-dimensional. Its \(S_4\) character is
trivial because \(\det W=1\). The same four equalities hold for \(k=0\)
with the negative symmetric power interpreted as zero, directly from
\(\mathcal I_Z^0=\mathcal O\).

The cohomology sequence of (V.19) now gives the short exact sequence \[
0\to
\operatorname{Sym}^kU\otimes\operatorname{Sym}^bW
\to
H^0(\mathcal I_Z^k(2k)\otimes\operatorname{Sym}^b\mathcal Q)
\to
\det U\otimes\operatorname{Sym}^{k-1}U\otimes
\operatorname{Sym}^{b-1}W
\to0.
\tag{V.21}
\] Every rational \(S_4\)-module is semisimple: averaging a projection
over its 24 elements gives an equivariant splitting. Thus (V.21) is the
direct sum of its end terms as an \(S_4\)-module. Apply (V.18), then
(V.4), where \(c=2k\) is even and \(\det U=\operatorname{sgn}\). This
proves (V.16). The splitting need not be canonical; the injection,
quotient, and resulting isomorphism type are exact.

\subsection{5. Fixed positive polytopes for every outer Schur
label}\label{fixed-positive-polytopes-for-every-outer-schur-label}

Let \(F_{(4)}=\{0\}\), \(F_{(1^4)}=\{3\}\), and \(F_{(2,2)}=\{1,2\}\).
The complete symmetric-power decomposition of \(U\) is \[
[\rho:\operatorname{Sym}^kU]
=\sum_{f\in F_\rho}
\#\{(a_2,a_3)\in\mathbf Z_{\ge0}^2:2a_2+3a_3=k-f\},
\tag{V.22}
\] for \(\rho=(4),(1^4),(2,2)\), with the other two multiplicities zero.
This is the \(S_3\) formula (P.10) transported along the same matching
quotient.

By the exact harmonic quotient (H.8), \[
[\tau:\operatorname{Sym}^{b}W]
=\sum_{e\in E_\tau}
\#\{(r_2,r_4,r_6)\in\mathbf Z_{\ge0}^3:
2r_2+4r_4+6r_6=b-e\}.
\tag{V.23}
\] The tag sets \(E_\tau\) are listed in (H.3).

The necessary tensor products have multiplicity one. Multiplication by
\((4)\) leaves the label unchanged, multiplication by \((1^4)\)
conjugates it, and multiplication by \((2,2)\) has \[
\begin{array}{c|c}
\tau&[(2,2)]\otimes[\tau]\\ \hline
(4),(1^4)&[(2,2)]\\
(2,2)&[(4)]\oplus[(1^4)]\oplus[(2,2)]\\
(3,1),(2,1,1)&[(3,1)]\oplus[(2,1,1)].
\end{array}
\tag{V.24}
\] Each equality follows by entrywise multiplication of the character
rows (H.11).

For a label \(\nu\vdash4\), define \(\mathcal C_\nu(k,b)\) to be the
following explicit sum. For every \(\rho\in\{(4),(1^4),(2,2)\}\) and
\(\tau\vdash4\) whose tensor product in (V.24) contains \(\nu\), and
every \(f\in F_\rho,e\in E_\tau\), count the nonnegative integer
solutions of \[
2a_2+3a_3=k-f,\qquad
2r_2+4r_4+6r_6=b-e.
\tag{V.25}
\] Each tag contributes one bounded five-variable polytope. Set the
whole count to zero if \(k<0\) or \(b<0\). Write \(\mu^{[d]}=\mu\) for
even \(d\) and \(\mu'\) for odd \(d\). Then \[
\boxed{
p_{\mu,(n)}^\lambda
=\mathcal C_{\mu^{[d]}}(k,b)
+\mathcal C_{(\mu^{[d]})'}(k-1,b-1).}
\tag{V.26}
\] This is (V.16) with every finite-group and symmetric-power
multiplicity replaced by its explicit nonnegative count.

For any fixed \(\nu\), the first count has at most 36 tagged polytopes:
the maximum occurs for \(\nu=(3,1)\) or \((2,1,1)\), with
\(6+6+2(6+6)=36\) tags. Thus (V.26) uses at most 72 fixed five-variable
polytopes for all five outer labels and all \(b,k,d\). No
degree-dependent invariant-ring computation occurs. The inequalities are
just nonnegativity and the two displayed equations; each witness has bit
length \(O(\log(n+2))\).

\subsection{Provenance and remaining
scope}\label{provenance-and-remaining-scope}

The exact morphisms used here are the chart (C.4), the transverse
equations (V.12), the complete-intersection resolution (V.14), and the
tautological quotient map (V.19). Their proofs are included in full. For
the elementary projective-space cohomology in (V.20), the primary
reference is \href{https://stacks.math.columbia.edu/tag/01XS}{Stacks
Project, Section 30.8, Lemma 30.8.1, tag 01XS}; its complete ordered
Čech proof was opened and read in this run. The determinant character is
fixed here by the specified \(SO(W)\) action. No numerical plethysm
value is used as input.

An independent cycle-index/Murnaghan--Nakayama calculation passed all
420 coefficient cases in (V.15) with \(n\le6\), every outer partition of
four, and all nonnegative \(b,k,d\); 175 are strict middle-chamber
cases. The independent side computed the right-hand module from
reciprocal characteristic-polynomial series, without the fat-point
derivation. Receipt:
\texttt{\small work/\allowbreak{}rectangular/\allowbreak{}OUTER4\_\allowbreak{}SELF\_\allowbreak{}COMPLEMENT\_\allowbreak{}INDEPENDENT\_\allowbreak{}CHECK.tsv};
code:
\texttt{\small work/\allowbreak{}rectangular/\allowbreak{}verify\_\allowbreak{}outer4\_\allowbreak{}self\_\allowbreak{}complement.py}.
A separate independent reading checked the complete valuation,
resolution, cohomology, and polytope proofs.

The family includes the paired-row case T-030 when \(b=0\), and the
equal-middle-row case T-027 when \(k=0\). For \(b,k>0\), it supplies a
direct positive formula inside the strict central chamber left by T-034.
The general four-point section model (V.4) also applies to unequal
opposite gaps. T-037 subsequently completes those cases by separating
the two-quadric ideal powers inside the incidence ring and gives a fixed
positive polytope formula for every outer partition of four.
G-009--G-011 supply terminating positive counts in every rank; the next
efficient direct-formula question in this route lies beyond outer degree
four.

\section{A fixed positive polytope formula for every outer Schur
partition of
four}\label{a-fixed-positive-polytope-formula-for-every-outer-schur-partition-of-four}

\subsection{T-037. Full statement and
conventions}\label{t-037.-full-statement-and-conventions}

Let \(\mu\vdash4\), \(n\ge0\), and \(\lambda\vdash4n\) have at most four
parts. There is an explicit formula for \(p_{\mu,(n)}^\lambda\) using at
most \textbf{108 tagged bounded rational polytopes, each with at most
five variables}. For the symmetric or exterior outer partition the bound
is 36, and for \(\mu=(2,2)\) it is 72. All tags and equations are listed
in §5. No character evaluation, invariant-ring elimination, or
degree-dependent precomputation occurs in this formula.

For a valid output partition with more than four parts, the coefficient
is zero. Indeed the Schur symmetrizer realizes
\(S_\mu(\operatorname{Sym}^nV)\) inside
\((\operatorname{Sym}^nV)^{\otimes4}\). Starting with the one-row
representation \(\operatorname{Sym}^nV\), each of the next three tensor
factors adds a horizontal strip by Pieri's rule, increasing the number
of rows by at most one. Thus every output in that ambient tensor product
has at most four rows. The evaluator includes this zero case and rejects
inputs of the wrong size or order.

The proof uses the original action \[
q\in\mathfrak h_4=[(3,1)],\qquad
t\in W=\mathfrak h_4\otimes\operatorname{sgn}=[(2,1,1)],
\qquad
R=\mathbf Q[q,t]/(q\cdot t),
\tag{A.1}
\] with bidegrees \((1,0),(0,1)\). Thus the incidence equation
transforms by sign. We do not retwist the \(q\)-action in the separation
of variables below. This convention makes its evaluation at the four
affine vectors \(\ell_i=e_i-\frac14\mathbf1\) the ordinary permutation
representation in every degree at least two.

Let \(U=[(2,2)]\subset\operatorname{Sym}^2\mathfrak h_4\) be the two
quadrics vanishing at the four points (V.3), and put \[
J_R=UR,\qquad \mathcal H=R/J_R.
\tag{A.2}
\] Then multiplication, after the explicitly specified equivariant
choice of sections in §1, gives the graded \(S_4\)-module isomorphism \[
\boxed{
R\simeq\operatorname{Sym}(U)\otimes\mathcal H.}
\tag{A.3}
\] The degree of \(U\) is \((2,0)\). In particular \[
\boxed{
(J_R^k)_{c,b}
\simeq\bigoplus_{j=k}^{\lfloor c/2\rfloor}
\operatorname{Sym}^j U\otimes\mathcal H_{c-2j,b}.}
\tag{A.4}
\] The finite-group modules in this decomposition are explicit: \[
\boxed{
\begin{array}{c|l}
C&\mathcal H_{C,b}\\ \hline
0&\operatorname{Sym}^b W\\
1&\displaystyle\operatorname{sgn}\otimes
       \bigoplus_{L=1}^{b+1}H_L\\
2&\displaystyle D_b\ \oplus\
       \operatorname{sgn}\otimes\operatorname{Sym}^{b-1}W\\
C\ge3&D_b,
\end{array}}
\qquad
D_b=\operatorname{Ind}_{S_3}^{S_4}\operatorname{Sym}^b\mathfrak h_3.
\tag{A.5}
\] A negative symmetric power is zero, and \(H_L\) is the \(SO(W)\)
harmonic module restricted along the actual rotational \(S_4\) map
(H.1).

\subsubsection{Orienting the partition by the exact complement
map}\label{orienting-the-partition-by-the-exact-complement-map}

Set \(a=\lambda_1-\lambda_2\) and \(b=\lambda_3-\lambda_4\). When
\(a<b\), replace the input by \[
H=\lambda_1,\quad
\lambda\longleftarrow(H-\lambda_4,H-\lambda_3,H-\lambda_2,0),
\quad n\longleftarrow H-n,
\quad \mu\longleftarrow\mu^{[H]},
\tag{A.6}
\] where \(\eta^{[H]}=\eta\) for even \(H\) and \(\eta'\) for odd \(H\).
The exact perfect pairing (F.12)--(F.15) proves that this replacement
leaves the coefficient unchanged. It interchanges \(a,b\), so after at
most one replacement \(a\ge b\).

For the resulting input put \[
a=\lambda_1-\lambda_2,\qquad b=\lambda_3-\lambda_4,\qquad
c=\lambda_2-\lambda_3,\qquad d=\lambda_4,\qquad
k=\max(0,\lambda_2-n).
\tag{A.7}
\] Then \(c\ge2k\). For \(k=0\) this is immediate. For \(k>0\), the size
equation \(|\lambda|=4n\) gives \[
c-2k=\frac{a-b}{2}\ge0.
\tag{A.8}
\] The exact four-fiber model T-035 now gives the following positive
module decomposition: \[
\boxed{
(S_\lambda\mathbf Q^4)_{(n^4)}
\simeq \operatorname{sgn}^d\otimes
\left(
\bigoplus_{j=k}^{\lfloor c/2\rfloor}
\operatorname{Sym}^jU\otimes\mathcal H_{c-2j,b}
\ \oplus\
\mathbf1_{c=2k}\,
\operatorname{sgn}\otimes
\operatorname{Sym}^{k-1}U\otimes\operatorname{Sym}^{b-1}W
\right).}
\tag{A.9}
\] The term with \(c=2k\) is the fully proved boundary term of T-036.
The proof of (A.3)--(A.9) follows before the explicit polytope formula.

\subsection{1. Four lines, a nonzerodivisor, and actual module
freeness}\label{four-lines-a-nonzerodivisor-and-actual-module-freeness}

Write \[
S=\mathbf Q[q],\qquad A=S/(U).
\] In the coordinates (V.5), \[
A=\mathbf Q[X,Y,Z_0]/(X^2-Y^2,\ Y^2-Z_0^2).
\tag{A.10}
\] This is the reduced coordinate ring of the four affine lines through
the four projective points. Indeed its two equations form a complete
intersection, the four projective points are reduced by the calculation
in T-035, and the same ring is free over \(\mathbf Q[Z_0]\) with basis
\(1,X,Y,XY\), using the monic equations \(X^2=Z_0^2,Y^2=Z_0^2\). It
embeds in the product of the four line rings: localizing at \(Z_0\)
makes that claim the four distinct roots \(X/Z_0,Y/Z_0=\pm1\);
injectivity before localization follows from its freeness over
\(\mathbf Q[Z_0]\). In particular no component supported at the origin
is being discarded.

The image of \(\ell=q\cdot t\) is a nonzerodivisor in \(A[t]\). On each
of the four line components it is a nonzero scalar line coordinate times
the nonzero linear form \(\ell_i\cdot t\), so it is a nonzerodivisor in
every component domain. The just-proved embedding into their product
makes multiplication injective on \(A[t]\) itself.

Consequently the ordinary bigraded Hilbert series satisfy \[
\operatorname{Hilb}(\mathcal H;s,t)
=\frac{(1-st)(1-s^2)^2}{(1-s)^3(1-t)^3}
=(1-s^2)^2\operatorname{Hilb}(R;s,t).
\tag{A.11}
\] The first equality follows from the two regular quadrics in \(S\) and
the injective multiplication by \(\ell\) just proved; the second follows
from the one nonzero incidence equation in the polynomial ring.

For each bidegree, choose a rational linear right inverse \(s_0\) to
\(R\to\mathcal H\) and replace it by \[
s(h)=\frac1{24}\sum_{g\in S_4}g\,s_0(g^{-1}h).
\tag{A.12}
\] This is an actual \(S_4\)-equivariant right inverse. Define the
multiplication map \[
\operatorname{Sym}(U)\otimes\mathcal H\to R,\qquad
p\otimes h\longmapsto p(U)\,s(h).
\tag{A.13}
\] It is surjective: subtract the chosen lift of a quotient class,
express the remainder in \(UR\), and repeat with strictly smaller
\(q\)-degree. The source has the same Hilbert series as the target by
(A.11). Each bidegree is finite-dimensional, so (A.13) is injective in
each bidegree. This proves (A.3) as an equivariant
\(\operatorname{Sym}(U)\)-module isomorphism. The ideal powers
correspond to polynomials in \(U\) of total degree at least \(k\),
proving (A.4). In particular the two elements of \(U\) form a regular
sequence in \(R\).

The section (A.12) proves existence of the specified multiplication
isomorphism; the explicit counts in §5 require neither its computation
nor any averaging of degree-dependent vectors.

\subsection{\texorpdfstring{2. Every quotient piece
\(\mathcal H_{C,b}\)}{2. Every quotient piece \textbackslash mathcal H\_\{C,b\}}}\label{every-quotient-piece-mathcal-h_cb}

The evaluation map on the four affine vectors \(\ell_i\) gives \[
A_0=\mathbf1,\qquad A_1=\mathfrak h_4,\qquad
A_C\simeq\operatorname{Perm}_4\quad(C\ge2).
\tag{A.14}
\] Here evaluation has no parity twist, since the original standard
action sends \(\ell_i\) exactly to \(\ell_{\sigma(i)}\). For
completeness, in the coordinates (V.5), the four functions
\(X^2,XY,XZ_0,YZ_0\) have linearly independent evaluations, so degree
two evaluation is onto. In degree three use \(XYZ_0,X^3,Y^3,Z_0^3\). The
invariant quadratic \(X^2+Y^2+Z_0^2\) is a nonzero constant on the four
selected vectors; multiplication by it extends surjectivity by two
degrees. The Hilbert series \[
\operatorname{Hilb}(A;s)=\frac{(1-s^2)^2}{(1-s)^3}
=1+3s+4s^2+4s^3+\cdots
\] then proves that every stated surjection is an isomorphism.

Since \(\ell=q\cdot t\) has sign character and is a nonzerodivisor on
\(A[t]\), its homogeneous exact sequence is \[
0\to
\operatorname{sgn}\otimes A_{C-1}\otimes
\operatorname{Sym}^{b-1}W
\xrightarrow{\ \ell\ }
A_C\otimes\operatorname{Sym}^{b}W
\to\mathcal H_{C,b}\to0.
\tag{A.15}
\]

For \(C\ge3\), both \(A_C,A_{C-1}\) are the four-point permutation
module. At the point \(\ell_i\), the map in (A.15) is multiplication by
\(\ell_i\cdot t\). The stabilizer \(S_3\) fixes the line
\(\mathbf Q\ell_i\) in \(\mathfrak h_4\), so its corresponding line in
\(W\) is sign. The quotient \(W/\operatorname{sgn}\) is
\(\mathfrak h_3\otimes\operatorname{sgn}\). Thus evaluation identifies
the cokernel in (A.15) with \[
\operatorname{Ind}_{S_3}^{S_4}
\operatorname{Sym}^b(\mathfrak h_3\otimes\operatorname{sgn}).
\] The explicit rational area-form intertwiner \(J\) of (P.8)--(P.9)
identifies \(\mathfrak h_3\otimes\operatorname{sgn}\) with
\(\mathfrak h_3\); taking its symmetric power gives precisely \(D_b\).

For \(C=2\), the source of (A.15) has \(A_1=\mathfrak h_4\), a
codimension-one submodule of \(\operatorname{Perm}_4\). Extend the
source to the entire permutation module. That larger multiplication map
is injective point by point and has cokernel \(D_b\), just as above. Its
image modulo the smaller image is
\(\operatorname{sgn}\otimes\operatorname{Sym}^{b-1}W\). Hence \[
0\to\operatorname{sgn}\otimes\operatorname{Sym}^{b-1}W
\to\mathcal H_{2,b}\to D_b\to0.
\] Averaging a projection splits this exact sequence, giving the \(C=2\)
row of (A.5).

For \(C=0\) there is no incidence relation, giving
\(\operatorname{Sym}^bW\). For \(C=1\), the quadrics have no component
and \(\mathcal H_{1,b}=R_{1,b}\). The exact incidence/Schur
identification gives \[
R_{1,b}=S_{(b+1,b)}\mathfrak h_4
=\operatorname{sgn}\otimes S_{(b+1,b)}W.
\] The Elliott formula (H.2) with Dynkin labels \((1,b)\) is \[
S_{(b+1,b)}W=\bigoplus_{L=1}^{b+1}H_L,
\] including \(b=0\). This completes the proof of every row of (A.5).

\subsection{3. Exact saturation in the oriented
range}\label{exact-saturation-in-the-oriented-range}

Retwisting the \(q\)-action in (V.4) back to the original action (A.1)
gives \[
(S_\lambda\mathbf Q^4)_{(n^4)}
\simeq\operatorname{sgn}^{d}\otimes
H^0(\mathcal F,\pi^*\mathcal I_Z^k\otimes\mathcal O_{\mathcal F}(c,b)).
\tag{A.16}
\] The extra sign \(\operatorname{sgn}^{c}\) of (V.4) is precisely
removed by that retwist; no determinant or parity factor is omitted.

For \(c\ge2k+1\), this section space is exactly \((J_R^k)_{c,b}\). Here
is a direct proof that does not assume an equality of saturated and
unsaturated ideals. Pull back the resolution (V.14) along the flat map
\(\pi\), and tensor by \(\mathcal O_{\mathcal F}(c,b)\). Its two
line-bundle terms have first degrees \[
r=c-2k\ge1,\qquad r-2\ge-1.
\] On the incidence flag, \[
H^1(\mathcal F,\mathcal O_{\mathcal F}(a,b))=0
\quad\text{for }a\ge-1,\ b\ge0.
\tag{A.17}
\] To check it, use the incidence hypersurface sequence on
\(\mathbf P^2\times\mathbf P^2\). The first cohomology of the ambient
line bundle is zero. The potentially contributing second cohomology of
the preceding term is also zero: its first degree is \(a-1\ge-2\), its
second degree is \(b-1\ge-1\), and neither factor has second cohomology
in those ranges. The one-dimensional sign character of the incidence
equation does not affect vanishing. These projective-space cohomology
statements are exactly the Čech calculation cited and read in T-035.

It follows that the section space in (A.16) is the cokernel \[
\operatorname{Sym}^{k-1}U\otimes\det U\otimes R_{r-2,b}
\longrightarrow
\operatorname{Sym}^{k}U\otimes R_{r,b}.
\tag{A.18}
\] On the other hand the regular sequence \(U\) in \(R\), established by
(A.3), has the power-ideal resolution obtained by extending the
two-variable resolution in (V.14) to this free
\(\operatorname{Sym}(U)\)-module. Its degree-\((c,b)\) cokernel is
exactly \((J_R^k)_{c,b}\), with the same multiplication map as (A.18).
This proves the claimed identification by the actual map, rather than a
dimension comparison.

For \(c=2k\), T-036 gives exactly \[
H^0(\mathcal F,\pi^*\mathcal I_Z^k\otimes\mathcal O_{\mathcal F}(2k,b))
\simeq
(J_R^k)_{2k,b}\oplus
\operatorname{sgn}\otimes\operatorname{Sym}^{k-1}U
\otimes\operatorname{Sym}^{b-1}W,
\] since (A.4) makes the first summand
\(\operatorname{Sym}^kU\otimes\operatorname{Sym}^bW\). The case
\(k=0,c=0\) has zero extra term. For \(k=0\) in all degrees, the ideal
is the whole ring and the same conclusion follows directly. Combining
these equalities with (A.4), (A.5), and (A.16) proves (A.9).

\subsection{4. The three fixed tag
tables}\label{the-three-fixed-tag-tables}

For \(\rho\in\{(4),(1^4),(2,2)\}\), let \(F_\rho\) be \[
F_{(4)}=\{0\},\qquad F_{(1^4)}=\{3\},\qquad F_{(2,2)}=\{1,2\}.
\tag{A.19}
\] Then \[
[\rho:\operatorname{Sym}^jU]
=\sum_{f\in F_\rho}
\#\{(x_2,x_3)\ge0:2x_2+3x_3=j-f\},
\tag{A.20}
\] where integer coordinates are understood.

Let \(E_\tau\) be the harmonic tags of (H.3), and let \(G_\tau\) be the
following induction tags: \[
\begin{array}{c|c|c}
\tau&E_\tau&G_\tau\\ \hline
(4)&\{0,9\}&\{0\}\\
(1^4)&\{3,6\}&\{3\}\\
(2,2)&\{2,4,5,7\}&\{1,2\}\\
(3,1)&\{2,3,4,5,6,7\}&\{0,1,2\}\\
(2,1,1)&\{1,3,4,5,6,8\}&\{1,2,3\}.
\end{array}
\tag{A.21}
\] Their exact meanings are \[
\begin{aligned}
[\tau:\operatorname{Sym}^bW]
&=\sum_{e\in E_\tau}
\#\{(y_2,y_4,y_6)\ge0:2y_2+4y_4+6y_6=b-e\},\\
[\tau:\bigoplus_{L=1}^{b+1}H_L]
&=\sum_{e\in E_\tau}
\#\{(y_4,y_6)\ge0:1\le4y_4+6y_6+e\le b+1\},\\
[\tau:D_b]
&=\sum_{g\in G_\tau}
\#\{(y_2,y_3)\ge0:2y_2+3y_3=b-g\}.
\end{aligned}
\tag{A.22}
\] The first two equations follow from (H.4), (H.8). For the third, the
exact induction branchings are \[
\begin{aligned}
\operatorname{Ind}_{S_3}^{S_4}[(3)]&=[(4)]\oplus[(3,1)],\\
\operatorname{Ind}_{S_3}^{S_4}[(2,1)]&=[(3,1)]\oplus[(2,2)]\oplus[(2,1,1)],\\
\operatorname{Ind}_{S_3}^{S_4}[(1^3)]&=[(2,1,1)]\oplus[(1^4)].
\end{aligned}
\tag{A.23}
\] The first is the four-coset permutation representation, the third is
its sign twist, and the second follows by subtracting these two from
induction of the regular representation of \(S_3\), or by the character
rows (H.11). Combining these branchings with the three \(S_3\) tag
series (P.10) gives the displayed \(G_\tau\).

Write \(K_{\rho\tau}^{\nu}\in\{0,1\}\) for the following fixed tensor
table: \((4)\) is the identity, \((1^4)\) conjugates a partition, and \[
\begin{array}{c|c}
\tau&\{\nu:K_{(2,2),\tau}^{\nu}=1\}\\ \hline
(4),(1^4)&\{(2,2)\}\\
(2,2)&\{(4),(1^4),(2,2)\}\\
(3,1),(2,1,1)&\{(3,1),(2,1,1)\}.
\end{array}
\tag{A.24}
\] This is the exact character multiplication (V.24). Every other entry
is zero.

\subsection{5. The complete positive polytope
evaluator}\label{the-complete-positive-polytope-evaluator}

The following definitions are finite sums over the fixed tables (A.19),
(A.21), (A.24). They contain no representation-theoretic multiplicities
as unknown inputs.

For a label \(\nu\), integers \(L,U_0,b\), define
\(\mathsf D_\nu(L,U_0;b)\) by summing, for every
\(K_{\rho\tau}^{\nu}=1\), \(f\in F_\rho\), \(g\in G_\tau\), the lattice
count in \[
\left\{(x_2,x_3,y_2,y_3)\in\mathbf R_{\ge0}^4:
\quad L\le f+2x_2+3x_3\le U_0,\quad
2y_2+3y_3=b-g\right\}.
\tag{A.25}
\] This count is zero when \(U_0<L\) or \(b<0\).

Define \(\mathsf W_\nu(j,b)\) by summing, for every
\(K_{\rho\tau}^{\nu}=1\), \(f\in F_\rho\), \(e\in E_\tau\), the lattice
count in \[
\left\{(x_2,x_3,y_2,y_4,y_6)\in\mathbf R_{\ge0}^5:
\quad 2x_2+3x_3=j-f,\quad
2y_2+4y_4+6y_6=b-e\right\}.
\tag{A.26}
\] Define \(\mathsf L_\nu(j,b)\) using the same tags but the
four-variable polytope \[
\left\{(x_2,x_3,y_4,y_6)\in\mathbf R_{\ge0}^4:
\quad 2x_2+3x_3=j-f,\quad
1\le4y_4+6y_6+e\le b+1\right\}.
\tag{A.27}
\] The last two functions are zero for \(j<0\) or \(b<0\).

After the exact orientation (A.6), use \(b,c,d,k\) in (A.7) and set
\(\nu=\mu^{[d]}\). The full coefficient is \[
\boxed{
\begin{aligned}
p_{\mu,(n)}^\lambda
={}&\mathsf D_\nu\!\left(k,\left\lfloor\frac{c-3}{2}\right\rfloor;b\right)\\
&+\begin{cases}
\displaystyle
\mathsf L_{\nu'}\!\left(\frac{c-1}{2},b\right),
&c\text{ odd},\\[2mm]
\displaystyle
\mathsf W_\nu\!\left(\frac c2,b\right)
+\mathbf1_{c\ge2k+2}
\left[
\mathsf D_\nu\!\left(\frac{c-2}{2},\frac{c-2}{2};b\right)
+\mathsf W_{\nu'}\!\left(\frac{c-2}{2},b-1\right)
\right]\\
\displaystyle\hspace{13mm}
+\mathbf1_{c=2k}\mathsf W_{\nu'}(k-1,b-1),
&c\text{ even}.
\end{cases}
\end{aligned}}
\tag{A.28}
\] Here the symbols on the left refer to the oriented input; (A.6)
proves equality with the original requested coefficient.

To verify (A.28), partition the positive sum (A.9) according to
\(C=c-2j\). All \(C\ge3\) terms are (A.25), where the value
\(j=f+2x_2+3x_3\) is uniquely recovered from the lattice point. For odd
\(c\), the remaining term has \(C=1\), yielding (A.27) with the sign
twist \(\nu'\). For even \(c\), the \(C=0\) term gives (A.26); the
\(C=2\) term exists exactly when \(c\ge2k+2\), and its two summands give
(A.25) and the sign-shifted (A.26). Finally the boundary correction is
present exactly when \(c=2k\), as displayed. Thus every summand in the
module decomposition has one and only one tagged polytope, and every
lattice point counts one coefficient of its two explicit symmetric-power
series.

All polytopes are bounded: the \(x\)-coordinates have an upper bound
from \(U_0\) or \(j\); the \(y\)-coordinates have an upper bound from
\(b\). Their matrices use only the fixed integers \(1,2,3,4,6\), and
their constants have bit length \(O(\log(n+2))\).

For a fixed label \(\nu\), (A.25) has at most 18 tags; (A.26), (A.27)
have at most 36. When \(c\) is odd, at most \(18+36=54\) polytopes
occur. When \(c\ge2k+2\) is even, the bound is \(18+36+18+36=108\). At
\(c=2k\), only the two (A.26) terms survive, giving at most 72. For
\(\nu=(4)\) or \((1^4)\), the corresponding tag maxima are 6 and 12,
yielding the bound 36; for \(\nu=(2,2)\), they are 12 and 24, yielding
72. Sign and complement twists preserve these three categories.

\subsection{6. Constant-size arithmetic and the standalone
implementation}\label{constant-size-arithmetic-and-the-standalone-implementation}

The fixed polytope counts can also be evaluated by a fixed number of
integer operations. Define \[
C(N)=\#\{(x,y)\ge0:2x+3y=N\},\qquad
P(N)=\#\{(x,y)\ge0:2x+3y\le N\},
\] and set both to zero for \(N<0\). For \(N\ge0\), \[
C(N)=1+\left\lfloor\frac N6\right\rfloor
-\mathbf1_{N\equiv1\pmod6}.
\tag{A.29}
\] For \(e=0,1\), put \(T_e=\lfloor(N-3e)/6\rfloor\). Splitting
\(y=2t+e\) gives \[
P(N)=\sum_{e=0}^1
\left[
(T_e+1)\left(\left\lfloor\frac{N-3e}{2}\right\rfloor+1\right)
-\frac{3T_e(T_e+1)}2
\right].
\tag{A.30}
\] When \(N=0,1,2\), the \(e=1\) value has \(T_e=-1\), so its summand is
zero, as required. Formula (A.29) counts the allowable parity of \(y\).
Formula (A.30) sums the arithmetic progression of allowable \(x\)
values; these are explicit counting proofs.

One tagged polytope in (A.25) therefore has count \[
[P(U_0-f)-P(L-f-1)]\,C(b-g).
\] One in (A.26) has count \[
C(j-f)\,
\begin{cases}
P((b-e)/2),&b-e\ge0\text{ and even},\\
0,&\text{otherwise}.
\end{cases}
\] One in (A.27) has count \[
C(j-f)\left[
P\!\left(\left\lfloor\frac{b+1-e}{2}\right\rfloor\right)
-P\!\left(\left\lfloor\frac{-e}{2}\right\rfloor\right)
\right].
\tag{A.31}
\] The first subtraction counts an interval of lattice points, and the
last excludes angular momentum zero; both quantities are nonnegative by
their defining sets.

The standalone file \url{outer4_positive.py} implements (A.19)--(A.31)
using only Python's standard library and the fixed integer tag tables.
It can also export every tagged polytope as exact equality and
inequality matrices, and enumerate its integer witnesses. Its counting
routine has no loop whose length grows with \(n\), and uses a fixed
number of integer operations. All values have \(O(\log(n+2))\) bits:
each template count is bounded by a constant times \((n+1)^3\).
Consequently ordinary schoolbook integer arithmetic gives an
\(O((\log(n+2))^2)\) bit bound for the arithmetic part, after reading
the input. Output partitions of arbitrary length are first checked by
the stated zero-support rule.

The direct implementation was compared with all 1,980 independent
cycle-index cases and passed. Its witness enumerator separately checked
64,260 exact tagged polytope descriptions, with 2,588 lattice witnesses
across the tested coefficients; every point satisfied the exported
matrices and every template count equaled its explicit enumeration. The
receipt is \url{OUTER4_DIRECT_POLYTOPE_CHECK.tsv}, with summary
\url{OUTER4_DIRECT_POLYTOPE_CHECK.txt} and reproduction code
\url{check_outer4_positive.py}. The exported
\url{EXAMPLE_642_POLYTOPES.json} gives the complete positive count for
\(p_{(4),(3)}^{(6,4,2)}=1\).

\subsection{Verification and scope}\label{verification-and-scope}

The independent calculation in
work/rectangular/verify\_outer4\_all\_formula.py compared (A.9), using
literal class traces and fixed-point induction, against power-sum
plethysm with the Murnaghan--Nakayama rule. Every one of the 1,980 cases
passed: all \(\mu\vdash4\), all \(\lambda\vdash4n\) of length at most
four, and \(0\le n\le6\). This includes 750 strict central-chamber
cases. The receipt is
work/rectangular/OUTER4\_ALL\_FORMULA\_INDEPENDENT\_CHECK.tsv. A
separate direct evaluator of the tagged polytope formula (A.28)
accompanies this report.

The result supplies a direct fixed-template positive formula for every
outer Schur label of size four. The proof uses the specified flag chart,
the four-point complete intersection, actual equivariant multiplication,
and the exact boundary cohomology map. The general terminating
construction G-009--G-011 remains valid in every rank; the efficient
direct-formula frontier after this result starts beyond outer degree
four. No assertion of historical novelty is made for the coefficient
family or for its classical ingredients.

\section{A projective-frame vanishing model in every outer
degree}\label{a-projective-frame-vanishing-model-in-every-outer-degree}

\subsection{G-015. The exact flag-sheaf model and its comparison
map}\label{g-015.-the-exact-flag-sheaf-model-and-its-comparison-map}

Let \(m\ge3\), let \(n\ge0\), and let \(\lambda\vdash mn\) have at most
\(m\) parts. Put \[
E=\mathfrak h_m=\{(q_1,\ldots,q_m):\sum_iq_i=0\},\qquad
\tau=(\lambda_2,\ldots,\lambda_m),\qquad
k=\max(0,\lambda_2-n).
\tag{GF.1}
\] Use the nondegenerate invariant form \(\sum_iq_iq'_i\) to identify
\(E\) and \(E^*\); the restriction is nondegenerate over \(\mathbf Q\),
since its orthogonal complement in \(\mathbf Q^m\) is the constant line.
Let \(\mathcal F=\operatorname{Flag}(E)\), with tautological subbundles
\(\mathcal S_j\), and let \[
\pi:\mathcal F\longrightarrow\mathbf P(E)
\] remember the first line. For \(r=m-1\), set \(\tau_{r+1}=0\) and
define the \(S_m\)-linearized invertible sheaf \[
\mathcal L_\tau=
\bigotimes_{j=1}^{r}
(\det\mathcal S_j^*)^{\otimes(\tau_j-\tau_{j+1})}.
\tag{GF.2}
\] The last factor is the constant determinant line with its actual
group action. In particular its contribution includes
\(\operatorname{sgn}^{\lambda_m}\); it has not been discarded.

Let \(Z_m\) be the following reduced, \(S_m\)-stable set of \(m\)
projective points: \[
Z_m=\{[\ell_1],\ldots,[\ell_m]\},\qquad
\ell_i=e_i-\tfrac1m\mathbf1.
\tag{GF.3}
\] They are distinct for \(m\ge3\). Every \(m-1\) of the vectors are
independent and their sum is zero, so they form a projective frame.

There is an actual \(S_m\)-equivariant vector-space isomorphism \[
\boxed{
(S_\lambda\mathbf Q^m)_{(n^m)}
\simeq
H^0\bigl(\mathcal F,
\mathcal L_\tau\otimes\pi^*\mathcal I_{Z_m}^{\,k}\bigr).}
\tag{GF.4}
\] Thus for every \(\mu\vdash m\), \[
\boxed{
p_{\mu,(n)}^\lambda=
\dim\operatorname{Hom}_{S_m}\!\left(
[\mu],H^0(\mathcal F,
\mathcal L_\tau\otimes\pi^*\mathcal I_{Z_m}^{\,k})\right).}
\tag{GF.5}
\] The map in (GF.4) is exactly the balanced flag chart T-032, followed
by the polynomial-extension criterion proved below. Its inverse on a
section represented by a flag polynomial \(F\) is \[
F\longmapsto u^{n-\lambda_2}F(v(X)),
\qquad u=\prod_iX_{i1},
\tag{GF.6}
\] with the polynomial minors \(v\) of (C.3). Formula (GF.4) therefore
gives a specified scheme-theoretic vanishing condition for the chart
extension in every outer degree.

\subsubsection{1. Sections and their flag-polynomial
realization}\label{sections-and-their-flag-polynomial-realization}

The usual Schur/flag identification is \[
H^0(\mathcal F,\mathcal L_\tau)\simeq S_\tau E^*
\simeq S_\tau E.
\tag{GF.7}
\] The second map is the Schur functor of the invariant form fixed
above. For the first map, a Plücker coordinate is evaluated on each
exterior column of a flag frame, and products of these coordinates have
precisely the degrees \(\tau_j-\tau_{j+1}\) in (GF.2). This fixes the
linearization and the comparison with the polynomial flag algebra,
including the full determinant factor. We use the known Borel--Weil
identification in this exact form: P. Magyar, \emph{A Borel--Weil
Construction for Schur Modules}, arXiv:alg-geom/9411014v1, §2.2,
\textbf{Proposition 4(a)--(c), printed p.~10}, together with
\textbf{Proposition 1, pp.~5--6}, relating the section realization to
the dual Schur module. The source statements and these explicit
comparison maps were read. The standard full-flag pullback in
Proposition 4(b) supplies (GF.2) when some adjacent parts agree.

By T-032, the left side of (GF.4) is identified with the subspace of
(GF.7) satisfying \[
u^k\mid F(v(X))\quad\hbox{in }\mathbf Q[X_{ij}].
\tag{GF.8}
\] It remains to identify this condition with the sheaf in (GF.4). This
is a local calculation along the \(m\) original coordinate divisors, and
we give it for arbitrary \(m\).

\subsubsection{2. The first-line coordinates and all transverse
parameters}\label{the-first-line-coordinates-and-all-transverse-parameters}

Write the first two matrix columns as \(x=(x_i)\), \(s=(s_i)\). The
first smaller-flag vector in the chart is \[
q=u\left((s_i/x_i)_{i=1}^m
-\tfrac1m\sum_j(s_j/x_j)\mathbf1\right).
\tag{GF.9}
\] Every coordinate is polynomial: each denominator divides \(u\).

Treat the divisor \(x_m=0\), on the open set where
\(x_1,\ldots,x_{m-1},s_m\) are nonzero. Put \[
x_m=\epsilon,\quad s_m=S,\quad r_j=s_j/x_j\quad(1\le j<m).
\] At \(\epsilon=0\), (GF.9) is a nonzero scalar multiple of \(\ell_m\).
On the first-projective-factor chart \(q_m-q_1\ne0\), the point
\([\ell_m]\) has the regular coordinate parameters \[
a_j=\frac{q_j-q_1}{q_m-q_1},\qquad2\le j\le m-1.
\] Their exact pullbacks are \[
\boxed{a_j=\epsilon\frac{r_j-r_1}{S-\epsilon r_1}.}
\tag{GF.10}
\] They are \(m-2\) independent parameters: the hyperplane equation
\(\sum_iq_i=0\), together with \(q_m-q_1=1\), recovers every remaining
affine coordinate from them. Because \(\pi\) is a flag bundle, its
inverse image \(\mathcal F_m=\pi^{-1}([\ell_m])\) is smooth and the same
\(a_j\) are regular parameters transverse to that fiber.

\subsubsection{3. Independence from every coordinate of the image
fiber}\label{independence-from-every-coordinate-of-the-image-fiber}

We verify the needed generic independence explicitly, including every
higher flag coordinate. Restrict the original matrices at \(\epsilon=0\)
to the subfamily \[
X_{mj}=0\quad(3\le j\le m).
\tag{GF.11}
\] For a minor \(v_{h+1,I}\) from (C.3), its value is zero when
\(m\notin I\), because its prefactor then contains \(x_m\). When
\(m\in I\), expand the determinant along its last row. Under (GF.11),
that row has its sole nonzero entry \(S\) in column two. The result,
with its ordinary determinant sign, is \[
\left(\prod_{i\notin I}x_i\right)S
\det(X_{i,j})_{i\in I\setminus\{m\},\ j\in\{1,3,\ldots,h+1\}}.
\tag{GF.12}
\] This expression is independent of every \(r_j=s_j/x_j\).

The quotient \(E/\mathbf Q\ell_m\) is identified with
\(\mathbf Q^{m-1}/\mathbf Q\mathbf1\) by restricting to the first
\(m-1\) coordinates and then passing to the quotient. This is an
isomorphism: its kernel in \(E\) is exactly \(\mathbf Q\ell_m\), and the
dimensions agree. In these quotient coordinates, (GF.12) is the exterior
flag generated by the arbitrary ratio columns \[
(X_{ij}/x_i)_{1\le i<m},\qquad3\le j\le m,
\] modulo their constant line. Their first \(m-2\) quotient columns may
be any independent ordered frame. Hence (GF.12) realizes a dense open
set of every flag in \(E/\mathbf Q\ell_m\), which is precisely the fiber
\(\mathcal F_m\). Rescaling the successive columns realizes its nonzero
Plücker frames, including the full determinant frame. All these values
remain independent of \(r_1,\ldots,r_{m-1}\).

Consequently the \(m-2\) leading coefficients \[
(r_j-r_1)/S\quad(2\le j\le m-1)
\] in (GF.10) are algebraically independent over the function field of
the image fiber, and remain independent after choosing any nonzero local
Plücker-frame trivialization. This follows from the exhibited subfamily:
every fiber flag and frame is attained while those \(m-2\) differences
vary freely.

\subsubsection{4. Equality of orders and global ideal-power
membership}\label{equality-of-orders-and-global-ideal-power-membership}

Take a nonzero section \(F\) of \(\mathcal L_\tau\). At the generic
point of \(\mathcal F_m\), trivialize all its Plücker factors by
coordinates nonzero there. Their pulled-back values in (GF.12) are
nonzero, so these trivializations have valuation zero along \(x_m=0\).
Let \(h\) be the section's order in the fiber ideal. Its first nonzero
transverse part belongs to \[
\mathbf Q(\mathcal F_m)[a_2,\ldots,a_{m-1}]_h.
\] Under (GF.10) its coefficient of \(\epsilon^h\) is a nonzero
homogeneous polynomial in the algebraically independent leading
coefficients established above. Therefore it does not vanish. We have
proved \[
\operatorname{ord}_{x_m=0}F(v(X))
=\operatorname{ord}_{\mathcal F_m}F.
\tag{GF.13}
\] Row permutations give the same equality for every \(i\).

For clarity, order at the generic point of this fiber detects its
ordinary sheaf ideal power everywhere. The fiber is smooth and is a
regular embedding in the smooth flag variety. Its associated graded
ideal sheaf is
\(\operatorname{Sym}_{\mathcal O_{\mathcal F_i}}(\mathcal I_i/\mathcal I_i^2)\),
a locally free sheaf on the integral fiber. A section of any such graded
piece that is zero at the generic point is zero everywhere. Starting
with the restriction modulo \(\mathcal I_i\) and proceeding through
degrees \(0,\ldots,k-1\) proves that generic order at least \(k\) is
equivalent to local membership in \(\mathcal I_i^k\) throughout the
fiber. Outside the fiber this ideal is the whole structure sheaf.

The \(x_i\) are pairwise nonassociate primes in the original polynomial
ring. Thus (GF.8) is equivalent to order at least \(k\) along every
original divisor. By (GF.13) and the preceding paragraph it is
equivalent to membership in every \(\mathcal I_i^k\). The fibers are
disjoint, so the intersection of these powers equals their product. The
flag bundle \(\pi\) is flat, and the reduced point ideals are pairwise
comaximal; hence this product is precisely \(\pi^*\mathcal I_{Z_m}^k\).
This proves (GF.4). Taking the explicitly specified \([\mu]\)-covariant
summand gives (GF.5), by T-003.

\subsection{G-016. The full ideal of the projective frame and its
permutation
quotient}\label{g-016.-the-full-ideal-of-the-projective-frame-and-its-permutation-quotient}

For every \(m\ge4\), the saturated homogeneous ideal of
\(Z_m\subset\mathbf P(E)\) is generated by its quadratic part \[
U_m=\ker\bigl(\operatorname{Sym}^2 E\longrightarrow\mathbf Q^{Z_m}\bigr),
\qquad\dim U_m=\frac{m(m-3)}2.
\tag{GF.14}
\] The evaluation map uses the exact affine representatives \(\ell_i\).
In every degree \(C\), its homogeneous coordinate quotient \(A\) has the
actual \(S_m\)-modules \[
\boxed{A_0=\mathbf1,\qquad A_1=E,\qquad
A_C=\operatorname{Perm}_m\quad(C\ge2).}
\tag{GF.15}
\] We prove generation, the scheme statement, and these actions
directly.

Choose \(r=m-1\) projective frame points as the coordinate points of
\(\mathbf P^{r-1}\). The last point has every coordinate nonzero, and
rescaling the coordinates makes it \([1:\cdots:1]\). The quadratic
polynomials vanishing at these \(r+1\) points are precisely the linear
combinations of \(z_iz_j\), \(i<j\), whose coefficient sum is zero. Thus
their ideal is generated by \[
z_i z_j-z_1z_2\qquad(i<j,\ (i,j)\ne(1,2)).
\tag{GF.16}
\] Let \(B\) be the quotient by (GF.16), and write \(z=z_1z_2\). All
cross-products in \(B\) equal \(z\). For distinct \(i,j\), choose an
index \(a\) different from both; this is possible because \(r\ge3\).
Then \[
z_i z=z_i(z_jz_a)=z_j(z_i z_a)=z_j z.
\tag{GF.17}
\] Multiplying this identity repeatedly shows that every degree-\(C\)
mixed monomial, for \(C\ge2\), has the same image \(z z_1^{C-2}\).
Therefore \(B_C\) is spanned by \[
z_1^C,\ldots,z_r^C,\quad z z_1^{C-2}.
\tag{GF.18}
\] Their evaluations at the \(r\) coordinate points and \([1:\cdots:1]\)
are linearly independent: the first \(r\) point rows isolate the first
\(r\) coefficients, and the last point then isolates the mixed
coefficient. Thus the displayed spanning set is a basis and the
evaluation kernel is zero in every degree \(C\ge2\). It is zero in
degrees zero and one as well. The quotient by (GF.16) consequently
embeds in the product of the coordinate rings of the \(m\) lines through
the points, degree by degree. Its kernel is exactly the homogeneous
ideal of their union. This proves radicality, saturation, and quadratic
generation, without an unsupported claim about a projective set
determining its affine scheme.

Return to the canonical affine representatives \(\ell_i\). The
evaluation map in degree \(C\) is equivariant, since literal permutation
matrices take \(\ell_i\) to \(\ell_{\sigma(i)}\). The degree-two
functions \(q_i^2\), \(1\le i\le m\), evaluate to a matrix with diagonal
\((1-1/m)^2\) and off-diagonal \(1/m^2\). Its two eigenvalues are
\((m-2)/m\) and \((m-1)/m\), both nonzero. The analogous degree-three
functions \(q_i^3\) give eigenvalues \(1-3/m+3/m^2\) and
\((m-1)(m-2)/m^2\), also nonzero. Multiplication by the invariant
quadratic \(\sum q_i^2\), whose value is the same nonzero number at
every \(\ell_i\), propagates surjectivity by two degrees. The
coordinate-ring dimensions proved by (GF.18) make each evaluation an
isomorphism in degrees \(C\ge2\). Degree one is the standard submodule
of the permutation module, and degree zero its constant submodule. This
proves every map and action in (GF.15), and dimension subtraction gives
(GF.14).

For \(m=3\), the three projective points in \(\mathbf P(E)=\mathbf P^1\)
have the principal ideal \[
\left(\prod_{1\le i<j\le3}(q_i-q_j)\right).
\tag{GF.19}
\] Each factor vanishes at exactly one of the three points, and their
distinct linear factors generate the reduced degree-three divisor ideal.
The generator has the sign character under literal row permutations.
This also explains the degree-three correction in the rank-three flag
formulas through the same map (GF.4).

\subsection{Geometric dictionary, scope, and the next exact
operation}\label{geometric-dictionary-scope-and-the-next-exact-operation}

All objects in (GF.1)--(GF.19) are explicit schemes, locally free
sheaves, or coherent ideal sheaves. The invariant form and the points
are defined over \(\mathbf Z[1/m]\); their characteristic-zero
global-section representations are those used in the stated equalities.
The flag scheme is the successive locally free quotient construction,
and \(\pi\) is the corresponding relative flag morphism. The comparison
is the coordinate-ring morphism (C.3), localized at \(u\), followed by
the exact regular-embedding calculation (GF.10)--(GF.13). The affine
coordinate dictionary is EGA I 9.5.2--9.5.3, and the group actions are
the constant-group comodules of SGA 3 I 4.7.2--4.7.3. Applying the Schur
and determinant operations gives the exact representation-group classes
of SGA 6 VI 3.3. The complete comparison maps, rather than only equality
of dimensions, carry the proof into these frameworks.

For \(m=4\), (GF.16) consists of two quadrics and the full section
evaluation is T-037. For \(m\ge5\), (GF.4) and (GF.16) leave a precise
new task: evaluate the \(S_m\)-covariant global sections of the
displayed powers of this explicitly generated projective-frame ideal.
The general finite positive construction G-010 already evaluates every
coefficient; the remaining direct-formula problem is to obtain small
uniform positive descriptions for these particular section spaces. No
vanishing of their higher cohomology or positive splitting of their
ideal powers is assumed.

\section{The five projective-frame points: explicit equivariant
resolution}\label{the-five-projective-frame-points-explicit-equivariant-resolution}

\subsection{T-038. Statement and exact
matrices}\label{t-038.-statement-and-exact-matrices}

Work over \(\mathbf Q\). Let \[
E=\mathfrak h_5=\{(q_1,\ldots,q_5):\sum q_i=0\},\qquad
\ell_i=e_i-\tfrac15\mathbf1,
\] and let \(Z\subset\mathbf P(E)\) be the reduced set of the five lines
\([\ell_i]\). The action is literal coordinate permutation. In the basis
\(\ell_1,\ldots,\ell_4\), these lines are the four coordinate points and
\([1:1:1:1]\), since \(\ell_5=-\sum_{i=1}^4\ell_i\). Put
\(S=\mathbf Q[x,y,z,w]=\operatorname{Sym}(E^*)\). The saturated ideal is
\[
I=(f_1,\ldots,f_5),\quad
f=(xz-xy, xw-xy, yz-xy, yw-xy, zw-xy)^t.
\tag{5R.1}
\] The following complex is an exact minimal graded free resolution: \[
0\longrightarrow S(-5)\mathrel{\mathop{\longrightarrow}^{N}}
 S(-3)^5\mathrel{\mathop{\longrightarrow}^{M}}
 S(-2)^5\mathrel{\mathop{\longrightarrow}^{f^t}}I
\longrightarrow0,
\tag{5R.2}
\] where \[
M=\begin{pmatrix}
y-w&y&y&y-w&0\\
z-y&-y&-y&-y&-y\\
0&-x&-w&-x&-w\\
0&x&z&0&0\\
0&0&0&x&y
\end{pmatrix},\qquad
N=\begin{pmatrix}
y(x-w)\\z(y-w)\\x(w-y)\\y(w-z)\\x(z-w)
\end{pmatrix}.
\tag{5R.3}
\] Let \(U=[3,2]\) and \(\epsilon=\operatorname{sgn}\). There are
uniquely specified actions on the spaces of minimal generators and
minimal syzygies, inherited from the action on \(I\). With these actions
the resolution is \[
\boxed{
0\longrightarrow\epsilon\otimes S(-5)
\longrightarrow(\epsilon\otimes U)\otimes S(-3)
\longrightarrow U\otimes S(-2)
\longrightarrow I\longrightarrow0.}
\tag{5R.4}
\] In particular, sheafification and tensoring by any specified locally
free sheaf give an exact equivariant complex. These are concrete maps:
(5R.3) gives their matrices, and the finite group acts on each syzygy
space by applying its action to its polynomial vectors.

\subsubsection{Proof of the scheme and of
exactness}\label{proof-of-the-scheme-and-of-exactness}

G-016 proves the projective-frame ideal statement for every outer
degree. Here it also has the following short direct verification. In
\(S/I\) every cross-product is equal to \(p=xy\). For distinct indices
\(i,j\), choose a third index \(a\); then
\(z_i p=z_i(z_jz_a)=z_j(z_i z_a)=z_jp\). Hence in degree \(C\ge2\) every
mixed monomial equals \(p x^{C-2}\). The five elements \[
x^C,y^C,z^C,w^C,px^{C-2}
\tag{5R.5}
\] are independent by evaluation at the four coordinate points and
\([1:1:1:1]\). The same evaluation is injective in degrees zero and one.
Thus \(S/I\) embeds in the product of the coordinate rings of the five
lines, and its ideal is precisely their reduced union ideal. In
particular \(I\) is saturated and \[
H_{S/I}(t)=1+4t+\frac{5t^2}{1-t}
=\frac{1+3t+t^2}{1-t}.
\tag{5R.6}
\]

Direct multiplication gives \(f^tM=0\) and \(MN=0\). The upper-left
\(4\times4\) minor of \(M\) is \[
x(z-w)(xy-wz),
\tag{5R.7}
\] which is nonzero. Thus the fraction-field kernel of \(M\) is the line
spanned by \(N\). The entries of \(N\) have greatest common divisor one:
already the irreducible factors of \(y(x-w)\) have no common factor with
\(z(y-w)\). If a fraction \(a/b\), in lowest terms, has
\((a/b)N\in S^5\), every irreducible divisor of \(b\) divides all the
entries of \(N\). Therefore \(b\) is a unit. This proves \(\ker M=S N\),
and the left map is injective.

Subtracting (5R.6) from \(H_S(t)=(1-t)^{-4}\) gives \[
H_I(t)=\frac{5t^2-5t^3+t^5}{(1-t)^4}.
\tag{5R.8}
\] It follows that \(\ker f^t\) has Hilbert series
\((5t^3-t^5)/(1-t)^4\). The established kernel of \(M\) gives exactly
that series for \(\operatorname{im}M\). Since the latter is contained in
\(\ker f^t\), their degree-by-degree dimensions prove equality. All
matrix entries have positive degree, so the resolution is minimal.

\subsubsection{The actual finite-group
labels}\label{the-actual-finite-group-labels}

The canonical affine representatives \(\ell_i\) are permuted without
rescaling. Consequently evaluation gives \[
(S/I)_0=\mathbf1,\qquad(S/I)_1=E^*,\qquad
(S/I)_C=\operatorname{Perm}_5\quad(C\ge2)
\tag{5R.9}
\] as actual \(S_5\)-modules. An invariant nondegenerate form identifies
\(E^*\) with \(E\). Write the conjugacy classes in the order \[
1^5,\quad21^3,\quad2^21,\quad31^2,\quad32,\quad41,\quad5.
\tag{5R.10}
\] The character of \(I_2=\operatorname{Sym}^2E-\operatorname{Perm}_5\)
is \[
(5,1,1,-1,1,-1,0),
\tag{5R.11}
\] which is the irreducible character \([3,2]\). For example the
character can be computed directly from
\(\chi_{\operatorname{Sym}^2E}(g)=
(\chi_E(g)^2+\chi_E(g^2))/2\); the irreducible labels agree with the
primary character table in Dolgachev--Farb--Looijenga, \emph{Geometry of
the Wiman--Edge pencil I}, §2.3, Table 1, printed p.~10.

The degree-three syzygy character is \[
\chi_E\chi_U-\chi_{\operatorname{Sym}^3E}
 +\chi_{\operatorname{Perm}_5}
=(5,-1,1,-1,-1,1,0)=\chi_{[2,2,1]}.
\tag{5R.12}
\] This follows either by the degree-three exact sequence, or by
expanding \(\prod_{\xi\in\operatorname{eig}(g|E)}(1-t\xi)^{-1}\) to
degree three. It fixes the action on the columns of \(M\).

There is one final generator. At a transposition \(s\), the equivariant
Hilbert series of \(S/I\) is \((1+t+t^2)/(1-t)\), whereas that of \(S\)
is \(((1-t)^3(1+t))^{-1}\). Their ratio is \[
1-t^2-t^3+t^5.
\tag{5R.13}
\] In the alternating numerator of (5R.2) the final coefficient is
\(-\chi_K(s)t^5\). Thus \(\chi_K(s)=-1\). Every one-dimensional
representation of \(S_5\) is determined by its common value on the
transpositions, which generate the group, so \(K=\epsilon\). This proves
(5R.4), including its final determinant character.

For an alternating form of the middle map, set \[
P=\begin{pmatrix}
0&0&0&-1&0\\0&0&1&0&-1\\0&1&0&0&0\\
0&0&-1&1&0\\1&-1&0&0&0
\end{pmatrix}.
\] Then \(\det P=-1\), \(N=Pf\), and \[
MP=\begin{pmatrix}
0&y&w&0&-y\\-y&0&0&-z&y\\-w&0&0&-x&x\\
0&z&x&0&-x\\y&-y&-x&x&0
\end{pmatrix}.
\tag{5R.14}
\] Its signed maximal Pfaffians are \(f\). Thus the five-point ideal
also comes with a specified alternating presentation, as used in the
separate application of Kustin--Polini--Ulrich in the archive report.
The elementary exactness proof above does not need an unverified generic
perfection claim.

\subsection{The quadratic subalgebra and its six-coordinate
morphism}\label{the-quadratic-subalgebra-and-its-six-coordinate-morphism}

This subsection records an exact comparison useful for the higher
powers. Let \(A,B,C,D,E_0\) be independent indeterminates mapping to the
five entries of \(f\), in that order. Then \[
\boxed{\mathbf Q[f_1,\ldots,f_5]
\simeq\mathbf Q[A,B,C,D,E_0]/(F),\qquad
F=AD(E_0-B-C)-BC(E_0-A-D).}
\tag{5R.15}
\] The variable \(E_0\) is distinguished from the standard module \(E\).
Substitution verifies \(F(f)=0\). In the fraction field of the image, \[
xy=\frac{AD}{E_0-A-D},\qquad
xz=xy+A,\quad xw=xy+B,\quad yz=xy+C,
\tag{5R.16}
\] and similarly for the other cross-products. In particular
\(x^2=(xy)(xz)/(yz)\). The full rational function field
\(\mathbf Q(x,y,z,w)\) is algebraic over the image's fraction field:
adjoin a square root \(x\) of this displayed \(x^2\), then recover
\(y,z,w\) by division. The image has transcendence degree four.

As a polynomial in \(E_0\), \(F\) has coefficient \(AD-BC\), an
irreducible determinant. This coefficient does not divide the constant
term: at \((A,B,C,D)=(1,2,3,6)\) the coefficient is zero and the
constant term is nonzero. The coefficient and constant term are
therefore coprime. Gauss's lemma and the degree-one polynomial over
\(\mathbf Q(A,B,C,D)\) prove that \(F\) is irreducible. The kernel has
height one and contains \(F\); since the target is a domain this proves
(5R.15). The transposition \(x\leftrightarrow y\) sends
\((A,B,C,D,E_0)\) to \((C,D,A,B,E_0)\) and sends \(F\) to \(-F\). Its
relation line is therefore \(\epsilon\).

Here is a rational linear comparison with the usual six-coordinate Segre
cubic. Set \[
\begin{aligned}
t_1&= A+B-C+D-E_0,&t_2&=-A+B+C-D-E_0,\\
t_3&=-A+B+C-D+E_0,&t_4&=-A-B-C+D+E_0,\\
t_5&= A-B-C-D+E_0,&t_6&= A-B+C+D-E_0.
\end{aligned}
\tag{5R.17}
\] These forms have rank five and satisfy \[
\sum_{i=1}^6t_i=0,\qquad\sum_{i=1}^6t_i^3=-24F.
\tag{5R.18}
\] For precision, the associated action on six letters is the
conjugation action of \(S_5\) on its six subgroups of order five. Index
those subgroups by lexicographically ordering their permutation tuples,
as in the certificate script. The six forms in (5R.17) obey \[
g(t_i)=\epsilon(g)t_{g(i)}.
\tag{5R.19}
\] There is also a group-theoretic way to specify these exact forms
without guessing an outer automorphism. If \(H_0\) is the first such
subgroup, average \(\epsilon(g)\rho_U(g)\) over its twenty-element
normalizer. Its image is a line. Take the first nonzero column, clear
denominators and divide by the positive gcd of its integer entries. For
a subgroup \(gH_0g^{-1}\), transport this vector by
\(\epsilon(g)\rho_U(g)\). This is independent of the chosen transporter
because of the normalizer character. It gives precisely (5R.17).
Multiplication by every one of the 120 group elements was checked
against (5R.19).

The conjugacy-class map, in the order (5R.10), is \[
1^6,\quad2^3,\quad2^21^2,\quad3^2,\quad6,\quad41^2,\quad51.
\tag{5R.20}
\] Equations (5R.17)--(5R.20) specify both the cubic isomorphism and the
finite-group intertwiner. This agrees with the classical five-point
quadratic map: Dolgachev--Farb--Looijenga, Remark 2.4, printed p.~13,
describes the five-point map and its Segre cubic image; Dolgachev,
\emph{Corrado Segre and nodal cubic threefolds}, §2, equation (2.1),
gives the six-coordinate equations. We do not attach a novelty claim to
this classical comparison.

\subsection{Sources and checks}\label{sources-and-checks}

\begin{itemize}
\tightlist
\item
  \href{https://www.math.uchicago.edu/~farb/papers/wiman.pdf}{Dolgachev--Farb--Looijenga,
  \emph{Geometry of the Wiman--Edge pencil I}}, §2.3, Table 1, printed
  p.~10, and Remark 2.4, printed p.~13. Both passages were opened and
  read in this run.
\item
  \href{https://sites.lsa.umich.edu/idolga-new/wp-content/uploads/sites/1467/2024/08/segre15.pdf}{Dolgachev,
  \emph{Corrado Segre and nodal cubic threefolds}}, §2, especially
  equation (2.1) and the five-point quadratic map.
\item
  \texttt{\small check\_\allowbreak{}five\_\allowbreak{}points.py}
  verifies the two zero matrix products, the primitive kernel, the
  nonzero minor, cubic substitution, six-form rank, all 120 group
  intertwining identities and the class map. The full receipt is
  \texttt{\small FIVE\_\allowbreak{}POINT\_\allowbreak{}EXACT\_\allowbreak{}MAPS.txt}.
\item
  \texttt{\small check\_\allowbreak{}five\_\allowbreak{}point\_\allowbreak{}positive.py}
  additionally checks the entire equivariant Hilbert numerator on all
  seven classes. The matrices and the degree-by-degree exactness proof
  specify the complex independently of those finite checks.
\end{itemize}

The next result, T-039, evaluates a family of quotient-bundle sections
from this actual resolution. General arbitrary powers with arbitrary
Schur quotient bundles remain a separate direct-evaluation task; G-010
already supplies a terminating finite positive count for all plethysm
coefficients.

The rectangular agent independently read the complete proof and checked
the actual \(M,N,P\) products, signed Pfaffians, scalar \(-24\), and all
120 signed actions using the five affine frame vectors. Every check
passed. Its independent code and receipt are
\texttt{\small work/\allowbreak{}rectangular/\allowbreak{}verify\_\allowbreak{}five\_\allowbreak{}point\_\allowbreak{}matrices.py}
and
\texttt{\small work/\allowbreak{}rectangular/\allowbreak{}FIVE\_\allowbreak{}POINT\_\allowbreak{}MATRIX\_\allowbreak{}INDEPENDENT\_\allowbreak{}CHECK.txt}.

\section{A positive higher-flag family in outer degree
five}\label{a-positive-higher-flag-family-in-outer-degree-five}

\subsection{T-039. The coefficient
formula}\label{t-039.-the-coefficient-formula}

For integers \(a,d\ge0\), set \[
n=d+a+1,\qquad
\lambda=(d+2a+3,\ d+a+2,\ d+a,\ d+a,\ d).
\tag{5F.1}
\] This is a partition of \(5n\). Let \(E=\mathfrak h_5\), \(U=[3,2]\),
and \(\epsilon=\operatorname{sgn}_{S_5}\). There is an actual
equivariant isomorphism \[
\boxed{
(S_\lambda\mathbf Q^5)_{(n^5)}
\simeq\epsilon^{d+a}\otimes
\left(U\otimes\operatorname{Sym}^a E
\;\oplus\;\operatorname{Sym}^{a-1}E\right),}
\tag{5F.2}
\] where \(\operatorname{Sym}^{-1}E=0\). The direct sum is a splitting
of a specified equivariant short exact sequence; it need not be a
canonical splitting. Consequently every outer Schur coefficient
\(p_{\mu,(n)}^\lambda\), \(\mu\vdash5\), is a sum of at most 36 fixed,
four-variable lattice-point counts. The numbers and inequalities are
specified in §4 below. There is no coefficient precomputation depending
on \(a\), \(d\), or \(n\).

The passage from the balanced weight space to (5F.2) is the polynomial
chart and extension morphism of T-032/G-015, followed by explicit
exterior pairings and the exact resolution (5R.3). We prove each
remaining step.

\subsection{1. From the flag model to a quotient-bundle
section}\label{from-the-flag-model-to-a-quotient-bundle-section}

Use the canonical \(S_5\)-invariant form on \(E\) to identify \(E\) and
\(E^*\). On \(\mathbf P(E)\), which parametrizes lines, let \[
0\longrightarrow\mathcal O(-1)\longrightarrow E\otimes\mathcal O
\longrightarrow Q\longrightarrow0
\tag{5F.3}
\] be the tautological quotient sequence. Thus \(Q\) has rank three. Let
\(Z=\{[e_i-\frac15\mathbf1]:1\le i\le5\}\), and write
\(\mathcal I=\mathcal I_Z\).

The numbers in (5F.1) satisfy \[
\lambda_2-n=1,\qquad \lambda_2-\lambda_3=2,
\qquad\lambda_3-\lambda_4=0,\qquad\lambda_3-\lambda_5=a.
\tag{5F.4}
\] G-015 therefore identifies the balanced module with \[
H^0\!\left(\operatorname{Flag}(E),
\mathcal L_{(d+a+2,d+a,d+a,d)}\otimes\pi^*\mathcal I\right),
\tag{5F.5}
\] where \(\pi\) remembers the first line. This is the actual comparison
map: a flag polynomial \(F\) is sent back to the original matrix ring by
\(u^{-1}F(v(X))\). G-015 proves that membership in (5F.5) is precisely
what makes this rational expression a polynomial, using the exact
transverse ratios along every first-column coordinate divisor.

We describe the pushforward in (5F.5), including all determinant
factors. The full determinant factor of \(\mathcal L\) is
\((\det E^*)^d\). After removing it, the first weight is \(a+2\) and the
relative quotient-flag weight is \((a,a,0)\). The standard Plücker
evaluation on the flag of the rank-three quotient gives \[
\pi_*\mathcal L_{(d+a+2,d+a,d+a,d)}
\simeq (\det E^*)^d\otimes\mathcal O(a+2)
                  \otimes S_{(a,a)}Q^*.
\tag{5F.6}
\] The relative identification can be checked on every open set
trivializing \(Q\): there it is the usual Schur realization by the
degree-\((a,a)\) flag minors, and changes of quotient frame act by that
same Schur functor. The local maps therefore glue. This is the
Borel--Weil/Schur comparison in Magyar, Proposition 1, printed pp.~5--6,
and Proposition 4(a)--(c), printed p.~10, with the bundle linearization
specified by its exterior columns.

For a rank-three vector space or bundle, the wedge pairing is the actual
map \[
\Lambda^2Q^*\simeq Q\otimes\det Q^*.
\tag{5F.7}
\] Its \(a\)-th symmetric power identifies
\(S_{(a,a)}Q^*=\operatorname{Sym}^a(\Lambda^2Q^*)\) with
\((\det Q^*)^a\otimes\operatorname{Sym}^aQ\). The asserted equality with
the Schur functor follows as well from the map on the highest vector
\((q_1^*\wedge q_2^*)^a\): its irreducible span has weight \((a,a,0)\)
and the dimensions on both sides are \(\binom{a+2}{2}\). Taking
determinants in (5F.3) gives \[
\det Q^*\simeq\det E^*\otimes\mathcal O(-1).
\tag{5F.8}
\] Thus (5F.6) becomes \[
\pi_*\mathcal L\simeq
(\det E^*)^{d+a}\otimes\mathcal O(2)\otimes\operatorname{Sym}^aQ.
\tag{5F.9}
\]

Projection with the ideal factor causes no hidden flatness assumption.
Locally on the base, the flag bundle is a product with the fixed flag
variety. A finite affine Čech complex for its line bundle consists of
modules obtained from vector spaces over \(\mathbf Q\), tensored with
the base ring. Tensoring this complex with the local ideal commutes with
its degree-zero cohomology: every exact sequence of vector spaces over
\(\mathbf Q\) splits. These local maps agree on overlaps. Equivalently
one can use the usual relative Schur section map and its base-change
property. Hence (5F.5)--(5F.9) give the specified isomorphism \[
\boxed{
(S_\lambda\mathbf Q^5)_{(n^5)}
\simeq \epsilon^{d+a}\otimes
H^0(\mathbf P(E),\mathcal I(2)\otimes\operatorname{Sym}^aQ).}
\tag{5F.10}
\] Here \(\det E=\epsilon\), because the full five-dimensional
permutation module is \(\mathbf1\oplus E\). This determines the twist
directly.

\subsection{2. The three bundle cohomology
calculations}\label{the-three-bundle-cohomology-calculations}

Write \(B_a=\operatorname{Sym}^aQ\). Multiplication by the tautological
line in (5F.3) is the first map of the exact sequence \[
0\longrightarrow\mathcal O(t-1)\otimes\operatorname{Sym}^{a-1}E
\longrightarrow\mathcal O(t)\otimes\operatorname{Sym}^{a}E
\longrightarrow B_a(t)\longrightarrow0.
\tag{5F.11}
\] The sequence is exact fiber by fiber: the kernel of
\(\operatorname{Sym}^aE\to\operatorname{Sym}^a(E/L)\) is exactly
\(L\operatorname{Sym}^{a-1}E\), and multiplication by a nonzero vector
of the line \(L\) is injective in a symmetric algebra. It is therefore
an exact sequence of locally free sheaves. For \(a=0\), its left term is
zero.

The projective-space cohomology used here is \[
H^i(\mathbf P^3,\mathcal O(t))=0\quad(0<i<3),
\] with \(\mathcal O(-1),\mathcal O(-2),\mathcal O(-3)\) acyclic and
\(H^3(\mathcal O(-4))=\det E\). The scalar vanishing ranges follow from
the full monomial Čech calculation in Stacks Project, Lemma 30.8.1, tag
01XS. The determinant action can be read from its top Čech generator:
for an ordered coordinate basis \(x_1,\ldots,x_4\), it is the class of
\(1/(x_1x_2x_3x_4)\); the ordered Čech orientation and change of basis
make its action \(\det E\). The same conclusion follows from the
canonical-bundle identification
\(\omega_{\mathbf P(E)}=\mathcal O(-4)\otimes\det E^*\).

Apply (5F.11) with \(t=0,-1,-3\). It gives all the following groups, not
merely their Euler characteristics: \[
\begin{array}{c|c|c}
\text{bundle}&\text{only possible nonzero cohomology}&\text{module}\\\hline
B_a&H^0&\operatorname{Sym}^aE\\
B_a(-1)&\text{none}&0\\
B_a(-3)&H^2&\epsilon\otimes\operatorname{Sym}^{a-1}E.
\end{array}
\tag{5F.12}
\] The last row is zero when \(a=0\), as required.

\subsection{3. The positive exact
sequence}\label{the-positive-exact-sequence}

Sheafify the exact resolution (5R.4), twist by \(\mathcal O(2)\), and
tensor with the locally free sheaf \(B_a\). This gives \[
0\longrightarrow\epsilon\otimes B_a(-3)
\longrightarrow(\epsilon\otimes U)\otimes B_a(-1)
\longrightarrow U\otimes B_a
\longrightarrow\mathcal I(2)\otimes B_a\longrightarrow0.
\tag{5F.13}
\] Let \(K\) be the image of its middle map into \(U\otimes B_a\). The
first short exact sequence and (5F.12) give \[
H^0(K)=0,\qquad
H^1(K)=\operatorname{Sym}^{a-1}E,
\qquad H^i(K)=0\ (i\ge2).
\tag{5F.14}
\] The two sign factors in \(H^1(K)\) are the last syzygy character of
(5R.4) and the top projective cohomology character of (5F.12); their
tensor product is trivial. The second short exact sequence therefore
gives \[
\boxed{
0\longrightarrow U\otimes\operatorname{Sym}^a E
\longrightarrow H^0(\mathcal I(2)\otimes B_a)
\longrightarrow\operatorname{Sym}^{a-1}E
\longrightarrow0.}
\tag{5F.15}
\] The first arrow multiplies a quadratic vanishing at \(Z\) by a
constant section of \(\operatorname{Sym}^aE\), followed by the quotient
\(\operatorname{Sym}^aE\to B_a\). The last arrow is the connecting
homomorphism of the actual resolution. This specifies the morphisms
underlying the positive decomposition.

Every sequence of finite-dimensional \(\mathbf Q[S_5]\)-modules splits:
average any linear section over the 120 group elements. Applying this to
(5F.15) and composing with (5F.10) proves (5F.2). It also proves \[
\dim(S_\lambda\mathbf Q^5)_{(n^5)}
=5\binom{a+3}{3}+\binom{a+2}{3}.
\tag{5F.16}
\] The equality for \(a=0\) is the single module \(U\), with no extra
cohomology term.

\subsection{4. Fixed positive polytopes for every outer
partition}\label{fixed-positive-polytopes-for-every-outer-partition}

Define the bounded rational polytope \[
P(R)=\{(x_2,x_3,x_4,x_5)\in\mathbf R_{\ge0}^4:
2x_2+3x_3+4x_4+5x_5=R\},
\qquad C(R)=\#(P(R)\cap\mathbf Z^4).
\tag{5F.17}
\] For \(R<0\) it is empty. Every integer coordinate is at most
\(\max(0,R)/2\), so no unbounded component is being counted.

The seven multisets below are the comajor indices of standard tableaux
of the indicated shape. For a tableau of size five, our convention is
\(\operatorname{comaj}T=\sum_{i:\,i+1\text{ is below }i}(5-i)\).
Repeated entries are distinct tableau tags.

\begin{longtable}[]{@{}llr@{}}
\toprule\noalign{}
\(\rho\) & multiset \(D_\rho\) & dimension \\
\midrule\noalign{}
\endhead
\bottomrule\noalign{}
\endlastfoot
\((5)\) & \(0\) & 1 \\
\((4,1)\) & \(1,2,3,4\) & 4 \\
\((3,2)\) & \(2,3,4,5,6\) & 5 \\
\((3,1,1)\) & \(3,4,5,5,6,7\) & 6 \\
\((2,2,1)\) & \(4,5,6,7,8\) & 5 \\
\((2,1,1,1)\) & \(6,7,8,9\) & 4 \\
\((1^5)\) & \(10\) & 1 \\
\end{longtable}

The exact standardization argument of T-027 gives \[
\dim\operatorname{Hom}_{S_5}([\rho],\operatorname{Sym}^aE)
=\sum_{r\in D_\rho}C(a-r).
\tag{5F.18}
\] In that proof, weak entries in a semistandard tableau become a
standard tableau plus nonnegative increments, and the forced strict
inequalities contribute precisely its comajor. The constant-line factor
in the permutation module removes the denominator of degree one, leaving
the four denominators \(2,3,4,5\) in (5F.17). Thus (5F.18) has an actual
tableau-to-lattice bijection. Here all tableau tags are already listed.

Let \(K_{\rho\nu}\) be the multiplicity of \([\nu]\) in
\(U\otimes[\rho]\). It is completely specified by the following table;
repetition denotes multiplicity two.

\begin{longtable}[]{@{}ll@{}}
\toprule\noalign{}
\(\rho\) & summands of \(U\otimes[\rho]\) \\
\midrule\noalign{}
\endhead
\bottomrule\noalign{}
\endlastfoot
\((5)\) & \((3,2)\) \\
\((4,1)\) & \((4,1),(3,2),(3,1,1),(2,2,1)\) \\
\((3,2)\) & \((5),(4,1),(3,2),(3,1,1),(2,2,1),(2,1,1,1)\) \\
\((3,1,1)\) & \((4,1),(3,2),(3,1,1),(3,1,1),(2,2,1),(2,1,1,1)\) \\
\((2,2,1)\) & \((1^5),(2,1,1,1),(2,2,1),(3,1,1),(3,2),(4,1)\) \\
\((2,1,1,1)\) & \((2,1,1,1),(2,2,1),(3,1,1),(3,2)\) \\
\((1^5)\) & \((2,2,1)\) \\
\end{longtable}

This finite tensor table follows by multiplying the seven character rows
of the primary Table 1 quoted in T-038 and taking the
class-size-weighted inner products. In particular it requires no
variable-size Kronecker coefficient computation. All 49 entries were
also recomputed by an independent Murnaghan--Nakayama implementation.

For \(\mu\vdash5\), put \[
\nu=\begin{cases}\mu&d+a\text{ even},\\\mu^t&d+a\text{ odd}.
\end{cases}
\tag{5F.19}
\] Taking the \([\mu]\)-covariant part in the balanced-space dictionary
T-003 and then using (5F.2) and (5F.18) yields \[
\boxed{
p_{\mu,(n)}^\lambda
=\sum_{\rho\vdash5}K_{\rho\nu}
       \sum_{r\in D_\rho}C(a-r)
 +\sum_{r\in D_\nu}C(a-1-r).}
\tag{5F.20}
\] Each summand has a fixed tag \((\rho,\text{copy},\text{tableau})\) or
\((\nu,\text{tableau})\), followed by one point of (5F.17). This is a
finite disjoint union of bounded lattice polytopes. Negative right sides
give empty polytopes and may be retained in the uniform list.

All these modules are self-dual. Hence
\(\sum_\rho K_{\rho\nu}\dim[\rho]=5\dim[\nu]\), by the dimension of
\(U\otimes[\nu]\). The second sum has \(\dim[\nu]\) tags. The total
number is exactly \(6\dim[\nu]\le36\), including empty tags. Each
polytope has four variables, one displayed equality, and four
nonnegativity inequalities. Its coefficients are fixed and its right
side has \(O(\log(a+2))\) bits. This gives a fixed-dimensional positive
formula in the ordinary binary input encoding for \(a,d\).

\subsection{5. The symmetric outer power and its complete support
criterion}\label{the-symmetric-outer-power-and-its-complete-support-criterion}

For \(\mu=(5)\), (5F.20) is particularly short: \[
\boxed{
p_{(5),(n)}^\lambda=
\begin{cases}
\displaystyle\sum_{j=1}^{6}C(a-j),&d+a\text{ even},\\[2mm]
\displaystyle\sum_{j=4}^{8}C(a-j)+C(a-11),&d+a\text{ odd}.
\end{cases}}
\tag{5F.21}
\] Both lists have six tags. The numerical semigroup generated by
\(2,3,4,5\) is \(\{0\}\cup\{r:r\ge2\}\): every even \(r\ge2\) uses twos
and every odd \(r\ge3\) uses one three and twos. Therefore \[
\boxed{
p_{(5),(n)}^\lambda>0
\quad\Longleftrightarrow\quad
\begin{cases}
a\ge1,&d+a\text{ even},\\
a\ge4,&d+a\text{ odd}.
\end{cases}}
\tag{5F.22}
\] To check the endpoints explicitly, in the first case \(a=1\) uses the
tag \(j=1\), and \(a=2\) uses \(j=2\); larger \(a\) use either a zero
residual or a residual at least two. In the second case the tags
\(j=4,5\) handle \(a=4,5\), and \(j=4\) handles every \(a\ge6\). Below
these thresholds all listed right sides are negative.

\subsection{6. Exact implementation, independent checks, and
scope}\label{exact-implementation-independent-checks-and-scope}

\texttt{\small five\_\allowbreak{}point\_\allowbreak{}positive.py} is a
standalone Python standard-library evaluator. It contains exactly the
two fixed tables above, exports the tagged equations, and enumerates
their integer witnesses. It accepts \(a,d,\mu\), checks that \(a,d\) are
nonnegative integers and that \(\mu\) is a partition of five, and
returns (5F.1) and (5F.20). For example:

\begin{Shaded}
\begin{Highlighting}[]
\NormalTok{python work/theta/outer5/five\_point\_positive.py {-}{-}a 4 {-}{-}d 0 {-}{-}outer 5 {-}{-}polytopes}
\end{Highlighting}
\end{Shaded}

The implementation uses the positive coin recurrence for evaluation; its
polytope description has fixed size, regardless of that elementary
counting implementation. No character, elimination, or plethysm oracle
is called by the evaluator.

\texttt{\small check\_\allowbreak{}five\_\allowbreak{}point\_\allowbreak{}positive.py}
recomputed all seven tableau degree lists, all 49 tensor entries, the
seven exact rational character identities behind (5F.18), and every
equivariant resolution numerator of T-038. It also compared the positive
count with literal tagged lattice enumeration in 434 cases, producing
118,334 explicit witnesses. Every check passed. Receipts:
\texttt{\small FIVE\_\allowbreak{}POINT\_\allowbreak{}POSITIVE\_\allowbreak{}CHECK.tsv}
and \texttt{\small .txt}.

Independently, the rectangular agent computed all seven outer Schur
cycle indices and used Murnaghan--Nakayama on the output partitions. All
147 cases with \(a,d\ge0\), \(d+a+1\le6\), matched (5F.2). On the
proposed side, that calculation used permutation-cycle traces for
\(\operatorname{Sym}^a\mathfrak h_5\) and the character
\(\binom{m_1}{2}+m_2-m_1\) for \([3,2]\), independently of the
resolution and the fixed tag table. Code and receipt:
\texttt{\small work/\allowbreak{}rectangular/\allowbreak{}verify\_\allowbreak{}outer5\_\allowbreak{}fivepoint\_\allowbreak{}family.py}
and
\texttt{\small work/\allowbreak{}rectangular/\allowbreak{}OUTER5\_\allowbreak{}FIVEPOINT\_\allowbreak{}INDEPENDENT\_\allowbreak{}CHECK.tsv}.

The same independent reader subsequently read the complete T-038 and
T-039 proofs, including every bundle map, the two sign cancellations,
the finite tag lists and all support endpoints, and found no gap.

Primary sources read for this proof:

\begin{itemize}
\tightlist
\item
  \href{https://arxiv.org/pdf/alg-geom/9411014}{Magyar, \emph{A
  Borel--Weil Construction for Schur Modules}}, Proposition 1, printed
  pp.~5--6, and Proposition 4(a)--(c), printed p.~10. The
  exterior-column comparison fixes the precise dual convention in
  (5F.6); the maps (5F.7)--(5F.9) then determine all twists.
\item
  \href{https://stacks.math.columbia.edu/tag/01XS}{Stacks Project, Lemma
  30.8.1, tag 01XS}, including its monomial Čech proof. Only the three
  exact bundle cohomology computations in (5F.12) are imported from
  these scalar vanishings.
\item
  \href{https://www.math.uchicago.edu/~farb/papers/wiman.pdf}{Dolgachev--Farb--Looijenga,
  \emph{Geometry of the Wiman--Edge pencil I}}, §2.3, Table 1, printed
  p.~10, for the fixed \(S_5\) labels and table.
\end{itemize}

For \(a>0\), the section space in (5F.10) has the nontrivial quotient
bundle \(\operatorname{Sym}^aQ\), and (5F.15) evaluates it in full. The
result is an explicit bounded family within the general projective-frame
model. It makes no global novelty claim and assumes no vanishing for
arbitrary point powers or arbitrary Schur quotient bundles. G-010
provides the general terminating positive count; T-037 already gives a
fixed finite formula for every outer partition of four.

\section{A positive five-point quotient-bundle family in every
multiplicity}\label{a-positive-five-point-quotient-bundle-family-in-every-multiplicity}

\subsection{G-017. A fixed positive formula on a varying five-row
family}\label{g-017.-a-fixed-positive-formula-on-a-varying-five-row-family}

For integers \(k,d\ge0\), put \[
n=d+k+1,\qquad
\lambda=(d+3k+2,d+2k+1,d+1,d+1,d).
\tag{Q.1}
\] The parts sum to \(5n\) and are weakly decreasing. Let
\(\mu\vdash5\), and put \(\nu=\mu\) for odd \(d\) and \(\nu=\mu'\) for
even \(d\). For an integer \(R\), define the bounded polytope \[
\mathscr P(R)=\{x\in\mathbf R_{\ge0}^4:
2x_1+4x_2+6x_3+10x_4=R\},\qquad
N(R)=\#(\mathscr P(R)\cap\mathbf Z^4).
\tag{Q.2}
\] It is empty for \(R<0\). The following formula is positive, with
fixed nonnegative integers \(h_{\nu,j}\) displayed in §6: \[
\boxed{p_{\mu,(n)}^\lambda
=\sum_{j=0}^{19}h_{\nu,j}N(k-j).}
\tag{Q.3}
\] Equivalently, use \(h_{\nu,j}\) distinct tags for the polytope
\(\mathscr P(k-j)\). There are at most 216 tagged polytopes, and exactly
36 tags for the symmetric or alternating outer partition. Their
dimension, coefficients, and number are independent of both parameters.
The evaluator in §7 uses a fixed number of integer arithmetic
operations.

The proof constructs the relevant geometric section module, its explicit
matrix presentation, and an equivariant free multiplication map over
four invariant parameters. The finite character calculation only
determines the graded types of the actual quotient module at the end of
this construction.

\subsection{1. The exact flag comparison and its determinant
twist}\label{the-exact-flag-comparison-and-its-determinant-twist}

Let \(E=\mathfrak h_5\), with the invariant pairing and projective-frame
vectors \(\ell_i=e_i-\frac15\mathbf1\) fixed in G-015. On
\(\mathbf P(E)\) let \(Q\) be the tautological rank-three quotient, and
let \(\mathcal I\) be the ideal sheaf of the five points \([\ell_i]\).
Define \[
M_k=H^0\bigl(\mathbf P(E),\mathcal I^k(2k)\otimes Q\bigr),
\qquad M=\bigoplus_{k\ge0}M_k.
\tag{Q.4}
\] G-015 gives the actual balanced-chart map for (Q.1), since
\(\lambda_2-n=k\). After removing the full determinant factor
\((\det E^*)^d\), the relative quotient-flag weight is \((1,1,0)\) and
the first-line weight is \(2k+1\). The relative Plücker evaluation, the
exterior pairing \(\Lambda^2 Q^*\simeq Q\otimes\det Q^*\), and
\(\det Q^*\simeq\det E^*\otimes\mathcal O(-1)\) therefore identify the
pushforward with \[
(\det E^*)^{d+1}\otimes\mathcal O(2k)\otimes Q.
\] These maps are exactly those in (5F.6)--(5F.9) for relative weight
\((1,1,0)\); their local quotient-frame definitions glue and commute
with tensoring the base ideal, as proved in T-039 §1. Consequently \[
\boxed{(S_\lambda\mathbf Q^5)_{(n^5)}
\simeq\operatorname{sgn}^{d+1}\otimes M_k.}
\tag{Q.5}
\] The inverse chart sends a flag polynomial \(f\) to \(u^{-k}f(v(X))\).
Its extension to a polynomial is equivalent to the ideal membership in
(Q.4) by the transverse-order proof G-015. Thus (Q.5) includes the map,
the sheaf condition, and its group linearization.

\subsection{2. An invariant section and five exact
relations}\label{an-invariant-section-and-five-exact-relations}

Use \(\ell_1,\ldots,\ell_4\) as a basis, and write its coordinates as
\(z_1,\ldots,z_4\). The fifth point has coordinates \([1:1:1:1]\). Write
\[
\begin{aligned}
f_A&=z_1z_3-z_1z_2,& f_B&=z_1z_4-z_1z_2,\\
f_C&=z_2z_3-z_1z_2,& f_D&=z_2z_4-z_1z_2,&
f_E&=z_3z_4-z_1z_2.
\end{aligned}
\tag{Q.6}
\] Let \(U=\langle f_A,f_B,f_C,f_D,f_E\rangle=[3,2]\), and set \[
S=\mathbf Q[A,B,C,D,E],\quad
F=AD(E-B-C)-BC(E-A-D),\quad C=S/(F).
\tag{Q.7}
\] Each variable in \(S\) has degree one. By T-038, substitution
\(A\mapsto f_A,\ldots,E\mapsto f_E\) has kernel exactly \((F)\), and
\(F\) has sign character. Hence \(C\) acts on (Q.4) by multiplication.

The class of \((z_1^2,z_2^2,z_3^2,z_4^2)\) modulo the tautological line
defines \[
\theta\in M_1.
\tag{Q.8}
\] At every coordinate point and at \((1,1,1,1)\) this vector is
proportional to the point vector, so it vanishes in the quotient fiber.
To verify its linearization, write \(q=\sum z_i\ell_i\) and
\(s=\sum z_i\). The \(S_5\)-equivariant quadratic map
\(q\mapsto(q_i^2-\frac15\sum_jq_j^2)_i\) has coordinates
\(z_i^2-\frac25sz_i\) in this basis. Its quotient class is (Q.8).
Therefore \(\theta\) is invariant under the actual \(S_5\) action.

Here is a small polynomial matrix that specifies all the relations. Put
\[
K=\begin{pmatrix}
0&E-B&0&-C&-BC+BD+CD-DE\\
0&A-E&D&0&-AC+AD-CD+CE\\
D-E&0&0&A&-AB+AD-BD+BE\\
E-C&0&-B&0&AB+AC-AE-BC\\
C-D&B-A&0&0&-AD+BC
\end{pmatrix},\qquad
P=\frac{\operatorname{adj}K}{F}.
\tag{Q.9}
\] All entries of \(\operatorname{adj}K\) are divisible by \(F\). The
first four rows of \(P\), denoted by \(\Phi\), have degree two; its last
row is \((-A,-B,-C,-D,-E)\). Expansion of the determinants gives \[
\det K=3F^2,\qquad PK=KP=3F I_5,\qquad\det P=81F^3.
\tag{Q.10}
\] These are identities in the polynomial ring \(S\). They are also
reproduced by exact symbolic arithmetic in
\texttt{\small five\_\allowbreak{}point\_\allowbreak{}q\_\allowbreak{}module.py}
and
\texttt{\small FIVE\_\allowbreak{}POINT\_\allowbreak{}Q\_\allowbreak{}RELATIONS.json};
(Q.9) specifies the matrix independently of those files.

For complete verification of the substitution relations, put \[
\begin{aligned}
h_A&=f_A(z_1+z_2+z_3-2z_4),&
h_C&=f_C(z_1+z_2+z_3-2z_4),\\
h_B&=f_B(z_1+z_2+z_4-2z_3),&
h_D&=f_D(z_1+z_2+z_4-2z_3),\\
h_E&=z_3z_4(z_3+z_4-2z_1-2z_2)
      -z_1z_2(z_1+z_2-2z_3-2z_4).
\end{aligned}
\tag{Q.11}
\] The twenty polynomial identities, for \(1\le i\le4\) and
\(j\in\{A,B,C,D,E\}\), are \[
f_jz_i^2-\Phi_{ij}(f_A,f_B,f_C,f_D,f_E)=z_i h_j.
\tag{Q.12}
\] They follow by substituting (Q.6) into the minors in (Q.9) and
expanding; the right sides (Q.11) give their explicit line corrections.
Thus the columns of \(P\) are actual relations on the constant sections
\(E\) and the section \(\theta\), not just relations on their
dimensions.

The relation space carries the representation \(U\). Indeed the map
\(E\otimes\operatorname{Sym}^2U\to M_2\) is injective. A kernel would be
a vector polynomial in the tautological line, hence a multiple \(z h\)
with \(h\in H^0(\mathcal I^2(3))\). A cubic double at all four
coordinate points is a linear combination of the four squarefree cubics.
At \((1,1,1,1)\), let their coefficients be \(c_i\), indexed by the
missing variable. Its four first derivatives are \(\sum_jc_j-c_i\).
Their simultaneous vanishing implies \(c_i=\sum_jc_j\) for every \(i\),
whence \(3\sum_jc_j=0\) and every coefficient is zero. Also
\(\operatorname{Sym}^2U\to C_2\) is injective because the first relation
\(F\) has degree three. This proves the asserted injectivity. It makes
the coefficient \(\Phi(f)\) in \(f\theta=\Phi(f)\) unique. Applying an
element of \(S_5\) to that equality shows that its five-dimensional
relation space is the image of the equivariant map from \(U\).

Define the graded module \[
\mathcal M=\operatorname{coker}\left[
S(-2)\otimes U\xrightarrow{P}
(S\otimes E)\oplus S(-1)\right].
\tag{Q.13}
\] The determinant (Q.10) makes the displayed map injective. Equations
(Q.12) define an equivariant graded map \(\mathcal M\to M\), and
\(PK=3FI\) shows that \(F\mathcal M=0\).

\subsection{3. The presentation injects into the section
module}\label{the-presentation-injects-into-the-section-module}

We first verify that \(\mathcal M\) has no torsion over the integral
domain \(C\). Let \(g\in S\) have nonzero image in \(C\), and suppose
\(gv=Pw\) for polynomial column vectors. Multiplication by
\(\operatorname{adj}P\) gives \[
g\operatorname{adj}P\,v=81F^3w.
\] The irreducible polynomial \(F\) does not divide \(g\). Unique
factorization implies that every entry of \(\operatorname{adj}P\,v\) is
divisible by \(F^3\). Therefore \(u=\operatorname{adj}P\,v/(81F^3)\) is
a polynomial vector and \(Pu=v\). The class of \(v\) in the cokernel is
zero, which proves torsion-freeness.

Its rank over \(C\) is three. Here is an elementary local calculation of
that rank. Localize \(S\) at the prime \((F)\), a discrete valuation
ring with uniformizer \(F\). Successive row and column operations clear
an entry of minimum valuation and then its row and column, giving the
diagonal elementary-divisor form of \(P\). The identity \(PK=3FI\)
implies that each diagonal exponent is at most one, and \(\det P=81F^3\)
implies that their sum is three. Thus exactly three diagonal entries
have exponent one. Reducing modulo \(F\) gives rank two for \(P\), and
dimension three for its cokernel over \(\operatorname{Frac}C\). For an
additional finite check, all 100 three-by-three minors of \(P\) are
divisible by \(F\), and its upper left two-by-two minor is
\(-3AB(AD-BC)\), which is not divisible by \(F\).

Embed \(C\) into \(\mathbf Q[z_1,z_2,z_3,z_4]\) by (Q.6), and pass to
the latter fraction field \(L\). Section evaluation embeds \(M\)
degreewise in the generic quotient \(L^4/Lz\), which has dimension
three. The four constant generators in (Q.13) span this quotient.
Consequently \(\mathcal M\otimes_C L\to L^4/Lz\) is a surjection between
three-dimensional spaces and is an isomorphism. The kernel of
\(\mathcal M\to M\) is therefore \(C\)-torsion and, by the preceding
paragraph, is zero. This proves the actual injection \[
\mathcal M\hookrightarrow M.
\tag{Q.14}
\]

\subsection{4. The cohomological dimension proves surjectivity in every
degree}\label{the-cohomological-dimension-proves-surjectivity-in-every-degree}

The precise regularity result needed here is proved in A-029: for the
ideal (Q.6), every \(I^k\), \(k\ge2\), has a \(2k\)-linear free
resolution. A-029 verifies its height-three Gorenstein presentation,
every local generator condition, and the exceptional degree in the known
symmetric-power complexes. Its cited inputs are Kustin--Polini--Ulrich,
\emph{The equations defining blowup algebras of height three Gorenstein
ideals}, Algebra \& Number Theory 11 (2017), formula (6.1.4), p.~1509,
Lemma 6.3 and Corollary 6.4, p.~1510, with the proof on p.~1511. The
ordinary-power resolution follows there by the actual torsion sequence,
whose kernel is the degree-zero module
\(\operatorname{sgn}\otimes\operatorname{Sym}^{k-3}U\).

In particular the saturation agrees with the ordinary power in degree
\(2k\), and \(H^1(\mathcal I^k(2k-1))=0\) for \(k\ge2\). Thus \[
\begin{aligned}
h^0(\mathcal I^k(2k))
 &=\dim C_k=\binom{k+4}{4}-\binom{k+1}{4},\\
h^0(\mathcal I^k(2k-1))
 &=\binom{2k+2}{3}-5\binom{k+2}{3}
 =\frac{k(k+1)(k-2)}2 \quad(k\ge2).
\end{aligned}
\tag{Q.15}
\] The second equality uses the length of five order-\(k\) fat points,
which is \(5\binom{k+2}{3}\). Tensoring the tautological sequence with
\(\mathcal I^k(2k)\) and taking global sections now gives \[
\dim M_k=4\dim C_k-\frac{k(k+1)(k-2)}2
=\frac{(k+1)(k+2)(3k+4)}2\quad(k\ge2).
\tag{Q.16}
\] At \(k=1\), the degree-one evaluation \(E\to\operatorname{Perm}_5\)
has a one-dimensional trivial cokernel, and \(H^0(\mathcal I(1))=0\).
The same cohomology sequence gives \(\dim M_1=4\cdot5+1=21\). At
\(k=0\), \(M_0=E\) has dimension four. These values agree with (Q.16).

The injective free presentation (Q.13) has Hilbert series \[
\operatorname{Hilb}_{\mathcal M}(t)
=\frac{4+t-5t^2}{(1-t)^5}
=\frac{4+5t}{(1-t)^4}
=\sum_{k\ge0}\frac{(k+1)(k+2)(3k+4)}2t^k.
\tag{Q.17}
\] Equality of every graded dimension in (Q.14) proves
\(\mathcal M\simeq M\). This establishes generation and the entire
relation module; neither was inferred solely from a low-degree test.

\subsection{5. Four invariant parameters and the actual free
multiplication
map}\label{four-invariant-parameters-and-the-actual-free-multiplication-map}

Use the six linear forms of T-038: \[
\begin{array}{lll}
t_1=A+B-C+D-E,&t_2=-A+B+C-D-E,&t_3=-A+B+C-D+E,\\
t_4=-A-B-C+D+E,&t_5=A-B-C-D+E,&t_6=A-B+C+D-E.
\end{array}
\tag{Q.18}
\] They satisfy \(\sum t_i=0\) and \(\sum t_i^3=-24F\), so their third
elementary symmetric polynomial is \(e_3=-8F\). The verified action is
sign times the permutation action on the six Sylow-five subgroups. For a
homogeneous elementary symmetric polynomial \(e_j\), the sign twist
consequently multiplies it by \(\operatorname{sgn}^j\). Hence \[
B=\mathbf Q[e_2,e_4,e_6,e_5^2]
\tag{Q.19}
\] consists of invariant elements, of degrees \(2,4,6,10\).

We record an algebraic freeness proof sufficient for the regularity
assertion about this module. For six independent variables, start with
the universal monic polynomial having coefficients \(e_1,\ldots,e_6\).
Adjoin its first ordered root, divide the polynomial by that root
factor, and successively adjoin a root of the remaining monic polynomial
of degrees five through one. At each step the quotient by a monic
polynomial is free over the preceding coefficient ring, with the powers
below that degree as a basis. The resulting ring is
\(\mathbf Q[t_1,\ldots,t_6]\): its defining equations identify the
coefficients with the elementary symmetric polynomials and allow those
coefficients to be eliminated. This proves freeness of rank \(6!\) over
\(\mathbf Q[e_1,\ldots,e_6]\), including algebraic independence.

Quotient by \(e_1\). The linear forms in (Q.18) have rank five, so
\(S\simeq\mathbf Q[t_1,\ldots,t_6]/(e_1)\). Replace the polynomial
variable \(e_5\) by \(e_5^2\), retaining the free basis \(1,e_5\).
Finally substitute \(e_3=-8F\). It follows that \(S\) is free of rank
1440 over the polynomial ring \(B[F]\). In particular, after quotienting
by any initial segment of the four parameters in (Q.19), both the next
parameter and \(F\) are nonzerodivisors in the resulting free
coefficient module.

The four parameters are a regular sequence on \(M\). To prove this
directly, work modulo any preceding parameters and denote the next by
\(p\). The matrix \(P\) remains injective even after quotient by \(p\),
because its determinant is \(81F^3\), a nonzerodivisor there;
multiplying by its adjugate proves injectivity. In a relation \(pv=Pw\),
reduce modulo \(p\) to conclude \(w=pw'\). Since \(p\) is a
nonzerodivisor on the ambient free modules, it follows that \(v=Pw'\).
This proves that multiplication by \(p\) has zero kernel on the
cokernel. The argument applies in turn to all four parameters.

Set \(H=M/B_+M\). It is finite-dimensional: \(M\) is finite over
\(C=S/(F)\), which is finite over \(B\). The regular sequence gives \[
\operatorname{Hilb}_H(t)
=(4+5t)\prod_{r\in\{2,4,6,10\}}(1+t+\cdots+t^{r-1}).
\tag{Q.20}
\] It has degree 19 and total dimension 4320. Choose any graded linear
section \(H\to M\) and average it over \(S_5\). Multiplication defines
an equivariant map \[
B\otimes H\longrightarrow M.
\tag{Q.21}
\] It is surjective by induction on degree, because the section lifts
the quotient modulo \(B_+M\). The regular-sequence Hilbert identity
makes source and target dimensions equal in every degree. Thus (Q.21) is
an actual isomorphism. All four variables in \(B\) have trivial action.

\subsection{6. The fixed positive tag
table}\label{the-fixed-positive-tag-table}

Here is the complete table \(h_{\rho,j}=[H_j:[\rho]]\). Each row is
listed in increasing degree \(j=0,\ldots,19\):

\begin{longtable}[]{@{}
  >{\raggedright\arraybackslash}p{(\columnwidth - 4\tabcolsep) * \real{0.3000}}
  >{\raggedright\arraybackslash}p{(\columnwidth - 4\tabcolsep) * \real{0.3000}}
  >{\raggedleft\arraybackslash}p{(\columnwidth - 4\tabcolsep) * \real{0.4000}}@{}}
\toprule\noalign{}
\begin{minipage}[b]{\linewidth}\raggedright
\(\rho\)
\end{minipage} & \begin{minipage}[b]{\linewidth}\raggedright
\((h_{\rho,0},\ldots,h_{\rho,19})\)
\end{minipage} & \begin{minipage}[b]{\linewidth}\raggedleft
total
\end{minipage} \\
\midrule\noalign{}
\endhead
\bottomrule\noalign{}
\endlastfoot
\((5)\) & \(0,1,1,1,2,2,3,3,3,4,3,3,3,2,2,1,1,1,0,0\) & 36 \\
\((4,1)\) & \(1,1,2,5,6,9,11,12,14,14,14,13,12,10,8,6,3,2,1,0\) & 144 \\
\((3,2)\) & \(0,1,3,5,8,11,13,16,17,17,18,17,15,13,10,7,5,2,1,1\) &
180 \\
\((3,1,1)\) & \(0,1,3,5,9,12,16,19,20,22,21,21,19,16,13,9,6,3,1,0\) &
216 \\
\((2,2,1)\) & \(0,1,2,4,7,10,13,15,17,18,18,17,16,14,11,8,5,3,1,0\) &
180 \\
\((2,1,1,1)\) & \(0,0,1,3,5,8,10,12,14,14,15,14,13,12,9,7,4,2,1,0\) &
144 \\
\((1^5)\) & \(0,0,0,1,1,1,3,3,3,4,3,4,4,2,3,2,1,1,0,0\) & 36 \\
\end{longtable}

For reproducibility, the seven rational characters of \(M\) calculated
from (Q.13), in the indicated class order, are \[
\begin{array}{c|c|r}
\text{cycle type}&\sum_k\operatorname{tr}(g\mid M_k)t^k&\text{class size}\\\hline
1^5&(4+5t)/(1-t)^4&1\\
2\,1^3&(2-t)/((1-t)^3(1+t))&10\\
2^2\,1&t/(1-t^2)^2&15\\
3\,1^2&1/(1-t^3)&20\\
3\,2&-1/(1-t^3)&20\\
4\,1&t/(1-t^4)&30\\
5&-(1-t)/(1-t^5)&24.
\end{array}
\tag{Q.22}
\] Indeed the character before cancellation is \[
\frac{\chi_E(g)+t-\chi_U(g)t^2}{\det(1-tg\mid U)},
\quad \chi_E(g)=m_1-1,\quad
\chi_U(g)=\binom{m_1}{2}+m_2-m_1,
\tag{Q.23}
\] where \(m_i\) counts \(i\)-cycles. The latter character follows from
the permutation action on two-element subsets minus the five-point
permutation module, as identified in T-038. Newton's identities applied
to \(\chi_U(g^r)\), \(1\le r\le5\), give every denominator in (Q.22).
Multiplying by \(D(t)=(1-t^2)(1-t^4)(1-t^6)(1-t^{10})\) gives the
polynomial character of the actual module \(H\), by (Q.21). Inner
products with the fixed seven irreducible characters give the displayed
table. The exact scripts
\texttt{\small five\_\allowbreak{}point\_\allowbreak{}q\_\allowbreak{}tags.py}
and
\texttt{\small check\_\allowbreak{}q\_\allowbreak{}module\_\allowbreak{}characters.py}
perform this finite calculation by independent expressions, and compare
all 140 entries. No table entry depends on \(k\) or \(d\).

At \(t=1\), the six nonidentity polynomial characters vanish and the
identity value is 4320. Thus, as an ungraded module, \(H\) has character
\(36\) times the regular character of \(S_5\). This checks all row
totals: \(\sum_jh_{\rho,j}=36\dim[\rho]\). Taking the \([\nu]\)-summand
of (Q.21) and using (Q.5) proves (Q.3).

\subsection{7. Exact support and a fixed-arithmetic
evaluator}\label{exact-support-and-a-fixed-arithmetic-evaluator}

All parameter weights in (Q.2) are even, and one is two. The first two
consecutive nonzero entries in each row of the table therefore determine
the complete support: \[
\begin{array}{c|c}
\nu& p_{\mu,(n)}^\lambda>0\\\hline
(4,1)&k\ge0\\
(5),(3,2),(3,1,1),(2,2,1)&k\ge1\\
(2,1,1,1)&k\ge2\\
(1^5)&k\ge3.
\end{array}
\tag{Q.24}
\] Before the indicated first degree every tag vanishes. At and above
it, choose a tag of the correct parity among those first two degrees and
use the weight-two variable to fill the remainder. This proves both
directions. In particular, for symmetric outer degree five, positivity
holds exactly for odd \(d\) with \(k\ge1\), or even \(d\) with
\(k\ge3\).

For completeness the count in (Q.2) needs no loop of length depending on
the input. It is zero for odd \(R\). For \(R=2N\ge0\), its weights are
\(1,2,3,5\). Write each variable uniquely as \[
x_w=(30/w)q_w+r_w,\qquad0\le r_w<30/w.
\] Define the fixed positive polynomial \[
G(t)=\prod_{w\in\{1,2,3,5\}}
\left(\sum_{r=0}^{30/w-1}t^{wr}\right)
=\sum_{j=0}^{109}g_jt^j.
\tag{Q.25}
\] Its coefficients count exactly the residue tuples. For a tuple of
residue weight \(j\), the four nonnegative \(q_w\) sum to \((N-j)/30\),
so stars and bars gives the positive identity \[
N(2N)=
\sum_{\substack{0\le j\le\min(N,109)\\j\equiv N\pmod{30}}}
g_j\binom{(N-j)/30+3}{3}.
\tag{Q.26}
\] There are at most four terms. The notation on the left is the
counting function of (Q.2); the integer \(N\) on the right is half its
argument. The factor degrees in (Q.25) are \(29,28,27,25\), whose sum is
109. Thus (Q.3) uses a fixed table, at most 80 cubic-binomial
evaluations, and integer inputs of bit length
\(O(\log(k+2)+\log(d+2))\). Ordinary integer arithmetic gives a
polynomial bit-cost bound; with schoolbook multiplication the count
itself costs \(O(\log^2(k+2))\).

\subsection{8. Algebraic-geometric and lambda-ring
dictionary}\label{algebraic-geometric-and-lambda-ring-dictionary}

The object counted by (Q.3) is the \([\mu]\)-covariant summand of the
section space of the explicit ideal-twisted flag line bundle in G-015.
The morphism (Q.5) carries that section space to the balanced weight of
the Schur functor, and the normalizer projector T-003 carries it to the
plethysm coefficient. In the representation category its class is
\([S_\mu(\operatorname{Sym}^n)]\), with composition interpreted in the
lambda-ring of polynomial representations. The equality is lifted here
to actual modules and maps by (Q.5), (Q.13), and (Q.21).

The geometric construction is defined over \(\mathbf Q\); the frame,
quadrics, cubic, and polynomial matrices are defined over
\(\mathbf Z[1/30]\) after the stated invariant-form identifications. The
averaging section additionally uses \(1/120\), already in that ring. The
characteristic-zero cohomology and representation conclusions have been
proved over \(\mathbf Q\) and by extension over characteristic-zero
fields; no assertion of all-fiber exactness in positive characteristic
is inferred from the rational proof. The exact base-extension morphism
is tensoring the displayed complexes and maps.

The ideal powers are those of a specified closed subscheme of projective
space, and all bundles are Schur functors of tautological locally free
sheaves. The determinant pairings and exact-sequence classes are the
operations in the SGA 6 dictionary, and the concrete ideal and its
graded quotients define the associated Hilbert functors in the FGA
dictionary. The exact archive/SGA/FGA locators and coefficient-category
maps are recorded in the separate geometric dictionary; this result uses
those same maps, with the additional explicit matrices (Q.9). No new
weight, purity, or motivic comparison is required by this proof.

\section{The exterior quotient bundle on the five-point
ray}\label{the-exterior-quotient-bundle-on-the-five-point-ray}

\subsection{G-018. An exact kernel and positive
formula}\label{g-018.-an-exact-kernel-and-positive-formula}

For \(k,d\ge0\), put \[
n=d+k+1,\qquad
\lambda=(d+3k+3,d+2k+1,d+1,d,d).
\tag{W.1}
\] For every \(\mu\vdash5\), the coefficient \(p_{\mu,(n)}^\lambda\) is
a count in at most 240 fixed bounded rational polytopes, with at most
five integer variables. The fixed table and every inequality are given
below. Its support is particularly simple: after setting \(\nu=\mu\) for
odd \(d\) and \(\nu=\mu'\) for even \(d\), it is positive exactly in the
following ranges: \[
\begin{array}{c|c}
\nu&\text{range}\\\hline
(3,1,1)&k\ge0\\
(4,1),(3,2),(2,2,1),(2,1,1,1)&k\ge1\\
(5),(1^5)&k\ge3.
\end{array}
\tag{W.2}
\] In particular the symmetric outer coefficient is positive exactly
when \(k\ge3\), independently of \(d\).

\subsection{1. The flag morphism and the multiplication
kernel}\label{the-flag-morphism-and-the-multiplication-kernel}

Use \(E=\mathfrak h_5\), its invariant form, the five-point ideal
\(\mathcal I\), and the tautological quotient bundle \(Q\) from G-015
and T-038. Define \[
W_k=H^0(\mathbf P(E),\mathcal I^k(2k)\otimes\Lambda^2Q).
\tag{W.3}
\] For (W.1), the point order is \(\lambda_2-n=k\), the first-line twist
after relative duality is \(\lambda_2-\lambda_3=2k\), and the quotient
Schur weight is \((\lambda_3-\lambda_5,\lambda_3-\lambda_4)=(1,1)\).
Explicitly, after the determinant factor \((\det E^*)^d\) the relative
flag weight is \((1,0,0)\), with first-line twist \(2k+1\). The exterior
pairing \(Q^*\simeq\Lambda^2Q\otimes\det Q^*\) and
\(\det Q^*=\det E^*\otimes\mathcal O(-1)\) therefore give \[
(S_\lambda\mathbf Q^5)_{(n^5)}\simeq\operatorname{sgn}^{d+1}\otimes W_k.
\tag{W.4}
\] These relative Schur evaluations and the inverse balanced chart
\(f\mapsto u^{-k}f(v(X))\) are the actual maps in G-015 and T-039 §1.

Set \(C_{j,k}=H^0(\mathcal I^k(2k+j))\), for \(j=1,2\). A-029 proves
\(C_{j,k}=(I^k)_{2k+j}\) for \(k\ge2\); saturation of \(I\) proves the
same for \(k=1\), and \(k=0\) is immediate. The ideal \(I^k\) is
generated in degree \(2k\). Hence multiplication by linear forms gives a
surjection \[
E^*\otimes C_{1,k}\longrightarrow C_{2,k}.
\tag{W.5}
\] In coordinates it sends \(\sum e_i^*\otimes f_i\) to \(\sum z_if_i\).
The dual tautological sequence, tensored by \(\mathcal I^k(2k+1)\),
identifies its kernel with \(H^0(\mathcal I^k(2k+1)\otimes Q^*)\). The
exterior isomorphism \(\Lambda^2Q\simeq Q^*\otimes\det Q\), where
\(\det Q=\operatorname{sgn}\otimes\mathcal O(1)\), consequently gives
the actual exact sequence \[
\boxed{0\longrightarrow\operatorname{sgn}\otimes W_k
\longrightarrow E^*\otimes C_{1,k}
\longrightarrow C_{2,k}\longrightarrow0.}
\tag{W.6}
\] The first map is the dual tautological inclusion followed by the
specified exterior pairing. Thus the character difference used next is
the trace of a concrete kernel of a surjective map.

\subsection{2. Exact traces from the five fat-point
jets}\label{exact-traces-from-the-five-fat-point-jets}

For a permutation \(g\in S_5\), let \(m_1(g)\) be its number of fixed
letters and let
\(s_r(g)=\operatorname{tr}(g\mid\operatorname{Sym}^rE)\), with \(s_r=0\)
for \(r<0\). These traces are completely specified by \[
\sum_{r\ge0}s_r(g)z^r
=\frac{1-z}{\prod_{\ell\in\operatorname{cyc}(g)}(1-z^\ell)}.
\tag{W.7}
\] A-029 gives \(H^1(\mathcal I^k(2k+j))=0\) for \(j\ge0\), including
the separately verified \(k=0,1\). Therefore the evaluation sequence for
the fat points is exact on global sections in the degrees used here.

We specify the jet action, including its trivializations. The stabilizer
of \([\ell_i]\) fixes the affine vector \(\ell_i\) and decomposes
\(E=\mathbf Q\ell_i\oplus\mathfrak h_4\). The linear functional equal to
one on \(\ell_i\) and zero on \(\mathfrak h_4\) is invariant under that
stabilizer. Its affine chart trivializes every \(\mathcal O(c)\) with
trivial fiber character, and its normal coordinates transform by
\(\mathfrak h_4^*\). Thus the order-\(k\) jet module at that point is
\(\bigoplus_{r=0}^{k-1}\operatorname{Sym}^r\mathfrak h_4^*\). Transport
by the canonical affine representatives gives the induced \(S_5\)-module
of these jets. Only fixed points contribute to the trace of \(g\); each
contributes the same value. For a fixed point the restriction
\(E=\mathbf1\oplus\mathfrak h_4\) identifies the sum of the first \(k\)
symmetric powers with the trace of \(\operatorname{Sym}^{k-1}E\). Hence
the full jet trace is \[
\operatorname{tr}(g\mid J_k)=m_1(g)s_{k-1}(g).
\tag{W.8}
\] This proves, through the actual evaluation sequence, \[
\operatorname{tr}(g\mid C_{j,k})=s_{2k+j}(g)-m_1(g)s_{k-1}(g).
\tag{W.9}
\] Combining (W.6) and (W.9) gives \[
w_k(g)=\operatorname{sgn}(g)\left[
(m_1-1)\{s_{2k+1}-m_1s_{k-1}\}
-\{s_{2k+2}-m_1s_{k-1}\}\right].
\tag{W.10}
\] Every scalar on the right is specified by the finite rational
function (W.7). At the identity this gives \[
\dim W_k=\frac{(k+2)(k+3)(3k+2)}2,
\qquad
\sum_{k\ge0}\dim W_k t^k
=\frac{6+6t-4t^2+t^3}{(1-t)^4}.
\tag{W.11}
\] The initial value \(W_0=\Lambda^2E\) has label \((3,1,1)\), as
required.

\subsection{3. Fixed positive Newton data and its complete finite
certificate}\label{fixed-positive-newton-data-and-its-complete-finite-certificate}

Let \[
b_\nu(k)=\frac1{120}\sum_{g\in S_5}\chi_\nu(g)w_k(g).
\tag{W.12}
\] This is the multiplicity of \([\nu]\) in the actual module \(W_k\).
The sum can equivalently be taken over the seven cycle types with sizes
\(1,10,15,20,20,30,24\), in the order of (Q.22).

For \(0\le r<60\) and \(0\le j\le3\), define the fixed integer \[
c_{\nu,r,j}
=\sum_{h=0}^j(-1)^{j-h}\binom jh b_\nu(r+60h).
\tag{W.13}
\] These 1680 integers are specified by (W.7), (W.10), (W.12), and
(W.13), using only symmetric groups of order at most 120 and indices at
most 239.
\texttt{\small FIVE\_\allowbreak{}POINT\_\allowbreak{}WEDGE2\_\allowbreak{}NEWTON\_\allowbreak{}TABLE.json}
lists every one.

Here is the exact verification that extends this fixed data to every
input. Taking even and odd parts of (W.7) gives the rational series
\(\sum_ks_{2k}t^k\) and \(\sum_ks_{2k+1}t^k\). Subtracting its constant
term and dividing by \(t\) gives the series for \(s_{2k+2}\). For each
of the seven explicitly listed cycle types, substitution into (W.10)
shows that its denominator divides \((1-t^{60})^4\) and its numerator
has degree at most 239. These seven exact polynomial divisibilities are
checked in
\texttt{\small five\_\allowbreak{}point\_\allowbreak{}wedge2\_\allowbreak{}certificate.py},
with their reduced rational functions retained in the JSON certificate.
The script also computes (W.7) directly by the seven cycle recurrences,
independently of the even/odd rational extraction.

A numerator degree at most 239 implies that, on a fixed residue \(r\),
the coefficient is a polynomial of degree at most three in
\(q=(k-r)/60\), for every \(q\ge0\). To check the initial endpoints,
each numerator exponent is \(r+60h\), \(0\le h\le3\), and its
contribution is \(\binom{q-h+3}{3}\). The polynomial expression is zero
also for \(0\le q<h\le3\), since its upper integer is then zero, one, or
two. Thus no delayed-onset term is lost. Newton interpolation at
\(q=0,1,2,3\) proves the identity \[
\boxed{b_\nu(60q+r)=\sum_{j=0}^3 c_{\nu,r,j}\binom qj\quad(q\ge0).}
\tag{W.14}
\]

All 1680 integers (W.13) were computed exactly and are nonnegative; the
following minima over all 60 residues make their signs reviewable.

\begin{longtable}[]{@{}lrrrr@{}}
\toprule\noalign{}
\(\nu\) & min \(c_0\) & min \(c_1\) & min \(c_2\) & min \(c_3\) \\
\midrule\noalign{}
\endhead
\bottomrule\noalign{}
\endlastfoot
\((5)\) & 0 & 2877 & 16560 & 16200 \\
\((4,1)\) & 0 & 11693 & 66540 & 64800 \\
\((3,2)\) & 0 & 14725 & 83400 & 81000 \\
\((3,1,1)\) & 1 & 17787 & 100260 & 97200 \\
\((2,2,1)\) & 0 & 14880 & 83700 & 81000 \\
\((2,1,1,1)\) & 0 & 12003 & 67140 & 64800 \\
\((1^5)\) & 0 & 3032 & 16860 & 16200 \\
\end{longtable}

The accompanying TSV has one row for each label and residue, with all
four entries. Its generating code contains the full exact finite check,
not a random or asymptotic test. Direct evaluation of (W.12) through
\(k=359\) also agrees with (W.14) in every case. This extra finite check
is recorded separately from the rational identity that proves the
formula in all degrees.

For each tag \((r,j)\) with \(c=c_{\nu,r,j}>0\), use integer variables
\(q,y,t_1,\ldots,t_j\) and inequalities \[
60q+r=k,\quad q\ge0,\quad0\le y<c,\quad
1\le t_1<t_2<\cdots<t_j\le q.
\tag{W.15}
\] For \(j=0\), the last list is empty. Replace strict integer
inequalities by \(y\le c-1\) and \(t_i+1\le t_{i+1}\) to obtain the
rational polytope. Its lattice count is exactly \(c\binom qj\), by the
bijection with a choice of \(y\) and a \(j\)-element subset of
\(\{1,\ldots,q\}\). Every variable is bounded. Summing these 240 tagged
sets proves the positive formula for (W.1) using the sign twist (W.4).

The fixed zeroth entries in residues zero through 59 have zeros exactly
at the excluded degrees in (W.2). For \(k\ge60\), the positive first
differences in the table and \(q\ge1\) prove positivity. This proves the
complete support statement, including all endpoints.

\subsection{Geometric transport}\label{geometric-transport}

This result is the global-section representation of the ideal-twisted
exterior quotient bundle (W.3). The flag chart (W.4), dual tautological
inclusion (W.6), and affine jet trivializations (W.8) specify every
comparison. Associated bundles along the frame torsor carry the Schur
operations to the SGA 6 representation-group dictionary; the point
scheme, ideal powers, and their evaluation maps give the EGA/FGA
objects. The exact framework locators are listed in the separate
geometric dictionary. All representation assertions here are over
characteristic-zero fields.

\section{Every symmetric quotient power on the five-point quadratic
ray}\label{every-symmetric-quotient-power-on-the-five-point-quadratic-ray}

\subsection{G-019. A three-parameter positive plethysm
family}\label{g-019.-a-three-parameter-positive-plethysm-family}

Let \(a,k,d\ge0\). Define \[
n=d+a+k,\qquad
\lambda=(d+2a+3k,d+a+2k,d+a,d+a,d).
\tag{SQ.1}
\] Zero parts are removed; at \(a=k=d=0\) the partition is empty. For
every \(\mu\vdash5\), the coefficient \(p_{\mu,(n)}^\lambda\) has a
positive formula in a fixed family of at most 57,600 bounded rational
polytopes of dimension at most nine. At most sixteen tags contribute
after the two residues \(a\bmod60\) and \(k\bmod60\) are selected. The
evaluator requires a fixed number of integer arithmetic operations on
the input parameters. All constants are explicitly defined in
(SQ.10)--(SQ.12) and supplied in a complete exact table; no
input-dependent character or plethysm computation is required.

This evaluates all symmetric quotient powers together with all point
orders in the geometric model, and includes T-039 at \(k=1\), G-017 at
\(a=1\), and the full scalar ray at \(a=0\). The underlying module
comparison is proved for every parameter, before the finite positive
coefficient conversion.

\subsection{1. The actual sheaf module and all cohomological
endpoints}\label{the-actual-sheaf-module-and-all-cohomological-endpoints}

Use the five-point scheme \(Z\subset\mathbf P(E)\), with
\(E=\mathfrak h_5\), ideal \(\mathcal I\), and quotient bundle \(Q\).
Put \[
M_{a,k}=H^0\bigl(\mathbf P(E),\mathcal I^k(2k)
                                      \otimes\operatorname{Sym}^aQ\bigr),
\qquad C_k=H^0(\mathcal I^k(2k)).
\tag{SQ.2}
\] For (SQ.1), the flag model has point order \(k\), first-line gap
\(2k\), and relative weight \((a,a,0)\). The exterior pairing used in
(5F.7) identifies
\(S_{(a,a)}Q^*\simeq(\det Q^*)^a\otimes\operatorname{Sym}^aQ\). The full
determinant and relative determinants therefore give the actual
equivariant isomorphism \[
\boxed{(S_\lambda\mathbf Q^5)_{(n^5)}
\simeq\operatorname{sgn}^{d+a}\otimes M_{a,k}.}
\tag{SQ.3}
\] Its relative frame maps and inverse chart are the maps in G-015 and
T-039 §1, with \(u^{-k}\) replacing \(u^{-1}\). In particular the
point-power condition has already been proved to be exactly the
polynomial extension condition; (SQ.3) retains its full linearization.

For \(a\ge1\), the symmetric quotient sequence is the actual exact
sequence of sheaves \[
0\to\mathcal I^k(2k-1)\otimes\operatorname{Sym}^{a-1}E
\longrightarrow\mathcal I^k(2k)\otimes\operatorname{Sym}^aE
\longrightarrow\mathcal I^k(2k)\otimes\operatorname{Sym}^aQ\to0.
\tag{SQ.4}
\] Its first map multiplies by the tautological vector, and its last map
is the symmetric power of the quotient. Before tensoring the ideal, it
is an exact locally split sequence of vector bundles: on a line
\(L\subset E\), the kernel in the symmetric algebra is
\(L\operatorname{Sym}^{a-1}E\), and multiplication by the nonzero line
vector is injective. Tensoring by \(\mathcal I^k(2k)\) preserves this
locally split exactness.

Let \(A_k=H^0(\mathcal I^k(2k-1))\). A-029 proves
\(H^1(\mathcal I^k(2k-1))=0\) for \(k\ge2\), using the verified
Kustin--Polini--Ulrich complexes, formula (6.1.4), p.~1509, Lemma 6.3
and Corollary 6.4, p.~1510. The complete applied proof checks the
ideal's local hypotheses and the passage to ordinary powers. At \(k=0\),
\(\mathcal O(-1)\) is acyclic. At \(k=1\), \(A_1=0\),
\(H^1(\mathcal I(1))=\mathbf1\), and \(H^1(\mathcal I(2))=0\), by the
actual five-point evaluation maps. Consequently the long cohomology
sequence of (SQ.4) is short exact for \(k\ne1\); at \(k=1\) its sole
additional term is \(\operatorname{Sym}^{a-1}E\), exactly as in T-039
(5F.15).

The case \(a=1\) is the module \(M_k\) proved in G-017. Its own
tautological sequence gives the equality in the exact representation
group \[
[M_k]=[E][C_k]-[A_k]+\mathbf1_{k=1}[\mathbf1].
\tag{SQ.5}
\] Combining the actual sequences just described, including their
identical \(k=1\) connecting terms, proves for every \(a,k\ge0\) \[
\boxed{[M_{a,k}]
=[\operatorname{Sym}^aE][C_k]
+[\operatorname{Sym}^{a-1}E]\bigl([M_k]-[E][C_k]\bigr).}
\tag{SQ.6}
\] Here \(\operatorname{Sym}^{-1}E=0\); this convention proves the
formula also at \(a=0\). At \(k=0\) it reduces to
\(M_{a,0}=\operatorname{Sym}^aE\). Thus no unverified vanishing for
varying Schur bundles is used.

\subsection{2. A completely specified bivariate character
series}\label{a-completely-specified-bivariate-character-series}

Write \(e_g=\chi_E(g)\), \(u_g=\chi_{[3,2]}(g)\), and
\(\epsilon_g=\operatorname{sgn}(g)\). T-038 and A-029 identify
\(\bigoplus C_k\) with the cubic algebra \(\operatorname{Sym}(U)/(F)\),
whose relation has sign character. G-017 supplies its quotient-bundle
module with the proved presentation. Their exact trace series are \[
C_g(t)=\frac{1-\epsilon_gt^3}{\det(1-tg\mid U)},\qquad
M_g(t)=\frac{e_g+t-u_gt^2}{\det(1-tg\mid U)}.
\tag{SQ.7}
\] Taking the traces of (SQ.6) gives the exact rational identity \[
\boxed{\sum_{a,k\ge0}\chi_{M_{a,k}}(g)s^at^k
=\frac{1+st-u_gst^2-\epsilon_gt^3+
               \epsilon_ge_gst^3}
 {\det(1-sg\mid E)\det(1-tg\mid U)}.}
\tag{SQ.8}
\] The numerator follows directly by summing
\(C_g\operatorname{Sym}^aE+(M_g-e_gC_g)\operatorname{Sym}^{a-1}E\).

For a complete finite specification, use the following seven-column
data. Columns correspond to the ordinary cycle types in the first row.

\begin{longtable}[]{@{}
  >{\raggedright\arraybackslash}p{(\columnwidth - 14\tabcolsep) * \real{0.0968}}
  >{\raggedleft\arraybackslash}p{(\columnwidth - 14\tabcolsep) * \real{0.1290}}
  >{\raggedleft\arraybackslash}p{(\columnwidth - 14\tabcolsep) * \real{0.1290}}
  >{\raggedleft\arraybackslash}p{(\columnwidth - 14\tabcolsep) * \real{0.1290}}
  >{\raggedleft\arraybackslash}p{(\columnwidth - 14\tabcolsep) * \real{0.1290}}
  >{\raggedleft\arraybackslash}p{(\columnwidth - 14\tabcolsep) * \real{0.1290}}
  >{\raggedleft\arraybackslash}p{(\columnwidth - 14\tabcolsep) * \real{0.1290}}
  >{\raggedleft\arraybackslash}p{(\columnwidth - 14\tabcolsep) * \real{0.1290}}@{}}
\toprule\noalign{}
\begin{minipage}[b]{\linewidth}\raggedright
datum
\end{minipage} & \begin{minipage}[b]{\linewidth}\raggedleft
\(1^5\)
\end{minipage} & \begin{minipage}[b]{\linewidth}\raggedleft
\(2\,1^3\)
\end{minipage} & \begin{minipage}[b]{\linewidth}\raggedleft
\(2^2\,1\)
\end{minipage} & \begin{minipage}[b]{\linewidth}\raggedleft
\(3\,1^2\)
\end{minipage} & \begin{minipage}[b]{\linewidth}\raggedleft
\(3\,2\)
\end{minipage} & \begin{minipage}[b]{\linewidth}\raggedleft
\(4\,1\)
\end{minipage} & \begin{minipage}[b]{\linewidth}\raggedleft
\(5\)
\end{minipage} \\
\midrule\noalign{}
\endhead
\bottomrule\noalign{}
\endlastfoot
class size & 1 & 10 & 15 & 20 & 20 & 30 & 24 \\
\(\chi_{(5)}\) & 1 & 1 & 1 & 1 & 1 & 1 & 1 \\
\(\chi_{(4,1)}\) & 4 & 2 & 0 & 1 & -1 & 0 & -1 \\
\(\chi_{(3,2)}\) & 5 & 1 & 1 & -1 & 1 & -1 & 0 \\
\(\chi_{(3,1,1)}\) & 6 & 0 & -2 & 0 & 0 & 0 & 1 \\
\(\chi_{(2,2,1)}\) & 5 & -1 & 1 & -1 & -1 & 1 & 0 \\
\(\chi_{(2,1,1,1)}\) & 4 & -2 & 0 & 1 & 1 & 0 & -1 \\
\(\chi_{(1^5)}\) & 1 & -1 & 1 & 1 & -1 & -1 & 1 \\
\end{longtable}

The character table is the fixed table already verified in T-038 and
T-039. It can be checked by the ordinary permutation actions for
\(E,U\), their exterior powers and sign twists, followed by the
class-size inner products. In particular every entry in (SQ.8) is fixed.
The two determinants can be written without eigenvalue choices: \[
\det(1-sg\mid E)=
\frac{\prod_{\ell\in\operatorname{cyc}(g)}(1-s^\ell)}{1-s},
\quad
\det(1-tg\mid U)=
\frac{\prod_{\ell\in\operatorname{cyc}_6(g)}
                      (1-\epsilon_g^\ell t^\ell)}{1-\epsilon_gt}.
\tag{SQ.9}
\] Here the exact six-letter cycle list, by T-038's intertwiner, is
\(1^6,2^3,2^21^2,3^2,6,41^2,51\), in the same column order.

\subsection{3. Fixed integer constants and the proof of their
positivity}\label{fixed-integer-constants-and-the-proof-of-their-positivity}

Define the fixed-degree coefficient function \[
b_\nu(a,k)=\frac1{120}\sum_{\text{seven classes }g}
|g^{S_5}|\,\chi_\nu(g)\,[s^at^k]\text{(SQ.8)}.
\tag{SQ.10}
\] For \(0\le r,v<60\) and \(0\le i,j\le3\), set \[
\boxed{c_{\nu,r,v,i,j}=
\sum_{h=0}^i\sum_{l=0}^j
(-1)^{i+j-h-l}\binom ih\binom jl\,
b_\nu(r+60h,v+60l).}
\tag{SQ.11}
\] Equations (SQ.8)--(SQ.11), together with the displayed character
table, are explicit definitions of 403,200 fixed integers. Only
coefficients with indices below 240 occur. The complete table is stored
as
\texttt{\small FIVE\_\allowbreak{}POINT\_\allowbreak{}SYMMETRIC\_\allowbreak{}QUOTIENT\_\allowbreak{}TABLE.json},
in the order
\texttt{\small labels[\allowbreak{}nu]\allowbreak{}[\allowbreak{}r]\allowbreak{}[\allowbreak{}v]\allowbreak{}[\allowbreak{}4*i+\allowbreak{}j]\allowbreak{}}.
It is a single fixed table, independent of every input parameter in
(SQ.1).

The following exact finite certificate proves all their signs. For each
of the seven classes, compute coefficients of \((1-s)/\prod(1-s^\ell)\)
by multiplying the cycle geometric series. Similarly multiply the six
signed geometric series in (SQ.9), then multiply by \(1-\epsilon_gt\),
to obtain the coefficients \(u_k\) of \(\det(1-tg\mid U)^{-1}\). The two
required \(t\)-coefficient arrays are exactly \[
C_k=u_k-\epsilon_gu_{k-3},\qquad
X_k=u_{k-1}-u_gu_{k-2}+\epsilon_ge_gu_{k-3},
\tag{SQ.12}
\] with negative indices zero. The class coefficient in (SQ.10) is
\(s_aC_k+s_{a-1}X_k\). Multiplication by the displayed integer class
sizes and characters, summation and division by 120 give (SQ.10)
exactly; every division in the executed certificate had zero remainder.
Successive forward differences of the four-by-four blocks then give
(SQ.11).
\texttt{\small five\_\allowbreak{}point\_\allowbreak{}symmetric\_\allowbreak{}quotient\_\allowbreak{}certificate.py}
contains this entire algorithm, with all iteration limits fixed at 240,
60, four, and seven. All 403,200 resulting integers were checked
exactly.

The \((i,j)=(3,3)\) coefficient is zero in every case. Every other
coefficient with \(i+j>0\) is strictly positive. The following table
records the minimum over all 3600 residue pairs and all fourteen such
index pairs. The zeroth coefficients are nonnegative, both by this
calculation and because (SQ.10) is the multiplicity of the actual
module.

\begin{longtable}[]{@{}
  >{\raggedright\arraybackslash}p{(\columnwidth - 4\tabcolsep) * \real{0.2727}}
  >{\raggedleft\arraybackslash}p{(\columnwidth - 4\tabcolsep) * \real{0.3636}}
  >{\raggedleft\arraybackslash}p{(\columnwidth - 4\tabcolsep) * \real{0.3636}}@{}}
\toprule\noalign{}
\begin{minipage}[b]{\linewidth}\raggedright
\(\nu\)
\end{minipage} & \begin{minipage}[b]{\linewidth}\raggedleft
minimum of the fourteen positive difference coefficients
\end{minipage} & \begin{minipage}[b]{\linewidth}\raggedleft
number of zero zeroth entries
\end{minipage} \\
\midrule\noalign{}
\endhead
\bottomrule\noalign{}
\endlastfoot
\((5)\) & 418 & 2 \\
\((4,1)\) & 1487 & 2 \\
\((3,2)\) & 1735 & 2 \\
\((3,1,1)\) & 1978 & 5 \\
\((2,2,1)\) & 1575 & 5 \\
\((2,1,1,1)\) & 1167 & 9 \\
\((1^5)\) & 258 & 19 \\
\end{longtable}

The summary certificate retains the separate minima and maxima of all
sixteen difference positions, for each label. These are exhaustive
finite integer checks with the complete reproducible recurrences
(SQ.9)--(SQ.12), rather than evidence inferred from sampled parameters.

\subsection{4. The rational identity that extends the formula to all
parameters}\label{the-rational-identity-that-extends-the-formula-to-all-parameters}

We verify the range of the Newton interpolation exactly. For each class,
the four one-variable factors \[
\frac1{\det(1-sg\mid E)},\quad
\frac{s}{\det(1-sg\mid E)},\quad
C_g(t),\quad M_g(t)-e_gC_g(t)
\tag{SQ.13}
\] have denominator dividing \((1-s^{60})^4\) or \((1-t^{60})^4\),
respectively, and common numerator of degree at most 239. Substitution
of the seven cycle lists in (SQ.9), cancellation and polynomial division
verify these 28 identities. All numerator remainders are zero in the
exact certificate, and each numerator degree is checked explicitly.

The argument of G-018 §3 now applies in each variable. On a fixed
residue, a numerator term indexed by \(h\le3\) contributes
\(\binom{q-h+3}{3}\); it equals zero in all early cases \(0\le q<h\).
Thus the full residue sequence is a polynomial from its first value,
with degree at most three. The two products in (SQ.8) consequently give
a polynomial of degree at most three in each quotient parameter. Its
unique Newton expansion, fixed by its four-by-four grid of values, is
therefore \[
\boxed{b_\nu(60p+r,60q+v)
=\sum_{i=0}^3\sum_{j=0}^3
c_{\nu,r,v,i,j}\binom pi\binom qj
\quad(p,q\ge0).}
\tag{SQ.14}
\] This proves the identity for every input. The finite sign check in §3
and this rational identity perform distinct, specified jobs.

At the identity the dimension formula offers a useful independent
algebra check: \[
\dim M_{a,k}
=\binom{a+2}{2}\frac{(k+1)(k^2+2k+2)}2
+\binom{a+2}{3}(k+1)(2k+1).
\tag{SQ.15}
\] It follows by expanding (SQ.6) using (Q.16) and (Q.15), with the
\(k=1\) correction retained. Its bidegree is at most \((3,3)\), and its
coefficient of \(a^3k^3\) is zero, consistent with the vanishing mixed
third difference. The full character identity, rather than only this
dimension check, proves (SQ.14).

\subsection{5. Positive polytopes, evaluation cost, and complete
support}\label{positive-polytopes-evaluation-cost-and-complete-support}

For the input (SQ.1), set \(\nu=\mu\) when \(d+a\) is even and
\(\nu=\mu'\) when \(d+a\) is odd. For each fixed tag \((r,v,i,j)\) whose
table entry is \(c>0\), use integer variables
\(p,q,y,x_1,\ldots,x_i,z_1,\ldots,z_j\) satisfying \[
\begin{aligned}
60p+r&=a,&60q+v&=k,&p,q&\ge0,&0\le y&\le c-1,\\
1\le x_1&<\cdots<x_i\le p,
&&1\le z_1<\cdots<z_j\le q.
\end{aligned}
\tag{SQ.16}
\] An empty list of ordered variables contributes one choice. Each
strict integer inequality is written equivalently with a difference at
least one. The equations fix \(p,q\); every other variable lies in a
displayed finite interval. This is a bounded rational polytope with at
most nine variables. Its lattice points are in bijection with a choice
of \(y\), an \(i\)-element subset of \(\{1,\ldots,p\}\), and a
\(j\)-element subset of \(\{1,\ldots,q\}\). Its count is exactly the
corresponding positive term in (SQ.14). The disjoint tagged union of
(SQ.16), together with (SQ.3), proves the announced plethysm formula.

There are at most \(60^2\cdot16=57,600\) fixed tags. Only one residue
pair admits integer \(p,q\), so direct evaluation selects that pair and
uses at most sixteen products of binomial polynomials of degree at most
three. The table lookup is fixed. Integer bit lengths are
\(O(\log(a+2)+\log(k+2)+\log(d+2))\), and schoolbook arithmetic gives a
polynomial bit-cost bound, with the counting arithmetic quadratic in the
combined input bit length.

There is a complete small exceptional set for vanishing. All
coefficients are positive for \(a\ge60\) or \(k\ge60\), by the strictly
positive \((1,0)\) and \((0,1)\) differences. Inside the 60-by-60
square, the only zeros are the following 44 pairs \((a,k)\), with the
effective label \(\nu\) just specified:

\begin{longtable}[]{@{}
  >{\raggedright\arraybackslash}p{(\columnwidth - 2\tabcolsep) * \real{0.5000}}
  >{\raggedright\arraybackslash}p{(\columnwidth - 2\tabcolsep) * \real{0.5000}}@{}}
\toprule\noalign{}
\begin{minipage}[b]{\linewidth}\raggedright
\(\nu\)
\end{minipage} & \begin{minipage}[b]{\linewidth}\raggedright
all zero pairs \((a,k)\)
\end{minipage} \\
\midrule\noalign{}
\endhead
\bottomrule\noalign{}
\endlastfoot
\((5)\) & \((0,1),(1,0)\) \\
\((4,1)\) & \((0,0),(0,1)\) \\
\((3,2)\) & \((0,0),(1,0)\) \\
\((3,1,1)\) & \((0,0),(0,1),(0,2),(1,0),(2,0)\) \\
\((2,2,1)\) & \((0,0),(0,1),(1,0),(2,0),(3,0)\) \\
\((2,1,1,1)\) &
\((0,0),(0,1),(0,2),(1,0),(1,1),(2,0),(3,0),(4,0),(5,0)\) \\
\((1^5)\) & \((0,k)\) for \(0\le k\le4\); \((1,k)\) for \(0\le k\le2\);
\((2,k),(3,k)\) for \(0\le k\le1\); \((a,0)\) for \(4\le a\le9\) or
\(a=11\) \\
\end{longtable}

Every other pair is positive. This statement follows from the exhaustive
zeroth-entry check and the positive Newton terms; it includes both axes
and the empty partition. In particular it gives a complete nonvanishing
decision for every outer label in the three-parameter family (SQ.1).

\subsection{6. Exact Grothendieck-framework
transport}\label{exact-grothendieck-framework-transport}

The sheaf in (SQ.2) is a symmetric Schur functor of the tautological
quotient, tensored with an ordinary ideal power of the specified reduced
point scheme. The balanced chart, relative exterior pairing and
determinant line in (SQ.3) identify its global-section representation
with the normalizer module of the Schur functor. Taking the constant
\(S_5\)-covariant summand is the exact projector T-003.

The equality (SQ.6) lives in exact representation \(K_0\), and is proved
by the displayed exact sheaf and cohomology sequences. Its lambda-ring
interpretation is composition of the Schur operations, transported by
the frame torsor as in SGA 3 I 4.5, 4.7.2--4.7.3 and SGA 6 VI 3.3. The
universal lambda-ring locators are SGA 6 V 4.2, 4.4, 4.8 and 4.9. The
projective ideal and its multiplication maps are the explicitly graded
EGA I 9.5.1--9.5.3 construction, and the corresponding flat-family
Hilbert and Quot interpretation is FGA Exposé 221, Theorems 3.1--3.2.
These are the same verified exact locators and associated-bundle maps in
the separate dictionary. The positive lattice formula follows by the
proved finite rational identities, without an additional unconstructed
purity or period map.

\section{The five-point quotient bundle one degree above the quadratic
ray}\label{the-five-point-quotient-bundle-one-degree-above-the-quadratic-ray}

\subsection{T-040. Statement}\label{t-040.-statement}

Let \(k,d\ge0\), and put \[
n=d+k+2,\qquad
\lambda=(d+3k+6,\ d+2k+2,\ d+1,\ d+1,\ d).
\tag{QP.1}
\] These are partitions of \(5n\). For every outer partition
\(\mu\vdash5\), the plethysm coefficient \(p_{\mu,(n)}^\lambda\) has a
positive count by at most 240 fixed tagged rational polytopes, each
using at most five integer variables. Only four tags can contribute
after the residue of \(k\) modulo 60 is fixed. Equations
(QP.19)--(QP.22) below give every constant and inequality. The full
fixed table is
\texttt{\small Q\_\allowbreak{}PLUS\_\allowbreak{}NEWTON\_\allowbreak{}TABLE.json};
a standalone evaluator is
\texttt{\small q\_\allowbreak{}plus\_\allowbreak{}positive.py}.

For the symmetric outer power the exact support criterion is \[
\boxed{
p_{(5),(n)}^\lambda>0\quad\Longleftrightarrow\quad
\begin{cases}
k\ge3,&d\text{ even},\\
k\ge1,&d\text{ odd}.
\end{cases}}
\tag{QP.2}
\] We prove the actual module presentation first and derive the positive
count from its seven finite-group traces. No character calculation or
elimination is required when evaluating the final fixed-table formula.

\subsection{1. The scalar Rees pieces and the precise geometric
map}\label{the-scalar-rees-pieces-and-the-precise-geometric-map}

Retain T-038's five-point ideal \[
I=(xz-xy,xw-xy,yz-xy,yw-xy,zw-xy)
\subset S=\mathbf Q[x,y,z,w],
\] with its literal \(S_5\) action, and put \(E=\mathfrak h_5\),
\(U=I_2=[3,2]\), \(\epsilon=\operatorname{sgn}\). Use \(Q\) for the
tautological rank-three quotient bundle on \(\mathbf P(E)\). Write \[
R=\operatorname{Sym}(U)=\mathbf Q[A,B,C,D,E_0],
\quad T=(A,B,C,D,E_0)^t,
\quad F=AD(E_0-B-C)-BC(E_0-A-D).
\tag{QP.3}
\] Each variable of \(R\) has degree one, corresponding to polynomial
degree two in \(S\). The names \(E\) and \(E_0\) distinguish the
representation from the fifth scalar variable.

Define the graded spaces \[
\mathcal A_k=H^0(\mathcal I_Z^k(2k)),\qquad
\mathcal C_k=H^0(\mathcal I_Z^k(2k+1)),\qquad
\mathcal N_k=H^0(\mathcal I_Z^k(2k+1)\otimes Q),
\tag{QP.4}
\] and take their direct sums over \(k\ge0\). Multiplication by the five
quadrics makes them graded \(R\)-modules.

The primary-source application A-029 establishes the following exact
facts for this ideal: \[
\mathcal A_k=(I^k)_{2k},\qquad
\mathcal C_k=(I^k)_{2k+1},\qquad
H^1(\mathcal I_Z^k(2k))=0.
\tag{QP.5}
\] For \(k\ge2\), these follow from its proved \(2k\)-linear resolution
and saturation bound. For \(k=1\), the ideal is saturated and degree-two
evaluation is onto the five point values. For \(k=0\), these are the
ordinary projective-space statements. The source input is
Kustin--Polini--Ulrich, formula (6.1.4), printed p.~1509, Lemma 6.3 and
Corollary 6.4, p.~1510, with the corollary proof on p.~1511. A-029
checks the actual alternating presentation and the local generator
condition at every prime, and proves the Tor and saturation deductions.
These passages and the complete applied proof were independently read
for this result.

In the coordinates (QP.3), the exact Rees presentation from that same
application is \[
\mathcal R(I)=S[T]/(T^t\varphi,F),
\tag{QP.6}
\] where \(\varphi\) is (5R.14). The map sends each \(T_i\) to its
displayed quadratic generator times the Rees parameter. Its component of
\(S\)-degree zero is \(\mathcal A=R/(F)\). Its component of \(S\)-degree
one is \(\mathcal C\). Thus the relationship of the ordinary powers,
saturated sections and the presentation is through specified maps, with
the degree ranges in (QP.5).

For (QP.1), the flag parameters of G-015 are \[
\lambda_2-n=k,\qquad \lambda_2-\lambda_3=2k+1,
\qquad(\lambda_3-\lambda_5,\lambda_3-\lambda_4)=(1,0).
\] The same relative exterior-column calculation as (5F.6)--(5F.9), now
with \(a=1\), starts with first-line weight \(2k+2\). The determinant
pairing removes one copy of \(\mathcal O(1)\), leaving twist \(2k+1\),
and gives \[
\boxed{(S_\lambda\mathbf Q^5)_{(n^5)}
\simeq\epsilon^{d+1}\otimes\mathcal N_k.}
\tag{QP.7}
\] Explicitly, the relative quotient-flag module is
\(S_{(1,1)}Q^*=\Lambda^2Q^*=Q\otimes\det Q^*\), and
\(\det Q^*=\det E^*\otimes\mathcal O(-1)\). The original full
determinant contributes \(\epsilon^d\), so their product is exactly
\(\epsilon^{d+1}\). The inverse to the flag chart is
\(F_{\rm flag}\mapsto u^{-k}F_{\rm flag}(v(X))\), and G-015 proves its
polynomial extension precisely from the order-\(k\) five-point ideal.
This proves the transport without suppressing its linearization.

\subsection{2. A four-by-five scalar-piece
presentation}\label{a-four-by-five-scalar-piece-presentation}

Write \(T^t\varphi=(x,y,z,w)\mathsf B(T)\). The actual matrix is \[
\mathsf B=\begin{pmatrix}
0&0&D-E_0&E_0-C&C-D\\
E_0-B&A-E_0&0&0&B-A\\
0&D&0&-B&0\\
-C&0&A&0&0
\end{pmatrix}.
\tag{QP.8}
\] It satisfies \[
\mathsf B T=0,\qquad
(-1)^{i-1}\det\mathsf B_{\widehat i}=-FT_i\quad(1\le i\le5).
\tag{QP.9}
\] Here \(\mathsf B_{\widehat i}\) deletes column \(i\). The maximal
minors are nonzero, so the fraction-field kernel has dimension one. The
five entries of \(T\) have gcd one; the UFD denominator argument of
T-038 proves that the entire polynomial kernel is \(RT\).

The matrix \[
\mathsf Q=\begin{pmatrix}
AB&AD&-A(A+C-E_0)&-B(A+D-E_0)\\
AB&BC&-A(B+C-E_0)&-B(B+D-E_0)\\
BC&CD&-C(A+C-E_0)&-D(B+C-E_0)\\
AD&CD&-C(A+D-E_0)&-D(B+D-E_0)\\
AB&CD&-AC&-BD
\end{pmatrix}
\tag{QP.10}
\] obeys \(\mathsf B\mathsf Q=F I_4\), as verified by direct
multiplication. Consequently every degree-one-in-\(S\) consequence of
the relation \(F\) is already in the linear-syzygy image. Equations
(QP.5)--(QP.6) therefore give an exact presentation and resolution \[
\boxed{
0\longrightarrow\epsilon\otimes R(-2)
\mathrel{\mathop{\longrightarrow}^{T}}
(\epsilon\otimes U)\otimes R(-1)
\mathrel{\mathop{\longrightarrow}^{\mathsf B}}
E^*\otimes R\longrightarrow\mathcal C\longrightarrow0.}
\tag{QP.11}
\] The action on the first syzygy space is the one inherited from
T-038's five linear syzygies. The final kernel character is also fixed
by that resolution: after its constant basis change \(P\), its last
column is \(f\), an equivariant map from the sign line. Replacing the
five quadratic coefficients \(f_i\) by the five variables \(T_i\) gives
the column \(T\) in (QP.11), with the same sign source. This explains
both actions directly. The invariant form identifies \(E^*\) with \(E\)
whenever character notation is used below.

\subsection{3. The quotient-bundle exact sequence and a full
resolution}\label{the-quotient-bundle-exact-sequence-and-a-full-resolution}

Tensor the tautological sequence by \(\mathcal I_Z^k(2k+1)\). It remains
exact because its quotient \(Q\) is locally free. By (QP.5), taking
global sections gives the actual exact sequence \[
0\longrightarrow\mathcal A_k
\mathrel{\mathop{\longrightarrow}^{\iota}}
E\otimes\mathcal C_k\longrightarrow\mathcal N_k\longrightarrow0,
\qquad
\iota(s)=\sum_{i=1}^4 e_i\otimes x_i s.
\tag{QP.12}
\] The tensor \(\sum e_i\otimes x_i\) is the identity, independent of
the choice of dual bases. Summing over \(k\) gives an exact sequence of
graded equivariant \(R\)-modules.

There is also the following explicit resolution of the quotient: \[
\boxed{
0\longrightarrow
\bigl(\epsilon E\otimes R(-2)\bigr)\oplus
\bigl(\epsilon\otimes R(-3)\bigr)
\longrightarrow
(\epsilon E\otimes U)\otimes R(-1)
\longrightarrow
(\operatorname{End}E/\mathbf1)\otimes R
\longrightarrow\mathcal N\longrightarrow0.}
\tag{QP.13}
\] To specify the maps in the chosen bases, represent the middle free
module by \(4\times5\) matrices \(Y\). Its map is \[
Y\longmapsto[Y\mathsf B^t]\in
\operatorname{Mat}_{4\times4}(R)/(R I_4),
\] and a lift of the first map is \[
(r,s)\longmapsto rT^t+s\mathsf Q^t,
\qquad r\in R^4, s\in R.
\tag{QP.14}
\] The cokernel is exactly (QP.12): the constant identity generator has
been removed from \(E\otimes E^*\), while the remaining relations are
\(E\) tensored with (QP.11).

Here is a direct check of exactness at the two preceding terms. If
\(Y\mathsf B^t=hI_4\), reduce this identity modulo \(F\). The matrix
\(\mathsf B\) has rank exactly three over the fraction field of
\(R/(F)\): all its maximal minors vanish by (QP.9), whereas its minor
using rows \(2,3,4\) and columns \(1,2,3\) is \(AD(E_0-B)\), not a
multiple of the irreducible cubic \(F\). A scalar matrix of rank four
cannot factor through a matrix of rank three. Therefore \(h=Fs\) for
some \(s\in R\). Subtracting \(s\mathsf Q^t\) leaves a matrix whose
every row lies in \(\ker\mathsf B\), which is \(RT\). Thus
\(Y=rT^t+s\mathsf Q^t\), proving exactness. If this last expression is
zero, multiply by \(\mathsf B^t\): it gives \(sF I_4=0\). Since \(R\) is
a domain, \(s=0\), and then \(rT^t=0\) implies \(r=0\). The first map is
injective.

The submodule \(rT^t\) has character \(\epsilon E\), as in (QP.11). The
quotient of the kernel by that submodule is the line of coefficients
\(h/F\), so its generator has character \(\epsilon\) because \(F\) does.
A graded equivariant choice of the displayed quadratic lift is obtained
by averaging its translates over \(S_5\); the identity
\(\mathsf B\mathsf Q=FI_4\) is preserved. This identifies the
equivariant free modules in (QP.13), even though a particular unaveraged
matrix \(\mathsf Q\) need not itself be an equivariant choice of basis.

The dimension series immediately follows: \[
\sum_{k\ge0}\dim\mathcal N_k t^k
=\frac{15-20t+4t^2+t^3}{(1-t)^5}
=\frac{15-5t-t^2}{(1-t)^4},
\] \[
\boxed{\dim\mathcal N_k
=\frac{(k+1)(3k^2+22k+30)}2.}
\tag{QP.15}
\] In particular the four invariant parameters of degrees \(2,4,6,10\)
used for G-017 are not a regular sequence on this module. Multiplying
(QP.15) by their four factors gives a polynomial with coefficient \(-1\)
in degree 20. A regular sequence would give the Hilbert series of its
actual finite quotient by the successive exact multiplication sequences,
and such a coefficient is impossible. This is a concrete obstruction to
that parameter route for this specific module.

\subsection{4. Seven exact character
series}\label{seven-exact-character-series}

For a group element \(g\), write \(e_g=\chi_E(g)\), \(u_g=\chi_U(g)\),
and \(\epsilon_g=\epsilon(g)\). Taking the graded trace in the actual
resolution (QP.13) gives \[
\sum_{k\ge0}\chi_{\mathcal N_k}(g)t^k
=\frac{e_g^2-1-\epsilon_g e_gu_g t
             +\epsilon_g e_g t^2+\epsilon_g t^3}
        {\det(1-tg\mid U)}.
\tag{QP.16}
\] T-038's signed six-coordinate intertwiner gives every denominator.
After cancellation the seven series and class sizes are as follows.

\begin{longtable}[]{@{}lrl@{}}
\toprule\noalign{}
class & size & series in (QP.16) \\
\midrule\noalign{}
\endhead
\bottomrule\noalign{}
\endlastfoot
\(1^5\) & 1 & \((15-5t-t^2)/(1-t)^4\) \\
\(21^3\) & 10 & \((3-t-t^2)/((1-t)^3(1+t))\) \\
\(2^21\) & 15 & \(-(1+t+t^2)/(1-t^2)^2\) \\
\(31^2\) & 20 & \(t/(1-t^3)\) \\
\(32\) & 20 & \(-t/(1-t^3)\) \\
\(41\) & 30 & \(-(1-t+t^2)/(1-t^4)\) \\
\(5\) & 24 & \(-t^2(1-t)/(1-t^5)\) \\
\end{longtable}

Every denominator divides \((1-t^{60})^4\). The numerator after
multiplication by this common denominator has degree at most 238; this
follows either from the displayed fractions, or from the proper fraction
in (QP.16). Hence the multiplicity of any fixed \([\nu]\) in
\(\mathcal N_{r+60q}\), for fixed \(0\le r<60\), is a polynomial of
degree at most three in \(q\), valid for every \(q\ge0\). Indeed,
grouping the common numerator by residues gives \(h_r(s)/(1-s)^4\),
where \(\deg h_r\le3\). Its coefficients are linear combinations of
\(\binom{q-j+3}{3}\), \(0\le j\le3\). For \(q\ge0\), the values with
\(q<j\) are zero even as polynomial binomial values, so there is no
initial exceptional range.

For completeness all fixed character rows used in taking the inner
product are listed here, in the same class order as the preceding table.

\begin{longtable}[]{@{}ll@{}}
\toprule\noalign{}
\(\nu\) & character row \\
\midrule\noalign{}
\endhead
\bottomrule\noalign{}
\endlastfoot
\((5)\) & \(1,1,1,1,1,1,1\) \\
\((4,1)\) & \(4,2,0,1,-1,0,-1\) \\
\((3,2)\) & \(5,1,1,-1,1,-1,0\) \\
\((3,1,1)\) & \(6,0,-2,0,0,0,1\) \\
\((2,2,1)\) & \(5,-1,1,-1,-1,1,0\) \\
\((2,1,1,1)\) & \(4,-2,0,1,1,0,-1\) \\
\((1^5)\) & \(1,-1,1,1,-1,-1,1\) \\
\end{longtable}

This is the same finite primary character table used and checked in
T-038/T-039.

\subsection{5. Explicit positive Newton
constants}\label{explicit-positive-newton-constants}

Let \(m_\nu(k)\) be the multiplicity of \([\nu]\) in \(\mathcal N_k\).
Write \[
f=\dim[\nu],\quad a=\chi_\nu(21^3),\quad
b=\chi_\nu(2^21).
\] For each \(r\in\{0,\ldots,59\}\), define four constants
\(c_{\nu,r,j}\), \(0\le j\le3\). The first is the following entirely
explicit fixed integer: \[
\begin{aligned}
c_{\nu,r,0}=\frac1{120}\bigl(&fP(r)+10aT(r)+15bD(r)\\
&+20(\chi_\nu(31^2)-\chi_\nu(32))\mathbf1_{r\equiv1\ (3)}\\
&+30\chi_\nu(41)L(r)+24\chi_\nu(5)M(r)\bigr),
\end{aligned}
\tag{QP.17}
\] where \[
\begin{aligned}
P(r)&=(r+1)(3r^2+22r+30)/2,\\
T(2h)&=h^2+5h+3,&T(2h+1)&=(h+1)(h+5),\\
D(2h)&=-(2h+1),&D(2h+1)&=-(h+1),\\
(L(0),L(1),L(2),L(3))&=(-1,1,-1,0)\quad\text{periodically modulo }4,\\
(M(0),\ldots,M(4))&=(0,0,-1,1,0)\quad\text{periodically modulo }5.
\end{aligned}
\tag{QP.18}
\] The other three constants have short closed forms: \[
\boxed{\begin{aligned}
c_{\nu,r,1}
&=\frac{f(9r^2+590r+12352)+a(10r+350)
      -b\bigl(30\mathbf1_{r\text{ even}}+15\mathbf1_{r\text{ odd}}\bigr)}4,\\
c_{\nu,r,2}&=f(270r+16950)+150a,\\
c_{\nu,r,3}&=16200f.
\end{aligned}}
\tag{QP.19}
\] They are the first, second and third forward differences of
\(m_\nu(r+60q)\) at \(q=0\). This follows by applying the difference
operator of step 60 to the coefficients of the seven fractions in
(QP.16). In particular \(P(k)=(3k^3+25k^2+52k+30)/2\),
\(T(k)=k^2/4+5k/2\) plus a parity constant, \(D(k)\) is linear on each
parity, and the last four traces are constant on each residue. Taking
these differences gives (QP.19) directly.

All four constants are nonnegative integers. Integrality follows from
their forward-difference interpretation in actual finite-group modules.
For \(j=0\), the constant is the actual multiplicity \(m_\nu(r)\), and
the 420 constants (QP.17) have also been evaluated in the fixed table.
For \(j\ge1\), positivity follows without a table search from
\(|a|,|b|\le f\). In the first row of (QP.19), the numerator is at least
\(f(9r^2+580r+11972)>0\). In the second it is at least
\(f(270r+16800)>0\), and the third is positive.

Newton's polynomial identity now gives the positive formula \[
\boxed{
m_\nu(r+60q)=\sum_{j=0}^3 c_{\nu,r,j}\binom qj
\qquad(q\ge0),}
\tag{QP.20}
\] where \(\binom qj=0\) for \(j>q\). This identity covers all \(k\),
including \(k=0\); it is not just an asymptotic quasipolynomial.

\subsection{6. The polytopes, coefficient, and
support}\label{the-polytopes-coefficient-and-support}

Put \[
\nu=\begin{cases}\mu,&d+1\text{ even},\\\mu^t,&d+1\text{ odd}.
\end{cases}
\tag{QP.21}
\] For each fixed tag \((r,j)\), \(0\le r<60\), \(0\le j\le3\), take the
polytope in variables \(q,z,x_1,\ldots,x_j\) defined by \[
\boxed{
60q=k-r,\quad q\ge0,\quad
0\le z\le c_{\nu,r,j}-1,\quad
1\le x_1<x_2<\cdots<x_j\le q.}
\tag{QP.22}
\] The final chain is omitted when \(j=0\). Its strict inequalities are
the closed rational inequalities \(x_{i+1}-x_i\ge1\) on integer points.
A zero constant makes the polytope empty. Every coordinate is bounded by
the equality or the displayed bounds, so these are bounded rational
polytopes, each with at most five variables.

There is at most one contributing residue. For it, the choices of the
strictly increasing chain are exactly the \(j\)-subsets of
\(\{1,\ldots,q\}\), and the independent coordinate \(z\) supplies the
factor \(c_{\nu,r,j}\). Consequently the disjoint tagged lattice count
of (QP.22) is \[
\boxed{p_{\mu,(n)}^\lambda=m_\nu(k),}
\tag{QP.23}
\] by (QP.7), the balanced-weight dictionary T-003, and (QP.20). This
specifies a bijection from every tag's lattice witnesses to the positive
summands of the formula. It uses at most 240 fixed tags uniformly, and
at most four after computing \(k\bmod60\).

The fixed constants have at most six decimal digits. The largest entry
is 197280; thus their bit lengths are absolute constants. The remaining
right sides and integer witnesses have \(O(\log(k+2))\) bits. Evaluating
(QP.20) uses four binomial polynomials of degree at most three, so it
requires a constant number of arithmetic operations on integers of
\(O(\log(k+2))\) bits, after the partition and parity checks.

For all seven labels the exact support is particularly simple:

\begin{longtable}[]{@{}ll@{}}
\toprule\noalign{}
\(\nu\) & condition for \(m_\nu(k)>0\) \\
\midrule\noalign{}
\endhead
\bottomrule\noalign{}
\endlastfoot
\((5)\) & \(k\ge1\) \\
\((4,1),(3,2),(3,1,1)\) & all \(k\ge0\) \\
\((2,2,1),(2,1,1,1)\) & \(k\ge1\) \\
\((1^5)\) & \(k\ge3\) \\
\end{longtable}

For \(0\le k<60\), these statements are direct substitution into the
fixed 420 constants (QP.17). The complete values are in the table file,
and the listed zero sets were checked exactly. For \(k\ge60\), the
\(j=1\) term of (QP.20) is positive because \(q\ge1\) and
\(c_{\nu,r,1}>0\). This proves the table for every integer \(k\), and
(QP.2) follows by the sign twist (QP.21).

\subsection{7. Checks, source use, and
scope}\label{checks-source-use-and-scope}

\texttt{\small check\_\allowbreak{}q\_\allowbreak{}plus\_\allowbreak{}family.py}
verifies (QP.9), all five signed minors and their common primitive
kernel. It also checks that each exact class series has the asserted
common denominator and degree bound, computes all 1680 nonnegative
Newton constants, checks their closed formulas (QP.19), and compares all
2520 coefficients through \(k=359\) with the Newton polynomials. These
last comparisons are additional finite checks; the rational-function
degree argument proves the identity for all \(k\). Receipts:
\texttt{\small Q\_\allowbreak{}PLUS\_\allowbreak{}NEWTON\_\allowbreak{}CHECK.tsv}
and \texttt{\small .txt}.

The rectangular agent independently expanded the outer Schur cycle
indices and applied Murnaghan--Nakayama to every output in this family
with \(k,d\ge0\), \(d+k+2\le6\), for all seven outer partitions. All 105
comparisons passed. The proposed traces were computed from the literal
unordered-pair character of \(U\), independently of the Newton table.
Code:
\texttt{\small work/\allowbreak{}rectangular/\allowbreak{}verify\_\allowbreak{}outer5\_\allowbreak{}odd\_\allowbreak{}q\_\allowbreak{}family.py};
receipt:
\texttt{\small work/\allowbreak{}rectangular/\allowbreak{}OUTER5\_\allowbreak{}ODD\_\allowbreak{}Q\_\allowbreak{}INDEPENDENT\_\allowbreak{}CHECK.tsv}.

The source-derived facts (QP.5)--(QP.6) are the precisely applied
results in A-029,
\texttt{\small work/\allowbreak{}archive/\allowbreak{}FIVE\_\allowbreak{}POINT\_\allowbreak{}POWERS.md}.
The primary source is
\href{https://msp.org/ant/2017/11-7/ant-v11-n7-p01-p.pdf}{Kustin--Polini--Ulrich,
\emph{The equations defining blowup algebras of height three Gorenstein
ideals}}, \emph{Algebra \& Number Theory} 11 (2017), 1489--1525, DOI
10.2140/ant.2017.11.1489; the exact locators are given in §1. The
explicit five-point matrices and sign labels are T-038, and the relative
Schur/wedge maps in (QP.7) are spelled out here and in T-039.

The presentation (QP.13) and the positive count (QP.22) evaluate this
particular higher-flag family in full. The parameter obstruction after
(QP.15) records exactly why G-017's four-parameter free-module proof
does not extend to this module with those parameters. The new count uses
its explicit resolution and fixed quasipolynomial coefficients; it does
not assume a new freeness or cohomology-vanishing assertion.

\section{A five-point hook-bundle
family}\label{a-five-point-hook-bundle-family}

\subsection{T-041. Statement and exact flag
comparison}\label{t-041.-statement-and-exact-flag-comparison}

For \(a\ge1\), \(d\ge0\), put \[
n=d+a+1,\qquad
\lambda=(d+2a+4,d+a+2,d+a,d+a-1,d).
\tag{5H.1}
\] Let \(E=\mathfrak h_5\), \(U=[3,2]\), and
\(\epsilon=\operatorname{sgn}\). There is an actual equivariant
isomorphism \[
\boxed{(S_\lambda\mathbf Q^5)_{(n^5)}
\simeq\epsilon^{d+a}\otimes U\otimes S_{(a,1)}E.}
\tag{5H.2}
\] Every outer \(\mu\vdash5\) has the positive formula below with at
most 90 fixed polytopes of dimension at most four. For symmetric outer
power, \[
\boxed{p_{(5),(n)}^\lambda>0\quad\Longleftrightarrow\quad a\ge2.}
\tag{5H.3}
\]

Indeed \(\lambda_2-n=1\), \(\lambda_2-\lambda_3=2\), and
\((\lambda_3-\lambda_5,\lambda_3-\lambda_4)=(a,1)\). The actual chart
and extension map of G-015 therefore identifies the left side of (5H.2)
with \[
\epsilon^{d+a}\otimes
H^0(\mathbf P(E),\mathcal I_Z(2)\otimes S_{(a,1)}Q),
\tag{5H.4}
\] where \(Q=E/\mathcal O(-1)\) and \(Z\) is the five canonical frame
points. To check every determinant, the relative quotient-flag weight
after removing the full determinant shift is \((a,a-1,0)\).
Complementary exterior-column pairings in the rank-three quotient give
\[
S_{(a,a-1)}Q^*\simeq(\det Q^*)^a\otimes S_{(a,1)}Q.
\] This is the restriction of the perfect wedge pairings of
complementary columns in the height-three, width-\(a\) rectangle to
their Schur summands. The first-line twist is \(a+2\); the identity
\(\det Q^*=\det E^*\otimes\mathcal O(-1)\) leaves \(\mathcal O(2)\). The
full determinant contributes \(\epsilon^d\), so the total character is
\(\epsilon^{d+a}\), as claimed. The inverse flag chart is
\(F\mapsto u^{-1}F(v(X))\); G-015 proves that (5H.4)'s ideal condition
is exactly its polynomial extension condition. Thus no determinant or
extension condition is suppressed.

\subsection{1. An exact hook resolution}\label{an-exact-hook-resolution}

For a vector space \(V\), write \[
H_a(V)=\ker(V\otimes\operatorname{Sym}^aV
                    \longrightarrow\operatorname{Sym}^{a+1}V).
\tag{5H.5}
\] Multiplication has its equivariant polarization section divided by
\(a+1\). The kernel is \(S_{(a,1)}V\): it contains the highest vector
\(v_1\otimes v_1^{a-1}v_2-v_2\otimes v_1^a\), and the irreducible
highest-weight summand it generates has the kernel's dimension by the
Weyl dimension formula. This specifies the hook functor used here.

The line inclusion \(L=\mathcal O(-1)\subset E\otimes\mathcal O\) gives
the exact locally free resolution \[
0\to L^2\otimes\operatorname{Sym}^{a-1}E
\mathrel{\mathop{\longrightarrow}^{d_2}}
L\otimes E\otimes\operatorname{Sym}^{a-1}E
\mathrel{\mathop{\longrightarrow}^{d_1}}
H_a(E)\otimes\mathcal O\to H_a(Q)\to0,
\tag{5H.6}
\] with actual maps \[
d_2(l^2\otimes v)=l\otimes l\otimes v,\qquad
d_1(l\otimes e\otimes v)=e\otimes(lv)-l\otimes(ev).
\tag{5H.7}
\] Their composition is zero. For exactness choose a local splitting
\(E=L\oplus W\). In \(d_1\), the first tensor factors in \(W\) record
\(e\otimes lv\). Multiplication by \(l\) is injective, so \(\ker d_1\)
is exactly the subspace with \(e\in L\), the image of \(d_2\). The map
\(H_a(E)\to H_a(Q)\) is onto by the splitting. An element of its kernel
has the form \(\sum e_i\otimes lv_i+l\otimes w\), with \(e_i\in W\). The
zero-product condition says \(w=-\sum e_iv_i\), so this element is
\(d_1(l\otimes\sum e_i\otimes v_i)\). This proves exactness at every
term of (5H.6).

\subsection{2. All cohomology used in the
formula}\label{all-cohomology-used-in-the-formula}

Put \(B_a=H_a(Q)\). The scalar projective-space calculation of Stacks
Project, Lemma 30.8.1, applied to (5H.6) gives \[
H^0(B_a)=H_a(E),\quad H^{i>0}(B_a)=0,
\qquad H^i(B_a(-1))=0\quad\text{for all }i.
\tag{5H.8}
\] The negative twists involved are respectively \(-2,-1\) and
\(-3,-2,-1\), all acyclic on \(\mathbf P^3\).

For \(B_a(-3)\), the three terms of (5H.6) have twists \(-5,-4,-3\). The
only possible cohomology is the map \[
H^3(\mathcal O(-5))\otimes\operatorname{Sym}^{a-1}E
\longrightarrow H^3(\mathcal O(-4))\otimes E
                                      \otimes\operatorname{Sym}^{a-1}E.
\tag{5H.9}
\] It is the identity under the natural identifications. The ordered
Čech generator indexed by \(i\) in the source is
\(1/(x_1x_2x_3x_4x_i)\). Multiplication by the \(j\)-th coordinate of
the line inclusion gives the top class \(1/(x_1x_2x_3x_4)\) for \(j=i\);
for \(j\ne i\), one exponent is zero, so the top class vanishes. Thus
the matrix is exactly \(\delta_{ij}\). Both spaces carry the same factor
\(\det E=H^3(\mathcal O(-4))\), so this is also the correct equivariant
map. The two short exact sequences of (5H.6) consequently prove \[
\boxed{H^i(B_a(-3))=0\quad\text{for all }i,\ a\ge1.}
\tag{5H.10}
\] This includes the boundary \(a=1\), where the hook is \(\Lambda^2Q\).

Now tensor T-038's actual five-point resolution by \(B_a(2)\): \[
0\to\epsilon B_a(-3)\to\epsilon U\otimes B_a(-1)
\to U\otimes B_a\to\mathcal I_Z(2)\otimes B_a\to0.
\tag{5H.11}
\] Tensoring is exact because \(B_a\) is locally free. Both left terms
are acyclic by (5H.8) and (5H.10), hence multiplication gives the actual
isomorphism \[
U\otimes H_a(E)\longrightarrow
H^0(\mathcal I_Z(2)\otimes B_a).
\tag{5H.12}
\] It multiplies the five specified quadrics by constant hook sections
and then applies the quotient \(E\to Q\). Together with (5H.4), this
proves (5H.2). Its dimension is \(5a(a+2)(a+3)/2\).

\subsection{3. The complete fixed positive
formula}\label{the-complete-fixed-positive-formula}

For positive weights \(W\), let \[
C_W(R)=\#\{x\in\mathbf Z_{\ge0}^{|W|}:\sum_{w\in W}wx_w=R\}.
\tag{5H.13}
\] This is a bounded rational-polytope count, zero for \(R<0\). Use
\(C_4=C_{(2,3,4,5)}\) and \(C_3=C_{(2,3,5)}\). The multiplicity
\(L_\rho(a)\) of \([\rho]\) in \(H_a(E)\) is
\(\sum_{r\in A_\rho}C_4(a-r)+\sum_{r\in B_\rho}C_3(a-r)\), with all
shifts listed below. Repetitions are separate tags.

\begin{longtable}[]{@{}lll@{}}
\toprule\noalign{}
\(\rho\) & \(A_\rho\) & \(B_\rho\) \\
\midrule\noalign{}
\endhead
\bottomrule\noalign{}
\endlastfoot
\((5)\) & \(4,6,13\) & \(5,6,7\) \\
\((4,1)\) & \(2,3,3,4,4,4,5,5,5,6,6,7\) & empty \\
\((3,2)\) & \(2,3,3,4,4,4,5,5,5,6,6,6,7,7,8\) & empty \\
\((3,1,1)\) & \(1,2,3,3,4,4,5,5,5,6,6,6,7,7,7,8,8,9\) & empty \\
\((2,2,1)\) & \(2,3,4,4,5,5,5,6,6,6,7,7,8,8,9\) & empty \\
\((2,1,1,1)\) & \(3,4,4,5,5,6,6,7,7,8,9,10\) & empty \\
\((1^5)\) & \(6,7,8\) & empty \\
\end{longtable}

To prove the table, the polarization splitting in (5H.5) gives the
actual decomposition
\(E\otimes\operatorname{Sym}^aE=H_a(E)\oplus\operatorname{Sym}^{a+1}E\).
Writing \(H_g(t)=\det(1-tg\mid E)^{-1}\), one obtains \[
\sum_{a\ge0}\chi_{H_a(E)}(g)t^a
=(\chi_E(g)-t^{-1})H_g(t)+t^{-1},\qquad H_0(E)=0.
\tag{5H.14}
\] Take inner products using the seven fixed character rows of QP.16's
following table. T-027's tableau standardization gives the common
denominator \(D(t)=(1-t^2)(1-t^3)(1-t^4)(1-t^5)\). For each nontrivial
row, its numerator has exactly the nonnegative coefficients recorded as
\(A_\rho\). For the trivial row the exact polynomial identity is \[
\begin{aligned}
&t^4+t^5+2t^6+t^7-t^9-t^{10}-t^{11}+t^{13}\\
&\hspace{8mm}=t^4+t^6+t^{13}+(1-t^4)(t^5+t^6+t^7).
\end{aligned}
\tag{5H.15}
\] After division by \(D(t)\), this is precisely its three four-variable
and three three-variable counts. Seven exact rational identities verify
the full table in every degree, including its zero initial terms.

For \(\mu\vdash5\), put \(\nu=\mu\) when \(d+a\) is even and
\(\nu=\mu^t\) when \(d+a\) is odd. Let \(K_{\rho\nu}\) be the completely
listed fixed \([3,2]\)-tensor table in §4 of T-039. Then \[
\boxed{p_{\mu,(n)}^\lambda=\sum_{\rho\vdash5}K_{\rho\nu}L_\rho(a).}
\tag{5H.16}
\] Each summand is tagged by a tensor copy, a table row and a shift. For
\(\rho\ne(5)\) there are exactly \(3\dim[\rho]\) tags; the trivial row
has three additional tags. Self-duality and the tensor dimension
identity give the total \[
15\dim[\nu]+3\mathbf1_{\nu=(3,2)}\le90.
\tag{5H.17}
\] Each tag has one displayed equality and nonnegative coordinates with
weights \((2,3,4,5)\) or \((2,3,5)\). Every witness has \(O(\log(a+2))\)
bits. All tables have fixed size.

For symmetric outer power, only row \((3,2)\) contributes in even parity
and only row \((2,2,1)\) in odd parity. Both have shifts two and three
and no smaller shift. Since their positive weights generate
\(\{0\}\cup\{r:r\ge2\}\), they are positive exactly for \(a\ge2\), which
proves (5H.3). The sign outer partition has the same threshold. Every
other outer partition receives row \((3,1,1)\), with shifts one and two,
and is positive for every \(a\ge1\).

\subsection{Checks and primary
sources}\label{checks-and-primary-sources}

\texttt{\small five\_\allowbreak{}point\_\allowbreak{}hook\_\allowbreak{}positive.py}
contains the entire fixed positive evaluator.
\texttt{\small check\_\allowbreak{}five\_\allowbreak{}point\_\allowbreak{}hook.py}
verifies all seven exact rational identities, the tag bounds, and 560
finite coefficient/support cases. Receipts:
\texttt{\small FIVE\_\allowbreak{}POINT\_\allowbreak{}HOOK\_\allowbreak{}CHECK.tsv}
and \texttt{\small .txt}. The semigroup argument proves support in every
degree; the finite checks supplement it.

The scalar cohomology and explicit Čech classes are those of
\href{https://stacks.math.columbia.edu/tag/01XS}{Stacks Project, Lemma
30.8.1}, whose full calculation was read. The relative Schur and
exterior-column convention uses
\href{https://arxiv.org/pdf/alg-geom/9411014}{Magyar, Proposition 1 and
Proposition 4}, as specified in T-039. The hook resolution and all its
maps are proved here. The finite-group table was checked against
\href{https://www.math.uchicago.edu/~farb/papers/wiman.pdf}{Dolgachev--Farb--Looijenga,
§2.3, Table 1}.

At \(a=1\) this meets the exterior-square quotient-bundle family at
point power one. For \(a>1\), it evaluates the actual rank-three hook
section space in (5H.4). No vanishing for arbitrary Schur quotient
bundles or higher point powers is assumed.

\part{Checked canonical-map and rectangular results}

\section{Rectangular cases and the normalized Hadamard--Howe
period}\label{rectangular-cases-and-the-normalized-hadamardhowe-period}

Run:
\texttt{\small /\allowbreak{}root/\allowbreak{}rectangular\_\allowbreak{}cases},
21 November 2026 (Europe/Berlin). Model: GPT-6 Astra Pro. The inherited
effort setting is not exposed to this agent.

This report supplies exact reformulations, proofs, source locators, and
finite certificates for the parent paper. The literature results remain
attributed to their authors. The labels R-001--R-006 identify findings
and explicit dictionaries established in this run; they do not claim
priority over the cited classical results.

\subsection{Notation}\label{notation-1}

Write

\[
a_\lambda(d[n])=[S_\lambda V:\operatorname{Sym}^d(\operatorname{Sym}^nV)],
\qquad R_{r,s}=(s^r),
\]

with \(\dim V\ge\ell(\lambda)\). Thus \(r\times s\) in the supplied
Ikenmeyer--Panova source means \(R_{r,s}\). All decompositions of
complex representations below use characteristic zero. The base ring
used for normalized integral tensors will be specified separately.

The Hadamard--Howe map is normalized by averaging permutations, so that

\[
h_{d,n}(v_1^n\cdots v_d^n)=(v_1\cdots v_d)^n.
\]

This is the convention of Kumar--Landsberg, §1.2, arXiv:1410.8585v1.

\subsection{R-001. A rectangular plethysm coefficient as a one-point
lattice
count}\label{r-001.-a-rectangular-plethysm-coefficient-as-a-one-point-lattice-count}

For every \(r\ge2\), \(n\ge1\),

\[
a_{(n^r)}(r[n])
=\#\{z\in\mathbb Z:2z=n\}
=\begin{cases}1&n\text{ even},\\0&n\text{ odd}.\end{cases}
\tag{R-001}
\]

The first equality presents a positive lattice-point formula on the
rational zero-dimensional polytope \(\{n/2\}\). The multiplicity-one
statement is classical: Howe,
\emph{\((\mathrm{GL}_n,\mathrm{GL}_m)\)-duality and symmetric plethysm},
Proposition 4.3, as located explicitly in Kumar--Landsberg §2; also
supplied BIP §2(b).

\subsubsection{Proof of the precise lattice-count
form}\label{proof-of-the-precise-lattice-count-form}

Take \(V=\mathbb C^r\). A tensor in
\((\operatorname{Sym}^nV)^{\otimes r}\) is a polynomial on \(r\)
covectors, separately of degree \(n\) in each covector. Arrange the
covectors as columns of an \(r\times r\) matrix \(X\). On the dense open
subset \(\det X\ne0\), an \(\mathrm{SL}_r\)-invariant polynomial depends
only on \(\det X\): two invertible matrices with the same determinant
differ by left multiplication by an element of \(\mathrm{SL}_r\).
Separate homogeneity forces the invariant to be a scalar multiple of
\((\det X)^n\). Equality on the dense open subset is equality of
polynomials. Thus the invariant space in the ordered tensor product is
one-dimensional.

Permutation of the \(r\) covectors acts on this line by
\(\operatorname{sgn}^n\). Passing to the outer symmetric power is taking
the invariants of this \(S_r\)-action. Since \(r\ge2\), the invariant
dimension is one for even \(n\) and zero for odd \(n\). The
\(\mathrm{SL}_r\)-invariant summand of homogeneous degree \(rn\) has
\(\mathrm{GL}_r\)-highest weight \((n^r)\). Finally, the equation
\(2z=n\) has exactly one integral solution for even \(n\) and none for
odd \(n\). This proves all displayed equalities.

The same coefficient holds in every \(\dim V\ge r\): Schur-functor
multiplicities are stable once the indexed partition fits, by the
symmetric-group multiplicity-space realization.

\subsection{R-002. The exact square Hadamard--Howe
scalar}\label{r-002.-the-exact-square-hadamardhowe-scalar}

Fix an even integer \(n\ge2\), let \(V=\mathbb C^n\), and put
\(H=S_n\wr S_n\). Index the tensor factors of \(V^{\otimes n^2}\) by the
cells \((i,j)\) of an \(n\times n\) array. Let \(e_H\) be the average
over permutations within rows and permutations of rows. Let \(\tau\)
transpose the array. Then

\[
\operatorname{Sym}^n(\operatorname{Sym}^nV)=e_HV^{\otimes n^2},
\qquad h_{n,n}=e_H\tau e_H.
\]

Define

\[
\delta=\sum_{\sigma\in S_n}\operatorname{sgn}(\sigma)
e_{\sigma(1)}\otimes\cdots\otimes e_{\sigma(n)},
\]

put one copy of \(\delta\) in each column, call the resulting tensor
\(D\), and put \(I=e_HD\). This \(I\) spans the determinant line from
R-001. For a Latin square \(L\), let \(c(L)\) denote the product of its
column signs, and let \(E_n-O_n\) be its signed enumeration using the
product of all row and column signs. Set

\[
\Delta_n=\sum_{L\text{ Latin}}c(L),\qquad
B_n=\prod_{j=0}^{n-1}\frac{(n+j)!}{j!}.
\]

The normalized scalar is

\[
\boxed{
h_{n,n}(I)=c_n I,\qquad
c_n=\frac{\Delta_n}{(n!)^n}
=\frac{(-1)^{n(n-1)/2}(E_n-O_n)}{(n!)^n}
=\frac{B_n}{(n!)^n}\int_{\mathrm{SU}(n)}\prod_{i,j=1}^n u_{ij}\,du.
}
\tag{R-002}
\]

Here Haar measure has total mass one. Kumar--Landsberg Theorem 1.9
proves the associated nonvanishing equivalences; Proposition 2.1 gives
the determinant contraction. The present calculation fixes every scalar
under the displayed averaging convention. Alpöge, \emph{Square-root
cancellation for the signs of Latin squares}, Corollary 5, supplies the
last coefficient-to-Haar identity. That paper is separate bibliographic
material from the supplied
\texttt{\small ALPOGE\_\allowbreak{}remark.pdf}.

\subsubsection{Proof of the scalar}\label{proof-of-the-scalar}

Equip the standard tensor basis with its orthonormal inner product. Let
\(t_0\) have entry \(e_i\) in every cell of row \(i\). Although \(t_0\)
is invariant under the within-row permutations, it need not be invariant
under permutations of rows. The needed pairing still has the value one:
every row permutation of \(t_0\) has coefficient
\(\operatorname{sgn}(\pi)^n=1\) in \(D\), so

\[
\langle I,t_0\rangle=\langle D,e_Ht_0\rangle=1.
\tag{R-002a}
\]

Because \(I\) spans the determinant line, write \(h_{n,n}I=c_nI\).
Self-adjointness of the averages and the transpose permutation gives

\[
c_n=\langle e_H\tau I,t_0\rangle
=\langle I,\tau e_Ht_0\rangle.
\]

The tensor \(e_Ht_0\) is the average of \(t_0\) over permutations of
rows. After applying \(\tau\), those permutations become simultaneous
within-row permutations of \(t_1:=\tau t_0\). Since \(I\) is invariant
under the within-row permutations,

\[
c_n=\langle I,t_1\rangle=\langle D,e_Ht_1\rangle.
\]

Every row of \(t_1\) is \((e_1,\ldots,e_n)\). Thus \(e_Ht_1\) averages
independently over the \(n!\) reorderings in each of its \(n\) rows. A
tensor in this average pairs nontrivially with \(D\) exactly when its
columns also contain each basis vector once, which is exactly the
Latin-square condition. Its coefficient in \(D\) is the product of
column signs. Therefore the pairing is \(\Delta_n/(n!)^n\).

\subsubsection{Proof of the parity
conversion}\label{proof-of-the-parity-conversion}

This conversion is an identity of total signed sums, not an equality of
the two sign functions on each square. Let \(r(L),c(L),s(L)\) be
respectively the products of row, column, and symbol-position signs. On
the \(n^2\) ordered pairs, define

\[
R:(i,j)\mapsto(i,L_{ij}),\quad
S:(i,s)\mapsto(j,s)\text{ where }L_{ij}=s,
\quad T:(j,s)\mapsto(s,j).
\]

The column permutation \(C:(i,j)\mapsto(L_{ij},j)\) satisfies \(C=TSR\).
The transpose permutation \(T\) has sign \(\epsilon=(-1)^{n(n-1)/2}\),
hence \(r(L)c(L)s(L)=\epsilon\). Form the conjugate Latin square \(L'\),
with row \(i\), column \(s\), and entry \(j\) specified by \(L_{ij}=s\).
Its row-sign product is \(r(L)\), and its column-sign product is
\(s(L)\). Its ordinary sign is therefore \(\epsilon c(L)\). The
correspondence \(L\mapsto L'\) is a bijection. Summing proves
\(\Delta_n=\epsilon(E_n-O_n)\).

\subsubsection{Proof of the Haar
normalization}\label{proof-of-the-haar-normalization}

The coefficient \(C_n\) of \(\prod_{i,j}X_{ij}\) in \(\det(X)^n\) is the
sum of the products of the signs of the symbol-position permutations in
Latin squares. Parastrophy exchanges this with the column-sign sum, so
\(C_n=\Delta_n\). Apply Alpöge Corollary 5 to the all-ones multiindex
\(\alpha\), with \(k=n\). Its factorial \(\alpha!\) is one, and the
displayed factor in that corollary is \(B_n\). Consequently
\(\Delta_n=B_n\int\prod u_{ij}\,du\), giving the last equality.

\subsubsection{Exact finite values and a checked
bound}\label{exact-finite-values-and-a-checked-bound}

The enumeration script
\texttt{\small verify\_\allowbreak{}rectangular.py} visits every reduced
Latin square for \(n=2,4,6\), with first row and first column fixed. It
found:

\begin{longtable}[]{@{}
  >{\raggedleft\arraybackslash}p{(\columnwidth - 10\tabcolsep) * \real{0.1667}}
  >{\raggedleft\arraybackslash}p{(\columnwidth - 10\tabcolsep) * \real{0.1667}}
  >{\raggedleft\arraybackslash}p{(\columnwidth - 10\tabcolsep) * \real{0.1667}}
  >{\raggedleft\arraybackslash}p{(\columnwidth - 10\tabcolsep) * \real{0.1667}}
  >{\raggedleft\arraybackslash}p{(\columnwidth - 10\tabcolsep) * \real{0.1667}}
  >{\raggedleft\arraybackslash}p{(\columnwidth - 10\tabcolsep) * \real{0.1667}}@{}}
\toprule\noalign{}
\begin{minipage}[b]{\linewidth}\raggedleft
\(n\)
\end{minipage} & \begin{minipage}[b]{\linewidth}\raggedleft
Reduced squares
\end{minipage} & \begin{minipage}[b]{\linewidth}\raggedleft
Reduced column-sign sum
\end{minipage} & \begin{minipage}[b]{\linewidth}\raggedleft
All squares
\end{minipage} & \begin{minipage}[b]{\linewidth}\raggedleft
\(\Delta_n\)
\end{minipage} & \begin{minipage}[b]{\linewidth}\raggedleft
\(c_n\)
\end{minipage} \\
\midrule\noalign{}
\endhead
\bottomrule\noalign{}
\endlastfoot
2 & 1 & -1 & 2 & -2 & \(-1/2\) \\
4 & 4 & 4 & 576 & 576 & \(1/576\) \\
6 & 9408 & -2304 & 812851200 & -199065600 & \(-1/699840000\) \\
\end{longtable}

Every reduced square has \(n!(n-1)!\) completions under row and symbol
normalization, and the column sign is unchanged for even \(n\). The
recursion enumerates all row permutations with their first entry
prescribed, rejects exactly repeated column symbols, and at depth \(n\)
evaluates the product of column permutation signs. Thus this is a finite
exhaustive certificate, not sampling. The certificate asserts only these
three orders.

Alpöge Theorem 2, evaluated at \(k=n\), gives the exact inequality

\[
0\le |c_n|^2\le\frac{B_n}{(n!)^{2n}}.
\tag{R-002b}
\]

The value zero belongs to this interval. The unresolved scalar in the
square rectangular component is precisely \(\Delta_n\), or equivalently
the displayed Haar integral. No lower bound for general even \(n\) is
established in this report.

\subsection{R-003. A concrete failure of positivity transfer from
rectangular Kronecker
coefficients}\label{r-003.-a-concrete-failure-of-positivity-transfer-from-rectangular-kronecker-coefficients}

For every \(d,n\ge1\) and \(1\le k\le dn-1\),

\[
a_{(dn-k,1^k)}(d[n])=0.
\tag{R-003a}
\]

This vanishing is known: Eugenio Giannelli, \emph{Representations of
symmetric groups}, PhD thesis (2015), Theorem 2.1.2, p.38, with proof in
§2.3, pp.42--43; Giannelli attributes the earlier Schur-function proof
to Langley--Remmel. The short argument below records the exact
wreath-product/exterior-power dictionary used in the counterexample. The
hook/exterior-power identification is Giannelli Lemma 2.3.3 (also
Fulton--Harris, Proposition 3.12 for irreducibility).

In particular,

\[
\boxed{a_{(13,1,1,1)}(4[4])=0,\qquad
g((13,1,1,1),(4^4),(4^4))=1.}
\tag{R-003b}
\]

The parameters in R-003b lie in the stable range of Ikenmeyer--Panova
Proposition 2.8: the tail \(\rho=(1,1,1)\) has size three and the inner
degree is four. Thus the actual inequality at this point is \(1\ge0\).

\subsubsection{Proof}\label{proof-1}

The finite-group multiplicity formula is

\[
a_\lambda(d[n])=\dim[\lambda]^{S_n\wr S_d}.
\]

For \(N=dn\), the hook Specht module \([N-k,1^k]\) is
\(\bigwedge^kU_N\), where \(U_N\) is the standard, sum-zero
representation of \(S_N\). Restrict to the base group \(K=S_n^d\). There
is an explicit decomposition

\[
U_N=A\oplus U_n^{(1)}\oplus\cdots\oplus U_n^{(d)},
\]

where \(A\) consists of vectors constant within each block and with
total sum zero. The base group fixes \(A\) pointwise, and its quotient
\(S_d\) acts on \(A\) as \(U_d\). For every \(j=1,\ldots,n-1\),
\(\bigwedge^jU_n\) is the nontrivial irreducible hook \([n-j,1^j]\), so
it has no \(S_n\)-invariants. The exterior-power decomposition of a
direct sum therefore gives

\[
(\bigwedge^kU_N)^K=\bigwedge^kA.
\]

Taking the remaining \(S_d\)-invariants yields zero for every \(k>0\):
the exterior powers in their nonzero range are the nontrivial hooks
\([d-k,1^k]\), and the remaining exterior powers vanish. The cases
\(n=1\) or \(d=1\) are immediate from the same decomposition with
zero-dimensional summands. This proves R-003a.

Ikenmeyer--Panova Proposition 6.4 gives

\[
\sum_k g((16-k,1^k),(4^4),(4^4))q^k
=(1+q^3)(1+q^5)(1+q^7).
\]

The coefficient of \(q^3\) is one, proving R-003b. Independently,
\texttt{\small verify\_\allowbreak{}rectangular.py} computed the
plethysm by the power-sum cycle index \(h_4[h_4]\), paired with the hook
character, and the Kronecker coefficient by the Murnaghan--Nakayama
character rule followed by the exact conjugacy-class sum. Both
computations use rational arithmetic and give the displayed values.

\subsection{R-004. The rectangular-hook
polytope}\label{r-004.-the-rectangular-hook-polytope}

For \(1\le r\le s\), \(0\le k\le rs-1\), let

\[
P_{r,k}=\{(x_2,\ldots,x_r)\in[0,1]^{r-1}:
\sum_{i=2}^r(2i-1)x_i=k\}.
\]

Ikenmeyer--Panova Proposition 6.4 gives exactly

\[
g((rs-k,1^k),(s^r),(s^r))=\#(P_{r,k}\cap\mathbb Z^{r-1}).
\tag{R-004}
\]

The proof of this reformulation is the bijection sending a partition
into distinct allowed parts \(3,5,\ldots,2r-1\) to its indicator vector.
It preserves the sum. R-003 explicitly computes the plethysm coefficient
for one of the positive points of this polytope.

\subsection{R-005. An explicit vector and multiplicity one for
row-extended even
rectangles}\label{r-005.-an-explicit-vector-and-multiplicity-one-for-row-extended-even-rectangles}

Take \(k\ge2\), an even integer \(\ell\ge2\), and \(n\ge\ell\). Put

\[
\lambda=(kn-(k-1)\ell,\ell^{k-1}).
\]

Then

\[
a_\lambda(k[n])=1.
\tag{R-005a}
\]

For \(n\ge k\ell\), the same multiplicity in
\(\mathbb C[\overline{\mathrm{GL}_{n^2}\det_n}]_k\) is exactly one. This
combines the explicit vector below with BIP Propositions 2.3 and 5.6(1),
and Theorem 2.5. The equality of multiplicities refines the occurrence
statement to the one-dimensional ambient component.

\subsubsection{Proof and normalized nonzero
evaluation}\label{proof-and-normalized-nonzero-evaluation}

Use \(k\) basis vectors and form \(I_{k,\ell}\) by placing the
alternating \(k\)-tensor in each of \(\ell\) columns and averaging the
within-row and outer permutations. It is a highest weight vector of
weight \((\ell^k)\). For nonzero coefficients \(c_1,\ldots,c_k\), direct
contraction gives

\[
\big\langle I_{k,\ell},(\sum_{i=1}^k c_iX_i^\ell)^{\otimes k}\big\rangle
=k!\prod_{i=1}^k c_i.
\tag{R-005b}
\]

Indeed, a nonzero determinant contraction assigns distinct basis labels
to the \(k\) rows. There are \(k!\) assignments, their signs are raised
to the even power \(\ell\), and every assignment has coefficient
\(\prod_i c_i\).

Let \(M:\operatorname{Sym}^\ell V\to\operatorname{Sym}^nV\) multiply by
\(e_1^{n-\ell}\), and put \(\kappa=\operatorname{Sym}^kM\). The upper
triangular unipotent group fixes \(e_1\), so \(\kappa(I_{k,\ell})\)
remains a highest weight vector and its weight is \(\lambda\). The map
\(M\), and therefore \(\kappa\), is injective. BIP Proposition 5.6(1)
says that \(\kappa\) is an isomorphism on these highest weight spaces,
since the second part of \((\ell^k)\) equals the initial inner degree
\(\ell\). R-001 now gives R-005a.

With normalized symmetric-tensor pairings, the dual map is

\[
M^*=\frac{\ell!}{n!}\partial_{X_1}^{n-\ell}.
\]

For

\[
q=X_1^{n-\ell}\sum_{i=1}^k c_iX_i^\ell
\]

one obtains

\[
M^*q=c_1X_1^\ell+\binom n\ell^{-1}\sum_{i=2}^k c_iX_i^\ell,
\]

and hence

\[
\boxed{\langle\kappa(I_{k,\ell}),q^{\otimes k}\rangle
=k!\binom n\ell^{-(k-1)}\prod_{i=1}^k c_i\ne0.}
\tag{R-005c}
\]

The pairing identity follows directly from the definition of a dual map,
also BIP Theorem 5.4. For \(n\ge k\ell\), BIP Theorem 2.5 places every
such padded power sum in the determinant orbit closure. Thus the
coordinate-ring quotient is nonzero on this one-dimensional highest
weight space. Its multiplicity is both at least one and at most the
ambient multiplicity one, proving the claimed equality.

\subsection{R-006. Scheme, locally free sheaf, and Grothendieck-group
dictionary}\label{r-006.-scheme-locally-free-sheaf-and-grothendieck-group-dictionary}

The following constructions specify the morphisms transporting
R-001--R-005.

\begin{enumerate}
\def\labelenumi{\arabic{enumi}.}
\item
  \textbf{Group schemes and Schur operations.} The representations are
  obtained from the standard locally free sheaf of
  \(\mathrm{GL}_{r,\mathbb Z}\) by tensor, exterior, and symmetric
  powers. Over \(\mathbb Q\), their identities are identities in
  \(K_0(\operatorname{Rep}_{\mathbb Q}\mathrm{GL}_r)\). A principal
  \(\mathrm{GL}_r\)-bundle on a \(\mathbb Q\)-scheme gives an associated
  locally free sheaf \(E\). Applying the associated-bundle functor
  transports each actual natural tensor morphism and each equality in
  this Grothendieck group to the corresponding sheaf morphism and class
  in \(K_0(X)\). In particular, the plethystic operation is
  \([E]\mapsto[\operatorname{Sym}^d(\operatorname{Sym}^nE)]\).
\item
  \textbf{The determinant line and the normalized map.} For R-002, take
  \(R=\mathbb Z[1/n!]\). The group order \(|H|=(n!)^{n+1}\) is
  invertible, so all displayed averages are defined over \(R\). Frame
  change multiplies \(I\) by \(\det(g)^n\). The resulting locally free
  line is \((\det E)^{\otimes n}\), and the restriction of the natural
  tensor map to it is multiplication by the exact scalar
  \(\Delta_n/(n!)^n\). The scalar and the line therefore descend from
  frames over \(R\) to every locally free \(E\) of rank \(n\). Over the
  complex fiber, the coordinate multiplicity in the image of this map is
  the rank of this scalar map.
\item
  \textbf{The Chow scheme and its coordinate map.} The morphism

  \[
  \mu:(E^*)^{\times n}\longrightarrow\operatorname{Sym}^n E^*,
  \qquad(\ell_1,\ldots,\ell_n)\longmapsto\ell_1\cdots\ell_n
  \]

  is defined over \(\mathbb Z\). Its pullback on degree-\(d\) coordinate
  functions lands in separately degree-\(d\) functions, invariant under
  \(S_n\), and over \(\mathbb Q\) identifies with \(h_{d,n}\). The
  affine scheme-theoretic image has coordinate algebra
  \(\operatorname{im}\mu^*\); this realizes the Chow cone and gives
  \(\ker h_{d,n}=I(\mathrm{Chow}_n)_d\). Thus R-002 is an exact equation
  for the rank of the rectangular component of a morphism of coordinate
  algebras. For \(n=2,4,6\), the finite certificates prove rank one in
  that component. This coordinate-map assertion is also Kumar--Landsberg
  Lemma 1.5.
\item
  \textbf{The exact \(K_0\) relation.} For a line \(L\) over a field and
  an endomorphism \(f:L\to L\), its kernel-image-cokernel exact sequence
  gives \([\ker f]=[\operatorname{coker}f]\). For the map \(f=0\), both
  classes equal \([L]\); for \(f=1\), both classes are zero. Both
  endomorphisms have domain and codomain \(L\). This computes exactly
  the information supplied by the class identity \([L]-[L]=0\). R-002
  records the additional morphism \(c_n\operatorname{id}_L\). The
  associated \(\lambda\)-operation has \(\lambda_t([L])=1+[L]t\), while
  \(\gamma_t([L]-1)=1+([L]-1)t\); neither formula changes with the
  chosen endomorphism of the line, as is verified by the two explicit
  maps.
\item
  \textbf{Finite étale schemes and weights.} Let \(\mathcal L_n^+\) and
  \(\mathcal L_n^-\) be the constant finite schemes over \(\mathbb Z\)
  indexed by column-even and column-odd Latin squares. Their virtual
  class is
  \([\mathcal L_n^+]-[\mathcal L_n^-]=\Delta_n[\operatorname{Spec}\mathbb Z]\).
  Over every finite field, each geometric point contributes one
  eigenvalue one to \(H_c^0\), and higher cohomology vanishes. Thus the
  difference of Frobenius traces is exactly \(\Delta_n\). This gives a
  direct finite-scheme and weight-zero dictionary for the same signed
  period numerator. For odd \(n>1\), swapping two rows reverses the
  column sign and gives an involution pairing all squares, so this
  virtual class and trace are zero. For even \(n=2,4,6\), their values
  are the three nonzero integers in the certificate. These are exact
  computations of the maps and classes involved, with no assertion
  excluding other geometric relations.
\end{enumerate}

The parent paper is to supply the exact SGA/EGA/FGA locators for these
standard constructions from its independent archive/SGA audit. This
report does not invent unverified exposé or proposition numbers.

\subsection{Exact inventory of supplied literature relevant to item
4}\label{exact-inventory-of-supplied-literature-relevant-to-item-4}

All statements in this inventory are attributed results, not new
findings. Parameter domains are written explicitly.

\subsubsection{Ikenmeyer--Panova (IP),
arXiv:1512.03798v2}\label{ikenmeyerpanova-ip-arxiv1512.03798v2}

Let \(\rho(N)=(N-|\rho|,\rho)\), whenever it is a partition, and
\(\bar\lambda\) be the tail after removing the first row. Let

\[
\mathcal X=\{(1),(1,1),(1,1,1,1),(1^6),(2,1),(3,1)\}.
\]

\begin{longtable}[]{@{}
  >{\raggedright\arraybackslash}p{(\columnwidth - 2\tabcolsep) * \real{0.5000}}
  >{\raggedright\arraybackslash}p{(\columnwidth - 2\tabcolsep) * \real{0.5000}}@{}}
\toprule\noalign{}
\begin{minipage}[b]{\linewidth}\raggedright
Locator
\end{minipage} & \begin{minipage}[b]{\linewidth}\raggedright
Exact content used
\end{minipage} \\
\midrule\noalign{}
\endhead
\bottomrule\noalign{}
\endlastfoot
Theorem 1.6, p.3 & For \(n>3m^4\), support of the padded-permanent
coordinate ring in degree \(d\) is contained in
\(\{\lambda:g(\lambda,(d^n),(d^n))>0\}\). \\
Theorem 1.8, pp.3--4 (Kadish--Landsberg) & That support has
\(\ell(\lambda)\le m^2\), \(|\bar\lambda|\le md\). \\
Corollary 1.9, p.4 & On \(|\bar\lambda|\le md\), a strict inequality
\(a_\lambda(d[n])>g(\lambda,(d^n),(d^n))\) requires \(d>n/m\). \\
Theorem 1.10(a), p.4 & Every \(\lambda\) with
\(\bar\lambda\in\mathcal X\) has \(a_\lambda(d[n])=0\). \\
Theorem 1.10(b), p.4 & For \(m\ge3\), \(\bar\lambda\notin\mathcal X\),
\(\ell(\lambda)\le m^2\), \(|\bar\lambda|\le md\), \(d>3m^3\),
\(n>3m^4\), the rectangular Kronecker coefficient is positive. \\
Theorem 2.1, p.5 (Manivel, Theorem 1) & For \(n\ge|\rho|\),
\(g(\rho(nd),(d^n),(d^n))=\dim S_\rho(\mathfrak{sl}_d)^{\mathrm{SL}_d}\).
The function is symmetric and nondecreasing in \(n,d\). \\
Proposition 2.8, p.7 & For \((n,d)\in\mathrm{St}^1(\rho)\), in
particular \(n\ge|\rho|\),
\(g(\rho(nd),(d^n),(d^n))\ge a_{\rho(nd)}(d[n])\). \\
Theorem 4.7, p.11 & For \(\nu\notin\mathcal X\),
\(L=\max(\ell(\nu)+1,9)\), \(a>3L^{3/2}\), \(b\ge3L^2\), and
\(|\nu|\le ab/6\), \(g(\nu(ab),(b^a),(b^a))>0\). \\
Corollary 5.1, p.13 & The stable limit \(A_\rho\) vanishes exactly for
\(\rho\in\mathcal X\). Here \(A_\rho\) is a stable rectangular Kronecker
coefficient, with its invariant-Schur-functor realization in Theorem
2.1. \\
Proposition 5.2, p.13 & For \(i\ge2\), \(m\ge7\), \(0\le k<m^2\),
\(g(i(m^2-k,1^k),((im)^m),((im)^m))>0\). \\
Corollary 5.3, p.14 & For \(m\ge7\), every even-row partition \(\rho\)
of length less than \(m^2\) has
\(\dim S_\rho(\mathfrak{sl}_m)^{\mathrm{SL}_m}>0\). \\
Corollary 5.4, p.14 & For \(d\ge7\), the group generated by the
rectangular Kronecker support semigroup in \(\mathbb Z^{d^2}\) is
\(\{z:\sum z_i\equiv0\pmod d\}\). \\
Proposition 6.4, p.15 & The hook coefficients have the subset-count
formula used in R-004. \\
\end{longtable}

\subsubsection{Bürgisser--Ikenmeyer--Panova (BIP),
arXiv:1604.06431v3}\label{buxfcrgisserikenmeyerpanova-bip-arxiv1604.06431v3}

\begin{longtable}[]{@{}
  >{\raggedright\arraybackslash}p{(\columnwidth - 2\tabcolsep) * \real{0.5000}}
  >{\raggedright\arraybackslash}p{(\columnwidth - 2\tabcolsep) * \real{0.5000}}@{}}
\toprule\noalign{}
\begin{minipage}[b]{\linewidth}\raggedright
Locator
\end{minipage} & \begin{minipage}[b]{\linewidth}\raggedright
Exact content used
\end{minipage} \\
\midrule\noalign{}
\endhead
\bottomrule\noalign{}
\endlastfoot
Theorem 1.4, p.3 & For \(n\ge m^{25}\), support of the padded-permanent
coordinate ring is contained in support of the determinant-orbit-closure
coordinate ring. \\
Proposition 2.3, p.5 & \(n\ge k\ell\), \(\ell\) even:
\((k\times\ell)^{\sharp nk}\) occurs in determinant-orbit-closure degree
\(k\). \\
Proposition 2.4, p.5 & \(|\bar\lambda|\le md\), \(md^2\le n\):
restriction to the determinant closure is injective on the
\(\lambda\)-highest-weight space. The natural reading of ``every highest
weight vector'' is every nonzero such vector. \\
Theorem 2.5, p.5 & \(n\ge sk\): every \(X^{n-s}\sum_{i=1}^k\varphi_i^s\)
belongs to the determinant orbit closure. \\
Corollary 4.8, p.13 & Explicit determinant-contraction nonvanishing on
power sums for even inner degree. \\
Theorem 5.4, p.15 &
\(\langle\kappa^d_{m,n}(f),q^d\rangle=\langle f,(M^*q)^d\rangle\). \\
Proposition 5.6(1), p.17 & \(n\ge m\), \(\mu\vdash md\), \(\mu_2\le m\):
\(\kappa\) induces an isomorphism of the highest-weight spaces of
weights \(\mu\) and \(\mu+(d(n-m))\). \\
\end{longtable}

The BIP support inclusion constrains occurrence obstructions. The
morphism-level question is retained by the restriction map and by its
ranks on multiplicity spaces. No conclusion on P versus NP is claimed
here.

\subsection{Source-proof audit notes and boundaries of
checking}\label{source-proof-audit-notes-and-boundaries-of-checking}

\begin{enumerate}
\def\labelenumi{\arabic{enumi}.}
\item
  IP Theorem 6.1's displayed exception
  \texttt{\small d\textasciicircum{}1-\allowbreak{}2} is a typographical
  error: Proposition 6.4 and the proof's table use \(d^2-2\). Its
  additionally stated strict inequality \(g_k>g_{k-2}\) for all
  \(k\le d^2/2\) also fails literally, for example \(d=7,k=2\) gives
  \(0>1\). We use the exact Proposition 6.4 instead. This discrepancy is
  proved by the generating function, rather than silently correcting the
  printed theorem.
\item
  In the supplied proof of IP Theorem 4.7, the line
  \(w<\sqrt{L+7}+1\le(L+1)/2\), with
  \(w=\max(\lceil\sqrt{L+7}\rceil,7)\), is false at \(L=9\): it asserts
  \(7<5\). This is a gap in that displayed estimate, not a
  counterexample to the theorem. The broad theorem is cited as
  literature, and is not certified here as an independently checked
  proof. No new proof in this report depends on that estimate. The
  parent can record this explicitly or omit reliance on Theorem 4.7.
\item
  The hook counterexample, parity conversion, averaging normalization,
  and row-extended evaluation were read and checked in full. A first
  informal scalar derivation had called \(t_0\) invariant under all of
  \(H\); that statement fails for row permutations. R-002's final proof
  explicitly averages \(t_0\), uses evenness to compute its pairing, and
  moves the row permutations through \(\tau\). This repairs the issue
  without changing the scalar.
\item
  Computational certificates are exact finite certificates for their
  stated inputs. They do not certify a universal nonvanishing theorem or
  the entire supplied literature.
\end{enumerate}

\subsection{Primary-source URLs for root
verification}\label{primary-source-urls-for-root-verification}

\begin{itemize}
\tightlist
\item
  IP: \url{https://arxiv.org/pdf/1512.03798} (v2).
\item
  BIP: \url{https://arxiv.org/pdf/1604.06431} (v3).
\item
  Kumar--Landsberg: \url{https://arxiv.org/pdf/1410.8585} (v1), Theorem
  1.9, Proposition 2.1, Lemma 1.5; Howe Proposition 4.3 cited in §2.
\item
  Alpöge:
  \url{https://people.math.harvard.edu/~alpoge/papers/combinatorica\%20paper.pdf},
  Theorem 2 and Corollary 5.
\item
  Bürgisser--Christandl--Ikenmeyer, \emph{Even partitions in plethysms}:
  \url{https://www3.math.tu-berlin.de/algebra/work/weintraub.pdf}, §1.1
  unnumbered main theorem; arXiv:1003.4474v2. Its domain
  \(\lambda\vdash kn,\ell(\lambda)\le k\) yields
  \(a_{2\lambda}(k[2n])>0\).
\item
  Giannelli, \emph{Representations of symmetric groups}:
  \url{https://eugeniomaths.wordpress.com/wp-content/uploads/2013/10/tesi-final-version-2015.pdf},
  Theorem 2.1.2 and Lemma 2.3.3.
\end{itemize}

Web references available in this child run:
\texttt{\small turn614803048view0} (IP),
\texttt{\small turn614803048view1} (BIP),
\texttt{\small turn768780200view0} (Kumar--Landsberg),
\texttt{\small turn245878279view0} (Alpöge),
\texttt{\small turn517469220search12} (BCI). The root must open them
independently before citing them in chat.

\section{Canonical-map boundary checks and prime rectangular
blocks}\label{canonical-map-boundary-checks-and-prime-rectangular-blocks}

Run: /root/rectangular\_cases, 21 November 2026 (Europe/Berlin). Model:
GPT-6 Astra Pro. The inherited reasoning-effort setting is not exposed.

This supplements rectangular\_report.md. The scalar normalization and
the determinant tensors \(I_{a,b}\) use that report's R-001 and R-002.
The new formulas below are explicit consequences of established results
of Glynn, Drisko, Howe, and the algebra morphism defining the canonical
map. No claim of priority is made for their nonvanishing content. The
full surjectivity bound proved in Sra III remains \(b\ge(a-1)^2\); the
results below prove specified components at smaller degrees.

\subsection{Exact map and coefficient
conventions}\label{exact-map-and-coefficient-conventions}

Put \(V=\mathbb C^a\), with its orthonormal basis \(e_1,\ldots,e_a\),
and

\[
 R_a=\bigoplus_{b\ge0}\left((\operatorname{Sym}^b V)^{\otimes a}\right)^{S_a}.
\]

The product is componentwise multiplication. Symmetric powers are
realized as invariant tensors by normalized averaging. There is a graded
algebra morphism

\[
 \mu_a:\operatorname{Sym}(\operatorname{Sym}^aV)\longrightarrow R_a,
 \qquad \mu_a(v^a)=v^{\otimes a}.
 \tag{B1}
\]

Its degree-\(b\) component is \(h_{b,a}\). Thus

\[
 h_{b,a}(u_1^a\cdots u_b^a)=(u_1\cdots u_b)^{\otimes a}.
 \tag{B2}
\]

For even \(b\), let \(I_{a,b}\) be obtained by putting the unnormalized
alternating \(a\)-tensor in each of \(b\) columns, and averaging within
the \(a\) rows. Permuting the rows has sign \(\operatorname{sgn}^b=1\),
so further outer averaging makes no change. This tensor satisfies

\[
 \langle I_{a,b},v_1^b\otimes\cdots\otimes v_a^b\rangle
       =\det(v_1,\ldots,v_a)^b.
 \tag{B3}
\]

The inner product is conjugate-linear in its first argument. All tensors
appearing in the coefficient calculations have rational real
coefficients. The coefficient at the tensor having row \(i\) constantly
\(e_i\) is one. The determinant line in \(R_{a,b}\) is one-dimensional
by R-001. Products of the determinant tensors satisfy the actual
equality in the algebra

\[
 I_{a,b}I_{a,c}=I_{a,b+c}\qquad(b,c\text{ even}).
 \tag{B4}
\]

Indeed, evaluation on arbitrary \(a\) covectors turns both sides into
the product of determinant powers. Such evaluations determine the
tensor.

\subsection{R-007. Prime-adjacent square scalars and
periods}\label{r-007.-prime-adjacent-square-scalars-and-periods}

Let \(p\) run through the odd primes, and write

\[
 \mathcal P_n=\int_{\operatorname{SU}(n)}\prod_{i,j=1}^{n}u_{ij}\,du,
 \qquad
 B_n=\prod_{j=0}^{n-1}\frac{(n+j)!}{j!}.
\]

The normalized square scalar from R-002 satisfies

\[
 \boxed{c_{p-1}\in1+p\mathbb Z_{(p)},\qquad
        p^{p-1}c_{p+1}\in1+p\mathbb Z_{(p)}.}
 \tag{R-007a}
\]

In particular these are nonzero scalars, with exact valuations

\[
 v_p(c_{p-1})=0,\qquad v_p(c_{p+1})=1-p,
 \qquad
 v_p(\mathcal P_{p-1})=2-p,\qquad
 v_p(\mathcal P_{p+1})=-p.
 \tag{R-007b}
\]

Here \(\mathbb Z_{(p)}\) denotes the rational numbers whose denominator
is coprime to \(p\); the integral is the rational number identified by
the coefficient formula of R-002.

\subsubsection{Source statements and proof of this
normalization}\label{source-statements-and-proof-of-this-normalization}

Glynn, \emph{The conjectures of Alon--Tarsi and Rota in dimension prime
minus one}, Theorem 3.2, states

\[
 E_{p-1}-O_{p-1}\equiv(-1)^{(p-1)/2}\pmod p.
\]

Drisko, \emph{On the number of even and odd Latin squares of order
\(p+1\)}, Theorem 9, pp.~33--34, states

\[
 E_{p+1}-O_{p+1}\equiv(-1)^{(p+1)/2}p^2\pmod {p^3}.
\]

Both use row-times-column parity. R-002 proves

\[
 \Delta_n=(-1)^{n(n-1)/2}(E_n-O_n),\qquad
 c_n=\frac{\Delta_n}{(n!)^n},\qquad \mathcal P_n=\frac{\Delta_n}{B_n}.
\]

For even \(n\), \(n(n-1)/2\equiv n/2\pmod2\). Consequently the two
source congruences become

\[
 \Delta_{p-1}\equiv1\pmod p,
 \qquad \Delta_{p+1}\equiv p^2\pmod {p^3}.
\]

Wilson's theorem gives \((p-1)!\equiv-1\pmod p\), and

\[
 \frac{(p+1)!}{p}=(p+1)(p-1)!\equiv-1\pmod p.
\]

The exponents \(p-1\) and \(p+1\) are even. Dividing by the
corresponding factorial units proves (R-007a), and its first two
valuation assertions. In \(B_{p-1}\), the numerator factorials from
\(p!\) through \((2p-3)!\) each contain exactly one factor of \(p\),
while all denominator factorials are prime to \(p\). Thus
\(v_p(B_{p-1})=p-2\). In \(B_{p+1}\), the \(p-1\) factorials
\((p+1)!,\ldots,(2p-1)!\) have valuation one, the two factorials
\((2p)!\) and \((2p+1)!\) have valuation two, and the denominator
contributes the one factor in \(p!\). Thus
\(v_p(B_{p+1})=(p-1)+4-1=p+2\). This proves the period valuations. The
ranges include \(p=3\).

These calculations use the established prime-adjacent Alon--Tarsi cases;
they do not supply further orders for that conjecture.

\subsection{R-008. A canonical determinant block for every
prime-minus-one
width}\label{r-008.-a-canonical-determinant-block-for-every-prime-minus-one-width}

For every odd prime \(p\), every \(2\le a\le p-1\), and \(b=p-1\), put

\[
 m_a=e_1\cdots e_a\in\operatorname{Sym}^aV,
 \qquad
 q_{a,b}=\langle I_{a,b},h_{b,a}(m_a^b)\rangle.
\]

Then

\[
 \boxed{q_{a,p-1}\in1+p\mathbb Z_{(p)}.}
 \tag{R-008a}
\]

In particular the canonical map \(h_{p-1,a}\) has rank one on the target
determinant isotypic component \(S_{((p-1)^a)}V\). The tensor \(m_a^b\)
is an explicit test vector for that component.

More explicitly, in tensor degree \(ab\) set

\[
 e_{(b^a)}=\frac{f^{(b^a)}}{(ab)!}
       \sum_{\sigma\in S_{ab}}\chi^{(b^a)}(\sigma)\sigma,
 \qquad N_{a,b}=\langle I_{a,b},I_{a,b}\rangle>0.
\]

This is the central Schur projector, with \(f^{(b^a)}\) the dimension of
the indicated Specht module. Define the concrete source vector

\[
 u_{a,p}=\frac{N_{a,p-1}}{q_{a,p-1}}
             e_{((p-1)^a)}m_a^{p-1}.
 \tag{B5}
\]

It satisfies the exact equality

\[
 h_{p-1,a}(u_{a,p})=I_{a,p-1}.
 \tag{R-008b}
\]

\subsubsection{Proof of the unit
contraction}\label{proof-of-the-unit-contraction}

Glynn's Theorem 3.1, together with the formula defining his modular
hyperdeterminant immediately before that theorem, gives the following
coefficient congruence. For an \(a\times a\) matrix \(E=(e_{ij})\) of
nonnegative integers having every row and every column sum \(b=p-1\),

\[
 C_E=[X^E]\det(X)^b
       \equiv\frac{(-1)^a}{\prod_{i,j}e_{ij}!}\pmod p.
 \tag{B6}
\]

Every denominator in (B6) is a unit modulo \(p\). Write a general
degree-\(b\) vector as

\[
 P=\sum_{|\alpha|=b}P_\alpha e^\alpha,
 \qquad
 F(P)=\langle I_{a,b},P^{\otimes a}\rangle.
\]

Let \(M(\alpha)=b!/\prod_j\alpha_j!\). Evaluation of (B3), first on
independent row vectors and then on their monomial expansions, gives the
exact coefficient formula

\[
 F(P)=\sum_E \frac{C_E}{\prod_i M(E_i)}\prod_i P_{E_i},
 \tag{B7}
\]

where \(E_i\) is row \(i\) of \(E\). This accounts for the normalized
embedding of each symmetric monomial. Wilson's theorem gives

\[
 \prod_iM(E_i)\equiv
 \frac{(-1)^a}{\prod_{i,j}e_{ij}!}\pmod p.
\]

Thus every ratio in (B7) is one modulo \(p\). Equivalently, with
\(P(x)=\sum P_\alpha x^\alpha\),

\[
 F(P)\equiv[x_1^b\cdots x_a^b]P(x)^a\pmod p.
 \tag{B8}
\]

Now introduce independent \(z_{ij}\), \(1\le i\le a\), \(1\le j\le b\),
put \(\ell_j=\sum_i z_{ij}e_i\), and substitute \(P=\ell_1\cdots\ell_b\)
in (B8). By (B2), the left side becomes

\[
 \left\langle I_{a,b},h_{b,a}(\ell_1^a\cdots\ell_b^a)\right\rangle.
\]

Its coefficient at the squarefree monomial \(\prod_{i,j}z_{ij}\) is
\((a!)^b q_{a,b}\), since the squarefree coefficient in \(\ell_j^a\) is
\(a!m_a\). On the right side of (B8), the same coefficient in
\(\prod_j\ell_j(x)^a\) is \((a!)^b(x_1\cdots x_a)^b\), so it becomes
\((a!)^b\) after coefficient extraction in \(x\). The integer \(a!\) is
prime to \(p\). Cancelling it proves (R-008a).

All tensor averages used here have denominators prime to \(p\), because
\(a,b<p\), so the stated congruences are in the specified local ring.

\subsubsection{Proof of the actual
preimage}\label{proof-of-the-actual-preimage}

The central projector \(e_{(b^a)}\) commutes with every tensor-position
permutation and every averaging idempotent. In particular it preserves
both nested symmetric powers and commutes with \(h_{b,a}\). On the
target it is orthogonal projection onto the line spanned by \(I_{a,b}\).
Character orthogonality gives this projector formula, and R-001 supplies
the target multiplicity one. Therefore

\[
 h_{b,a}(e_{(b^a)}m_a^b)
   =e_{(b^a)}h_{b,a}(m_a^b)
   =\frac{q_{a,b}}{N_{a,b}}I_{a,b}.
\]

The scalar \(q_{a,b}\) is nonzero by (R-008a). Multiplication by the
factor in (B5) proves (R-008b). Thus the proof identifies the particular
canonical morphism and an explicit source vector.

\subsection{R-009. Products and first-row extensions below the full
bound}\label{r-009.-products-and-first-row-extensions-below-the-full-bound}

For every odd prime \(p\), \(2\le a\le p-1\), \(t\ge1\), and \(s\ge0\),
set

\[
 b=t(p-1)+s,\qquad
 \lambda=(as+t(p-1),\;\underbrace{t(p-1),\ldots,t(p-1)}_{a-1}).
\]

The following source and image are explicit:

\[
 \boxed{
 h_{b,a}\bigl(u_{a,p}^{\,t}(e_1^a)^s\bigr)
   =I_{a,t(p-1)}(e_1^{\otimes a})^s\ne0.
 }
 \tag{R-009a}
\]

Both sides are highest weight vectors of weight \(\lambda\). Moreover

\[
 [S_\lambda V:\operatorname{im}h_{b,a}]
   =[S_\lambda V:\operatorname{Sym}^a(\operatorname{Sym}^bV)]=1
 \tag{R-009b}
\]

in every ambient dimension at least \(a\).

\subsubsection{Proof}\label{proof-2}

The image equality follows by applying the algebra morphism (B1), using
(R-008b), (B4), and \(\mu_a(e_1^a)=e_1^{\otimes a}\). Evaluation on
\(a\) covectors gives the nonzero polynomial

\[
 \det(X)^{t(p-1)}\prod_{j=1}^a x_{1j}^{s}.
\]

This proves the asserted nonvanishing without a cancellation assertion.
The determinant vector is highest weight of weight \((t(p-1)^a)\); the
added factor has weight \((as)\) and is killed by the upper triangular
raising operators. The same is true of the source vector (B5), which
lies in a determinant representation over \(\mathbb C^a\). Embedding
\(\mathbb C^a\subseteq V\) gives highest weight vectors in every larger
ambient space, since the additional raising operators differentiate
basis directions absent from these tensors.

For the target multiplicity, set \(\ell=t(p-1)\). The inner-degree
lifting map \(\kappa^a_{\ell,b}\), obtained by multiplying each inner
factor by \(e_1^s\), maps the unique highest weight vector of weight
\((\ell^a)\) to the displayed image. BIP Proposition 5.6(1) applies
because the second part of \((\ell^a)\) is \(\ell\). It identifies the
two highest weight spaces. Howe Proposition 4.3, equivalently R-001,
gives dimension one because \(\ell\) is even. This proves (R-009b).

\subsubsection{Actual boundary families and
examples}\label{actual-boundary-families-and-examples}

For each prime \(p\equiv1\pmod4\), \(p\ge13\), put

\[
 a=(p-1)/2,\qquad t=(a-2)/2,\qquad s=0.
\]

Then \(a\) is even, and (R-009) proves the entire determinant component
at

\[
 b=t(p-1)=a(a-2)=(a-1)^2-1.
 \tag{B9}
\]

Examples are \((a,b)=(6,24),(8,48),(14,168),(18,288),(20,360),
(26,624),(30,840)\). The case \(a=26\) uses \(p=53\) and \(t=12\); its
base rectangle has width \(52\).

Boundary examples in odd outer degree are obtained with \(t=1\) and
positive \(s\). For \((a,b,p)=(7,35,31)\), the image has its unique
component \(S_{(65,30^6)}V\). For \((a,b,p)=(9,63,61)\), the image has
its unique component \(S_{(87,60^8)}V\). The degree sums are
respectively \(65+6\cdot30=7\cdot35\) and \(87+8\cdot60=9\cdot63\).

The square scalar yields a second exact form. For every
\(n\in\{p-1,p+1:p\text{ odd prime}\}\) and \(k\ge1\),

\[
 h_{kn,n}(I_{n,n}^{\,k})=c_n^{\,k}I_{n,kn}\ne0.
 \tag{B10}
\]

This follows directly from R-002 and (B1). For \(n=6,k=4\), the scalar
is \(699840000^{-4}\), since \(c_6=-699840000^{-1}\). This verifies the
determinant component at \((a,b)=(6,24)\) by an explicit rational
scalar.

\subsection{R-010. What the last common-power step proves at the
boundary}\label{r-010.-what-the-last-common-power-step-proves-at-the-boundary}

Write \(r=a-1\), and consider the actual differential operator

\[
 \Psi_{r,b}=\prod_{i=1}^r(x\cdot\partial_{v_i})^r
\]

on polynomials initially independent of \(x\), separately of degree
\(b\) in each \(v_i\). Sra III Lemma 2.1 and the proof of Theorem 1.3
give its injectivity for \(b\ge r^2\). The supplement's Corollary S8.1
records the quantitative estimate; immediately after it, p.~36, the
author gives the example \(r=5,b=24\).

For every odd \(r\ge3\), put \(b=r^2-1\) and \(V=\mathbb C^r\). Then

\[
 f(v_1,\ldots,v_r)=\det(v_1,\ldots,v_r)^b
 \quad\text{is nonzero and symmetric,}\qquad \Psi_{r,b}f=0.
 \tag{R-010a}
\]

Indeed \(b\) is even, so the determinant power is symmetric. Every term
in the differentiated expression inserts \(r^2=b+1\) copies of the same
vector \(x\) into \(b\) determinant factors. A determinant receiving two
copies of \(x\) vanishes, and every term has such a factor. For even
\(r\), the same determinant polynomial is alternating because \(b\) is
odd. Thus its projection to the symmetric \(r\)-slot input is zero. This
last parity calculation does not classify the other kernel components.

The arbitrary-common-power full-tensor kernel admits nonrectangular
components. A concrete example of the analogous operator with \(r=3\),
power one, and input degree two per slot has an additional kernel vector
on \(V=\mathbb C^2\). Write \([ij]=x_i y_j-y_i x_j\) and set

\[
 g=[12][13][23].
\]

It has degree two in each vector and transforms by the sign of the slot
permutation. Under \(v_i\mapsto v_i+t_i x\), each bracket is at most
linear in the \(t_i\), and direct expansion gives

\[
 [t_1t_2t_3]\,
 [v_1+t_1x,v_2+t_2x]
 [v_1+t_1x,v_3+t_3x]
 [v_2+t_2x,v_3+t_3x]=0.
 \tag{R-010b}
\]

The two possible terms are \(-[x,v_2][x,v_3][x,v_1]\) and
\(+[x,v_1][x,v_2][x,v_3]\), which cancel. The vector has
\(\operatorname{GL}_2\)-weight \((3,3)\), whereas the height-three
rectangle of degree two vanishes on \(\mathbb C^2\). This is an exact
counterexample to a uniform rectangular-kernel classification across
arbitrary common powers. It does not decide the symmetric-input kernel
at \(r\) even and power \(r\).

The current proof does not establish full canonical-map surjectivity at
\(b=(a-1)^2-1\) for general odd \(a\). The unconditional improvements in
this report are the specified canonical-map components (R-008--R-009).

\subsection{Independent finite coefficient
checks}\label{independent-finite-coefficient-checks}

Let \(D_E=[X^E]\operatorname{per}(X)^b\). The coefficient of
\(\prod_{i,j}z_{ij}\) in \(\prod_iP_{E_i}\), after
\(P=\prod_{j=1}^b\ell_j\), is \(D_E\): each of its \(b\) columns assigns
the \(a\) symbols bijectively to the \(a\) tensor rows, and the
resulting row multiplicities are \(E\). Combining this with (B7) and its
\((a!)^b\) squarefree coefficient gives

\[
 q_{a,b}=
 \frac{\sum_E C_ED_E\prod_{i,j}e_{ij}!}
      {(a!)^b(b!)^a}.
 \tag{B11}
\]

The numerator is the Fischer pairing of the determinant power and
permanent power. The program verify\_prime\_blocks.py computes their
coefficient tables by exact monomial multiplication, checks (B6) at
every exponent matrix, and evaluates (B11). It imports no source-paper
implementation.

\begin{longtable}[]{@{}rrrr@{}}
\toprule\noalign{}
\(a\) & \(b=p-1\) & Exponent matrices checked & Exact \(q_{a,b}\) \\
\midrule\noalign{}
\endhead
\bottomrule\noalign{}
\endlastfoot
2 & 2 & 3 & \(1/4\) \\
2 & 4 & 5 & \(1/16\) \\
3 & 4 & 120 & \(1/1296\) \\
3 & 6 & 406 & \(1/87480\) \\
4 & 4 & 10147 & \(1/331776\) \\
3 & 10 & 2211 & \(1/88179840\) \\
3 & 12 & 4186 & \(281/838061199360\) \\
\end{longtable}

All seven residues are one modulo the indicated prime. The two square
cases also equal \(c_a^2\), using the independently enumerated Latin
scalars in rectangular\_report.md. The nondivisible case \((a,b)=(3,4)\)
checks that the general coefficient argument does not require
\(a\mid b\).

\subsection{Source audit and exact
locators}\label{source-audit-and-exact-locators}

\begin{itemize}
\tightlist
\item
  Suvrit Sra, \emph{The q-Foulkes problem III: On Foulkes' conjecture
  with a graded inner factor}, v1, Theorem 1.3, Lemma 2.1, and §2.2,
  pp.~3--5. Author's repository:
  \url{https://github.com/suvrit/preprints/blob/main/preprints/qfoulkes3/v1/qfoulkes3-v1.pdf}.
\item
  Sra, \emph{Supplement to the q-Foulkes project}, S8, p.~36, Corollary
  S8.1 and its following paragraph. Author's repository:
  \url{https://github.com/suvrit/preprints/blob/main/preprints/qfoulkes3/v1/qfoulkes3-v1-supplement.pdf}.
\item
  David G. Glynn, \emph{The conjectures of Alon--Tarsi and Rota in
  dimension prime minus one}, SIAM J. Discrete Math. \textbf{24} (2010),
  394--399, Theorems 3.1--3.2 and the definition immediately before
  Theorem 3.1. DOI: \url{https://doi.org/10.1137/090773751}. The
  author's full manuscript was read through its publicly indexed
  transcript at
  \url{https://www.academia.edu/25562085/The_Conjectures_of_Alon_Tarsi_and_Rota_in_Dimension_Prime_Minus_One}.
\item
  Arthur A. Drisko, \emph{On the number of even and odd Latin squares of
  order p+1}, Adv. Math. \textbf{128} (1997), 20--35, Theorem 9,
  pp.~33--34. Publisher record:
  \url{https://www.sciencedirect.com/science/article/pii/S0001870897916236}.
  The original article's full transcript, including its page numbers,
  was read at
  \url{https://studyres.com/doc/19593200/on-the-number-of-even-and-odd-latin-squares-of}.
\item
  The 1998 assertion extending prime cases by powers of two is not used.
  Glynn's introduction explicitly identifies the error in the preceding
  doubling argument. No such enlargement is included in these parameter
  families.
\item
  Howe Proposition 4.3 and BIP Proposition 5.6(1) are located and
  checked in rectangular\_report.md, R-001 and R-005.
\end{itemize}

\section{Two exact symmetric common-polarization
kernels}\label{two-exact-symmetric-common-polarization-kernels}

Statement ID \textbf{R-011}. Written and checked on 2026-11-20 UTC by
GPT-6 Astra Pro, agent
\texttt{\small /\allowbreak{}root/\allowbreak{}rectangular\_\allowbreak{}cases};
the inherited effort setting is not exposed. The statements below are
finite, unconditional computational theorems. They do not assert a
classification at arbitrary parameters.

\subsection{The actual maps and
results}\label{the-actual-maps-and-results}

For a complex vector space \(V\), use four separately homogeneous vector
slots \(v_1,\ldots,v_4\), of degree \(b\) each, and an additional vector
slot \(x\). Define the natural linear map

\[
 T_{L,b}:\operatorname{Sym}^4(\operatorname{Sym}^bV)
 \longrightarrow \operatorname{Sym}^{4L}V\otimes
 \operatorname{Sym}^4(\operatorname{Sym}^{b-L}V),
 \qquad
 T_{L,b}(f)=[t_1^L\cdots t_4^L]
 f(v_1+t_1x,\ldots,v_4+t_4x).
\]

Equivalently,
\(T_{L,b}=(L!)^{-4}\prod_{i=1}^4(x\cdot\partial_{v_i})^L\). The
polynomial description means the canonical multihomogeneous tensor
model; using coordinate functions replaces \(V\) throughout by its dual
and gives the same natural map after that explicit substitution.
Permuting the four slots commutes with coefficient extraction.
Consequently the displayed target has its stated fourth symmetric power.
Each term has \(x\)-degree \(4L\) and residual degree \(b-L\) in every
slot.

\textbf{Theorem R-011.} The two maps \(T_{1,3}\) and \(T_{2,7}\) are
injective for every finite-dimensional complex vector space \(V\).

Both parameters satisfy \(b=4L-1\). The ranks were established on every
source highest-weight multiplicity space, not inferred from source and
target multiplicities.

\begin{longtable}[]{@{}
  >{\raggedright\arraybackslash}p{(\columnwidth - 10\tabcolsep) * \real{0.1304}}
  >{\raggedleft\arraybackslash}p{(\columnwidth - 10\tabcolsep) * \real{0.1739}}
  >{\raggedleft\arraybackslash}p{(\columnwidth - 10\tabcolsep) * \real{0.1739}}
  >{\raggedleft\arraybackslash}p{(\columnwidth - 10\tabcolsep) * \real{0.1739}}
  >{\raggedleft\arraybackslash}p{(\columnwidth - 10\tabcolsep) * \real{0.1739}}
  >{\raggedleft\arraybackslash}p{(\columnwidth - 10\tabcolsep) * \real{0.1739}}@{}}
\toprule\noalign{}
\begin{minipage}[b]{\linewidth}\raggedright
\(L,b\)
\end{minipage} & \begin{minipage}[b]{\linewidth}\raggedleft
Source dimension on \(\mathbb C^4\)
\end{minipage} & \begin{minipage}[b]{\linewidth}\raggedleft
Distinct source constituents
\end{minipage} & \begin{minipage}[b]{\linewidth}\raggedleft
Sum of source multiplicities
\end{minipage} & \begin{minipage}[b]{\linewidth}\raggedleft
Largest source multiplicity
\end{minipage} & \begin{minipage}[b]{\linewidth}\raggedleft
Certified image rank on highest-weight spaces
\end{minipage} \\
\midrule\noalign{}
\endhead
\bottomrule\noalign{}
\endlastfoot
\(1,3\) & 8,855 & 12 & 12 & 1 & 12 \\
\(2,7\) & 9,078,630 & 173 & 319 & 6 & 319 \\
\end{longtable}

This verifies the first two requested common-power boundary cases for
four symmetric slots. It supplies finite evidence relevant to, but no
theorem about, the general even-\(r\) boundary
\(\prod_{i=1}^r(x\cdot\partial_{v_i})^r\) at \(b=r^2-1\). The actual
case \(r=L=4,b=15\) is now separately proved and independently certified
as R-012 in polarization\_r4\_boundary\_and\_geometry.md.

\subsection{Complete proof by exact
certificates}\label{complete-proof-by-exact-certificates}

\subsubsection{1. Reduction to multiplicity spaces in dimension
four}\label{reduction-to-multiplicity-spaces-in-dimension-four}

The source is the invariant part of
\((\operatorname{Sym}^b V)^{\otimes4}\) under permutation of its four
factors. Iterated Pieri multiplication of four one-row Schur characters
has at most four rows: adding a horizontal strip to a diagram with \(k\)
rows creates at most one additional row. Hence every source Schur
functor is \(S_\lambda V\) with \(\ell(\lambda)\le4\).

The construction of \(T_{L,b}\) uses permutations, symmetrizations and
contraction of slot labels in tensor powers, so its matrices on the
multiplicity spaces of the Schur functors are independent of \(\dim V\)
once \(\dim V\ge4\). This can also be seen by the tensor model
\(V^{\otimes4b}\): the same element of the rational permutation algebra
acts on each \(V^{\otimes4b}\), and Schur--Weyl decomposition identifies
its multiplicity matrices with matrices between the corresponding fixed
symmetric-group modules. Since every source \(S_\lambda\mathbb C^4\) is
nonzero, injectivity for \(V=\mathbb C^4\) proves injectivity for every
dimension at least four. Inclusion of a smaller \(V\) in \(\mathbb C^4\)
and the naturality square proves it for dimensions below four.

Over \(\mathbb C\), a nonzero kernel of a \(GL_4\)-equivariant map
contains an irreducible summand and hence a nonzero highest-weight
vector. It therefore suffices to exhibit, for each source highest weight
\(\lambda\), exactly its source multiplicity many vectors with linearly
independent images.

\subsubsection{2. Exact source multiplicities, checked by two
constructions}\label{exact-source-multiplicities-checked-by-two-constructions}

Let \(W_b\) be the set of all quadruples of nonnegative integers summing
to \(b\). The coefficient \(M(\alpha)\) of the weight \(\alpha\) in the
source character is exactly the number of unordered multisets of four
elements of \(W_b\) having sum \(\alpha\). This follows directly by
choosing the monomial basis of \(\operatorname{Sym}^b\mathbb C^4\) and
then its degree-four monomial basis. With \(\rho=(3,2,1,0)\), the
multiplicity of \(S_\lambda\mathbb C^4\) equals

\[
 m_\lambda=\sum_{\sigma\in S_4}\operatorname{sgn}(\sigma)
 M(\lambda+\rho-\sigma\rho).
\]

To derive the formula, multiply the Schur-character expansion by
\(\prod_{i<j}(1-x_j/x_i)\). The Weyl alternant formula turns \(s_\mu\)
into the signed orbit of \(\mu+\rho\), translated by \(-\rho\). The
coefficient at a dominant \(\lambda\) is \(1\) for \(\mu=\lambda\) and
\(0\) otherwise because \(\mu+\rho\) is strictly decreasing and has a
unique dominant permutation.

The producer \texttt{\small polarization\_\allowbreak{}candidates.py}
obtains \(M\) from Newton's exact recurrence
\(h_i=i^{-1}\sum_{k=1}^i h_{i-k}p_k\), where
\(p_k=\sum_{\alpha\in W_b}x^{k\alpha}\). The C++ verifier independently
enumerates all unordered four-multisets of \(W_b\), with no use of
Newton's recurrence, and recomputes every required \(m_\lambda\). It
checks equality with the supplied table in each block. The producer also
checks

\[
 \sum_\lambda m_\lambda\prod_{1\le i<j\le4}
 \frac{\lambda_i-\lambda_j+j-i}{j-i}
 =\binom{\binom{b+3}{3}+3}{4}.
\]

All partitions of \(4b\) of length at most four are enumerated, so zero
source multiplicities are included in the extraction process. The
resulting positive-multiplicity lists and dimensions are those in the
table above.

\subsubsection{3. Explicit integral highest-weight vectors in the
source}\label{explicit-integral-highest-weight-vectors-in-the-source}

Use a \(4\times4\) variable matrix \(X\), whose columns are the slots
\(v_i\). For a nonempty subset \(I\subseteq\{1,2,3,4\}\), let \(d_I(X)\)
be the determinant of the submatrix with the first \(|I|\) rows and the
columns \(I\), in increasing order. There are fifteen such flag minors.
A product

\[
 F_c(X)=\prod_{I\ne\varnothing}d_I(X)^{c_I}
\]

has highest weight \(\lambda_i=\sum_{|I|\ge i}c_I\), and its degree in
slot \(j\) is \(\sum_{I\ni j}c_I\). The highest-weight assertion follows
explicitly by applying a raising operator: replacing a row indexed
\(j\le |I|\) by an earlier row creates two equal rows in its
determinant; a row \(j>|I|\) is absent. Thus every raising operator
annihilates every factor and hence the product.

For each tableau of shape \(\lambda\) with content \((b,b,b,b)\), its
columns give such a product with slot degrees \(b\). The producer sums
its twenty-four slot permutations, including the signs caused by sorting
the columns within each determinant. Dividing all integral coefficients
by their common divisor gives a nonzero integral symmetric
highest-weight vector, or the zero sum is discarded. It removes repeated
resulting coefficient lists. The input files record each remaining
vector as an explicit integer linear combination of the fifteen-minor
monomials. No claim that these lists form a basis is needed: exhibiting
\(m_\lambda\) independent images already suffices.

The independent checker verifies the weight and all four degrees of
every recorded monomial. A separate independent adjacent-transposition
check verifies that each selected vector is fixed by the three
generators of \(S_4\); it uses the explicit determinant signs rather
than the producer's permutation table. It passed on all 12 and 319
selected vectors respectively.

\subsubsection{4. Exact image evaluation and a nonzero integer
minor}\label{exact-image-evaluation-and-a-nonzero-integer-minor}

Work modulo the prime \(P=1,000,003\). Each certificate includes its
integral evaluation points \(X,x\), its selected candidate indices, the
selected evaluation rows, the entire square evaluation matrix, and its
nonzero determinant modulo \(P\). The points were chosen by a fixed
deterministic generator; the proof depends only on the printed points
and their nonzero determinants.

For every flag minor,

\[
 d_I(v_1+t_1x,\ldots,v_4+t_4x)
 =d_I(X)+\sum_{j\in I}t_jd_{I,j}(X,x),
\]

where \(d_{I,j}\) replaces column \(j\) by \(x\). Terms involving two
\(t\)'s vanish because two columns are then equal to \(x\). Multiplying
these explicitly linear polynomials, truncating each \(t_i\)-exponent
above \(L\), and reading coefficient \((L,L,L,L)\) therefore evaluates
exactly \(T_{L,b}(F_c)\). Truncation cannot affect that coefficient
because all exponents are nonnegative. The calculation then adds the
candidate's signed integral terms.

The C++ program finds a nonsingular \(m_\lambda\times m_\lambda\)
submatrix in every block, prints all entries and verifies its
determinant. A matrix of integral polynomial evaluations with nonzero
determinant modulo a prime has a nonzero integer determinant. Hence the
corresponding image polynomials are linearly independent over
\(\mathbb Q\) and over \(\mathbb C\). All source multiplicity spaces are
thus injective, proving R-011.

\subsubsection{5. Independent entry-by-entry
recheck}\label{independent-entry-by-entry-recheck}

\texttt{\small recheck\_\allowbreak{}polarization\_\allowbreak{}certificates.py}
never calls the C++ evaluator. It computes the small flag determinants
by recursive Laplace expansion, and computes Taylor coefficients using
four-dimensional NumPy arrays and axis shifts. It recomputes every entry
of every selected square matrix, compares with the certificate, and
recomputes each selected determinant by Laplace expansion. All selected
determinants have size at most six. Integer intermediates in each array
step are bounded by \(5(P-1)^2<2^{63}\), so the \texttt{\small int64}
operations are exact before reduction.

The independent receipt is:

\begin{verbatim}
POLARIZATION_R4_L1_B3: PASS; 12 blocks; 12 independently recomputed entries; 95 Taylor evaluations
POLARIZATION_R4_L2_B7: PASS; 173 blocks; 787 independently recomputed entries; 13865 Taylor evaluations
\end{verbatim}

The prime was independently checked by exact primality testing. The
calculations use no probabilistic inference: randomness only selected
points, and every determinant value is an exact certificate.

\subsection{Reproduction and artifact
inventory}\label{reproduction-and-artifact-inventory}

Run from this directory:

\begin{verbatim}
python polarization_candidates.py
g++ -O3 -std=c++17 check_polarization_ranks.cpp -o check_polarization_ranks
./check_polarization_ranks POLARIZATION_R4_L1_B3.input
./check_polarization_ranks POLARIZATION_R4_L2_B7.input
python recheck_polarization_certificates.py
\end{verbatim}

The two \texttt{\small .tsv} tables list the exact source and target
multiplicities and candidate counts. The two \texttt{\small .input}
files specify all candidate polynomials. The
\texttt{\small \_\allowbreak{}RANKS.tsv} files list every achieved rank
and determinant. The \texttt{\small \_\allowbreak{}CERTIFICATES.txt}
files give all points, selected candidates and square matrices.
\texttt{\small POLARIZATION\_\allowbreak{}INDEPENDENT\_\allowbreak{}RECHECK.txt}
records the independent entry-by-entry checks. These files, the source
code and this proof together provide a finite certificate of R-011.

The source of the motivating general polarization bound and determinant
obstruction is Sra, \emph{The q-Foulkes problem III: On Foulkes'
conjecture with a graded inner factor}, Theorem 1.3 and supplement S8,
p.~36, audited in
\texttt{\small boundary\_\allowbreak{}and\_\allowbreak{}prime\_\allowbreak{}blocks.md},
R-010. R-011 is a finite computation of these explicit maps; it does not
attribute its rank calculation to that source or claim a literature
priority search.

\section{The actual four-slot boundary and its coordinate-ring
morphism}\label{the-actual-four-slot-boundary-and-its-coordinate-ring-morphism}

Statements \textbf{R-012} and \textbf{R-013}. Completed and
independently checked on 2026-11-20 UTC by GPT-6 Astra Pro, agent
/root/rectangular\_cases; inherited effort setting not exposed.

\subsection{R-012. The finite boundary
map}\label{r-012.-the-finite-boundary-map}

For every finite-dimensional complex vector space \(V\), the natural map

\[
 T_{4,15}:\operatorname{Sym}^4(\operatorname{Sym}^{15}V)
 \longrightarrow
 \operatorname{Sym}^{16}V\otimes
 \operatorname{Sym}^4(\operatorname{Sym}^{11}V),
 \qquad
 f\longmapsto[t_1^4t_2^4t_3^4t_4^4]\,
 f(v_1+t_1x,\ldots,v_4+t_4x)
\]

is injective. This is the actual common-power boundary
\(r=L=4,\ b=r^2-1=15\). Multiplication by \((4!)^4\) gives the undivided
differential operator \(\prod_i(x\cdot\partial_{v_i})^4\), with the same
kernel.

\subsubsection{Complete proof}\label{complete-proof}

The dimension reduction and flag-minor highest-weight construction are
proved in full in R-011, sections 1--4 of
polarization\_boundary\_report.md. They apply unchanged here: the source
has only partitions of length at most four, the tensor permutation map
has dimension-independent multiplicity matrices, and a nonzero
equivariant kernel would contain a highest-weight vector. The following
exhaustive finite certificate proves that every source multiplicity
space maps injectively.

The exact source character has \textbf{1,713} distinct constituents,
highest-weight multiplicity sum \textbf{14,018}, maximum multiplicity
\textbf{35}, and full representation dimension on \(\mathbb C^4\)

\[
 \dim\operatorname{Sym}^4(\operatorname{Sym}^{15}\mathbb C^4)
 =\binom{819}{4}=18,609,648,876.
\]

For each constituent \(\lambda\) of multiplicity \(m_\lambda\), the
certificate provides \(m_\lambda\) integral symmetric highest-weight
polynomials, specified as signed sums of products of the fifteen flag
minors. Their row weights and equal column degrees are verified
directly. Invariance under the three adjacent slot transpositions is
verified directly, including all determinant signs. The certificate then
gives \(m_\lambda\) evaluation points and the full square evaluation
matrix of their images under \(T_{4,15}\).

Every one of those 1,713 square matrices has a nonzero determinant
modulo the prime \(1,000,003\). In particular, every determinant is a
nonzero integer. Its columns are therefore linearly independent over
\(\mathbb Q\), and hence over \(\mathbb C\). Thus the map is injective
on every highest-weight multiplicity space. The highest-weight argument
proves injectivity for \(\mathbb C^4\), and the dimension-independent
tensor permutation argument proves it for all \(V\).

The certificate is exhaustive in the following independently verified
sense.

\begin{enumerate}
\def\labelenumi{\arabic{enumi}.}
\item
  The source multiplicities were first extracted from the exact Newton
  recurrence for \(h_4[h_{15}]\). A separate C++ program,
  verify\_boundary\_character.cpp, independently computes

  \[
  [t^4]\prod_{\substack{\alpha\in\mathbb Z_{\ge0}^4\\|\alpha|=15}}
  (1-tx^\alpha)^{-1}
  \]

  by multiplying the 816 individual geometric-series factors. It applies
  the Weyl alternant to every partition of 60 of length at most four,
  and obtains exactly the same positive multiplicities. Its
  Weyl-dimension sum and sum of all monomial weight coefficients both
  equal 18,609,648,876.
\item
  The original evaluator expands each shifted flag minor as its constant
  term plus its four possible linear terms, then uses forward
  coefficient propagation and modular Gaussian elimination.
\item
  A separately written evaluator, recheck\_large\_polarization.cpp,
  recomputes \textbf{all 203,840 entries} of the selected square
  matrices. It processes the flag-minor factors in reverse order, uses
  descending in-place coefficient recurrence instead of forward
  propagation, and independently computes every small determinant and
  selected matrix determinant by power traces and Newton identities. It
  also checks every selected polynomial's highest weight, four slot
  degrees, and exact invariance under the three generators of \(S_4\).
\end{enumerate}

For clarity, the independent coefficient recurrence replaces a current
coefficient array \(c_\alpha\), in descending index order, by

\[
 c_\alpha\ \leftarrow\
 d\,c_\alpha+\sum_{j=1}^4a_jc_{\alpha-e_j}
 \pmod {1,000,003}
\]

for each linear minor factor \(d+\sum_j a_jt_j\), retaining only
\(0\le\alpha_j\le4\). Descending order guarantees that every referenced
lower index still has its old value. Truncation cannot change the target
coefficient because exponents are nonnegative. The determinant recheck
uses

\[
 e_0=1,\qquad
 e_k=\frac1k\sum_{i=1}^k(-1)^{i-1}e_{k-i}\operatorname{tr}(A^i),
 \qquad\det A=e_m,
\]

for an \(m\times m\) matrix. Here \(m\le35<1,000,003\), so every
division by \(k\) is valid. All products and unreduced sums fit signed
64-bit integers: the largest matrix-product sums have at most 35 terms,
each smaller than \(1,000,003^2\).

The independent run made 4,517,197 Taylor evaluations, checked all
14,018 selected vectors, and passed every entry and determinant. The
underlying finite-field modulus is checked by trial division.
Consequently no probabilistic assertion enters the proof; the point
generator merely supplied the printed points.

\subsubsection{Compact certificate}\label{compact-certificate}

The file POLARIZATION\_R4\_L4\_B15\_SELECTED.compact.gz contains the
complete selected source polynomials, points, matrices and determinants.
Its uncompressed size is 11,569,652 bytes; compressed size is 2,761,977
bytes. Its SHA256 is

\begin{verbatim}
fa2bbbeb5e335d08cda726359ac2b1ffaf1fe1270bba56e3f29841b9957db107
\end{verbatim}

The character table is POLARIZATION\_R4\_L4\_B15.tsv. The rank summary
is POLARIZATION\_R4\_L4\_B15\_ALL\_RANKS.tsv. The independent receipt is
POLARIZATION\_R4\_L4\_B15\_INDEPENDENT\_RECHECK.txt. The complete
full-candidate input and parallel batches are intermediate files; they
are unnecessary for checking the selected certificate.

To recheck the selected certificate, decompress the gzip file, compile
recheck\_large\_polarization.cpp with C++17 and optimization, and pass
the resulting compact file as its sole argument. Compile
verify\_boundary\_character.cpp and pass POLARIZATION\_R4\_L4\_B15.tsv
to recheck the source multiplicities independently. The compressed
certificate contains only selected vectors, so the 124 MB full candidate
input is not needed.

No priority claim is made for this finite rank computation. The supplied
proof establishes this map and these parameters, without a
classification for larger even numbers of slots.

\subsection{R-013. Exact morphism to the variety of padded
forms}\label{r-013.-exact-morphism-to-the-variety-of-padded-forms}

The coordinate-ring description is a known generalized Foulkes--Howe
construction. A precise primary locator is Kadish--Landsberg,
\emph{Padded polynomials, their cousins, and geometric complexity
theory}, \textbf{Theorem 1.7}, p.~3, with the full tensor construction
and evaluation proof in \textbf{section 3, pp.~5--6}:

\url{https://arxiv.org/pdf/1204.4693}

The published article is \emph{Communications in Algebra} 42 (2014),
2171--2180, DOI 10.1080/00927872.2012.758268. The statement and proof in
the cited arXiv version were read in this run. The following proof fixes
the exact scalar relating that coordinate-ring map to the
coefficient-extraction operator used in our certificates.

Let \(r\ge1\), \(b\ge1\), and \(0\le L\le b\). Fix the normalized
perfect pairing on symmetric powers by

\[
 \langle v^k,\ell_1\cdots\ell_k\rangle
 =\prod_{j=1}^k\ell_j(v).
\]

Define the polynomial morphism

\[
 \nu_{L,b}:V^*\times\operatorname{Sym}^{b-L}V^*
 \longrightarrow\operatorname{Sym}^bV^*,
 \qquad(\ell,q)\longmapsto\ell^Lq.
\]

Its pullback in coordinate degree \(r\) is

\[
 (\nu_{L,b}^*)_r:
 \operatorname{Sym}^r(\operatorname{Sym}^bV)
 \longrightarrow
 \operatorname{Sym}^{rL}V\otimes
 \operatorname{Sym}^r(\operatorname{Sym}^{b-L}V).
\]

Define \(T_{r,L,b}\) by extracting \(t_1^L\cdots t_r^L\) after
substituting \(v_i+t_ix\) in the symmetric \(r\)-slot model. Then the
exact equality is

\[
 \boxed{\displaystyle
 (\nu_{L,b}^*)_r=\binom bL^{-r}T_{r,L,b}.}
 \tag{R013.1}
\]

\textbf{Proof.} Products \(w_1^b\cdots w_r^b\) span
\(\operatorname{Sym}^r(\operatorname{Sym}^bV)\), because pure powers
span each symmetric power. As a coordinate polynomial, such a product
sends \(P\) to \(\prod_i\langle w_i^b,P\rangle\). Its pullback along
\(\nu_{L,b}\) is

\[
 \prod_i\ell(w_i)^L
 \prod_i\langle w_i^{b-L},q\rangle.
\]

This is represented in the displayed target by

\[
 (w_1^L\cdots w_r^L)\otimes
 (w_1^{b-L}\cdots w_r^{b-L}).
\]

The symmetric-slot representative of the original product is the average
over all assignments of the \(w_i\) to the \(r\) slots. In each assigned
row the binomial theorem gives the factor

\[
 [t_i^L]\langle w,v_i+t_ix\rangle^b
 =\binom bL\langle w,x\rangle^L
 \langle w,v_i\rangle^{b-L}.
\]

Multiplying the \(r\) factors and taking the same slot average therefore
gives \(\binom bL^r\) times the tensor above. Equality on a spanning set
proves (R013.1).

Let \(Z_{L,b}(V)\) be the Zariski closure of the image of \(\nu_{L,b}\).
A homogeneous polynomial vanishes on that image precisely when its
pullback is zero; vanishing on the image and its closure are equivalent
by the definition of Zariski closure. Thus

\[
 \boxed{\displaystyle
 I_r(Z_{L,b}(V))
 =\ker(\nu_{L,b}^*)_r
 =\ker T_{r,L,b}.}
 \tag{R013.2}
\]

Moreover the maps \(\binom bL^{-r}T_{r,L,b}\), as \(r\) varies, are the
graded pieces of the one algebra homomorphism \(\nu_{L,b}^*\). This
identifies their multiplication exactly, including the normalization.

\subsubsection{Actual geometric
corollaries}\label{actual-geometric-corollaries}

Combining R-011 and R-012 with (R013.2) gives, for every
finite-dimensional complex \(V\),

\[
 I_4(Z_{1,3}(V))=
 I_4(Z_{2,7}(V))=
 I_4(Z_{4,15}(V))=0.
\]

The homogeneous ideals in all degrees \(0\le d\le4\) are zero as well.
Indeed, a nonzero equation of smaller positive degree could be
multiplied by a nonzero polynomial to degree four in the ambient
polynomial coordinate ring, which is an integral domain; the product
would contradict the displayed vanishing. Degree zero contains no
equation because each image contains the origin. When the ambient vector
space is zero-dimensional, the assertion follows directly from its
coordinate ring.

These are explicit low-degree assertions for the varieties of forms
\(\ell q\) of degree 3, \(\ell^2q\) of degree 7, and \(\ell^4q\) of
degree 15. They give a geometric interpretation of the verified finite
kernels, with the complete morphism and scalar proved above.

\section{Source-relative endpoint proof
check}\label{source-relative-endpoint-proof-check}

ID T-015; run GCT\_PLETH-20261121-theta. This observation is included as
a checked comparison with a source, not as a new best global bound.

The full proof in
\texttt{\small work/\allowbreak{}external/\allowbreak{}quadratic\_\allowbreak{}stabilization\_\allowbreak{}2026.tex},
sections 2--4, lines 178--436, was read, including its
diagonal-annihilator criterion, Fischer inner-product lemma, and
four-step induction.

For \texttt{\small a>=\allowbreak{}3}, set
\texttt{\small r=\allowbreak{}a-\allowbreak{}1} and
\texttt{\small b=\allowbreak{}a(\allowbreak{}a-\allowbreak{}1)\allowbreak{}-\allowbreak{}1=\allowbreak{}r(\allowbreak{}r+\allowbreak{}1)\allowbreak{}-\allowbreak{}1}.
The residual symmetric degree is
\texttt{\small m=\allowbreak{}b-\allowbreak{}r=\allowbreak{}r\textasciicircum{}2-\allowbreak{}1}.
Source Step 2 still has
\texttt{\small p-\allowbreak{}q=\allowbreak{}m-\allowbreak{}i+\allowbreak{}1>=\allowbreak{}1}
for every derivative insertion. In Step 3,
\texttt{\small m=\allowbreak{}r\textasciicircum{}2-\allowbreak{}1>=\allowbreak{}r(\allowbreak{}r-\allowbreak{}1)\allowbreak{}}
and \texttt{\small m>=\allowbreak{}1}, so the source's original
Proposition 3 for \texttt{\small r} slots applies; no stronger induction
claim is assumed. In Step 4 the shared destination degree is at most
\texttt{\small r\textasciicircum{}2-\allowbreak{}1}, while the source
degree before each derivative is at least
\texttt{\small b-\allowbreak{}r+\allowbreak{}1=\allowbreak{}m+\allowbreak{}1=\allowbreak{}r\textasciicircum{}2},
so the injectivity inequality remains strict. The final multilinear
spanning step is unchanged.

Thus the source proof also gives surjectivity of
\texttt{\small h\_\allowbreak{}\{b,\allowbreak{}a\}} at the one-lower
endpoint; its original theorem covers all larger \texttt{\small b}. At
\texttt{\small a=\allowbreak{}2,\allowbreak{}b=\allowbreak{}1} the map
is the identity under the canonical symmetric-power identifications. The
parent agent found a stronger recent bound and is separately checking
it; this event does not certify that other paper.

\section{Prime-adjacent square blocks, specialization, and exact
canonical-image
semigroups}\label{prime-adjacent-square-blocks-specialization-and-exact-canonical-image-semigroups}

\subsection{G-013. Explicit determinant and padded-rectangle image
ranges}\label{g-013.-explicit-determinant-and-padded-rectangle-image-ranges}

Work over \(\mathbf Q\), with the normalized symmetric-tensor and
canonical-map conventions of R-001, R-002, and R-008. For an integer
\(a\ge2\), let \[
\mathcal B_a=\{b\ge a: b\text{ is even and either }b-1\text{ or }b+1\text{ is an odd prime}\},
\qquad
\mathcal S_a=\left\{\sum_{b\in\mathcal B_a}n_b b:n_b\in\mathbf N,
\text{ finitely many }n_b\ne0\right\}.
\tag{P.1}
\] Every \(B\in\mathcal S_a\) has a rational, explicitly constructed
preimage of the determinant tensor \(I_{a,B}\) under the actual
canonical map \[
 h_{B,a}:\operatorname{Sym}^B(\operatorname{Sym}^a W)
 \longrightarrow\operatorname{Sym}^a(\operatorname{Sym}^B W),\qquad \dim W=a.
\tag{P.2}
\] In particular its rank on the target \((B^a)\) determinant component
is one. For every integer \(s\ge0\), multiplication of the constructed
preimage by \((e_1^a)^s\) gives a nonzero image of highest weight \[
 (as+B,B^{a-1})
\quad\text{under}\quad h_{B+s,a}.
\tag{P.3}
\] The target multiplicity of this weight is exactly one. These
component statements hold in every ambient dimension at least \(a\), by
the displayed subspace inclusion and highest-weight vector.

The set \(\mathcal S_a\) is cofinite among the even integers. More
explicitly, for every odd prime \(p>a\), \[
 B\text{ even},\qquad B\ge \frac{(p-1)(p-3)}2
 \quad\Longrightarrow\quad B\in\mathcal S_a.
\tag{P.4}
\] This is an established range with an explicit preimage; no
distributional assertion about primes is used in its proof. An ordinary
prime search supplies a suitable \(p\).

For the ranks \(2\le a\le100\), the attached finite Apéry certificates
determine the entire semigroup \(\mathcal S_a\), using only the widths
\(a\le b\le4a\) in (P.1). The exact even conductor \(C_a\) means that
every even \(B\ge C_a\) belongs to \(\mathcal S_a\), and, for \(C_a>0\),
the even integer \(C_a-2\) does not. It describes precisely this
constructed seed semigroup; a degree outside it may still have a nonzero
canonical determinant block by another construction.

\begin{longtable}[]{@{}
  >{\raggedleft\arraybackslash}p{(\columnwidth - 6\tabcolsep) * \real{0.2500}}
  >{\raggedleft\arraybackslash}p{(\columnwidth - 6\tabcolsep) * \real{0.2500}}
  >{\raggedleft\arraybackslash}p{(\columnwidth - 6\tabcolsep) * \real{0.2500}}
  >{\raggedleft\arraybackslash}p{(\columnwidth - 6\tabcolsep) * \real{0.2500}}@{}}
\toprule\noalign{}
\begin{minipage}[b]{\linewidth}\raggedleft
Rank \(a\)
\end{minipage} & \begin{minipage}[b]{\linewidth}\raggedleft
Exact even conductor \(C_a\), both prime-adjacent seeds
\end{minipage} & \begin{minipage}[b]{\linewidth}\raggedleft
Prime-minus-one seeds alone
\end{minipage} & \begin{minipage}[b]{\linewidth}\raggedleft
Known full-map surjectivity bound \((a-1)^2\)
\end{minipage} \\
\midrule\noalign{}
\endhead
\bottomrule\noalign{}
\endlastfoot
2 & 0 & 0 & 1 \\
3 & 4 & 4 & 4 \\
4 & 4 & 4 & 9 \\
5 & 6 & 16 & 16 \\
6 & 6 & 16 & 25 \\
7 & 8 & 16 & 36 \\
8 & 8 & 16 & 49 \\
9 & 10 & 16 & 64 \\
10 & 10 & 16 & 81 \\
12 & 12 & 28 & 121 \\
16 & 28 & 28 & 225 \\
20 & 36 & 56 & 361 \\
26 & 52 & 64 & 625 \\
32 & 58 & 92 & 961 \\
50 & 96 & 116 & 2401 \\
64 & 126 & 148 & 3969 \\
100 & 190 & 196 & 9801 \\
\end{longtable}

The last column is Sra III, Theorem 1.3, with its complete proof read in
this run. It concerns every target component and is included only with
that larger scope. The first columns certify the specified determinant
and row-extended determinant components.

\subsubsection{1. The exact square-to-smaller-rank
pullback}\label{the-exact-square-to-smaller-rank-pullback}

Let \(b\) be one of the even square orders with \(c_b\ne0\), where R-002
defines the normalized scalar by \[
 h_{b,b}(I_{b,b})=c_b I_{b,b}.
\tag{P.5}
\] Take \(2\le a\le b\). Choose a decomposition \[
 U^*=W^*\oplus\langle z_{a+1},\ldots,z_b\rangle,\qquad \dim U=b,
\] and the actual linear multiplication map \[
 j:\operatorname{Sym}^a W^*\longrightarrow\operatorname{Sym}^b U^*,
 \qquad F\longmapsto Fz_{a+1}\cdots z_b.
\tag{P.6}
\] Let \(\Pi_a:(W^*)^a\to\operatorname{Sym}^a W^*\) and
\(\Pi_b:(U^*)^b\to\operatorname{Sym}^b U^*\) multiply their ordered
linear forms. The affine map \(\iota\) appends the fixed forms
\(z_{a+1},\ldots,z_b\). Then \[
 j\circ\Pi_a=\Pi_b\circ\iota.
\tag{P.7}
\] Taking the degree-\(b\) coordinate pullbacks gives the commutative
equality \[
 h_{b,a}\circ\operatorname{Sym}^b(j^*)
 =\iota^*\circ h_{b,b}.
\tag{P.8}
\] Here \(j^*\) is the dual linear map. The usual perfect
characteristic-zero symmetric pairings identify it with a map
\(\operatorname{Sym}^b U\to\operatorname{Sym}^a W\); its coefficients
are determined by this duality. Thus all binomial factors of a monomial
tensor basis are included by the defined map, without an implicit
rescaling.

The normalized tensor \(I_{b,b}\) represents the polynomial
\(\det_b(\ell_1,\ldots,\ell_b)^b\) on the ordered linear forms.
Substituting the appended basis forms gives a block matrix with an
identity complement, so \[
 \iota^*(I_{b,b})=I_{a,b}.
\tag{P.9}
\] Consequently the explicit rational vector \[
 u_{a,b}=\operatorname{Sym}^b(j^*)\left(\frac{I_{b,b}}{c_b}\right)
 \quad\text{satisfies}\quad h_{b,a}(u_{a,b})=I_{a,b}.
\tag{P.10}
\] The map \(j\) commutes with the block inclusion \(GL(W)\to GL(U)\);
therefore \(u_{a,b}\) has determinant character \(\det_W^b\). Equation
(P.10) also proves its nonvanishing directly.

For \(b=p-1\), Glynn's Theorem 3.2, and for \(b=p+1\), Drisko's Theorem
9, give \(c_b\ne0\), with the exact rational normalizations and
prime-adic valuations proved in R-007. Applying (P.10) proves every seed
in \(\mathcal B_a\). R-008 gives an additional, sharper congruence and
an alternative source vector for the minus-one seeds. We use these known
square results with their exact locators and make no priority claim for
this specialization consequence.

\subsubsection{2. Multiplication gives the specified
preimages}\label{multiplication-gives-the-specified-preimages}

R-008's map \[
 \mu_a:\operatorname{Sym}(\operatorname{Sym}^a W)
 \longrightarrow \bigoplus_{d\ge0}\operatorname{Sym}^a(\operatorname{Sym}^d W)
\] is a graded algebra morphism; its degree-\(d\) component is
\(h_{d,a}\). Its target multiplication is componentwise multiplication
in the \(a\) tensor factors. The determinant tensors satisfy \[
 I_{a,b}I_{a,c}=I_{a,b+c}.
\tag{P.11}
\] Evaluation on arbitrary \(a\) covectors proves this equality by the
product \(\det^b\det^c=\det^{b+c}\), and these evaluations separate the
symmetric tensors. Hence, for any specified expression \(B=\sum n_b b\),
\[
 u_{a,B}:=\prod_b u_{a,b}^{n_b}
 \quad\text{satisfies}\quad h_{B,a}(u_{a,B})=I_{a,B}.
\tag{P.12}
\] No linear independence of different seed expressions is claimed or
needed.

For (P.3), multiply (P.12) by \((e_1^a)^s\). Its image is \[
 I_{a,B}\,(e_1^s)^{\otimes a},
\tag{P.13}
\] a nonzero polynomial on the ordered covectors. It is a highest-weight
vector of weight \((as+B,B^{a-1})\): the first factor has determinant
character and the second is the power of the highest line. The refined
balanced-weight realization T-002 identifies its target multiplicity
with the invariant part of \[
 \left(\det^B\otimes\operatorname{Sym}^{as}\mathbf Q^a\right)_{((B+s)^a)}.
\] After removing the determinant factor, the only balanced monomial is
\((x_1\cdots x_a)^s\), and its permutation character is
\(\operatorname{sgn}^B=1\). Thus this multiplicity is exactly one.
Appending zero weight coordinates under \(W\hookrightarrow V\) retains
the same highest-weight vector, proving the ambient-dimension statement.

\subsubsection{3. Cofiniteness and the explicit uniform
range}\label{cofiniteness-and-the-explicit-uniform-range}

For an odd prime \(p>a\), put \(g=(p-1)/2\). Both \(2g\) and \(2(g+1)\)
lie in \(\mathcal B_a\). For each nonnegative \(N\), write
\(r=N\bmod g\), \(0\le r<g\) (for \(g=1\), take \(r=0\)). The least
nonnegative member of the semigroup \(\langle g,g+1\rangle\) with
residue \(r\) is \[
 w_r=r(g+1).
\tag{P.14}
\] Indeed, in any expression \(ug+v(g+1)\), the second coefficient
satisfies \(v\equiv r\pmod g\), and its smallest possible value is
\(r\). Conversely \(w_r\) is represented. Every \(N\ge g(g-1)\)
satisfies \(N\ge w_r\): the least integer at least \(g(g-1)\) in that
residue is \(g(g-1)+r\), whose difference from \(w_r\) is
\(g(g-1-r)\ge0\). Therefore \[
 N=\frac{N-r(g+1)}g\,g+r(g+1)
\tag{P.15}
\] has nonnegative integral coefficients. Multiplying by two gives
(P.4), as well as an explicit two-factor construction in (P.12). The
existence of arbitrarily large primes follows from Euclid's
finite-product argument, so choosing \(p>a\) requires no unproved prime
statement.

\subsubsection{4. Exact finite semigroups by a residue
graph}\label{exact-finite-semigroups-by-a-residue-graph}

For a finite seed list \(2g_1,\ldots,2g_k\) with
\(\gcd(g_1,\ldots,g_k)=1\), let \(g=\min g_i\). Make a directed graph on
\(\mathbf Z/g\mathbf Z\), with an edge of weight \(g_i\) from \(r\) to
\(r+g_i\). The residues generated by the \(g_i\) form the whole cyclic
group: their subgroup has index the stated gcd, and the monoid they
generate in a finite group is a group, since a positive multiple
supplies each inverse. All edge weights are positive. Thus a shortest
path from zero exists to every residue, and its length \(w_r\) is
exactly the smallest nonnegative represented integer in that residue.
The path itself records a seed expression.

For every integer \(N\ge0\), the exact membership equivalence is \[
 N\in\langle g_1,\ldots,g_k\rangle
 \quad\Longleftrightarrow\quad N\ge w_{N\bmod g}.
\tag{P.16}
\] The forward implication is minimality; the reverse adds copies of the
generator \(g\). The greatest missing integer in residue \(r\) is
\(w_r-g\), so the normalized conductor is \[
 c=\max_r w_r-g+1.
\tag{P.17}
\] For \(g=1\), this formula gives \(c=0\). These equalities also prove
the shortest-path algorithm used in the attached code and explain
precisely what its certificates certify.

The script
\texttt{\small work/\allowbreak{}root/\allowbreak{}prime\_\allowbreak{}semigroup.py}
uses a direct sieve, enumerates every prime-adjacent width in
\([a,4a]\), computes the shortest paths and their explicit integer
witnesses, and then verifies all memberships through
\(\max w_r+g+\max g_i\) by a separate Boolean recurrence
\(D(N)=\bigvee_iD(N-g_i)\). It checks that all witnesses have the stated
weights and residues. For every \(2\le a\le100\), it additionally
verifies \(c\le2a\). Every omitted prime-adjacent width is greater than
\(4a\); its half is therefore greater than \(2a\), and hence already
belongs to the finite computed semigroup by (P.16). This proves that the
finite list generates the entire \(\mathcal S_a\), without searching
infinitely many primes. The same calculation and omission argument were
checked for the minus-one lists displayed for comparison.

All 198 finite seed-set calculations passed. The complete data are in
\texttt{\small PRIME\_\allowbreak{}SEMIGROUP\_\allowbreak{}RANGES.tsv}
and
\texttt{\small PRIME\_\allowbreak{}SEMIGROUP\_\allowbreak{}CERTIFICATES.json},
beside the reproducing script. The proof of the square specialization
was independently checked through the commutative multiplication diagram
by the rectangular agent. The whole proof and all computational checks
receive separate events in the ledger.

\subsubsection{Scheme and representation-ring
dictionary}\label{scheme-and-representation-ring-dictionary}

The maps \(\Pi_a,\Pi_b,j,\iota\) are morphisms of affine schemes defined
by integral coordinate polynomials. Equation (P.7) is equality of these
actual scheme morphisms, and (P.8) is its contravariant equality on
coordinate rings. The inverse scalar \(c_b^{-1}\) is taken over
\(\mathbf Q\); after inverting its numerator and the finitely many
symmetric-pairing denominators, the same preimage is defined over the
resulting localization of \(\mathbf Z\). This identifies the base
explicitly and makes no assertion of an integral inverse at an excluded
prime. The nonzero image line is the coefficient of the Schur operation
\(s_{(B^a)}\) in the image of the Chow coordinate morphism from G-003,
equivalently the corresponding isotypic rank in
\(K_0(\operatorname{Rep}_{\mathbf Q}GL_a)\). The tensor multiplication
is the product in that representation algebra and in the graded scheme
coordinate ring. The finite residue semigroup specifies a concrete
finitely generated subalgebra of the canonical image through (P.12); it
does not replace the entire coordinate ring by a numerical semigroup.
The general SGA 3, SGA 6, EGA I, and FGA locators for these exact
constructions are supplied in the dedicated dictionary part.

\section{An explicit canonical-map lift for every sufficiently long
first
row}\label{an-explicit-canonical-map-lift-for-every-sufficiently-long-first-row}

\subsection{G-014. Degree-tracked multisymmetric
lifting}\label{g-014.-degree-tracked-multisymmetric-lifting}

Let \(a\ge2\), \(n\ge0\), and let \(V\) be a finite-dimensional
characteristic-zero vector space. For every partition
\(\lambda\vdash an\) with \[
 d:=an-\lambda_1=\lambda_2+\lambda_3+\cdots\le n,
\tag{C.1}
\] the canonical map \[
 h_{n,a}:\operatorname{Sym}^n(\operatorname{Sym}^a V)
 \longrightarrow\operatorname{Sym}^a(\operatorname{Sym}^n V)
\tag{C.2}
\] is surjective on the entire target \(\lambda\)-isotypic component.
For the empty partition we put \(\lambda_1=0\); the zero-dimensional
ambient space is immediate. The proof gives a rational preimage for
every target highest-weight vector by explicit coefficient extraction,
Newton identities, and a finite sum over set partitions. Consequently
the image multiplicity equals \(a_{a,n}^{\lambda}\) throughout (C.1).
This includes all ambient dimensions: a Schur module whose partition
exceeds the dimension is zero on both sides of this component assertion.

The elementary-multisymmetric coordinate chart is classical:
Raicu--Sam--Weyman, \emph{On some modules supported in the Chow
variety}, Theorem 2.3, its proof (2.8), and Remark 2.9, printed pp.5--6,
give the chart and its characteristic-zero invariant generators. The
proof below tracks the ordinary number of generators as well as their
polynomial weights, to supply the stated finite-degree lift and the
exact canonical map. No priority claim is made for the resulting
first-row range.

\subsubsection{1. The actual coordinate
morphism}\label{the-actual-coordinate-morphism}

Choose a basis of \(V\). Write the \(i\)-th linear form as \[
 \ell_i=x_i t_1+\sum_{j=2}^{r}z_{ij}t_j,\qquad 1\le i\le a,
\] and put \(u=\prod_i x_i\). On the principal open set \(u\ne0\),
define \(w_{ij}=z_{ij}/x_i\). A target vector is a polynomial in the
\(a\) ordered forms, homogeneous of degree \(n\) in each and invariant
under permuting the forms. Any such polynomial of total degree \(d\) in
the variables \(z_{ij}\), \(j\ge2\), has the unique expression \[
 f=u^n F(w_{ij}),
\tag{C.3}
\] where \(F\) is an \(S_a\)-invariant polynomial, homogeneous of degree
\(d\). Indeed, in every monomial the exponent of \(x_i\) is
\(n-\sum_{j\ge2}\deg_{z_{ij}}\); substituting \(z_{ij}=x_iw_{ij}\) gives
exactly \(x_i^n\) in each row. Permutation invariance passes to \(F\)
because \(u\) is invariant.

For a vector of auxiliary variables \(y=(y_2,\ldots,y_r)\), set
\(L_i(y)=\sum_{j\ge2}w_{ij}y_j\) and \[
 \prod_{i=1}^a(1+sL_i(y))
 =\sum_{j=0}^a E_j(y)s^j,
 \qquad E_j(y)=\sum_{|\beta|=j}E_\beta(w)y^\beta.
\tag{C.4}
\] Let \(C_\beta\) denote the linear coordinate functional extracting
the coefficient of \(t_1^{a-|\beta|}\prod_{j\ge2}t_j^{\beta_j}\) from an
arbitrary degree-\(a\) form; let \(C_0\) extract \(t_1^a\). These are
actual elements of
\((\operatorname{Sym}^aV^*)^*\simeq\operatorname{Sym}^aV\), with the
dual symmetric pairing specifying any monomial-basis factorials.
Pullback under multiplication of the ordered forms gives \[
 \mu_a(C_0)=u,
 \qquad \mu_a(C_\beta)=uE_\beta(w).
\tag{C.5}
\] The homogeneous degree-\(n\) part of this algebra pullback is exactly
(C.2), with the conventions of G-003 and R-008.

\subsubsection{\texorpdfstring{2. An explicit expression with at most
\(d\) coefficient
factors}{2. An explicit expression with at most d coefficient factors}}\label{an-explicit-expression-with-at-most-d-coefficient-factors}

We prove that every degree-\(d\) invariant \(F\) is a rational linear
combination of monomials \[
 \prod_{q=1}^k E_{\beta^{(q)}},
 \qquad \sum_q|\beta^{(q)}|=d,
 \qquad k\le d,
\tag{C.6}
\] where all \(\beta^{(q)}\ne0\). The constant case \(d=0\) has \(k=0\).

For a nonzero exponent vector \(\alpha\), put
\(P_\alpha=\sum_iw_i^\alpha\). Every invariant polynomial is a linear
combination of monomial orbit sums. For a list of nonzero exponent
vectors \(\alpha_1,\ldots,\alpha_k\), the distinct-row sum has the exact
identity \[
 \sum_{i_1,\ldots,i_k\ \mathrm{distinct}}
 \prod_{j=1}^k w_{i_j}^{\alpha_j}
 =\sum_{\pi\in\Pi_k}
 \left(\prod_{B\in\pi}(-1)^{|B|-1}(|B|-1)!\right)
 \prod_{B\in\pi}P_{\sum_{j\in B}\alpha_j}.
\tag{C.7}
\] To check it, fix a map from the labels \(1,\ldots,k\) to row indices.
Its coefficient on the right is the sum of the displayed factors over
partitions refining its equality classes. On an equality class of size
\(s\), this is the signed sum of all permutations of \(s\) letters,
grouped by their cycles; it is one for \(s=1\) and zero for \(s\ge2\),
by pairing a permutation with its product with a fixed transposition.
Thus precisely the injective maps survive. Each ordinary orbit sum is
the left side divided by the factorials of the multiplicities of
repeated exponent vectors. This proves the asserted generation by the
\(P_\alpha\), while preserving the total exponent degree.

Now put \(P_j(y)=\sum_iL_i(y)^j\). The formal identity \[
 \frac{s\,\partial_s\prod_i(1+sL_i)}{\prod_i(1+sL_i)}
 =\sum_{j\ge1}(-1)^{j-1}P_j(y)s^j
\tag{C.8}
\] gives Newton's recurrence after multiplication by
\(\sum_{j=0}^a E_j(y)s^j\) and coefficient comparison. It expresses
\(P_j(y)\) as a polynomial in \(E_1(y),\ldots,E_a(y)\), homogeneous of
weighted degree \(j\) when \(E_i\) has weight \(i\). Its coefficient of
\(y^\alpha\), for \(|\alpha|=j\), is \(\binom{j}{\alpha}P_\alpha\).
Dividing by that nonzero integer therefore expresses \(P_\alpha\) in the
\(E_\beta\), preserving weighted degree \(|\alpha|\). Every factor has
positive weight, so every resulting monomial contains at most
\(|\alpha|\) factors. Substitute these expressions in (C.7). The total
weight of every term is \(d\), and its number of factors is at most
\(d\). This proves (C.6) constructively.

\subsubsection{3. Homogenization produces a preimage of the required
degree}\label{homogenization-produces-a-preimage-of-the-required-degree}

For each term of (C.6), the assumption \(n\ge d\) gives \(n-k\ge0\).
Equation (C.5) then gives \[
 \mu_a\left(C_0^{n-k}\prod_{q=1}^k C_{\beta^{(q)}}\right)
 =u^n\prod_{q=1}^k E_{\beta^{(q)}}.
\tag{C.9}
\] Apply (C.9) term by term to the explicitly computed expression for
\(F\). The resulting source polynomial is homogeneous of coordinate
degree \(n\) and maps to \(f\). The equality was derived on \(u\ne0\);
both sides are polynomials in the ordered linear forms, and this is a
dense principal open of their affine space. Hence they agree everywhere.

Every target highest-weight vector of weight \(\lambda\) has
first-coordinate degree \(\lambda_1\), so its total degree in the
remaining coordinates is precisely \(d=an-\lambda_1\). Thus (C.3)--(C.9)
lift every such vector under (C.1). The canonical map is
\(GL(V)\)-equivariant. Its image therefore contains the submodules
generated by all target highest-weight vectors of weight \(\lambda\),
which together are the full \(\lambda\)-isotypic component. This proves
G-014.

The argument gives image multiplicity rather than only a coefficient
inequality. A preimage need not initially be highest weight;
equivariance and the equality with the specified target vector suffice,
and the usual characteristic-zero isotypic projection provides a
preimage in the same irreducible block.

\subsubsection{Exact isomorphism on the common first-row
range}\label{exact-isomorphism-on-the-common-first-row-range}

For \(a,n\ge2\) and \(d\le\min(a,n)\), apply G-014 both to \(h_{n,a}\)
and to \(h_{a,n}\). In the common tensor space, normalized row and
column symmetrizers are self-adjoint orthogonal projections, and the
grid transpose identifies the two maps as adjoints, exactly as in M-001
and G-002. Their ranks are equal. Surjectivity in both directions
therefore makes the two component dimensions equal to that common rank.
Thus both maps are isomorphisms on their \(\lambda\)-isotypic components
in the displayed common range. This follows for the prescribed maps
themselves, with their specified tensor identifications.

\subsubsection{Geometric and lambda-ring
form}\label{geometric-and-lambda-ring-form}

The multiplication map on ordered linear forms is a morphism of affine
schemes, and (C.5) is its coordinate pullback. Localization at the
coefficient of the pure power \(t_1^a\) produces the displayed
multisymmetric chart. Equations (C.7)--(C.9) prove that its denominators
clear in coordinate degree \(n\) for every component satisfying (C.1).
Thus the Chow coordinate image from G-003 and its normalizing
representation algebra have identical \(\lambda\)-isotypic pieces in
this range. Applying the coefficient functional \([s_\lambda]\) in
\(K_0(\operatorname{Rep}_{\mathbf Q}GL(V))\) gives precisely the rank
equality asserted above. All coordinate equations are integral; the
rational divisions in (C.7)--(C.8) specify a localization of
\(\mathbf Z\) over which each particular preimage spreads out. The
general integral symmetric/divided-power map remains the one in G-003,
and no integral surjectivity at excluded primes is asserted.

Primary source: Raicu, Sam, Weyman, \emph{On some modules supported in
the Chow variety}, Theorem 2.3, equation (2.8), Remark 2.9, pp.5--6;
\url{https://par.nsf.gov/servlets/purl/10334188} . These source passages
were opened and read in this run. The new degree tracking and the full
lifting proof above are checked separately in
\texttt{\small CHECK\_\allowbreak{}EVENTS.tsv}.

\section{An explicit stable rectangular-Kronecker lower
bound}\label{an-explicit-stable-rectangular-kronecker-lower-bound}

Statement \textbf{R-014}. Derived on 2026-11-21 UTC by GPT-6 Astra Pro,
agent /root/rectangular\_cases; effort setting not exposed. This is a
corollary of the proved equal-middle-row formula T-025 and the known
Ikenmeyer--Panova inequality. No priority claim is made.

For \(d\ge3\), write

\[
 p_{2,\ldots,d}(s)
 =\#\{(c_2,\ldots,c_d)\in\mathbb Z_{\ge0}^{d-1}:
             2c_2+\cdots+dc_d=s\},
\]

with the value zero for \(s<0\).

\subsection{The actual parameter
family}\label{the-actual-parameter-family}

Let \(d\ge3\), \(\beta\ge\gamma\ge0\), and

\[
 t=(d-2)\beta+\gamma,\qquad
 n\ge\max(1,t),\qquad
 \lambda=(dn-t,\beta^{d-2},\gamma).
\]

Trailing zero parts are omitted. Put

\[
 R=\beta-\gamma-\binom d2(\beta\bmod2).
\]

Then

\[
 \boxed{\displaystyle
 g\bigl(\lambda,(d^n),(d^n)\bigr)
 \ \ge\
 [S_\lambda V:\operatorname{Sym}^d(\operatorname{Sym}^nV)]
 \ =\ p_{2,\ldots,d}(R).}
 \tag{R014.1}
\]

The plethysm coefficient is taken in any sufficiently large dimension.
In particular, throughout the explicitly specified subfamilies \(R=0\)
and \(R\ge2\),

\[
 g\bigl(\lambda,(d^n),(d^n)\bigr)>0.
 \tag{R014.2}
\]

The displayed lower bound is the lattice-point count of one bounded
polytope with \(d-1\) variables. It is uniform in all the parameters and
independent of \(n\) inside the stated stable range.

\subsection{Complete proof of the
corollary}\label{complete-proof-of-the-corollary}

First, \(t\ge\beta\), since \(d\ge3\) and \(\gamma\ge0\). Thus
\(n\ge\beta\). Also

\[
 dn-t\ge(d-1)n\ge\beta,
\]

so every displayed \(\lambda\) is indeed a partition of \(dn\). The
restriction \(n\ge1\) avoids an empty rectangle convention, including
the case \(\beta=\gamma=0\).

Apply T-025 with outer degree \(m=d\). Its correction term
\(d\max(0,\beta-n)\) is zero because \(n\ge\beta\). Its residual integer
is consequently exactly \(R\), and its explicit basis-to-polytope
bijection proves the equality in (R014.1).

Let \(\rho=(\beta^{d-2},\gamma)\), omitting zeros. Its size is \(t\),
and \(\lambda=\rho(dn)\) in the notation of Ikenmeyer--Panova. Their
\textbf{Proposition 2.8, ``Inequality of Multiplicities,'' p.~7},
applies for \(n\ge|\rho|\) and gives exactly the inequality in (R014.1).
This source proposition and its entire proof were read again for this
corollary. The exact primary source is:

Ikenmeyer--Panova, \emph{Rectangular Kronecker coefficients and
plethysms in geometric complexity theory}, arXiv:1512.03798v2:

\url{https://arxiv.org/pdf/1512.03798}

Finally the nonnegative semigroup generated by \(2,3,\ldots,d\) is
\(\{0\}\cup\{2,3,\ldots\}\): even positive integers are sums of 2's, and
an odd integer at least 3 is 3 plus a sum of 2's. This proves the stated
positive subfamilies.

\subsubsection{Exact maps behind the cited
inequality}\label{exact-maps-behind-the-cited-inequality}

The source proof relates coordinate rings through padding and
determinant realization; the comparison is not an identification of
their coefficients. For a precise finite chain, let

\[
 M=n\binom{n+n^2-1}{n},
 \qquad
 \lambda_M=(dM-t,\rho).
\]

Ikenmeyer--Panova Lemma 2.7 gives this finite uniform determinant-size
bound for all degree-\(n\) forms in \(n^2\) variables. Denote the
degree-\(M\) padded-form variety by \(\Gamma_n^M\), as defined in their
section 2(b), and the determinant orbit closure by \(\mathrm{Det}_M\).
Their construction gives the actual closed inclusion

\[
 \Gamma_n^M\subseteq\mathrm{Det}_M
\]

and hence the equivariant coordinate restriction

\[
 \mathbb C[\mathrm{Det}_M]_d
 \twoheadrightarrow\mathbb C[\Gamma_n^M]_d.
\]

Their Proposition 2.6(b) constructs the padding lift on highest-weight
polynomial spaces, preserving the evaluation pairing on padded forms and
changing weight \(\lambda\) to \(\lambda_M\). Thus it gives an injection
of the original highest-weight multiplicity space into the one in
\(\mathbb C[\Gamma_n^M]_d\). The Peter--Weyl comparison in the proof of
their Proposition 2.5 gives the final inequality in the finite chain

\[
 a_{d,n}^{\lambda}
 \le
 \operatorname{mult}_{\lambda_M}\mathbb C[\Gamma_n^M]_d
 \le
 \operatorname{mult}_{\lambda_M}\mathbb C[\mathrm{Det}_M]_d
 \le
 g\bigl(\lambda_M,(d^M),(d^M)\bigr).
\]

Their Theorem 2.1, citing Manivel's Theorem 1, identifies the last term
with \(g(\lambda,(d^n),(d^n))\), because both \(M\) and \(n\) are at
least \(|\rho|=t\). This is the precise chain used by Proposition 2.8.
The new content of (R014.1) is the explicit positive count supplied for
its plethysm term by T-025.

\subsection{Concrete infinite
subfamilies}\label{concrete-infinite-subfamilies}

Every line below holds for every integer \(n\) in the displayed range.

\begin{longtable}[]{@{}
  >{\raggedright\arraybackslash}p{(\columnwidth - 8\tabcolsep) * \real{0.1765}}
  >{\raggedright\arraybackslash}p{(\columnwidth - 8\tabcolsep) * \real{0.1765}}
  >{\raggedright\arraybackslash}p{(\columnwidth - 8\tabcolsep) * \real{0.1765}}
  >{\raggedleft\arraybackslash}p{(\columnwidth - 8\tabcolsep) * \real{0.2353}}
  >{\raggedleft\arraybackslash}p{(\columnwidth - 8\tabcolsep) * \real{0.2353}}@{}}
\toprule\noalign{}
\begin{minipage}[b]{\linewidth}\raggedright
\(d,\beta,\gamma\)
\end{minipage} & \begin{minipage}[b]{\linewidth}\raggedright
Range
\end{minipage} & \begin{minipage}[b]{\linewidth}\raggedright
Partition \(\lambda\)
\end{minipage} & \begin{minipage}[b]{\linewidth}\raggedleft
\(R\)
\end{minipage} & \begin{minipage}[b]{\linewidth}\raggedleft
Proven rectangular-Kronecker lower bound
\end{minipage} \\
\midrule\noalign{}
\endhead
\bottomrule\noalign{}
\endlastfoot
\(3,2,0\) & \(n\ge2\) & \((3n-2,2)\) & 2 & 1 \\
\(3,3,0\) & \(n\ge3\) & \((3n-3,3)\) & 0 & 1 \\
\(4,6,0\) & \(n\ge12\) & \((4n-12,6,6)\) & 6 & 3 \\
\(4,7,1\) & \(n\ge15\) & \((4n-15,7,7,1)\) & 0 & 1 \\
\(5,8,0\) & \(n\ge24\) & \((5n-24,8,8,8)\) & 8 & 5 \\
\end{longtable}

For the lower bound 3, the solutions to \(2c_2+3c_3+4c_4=6\) are
\((3,0,0),(0,2,0),(1,0,1)\). For the lower bound 5, the solutions to
\(2c_2+3c_3+4c_4+5c_5=8\) are
\((4,0,0,0),(1,2,0,0),(2,0,1,0),(0,0,2,0),(0,1,0,1)\).

The inequality can be strict even inside this family. Exact
symmetric-group character sums give

\[
 a_{4,4}^{(12,2,2)}=1,\qquad
 g\bigl((12,2,2),(4^4),(4^4)\bigr)=2.
\]

This calculation is reproduced by verify\_rectangular\_end\_bound.py and
RECTANGULAR\_END\_BOUND\_CHECK.tsv. They also verify equality \(a=g=1\)
at \((d,n,\beta,\gamma)=(3,2,2,0)\) and \((3,3,3,0)\). The cases \(R<0\)
and \(R=1\) give a zero value for the specified plethysm term; (R014.1)
makes no rectangular-Kronecker vanishing assertion there.

\section{Six symmetric slots at the common-power-one
boundary}\label{six-symmetric-slots-at-the-common-power-one-boundary}

Statement \textbf{R-015}. Completed and independently checked on
2026-11-21 UTC by GPT-6 Astra Pro, agent /root/rectangular\_cases;
inherited effort setting not exposed.

\subsection{Exact result}\label{exact-result}

For every finite-dimensional complex vector space \(V\), the natural map

\[
 T_{6,1,5}:\operatorname{Sym}^6(\operatorname{Sym}^5V)
 \longrightarrow
 \operatorname{Sym}^6V\otimes\operatorname{Sym}^6(\operatorname{Sym}^4V),
 \qquad
 f\longmapsto[t_1\cdots t_6]\,
 f(v_1+t_1x,\ldots,v_6+t_6x)
\]

is injective. Here the parameters satisfy \(b=rL-1\) with
\(r=6,L=1,b=5\).

By the exact coordinate-ring morphism R-013, the variety

\[
 Z_{1,5}(V)=\overline{\{\ell q:
        \ell\in V^*,\ q\in\operatorname{Sym}^4V^*\}}
       \subseteq\operatorname{Sym}^5V^*
\]

has no equations of degree at most six:

\[
 I_d(Z_{1,5}(V))=0\qquad(0\le d\le6).
\]

The pullback in degree six is precisely \(5^{-6}T_{6,1,5}\), not an
unspecified nonzero scalar. This follows from R013.1. Vanishing in lower
degrees follows by multiplying a nonzero lower-degree equation to degree
six in the ambient polynomial domain, exactly as in R-013.

No general even-slot boundary classification or historical priority
claim is made.

\subsection{Complete finite proof}\label{complete-finite-proof}

\subsubsection{Reduction and source
character}\label{reduction-and-source-character}

The source embeds in \((\operatorname{Sym}^5V)^{\otimes6}\). Iterated
Pieri multiplication therefore gives only Schur functors with at most
six rows. The natural tensor permutation and symmetrization formula for
\(T_{6,1,5}\) acts by the same multiplicity-space matrices in all
dimensions at least six. Checking \(\mathbb C^6\) consequently detects
every possible kernel summand; smaller dimensions follow by naturality
under an inclusion into \(\mathbb C^6\). This is the same Schur--Weyl
argument proved in R-011, section 1, with six replacing four.

The exact source decomposition on \(\mathbb C^6\) has:

\begin{longtable}[]{@{}lr@{}}
\toprule\noalign{}
Quantity & Value \\
\midrule\noalign{}
\endhead
\bottomrule\noalign{}
\endlastfoot
Distinct irreducible constituents & 893 \\
Sum of highest-weight multiplicities & 5,369 \\
Largest multiplicity & 32 \\
Number of monomial weights & 324,632 \\
Full representation dimension & 377,342,351,232 \\
\end{longtable}

The last dimension equals \(\binom{257}{6}\), because
\(\dim\operatorname{Sym}^5\mathbb C^6=\binom{10}{5}=252\).

The program character\_r6\_b5.cpp computes the complete character twice.
One calculation multiplies the 252 factors

\[
 \prod_{\substack{\alpha\in\mathbb Z_{\ge0}^6\\|\alpha|=5}}
       (1-tx^\alpha)^{-1}
\]

through \(t\)-degree six. The other uses the exact Newton recurrence for
\(h_j[h_5]\), \(j=1,\ldots,6\). It compares every monomial coefficient
for every one of these six outer degrees. At degree six all 324,632
weights agree. The Weyl-alternant coefficient formula then extracts the
multiplicity at every partition of 30 of length at most six. The sum of
multiplicities times Weyl dimensions equals both the sum of all weight
coefficients and \(\binom{257}{6}\). This proves the complete source
multiplicity list used in the certificate.

\subsubsection{Integral source vectors}\label{integral-source-vectors}

Let \(X\) be a \(6\times6\) variable matrix, with columns the six slots.
For every nonempty subset \(I\subseteq\{1,\ldots,6\}\), let \(d_I(X)\)
be the determinant using the first \(|I|\) rows and the columns \(I\) in
increasing order. There are 63 such flag minors. A seed is

\[
 F_c(X)=\prod_I d_I(X)^{c_I}.
\]

The certificate records the 63 nonnegative integers \(c_I\). They
satisfy

\[
 \sum_{I\ni j}c_I=5\quad(1\le j\le6),\qquad
 \lambda_k=\sum_{|I|\ge k}c_I\quad(1\le k\le6).
\]

Thus the seed has equal slot degrees five and highest weight
\(\lambda\). The highest-weight property follows by applying a raising
operator to each flag determinant: replacing a row by an earlier row
repeats a row, giving zero.

The associated symmetric vector is its Reynolds sum

\[
 R_c(X)=\sum_{\sigma\in S_6}F_c(X_{\sigma(1)},\ldots,X_{\sigma(6)}),
\]

including the signs from sorting each determinant's columns. For compact
integer storage, collect equal monomials in the 63 minors, divide their
integer coefficients by their greatest common divisor, and choose the
sign so that the first nonzero coefficient in lexicographic exponent
order is positive. Each selected sum is nonzero. Its divisor is at most
\(720\), since each coefficient is a signed sum of at most 720 terms; in
particular it is a unit modulo \(1,000,003\).

Every source vector is therefore specified exactly by one 63-entry seed
and this fixed normalization rule. The independent checker generates all
720 permutations anew, collects their signed contributions, and checks
the seed's weight and all six slot degrees. Invariance follows directly
by reindexing the complete group sum. No spanning assertion about the
candidate seeds is required: the certificate supplies exactly the
required number of independent images at each highest weight.

\subsubsection{Primary image evaluations by a squarefree
logarithm}\label{primary-image-evaluations-by-a-squarefree-logarithm}

For a fixed point \(X,x\), every flag minor after the shift is exactly
linear:

\[
 d_I(X+t x)=d_I(X)+\sum_{j\in I}a_{I,j}(X,x)t_j.
\]

Every term with two inserted columns vanishes because those columns are
equal to \(x\). To find the coefficient \(t_1\cdots t_6\), work in the
squarefree quotient ring

\[
 \mathbb F_P[t_1,\ldots,t_6]/(t_1^2,\ldots,t_6^2),
 \qquad P=1,000,003.
\]

The printed primary evaluation points have every \(d_I(X)\ne0\). Put
\(C=\prod_Id_I(X)^{c_I}\). For a nonempty subset \(S\), define

\[
 \ell_S=(-1)^{|S|-1}(|S|-1)!
       \sum_I c_I\prod_{j\in S}\frac{a_{I,j}}{d_I(X)},
\]

where the summand is zero for \(S\not\subseteq I\). The formal logarithm
of the product, after its constant factor \(C\) is removed, has
coefficient \(\ell_S\) at \(t_S=\prod_{j\in S}t_j\). Indeed, the
\(k\)-th term of \(\log(1+u)\) contributes \((-1)^{k-1}u^k/k\), and a
product of \(k\) distinct variables has coefficient \(k!\) in \(u^k\).
All relevant denominators are units because \(P>6\).

Exponentiating in this nilpotent ring gives the exact coefficient

\[
 [t_1\cdots t_6]F_c(X+t x)
 =C\sum_{\pi\text{ a set partition of }[6]}
          \prod_{S\in\pi}\ell_S.
\]

The primary evaluator computes this finite sum by the recurrence

\[
 H(\varnothing)=1,\qquad
 H(A)=\sum_{\substack{\varnothing\ne S\subseteq A\\\min A\in S}}
             \ell_SH(A\setminus S).
\]

The block containing the least element of \(A\) makes every set
partition occur exactly once. It then adds the signed normalized
Reynolds contributions.

\subsubsection{Independent verification and integer
lifting}\label{independent-verification-and-integer-lifting}

For every source constituent \(\lambda\) of multiplicity \(m_\lambda\),
the certificate supplies \(m_\lambda\) selected normalized Reynolds
vectors, \(m_\lambda\) evaluation rows, and the complete square matrix
of their image values. Every one of the 893 determinants is nonzero
modulo \(P\).

The independent program recheck\_polarization\_r6.cpp recomputes every
selected entry without logarithms, inverses of minor constants, or set
partitions. It expands the same linear minor factors by descending
in-place Taylor recurrence, in reverse factor order:

\[
 c_A\leftarrow d_Ic_A+
          \sum_{j\in I\cap A}a_{I,j}c_{A\setminus\{j\}}.
\]

Descending subset index leaves each lower-index old coefficient intact.
The independent program also evaluates all small flag determinants and
every selected square determinant using power traces and Newton
identities instead of the primary program's Gaussian elimination. For an
\(m\times m\) matrix, these identities are

\[
 e_0=1,\qquad
 e_k=\frac1k\sum_{i=1}^k(-1)^{i-1}e_{k-i}\operatorname{tr}(A^i),
 \qquad\det A=e_m.
\]

Here \(m\le32<P\). The prime is separately checked by trial division,
and the largest unreduced 64-bit sums are bounded by \(32P^2\).

All \textbf{65,481 selected entries} and all 893 determinants passed the
independent check. Since the source vectors and coefficient-extraction
map have integer coefficients, a nonzero determinant modulo \(P\) is a
nonzero integer determinant. The selected images are therefore linearly
independent over \(\mathbb Q\) and \(\mathbb C\).

Every source highest-weight multiplicity space is thus injective. Any
nonzero equivariant kernel would contain a nonzero highest-weight
vector. This proves the stated injectivity on \(\mathbb C^6\), and the
dimension transfer already proved gives the result for all \(V\).

Two first-pass random seed lists were insufficient to exhibit their full
ranks. Enlarged lists supplied full ranks four and seven in those
blocks. The final selected certificate uses only complete square
matrices; the two incomplete first-pass evaluations are not used in the
proof.

\subsection{Reproducible compact
artifacts}\label{reproducible-compact-artifacts}

The complete certificate is
POLARIZATION\_R6\_L1\_B5\_SELECTED.compact.gz. It contains every
selected seed, point, square matrix and nonzero determinant, and
occupies 329,843 compressed bytes. Its SHA256 is

\begin{verbatim}
6db560ac93c2f53850ddc547946b939556531a0e43aa0e52e18c2ae01e1326ca
\end{verbatim}

The source character is POLARIZATION\_R6\_L1\_B5\_SOURCE.tsv; the
complete rank summary is POLARIZATION\_R6\_L1\_B5\_ALL\_RANKS.tsv; the
independent receipt is
POLARIZATION\_R6\_L1\_B5\_INDEPENDENT\_RECHECK.txt.

To recheck, decompress the compact certificate, compile
recheck\_polarization\_r6.cpp with C++17 and optimization, and give the
decompressed file as its sole argument. The program
character\_r6\_b5.cpp independently reproduces and compares both
complete source-character calculations. The full candidate generator,
original evaluator, and merge script are polarization\_r6\_seeds.py,
check\_polarization\_r6.cpp, and merge\_polarization\_r6.py. The
parallel intermediate directory is unnecessary for verifying the final
certificate.

The next pair \(r=6,L=2,b=11\) has 13,019,909 monomial weights before
Schur extraction and source dimension 9,679,525,138,407,392,632 on
\(\mathbb C^6\). No rank certificate or theorem for that pair is
asserted here.

\section{Three complete finite boundary
kernels}\label{three-complete-finite-boundary-kernels}

Statement \textbf{R-016}, completed and independently checked on
2026-11-21 UTC by GPT-6 Astra Pro, /root/rectangular\_cases; effort
setting not exposed.

For the coefficient-extraction morphism

\[
 T_{r,L,b}:\operatorname{Sym}^r(\operatorname{Sym}^bV)
 \longrightarrow
 \operatorname{Sym}^{rL}V\otimes
 \operatorname{Sym}^r(\operatorname{Sym}^{b-L}V),
 \quad
 f\longmapsto[t_1^L\cdots t_r^L]\,
 f(v_1+t_1x,\ldots,v_r+t_rx),
\]

the following are the exact kernels for every finite-dimensional complex
\(V\):

\[
 \boxed{\ker T_{3,2,5}=0,}
 \qquad
 \boxed{\ker T_{3,3,8}=S_{(8,8,8)}V,}
 \qquad
 \boxed{\ker T_{5,1,4}=S_{(4,4,4,4,4)}V.}
\]

In the last two formulas the right-hand side denotes the unique embedded
Schur summand in the displayed source. It is specified concretely below.
Thus it is zero when the ambient dimension is smaller than its number of
rows.

\subsection{Exact embedded rectangular
kernel}\label{exact-embedded-rectangular-kernel}

In dimension \(r\), the polynomial \(\det(v_1,\ldots,v_r)^b\) has degree
\(b\) in each of the \(r\) slots. For even \(b\) it is invariant under
slot permutations. It is a nonzero highest-weight vector of weight
\((b^r)\). Its source multiplicity is one: the corresponding
\(GL_r\)-module is \(\det^b\), and the balanced-weight/permutation
realization gives precisely its sign character
\(\operatorname{sgn}^b=1\).

For the two pairs \((r,L,b)=(3,3,8)\) and \((5,1,4)\), one has
\(rL=b+1\). In a term contributing to \(t_1^L\cdots t_r^L\), exactly
\(b+1\) columns of the \(b\) determinant factors are replaced by the
common vector \(x\). Some determinant factor receives at least two
copies of \(x\), and that term is zero. Hence the entire displayed
determinant power lies in the kernel.

For any ambient \(V\) of dimension at least \(r\), choose the first
\(r\) basis directions and use the determinant of the corresponding
\(r\times r\) coordinate minor. The \(GL(V)\)-span of its \(b\)-th power
is the embedded Cartan module \(S_{(b^r)}V\): the polynomial is a
nonzero highest-weight vector of that weight, and the source
multiplicity is one by the stable Schur-functor multiplicity-space
realization. This defines the actual inclusion in the source.
Equivariance places its entire span in the kernel.

The selected-minor certificates below prove injectivity on every other
source summand, so this explicit inclusion accounts for the complete
kernel.

\subsection{Complete rank
certificates}\label{complete-rank-certificates}

\begin{longtable}[]{@{}
  >{\raggedright\arraybackslash}p{(\columnwidth - 10\tabcolsep) * \real{0.1304}}
  >{\raggedleft\arraybackslash}p{(\columnwidth - 10\tabcolsep) * \real{0.1739}}
  >{\raggedleft\arraybackslash}p{(\columnwidth - 10\tabcolsep) * \real{0.1739}}
  >{\raggedleft\arraybackslash}p{(\columnwidth - 10\tabcolsep) * \real{0.1739}}
  >{\raggedleft\arraybackslash}p{(\columnwidth - 10\tabcolsep) * \real{0.1739}}
  >{\raggedleft\arraybackslash}p{(\columnwidth - 10\tabcolsep) * \real{0.1739}}@{}}
\toprule\noalign{}
\begin{minipage}[b]{\linewidth}\raggedright
\((r,L,b)\)
\end{minipage} & \begin{minipage}[b]{\linewidth}\raggedleft
Source constituents
\end{minipage} & \begin{minipage}[b]{\linewidth}\raggedleft
Source highest-weight dimension
\end{minipage} & \begin{minipage}[b]{\linewidth}\raggedleft
Certified image highest-weight rank
\end{minipage} & \begin{minipage}[b]{\linewidth}\raggedleft
Largest source multiplicity
\end{minipage} & \begin{minipage}[b]{\linewidth}\raggedleft
Independent selected entries
\end{minipage} \\
\midrule\noalign{}
\endhead
\bottomrule\noalign{}
\endlastfoot
\((3,2,5)\) & 13 & 13 & 13 & 1 & 13 \\
\((3,3,8)\) & 36 & 39 & 38 & 2 & 44 \\
\((5,1,4)\) & 95 & 142 & 141 & 4 & 263 \\
\end{longtable}

The source dimensions in ambient dimensions \(r\) are respectively
1,771; 16,215; and 16,108,764.

The producer small\_boundary\_candidates.py computes the complete source
characters from Newton's recurrence, extracts every Schur multiplicity
by the Weyl alternant, and checks the full Weyl-dimension sum against
the dimension of the symmetric power. The independent program
recheck\_small\_boundary.py obtains every source multiplicity, including
all zeros, by power-sum plethysm expansion and Murnaghan--Nakayama
character sums. It checks all 27, 61 and 192 partitions of the
respective degrees and permitted lengths.

For each nonrectangular highest weight, the supplied source vectors are
signed, normalized Reynolds sums of products of flag minors. Their
complete definition is the same one proved in R-015, with \(r\) slots:
average all \(r!\) literal slot permutations, collect the determinant
signs, divide the integral coefficients by their gcd, and choose the
first nonzero coefficient positive. The seeds come from semistandard
tableaux of content \((b^r)\). Each seed's row weights and every slot
degree are independently verified.

The primary C++ program check\_small\_boundary.cpp evaluates
\(T_{r,L,b}\) by multiplying the exactly linear shifted minors and
retaining the exponents \(0\le\deg t_i\le L\). It selects a square
matrix of size equal to the full source multiplicity in each
nonrectangular block. Every selected determinant is nonzero modulo the
prime \(1,000,003\). The rectangular blocks are handled by the explicit
polynomial proof above, not by a numerical inference from sampled zeros.

The independent Python program recomputes every selected matrix entry
with a separate NumPy Taylor-array implementation, reverses the factor
order, generates every Reynolds sum independently, and computes all
small and selected determinants by recursive Laplace expansion. The
programs agree on all 320 selected entries. They made 40, 183 and 20,402
independent Taylor evaluations respectively. The largest selected matrix
has size four, so the Laplace checks are exhaustive exact calculations.
Array intermediates are bounded by \(6P^2\), inside the signed 64-bit
range.

A nonzero determinant modulo \(P\) proves that the corresponding integer
polynomial-image columns are linearly independent over \(\mathbb Q\) and
\(\mathbb C\). Thus every nonrectangular highest-weight multiplicity
space is injective. Complete reducibility then proves the stated kernels
in dimension \(r\).

Finally, the source has only Schur functors of length at most \(r\), by
iterated Pieri multiplication of \(r\) one-row representations. The
coefficient-extraction map is defined by the same tensor permutations
and averages in every dimension. Its Schur-multiplicity matrices are
therefore dimension-independent once all source partitions fit, exactly
as proved in R-011. This transfers the complete kernel descriptions to
every larger \(V\). Naturality under an inclusion handles smaller \(V\).

The data are fully specified in the three files
POLARIZATION\_R3\_L2\_B5.selected.gz,
POLARIZATION\_R3\_L3\_B8.selected.gz, and
POLARIZATION\_R5\_L1\_B4.selected.gz, together with their source and
rank TSV tables. SMALL\_BOUNDARY\_INDEPENDENT\_CHECK.tsv is the
independent receipt. Decompress the selected files and run
recheck\_small\_boundary.py from any directory to reproduce the
independent character and entry checks.

These finite kernels agree with the rectangular obstruction at these
actual parameter pairs. No assertion about all remaining boundary
parameters is inferred from the calculation.

\section{An exact polynomial model for every highest-weight boundary
kernel}\label{an-exact-polynomial-model-for-every-highest-weight-boundary-kernel}

Statement \textbf{R-017}. This is an unconditional reformulation of the
common-polarization kernel, not a classification of its remaining
multiplicities. Written and checked on 2026-11-21 UTC by GPT-6 Astra
Pro, /root/rectangular\_cases; effort setting not exposed. The
first-coordinate evaluation underlying the reformulation is also used in
Kadish--Landsberg, arXiv:1204.4693, Theorem 1.3 and its proof in section
2.

Fix \(r\ge1\), \(b\ge L\ge0\), and \(m\ge2\). In each of the \(r\) slots
write

\[
 v_i=(x_i,y_{i,2},\ldots,y_{i,m}),\qquad
 z_{i,j}=y_{i,j}/x_i.
\]

For a polynomial \(P(z_1,\ldots,z_r)\) whose degree in each block
\(z_i\) is at most \(b\), define its separate homogenization

\[
 H_b(P)=\prod_{i=1}^r x_i^b
 P(y_{1,2}/x_1,\ldots,y_{r,m}/x_r).
 \tag{R017.1}
\]

The block degree bounds make this a polynomial. Evaluation at every
\(x_i=1\) is its inverse. These are exact linear isomorphisms between
polynomials of block degrees at most \(b\) and polynomials separately
homogeneous of degree \(b\) in the \(v_i\). They commute with slot
permutations.

\subsection{Highest vectors become translated-polynomial
invariants}\label{highest-vectors-become-translated-polynomial-invariants}

The raising operator \(E_{1j}\), \(j>1\), acts on a separately
homogeneous polynomial as \(\sum_i x_i\partial_{y_{i,j}}\). Direct
differentiation gives

\[
 E_{1j}H_b(P)=
 \prod_i x_i^b\sum_i\partial_{z_{i,j}}P.
 \tag{R017.2}
\]

Thus its kernel corresponds exactly to invariance under simultaneous
translation \(z_{i,j}\mapsto z_{i,j}+c_j\). The equality follows by
exponentiating the displayed locally nilpotent derivations in
characteristic zero; differentiating the translation identity proves the
converse.

For \(2\le k<j\le m\), the remaining raising operators correspond
exactly to

\[
 \sum_i z_{i,k}\partial_{z_{i,j}}.
 \tag{R017.3}
\]

A source highest vector of weight \(\lambda\vdash rb\) therefore
corresponds to a symmetric, simultaneously translation-invariant
polynomial \(P\), of individual block degrees at most \(b\), with
coordinate weights \((\lambda_2,\ldots,\lambda_m)\), annihilated by
(R017.3). Its total \(z\)-degree is

\[
 D=\lambda_2+\cdots+\lambda_m=rb-\lambda_1.
 \tag{R017.4}
\]

Every polynomial satisfying these displayed properties homogenizes to
precisely such a highest vector; the formulas above prove both
directions.

\subsection{Polarization is weighted block
truncation}\label{polarization-is-weighted-block-truncation}

For a monomial

\[
 z^\alpha=\prod_{i=1}^r\prod_{j=2}^m z_{i,j}^{\alpha_{i,j}},
 \qquad d_i=\sum_{j=2}^m\alpha_{i,j},
\]

define the explicit diagonal operator

\[
 Q_{L,b}(z^\alpha)=
 \begin{cases}
 \left(\prod_i\binom{b-d_i}{L}\right)z^\alpha,
       &d_i\le b-L\text{ for every }i,\\
 0,&\text{otherwise}.
 \end{cases}
 \tag{R017.5}
\]

The coefficient-extraction map satisfies the exact identity

\[
 T_{r,L,b}(H_b(P))(v_1,\ldots,v_r;\,e_1)
 =H_{b-L}(Q_{L,b}P).
 \tag{R017.6}
\]

To prove it, a monomial in (R017.1) is
\(\prod_i x_i^{b-d_i}y_i^{\alpha_i}\). Substitution of \(v_i+t_ie_1\)
changes only \(x_i\) to \(x_i+t_i\). The coefficient of \(t_i^L\) is
\(\binom{b-d_i}{L}x_i^{b-d_i-L}y_i^{\alpha_i}\), with value zero for
\(b-d_i<L\). Multiplication over all slots gives (R017.6) on every
monomial.

Restriction of a highest image to common vector \(e_1\) detects its
vanishing. To see this directly, let \(G(v_1,\ldots,v_r;w)\) be such an
image. It is annihilated by every \(E_{1j}\), and is homogeneous of
degree \(rL\) in the common vector \(w=(w_1,\ldots,w_m)\). On the dense
open set \(w_1\ne0\), apply the unipotent substitutions

\[
 w_j\mapsto w_j-\frac{w_j}{w_1}w_1=0,\qquad
 y_{i,j}\mapsto y_{i,j}-\frac{w_j}{w_1}x_i
 \quad(j>1).
\]

Invariance under the exponentials of \(E_{1j}\) leaves \(G\) unchanged.
The common vector becomes \(w_1e_1\), and common-vector homogeneity
gives

\[
 G(v_1,\ldots,v_r;w)
 =w_1^{rL}
 G(v'_1,\ldots,v'_r;e_1).
 \tag{R017.7}
\]

Hence a zero restriction at \(e_1\) makes \(G\) zero on a dense open set
and therefore identically zero. The converse is immediate.

Combining (R017.1)--(R017.7) gives the exact highest-weight kernel
isomorphism

\[
 \boxed{
 H_b^{-1}\bigl((\ker T_{r,L,b})_\lambda^{\mathrm{highest}}\bigr)
 =
 \ker Q_{L,b}\ \cap\
 \{\text{polynomials satisfying (R017.2)–(R017.4) and slot symmetry}\}.
 }
 \tag{R017.8}
\]

All nonzero diagonal coefficients in (R017.5) are nonzero rational
integers. Its kernel is therefore exactly the kernel of ordinary
monomial truncation retaining only the monomials with every block degree
at most \(b-L\). This identifies the relevant subspace with explicit
linear equations, rather than inferring it from source and target
multiplicities.

\subsection{The precise boundary degree
gap}\label{the-precise-boundary-degree-gap}

At \(r\ge2,L\ge1,b=rL-1\), a source partition has \(\lambda_1\ge b\),
because its length is at most \(r\) and its size is \(rb\). Equality
forces the rectangle \((b^r)\). For that rectangle,

\[
 D=(r-1)b=r(b-L)+1,
\]

so every monomial exceeds the total degree allowed by the block
truncation. In dimensions \(m\ge r\), the determinant power provides the
corresponding source highest vector exactly when \(b\) is even; at this
boundary that parity is equivalent to \(r\) and \(L\) both odd. For
\(m<r\), that Schur functor is zero.

Every other source partition has \(\lambda_1\ge b+1=rL\), so (R017.4)
gives \(D\le r(b-L)\). The total-degree obstruction is then absent.
Formula (R017.8) specifies the actual remaining kernel problem as a
translation-invariant symmetric polynomial truncation problem; it does
not assert that its remaining intersection is zero. R-011, R-012, R-015
and R-016 supply the completely certified finite cases currently
established in this run.

\section{A complete second-row injectivity range for common
polarization}\label{a-complete-second-row-injectivity-range-for-common-polarization}

Statement \textbf{R-018}. Written and checked on 21 November 2026 UTC by
GPT-6 Astra Pro,
\texttt{\small /\allowbreak{}root/\allowbreak{}rectangular\_\allowbreak{}cases};
the effort setting is not exposed. This is an exact consequence of the
flag-minor algebra and the truncation morphism R-017. No historical
novelty claim is made.

\subsection{The actual maps and the unconditional
range}\label{the-actual-maps-and-the-unconditional-range}

Let \(r\ge1\), \(b\ge L\ge0\), and let \(V\) be a finite-dimensional
vector space over a field of characteristic zero. On polynomials in
\(r\) vector slots, each homogeneous of degree \(b\), define

\[
\widetilde T_{r,L,b}f(v_1,\ldots,v_r;w)
 =[t_1^L\cdots t_r^L]
 f(v_1+t_1w,\ldots,v_r+t_rw).
\tag{R018.1}
\]

This is the \(GL(V)\)-equivariant, \(S_r\)-equivariant map

\[
(\operatorname{Sym}^bV)^{\otimes r}
\longrightarrow
\operatorname{Sym}^{rL}V\otimes
(\operatorname{Sym}^{b-L}V)^{\otimes r}.
\tag{R018.2}
\]

Here the polynomial-generator realization of tensors is used, as in
R-017. The permutation group acts on the \(r\) slots in both
expressions. For \(\mu\vdash r\), application of
\(\operatorname{Hom}_{S_r}([\mu],-)\) gives the actual induced map

\[
T_{r,L,b}^{\mu}:S_\mu(\operatorname{Sym}^bV)
\longrightarrow
\operatorname{Sym}^{rL}V\otimes
S_\mu(\operatorname{Sym}^{b-L}V).
\tag{R018.3}
\]

The standard Schur-functor realization here is
\(S_\mu W=\operatorname{Hom}_{S_r}([\mu],W^{\otimes r})\). For
\(\mu=(r)\), (R018.3) is exactly the map \(T_{r,L,b}\) used in R-011
through R-017. Formula (R018.1) fixes the coefficient-extraction
normalization; the product of ordinary differential operators is
\((L!)^r\widetilde T_{r,L,b}\).

For every \(\lambda\vdash rb\) with \(\lambda_2\le b-L\), and for every
\(\mu\vdash r\),

\[
\boxed{
\bigl(\ker\widetilde T_{r,L,b}\bigr)_{\lambda}=0,
\qquad
\bigl(\ker T_{r,L,b}^{\mu}\bigr)_{\lambda}=0.}
\tag{R018.4}
\]

The subscript denotes the entire \(GL(V)\)-isotypic component. A missing
second part is zero. The formulas include all source multiplicities and
all ambient dimensions. They assert an actual injectivity range for the
specified map.

\subsection{1. Why every highest polynomial has the block-degree
bound}\label{why-every-highest-polynomial-has-the-block-degree-bound}

Take \(m=\dim V\), with polynomial generators \(X_{ji}\), where
\(1\le j\le m\) is the vector coordinate and \(1\le i\le r\) is the
slot. Write \(x_i=X_{1i}\). For a subset \(I\subseteq\{1,\ldots,r\}\) of
size \(k\le m\), put

\[
\Delta_{k,I}=\det(X_{ji})_{1\le j\le k,\ i\in I},
\tag{R018.5}
\]

with the columns in increasing order. This is a highest polynomial of
\(GL(V)\)-weight \((1^k)\). Its slot weight is one in the slots in \(I\)
and zero elsewhere.

The algebra of highest polynomials is generated by these minors. Here is
a representation-theoretic proof of that precise generation statement.
The Cauchy decomposition of the matrix polynomial ring, followed by
taking highest vectors for the row action, is

\[
\mathbf Q[X_{ji}]^{U_m}
 =\bigoplus_{\ell(\eta)\le\min(m,r)} S_\eta\mathbf Q^r,
\tag{R018.6}
\]

with the summand tagged by its row highest weight \(\eta\). The algebra
generated by (R018.5) is stable under the column \(GL_r\): the formula
for a minor of a matrix product expresses its transform as a linear
combination of minors of the same size. For every summand in (R018.6),
the nonzero product

\[
\prod_{k=1}^{\min(m,r)}
\Delta_{k,\{1,\ldots,k\}}^{\eta_k-\eta_{k+1}}
\tag{R018.7}
\]

belongs to that summand and to the generated algebra. The summand is an
irreducible column \(GL_r\)-module, so the column-stable algebra
contains it in full. This proves generation for every summand and hence
for the whole highest algebra. The same proof works after extension of
scalars.

Now let \(f\) be any highest vector of weight \(\lambda\) in
\((\operatorname{Sym}^bV)^{\otimes r}\). Express it using products of
the minors, and project that expression onto row weight \(\lambda\) and
slot multidegree \((b,\ldots,b)\). Each minor is homogeneous for both
gradings, so this projection retains precisely products of those two
degrees. If \(c_k\) is the number of height-\(k\) factors in any
retained product, then

\[
\lambda_j=\sum_{k\ge j}c_k,
\qquad
\sum_{k\ge2}c_k=\lambda_2.
\tag{R018.8}
\]

Set every \(x_i=1\) and put \(z_i=(X_{2i},\ldots,X_{mi})\). A height-one
factor becomes the constant one. Each factor of height at least two is a
polynomial of degree at most one in each block \(z_i\), because a
determinant uses a column exactly once. Equation (R018.8) therefore
proves

\[
P(z_1,\ldots,z_r)=f\bigl((1,z_1),\ldots,(1,z_r)\bigr)
\quad\Longrightarrow\quad
\deg_{z_i}P\le\lambda_2\quad(1\le i\le r).
\tag{R018.9}
\]

This conclusion concerns every highest vector, including arbitrary
linear combinations in a multiplicity space. Neither slot symmetry nor a
choice of highest-vector basis was used.

\subsection{2. The truncation map and its explicit inverse on this
range}\label{the-truncation-map-and-its-explicit-inverse-on-this-range}

Separate homogenization is

\[
H_b(P)=\prod_i x_i^b P(z_{i,j}=X_{ji}/x_i).
\tag{R018.10}
\]

It inverts evaluation at \(x_i=1\) on separately degree-\(b\)
polynomials. For a monomial \(z^\alpha\) of block degrees
\(d_i=\sum_{j\ge2}\alpha_{i,j}\), the exact calculation from R-017 is

\[
\widetilde T_{r,L,b}(H_b(P))(v_1,\ldots,v_r;e_1)
 =H_{b-L}(Q_{L,b}P),
\tag{R018.11}
\]

where

\[
Q_{L,b}(z^\alpha)
 =\left(\prod_i\binom{b-d_i}{L}\right)z^\alpha
\qquad(d_i\le b-L\text{ for all }i).
\tag{R018.12}
\]

Monomials exceeding the displayed block bound have image zero. On the
range (R018.9) with \(\lambda_2\le b-L\), every binomial factor in
(R018.12) is a positive integer. Thus the following inverse is an actual
linear map on that polynomial subspace:

\[
Q_{L,b}^{-1}(z^\alpha)
 =\left(\prod_i\binom{b-d_i}{L}^{-1}\right)z^\alpha.
\tag{R018.13}
\]

In particular a nonzero highest vector has nonzero restriction of its
image to common vector \(e_1\). The kernel of (R018.2) is a
\(GL(V)\)-submodule. A nonzero part in its \(\lambda\)-isotypic
component would contain a nonzero highest vector of weight \(\lambda\),
contradicting (R018.9)--(R018.13). This proves the first equality in
(R018.4).

The coefficient-extraction map commutes with slot permutations.
Application of \(\operatorname{Hom}_{S_r}([\mu],-)\) to this injective
restriction remains injective. Equivalently, any vector in the kernel of
the induced map gives an \(S_r\)-equivariant map into the already zero
kernel component. This proves the second equality in (R018.4), with the
exact source and target in (R018.3).

For \(m=0\) every positive-degree source is zero, and the degree-zero
map is the identity. For \(m=1\) the only source partition is a row, and
(R018.1) is multiplication by \(\binom bL^r\), which is nonzero. These
statements cover the low dimensions not requiring the matrix argument.

\subsection{3. Consequences for the boundary and the padded-form
ideal}\label{consequences-for-the-boundary-and-the-padded-form-ideal}

At the boundary \(r\ge2\), \(L\ge1\), \(b=rL-1\), (R018.4) gives

\[
\ker T_{r,L,rL-1}
\ \subseteq\!
\bigoplus_{\substack{\lambda\vdash r(rL-1)\\
                     \lambda_2\ge(r-1)L}}
\left(\operatorname{Sym}^r(\operatorname{Sym}^{rL-1}V)\right)_\lambda.
\tag{R018.14}
\]

Every component with \(\lambda_2\le(r-1)L-1\) is injective. In
particular, for the original common power \(L=r\), every possible kernel
component has second row at least \(r(r-1)\).

More generally, let

\[
Z_{L,b}=\overline{\{\ell^Lq:
\ell\in V^*,\ q\in\operatorname{Sym}^{b-L}V^*\}}.
\]

The explicitly normalized pullback in R-013 is
\((\nu_{L,b}^*)_r=\binom bL^{-r}T_{r,L,b}\). Combining that exact
identity with (R018.4) yields

\[
\boxed{I_r(Z_{L,b})_\lambda=0
\qquad(\lambda\vdash rb,\ \lambda_2\le b-L).}
\tag{R018.15}
\]

Equations (R018.14)--(R018.15) specify which whole isotypic components
survive the actual morphism. The still-unclassified components at the
boundary lie in the displayed complementary second-row range. R-016
classifies three finite boundary kernels by independent exact
certificates; no general rectangular-kernel classification is asserted
here.

\subsection{Relation to checked primary
material}\label{relation-to-checked-primary-material}

Kadish--Landsberg, \emph{Padded polynomials, their cousins, and
geometric complexity theory},
\href{https://arxiv.org/pdf/1204.4693}{arXiv:1204.4693v1}, Theorem 1.7,
printed p.~3, with its full proof in §3, gives the generalized
polarization morphism whose kernel is this padded-form ideal. Their
Theorem 1.3, printed p.~2, and complete proof in §2 use first-coordinate
evaluation to establish first-row ranges. R-013 and R-017 specify the
coefficient-extraction normalization and the precise polynomial
truncation. The flag-minor algebra argument above supplies the stated
second-row range directly. T-032 gives a separate explicit balanced
chart using the same flag generators.

\section{Exact common polarization on the first boundary row
band}\label{exact-common-polarization-on-the-first-boundary-row-band}

Statements \textbf{R-019} and \textbf{R-020}. Written and checked on 21
November 2026 UTC by GPT-6 Astra Pro,
\texttt{\small /\allowbreak{}root/\allowbreak{}rectangular\_\allowbreak{}cases};
effort setting not exposed. R-019 is an explicit coefficient
calculation. R-020 uses the complementary Schur pairing and the
multilinear Young-symmetrizer map to identify the entire indicated band.
No historical novelty claim is made.

Throughout, \(T_{r,L,b}\) means the coefficient-extraction map

\[
T_{r,L,b}f(v_1,\ldots,v_r;w)
=[t_1^L\cdots t_r^L]f(v_1+t_1w,\ldots,v_r+t_rw),
\tag{N.1}
\]

from \(\operatorname{Sym}^r(\operatorname{Sym}^bV)\) to
\(\operatorname{Sym}^{rL}V\otimes
\operatorname{Sym}^r(\operatorname{Sym}^{b-L}V)\). The product of
ordinary differential operators is \((L!)^rT_{r,L,b}\).

\subsection{R-019. A nonrectangular boundary component with an exact
coefficient}\label{r-019.-a-nonrectangular-boundary-component-with-an-exact-coefficient}

Let \(r\ge2\), \(L\ge1\), and \(rL\) be even. Put

\[
b=rL-1,\qquad C=r(L-1),\qquad
\lambda=((rL)^{r-1},C),\qquad
s=\lfloor r/2\rfloor,\quad t=r-s.
\tag{N.2}
\]

The source \(\lambda\)-isotypic component has multiplicity one, and
\(T_{r,L,b}\) is injective on that component in every ambient dimension.
Here a Schur functor of length greater than \(\dim V\) is zero.

An explicit highest vector in dimension \(r\) is as follows. Let
\(M=(v_1,\ldots,v_r)\) be the matrix of vector-slot generators. Write
\(D=\det M\), let \(B\) be its first \(r-1\) rows, and set

\[
y_i=(-1)^{r+i}\det B_{\widehat i},\qquad
f=D^C\prod_{i=1}^r y_i.
\tag{N.3}
\]

The hat deletes column \(i\). The same formula with \(D^0=1\) defines
the vector in dimension \(r-1\) when \(L=1\).

To specify an exact evaluation of its image, split the column indices
into \(I=\{1,\ldots,s\}\) and \(J=\{s+1,\ldots,r\}\). Let \(A\) be the
\((r-2)\times r\) matrix whose rows, in the indicated order, are

\[
e_i-e_s\quad(1\le i<s),\qquad
e_{s+j}-e_r\quad(1\le j<t),
\tag{N.4}
\]

and put \(c=-e_s+e_r\). Here \(e_i\) denotes a row coordinate vector in
this displayed matrix construction. Take the vector-slot matrix \(M_0\)
with first row zero, middle rows \(A\), and last row \(c\). The common
vector in (N.1) is the first vector-coordinate basis vector \(e_1\).
Then

\[
\boxed{
T_{r,L,b}f(M_0;e_1)
=(-1)^s
\binom{r(L-1)}{sL-t}
\frac{(sL)!\,(tL)!}{(L!)^r}
\ne0.}
\tag{N.5}
\]

For \(L=1\), the last row of \(M_0\) can be deleted: the polynomial
\(f\) does not use it. Thus the same evaluation proves the statement in
the minimum nonzero ambient dimension in that case.

\subsubsection{Proof of the highest vector and its slot
symmetry}\label{proof-of-the-highest-vector-and-its-slot-symmetry}

Every cofactor \(y_i\) is a flag minor of height \(r-1\), with row
highest weight \((1^{r-1},0)\). Their product has weight
\((r^{r-1},0)\). Multiplication by \(D^C\) adds \((C^r)\), giving
exactly \(\lambda\) in (N.2). The polynomial is nonzero because it is a
product of nonzero polynomials in an integral domain.

Each column belongs to \(r-1\) of the cofactor factors, so its total
degree in \(f\) is \(C+r-1=b\). A column permutation \(\sigma\) sends
the cofactor coordinates to \(\operatorname{sgn}(\sigma)y_{\sigma(i)}\),
with the inverse indexing when the opposite permutation convention is
used. Consequently its effect on their product is
\(\operatorname{sgn}(\sigma)^r\). The determinant power contributes
\(\operatorname{sgn}(\sigma)^C\). The combined character is
\(\operatorname{sgn}(\sigma)^{C+r}
=\operatorname{sgn}(\sigma)^{rL}=1\). Thus (N.3) belongs to the
symmetric-slot source of (N.1).

\subsubsection{Proof of the coefficient and all
signs}\label{proof-of-the-coefficient-and-all-signs}

Replace the first row of \(M_0\) by independent variables
\(x=(x_1,\ldots,x_r)\), and write

\[
X=\sum_{i\in I}x_i,\qquad Y=\sum_{i\in J}x_i.
\]

The kernel of the matrix \(A\) in (N.4) consists exactly of vectors
constant on each of the two groups. The cofactor vector \(y\) of
\(B=[x;A]\) lies in that kernel and satisfies \(x\cdot y=0\). Its last
coordinate can also be computed directly:

\[
y_r=\det B_{\widehat r}=(-1)^{s-1}X.
\]

Indeed the rows belonging to the second group leave an identity block on
columns \(s+1,\ldots,r-1\), and the remaining determinant has rows
\((x_1,\ldots,x_s)\) and \(e_i-e_s\), \(1\le i<s\). Expansion gives
\((-1)^{s-1}\sum_{i\le s}x_i\). The two kernel conditions now give the
full vector, as polynomial identities:

\[
y_i=(-1)^sY\quad(i\in I),\qquad
y_i=(-1)^{s+1}X\quad(i\in J).
\tag{N.6}
\]

Expansion of \(D\) in its last row gives

\[
D=c\cdot y=(-1)^{s+1}(X+Y).
\tag{N.7}
\]

Combining (N.6)--(N.7), the specialization of (N.3) is

\[
(-1)^{(s+1)C+sr+t}(X+Y)^C X^tY^s.
\]

Since \(C+r=rL\) is even, its sign exponent is congruent to
\(C+t\equiv r+t\equiv s\pmod2\). Thus it equals

\[
(-1)^s(X+Y)^C X^tY^s.
\tag{N.8}
\]

To obtain \(\prod_i x_i^L\), the total \(X\)-degree must be \(sL\) and
the total \(Y\)-degree must be \(tL\). The unique relevant term of
\((X+Y)^C\) has \(X\)-degree \(sL-t\). Its coefficient is
\(\binom C{sL-t}\). Expansion of \(X^{sL}\) and \(Y^{tL}\) then gives
the two multinomial coefficients \((sL)!/(L!)^s\) and \((tL)!/(L!)^t\).

For even \(r\), \(s=t\) and \(sL-t=s(L-1)\) lies between zero and \(C\).
For odd \(r\), the assumed evenness of \(rL\) makes \(L\ge2\) even; then
\(sL\ge t=s+1\) and \(tL\ge s\), which give the same two bounds. All
factorials are therefore defined and the coefficient is nonzero.

Finally the first row of \(M_0\) is zero, so substitution in (N.1)
replaces this row exactly by \((t_1,\ldots,t_r)\). Hence the computed
coefficient is precisely the value in (N.5), with no omitted derivative
or averaging factor. This proves nonvanishing of the actual map.

The multiplicity-one assertion follows from the exact pairing and
multilinear identification in R-020 below. A nonzero map on a single
irreducible source copy is injective. The tensor map (N.1) is natural
and is given by tensor-position symmetrizations and permutations
independent of dimension. Thus its map on a Schur multiplicity space is
the same in every ambient dimension where that Schur functor is nonzero.
Alternatively the above construction already uses the minimum possible
dimension: \(r\) for \(C>0\), and \(r-1\) for \(C=0\); embedding that
space preserves the displayed highest vector and its nonzero image.

\subsubsection{Actual examples and relation to the original
boundary}\label{actual-examples-and-relation-to-the-original-boundary}

\[
\begin{array}{c|c|c|c|r}
r&L&b&\lambda&Tf(M_0;e_1)\\\hline
3&2&5&(6,6,3)&-6\\
4&1&3&(4,4,4,0)&4\\
4&4&15&(16,16,16,12)&4527600\\
5&2&9&(10,10,10,10,5)&2700\\
6&1&5&(6,6,6,6,6,0)&-36
\end{array}
\tag{N.9}
\]

For every even \(r\ge2\), the original common power \(L=r\) gives the
explicit surviving component

\[
\lambda=((r^2)^{r-1},r^2-r)
\quad\text{at}\quad b=r^2-1.
\tag{N.10}
\]

Its second row is \(r^2\) for \(r\ge3\), outside the range
\(\lambda_2\le r(r-1)-1\) proved in R-018. This result concerns the
specified common-polarization morphism; it does not identify an
additional component of a different map without a further morphism
calculation.

\subsection{\texorpdfstring{R-020. The complete band
\(\lambda_1\le b+1\)}{R-020. The complete band \textbackslash lambda\_1\textbackslash le b+1}}\label{r-020.-the-complete-band-lambda_1le-b1}

Fix \(r\ge2,L\ge1,b=rL-1\), and put \(H=b+1=rL\). In the source
\(\operatorname{Sym}^r(\operatorname{Sym}^bV)\), the sum of all isotypic
components with \(\lambda_1\le H\) is

\[
\boxed{
\begin{cases}
S_{(H^{r-1},H-r)}V,&H\text{ even},\\
S_{((H-1)^r)}V,&H\text{ odd},
\end{cases}}
\tag{N.11}
\]

with multiplicity one in each case. The map \(T_{r,L,b}\) is injective
on the entire band in the first line. In the second line the entire
displayed component is its kernel in that band. This classifies the map
on every source component with first row at most \(b+1\), for every
\(r,L\).

\subsubsection{1. The exact complementary Schur
pairing}\label{the-exact-complementary-schur-pairing}

Work first in the \(r\)-dimensional slot space \(E=\mathbf Q^r\). Every
partition in \((\operatorname{Sym}^bV)^{\otimes r}\) has at most \(r\)
rows. For \(\lambda\vdash rb\) inside the rectangle \((H^r)\), define

\[
\nu=(H-\lambda_r,\ldots,H-\lambda_1).
\tag{N.12}
\]

It is a partition of \(rH-rb=r\). There is a nonzero perfect
\(GL(E)\)-equivariant pairing

\[
S_\lambda E\otimes S_\nu E\longrightarrow(\det E)^H.
\tag{N.13}
\]

An explicit construction uses the exterior columns. If the \(H\) column
heights of \(\lambda\), with zero heights included, are
\(h_1\ge\cdots\ge h_H\), those of \(\nu\) are \(r-h_H,\ldots,r-h_1\).
Embed each Schur module as its Cartan component in the corresponding
tensor product of exterior powers. Reverse the factors of the second
module and use, in each column, the actual exterior multiplication

\[
\Lambda^{h_j}E\otimes\Lambda^{r-h_j}E\longrightarrow\Lambda^rE.
\]

Their tensor product gives (N.13). It is nonzero: pair the highest
vector formed from \(e_1\wedge\cdots\wedge e_{h_j}\) in each first
factor with the lowest vector formed from
\(e_{h_j+1}\wedge\cdots\wedge e_r\) in its paired second factor. The
value is \((e_1\wedge\cdots\wedge e_r)^{\otimes H}\). Irreducibility
makes the resulting pairing perfect.

Torus weights in this pairing add to \((H^r)\). It therefore restricts
to a perfect pairing between the balanced weight \((b^r)\) on the first
factor and the multilinear weight \((1^r)\) on the second. Literal
permutation matrices act on the target by \(\operatorname{sgn}^H\).
Consequently

\[
(S_\lambda E)_{(b^r)}
\simeq\operatorname{sgn}^H\otimes
\bigl((S_\nu E)_{(1^r)}\bigr)^*
\quad\text{as actual }S_r\text{-modules}.
\tag{N.14}
\]

\subsubsection{2. The multilinear Young-symmetrizer
map}\label{the-multilinear-young-symmetrizer-map}

Let \(e_1,\ldots,e_r\) be the basis of \(E\). The map

\[
\mathbf Q[S_r]\longrightarrow(E^{\otimes r})_{(1^r)},
\qquad
\sigma\longmapsto
e_{\sigma(1)}\otimes\cdots\otimes e_{\sigma(r)}
\tag{N.15}
\]

is an isomorphism: the displayed tensors form the entire multilinear
basis. Literal basis permutations act by left multiplication on the
group algebra. Give tensor positions the right action
\((v_1\otimes\cdots\otimes v_r)\rho
=v_{\rho(1)}\otimes\cdots\otimes v_{\rho(r)}\). Then (N.15) also
intertwines right multiplication.

Choose a tableau of shape \(\nu\), take its row symmetrizer and column
antisymmetrizer, and divide their product by its nonzero
Young-symmetrizer scalar to obtain a primitive Young idempotent
\(e_\nu\). The image \((E^{\otimes r})e_\nu\) is the corresponding
Schur-module realization \(S_\nu E\). Restricting (N.15) to the left
ideal \(\mathbf Q[S_r]e_\nu\) thus gives the exact map

\[
[\nu]\simeq\mathbf Q[S_r]e_\nu
\xrightarrow{\ \sim\ }(S_\nu E)_{(1^r)}.
\tag{N.16}
\]

Specht modules over \(\mathbf Q\) are self-dual: average any positive
definite rational form over the finite group and use it to identify the
module with its dual. Equations (N.14)--(N.16) prove

\[
(S_\lambda E)_{(b^r)}\simeq
\operatorname{sgn}^H\otimes[\nu].
\tag{N.17}
\]

The Cauchy balanced-weight identification identifies the highest
\(\lambda\)-multiplicity space of
\(\operatorname{Sym}^r(\operatorname{Sym}^bV)\) with the invariant
subspace of (N.17). Its dimension is one for \(\nu=(r)\) when \(H\) is
even, one for \(\nu=(1^r)\) when \(H\) is odd, and zero for every other
\(\nu\). Taking complements in (N.12) gives exactly (N.11). This proves
the complete source-band description, including all zero multiplicities.

\subsubsection{3. The map on the two actual
bands}\label{the-map-on-the-two-actual-bands}

For even \(H=rL\), the sole source component is the component in R-019.
Formula (N.5) proves nonvanishing, so its multiplicity-one source is
mapped injectively.

For odd \(H=rL\), both \(r\) and \(L\) are odd and \(b=H-1\) is even.
The highest vector of the unique source component is \(D^b\), which is
symmetric under slot permutations. Every term in the coefficient
extraction (N.1) puts \(rL=b+1\) copies of the common vector into \(b\)
determinant factors. At least one factor receives two copies and is zero
by alternating multilinearity. Thus \(T_{r,L,b}(D^b)=0\). Equivariance
kills its entire irreducible component, proving the second line of the
band classification.

No other source component belongs to this band by (N.17). Combined with
R-018, the unclassified boundary components therefore satisfy both

\[
\boxed{\lambda_1\ge rL+1,\qquad
       \lambda_2\ge(r-1)L.}
\tag{N.18}
\]

The determinant component in the odd case is already explicitly
classified and is not included in this remaining range. Statements
R-011, R-012, R-015 and R-016 give additional complete finite kernels;
(N.18) is a uniform consequence of the exact maps proved here.

\section{Canonical-map ranges below the full square
bound}\label{canonical-map-ranges-below-the-full-square-bound}

Statements \textbf{R-021}, \textbf{R-022}, and \textbf{R-023}. Written
and checked on 21 November 2026 UTC by GPT-6 Astra Pro,
\texttt{\small /\allowbreak{}root/\allowbreak{}rectangular\_\allowbreak{}cases};
effort setting not exposed. The prime-interval input is Nagura's
established theorem. The square determinant scalars are the established
prime-adjacent Alon--Tarsi cases in R-007. The exact specialization and
multiplication maps are G-013. The common-polarization inputs are
R-018--R-020. No historical novelty claim is made.

\subsection{Exact canonical-map
convention}\label{exact-canonical-map-convention}

The normalized grid map is

\[
h_{b,a}:\operatorname{Sym}^{b}(\operatorname{Sym}^{a}V)
\longrightarrow\operatorname{Sym}^{a}(\operatorname{Sym}^{b}V),
\qquad
h_{b,a}(v_1^a\cdots v_b^a)=(v_1\cdots v_b)^a.
\tag{K.1}
\]

It is the degree-\(b\) part of the graded algebra map \(\mu_a\) of
R-008--R-009. Symmetric powers are embedded in tensors by normalized
averaging. Taking the transpose of the row and column orthogonal
averaging projectors gives the opposite grid map \(h_{a,b}\). Thus

\[
\mathcal K_{a,b}(W):=\ker\!\left(
h_{a,b}:\operatorname{Sym}^{a}(\operatorname{Sym}^{b}W)
\to\operatorname{Sym}^{b}(\operatorname{Sym}^{a}W)\right)
\tag{K.2}
\]

is, for \(W=V^*\), the annihilator of the image in (K.1). The standard
perfect symmetric-tensor pairings identify these maps without a scalar
change. In particular absence of the label \(\lambda\) from (K.2) proves
surjectivity of (K.1) onto its entire target \(\lambda\)-isotypic
component.

\subsection{R-021. Every even-rank boundary determinant is in the
canonical
image}\label{r-021.-every-even-rank-boundary-determinant-is-in-the-canonical-image}

For every even integer \(a\ge4\), put \(B=a(a-2)\). Then

\[
\boxed{
[S_{(B^a)}V:\operatorname{im}h_{B,a}]
=[S_{(B^a)}V:\operatorname{Sym}^a(\operatorname{Sym}^BV)]
=1\qquad(\dim V\ge a).}
\tag{K.3}
\]

The proof supplies an explicit rational preimage of the normalized
determinant tensor \(I_{a,B}\). In smaller dimensions that Schur functor
is zero.

\subsubsection{1. The exact seed
preimages}\label{the-exact-seed-preimages}

For every odd prime \(p\), R-007 gives nonzero square scalars
\(c_{p-1}\) and \(c_{p+1}\), with

\[
h_{w,w}(I_{w,w})=c_w I_{w,w},
\qquad w=p-1\text{ or }p+1.
\]

For \(w\ge a\), let \(E=\mathbf Q^a\), enlarge it to
\(U=E\oplus\mathbf Q^{w-a}\), and write \(z_{a+1},\ldots,z_w\) for the
added basis covectors. The linear map

\[
j:\operatorname{Sym}^aE^*\longrightarrow\operatorname{Sym}^wU^*,
\qquad F\longmapsto Fz_{a+1}\cdots z_w
\tag{K.4}
\]

has its ordinary dual \(j^*:\operatorname{Sym}^wU\to
\operatorname{Sym}^aE\) under the normalized symmetric pairings. Define

\[
u_{a,w}=\operatorname{Sym}^w(j^*)\bigl(I_{w,w}/c_w\bigr).
\tag{K.5}
\]

The exact equality is

\[
h_{w,a}(u_{a,w})=I_{a,w}.
\tag{K.6}
\]

To check the scalar, let \(\Pi_d\) multiply \(d\) linear forms.
Appending the fixed forms \(z_{a+1},\ldots,z_w\) gives a map \(\iota\)
of ordered factors with \(j\circ\Pi_a=\Pi_w\circ\iota\). Pullback of
this identity gives \(h_{w,a}\operatorname{Sym}^w(j^*)=\iota^*h_{w,w}\).
The determinant with the appended coordinate forms has a block identity
matrix, so \(\iota^*(\det_w^w)=\det_a^w\) with scalar one. These
equations prove (K.6), with every normalization contained in the actual
dual map \(j^*\).

The algebra identities already proved in R-009 and G-013 give

\[
h_{B,a}\left(\prod_w u_{a,w}^{e_w}\right)
=I_{a,B},\qquad B=\sum_w e_w w.
\tag{K.7}
\]

Indeed \(\mu_a\) is an algebra morphism and
\(I_{a,w}I_{a,w'}=I_{a,w+w'}\), as verified by evaluation on arbitrary
ordered covectors.

\subsubsection{\texorpdfstring{2. A prime and explicit nonnegative
exponents for every
\(a\)}{2. A prime and explicit nonnegative exponents for every a}}\label{a-prime-and-explicit-nonnegative-exponents-for-every-a}

For \(a=4,6,\ldots,24\), the following table provides a square seed of
width exactly \(a\):

\[
\begin{array}{c|rrrrrrrrrrr}
a&4&6&8&10&12&14&16&18&20&22&24\\\hline
p&3&5&7&11&11&13&17&17&19&23&23\\
a=p\pm1&+1&+1&+1&-1&+1&+1&-1&+1&+1&-1&+1.
\end{array}
\tag{K.8}
\]

Each displayed \(p\) is prime. Thus \(u_{a,a}^{a-2}\) is a preimage of
\(I_{a,a(a-2)}\) by (K.7).

For every even \(a\ge26\), Nagura's theorem supplies a prime

\[
a<p<\frac65a.
\tag{K.9}
\]

The exact primary locator is Jitsuro Nagura, \emph{On the interval
containing at least one prime number}, \textbf{Proceedings of the Japan
Academy 28} (1952), 177--181, the unnumbered \textbf{Theorem on printed
p.~180, case \(n=5\)}, with proof completed on p.~181. Official full
article:
\url{https://www.jstage.jst.go.jp/article/pjab1945/28/4/28_4_177/_pdf};
DOI \url{https://doi.org/10.3792/pja/1195570997}. Only its range
\(x>25\) is needed here, so no endpoint convention at \(x=25\) affects
this proof.

Put \(g=(p-1)/2\). The two permitted widths \(p-1,p+1\) are at least
\(a\), and their normalized values are \(g,g+1\). Their exact even
conductor is \(2g(g-1)\), as proved in G-013 by the explicit residue
representation. Here the needed inequality follows directly:

\[
2g(g-1)=\frac{(p-1)(p-3)}2
<\frac{p^2}2<\frac{18a^2}{25}<a(a-2)=B,
\tag{K.10}
\]

where the last strict inequality holds for every \(a\ge26\). For
complete explicitness, set

\[
N=B/2,\qquad y=N-g\lfloor N/g\rfloor,\qquad
x=\lfloor N/g\rfloor-y.
\tag{K.11}
\]

Then \(0\le y\le g-1\), and (K.10) implies
\(\lfloor N/g\rfloor\ge g-1\), so \(x\ge0\). The identity
\(N=xg+y(g+1)\) is immediate from (K.11). Consequently

\[
\boxed{
h_{B,a}(u_{a,p-1}^{x}u_{a,p+1}^{y})=I_{a,B}\ne0.}
\tag{K.12}
\]

The target determinant multiplicity is one for even \(B\) by R-001.
Equations (K.8) and (K.12) prove (K.3) for every even \(a\ge4\). This is
a transfer of the checked prime-interval and prime-adjacent scalar
results through explicit algebra morphisms.

\subsubsection{A uniform determinant-image range and its row
extensions}\label{a-uniform-determinant-image-range-and-its-row-extensions}

The same construction gives a uniform range for every integer rank
\(a\ge26\), with no parity restriction on \(a\): for every even

\[
B\ge2\left\lceil\frac{9a^2}{25}\right\rceil,
\qquad s\ge0,
\tag{K.12a}
\]

there is the explicit nonzero image

\[
\boxed{
h_{B+s,a}\left(u_{a,p-1}^{x}u_{a,p+1}^{y}(e_1^a)^s\right)
=I_{a,B}(e_1^{\otimes a})^s\ne0,}
\tag{K.12b}
\]

where \(p\) is any prime from (K.9), and \(x,y\) are given by (K.11)
with this value of \(B\). Indeed \(g(g-1)<9a^2/25\le B/2\), so the same
proof gives nonnegative integer exponents and \(B=x(p-1)+y(p+1)\).
Applying the graded algebra map proves (K.12b) exactly.

The target weight is \(\lambda=(as+B,B^{a-1})\). Its target and image
multiplicities are both one. For the target, the balanced Cauchy module
is \(S_\lambda\mathbf Q^a\simeq\det^B\otimes
\operatorname{Sym}^{as}\mathbf Q^a\). After removing the determinant
weight, its balanced weight is \((s^a)\), with sole monomial
\(x_1^s\cdots x_a^s\). Slot permutations fix this vector because \(B\)
is even. Nonvanishing of (K.12b) follows by evaluating its target on
covectors: it becomes a determinant power times the nonzero product of
first-coordinate powers. These are actual canonical images for the full
range (K.12a).

\subsection{R-022. The canonical surjectivity
ranges}\label{r-022.-the-canonical-surjectivity-ranges}

Let \(a\ge3\), and put \(r=a-1\). The following are unconditional
component ranges for the actual map (K.1), in every finite dimension.

For every integer

\[
r^2-r+1\le b\le r^2-1,
\]

the map is surjective onto every target isotypic component with

\[
\boxed{\lambda\vdash ab,\qquad \lambda_2\le b-r.}
\tag{K.13}
\]

At the upper boundary \(b=r^2-1\), the larger stated range is

\[
\boxed{\lambda\vdash ab,\qquad
\lambda_1\le r^2\quad\text{or}\quad
\lambda_2\le r(r-1)-1.}
\tag{K.14}
\]

Thus the rank on each such multiplicity space equals its full target
plethysm coefficient. In particular the coefficient inequality in that
component follows from this rank statement. The proof below tracks the
annihilator into the common-polarization kernel and then removes its
rectangular summand by an exact prolongation identity.

\subsubsection{3. An explicit injection from the canonical
annihilator}\label{an-explicit-injection-from-the-canonical-annihilator}

For all \(b\) in the stated interval, let \(m=b-r\). Then

\[
m\ge(r-1)^2\ge1.
\tag{K.15}
\]

Sra, \emph{The q-Foulkes problem III: On Foulkes' conjecture with a
graded inner factor}, v1, \textbf{Theorem 1.3}, gives surjectivity of
\(h_{m,r}\). Its complete proof in §2.2 and its power-polarization Lemma
2.1 were read in this run. The author's primary PDF is
\url{https://suvrit.github.io/papers/qfoulkes3/qfoulkes3-v1.pdf}. For
\(r=2,m=1\), this lower map is directly the identity
\(\operatorname{Sym}^2V\to\operatorname{Sym}^2V\), as the source also
explicitly records.

Let \(W=V^*\) and \(F\in\mathcal K_{a,b}(W)\). Its corresponding
symmetric \(a\)-linear form on \(\operatorname{Sym}^bV\) is denoted
\(\mathcal A_F\). The annihilator condition is the exact identity

\[
\mathcal A_F(P,\ldots,P)=0
\quad\text{for every }P=u_1\cdots u_b.
\tag{K.16}
\]

Indeed (K.1) has image spanned by precisely these diagonal tensors.

Fix a product \(Q=y_1\cdots y_m\). The polynomial \((x^r+c\,t^r)Q\) is a
product of \(b\) vectors over \(\mathbf C\), because its binary
degree-\(r\) factor splits into linear factors. The coefficient of \(c\)
in (K.16), divided by the nonzero integer \(a\), gives

\[
\mathcal A_F(x^rQ,\ldots,x^rQ,t^rQ)=0.
\tag{K.17}
\]

Define \(G_x(y_1,\ldots,y_m,t)=
\mathcal A_F(x^rQ,\ldots,x^rQ,t^b)\). It has degree \(b\) in \(t\) and
degree \(r\) in each \(y_i\). The exact differential identity is

\[
D_{y_m,t}\cdots D_{y_1,t}G_x
=\frac{b!}{r!}
\mathcal A_F(x^rQ,\ldots,x^rQ,t^rQ)=0.
\tag{K.18}
\]

Before the \(i\)-th differentiation, the relevant degrees are
\(p=b-i+1\) and \(s=r\), so \(p-s=m-i+1\ge1\). Sra's Lemma 2.1 with
power one proves injectivity of each step, with all other variables as
spectators. Hence \(G_x=0\).

For fixed \(x,t\), this vanishing is the diagonal of an \(r\)-linear
symmetric form on \(\operatorname{Sym}^mV\), evaluated on all products
\(Q\). The already established surjectivity of \(h_{m,r}\) therefore
gives

\[
\mathcal A_F(x^rQ_1,\ldots,x^rQ_r,t^b)=0
\quad(Q_i\in\operatorname{Sym}^mV).
\tag{K.19}
\]

Represent \(F\) by its polynomial \(f(v_1,\ldots,v_r;t)=
\mathcal A_F(v_1^b,\ldots,v_r^b,t^b)\). The exact coefficient-extraction
calculation now is

\[
(T_{r,r,b}\otimes\mathrm{id})f
=\binom{b}{r}^{r}
\mathcal A_F(x^rv_1^m,\ldots,x^rv_r^m,t^b)=0.
\tag{K.20}
\]

No injectivity assertion about the last common-polarization step has
been inserted into this argument.

The normalized last-slot inclusion is the natural injective map

\[
\begin{aligned}
\delta:\operatorname{Sym}^{r+1}(\operatorname{Sym}^bW)
&\longrightarrow
\operatorname{Sym}^r(\operatorname{Sym}^bW)
\otimes\operatorname{Sym}^bW,\\
w_0\cdots w_r&\longmapsto
\frac1{r+1}\sum_{i=0}^r
\left(\prod_{j\ne i}w_j\right)\otimes w_i.
\end{aligned}
\tag{K.21}
\]

Multiplying the two tensor factors is a left inverse, proving
injectivity. In the invariant-tensor realization (K.21) is simply
inclusion of the fully symmetric tensors into the tensors symmetric in
the first \(r\) slots. Equation (K.20) consequently proves the actual
natural injection

\[
\boxed{
\delta:\mathcal K_{a,b}(W)\hookrightarrow
\ker T_{r,r,b}(W)\otimes\operatorname{Sym}^bW.}
\tag{K.22}
\]

\subsubsection{4. Proof of the second-row
range}\label{proof-of-the-second-row-range}

R-018 proves that every constituent \(S_\nu W\) of \(\ker T_{r,r,b}\)
has \(\nu_2\ge b-r+1\). Pieri's rule for tensoring by
\(\operatorname{Sym}^bW\) gives

\[
\lambda_1\ge\nu_1\ge\lambda_2\ge\nu_2
\ge\lambda_3\ge\cdots.
\tag{K.23}
\]

In particular \(\lambda_2\ge\nu_2\). The injection (K.22) proves that no
annihilator component can have \(\lambda_2\le b-r\). By the exact
duality (K.2), this is precisely the component surjectivity assertion
(K.13).

\subsection{R-023. An exact rectangular prolongation
identity}\label{r-023.-an-exact-rectangular-prolongation-identity}

Let \(r\ge2\), let \(b\ge0\) be even, and set \(a=r+1\). Regard the
unique source copy of \(S_{(b^r)}W\) inside
\(\operatorname{Sym}^r(\operatorname{Sym}^bW)\) using the determinant
highest vector and its \(GL(W)\)-span. Inside the target of (K.21), the
following is an equality of actual submodules:

\[
\boxed{
\delta\bigl(\operatorname{Sym}^{a}(\operatorname{Sym}^bW)\bigr)
\cap\bigl(S_{(b^r)}W\otimes\operatorname{Sym}^bW\bigr)
=\delta\bigl(S_{(b^a)}W\bigr).}
\tag{K.24}
\]

\subsubsection{Proof}\label{proof-3}

Use the polynomial-generator realization with \(a\) vector slots. The
Cauchy decomposition for the first \(r\) slots has its \(GL(W)\)-type
\((b^r)\) paired with \(S_{(b^r)}\mathbf Q^r=(\det\mathbf Q^r)^b\).
Consequently every polynomial in the left-hand intersection is invariant
under every elementary transvection mixing two of the first \(r\) slots.
In coordinates, it is killed by all the slot derivations

\[
E_{ij}^{\mathrm{slot}}
=\sum_u X_{u,i}\frac{\partial}{\partial X_{u,j}},
\qquad i\ne j,\quad1\le i,j\le r.
\tag{K.25}
\]

The polynomial is fully symmetric in all \(a\) slots. Conjugating (K.25)
by slot permutations therefore shows that it is killed by the
corresponding derivation for every ordered pair \(i\ne j\) among all
\(a\) slots. Exponentiation of these locally nilpotent derivations
proves invariance under \(SL_a\) acting on the slot space. The Cauchy
decomposition for all \(a\) slots identifies its invariant part: only
the determinant characters on the slot factor can restrict trivially to
\(SL_a\). The total degree is \(ab\), so that character is exactly
\(\det^b\), paired with \(S_{(b^a)}W\). This proves inclusion in the
right side of (K.24).

Conversely, that Cauchy summand has slot action \(\det^b\). It is fixed
by all slot permutations because \(b\) is even. Restriction to the first
\(r\) slots is again \(\det^b\), so its image under (K.21) lies in the
displayed rectangular first factor. This proves the reverse inclusion
and the equality. The argument works in every dimension; a Schur functor
exceeding the ambient dimension contributes zero on both sides.

\subsubsection{5. The sharper critical-boundary injection and proof of
(K.14)}\label{the-sharper-critical-boundary-injection-and-proof-of-k.14}

Now take \(b=r^2-1\). For even \(r\), R-018--R-020 prove that every
constituent \(S_\nu W\) of the common-polarization kernel satisfies

\[
\nu_1\ge r^2+1,\qquad \nu_2\ge r(r-1).
\tag{K.26}
\]

For odd \(r\), those same bounds hold after removing the unique
rectangular summand \(S_{(b^r)}W\), which R-020 proves is in the kernel.
Let \(\mathcal N_r\) denote the common kernel in the even case, and this
complementary sum of isotypic components in the odd case. The projection
to \(\mathcal N_r\) is explicit and natural: on tensor degree \(rb\),
remove the central Schur projector

\[
e_{(b^r)}=
\frac{\dim[(b^r)]}{(rb)!}
\sum_{\sigma\in S_{rb}}\chi^{(b^r)}(\sigma)\sigma
\tag{K.27}
\]

in the odd case; use the identity in the even case. Centrality makes
this projector commute with the tensor symmetrizations defining the
source and the common-polarization map.

For odd \(r\), the composite of (K.22) with this projection can have
kernel only in the intersection (K.24), hence only in \(S_{(b^a)}W\).
Here \(a=r+1\) is even and \(b=r^2-1=a(a-2)\). R-021 proves that the
original canonical map is nonzero on its multiplicity-one target
determinant component. By (K.2), its annihilator has zero intersection
with \(S_{(b^a)}W\). The composite projection is therefore injective. In
the even case injectivity already follows from (K.22). Thus in both
parities there is the exact natural injection

\[
\boxed{
\mathcal K_{r+1,r^2-1}(W)
\hookrightarrow\mathcal N_r(W)\otimes
\operatorname{Sym}^{r^2-1}W,}
\tag{K.28}
\]

given explicitly by (K.21) and (K.27), whose first factor has only the
types (K.26).

Pieri interlacing (K.23) now forces every annihilator label to have
\(\lambda_1\ge r^2+1\) and \(\lambda_2\ge r(r-1)\). This is exactly the
complement of the surjectivity range (K.14). The proof removes the
rectangular contraction by its actual prolongation and canonical
preimage, rather than discarding it by a parity assertion.

\subsubsection{6. Exact integer-interval description of the candidate
support}\label{exact-integer-interval-description-of-the-candidate-support}

For completeness, pad \(\lambda\) to \(a=r+1\) parts. A partition
\(\nu\vdash rb\) satisfying both (K.26) and the Pieri interlacing (K.23)
exists exactly for the following explicit integer region:

\[
\begin{aligned}
\lambda_1&\ge b+2,\\
\lambda_2&\ge b-r+1,\\
\max(0,b+2-\lambda_2)+\max(0,b-r+1-\lambda_3)
&\le\lambda_1-b.
\end{aligned}
\tag{K.30}
\]

This gives a necessary support region for the canonical annihilator via
(K.28), not an assertion that every label in it occurs. To prove the
equivalence, choose independently \(\nu_i\in[\lambda_{i+1},\lambda_i]\),
\(1\le i\le r\), raising the first two lower bounds to \(b+2,b-r+1\).
Their individual feasibility is precisely the first two lines. Ordering
is automatic because \(\nu_i\ge\lambda_{i+1}\ge\nu_{i+1}\). The sum of
the lower bounds is \(ab-\lambda_1\) plus the two maxima in (K.30), and
the sum of the upper bounds is \(ab-\lambda_{r+1}\ge rb\), since
\(\lambda_{r+1}\le b\). Sums of independent integer intervals fill every
integer between their lower and upper sums. Existence of sum \(rb\) is
therefore exactly the last line of (K.30).

For \(r\ge3\), that last inequality follows already from the first two
and the partition size. Here is a complete check, which prevents
counting it as an additional range in those ranks. Set
\(x=\lambda_1-b\), \(y=\lambda_2-b\), and \(z=\lambda_3-b\). Then

\[
x\ge2,\quad y\ge1-r,\quad z\le y,\quad
x\ge y,\quad x\ge-y-(r-1)z.
\]

The desired inequality is \(\max(0,2-y)+\max(0,1-r-z)\le x\). For
\(y\ge2,z\ge1-r\), its left side is zero. For \(y\ge2,z=-q\le-r\), the
two lower bounds on \(x\) give \(x\ge(r-1)q/2\ge q-r+1\). For
\(y\le1,z\ge1-r\), the assertion is \(x+y\ge2\): this follows from
\(x\ge2\) at \(y=0,1\), and from \(x+y\ge-(r-1)z\ge-(r-1)y\ge r-1\ge2\)
at \(y\le-1\). Finally, for \(y\le1,z=-q\le-r\), the lower bound
\(x\ge(r-1)q-y\) exceeds the required \(q-y-r+3\) by \((r-2)q+r-3\ge0\).
These cases exhaust the integer possibilities.

At \(r=2,b=3\), no such \(\nu\) exists at all: (K.26) would give
\(\nu_1+\nu_2\ge5+2>rb=6\). Thus (K.28) proves that the full map
\(h_{3,3}\) is an isomorphism. This edge case agrees with the classical
result of Hadamard, stated with its complete proof in Landsberg,
\emph{Geometric complexity theory: an introduction for geometers},
\href{https://arxiv.org/pdf/1305.7387}{arXiv:1305.7387v3},
\textbf{Theorem 7.22, §7.8, printed p.~24}. That statement and proof
were independently read here.

\subsection{Parameter checks and concrete
families}\label{parameter-checks-and-concrete-families}

The smallest parameter in (K.13) is \(a=3,r=2,b=3\), where \(m=b-r=1\),
the lower grid map is the identity, and the power-one injectivity in
(K.18) has degree excess one. The critical second-row range is
\(\lambda_2\le1\), and the first-row range is \(\lambda_1\le4\). The
sole nonzero target component in that first-row band is \((4,4,1)\), by
the complement/Specht calculation of R-020 with \(a=3,H=4\). Thus no
unavailable induction base or negative residual degree is used.

At the critical boundary, the complete first-row target band itself has
a particularly simple description. Applying the same explicit
complementary pairing and multilinear map from R-020 with \(a\) slots
and \(H=b+1\) gives

\[
\bigoplus_{\lambda_1\le b+1}
\left(\operatorname{Sym}^a(\operatorname{Sym}^bV)\right)_\lambda
=
\begin{cases}
S_{((r^2)^r,r^2-r-1)}V,&r\text{ even},\\
S_{((r^2-1)^{r+1})}V,&r\text{ odd}.
\end{cases}
\tag{K.29}
\]

Each displayed component has multiplicity one and is entirely in the
canonical image by (K.14). Examples are \((a,b)=(5,15)\) with
\(\lambda=(16,16,16,16,11)\), and \((a,b)=(7,35)\) with
\(\lambda=(36,36,36,36,36,36,29)\). In even outer degree the determinant
examples are \((a,b)=(4,8),(6,24),(8,48)\), and (K.3) covers every even
\(a\ge4\).

The full canonical map remains governed by the remaining components in
(K.28). The statements above prove the displayed entire isotypic ranges,
without asserting vanishing of the remaining varying common-polarization
kernel.

\subsubsection{Explicit positive ranks on equal-middle-row
families}\label{explicit-positive-ranks-on-equal-middle-row-families}

The component surjectivity combines directly with the already proved
T-025 end-shape formula. For every \(a\ge3\), \(r=a-1\),
\(r^2-r+1\le b\le r^2-1\), and \(0\le\gamma\le\beta\le b-r\), put

\[
\lambda=(ab-(a-2)\beta-\gamma,\beta^{a-2},\gamma),
\qquad R=\beta-\gamma-\binom a2(\beta\bmod2).
\]

This is a partition: its first part is at least
\(ab-(a-1)\beta\ge b+r^2\ge\beta\). Its second part satisfies the exact
range (K.13), and \(b\ge\beta\). Therefore its actual canonical image
multiplicity is

\[
\boxed{
[S_\lambda V:\operatorname{im}h_{b,a}]
=\#\{(c_2,\ldots,c_a)\in\mathbf Z_{\ge0}^{a-1}:
2c_2+3c_3+\cdots+ac_a=R\}.}
\tag{K.31}
\]

For dimensions at least the length of \(\lambda\), this count is
positive exactly for \(R=0\) or \(R\ge2\), because parts two and three
already represent every integer at least two. Negative right sides and
\(R=1\) give zero. Examples are \(a=4,b=7\), \(\lambda=(20,4,4)\), with
image multiplicity two, and \(a=5,b=13\), \(\lambda=(41,8,8,8)\), with
image multiplicity five. Formula (K.31) records the exact map rank of
the known positive coefficient formula throughout these concrete
parameter families.

\section{A recursive second-row range for the canonical
map}\label{a-recursive-second-row-range-for-the-canonical-map}

Statement \textbf{R-024}. Written and checked on 21 November 2026 UTC by
GPT-6 Astra Pro, agent /root/rectangular\_cases; effort setting not
exposed. The canonical degree-five base is the established result of
Cheung--Ikenmeyer--Mkrtchyan and McKay. The component argument below
uses the explicit common-polarization injectivity R-018. No historical
novelty claim is made.

\subsection{1. The unconditional range and its exact
map}\label{the-unconditional-range-and-its-exact-map}

Let \(V\) be any finite-dimensional vector space over a field of
characteristic zero. For every integer \(a\ge6\), set

\[
C_a=\binom a2-10=5+6+\cdots+(a-1).
\tag{P.1}
\]

For every \(b\ge C_a+6\), and every partition

\[
\lambda\vdash ab,\qquad \lambda_2\le b-C_a,
\tag{P.2}
\]

the canonical map

\[
h_{b,a}:\operatorname{Sym}^{b}(\operatorname{Sym}^{a}V)
\longrightarrow\operatorname{Sym}^{a}(\operatorname{Sym}^{b}V),
\qquad
h_{b,a}(v_1^a\cdots v_b^a)=(v_1\cdots v_b)^a,
\tag{P.3}
\]

is surjective onto the entire target \(\lambda\)-isotypic component. The
assertion includes all multiplicities and all dimensions; a Schur
functor with too many rows evaluates to zero. A missing second part is
zero. Equivalently, the explicit parameter range is

\[
\boxed{
a\ge6,\qquad b\ge\binom a2-4,\qquad
\lambda_2\le b-\binom a2+10.}
\tag{P.4}
\]

For example the ranges for \(a=6,7,8,9\) respectively are

\[
\begin{array}{c|c|c}
a&\text{inner degree}&\text{second row}\\\hline
6&b\ge11&\lambda_2\le b-5\\
7&b\ge17&\lambda_2\le b-11\\
8&b\ge24&\lambda_2\le b-18\\
9&b\ge32&\lambda_2\le b-26.
\end{array}
\tag{P.5}
\]

All symmetric powers in (P.3) use normalized averaging in the tensor
realization. Each averaging projector is self-adjoint for the tensor
pairing. Writing the grid map as target averaging composed with grid
permutation on the source invariant subspace, its transpose is source
averaging composed with the inverse permutation on the target invariant
subspace. This is the opposite grid map, with scalar one. Therefore, for
\(W=V^*\), the image annihilator is exactly

\[
\mathcal K_{a,b}(W)
=\ker\left[h_{a,b}:\operatorname{Sym}^{a}(\operatorname{Sym}^{b}W)
\longrightarrow\operatorname{Sym}^{b}(\operatorname{Sym}^{a}W)\right].
\tag{P.6}
\]

We prove that its \(\lambda\)-isotypic component is zero in (P.4). This
proves the asserted surjectivity by the perfect symmetric-tensor
pairing, including the map on each multiplicity space.

\subsection{2. The established full degree-five
base}\label{the-established-full-degree-five-base}

The exact input is

\[
\mathcal K_{5,b}(W)=0\qquad\text{for every }b\ge6.
\tag{P.7}
\]

Cheung--Ikenmeyer--Mkrtchyan, \emph{Symmetrizing tableaux and the 5th
case of the Foulkes conjecture}, arXiv:1509.03944v1, \textbf{Theorem
6(a), printed p.~3}, proves injectivity of \(\Psi_{5,6}\). Their
displayed definition on printed p.~2 is the composition of the
symmetric-tensor inclusion, the transpose of the \(5\times6\) tensor
grid, and symmetric projection. It is precisely \(h_{5,6}\) under the
convention (P.3). Applying McKay's propagation theorem to this
established base gives (P.7). A complete primary proof of that
propagation is Christian Ikenmeyer, \emph{On McKay's propagation theorem
for the Foulkes conjecture}, arXiv:1509.04957, \textbf{Theorem 1.3,
printed p.~2}, with proof in §§2--6, printed pp.~3--8. Its definition of
\(\Psi_{a\times b}\), also on printed p.~2, is the same inclusion, grid
transpose, and projection, so its conclusion applies to these actual
maps. The official full texts are

\begin{itemize}
\tightlist
\item
  \url{https://arxiv.org/pdf/1509.03944}
\item
  \url{https://arxiv.org/pdf/1509.04957}.
\end{itemize}

The source input is injectivity of the precise transpose map in (P.6).
The source computation for Theorem 6(a) is an exact integer computation
described and justified in §§2--5; its original large evaluation
matrices were not regenerated in this run. The propagation proof was
read in full. The scalar audit in §6 below records the harmless
normalization corrections in its auxiliary raising-operator displays.

The sources work in sufficiently large dimension. Their injectivity
implies (P.7) in every smaller dimension by the actual inclusion
\(W\hookrightarrow W\oplus W'\): both symmetric-power functors preserve
this inclusion, and the grid maps commute with it. A vector in the
smaller kernel would map to a vector in the zero larger kernel. All maps
are rational, so extension to characteristic-zero fields preserves the
assertion.

\subsection{3. An exact intermediate-annihilator
morphism}\label{an-exact-intermediate-annihilator-morphism}

We now construct the map needed for the induction. For every

\[
a=r+1\ge3,\qquad b\ge r,\qquad m=b-r\ge0,
\tag{P.8}
\]

let

\[
\begin{aligned}
\delta:\operatorname{Sym}^{r+1}(\operatorname{Sym}^{b}W)
&\longrightarrow
\operatorname{Sym}^{r}(\operatorname{Sym}^{b}W)
\otimes\operatorname{Sym}^{b}W,\\
w_0\cdots w_r&\longmapsto
\frac1{r+1}\sum_{i=0}^{r}
\left(\prod_{j\ne i}w_j\right)\otimes w_i.
\end{aligned}
\tag{P.9}
\]

Multiplication of the two tensor factors is a left inverse. Thus
\(\delta\) is injective. In the normalized invariant-tensor model it is
the inclusion of fully symmetric tensors into tensors symmetric in the
first \(r\) slots.

The common-polarization map is the coefficient map

\[
T_{r,r,b}f(v_1,\ldots,v_r;x)
=[t_1^r\cdots t_r^r]
f(v_1+t_1x,\ldots,v_r+t_rx).
\tag{P.10}
\]

It has target \(\operatorname{Sym}^{r^2}W\otimes
\operatorname{Sym}^{r}(\operatorname{Sym}^{m}W)\). There is an actual
restriction of the composite map

\[
\boxed{
\Theta_{a,b}:=(T_{r,r,b}\otimes\mathrm{id})\delta:
\mathcal K_{a,b}(W)\longrightarrow
\operatorname{Sym}^{r^2}W\otimes
\mathcal K_{r,m}(W)\otimes\operatorname{Sym}^{b}W.}
\tag{P.11}
\]

The codomain in (P.11) retains the full lower annihilator. Here is a
complete proof of this membership.

Work over \(\mathbf C\), which suffices because the map is rational. For
\(F\in\mathcal K_{a,b}(W)\), denote its corresponding symmetric
\(a\)-linear form on \(\operatorname{Sym}^{b}V\) by \(\mathcal A_F\).
The kernel condition (P.6) says exactly

\[
\mathcal A_F(P,\ldots,P)=0
\quad\text{for all }P=u_1\cdots u_b.
\tag{P.12}
\]

Fix \(Q=y_1\cdots y_m\). The form \((x^r+c\,t^r)Q\) is a product of
\(b\) linear factors over \(\mathbf C\): the degree-\(r\) binary factor
splits, also when \(x,t\) are linearly dependent. Take the coefficient
of \(c\) in (P.12) and divide by \(a\). The result is

\[
\mathcal A_F(x^rQ,\ldots,x^rQ,t^rQ)=0.
\tag{P.13}
\]

Put

\[
G_x(y_1,\ldots,y_m;t)
=\mathcal A_F(x^rQ,\ldots,x^rQ,t^b).
\]

For \(D_{y,t}=y\cdot\partial_t\), ordinary differentiation gives the
exact equation

\[
D_{y_m,t}\cdots D_{y_1,t}G_x
=\frac{b!}{r!}
\mathcal A_F(x^rQ,\ldots,x^rQ,t^rQ)=0.
\tag{P.14}
\]

Before step \(i\), the degrees in \(t,y_i\) are \(b-i+1,r\). Their
difference is \(m-i+1\ge1\). The operator \(D_{y_i,t}\) is injective on
every such space. To verify the coefficient, put \(p=b-i+1\), \(q=r\),
and decompose under the \(GL_2\) action mixing \(t,y_i\). A summand
containing weight \((p,q)\) has the form
\(\det^j\otimes\operatorname{Sym}^{p+q-2j}\mathbf C^2\), with
\(j\le q<p\). Its weight vector is the determinant factor times
\(e_t^{p-j}e_y^{q-j}\). The ordinary raising operator sends it to
\((p-j)e_t^{p-j-1}e_y^{q-j+1}\) times the same determinant factor, which
is nonzero. This verifies injectivity on every multiplicity space as
well. Equivalently this is Sra's Lemma 2.1 with power one. All other
variables are spectators. Starting with the last map in the composition
and successively using these injective restrictions gives \(G_x=0\). For
\(m=0\), (P.13) already states \(G_x=0\), with no derivatives; the
degree-zero grid map is the identity of \(\mathbf Q\).

For fixed \(x,t\), the symmetric \(r\)-linear form

\[
(Q_1,\ldots,Q_r)\longmapsto
\mathcal A_F(x^rQ_1,\ldots,x^rQ_r,t^b)
\tag{P.15}
\]

has zero diagonal on every \(Q=y_1\cdots y_m\). By the exact definition
(P.6), with \(a,b\) replaced by \(r,m\), it belongs to
\(\mathcal K_{r,m}(W)\). The coefficient calculation in the original
vector slots is

\[
\Theta_{a,b}F(v_1,\ldots,v_r;x,t)
=\binom br^{r}
\mathcal A_F(x^rv_1^m,\ldots,x^rv_r^m,t^b).
\tag{P.16}
\]

Its degrees in \(x,t\) are \(r^2,b\). Evaluation at all \(x,t\)
separates both symmetric-power spaces. Applying the lower grid map to
the middle factor of (P.16) therefore gives zero as a tensor, because it
gives zero under every such evaluation by (P.15). This proves precisely
the membership (P.11). The factors \(b!/r!\) and \(\binom br^r\) in
(P.14) and (P.16) also fix all normalizations.

\subsection{4. Proof of the recursive component
range}\label{proof-of-the-recursive-component-range}

Set \(C_5=0\), so that for every \(a\ge6\), with \(r=a-1\),

\[
C_a=r+C_r.
\tag{P.17}
\]

Prove (P.4) by induction on \(a\), using the full base (P.7). Fix
\(a\ge6\), \(b\ge C_a+6\), and a partition \(\lambda\) satisfying (P.2).
Write \(s=b-C_a\ge6\) and \(m=b-r\). Then

\[
m\ge C_r+6,\qquad m-C_r=s,\qquad s\le m.
\tag{P.18}
\]

The last inequality uses \(C_r\ge0\), including \(r=5\).

First restrict (P.9) to the entire \(\lambda\)-isotypic component.
Pieri's rule says that every type \(\nu\) in its first tensor factor
satisfies \(\lambda/\nu\) horizontal of size \(b\), and in particular

\[
\nu_2\le\lambda_2\le s\le m=b-r.
\tag{P.19}
\]

R-018, applied with \(L=r\), proves that \(T_{r,r,b}\) is injective on
every entire \(\nu\)-isotypic component in (P.19). Their direct sum is
therefore an injective restriction too: images of distinct types lie in
distinct isotypic components of its target, and the statement for one
type already includes its full multiplicity space. Tensoring by
\(\operatorname{Sym}^{b}W\) preserves this injectivity. Since \(\delta\)
itself is injective, the composite in (P.11) is injective on the
\(\lambda\)-isotypic component of \(\mathcal K_{a,b}(W)\).

Next consider its codomain in (P.11). At \(r=5\) the middle factor is
zero by (P.7) and \(m\ge6\). At every \(r\ge6\), the induction statement
and (P.18) give

\[
\mathcal K_{r,m}(W)
\subseteq
\bigoplus_{\kappa_2\ge s+1}
\left(\operatorname{Sym}^{r}(\operatorname{Sym}^{m}W)\right)_\kappa.
\tag{P.20}
\]

Tensoring one of these types with the two row representations
\(\operatorname{Sym}^{r^2}W\) and \(\operatorname{Sym}^{b}W\) applies
Pieri twice. Each added horizontal strip only increases row lengths, so
every resulting type \(\xi\) contains \(\kappa\), and
\(\xi_2\ge\kappa_2\ge s+1\). The codomain of (P.11) consequently has no
\(\lambda\)-isotypic component with \(\lambda_2\le s\).

The injective restriction just proved thus maps the specified
\(\lambda\)-component into zero. That component is zero. Equation (P.6)
proves the claimed canonical surjectivity, completing the induction for
every \(a\ge6\).

The proof uses whole isotypic components at both stages and retains
every copy of each Schur functor.

\subsection{5. Explicit positive ranks in the new parameter
range}\label{explicit-positive-ranks-in-the-new-parameter-range}

Here is an unconditional family with a closed positive counting formula.
Take the parameters in (P.4), let

\[
0\le\gamma\le\beta\le b-C_a,\qquad
\lambda=(ab-(a-2)\beta-\gamma,\beta^{a-2},\gamma),
\qquad
R=\beta-\gamma-\binom a2(\beta\bmod2).
\tag{P.21}
\]

This is a partition. Its first part is at least
\(ab-(a-1)\beta\ge b+(a-1)C_a\ge\beta\). The exact end-sector formula
T-025 applies, since \(b\ge\beta\). It computes the target multiplicity
as

\[
p_{2,\ldots,a}(R)
=\#\{(c_2,\ldots,c_a)\in\mathbf Z_{\ge0}^{a-1}:
             2c_2+3c_3+\cdots+ac_a=R\}.
\tag{P.22}
\]

The set is empty for negative \(R\). Combining (P.4) with this exact
target multiplicity gives the actual rank on multiplicity spaces:

\[
\boxed{
\operatorname{rank}\operatorname{Hom}_{GL(V)}(S_\lambda V,h_{b,a})
=[S_\lambda V:\operatorname{im}h_{b,a}]
=p_{2,\ldots,a}(R).}
\tag{P.23}
\]

Here \(\dim V\ge\ell(\lambda)\); in smaller dimensions the Schur functor
vanishes. Because \(2,3\) are permitted parts, the number in (P.22) is
positive exactly for \(R=0\) or \(R\ge2\). The construction \(R=2u\), or
\(R=3+2u\), supplies a solution for each such positive \(R\).

For every \(a\ge6\) and every \(b\ge\binom a2-4\), the specific
partition

\[
\lambda=(ab-6a+12,6^{a-2})
\tag{P.24}
\]

has canonical image multiplicity \(4\), given by the four solutions of
(P.22) at \(R=6\): the partitions \(6\), \(4+2\), \(3+3\), and
\(2+2+2\). In particular \(h_{11,6}\) is surjective onto all four copies
of \(S_{(42,6,6,6,6)}V\) in its target.

\subsection{6. Source-normalization
audit}\label{source-normalization-audit}

Ikenmeyer's appendix defines its auxiliary tensor raising operator on
degree \(d\) by \(\phi_{ij}=d^{-1}\sum_{p=1}^{d}\zeta_{ij}^{(p)}\).
Consequently its two-operator composition is exactly

\[
\phi_{ij}\phi_{i'j'}
=d^{-2}\sum_{p\ne q}\zeta_{ij}^{(p)}\zeta_{i'j'}^{(q)};
\tag{P.25}
\]

the printed factor \(1/[d(d-1)]\) on p.~8 needs this correction. The sum
commutes under interchange of the two operators. Removing fixed tensor
positions in §4 changes the remaining raising product by the nonzero
scalar \(((b-1)/b)^{a(b-1)}\). Its injectivity is unchanged. In the §5
\(GL_2\) weight-vector display, the coefficient for the degree-\(d\)
normalization is \((b-\eta_2)/d\), which is nonzero in the displayed
range. Thus every injectivity step of the source propagation proof
survives with exact scalars. The grid map itself is fixed by its
independent inclusion--permutation--projection definition, already
identified with (P.3).

\subsection{7. One further degree using the exact square
kernel}\label{one-further-degree-using-the-exact-square-kernel}

For every \(a\ge6\), at the additional degree

\[
b=\binom a2-5=C_a+5,
\tag{P.26}
\]

the canonical map \(h_{b,a}\) is surjective on every target component
with \(\lambda_2\le4\). This is a separate unconditional boundary
family.

The source input is the exact decomposition in
Cheung--Ikenmeyer--Mkrtchyan, \textbf{Theorem 6(b), printed p.~3}, at
the same primary URL in §2:

\[
\begin{aligned}
\mathcal K_{5,5}={}&
S_{(14,7,2,2)}\oplus S_{(13,7,2,2,1)}
\oplus S_{(12,7,3,2,1)}\oplus S_{(12,6,3,2,2)}\\
&\oplus S_{(12,5,4,3,1)}\oplus S_{(11,5,4,4,1)}
\oplus S_{(10,8,4,2,1)}\oplus S_{(9,7,6,3)}.
\end{aligned}
\tag{P.27}
\]

Each summand has second row at least five, so the full degree-five base
has zero kernel on every type with second row at most four. At an
induction step \(a=r+1\ge6\), (P.26) gives \(m=b-r=C_r+5\ge5\). Thus
(P.19) holds with the fixed bound four, and the composite (P.11) is
injective on every such source component. The induction statement at
\(r\) says that the middle factor \(\mathcal K_{r,m}\) has only labels
with second row at least five; for \(r=5\) this is exactly (P.27). The
two Pieri steps in §4 keep that lower bound. The codomain of (P.11)
therefore has no label with second row at most four, and its injective
restriction has zero source. This proves the new boundary assertion at
every \(a\ge6\).

In particular every \(a\ge6\), with \(b\) as in (P.26), has actual
canonical image multiplicity two at

\[
\lambda=(ab-4a+8,4^{a-2}).
\tag{P.28}
\]

The exact end-sector count is \(p_{2,\ldots,a}(4)=2\), corresponding to
\(4\) and \(2+2\). All parts in (P.28) satisfy the established row
bound, so both target copies are in the canonical image.

\subsection{8. An unconditional triangular profile from the full square
bases}\label{an-unconditional-triangular-profile-from-the-full-square-bases}

For every pair of integers \(a>k\ge6\), the following is an additional
explicit canonical component range:

\[
\boxed{
b\ge\binom a2+\binom{k-1}{2},\qquad
\lambda\vdash ab,\qquad
\lambda_2\le b-\binom a2+\binom k2.}
\tag{P.29}
\]

The map \(h_{b,a}\) is surjective on the entire indicated target
component for every such triple \(a,k,b\).

The exact established lower maps are

\[
\mathcal K_{k,s}(W)=0\qquad
\text{for every }k\ge6,\ s\ge(k-1)^2.
\tag{P.30}
\]

These are the transposes of the surjective grid maps in Sra, \emph{The
q-Foulkes problem III: On Foulkes' conjecture with a graded inner
factor}, v1, \textbf{Theorem 1.3}, proved in §2.2 with its
power-polarization Lemma 2.1. The primary text is
\url{https://suvrit.github.io/papers/qfoulkes3/qfoulkes3-v1.pdf}. The
exact conventions and complete proof of this input were already checked
for R-022. Equation (P.30) is a statement about the actual maps, with no
coefficient-only replacement.

Here is the full induction proving (P.29). Fix \(k\ge6\) and
\(s\ge(k-1)^2\). For every \(j\ge k\), put

\[
b_j=s+\binom j2-\binom k2.
\tag{P.31}
\]

At \(j=k\) the annihilator is zero by (P.30). Suppose \(j>k\); write
\(r=j-1\). The exact numerical identities are

\[
b_j-r=b_r\ge s\ge(k-1)^2,\qquad b_j\ge r.
\tag{P.32}
\]

Consider a type \(\lambda\vdash j b_j\) with \(\lambda_2\le s\). The
first tensor-factor labels \(\nu\) in its image under (P.9) satisfy
\(\nu_2\le\lambda_2\le s\le b_r\), by Pieri. R-018 consequently makes
\(T_{r,r,b_j}\) injective on every such whole \(\nu\)-isotypic
component. The composite \(\Theta_{j,b_j}\) is injective on the original
\(\lambda\)-component.

By (P.11), this composite has codomain

\[
\operatorname{Sym}^{r^2}W\otimes
\mathcal K_{r,b_r}(W)\otimes\operatorname{Sym}^{b_j}W.
\tag{P.33}
\]

At \(r=k\) this codomain is zero. For \(r>k\), the induction assertion
states that every label \(\kappa\) of its middle factor has
\(\kappa_2\ge s+1\). Tensoring with either of the displayed row
representations adds a horizontal strip, so every resulting label still
has second row at least \(s+1\). Thus (P.33) has no type \(\lambda\)
under consideration. The injective restriction has zero domain. This
proves zero annihilator on every second row at most \(s\), at every
\(j\ge k\).

Set \(j=a\) in (P.31). Then \(s=b-\binom a2+\binom k2\); the identity \[
(k-1)^2-\binom k2=\binom{k-1}{2}
\tag{P.34}
\] converts the exact base range \(s\ge(k-1)^2\) into the first
inequality in (P.29). This proves all the stated parameters.

For a concrete comparison within the proved families, at \(a=8\) the
base-five assertion begins at \(b=24\) and gives \(\lambda_2\le b-18\).
Formula (P.29) with \(k=6\) begins at \(b=38\) and gives
\(\lambda_2\le b-13\); with \(k=7\) it begins at \(b=43\) and gives
\(\lambda_2\le b-7\). Every assertion refers to the same specified map
\(h_{b,8}\). In particular \(h_{38,8}\) is surjective on every component
with \(\lambda_2\le25\).

\subsection{Check status}\label{check-status}

The proof of (P.11)--(P.20), the base-map identification, and all
endpoint inequalities were checked in this run. The parent agent will
record its independent reading separately. The original degree-five
computational theorem is used with the source status stated in §2; it is
not represented as a replayed calculation.

\section{Current map ranges and the exact remaining
transfer}\label{current-map-ranges-and-the-exact-remaining-transfer}

\subsection{M-001. Canonical maps, coefficient inequalities, and the
checked parameter
ranges}\label{m-001.-canonical-maps-coefficient-inequalities-and-the-checked-parameter-ranges}

Use the actual normalized grid morphism \[
h_{b,a}:\operatorname{Sym}^{b}(\operatorname{Sym}^{a}V)
 \longrightarrow\operatorname{Sym}^{a}(\operatorname{Sym}^{b}V),
\qquad
h_{b,a}(v_1^a\cdots v_b^a)=(v_1\cdots v_b)^a,
\quad b\ge a.
\] Its adjoint, for the invariant tensor inner products, is the
transposed-grid map \(h_{a,b}\). This follows by taking the adjoint of
the orthogonal row and column averaging projectors:
\((P_KP_H)^*=P_HP_K\). Thus surjectivity of the displayed map and
injectivity of its adjoint are equivalent by finite-dimensional linear
algebra. The dimensions of the source and target multiplicity spaces are
the two plethysm coefficients, while their common matrix rank is
G-002/G-011. These are the exact maps connecting the several assertions
used below.

The strongest uniform bound located and checked in this run is Suvrit
Sra, \emph{The q-Foulkes problem III: On Foulkes' conjecture with a
graded inner factor}, v1, 29 October 2026, Theorem 1.3 and its proof in
Section 2.2: the displayed map is surjective for \[
b\ge(a-1)^2.
\tag{M1.1}
\] The proof, including Lemma 2.1 and the common-polarization induction,
was read in full. Its statement is used as a cited result. The September
2026 bound \(b\ge a(a-1)\) in OpenAI's \emph{Quadratic stabilization for
Foulkes' conjecture}, Theorem 1, is contained in (M1.1). The endpoint
observation T-015 is retained as work history and makes no improvement
to (M1.1).

The canonical binary map is an isomorphism on every output with at most
two rows, for every \(a,b\), in all ambient dimensions: this is G-006,
carrying Landsberg's Theorem 7.21 through the natural Schur-functor
multiplicity spaces. The actual morphism, and its rank on every such
block, are specified there.

Cheung--Ikenmeyer--Mkrtchyan, \emph{Symmetrizing tableaux and the 5th
case of the Foulkes conjecture}, arXiv:1509.03944v1, Theorem 6(a),
proves injectivity of \(h_{5,6}\); Theorem 6(b) records the kernel at
\((5,5)\). Theorem 6(c) supplies the exact component \[
\lambda=(20,7,6,1,1,1)
\] with multiplicity one in each
\(\operatorname{Sym}^6\operatorname{Sym}^6 V\), on which \(h_{6,6}\) is
zero. This is a concrete test of the rank step: both multiplicities
equal one while the prescribed map has rank zero.

The coefficient inequality for outer degree six is supplied by OpenAI,
\emph{Foulkes' conjecture for sixth powers}, 25 September 2026, Theorem
1.1. Its exact assertion is existence of an equivariant injection for
all inner degrees at least six. The current run read its statement,
reduction, and computation descriptions; it did not execute its full
certificates and does not label the source formally verified. The
specified zero component at \((6,6)\) remains a property of the
prescribed canonical map. The two assertions are related by the common
multiplicity spaces and the explicit zero matrix on that component.

Evseev--Paget--Wildon, \emph{Character deflations and a generalization
of the Murnaghan--Nakayama rule}, Corollary 5.2, supplies the
coefficient inequality for \(a\le b\) and \(a+b\le19\). Sra III, Theorem
1.4, supplies the stated graded inequalities in depths at most six for
\(b\ge a\ge4\), and the further depth-seven, eight, and ten ranges in
that theorem. Those finite graded certificates were inspected as source
evidence and were not replayed in this run. No canonical-map rank is
inferred from these dimension inequalities.

The new exact rank certificates of this run are the prime-width and
product families R-008--R-009, together with G-006. For example, R-009
proves the complete determinant component of \(h_{24,6}\), and the
unique components of weights \((65,30^6)\) at \((a,b)=(7,35)\) and
\((87,60^8)\) at \((9,63)\). Their proofs exhibit source vectors and
their nonzero images. Full surjectivity at general \(b=(a-1)^2-1\)
remains an undecided question in this run.

\subsection{M-002. The orbit, theta, and purity transfer
gates}\label{m-002.-the-orbit-theta-and-purity-transfer-gates}

The compact dual pair and its vacuum twist in T-001 identify the
polynomial oscillator module with the Cauchy decomposition. Balanced
weight extraction gives \((S_\lambda\mathbb C^m)_{(n^m)}\), and the
specified symmetric-group Reynolds idempotent gives its invariant part.
The ordinary Gelfand--Tsetlin polytope counts the whole weight space.
Its transfer to plethysm is exactly the finite projector, evaluated by
G-009 and, explicitly for \(m=3\), G-012. T-005's determinant sign and
G-001's nonmonomial example identify failures of the particular
unprojected count and permutation-basis construction, with their exact
morphisms.

The attached Alpöge note is cited as instructed. Its steps 7 and 8, page
4, discuss interlacing period multiplicities and iterated theta lifts.
The cited declarations \texttt{\small ALPOGE\_\allowbreak{}Main.lean},
\texttt{\small ALPOGE\_\allowbreak{}Statements.lean}, and
\texttt{\small ALPOGE\_\allowbreak{}LEAN\_\allowbreak{}SCAN.tsv} are
absent from the supplied ZIP; no declaration in them was inspected in
this run. The transfer needed for the present question is an explicit
equivariant map from the groups and representations of those steps to
the balanced-weight representation together with its normalizer
idempotent, or to the two grid idempotents defining \(h_{b,a}\). The
four-page note supplies no such object or map. The exact compact theta
realization in T-001--T-002 and the invariant presentations in
G-009--G-011 establish the formulas directly without such an
identification. No claim about the cited RH and GRH theorems is needed
for these proofs.

For the suggested purity route, G-004 constructs the two exact Artin
projectors on \(\operatorname{Spec}\mathbb F_{q^2}\). Both images have
weight zero, while their invariant dimensions are one and zero. It
computes the missing datum explicitly: the projector rank. This
calculation leaves available geometric constructions that include the
projector, such as the finite quotient and coherent-module constructions
proved here.

For width seven, T-016--T-020 identifies the open orbit with stabilizer
\(G_2\times\mu_3\), and the extension map is the inclusion of polynomial
functions into their discriminant localization. The end-node sector has
the exact valuation and positive formula of T-018. The remaining
valuation filtration for other highest weights is an explicit smaller
orbit calculation. G-009--G-010 already give terminating positive counts
for those coefficients; the unresolved question there is a small direct
orbit formula replacing the general elimination construction.

The requested Masullo locator is an input gate: no Masullo work,
statement, or identifier occurs in the supplied file inventory. A
theorem or a map from such a work is not introduced by inference.

\subsection{Source URLs}\label{source-urls}

\begin{itemize}
\tightlist
\item
  Sra III:
  \url{https://suvrit.github.io/papers/qfoulkes3/qfoulkes3-v1.pdf};
  supplement:
  \url{https://suvrit.github.io/papers/qfoulkes3/qfoulkes3-v1-supplement.pdf}.
\item
  Cheung--Ikenmeyer--Mkrtchyan: \url{https://arxiv.org/abs/1509.03944}.
\item
  Landsberg's supplied survey, Theorem 7.21:
  \url{https://arxiv.org/abs/1305.7387}.
\item
  Alpöge: \url{https://alpo.ge/remark.pdf}, attached four-page file,
  steps 7--8 on page 4.
\end{itemize}

This register is a current-source audit dated 21 November 2026,
Europe/Berlin. Its labels identify the audit and explicit morphisms; the
attributed literature theorems remain cited results.

\part{Checked archive findings and integral identities}

\section{Archive audit for
GCT\_PLETH}\label{archive-audit-for-gct_pleth}

Run:
\texttt{\small /\allowbreak{}root/\allowbreak{}archive\_\allowbreak{}audit},
21 November 2026 (Europe/Berlin). Model: GPT-6 Astra Pro; effort setting
is not exposed in the agent context.

\subsection{Evidence and scope}\label{evidence-and-scope}

All thirteen supplied \texttt{\small ARCHIVE\_\allowbreak{}*.tex} files
were read in full. They are transcriptions, not images. The
uncertainties, deletions, and corrections encoded in them remain
evidence about the transcription. The public-source follow-up in
\texttt{\small PRIMARY\_\allowbreak{}SCAN\_\allowbreak{}FOLLOWUP.md} now
records direct inspection of all thirteen original Montpellier target
leaves, from cotes2,4,12,16,43,79,89,119,133, together with their
defining neighbors where needed.
\texttt{\small SA17\_\allowbreak{}STUDY\_\allowbreak{}cote133\_\allowbreak{}lambda.md}
was read in full. Its statements about scans absent from the supplied
transcription set are prior-run reports; the lambda/finite-set
pages133/004 and133/019--029 have now also been inspected directly, with
current-run proofs distinguished in
\texttt{\small STUDY\_\allowbreak{}ADDITIONS.md}. Its claimed
equivariant check on scan133/019 is refuted below and checked against
the actual manuscript. The additional single-map-power summand assertion
is likewise present on original133/025 and is refuted in A-013.

The symbol \texttt{\small WG} on 133/006 is explicitly the group
completion of the orthogonal-sum monoid of isometry classes. Write it
\texttt{\small GW\textasciicircum{}⊕\allowbreak{}} here. Write
\texttt{\small GW} for the quotient that also imposes the metabolic
relation
\texttt{\small [\allowbreak{}M]\allowbreak{}=\allowbreak{}[\allowbreak{}H(\allowbreak{}L)\allowbreak{}]\allowbreak{}}
for a Lagrangian \texttt{\small L}. There is a concrete surjective ring
map

\[
q:GW^\oplus\longrightarrow GW,
\quad [M,b]\longmapsto[M,b],
\]

with kernel generated by the differences
\texttt{\small [\allowbreak{}M]\allowbreak{}-\allowbreak{}[\allowbreak{}H(\allowbreak{}L)\allowbreak{}]\allowbreak{}}.
The corresponding forgetful maps land respectively in the split Green
ring \texttt{\small K\_\allowbreak{}0\textasciicircum{}⊕\allowbreak{}}
and the exact Grothendieck ring \texttt{\small K\_\allowbreak{}0}.
Exterior powers commute with these forgetful maps. The map
\texttt{\small q} and its nonzero kernel are essential in A-003; the
general special-λ theorem for \texttt{\small GW} cannot be pulled back
through this noninjective map merely by citing it.

\subsection{A-001. Cote 133, scan 019: symmetric powers, subsets, and
the supplied study's
counterexample}\label{a-001.-cote-133-scan-019-symmetric-powers-subsets-and-the-supplied-studys-counterexample}

\textbf{Statements on the page.} The page displays
\texttt{\small σ\_\allowbreak{}t(\allowbreak{}E)\allowbreak{}λ\_\allowbreak{}\{−t\}(\allowbreak{}E)\allowbreak{}=\allowbreak{}1},
\texttt{\small λ\_\allowbreak{}t(\allowbreak{}L)\allowbreak{}=\allowbreak{}1+\allowbreak{}tL},
\texttt{\small σ\_\allowbreak{}t(\allowbreak{}L)\allowbreak{}=\allowbreak{}(\allowbreak{}1−tL)\allowbreak{}\textasciicircum{}\{-\allowbreak{}1\}},
\texttt{\small σ²(\allowbreak{}E)\allowbreak{}=\allowbreak{}E²−λ²(\allowbreak{}E)\allowbreak{}},
\texttt{\small σ²(\allowbreak{}E)\allowbreak{}=\allowbreak{}λ²(\allowbreak{}E)\allowbreak{}+\allowbreak{}E},
and
\texttt{\small 2λ²(\allowbreak{}E)\allowbreak{}+\allowbreak{}E=\allowbreak{}E²}.
The last numerical check is
\texttt{\small 2 binom(\allowbreak{}d,\allowbreak{}2)\allowbreak{}+\allowbreak{}d=\allowbreak{}d²}.
It also asks for
\texttt{\small λ\textasciicircum{}i(\allowbreak{}xy)\allowbreak{}}.

For a locally free sheaf, the first four formulas are the standard
exterior/symmetric-power identities in exact
\texttt{\small K\_\allowbreak{}0}. The short exact sequence

\[
0\longrightarrow\Lambda^2E\xrightarrow{u\wedge v\mapsto u\otimes v-v\otimes u}
E\otimes E\longrightarrow\operatorname{Sym}^2E\longrightarrow0
\]

is locally split as a sequence of underlying modules and proves the
degree-two identity. Exterior-power operations on this exact
\texttt{\small K\_\allowbreak{}0} satisfy the universal special-λ
identities; see SGA 6, V, Lemma 2.2.1 and VI, Theorem 3.3. The line
identities follow directly from the one-dimensional tensor powers. These
are known identities, not novelty claims.

\textbf{Refutation in the finite-set reading used by the supplied
study.} Let \texttt{\small G=\allowbreak{}C\_\allowbreak{}2} and let
\texttt{\small E=\allowbreak{}\{0,\allowbreak{}1\}} with the nonidentity
element interchanging the points. In the Burnside ring
\texttt{\small B(\allowbreak{}G)\allowbreak{}}, write \texttt{\small 1}
for the singleton. The unordered two-element subsets form
\texttt{\small 1}. The two-element multisets form
\texttt{\small E+\allowbreak{}1}: two doubled points are interchanged
and the mixed pair is fixed. Ordered pairs form two free orbits, so
\texttt{\small E²=\allowbreak{}2E}. Consequently

\[
\lambda^2_{\rm subset}(E)=1,\qquad
\sigma^2_{\rm multiset}(E)=E+1,\qquad E^2=2E.
\]

The boxed equality would assert
\texttt{\small 2+\allowbreak{}E=\allowbreak{}2E}. Apply the fixed-point
ring homomorphism

\[
\phi_G:B(G)\longrightarrow\mathbb Z,
\qquad [X]\longmapsto|X^G|.
\]

Its two sides are \texttt{\small 2} and \texttt{\small 0}. This is a
counterexample, including after group completion. The numerical
cardinality map \texttt{\small φ\_\allowbreak{}1} sends both sides to
\texttt{\small 4}; that explains exactly what the last numerical line
verifies.

\textbf{Exact dictionary.} Linearization
\texttt{\small \textbackslash{}ell:\allowbreak{}B(\allowbreak{}G)\allowbreak{}→\allowbreak{}R\_\allowbreak{}\textbackslash{}mathbb Q(\allowbreak{}G)\allowbreak{}}
sends \texttt{\small E} to \texttt{\small 1+\allowbreak{}ε}, with
\texttt{\small ε} the sign character. It commutes with symmetric powers
because the monomial basis of
\texttt{\small Sym\textasciicircum{}n\textbackslash{}mathbb Q[\allowbreak{}E]\allowbreak{}}
is the multiset basis. It does not carry the unsigned-subset operation
to exterior powers:
\texttt{\small \textbackslash{}ell(\allowbreak{}\textbackslash{}lambda\textasciicircum{}2\_\allowbreak{}\{subset\}E)\allowbreak{}=\allowbreak{}1},
whereas
\texttt{\small \textbackslash{}Lambda\textasciicircum{}2\textbackslash{}ell(\allowbreak{}E)\allowbreak{}=\allowbreak{}ε}.
The Burnside pre-λ operation determined by
\texttt{\small λ\_\allowbreak{}t=\allowbreak{}σ\_\allowbreak{}\{−t\}\textasciicircum{}\{−1\}}
has

\[
\lambda^2(E)=E^2-\sigma^2(E)=E-1,
\]

whose linearization is \texttt{\small ε}. Thus the discrepancy is
exhibited by actual maps and a nonzero element, rather than by a claimed
absence of mathematical relations.

\textbf{Plethysm consequence.} The page's valid
\texttt{\small K\_\allowbreak{}0} identities supply λ-polynomial
arithmetic. The unsigned finite-subset identity cannot be used to derive
equivariant multiplicities: the \texttt{\small C\_\allowbreak{}2}
counterexample already changes an irreducible constituent. The source
\texttt{\small SA17}, section ``Scan 19: λ- and σ-operations'',
incorrectly records the simultaneous equivariant equalities as checked.

\subsection{A-002. Cote 133, scan 006: the rank obstruction in
characteristic
two}\label{a-002.-cote-133-scan-006-the-rank-obstruction-in-characteristic-two}

\textbf{Statements on the page.} The group completion of strictly
nondegenerate quadratic modules is asserted to have a pre-λ structure
augmented by rank; on a connected base where \texttt{\small 2} is not
invertible the ranks are even; with \texttt{\small 2} invertible there
is a unit. It asks whether the universal product and iterate formulas
hold, and reduces their study to classifying topoi of orthogonal groups.

\textbf{Proof of the parity statement.} Over a characteristic-two
residue field, the polar bilinear form of a quadratic form is
alternating. A nondegenerate alternating matrix has even size
(successive hyperbolic-pair elimination proves this). On a connected
ringed space, the rank of a locally free module of finite constant rank
is constant. A point with characteristic-two residue field therefore
forces its rank to be even. The assertion needs such a point; saying
merely ``2 not invertible'' for an arbitrary locally ringed topos still
requires the usual stalk interpretation.

\textbf{Counterexample to the stated pre-λ augmentation.} Already over
\texttt{\small Spec\textbackslash{}mathbb F\_\allowbreak{}2}, every
object in the specified monoid has even rank, hence every element of its
group completion has even rank. The hyperbolic plane has rank two. A
pre-λ augmentation to the binomial λ-ring would give

\[
\operatorname{rk}(\lambda^2H)=\binom22=1,
\]

which is impossible in that group. This refutes the pre-λ assertion with
precisely the exterior/rank compatibility written on the page. Also the
standard unital definition has
\texttt{\small λ\textasciicircum{}0=\allowbreak{}1}, whose rank is one.
The nonunital parenthesis does not repair the displayed rank
compatibility.

With \texttt{\small 2} invertible, use the explicit inverse maps
\texttt{\small q↦b\_\allowbreak{}q/\allowbreak{}2} and
\texttt{\small b↦(\allowbreak{}v↦b(\allowbreak{}v,\allowbreak{}v)\allowbreak{})\allowbreak{}}
between quadratic forms and symmetric bilinear forms. Tensor products,
the form on exterior powers, and the rank-one unit are defined through
this equivalence. The quadratic hyperbolic form
\texttt{\small q\_\allowbreak{}H(\allowbreak{}v,\allowbreak{}φ)\allowbreak{}=\allowbreak{}φ(\allowbreak{}v)\allowbreak{}}
corresponds to the bilinear matrix with off-diagonal entries
\texttt{\small 1/\allowbreak{}2}; scaling one of the two summands by
\texttt{\small 2} identifies it with the customary bilinear hyperbolic
matrix with off-diagonal entries \texttt{\small 1}. This fixes the
convention used below. The remaining special-λ question for the literal
group completion is answered by A-003.

\subsection{A-003. Cote 133, scan 006: the literal split ring fails the
special-λ identities even with 2
invertible}\label{a-003.-cote-133-scan-006-the-literal-split-ring-fails-the-special-ux3bb-identities-even-with-2-invertible}

This is a complete counterexample to the page's general special-λ
question for \texttt{\small GW\textasciicircum{}⊕\allowbreak{}}.

Let
\texttt{\small k=\allowbreak{}\textbackslash{}overline\{\textbackslash{}mathbb F\}\_\allowbreak{}3},
\texttt{\small G=\allowbreak{}C\_\allowbreak{}3}, and
\texttt{\small S=\allowbreak{}BG} with structure sheaf
\texttt{\small k}. The associated vector-bundle category is
finite-dimensional \texttt{\small kG}-modules. Let
\texttt{\small J\_\allowbreak{}i},
\texttt{\small 1≤\allowbreak{}i≤\allowbreak{}3}, denote the
indecomposable module whose generator acts as a single unipotent Jordan
block of size \texttt{\small i}. Every module is a direct sum of these
three, uniquely, because
\texttt{\small (\allowbreak{}g−1)\allowbreak{}\textasciicircum{}3=\allowbreak{}0}.
Thus the additive group of the split Green ring is free on
\texttt{\small 1=\allowbreak{}[\allowbreak{}J\_\allowbreak{}1]\allowbreak{}},
\texttt{\small y=\allowbreak{}[\allowbreak{}J\_\allowbreak{}2]\allowbreak{}},
and
\texttt{\small x=\allowbreak{}[\allowbreak{}J\_\allowbreak{}3]\allowbreak{}}.

Choose
\texttt{\small J\_\allowbreak{}3=\allowbreak{}k[\allowbreak{}G]\allowbreak{}}
with its permutation basis and nondegenerate invariant dot-product
bilinear form; denote the symmetric form class by \texttt{\small X}. Let
\texttt{\small Z=\allowbreak{}H(\allowbreak{}J\_\allowbreak{}2)\allowbreak{}}.
Under forgetting forms,
\texttt{\small F(\allowbreak{}X)\allowbreak{}=\allowbreak{}x} and
\texttt{\small F(\allowbreak{}Z)\allowbreak{}=\allowbreak{}2y}.

The split Green-ring identities are

\[
x^2=3x,\qquad xy=2x,\qquad y^2=x+1,
\quad \lambda^2x=x,\quad\lambda^3x=1,
\quad\lambda^2y=1,\quad\lambda^3y=0.
\tag{A3.1}
\]

The first two follow from the equivariant isomorphism
\texttt{\small kG⊗\allowbreak{}M≅kG\textasciicircum{}\{dim M\}}: the map
on basis vectors is
\texttt{\small g⊗\allowbreak{}m↦g⊗\allowbreak{}g\textasciicircum{}\{-\allowbreak{}1\}m},
with the group acting on the first factor in the codomain. For the
third, take
\texttt{\small g=\allowbreak{}[\allowbreak{}[\allowbreak{}1,\allowbreak{}1]\allowbreak{},\allowbreak{}[\allowbreak{}0,\allowbreak{}1]\allowbreak{}]\allowbreak{}}.
On \texttt{\small J\_\allowbreak{}2⊗\allowbreak{}J\_\allowbreak{}2},
\texttt{\small N=\allowbreak{}g⊗\allowbreak{}g−1} has rank two,
\texttt{\small N²} has rank one, and \texttt{\small N³=\allowbreak{}0};
hence the Jordan type is
\texttt{\small (\allowbreak{}3,\allowbreak{}1)\allowbreak{}}. On
\texttt{\small \textbackslash{}Lambda\textasciicircum{}2J\_\allowbreak{}3},
the action is again a 3-cycle on the basis
\texttt{\small e\_\allowbreak{}1∧e\_\allowbreak{}2,\allowbreak{}e\_\allowbreak{}2∧e\_\allowbreak{}3,\allowbreak{}e\_\allowbreak{}3∧e\_\allowbreak{}1},
so it is \texttt{\small J\_\allowbreak{}3}. Both determinants are one.
These computations prove (A3.1);
\texttt{\small check\_\allowbreak{}archive.py} additionally verifies the
finite-field matrix ranks and exterior matrices exactly.

In every pre-λ ring, Newton's polynomial gives

\[
\psi^3(a)=a^3-3a\lambda^2(a)+3\lambda^3(a),
\]

and \texttt{\small ψ³} is additive; this follows by extracting the third
coefficient of
\texttt{\small t(\allowbreak{}d/\allowbreak{}dt)\allowbreak{}logλ\_\allowbreak{}t}.
The special-λ product polynomials force
\texttt{\small ψ³(\allowbreak{}ab)\allowbreak{}=\allowbreak{}ψ³(\allowbreak{}a)\allowbreak{}ψ³(\allowbreak{}b)\allowbreak{}}:
over universal line elements both sides are
\texttt{\small \textbackslash{}sum\_\allowbreak{}\{i,\allowbreak{}j\}u\_\allowbreak{}i\textasciicircum{}3v\_\allowbreak{}j\textasciicircum{}3},
an integral symmetric-polynomial identity.

Equation (A3.1) gives

\[
\psi^3(x)=3,\qquad \psi^3(y)=2x-2y.
\]

The forgetful map commutes with every exterior power and hence with
\texttt{\small ψ³}. Since
\texttt{\small F(\allowbreak{}XZ)\allowbreak{}=\allowbreak{}x(\allowbreak{}2y)\allowbreak{}=\allowbreak{}4x},
the two required sides after forgetting are

\[
F(\psi^3(XZ))=12,
\qquad F(\psi^3(X)\psi^3(Z))=12x-12y.
\]

Their difference is \texttt{\small 12(\allowbreak{}x−y−1)\allowbreak{}},
a nonzero element of the free abelian group on
\texttt{\small J\_\allowbreak{}1,\allowbreak{}J\_\allowbreak{}2,\allowbreak{}J\_\allowbreak{}3}.
The coefficients here are integers in a Grothendieck group, not scalars
in \texttt{\small k}. Consequently
\texttt{\small GW\textasciicircum{}⊕\allowbreak{}(\allowbreak{}BC\_\allowbreak{}3)\allowbreak{}}
with its exterior powers is not a special λ-ring.

The separate iterate identity is now refuted in \textbf{A-015},
\texttt{\small ITERATE\_\allowbreak{}COUNTEREXAMPLE.md}. For the actual
object \texttt{\small A=\allowbreak{}X⊥H(\allowbreak{}J₂)\allowbreak{}},
the underlying class is
\texttt{\small a=\allowbreak{}x+\allowbreak{}2y}. Its second exterior
power is \texttt{\small 6x+\allowbreak{}3}, and its iterated third
exterior power is \texttt{\small 441x+\allowbreak{}7}. The universal
polynomial
\texttt{\small Q\_\allowbreak{}(\allowbreak{}3,\allowbreak{}2)\allowbreak{}=\allowbreak{}e₁²e₄−2e₂e₄−e₁e₅+\allowbreak{}e₆+\allowbreak{}e₃²}
instead evaluates to \texttt{\small 445x−4y+\allowbreak{}3}, giving the
nonzero defect \texttt{\small −4(\allowbreak{}x−y−1)\allowbreak{}}. The
universal polynomial was verified coefficient by coefficient in six
variables; this is injective in symmetric-function degree six. Both the
product and iterate questions therefore have separate explicit
counterexamples for the literal split ring.

\textbf{Exact geometric transfer.} The representations with forms define
homomorphisms of group schemes
\texttt{\small C\_\allowbreak{}3→\allowbreak{}O\_\allowbreak{}3} and
\texttt{\small C\_\allowbreak{}3→\allowbreak{}O\_\allowbreak{}4}, and
hence a classifying-topos map
\texttt{\small BC\_\allowbreak{}3→\allowbreak{}B(\allowbreak{}O\_\allowbreak{}3×O\_\allowbreak{}4)\allowbreak{}}.
Restriction pulls back the two universal quadratic modules to
\texttt{\small X,\allowbreak{}Z}. This is the concrete instance of the
universal orthogonal-group reduction proposed in the margin.

\textbf{Relation with the modern theorem.} Zibrowius, \emph{Symmetric
representation rings are λ-rings}, NYJM 21 (2015), Definition 1.2,
p.1060, defines \texttt{\small GW} by orthogonal-sum and metabolic
relations; its main theorem, p.1056, establishes the special-λ structure
on that quotient. For the present \texttt{\small C\_\allowbreak{}3},
exact \texttt{\small K\_\allowbreak{}0} sends
\texttt{\small J\_\allowbreak{}i} to
\texttt{\small i[\allowbreak{}J\_\allowbreak{}1]\allowbreak{}}, and
therefore sends \texttt{\small x−y−1} to zero. The counterexample
establishes precisely why the split statement cannot be substituted for
that cited theorem. Zibrowius, \emph{The γ-filtration on the Witt ring
of a scheme}, arXiv:1411.3198v2, Theorem 2, extends the theorem to
schemes over a field of characteristic different from two; Proposition
2.15 supplies the \texttt{\small O\_\allowbreak{}n×O\_\allowbreak{}m}
torsor construction. These known theorems are cited, not reproved.

\textbf{Plethysm consequence.} Passing to exact
\texttt{\small K\_\allowbreak{}0} retains the usual Schur operations by
SGA 6 VI 3.3. Working in the uncorrected split quadratic ring would
invalidate even the degree-three universal product relation. The
correcting quotient is explicit.

\subsection{A-004. Cote 133, scan 008: hyperbolic sums and the proposed
equivalence
relation}\label{a-004.-cote-133-scan-008-hyperbolic-sums-and-the-proposed-equivalence-relation}

\subsubsection{A-004a: explicit isometry with 2
invertible}\label{a-004a-explicit-isometry-with-2-invertible}

For a strictly nondegenerate quadratic form \texttt{\small q} with polar
form \texttt{\small b} over a ringed topos in which \texttt{\small 2} is
invertible, define

\[
\Phi:(M,q)\perp(M,-q)\longrightarrow H(M),
\quad (x,y)\longmapsto(x+y,\tfrac12 b(x-y,-)).
\]

The hyperbolic quadratic form evaluates this image to
\texttt{\small b(\allowbreak{}x−y,\allowbreak{}x+\allowbreak{}y)\allowbreak{}/\allowbreak{}2=\allowbreak{}q(\allowbreak{}x)\allowbreak{}−q(\allowbreak{}y)\allowbreak{}}.
Nondegeneracy of \texttt{\small b} and invertibility of
\texttt{\small 2} give an inverse by solving for
\texttt{\small x+\allowbreak{}y} and \texttt{\small x−y}. This formula
is natural in isometries, hence works on classifying topoi as well. The
page's classifying-topos caveat therefore has no example in this range.
The symmetric-bilinear version appears as Zibrowius 2015, Lemma 1.1,
p.1059; the displayed formula here translates it into the page's
quadratic convention.

\subsubsection{A-004b: a characteristic-two classifying-topos
counterexample}\label{a-004b-a-characteristic-two-classifying-topos-counterexample}

Let
\texttt{\small k=\allowbreak{}\textbackslash{}overline\{\textbackslash{}mathbb F\}\_\allowbreak{}2},
\texttt{\small M=\allowbreak{}k²}, and
\texttt{\small q(\allowbreak{}x,\allowbreak{}y)\allowbreak{}=\allowbreak{}xy}.
Its polar form
\texttt{\small b(\allowbreak{}(\allowbreak{}x,\allowbreak{}y)\allowbreak{},\allowbreak{}(\allowbreak{}x',\allowbreak{}y')\allowbreak{})\allowbreak{}=\allowbreak{}xy'+\allowbreak{}x'y}
is nondegenerate. Let
\texttt{\small G=\allowbreak{}O(\allowbreak{}q)\allowbreak{}=\allowbreak{}\textbackslash{}mathbb G\_\allowbreak{}m\textbackslash{}rtimes C\_\allowbreak{}2},
with the torus acting by
\texttt{\small diag(\allowbreak{}t,\allowbreak{}t\textasciicircum{}\{-\allowbreak{}1\})\allowbreak{}}
and the involution interchanging coordinates. The module
\texttt{\small M} is simple: a torus-stable submodule is a sum of its
two distinct weight lines, and the involution exchanges those lines. Its
equivariant endomorphism ring is \texttt{\small k}. Its invariant
bilinear forms are exactly the multiples of \texttt{\small b}, hence
alternating.

For every \texttt{\small G}-module \texttt{\small E} and every
\texttt{\small G}-map
\texttt{\small f:\allowbreak{}M→\allowbreak{}H(\allowbreak{}E)\allowbreak{}=\allowbreak{}E⊕\allowbreak{}E\textasciicircum{}∨},
write
\texttt{\small f=\allowbreak{}(\allowbreak{}i,\allowbreak{}j)\allowbreak{}}.
The bilinear form
\texttt{\small B(\allowbreak{}u,\allowbreak{}v)\allowbreak{}=\allowbreak{}j(\allowbreak{}v)\allowbreak{}(\allowbreak{}i(\allowbreak{}u)\allowbreak{})\allowbreak{}}
on \texttt{\small M} is \texttt{\small G}-invariant, so
\texttt{\small B(\allowbreak{}v,\allowbreak{}v)\allowbreak{}=\allowbreak{}0}.
Therefore
\texttt{\small q\_\allowbreak{}H(\allowbreak{}f(\allowbreak{}v)\allowbreak{})\allowbreak{}=\allowbreak{}0}
for every such map.

The first-summand map
\texttt{\small M→\allowbreak{}(\allowbreak{}M,\allowbreak{}q)\allowbreak{}⊥(\allowbreak{}M,\allowbreak{}q)\allowbreak{}}
has pullback \texttt{\small q≠0}. Thus
\texttt{\small (\allowbreak{}M,\allowbreak{}q)\allowbreak{}⊥(\allowbreak{}M,\allowbreak{}q)\allowbreak{}}
cannot be isometric to \texttt{\small H(\allowbreak{}E)\allowbreak{}}
for any \texttt{\small E}. The same argument, using the first summand,
proves

\[
(M,q)\perp(M,q)\perp H(F)\not\simeq H(E)
\]

for all \texttt{\small E,\allowbreak{}F}. This proves the page's
suggested failure on \texttt{\small BO(\allowbreak{}2)\allowbreak{}}
with full characteristic and form data.

The proposed relation on scan 008 requires
\texttt{\small q⊥(\allowbreak{}−q')\allowbreak{}⊥H(\allowbreak{}F)\allowbreak{}≅H(\allowbreak{}E)\allowbreak{}}.
In this example \texttt{\small −q=\allowbreak{}q}. Taking
\texttt{\small q'=\allowbreak{}q}, the displayed impossibility says that
the proposed relation is not even reflexive. Consequently it cannot
describe a quotient in the stated generality. It does describe the usual
stable-hyperbolic relation in the range of A-004a: adding negative forms
removes the extra group-completion stabilizer because every
\texttt{\small P⊥(\allowbreak{}−P)\allowbreak{}} is hyperbolic.

\subsection{A-005. Cote 133, scan 008: the hyperbolic ideal and failure
of λ
descent}\label{a-005.-cote-133-scan-008-the-hyperbolic-ideal-and-failure-of-ux3bb-descent}

Over a base with \texttt{\small 2} invertible, let \texttt{\small F}
forget the form and \texttt{\small H} be the hyperbolic functor. The
exact identity

\[
(M,b)\otimes H(E)\simeq H(M\otimes E)
\tag{A5.1}
\]

uses \texttt{\small b:\allowbreak{}M≅M\textasciicircum{}∨} to identify
the second summand \texttt{\small M⊗\allowbreak{}E\textasciicircum{}∨}
with
\texttt{\small (\allowbreak{}M⊗\allowbreak{}E)\allowbreak{}\textasciicircum{}∨}.
The hyperbolic pairing is carried to the tensor-product pairing. Thus
the image of \texttt{\small H} is an ideal in the corresponding form
ring. The two available additive maps are
\texttt{\small H:\allowbreak{}K\_\allowbreak{}0\textasciicircum{}⊕\allowbreak{}→\allowbreak{}GW\textasciicircum{}⊕\allowbreak{}}
and \texttt{\small H:\allowbreak{}K\_\allowbreak{}0→\allowbreak{}GW}.
The map from exact \texttt{\small K\_\allowbreak{}0} to the split form
ring need not be well-defined: in characteristic three,
\texttt{\small J\_\allowbreak{}2=\allowbreak{}2J\_\allowbreak{}1} in
exact \texttt{\small K\_\allowbreak{}0}, while the split forgetful
images of \texttt{\small H(\allowbreak{}J\_\allowbreak{}2)\allowbreak{}}
and \texttt{\small H(\allowbreak{}2J\_\allowbreak{}1)\allowbreak{}} are
\texttt{\small 2J\_\allowbreak{}2} and
\texttt{\small 4J\_\allowbreak{}1}.

The rank of \texttt{\small H(\allowbreak{}E)\allowbreak{}} is even, so
rank modulo two descends to the Witt quotient. Exterior powers do not
descend, already over any field of characteristic different from two:
the hyperbolic plane is zero in the Witt ring, but its determinant form
is

\[
\lambda^2 H(1)=\langle-1\rangle.
\]

In the quadratic convention fixed in A-002, the determinant coefficient
is initially \texttt{\small −1/\allowbreak{}4}, and the change of basis
by \texttt{\small 2} gives the displayed isometry to
\texttt{\small \textbackslash{}langle-\allowbreak{}1\textbackslash{}rangle}.
The latter has odd rank, hence nonzero rank modulo two in the quotient.
Since \texttt{\small λ²(\allowbreak{}0)\allowbreak{}=\allowbreak{}0}, no
operation on the quotient can commute with both the quotient map and
this exterior-power operation. This completely answers the page's
descent question for the ordinary Witt quotient.

\subsection{A-006. Cote 133, scan 009: characteristic classes and the
struck-out
claim}\label{a-006.-cote-133-scan-009-characteristic-classes-and-the-struck-out-claim}

The page specifies naturality, normalization on symmetric lines, and
\texttt{\small w(\allowbreak{}x+\allowbreak{}y)\allowbreak{}=\allowbreak{}w(\allowbreak{}x)\allowbreak{}w(\allowbreak{}y)\allowbreak{}},
then states universal product and exterior-power formulas. For schemes
over
\texttt{\small \textbackslash{}mathbb Z[\allowbreak{}1/\allowbreak{}2]\allowbreak{}},
the exact modern input is the étale-cohomology splitting morphism of
Zibrowius 2016, Proposition 3.3: it splits a symmetric bundle into
symmetric lines and is injective on
\texttt{\small H\textasciicircum{}*\_\allowbreak{}\{et\}(\allowbreak{}-\allowbreak{},\allowbreak{}\textbackslash{}mathbb Z/\allowbreak{}2)\allowbreak{}}.
This is sufficient for the following precise restatement of the page.

Write the degree-one classes of the split lines of
\texttt{\small E,\allowbreak{}F} as
\texttt{\small u\_\allowbreak{}1,\allowbreak{}…,\allowbreak{}u\_\allowbreak{}r}
and
\texttt{\small v\_\allowbreak{}1,\allowbreak{}…,\allowbreak{}v\_\allowbreak{}s}.
Naturality and injectivity carry back the formulas

\[
w(E)=\prod_i(1+u_i),\qquad
w(E\otimes F)=\prod_{i,j}(1+u_i+v_j),
\] \[
w(\Lambda^a E)=\prod_{|I|=a}\left(1+\sum_{i\in I}u_i\right).
\tag{A6.1}
\]

The products are symmetric in each root family, so each coefficient is a
universal polynomial in the elementary symmetric functions and the
ranks. Tensoring symmetric lines adds their degree-one Kummer classes,
which proves the root substitution in (A6.1); exterior powers select
subsets. The injective pullback proves the formulas before splitting.
For virtual classes total \texttt{\small w} takes values in the
completed group
\texttt{\small 1+\allowbreak{}\textbackslash{}prod\_\allowbreak{}\{i>0\}H\textasciicircum{}i};
the completion is needed to invert total classes without a
cohomological-dimension bound. SGA 6, V §6 describes the associated
Chern-ring universal construction. The page does not specify such a ring
structure on its displayed target, and its assertion of a λ-homomorphism
is read through this universal construction.

The struck-out claim \texttt{\small w∘H=\allowbreak{}1} is false. Over
\texttt{\small Spec\textbackslash{}mathbb R},
\texttt{\small H(\allowbreak{}1)\allowbreak{}≅\textbackslash{}langle1,\allowbreak{}-\allowbreak{}1\textbackslash{}rangle},
and normalization gives
\texttt{\small w\_\allowbreak{}1(\allowbreak{}H(\allowbreak{}1)\allowbreak{})\allowbreak{}=\allowbreak{}(\allowbreak{}-\allowbreak{}1)\allowbreak{}≠0}
in
\texttt{\small H\textasciicircum{}1\_\allowbreak{}\{et\}(\allowbreak{}\textbackslash{}mathbb R,\allowbreak{}\textbackslash{}mathbb Z/\allowbreak{}2)\allowbreak{}}.
Thus total \texttt{\small w} does not factor through the Witt quotient.
This counterexample suffices without relying on the more elaborate
formula in the supplied SA17 study. The study's formula for hyperbolic
bundles is not used in this run.

\textbf{Plethysm consequence.} Applying (A6.1) to iterated Schur
operations computes characteristic classes of those bundles. The exact
morphism
\texttt{\small GW→\allowbreak{}H\textasciicircum{}0(\allowbreak{}\textbackslash{}mathbb Z)\allowbreak{}×(\allowbreak{}1+\allowbreak{}\textbackslash{}widehat H\textasciicircum{}+\allowbreak{})\allowbreak{}}
does not provide the full irreducible-character expansion; e.g.~the
characteristic-three defect in A-003 already maps to zero after exact
passage. No assertion about all possible geometric ways to recover that
expansion is made.

\subsection{A-007. Cote 133, scan 028: tensor-power identity and its
positivity
test}\label{a-007.-cote-133-scan-028-tensor-power-identity-and-its-positivity-test}

Let \texttt{\small X} be a commutatively ringed topos over
\texttt{\small \textbackslash{}mathbb Q}, \texttt{\small E} a finite
locally free object, and
\texttt{\small ν=\allowbreak{}(\allowbreak{}1\textasciicircum{}\{α\_\allowbreak{}1\}2\textasciicircum{}\{α\_\allowbreak{}2\}…)\allowbreak{}⊢n}.
Put

\[
\rho_i=[\mathbf1_{S_i}],\qquad
\rho_\nu=\prod_i\rho_i^{\alpha_i}
=[\operatorname{Ind}_{\prod_jS_{\nu_j}}^{S_n}\mathbf1].
\]

The product here is induction. With \texttt{\small m\_\allowbreak{}ν}
expressed as an integral polynomial in
\texttt{\small e\_\allowbreak{}1,\allowbreak{}…,\allowbreak{}e\_\allowbreak{}n},
the page's formula is

\[
[E^{\otimes n}]=\sum_{\nu\vdash n}m_\nu(E)\rho_\nu.
\tag{A7.1}
\]

\textbf{Proof in this fully specified setting.} The split semisimple
algebra
\texttt{\small \textbackslash{}mathbb Q[\allowbreak{}S\_\allowbreak{}n]\allowbreak{}}
decomposes into matrix algebras indexed by partitions. Its rational
idempotents act on
\texttt{\small E\textasciicircum{}\{⊗\allowbreak{}n\}} and give the
functorial decomposition

\[
E^{\otimes n}\simeq\bigoplus_{\lambda\vdash n}S_\lambda(E)\otimes[\lambda].
\tag{A7.2}
\]

Thus
\texttt{\small K\_\allowbreak{}0(\allowbreak{}BS\_\allowbreak{}\{n,\allowbreak{}X\})\allowbreak{}≅\textbackslash{}bigoplus\_\allowbreak{}\{λ⊢n\}K\_\allowbreak{}0(\allowbreak{}X)\allowbreak{}}.
The Frobenius-characteristic map sends
\texttt{\small [\allowbreak{}λ]\allowbreak{}} to
\texttt{\small s\_\allowbreak{}λ} and \texttt{\small ρ\_\allowbreak{}ν}
to \texttt{\small h\_\allowbreak{}ν}. Apply evaluation at
\texttt{\small E} to the first factor of the Cauchy identity

\[
\sum_\lambda s_\lambda(x)s_\lambda(y)
=\prod_{i,j}(1-x_i y_j)^{-1}
=\sum_\nu m_\nu(x)h_\nu(y).
\]

The equality of coefficients in this direct sum is precisely (A7.1).
This is a proof of the page's formulation from the known Cauchy and
Schur--Weyl theorems, not a new proof of those theorems. The exact
primary locators in Macdonald, \emph{Symmetric Functions and Hall
Polynomials}, 2nd ed., are I (4.2), p.62, and (4.3), p.63, for the two
Cauchy expansions; I (7.3)--(7.6), pp.113--114, for the characteristic
map and irreducible characters; Appendix A (5.6), p.156, for the
tensor-power decomposition. These equations were inspected in the
primary PDF.

The coefficients singled out in the supplied study are
\texttt{\small m\_\allowbreak{}(\allowbreak{}1\textasciicircum{}n)\allowbreak{}=\allowbreak{}e\_\allowbreak{}n}
and
\texttt{\small m\_\allowbreak{}(\allowbreak{}n)\allowbreak{}=\allowbreak{}p\_\allowbreak{}n},
hence exterior and Adams operations. The
\texttt{\small m\_\allowbreak{}ν} are integral and isobaric of weight
\texttt{\small n} by the fundamental theorem of symmetric polynomials.
This confirms the polynomial claim.

\textbf{Counterexample to the direct-summand reading of the margin.}
Take
\texttt{\small X=\allowbreak{}BGL\_\allowbreak{}\{2,\allowbreak{}\textbackslash{}mathbb Q\}}
and \texttt{\small E} the standard rank-two bundle. In degree two,

\[
\tau^2(E)=\psi^2(E)\rho_2+\Lambda^2(E)\rho_1^2,
\qquad \psi^2(E)=\operatorname{Sym}^2(E)-\Lambda^2(E).
\]

The character of \texttt{\small ψ²(\allowbreak{}E)\allowbreak{}} is
\texttt{\small x\_\allowbreak{}1²+\allowbreak{}x\_\allowbreak{}2²=\allowbreak{}s\_\allowbreak{}(\allowbreak{}2)\allowbreak{}−s\_\allowbreak{}(\allowbreak{}1,\allowbreak{}1)\allowbreak{}}.
A genuine characteristic-zero \texttt{\small GL\_\allowbreak{}2}
representation decomposes into irreducibles with nonnegative
multiplicities; the coefficient \texttt{\small −1} proves that this
virtual class is not such a representation. Therefore the first summand
in the displayed monomial expansion cannot be realized as an actual
direct summand of
\texttt{\small E\textasciicircum{}\{⊗\allowbreak{}2\}}. The exact change
of basis
\texttt{\small ρ\_\allowbreak{}1²=\allowbreak{}1+\allowbreak{}sgn}
cancels that negative term and gives the positive decomposition (A7.2):
\texttt{\small Sym²(\allowbreak{}E)\allowbreak{}⊗\allowbreak{}1 ⊕\allowbreak{}Λ²(\allowbreak{}E)\allowbreak{}⊗\allowbreak{}sgn}.

\textbf{Modular polynomial-ring question reported by SA17.} In exact
\texttt{\small R\_\allowbreak{}\{\textbackslash{}mathbb F\_\allowbreak{}2\}(\allowbreak{}S\_\allowbreak{}2)\allowbreak{}},
the regular module has a filtration with two trivial factors. Hence
\texttt{\small ρ\_\allowbreak{}1²=\allowbreak{}2ρ\_\allowbreak{}2},
refuting polynomial independence of the indicated generators. The map
from the split Green ring to the exact representation ring sends the
indecomposable regular module to this class; the relation is asserted in
the latter. The general integral/ringed-topos version of (A7.1) is now
proved in \textbf{A-012},
\texttt{\small TENSOR\_\allowbreak{}POWER\_\allowbreak{}INTEGRAL.md},
using the ABW integral Cauchy filtration and three exact torus-weight
extractions. It holds for exact K₀ of bounded-rank finite locally free
modules on every commutatively ringed topos; no semisimplicity or
polynomial independence of the Young classes is assumed.

The coefficient-ring question has two further explicit answers.
\textbf{A-021},
\texttt{\small MIXED\_\allowbreak{}CHAR\_\allowbreak{}LATTICES.md},
proves that rationalization gives
\texttt{\small K₀(\allowbreak{}O[\allowbreak{}S\_\allowbreak{}n]\allowbreak{}-\allowbreak{}lattices)\allowbreak{}≅R\_\allowbreak{}K(\allowbreak{}S\_\allowbreak{}n)\allowbreak{}}
for every characteristic-zero discrete valuation ring \texttt{\small O}
with fraction field \texttt{\small K}. The inverse is specified on the
Specht classes, and compatibility with induction yields the polynomial
ring
\texttt{\small Z[\allowbreak{}ρ₁,\allowbreak{}ρ₂,\allowbreak{}…]\allowbreak{}}.
This includes the study's expressly open
\texttt{\small Z\_\allowbreak{}p} case. \textbf{A-023},
\texttt{\small LOCAL\_\allowbreak{}RING\_\allowbreak{}COUNTEREXAMPLE.md},
takes the characteristic-zero local ring
\texttt{\small A=\allowbreak{}Z\_\allowbreak{}(\allowbreak{}2)\allowbreak{}[\allowbreak{}ε]\allowbreak{}/\allowbreak{}(\allowbreak{}2ε,\allowbreak{}ε²)\allowbreak{}}
and the rank-one \texttt{\small S₂} action
\texttt{\small g↦1+\allowbreak{}ε}. Its determinant character lies
outside the subgroup generated by the determinants of \texttt{\small ρ₂}
and \texttt{\small ρ₁²}, proving failure of generation. \textbf{A-020}
also verifies
\texttt{\small rank R\_\allowbreak{}(\allowbreak{}F\_\allowbreak{}p)\allowbreak{}(\allowbreak{}S\_\allowbreak{}p)\allowbreak{}=\allowbreak{}p(\allowbreak{}p)\allowbreak{}−1}
from James's all-field simple-module classification. OriginalPDFpage29
was inspected; the remaining question is the full classification of
coefficient rings.

\textbf{Plethysm consequence.} Schur--Weyl gives positive Schur objects;
the monomial expansion supplies signed universal λ expressions. Further
composition
\texttt{\small s\_\allowbreak{}μ[\allowbreak{}s\_\allowbreak{}ν]\allowbreak{}}
is well-defined in the same free λ-ring. A-007 does not identify a
lattice-point basis for its coefficients. It refutes specifically the
proposed direct-summand reading of (A7.1), while giving the exact basis
morphism that relates the formulas.

The subsequent constructions G-009--G-012 complete the finite
positive-count calculation by explicit invariant-algebra and
coherent-module presentations. In particular, G-010 treats every Schur
plethysm coefficient through the actual wreath-product covariant module,
and G-011 presents the image of the actual canonical-map projector.
G-012 evaluates the outer-degree-three algebra with six generators and
two relations, giving four explicit bounded polytopes. These
constructions and their independent complete proof reviews are recorded
in \texttt{\small ROOT\_\allowbreak{}PROOF\_\allowbreak{}REVIEW.md}; no
generic unevaluated-projector gate remains for terminating finite
positive counts. Bounds polynomial in varying input size and global open
injectivity families require additional results not asserted by those
constructions.

\subsection{A-008. Cote 43, scan 037: Witt-series
formulas}\label{a-008.-cote-43-scan-037-witt-series-formulas}

The page's series ring is represented by
\texttt{\small Spec\textbackslash{}mathbb Z[\allowbreak{}a\_\allowbreak{}1,\allowbreak{}a\_\allowbreak{}2,\allowbreak{}…]\allowbreak{}},
with addition given by multiplication of series. Use the explicit
isomorphism to the conventional big Witt series

\[
\alpha(f)(t)=f(-t)^{-1}.
\]

It sends the page's line element \texttt{\small 1+\allowbreak{}at} to
the conventional Teichmüller series
\texttt{\small (\allowbreak{}1−at)\allowbreak{}\textasciicircum{}\{-\allowbreak{}1\}}.
Transporting the conventional Witt multiplication gives
\texttt{\small (\allowbreak{}1+\allowbreak{}at)\allowbreak{}*f(\allowbreak{}t)\allowbreak{}=\allowbreak{}f(\allowbreak{}at)\allowbreak{}}.
This fixes the convention without identifying two products silently.

For an \texttt{\small \textbackslash{}mathbb F\_\allowbreak{}p}-algebra
\texttt{\small A}, define
\texttt{\small F(\allowbreak{}\textbackslash{}sum a\_\allowbreak{}i t\textasciicircum{}i)\allowbreak{}=\allowbreak{}\textbackslash{}sum a\_\allowbreak{}i\textasciicircum{}p t\textasciicircum{}i}
and
\texttt{\small V(\allowbreak{}f)\allowbreak{}=\allowbreak{}f(\allowbreak{}t\textasciicircum{}p)\allowbreak{}}.
The ring-scheme multiplication has integer coefficients, so coefficient
Frobenius preserves it. Direct multiplication of series in
characteristic \texttt{\small p} gives

\[
f(t)^p=F(f)(t^p),\qquad FV=VF=[p].
\]

Here \texttt{\small [\allowbreak{}p]\allowbreak{}} is multiplication by
\texttt{\small p} in the additive Witt group, hence the ordinary
\texttt{\small p}th power of a series. Also
\texttt{\small V(\allowbreak{}fg)\allowbreak{}=\allowbreak{}V(\allowbreak{}f)\allowbreak{}V(\allowbreak{}g)\allowbreak{}},
proving its additivity. Iterating and applying the line-element rule
yields

\[
(1+x t^{p^i})*f(t)=F^i(f)(x t^{p^i}).
\]

To justify the page's substitution
\texttt{\small x=\allowbreak{}a\textasciicircum{}\{p\textasciicircum{}i\}}
for arbitrary \texttt{\small x}, adjoin \texttt{\small a} using the
monic algebra
\texttt{\small A[\allowbreak{}a]\allowbreak{}/\allowbreak{}(\allowbreak{}a\textasciicircum{}\{p\textasciicircum{}i\}−x)\allowbreak{}}.
It is free as an \texttt{\small A}-module, so \texttt{\small A} injects
into it; the equality checked there descends coefficient by coefficient.
This fills the explicit algebraic step in the page.

The map to \texttt{\small p}-typical Witt vectors is the truncation map
associated with the divisor-stable set
\texttt{\small \{1,\allowbreak{}p,\allowbreak{}p²,\allowbreak{}…\}}
after \texttt{\small α}; its ghost coordinates retain the indices
\texttt{\small p\textasciicircum{}r}, and it sends
\texttt{\small 1+\allowbreak{}at} to the Teichmüller vector
\texttt{\small (\allowbreak{}a,\allowbreak{}0,\allowbreak{}…)\allowbreak{}}.
This is the standard big-to-typical construction, cited rather than
reproved; see Hazewinkel, \emph{Witt vectors. Part 1}, §§14.1--14.8,
pp.71--73. Those paragraphs explicitly establish that the operation
polynomials in coordinate \texttt{\small n} involve only coordinates
whose indices divide \texttt{\small n}, and hence that selecting the
\texttt{\small p}-power coordinates defines the quotient functor on
every commutative ring. The page's assertion that line values determine
the map needs naturality and continuity in the series filtration; those
are part of this specified natural Witt morphism.

\textbf{Plethysm consequence and Borger dictionary.} The representing
symmetric-function algebra is the free special λ-ring. In the character
algebra
\texttt{\small R(\allowbreak{}GL\_\allowbreak{}r)\allowbreak{}=\allowbreak{}\textbackslash{}mathbb Z[\allowbreak{}x\_\allowbreak{}1\textasciicircum{}\{±1\},\allowbreak{}…,\allowbreak{}x\_\allowbreak{}r\textasciicircum{}\{±1\}]\allowbreak{}\textasciicircum{}\{S\_\allowbreak{}r\}},
the Adams map is
\texttt{\small ψ\textasciicircum{}p(\allowbreak{}f)\allowbreak{}=\allowbreak{}f(\allowbreak{}x\_\allowbreak{}1\textasciicircum{}p,\allowbreak{}…,\allowbreak{}x\_\allowbreak{}r\textasciicircum{}p)\allowbreak{}}
and is congruent to \texttt{\small f\textasciicircum{}p} modulo
\texttt{\small p}. This proves the commuting Frobenius-lift property on
\texttt{\small Spec R(\allowbreak{}GL\_\allowbreak{}r)\allowbreak{}}
directly. Borger, supplied \texttt{\small lrfoe13.tex}, §§1.1--1.2,
identifies the free λ-ring with the arithmetic jet space of the affine
line and describes its topos adjunctions; §2.2 gives the monoid-algebra
Frobenius action. These are explicit morphisms connecting the page to
the supplied Borger framework. The character coefficient-extraction
problem remains present in the Schur-basis expansion.

The remaining uniqueness claim is now proved in A-018,
\texttt{\small WITT\_\allowbreak{}UNIQUENESS.md}: finite coordinate
dependence, additivity on universal products of line series, and the
injective map from elementary symmetric functions show that an additive
affine-scheme morphism is determined by its line values. This applies to
the big-to-p-typical map over every characteristic.

\subsection{A-009. Cote 2, scan 017: restriction, induction, external
tensor
products}\label{a-009.-cote-2-scan-017-restriction-induction-external-tensor-products}

For a homomorphism of finite groups
\texttt{\small u:\allowbreak{}H→\allowbreak{}G}, restriction preserves
tensor products. On all modules over a commutative ring
\texttt{\small R}, the induction projection formula has the isomorphism

\[
R[G]\otimes_{R[H]}(M\otimes\operatorname{Res}N)
\longrightarrow (R[G]\otimes_{R[H]}M)\otimes N,
\quad g\otimes(m\otimes n)\longmapsto(g\otimes m)\otimes gn.
\]

Its inverse sends
\texttt{\small (\allowbreak{}g⊗\allowbreak{}m)\allowbreak{}⊗\allowbreak{}n}
to
\texttt{\small g⊗\allowbreak{}(\allowbreak{}m⊗\allowbreak{}g\textasciicircum{}\{-\allowbreak{}1\}n)\allowbreak{}}.
The \texttt{\small H}-balancing relations prove both maps are
well-defined; their displayed composites are identities. Restriction
along \texttt{\small G→\allowbreak{}G×G},
\texttt{\small g↦(\allowbreak{}g,\allowbreak{}g)\allowbreak{}}, takes
\texttt{\small M\textbackslash{}boxtimes N} to
\texttt{\small M⊗\allowbreak{}N}. Extension of scalars identifies
iterated induction with induction along the composite. These formulas
also apply sheafwise.

OriginalPDFpage16/site15 now explicitly verifies the convention: it
defines the all-module category with group action and calls
\texttt{\small u\_\allowbreak{}*} extension, with \texttt{\small u*}
restriction. Originaltargetpage17 and followingpage18 were also
inspected; the displayed formulas and induction product match the
transcription.

For subgroup inclusions, \texttt{\small R[\allowbreak{}G]\allowbreak{}}
is free over \texttt{\small R[\allowbreak{}H]\allowbreak{}}, so
induction preserves finite locally free underlying objects and exact
sequences. For arbitrary group homomorphisms, the corresponding
preservation claim is false:
\texttt{\small C\_\allowbreak{}2→\allowbreak{}1} applied to the rank-one
sign module over \texttt{\small \textbackslash{}mathbb Z} gives
coinvariants \texttt{\small \textbackslash{}mathbb Z/\allowbreak{}2}.
Thus the page's arbitrary-group-map functors are unproblematic on all
modules, but an induced map on the exact
\texttt{\small K\_\allowbreak{}0} of locally free objects needs a stated
category or an admissibility restriction. The exterior/symmetric/Schur
constructions on tensor powers give the finite-operator realization of λ
operations mentioned also on cote12/041.

\subsection{A-010. The λ-ring and γ-filtration
dictionary}\label{a-010.-the-ux3bb-ring-and-ux3b3-filtration-dictionary}

The free-ring fact is SGA 6, V, Proposition 4.2 and Definition 4.4: for
every special λ-ring \texttt{\small R}, evaluation at
\texttt{\small x∈R} is the unique λ-map

\[
\operatorname{ev}_x:\Lambda=\mathbb Z[e_1,e_2,…]\longrightarrow R,
\qquad e_i\longmapsto\lambda^i(x).
\]

Composition of these operations is plethysm. Finite rank is the explicit
quotient
\texttt{\small Λ→\allowbreak{}\textbackslash{}mathbb Z[\allowbreak{}e\_\allowbreak{}1,\allowbreak{}…,\allowbreak{}e\_\allowbreak{}r]\allowbreak{}},
with kernel generated by \texttt{\small e\_\allowbreak{}i},
\texttt{\small i>r}; SGA 6 V, Proposition 4.8 and Corollary 4.9 verify
its λ-ideal property. Inverting \texttt{\small e\_\allowbreak{}r} gives
the full \texttt{\small GL\_\allowbreak{}r} character ring. A
rank-\texttt{\small r} bundle \texttt{\small E} with frame torsor
\texttt{\small P} supplies the homomorphism
\texttt{\small R\_\allowbreak{}\textbackslash{}mathbb Z(\allowbreak{}GL\_\allowbreak{}r)\allowbreak{}→\allowbreak{}K\_\allowbreak{}0(\allowbreak{}X)\allowbreak{}},
\texttt{\small [\allowbreak{}W]\allowbreak{}↦[\allowbreak{}P×\textasciicircum{}\{GL\_\allowbreak{}r\}W]\allowbreak{}},
used in SGA 6 VI, Theorem 3.3. It carries
\texttt{\small S\_\allowbreak{}λ} to \texttt{\small S\_\allowbreak{}λE},
so it carries the universal plethysm identity to the same
\texttt{\small K\_\allowbreak{}0} identity on sheaves.

For the rank augmentation, SGA 6 V 3.2--3.4 defines
\texttt{\small γ\_\allowbreak{}t(\allowbreak{}x)\allowbreak{}=\allowbreak{}λ\_\allowbreak{}\{t/\allowbreak{}(\allowbreak{}1−t)\allowbreak{}\}(\allowbreak{}x)\allowbreak{}}
and V 3.10 defines the filtration generated by products
\texttt{\small \textbackslash{}prod γ\textasciicircum{}\{i\_\allowbreak{}j\}(\allowbreak{}x\_\allowbreak{}j)\allowbreak{}}
with rank-zero \texttt{\small x\_\allowbreak{}j} and
\texttt{\small \textbackslash{}sum i\_\allowbreak{}j≥\allowbreak{}n}. On
a splitting pullback with line classes
\texttt{\small L\_\allowbreak{}1,\allowbreak{}…,\allowbreak{}L\_\allowbreak{}r},
direct substitution gives

\[
\gamma_t(E-r)=\prod_{j=1}^r(1+(L_j-1)t).
\]

The pullback injection and passage to associated graded give the
Chern-class interpretation in SGA 6. This supplies an exact translation
of every universal Schur expression into γ coordinates.

The map induced by
\texttt{\small Spec\textbackslash{}mathbb C→\allowbreak{}BGL\_\allowbreak{}r}
sends a representation to its dimension. In characters it is evaluation
\texttt{\small x\_\allowbreak{}i↦1}. The nonzero element
\texttt{\small [\allowbreak{}V]\allowbreak{}-\allowbreak{}r} lies in its
kernel. Thus ordinary rank at a point does not determine the character,
with the loss measured by this actual map. Equivariant
\texttt{\small K\_\allowbreak{}0(\allowbreak{}BGL\_\allowbreak{}r)\allowbreak{}}
retains the full character. This is the specific distinction needed when
transferring a coefficient statement to geometry; no claim about the
impossibility of other geometric constructions is involved.

\subsection{A-011. Cote 89, scan 018: the exterior pairing in the
rank-two
reconstruction}\label{a-011.-cote-89-scan-018-the-exterior-pairing-in-the-rank-two-reconstruction}

The public original PDFpages16--18 now supply the definitions:
\texttt{\small V} is free of rank two over a commutative ring
\texttt{\small Λ}, and the disjoint line generators
\texttt{\small u,\allowbreak{}v} form a basis. Write
\texttt{\small w=\allowbreak{}au+\allowbreak{}bv}. Then

\[
u\wedge v=v\wedge w=w\wedge u
\quad\Longleftrightarrow\quad a=b=-1
\quad\Longleftrightarrow\quad u+v+w=0.
\]

Indeed the middle two wedges are
\texttt{\small −a(\allowbreak{}u∧v)\allowbreak{}} and
\texttt{\small −b(\allowbreak{}u∧v)\allowbreak{}}, and
\texttt{\small u∧v} is a generator of the free determinant line, so its
coefficients can be compared over \texttt{\small Λ}. Over a field, the
zero vector, the scalar action on each punctured line torsor, and this
reconstruction of addition on independent pairs recover addition on all
pairs: dependent pairs lie on one scalar line. The target explicitly
discusses this field reconstruction. Although supplied scan018 omits
\texttt{\small =\allowbreak{}0} after
\texttt{\small u+\allowbreak{}v+\allowbreak{}w}, the preceding original
PDFpage17/site16, formula(8), writes the full equation. Thus the repair
is supplied by the neighboring manuscript itself. The precise geometric
map is
\texttt{\small \textbackslash{}Lambda² V→\allowbreak{}\textbackslash{}det V},
the identity in rank two. An arbitrary-ring reconstruction of every
nonunimodular vector is not claimed from the field argument.

\subsection{Complete page inventory and remaining
gates}\label{complete-page-inventory-and-remaining-gates}

\begin{longtable}[]{@{}
  >{\raggedright\arraybackslash}p{(\columnwidth - 6\tabcolsep) * \real{0.2500}}
  >{\raggedright\arraybackslash}p{(\columnwidth - 6\tabcolsep) * \real{0.2500}}
  >{\raggedright\arraybackslash}p{(\columnwidth - 6\tabcolsep) * \real{0.2500}}
  >{\raggedright\arraybackslash}p{(\columnwidth - 6\tabcolsep) * \real{0.2500}}@{}}
\toprule\noalign{}
\begin{minipage}[b]{\linewidth}\raggedright
Cote / scan
\end{minipage} & \begin{minipage}[b]{\linewidth}\raggedright
Mathematical content inspected
\end{minipage} & \begin{minipage}[b]{\linewidth}\raggedright
State in this audit
\end{minipage} & \begin{minipage}[b]{\linewidth}\raggedright
What it supplies for plethysm
\end{minipage} \\
\midrule\noalign{}
\endhead
\bottomrule\noalign{}
\endlastfoot
2 / 017 & Tensor pullback, induction projection formula, external
product, diagonal restriction, transitivity, symmetric-group categories
& A-009 proves the displayed maps; arbitrary-map locally-free
preservation refuted by sign coinvariants & The functorial operations
used in wreath-product and Schur constructions \\
4 / 007 & Crystalline vanishing cohomology; an exact sequence marked
``Faux''; a dimension formula; Λ-algebraicity qualification & Public
original PDFpages2--9 recover the definitions. Section8.3 explicitly
admits the crystalline Lefschetz/freeness input; A-014 proves the
stable-part and unit-root linear algebra. The geometric exactness and
Λ-algebraicity inputs remain identified gates & The inspected pages give
no specific plethysm-projector map; G-009--G-012 separately provide
explicit invariant-algebra/module maps and positive counts \\
12 / 041 & \texttt{\small x↦check x ξ\textasciicircum{}i} involution;
expected identity on \texttt{\small σ\_\allowbreak{}i}; duality
\texttt{\small x∨=\allowbreak{}ξ\textasciicircum{}\{−i\}x}; λ through
finite operator groups & OriginalPDFpp31--43 recover the weight
filtration, simple-object and primitive label sets, and Tate class.
Targetp41 expressly reserves identity of the involution to its projected
motivic theory; construction and self-duality morphisms remain the
precise gate & Actual Schur operations through finite symmetric groups
are supplied by A-007/A-009/A-012 \\
16 / 038 & Lefschetz lowering maps \texttt{\small Λ}, kernels, primitive
decompositions, hyperplane pull/push identities & OriginalPDFpp33,34,36
recover definitions. A-016 proves the complete typed string model,
corrects explicit index/type errors, and refutes the last projected
equality by P¹/point & Actual restriction, pushforward and
primitive-projector maps are supplied; no algebraicity assertion is
inferred \\
43 / 037 & Formal series ring, λ operations, coefficient Frobenius,
Verschiebung, \texttt{\small FV=\allowbreak{}VF=\allowbreak{}p}, twisted
line formula, big-to-typical map & A-008 proves the formulas and
conventional-series morphism; A-018 proves additive scheme-morphism
uniqueness from its line values & Adams/Frobenius operations on the
symmetric-function character scheme \\
79 / 051 & Scalar action \texttt{\small T\_\allowbreak{}λ}, quadratic
scaling of \texttt{\small τ}, triangulated-polyhedron automorphisms &
OriginalPDFpp46--51 recover the rank-two Z/n module and wedge model.
A-017 proves the scalar action, determinant-square formula and
compatible structure maps & Exact determinant map supplies the square;
reconstruction of the whole family from one abstract triangulation
remains the source's stated question \\
89 / 018 & Reconstruction from punctured line torsors and the
alternating pairing & OriginalPDFpp16--18 supply rank2, the determinant
torsor and the missing equation. A-011 proves the basis wedge formula
over the ring and full reconstruction over a field & Explicit
exterior-square geometry and scalar-line addition \\
119 / 217 (site216) & Historical account of K-invariants, λ operations,
symmetric-group operations, topoi & OriginalPDFpp215--217 identify
the1972 Esquisse thématique; the target wording matches its typed
reportp3 & The precise mathematical maps are instantiated by
A-007/A-009/A-012 \\
133 / 006 & Quadratic \texttt{\small WG}, rank augmentation, unit,
special-λ universal polynomials, orthogonal classifying-topos reduction
& Original image checked. Parity proved; pre-λ assertion in
characteristic2 refuted (A-002); literal split product identity refuted
(A-003) and iterate identity independently refuted (A-015); exact
quotient theorem cited & Correct exact-\texttt{\small K\_\allowbreak{}0}
target and precise quotient required for Schur identities \\
133 / 008 & Hyperbolic sum, stable-equivalence criterion, ideal, parity
augmentation, λ descent & A-004 proves the invertible2 formula and
refutes the characteristic2 equivalence relation; A-005 proves ideal and
refutes λ descent; A-019 checks field hyperbolic cancellation in every
characteristic & Controls which λ operations survive the quotient \\
133 / 009 & Stiefel--Whitney normalization, sum, tensor/exterior
formulas, completed target, struck-out factorization & A-006 proves the
precise scheme formulation and refutes the struck-out factorization;
arbitrary-topos extension remains gated by definitions & Characteristic
classes of composed Schur objects through explicit root formulas \\
133 / 019 & λ/σ inversion, line formulas, finite-set-looking square
identity, product question & A-001 proves/refutes each specified reading
and corrects SA17 & Universal λ arithmetic with the precise signed
linearization \\
133 / 028 & Tensor powers with symmetric-group action, monomial/Young
expansion, direct-summand margin & A-007 proves the Q-topos formula and
refutes positivity of the monomial summands; A-012 proves the identity
integrally on every commutatively ringed topos in exact K₀ & Exact
integral filtration bridge, while monomial summands can have negative
Schur coefficients \\
\end{longtable}

The study's additional originalPDFpage4 question is answered in
\textbf{A-022}, \texttt{\small HYPERBOLIC\_\allowbreak{}CLASSES.md}:
\texttt{\small w₂(\allowbreak{}H(\allowbreak{}L)\allowbreak{})\allowbreak{}=\allowbreak{}c̄₁(\allowbreak{}L)\allowbreak{}}
and the full higher-rank hyperbolic formula hold over every scheme over
\texttt{\small Z[\allowbreak{}1/\allowbreak{}2]\allowbreak{}}. The proof
constructs actual line-splitting idempotents on successive
nonisotropic-line covers and uses their verified injective cohomology
pullbacks. EKVProposition5.5, pp183--184, supplies the primary formula;
its original regular-scheme scope is recorded separately from this
explicit general proof. The register
\texttt{\small ARCHIVE\_\allowbreak{}QUESTIONS.tsv} now has123
rows:80proved,25refuted,18gates. The remaining gates concern the
expressly admitted crystalline geometric input, the projected motivic
category and self-duality morphisms, literal ill-typed fragments,
reconstruction of a family from one unlabeled triangulation, the
unspecified arbitrary-topos characteristic-class target, and
classification of all coefficient rings. They do not include a generic
unevaluated-plethysm-projector gate.

\subsection{Source notes for the parent
paper}\label{source-notes-for-the-parent-paper}

The user-facing paper should label these as checked findings of this run
or checked reformulations/counterexamples, not claim first publication.
They are independent of the arithmetic RH note. Cited background is
recorded separately in \texttt{\small CHECK\_\allowbreak{}EVENTS.tsv}.

Useful primary references inspected in this run:

\begin{itemize}
\tightlist
\item
  Marcus Zibrowius, \emph{Symmetric representation rings are λ-rings},
  New York J. Math. 21 (2015), 1055--1092. Main theorem p.1056; Lemma1.1
  p.1059; Definition1.2 p.1060; §3.1 p.1068--1069.
  \url{https://nyjm.albany.edu/j/2015/21-49v.pdf}
\item
  Marcus Zibrowius, \emph{The γ-filtration on the Witt ring of a
  scheme}, arXiv:1411.3198v2 (11 August 2016). Theorem2 p.3;
  Proposition2.15 pp.17--18; Proposition3.3 pp.19--20 in the paper's
  displayed pagination (PDF has a one-page-offset ambiguity; use
  statement numbers). \url{https://arxiv.org/pdf/1411.3198}
\item
  SGA6, V2.1, V2.2.1, V2.4.1, V3.2--3.4, V3.10, V4.2, V4.4, V4.8--4.9,
  V§6, VI3.3. Exact numbering checked in the French transcription at
  \texttt{\small grothendiecksga.com} and, for VI3.3/V2.2.1,
  independently confirmed in Zibrowius's primary papers. Original scan
  download from \texttt{\small iut-\allowbreak{}theory.org} failed with
  HTTP503; do not say an original SGA6 scan was visually checked.
\item
  Borger, supplied \texttt{\small lrfoe13.tex}, §§1.1--1.2, §2.2,
  matching the free-ring, Witt-space, and Frobenius-lift statements
  cited above.
\item
  Macdonald and Hazewinkel primary PDFs were retrieved for numbered
  source checks; their formula locators are appended in
  \texttt{\small SOURCE\_\allowbreak{}NOTES.md}.
\end{itemize}

\section{A-012. Integral completion of archive cote 133, scan
028}\label{a-012.-integral-completion-of-archive-cote-133-scan-028}

This completes the exact-K₀ gate left in A-007. The asserted identity
holds on every commutatively ringed topos
\texttt{\small (\allowbreak{}X,\allowbreak{}A)\allowbreak{}}, with
\texttt{\small K₀} formed from finite locally free modules of bounded
rank and short-exact-sequence relations. The integral construction below
specifies both the map and the exact category. It does not assert
polynomial independence of the Young generators in a modular
representation ring.

\subsection{Statement and conventions}\label{statement-and-conventions}

Let
\texttt{\small K\_\allowbreak{}n(\allowbreak{}X)\allowbreak{}=\allowbreak{}K₀(\allowbreak{}BS\_\allowbreak{}(\allowbreak{}n,\allowbreak{}X)\allowbreak{},\allowbreak{}A)\allowbreak{}},
the exact Grothendieck group of finite locally free
\texttt{\small A}-modules with \texttt{\small S\_\allowbreak{}n} action,
and put
\texttt{\small K\_\allowbreak{}*(\allowbreak{}X)\allowbreak{}=\allowbreak{}⊕\allowbreak{}\_\allowbreak{}(\allowbreak{}n≥\allowbreak{}0)\allowbreak{}K\_\allowbreak{}n(\allowbreak{}X)\allowbreak{}}.
Products from degrees \texttt{\small a,\allowbreak{}b} to
\texttt{\small a+\allowbreak{}b} are external tensor product followed by
induction from \texttt{\small S\_\allowbreak{}a×S\_\allowbreak{}b} to
\texttt{\small S\_\allowbreak{}(\allowbreak{}a+\allowbreak{}b)\allowbreak{}}.
A class of \texttt{\small K₀(\allowbreak{}X)\allowbreak{}} acts on each
\texttt{\small K\_\allowbreak{}n(\allowbreak{}X)\allowbreak{}} by tensor
product with trivial \texttt{\small S\_\allowbreak{}n} action.

For a partition
\texttt{\small ν=\allowbreak{}(\allowbreak{}ν\_\allowbreak{}1,\allowbreak{}…,\allowbreak{}ν\_\allowbreak{}l)\allowbreak{}}
of \texttt{\small n}, define the actual Young permutation module

\texttt{\small ρ\_\allowbreak{}ν =\allowbreak{} [\allowbreak{}A[\allowbreak{}S\_\allowbreak{}n/\allowbreak{}(\allowbreak{}S\_\allowbreak{}(\allowbreak{}ν\_\allowbreak{}1)\allowbreak{}×⋯×S\_\allowbreak{}(\allowbreak{}ν\_\allowbreak{}l)\allowbreak{})\allowbreak{}]\allowbreak{}]\allowbreak{}}.

The orbit basis identifies this with the product of the
degree-\texttt{\small ν\_\allowbreak{}i} trivial classes. Let
\texttt{\small m\_\allowbreak{}ν} be the ordinary monomial symmetric
function, expressed as its integral polynomial in the elementary
functions. For every \texttt{\small ξ∈K₀(\allowbreak{}X)\allowbreak{}},

\[
\tau^n(\xi)=\sum_{\nu\vdash n}m_\nu(\xi)\,\rho_\nu
\quad\text{in }K_n(X).
\tag{A12.1}
\]

On an effective class \texttt{\small [\allowbreak{}E]\allowbreak{}},
\texttt{\small τⁿ(\allowbreak{}[\allowbreak{}E]\allowbreak{})\allowbreak{}=\allowbreak{}[\allowbreak{}E\textasciicircum{}⊗\allowbreak{}n]\allowbreak{}}
with the place-permutation action. On virtual classes, \texttt{\small τ}
is the homomorphism from addition in
\texttt{\small K₀(\allowbreak{}X)\allowbreak{}} to the degree-completed
multiplicative group
\texttt{\small 1+\allowbreak{}∏\_\allowbreak{}(\allowbreak{}n>0)\allowbreak{}K\_\allowbreak{}n(\allowbreak{}X)\allowbreak{}}
constructed below.

\subsection{Primary integral input and its exact
translation}\label{primary-integral-input-and-its-exact-translation}

K. Akin, D. A. Buchsbaum, J. Weyman, \emph{Schur functors and Schur
complexes}, Adv. Math. 44 (1982), 207--278, DOI
10.1016/0001-8708(82)90039-1. The original paper was downloaded from the
author's current research page and inspected in full around the relevant
statements. Local copies: \texttt{\small ABW1982.pdf},
\texttt{\small ABW1982.txt}.

\begin{itemize}
\tightlist
\item
  Theorem II.2.16, printed p.232, supplies a universal free basis of
  each Schur module by standard tableaux.
\item
  Theorem III.1.4, printed pp.245--246, supplies a natural filtration of
  \texttt{\small Symⁿ(\allowbreak{}F⊗\allowbreak{}G)\allowbreak{}} by
  \texttt{\small GL(\allowbreak{}F)\allowbreak{}×GL(\allowbreak{}G)\allowbreak{}}
  submodules, with associated graded
  \texttt{\small ⊕\allowbreak{}\_\allowbreak{}(\allowbreak{}η⊢n)\allowbreak{}L\_\allowbreak{}η(\allowbreak{}F)\allowbreak{}⊗\allowbreak{}L\_\allowbreak{}η(\allowbreak{}G)\allowbreak{}}.
\end{itemize}

The paper uses exterior-row notation:
\texttt{\small L\_\allowbreak{}(\allowbreak{}a)\allowbreak{}=\allowbreak{}Λᵃ}.
To match the partition convention of the archive and Macdonald,
\textbf{define}
\texttt{\small S\_\allowbreak{}λ\textasciicircum{}int=\allowbreak{}L\_\allowbreak{}(\allowbreak{}λ′)\allowbreak{}}.
This is an explicit conjugation of the indexing partition. After this
relabeling, its standard tableau basis has weight multiplicities
\texttt{\small K\_\allowbreak{}(\allowbreak{}λν)\allowbreak{}}, the
usual Kostka integers, and the Cauchy associated graded is
\texttt{\small ⊕\allowbreak{}\_\allowbreak{}(\allowbreak{}λ⊢n)\allowbreak{}S\_\allowbreak{}λ\textasciicircum{}int(\allowbreak{}F)\allowbreak{}⊗\allowbreak{}S\_\allowbreak{}λ\textasciicircum{}int(\allowbreak{}G)\allowbreak{}}.

Every graded factor is finite free over Z. The filtration is defined by
natural exterior-diagonalization and symmetric-multiplication maps, so
it is stable under change of bases and scalar extension. Exactness over
Z with free successive quotients is preserved by arbitrary base change.
Applied locally to free sheaves and glued by this naturality, it gives
the same filtration on a commutatively ringed topos. This is the precise
transfer used here.

\subsection{Three weight extractions from the same
filtration}\label{three-weight-extractions-from-the-same-filtration}

Take \texttt{\small W=\allowbreak{}Z\textasciicircum{}n}, with its
standard diagonal torus and permutation-matrix
\texttt{\small S\_\allowbreak{}n}. Write

\texttt{\small C\_\allowbreak{}λ=\allowbreak{}(\allowbreak{}S\_\allowbreak{}λ\textasciicircum{}int(\allowbreak{}W)\allowbreak{})\allowbreak{}\_\allowbreak{}(\allowbreak{}1,\allowbreak{}…,\allowbreak{}1)\allowbreak{}}.

This is a finite free Z-lattice with \texttt{\small S\_\allowbreak{}n}
action: permuting the torus weights fixes the all-ones weight. No
identification of \texttt{\small C\_\allowbreak{}λ} with an irreducible
module is required. Taking a torus weight is exact over Z, since a torus
comodule is a direct sum of its weight modules, by SGA 3 I Proposition
4.7.3.

\textbf{First extraction: the tensor power.} The
\texttt{\small W}-weight
\texttt{\small (\allowbreak{}1,\allowbreak{}…,\allowbreak{}1)\allowbreak{}}
component of
\texttt{\small Symⁿ(\allowbreak{}E⊗\allowbreak{}W)\allowbreak{}} is
equivariantly isomorphic to
\texttt{\small E\textasciicircum{}⊗\allowbreak{}n} through

\texttt{\small e\_\allowbreak{}1⊗\allowbreak{}⋯⊗\allowbreak{}e\_\allowbreak{}n ↦ (\allowbreak{}e\_\allowbreak{}1⊗\allowbreak{}w\_\allowbreak{}1)\allowbreak{}⋯(\allowbreak{}e\_\allowbreak{}n⊗\allowbreak{}w\_\allowbreak{}n)\allowbreak{}}.

Permuting the \texttt{\small w\_\allowbreak{}i} is the place-permutation
action through this map. Extracting that weight from the Cauchy
filtration gives

\[
[E^{\otimes n}]=\sum_{\lambda\vdash n}[S_\lambda^{\mathrm{int}}(E)]\,[C_\lambda]
\quad\text{in }K_n(X).
\tag{A12.2}
\]

\textbf{Second extraction: the integral Young relation.} Apply (A12.2)
to \texttt{\small E=\allowbreak{}Z\textasciicircum{}n}, now equipped
with its own diagonal torus, and select the weight \texttt{\small ν}
padded by zeroes. The tensor-word basis on the left is the set of words
with each letter \texttt{\small j} appearing
\texttt{\small ν\_\allowbreak{}j} times. Its permutation action is the
Young coset action, since the stabilizer of the ordered word is
\texttt{\small ∏\_\allowbreak{}jS\_\allowbreak{}(\allowbreak{}ν\_\allowbreak{}j)\allowbreak{}}.
On the right, the standard tableau basis gives a free weight module of
rank \texttt{\small K\_\allowbreak{}(\allowbreak{}λν)\allowbreak{}}.
Therefore, already over Z,

\[
\rho_\nu=\sum_{\lambda\vdash n}K_{\lambda\nu}[C_\lambda].
\tag{A12.3}
\]

This proof uses the exact filtration; it does not infer integral
equality from rational characters.

\textbf{Third extraction: the class of the integral Schur functor.}
Return to
\texttt{\small Symⁿ(\allowbreak{}E⊗\allowbreak{}W)\allowbreak{}} and
take an arbitrary \texttt{\small W}-weight \texttt{\small ν}. The left
side is the tensor product
\texttt{\small ⊗\allowbreak{}\_\allowbreak{}jSym\textasciicircum{}(\allowbreak{}ν\_\allowbreak{}j)\allowbreak{}(\allowbreak{}E)\allowbreak{}},
and the graded factors on the right have multiplicities
\texttt{\small K\_\allowbreak{}(\allowbreak{}λν)\allowbreak{}}. Hence

\[
h_\nu(E)=\sum_{\lambda\vdash n}K_{\lambda\nu}[S_\lambda^{\mathrm{int}}(E)].
\tag{A12.4}
\]

Here
\texttt{\small [\allowbreak{}Symᵃ(\allowbreak{}E)\allowbreak{}]\allowbreak{}=\allowbreak{}h\_\allowbreak{}a(\allowbreak{}E)\allowbreak{}}
follows from the exterior/symmetric Koszul identity in exact K₀,
including in positive characteristic. The corresponding
symmetric-function identity is
\texttt{\small h\_\allowbreak{}ν=\allowbreak{}Σ\_\allowbreak{}λK\_\allowbreak{}(\allowbreak{}λν)\allowbreak{}s\_\allowbreak{}λ}.
The Kostka matrix is integral triangular with diagonal entries one, so
its transpose is invertible over Z. Comparing and inverting (A12.4)
gives

\texttt{\small [\allowbreak{}S\_\allowbreak{}λ\textasciicircum{}int(\allowbreak{}E)\allowbreak{}]\allowbreak{}=\allowbreak{}s\_\allowbreak{}λ(\allowbreak{}E)\allowbreak{}}.

This proves the needed Schur/λ polynomial dictionary without a
character-injectivity assumption over Z or a semisimplicity assumption.

Now substitute into (A12.2) and use
\texttt{\small s\_\allowbreak{}λ=\allowbreak{}Σ\_\allowbreak{}νK\_\allowbreak{}(\allowbreak{}λν)\allowbreak{}m\_\allowbreak{}ν},
followed by (A12.3):

\texttt{\small Σ\_\allowbreak{}λs\_\allowbreak{}λ(\allowbreak{}E)\allowbreak{}[\allowbreak{}C\_\allowbreak{}λ]\allowbreak{}=\allowbreak{}Σ\_\allowbreak{}νm\_\allowbreak{}ν(\allowbreak{}E)\allowbreak{}Σ\_\allowbreak{}λK\_\allowbreak{}(\allowbreak{}λν)\allowbreak{}[\allowbreak{}C\_\allowbreak{}λ]\allowbreak{}=\allowbreak{}Σ\_\allowbreak{}νm\_\allowbreak{}ν(\allowbreak{}E)\allowbreak{}ρ\_\allowbreak{}ν}.

This proves (A12.1) on effective classes for every commutatively ringed
topos in the declared exact category.

\subsection{Exact-sequence relations and virtual
classes}\label{exact-sequence-relations-and-virtual-classes}

For
\texttt{\small 0→\allowbreak{}E′→\allowbreak{}E→\allowbreak{}E″→\allowbreak{}0},
filter \texttt{\small E\textasciicircum{}⊗\allowbreak{}n} by the number
of factors in \texttt{\small E′}. The associated graded degree with
\texttt{\small a} such factors is

\texttt{\small Ind\_\allowbreak{}(\allowbreak{}S\_\allowbreak{}a×S\_\allowbreak{}(\allowbreak{}n−a)\allowbreak{})\allowbreak{}\textasciicircum{}(\allowbreak{}S\_\allowbreak{}n)\allowbreak{}(\allowbreak{}(\allowbreak{}E′)\allowbreak{}\textasciicircum{}⊗\allowbreak{}a ⊗\allowbreak{} (\allowbreak{}E″)\allowbreak{}\textasciicircum{}⊗\allowbreak{}(\allowbreak{}n−a)\allowbreak{})\allowbreak{}}.

The submodules and quotients are finite locally free: this is checked
after a local splitting of the original sequence, and the filtration
itself is natural and \texttt{\small S\_\allowbreak{}n} stable. Thus,
with the induction product,

\texttt{\small τ(\allowbreak{}E)\allowbreak{}=\allowbreak{}τ(\allowbreak{}E′)\allowbreak{}τ(\allowbreak{}E″)\allowbreak{}}.

Consequently \texttt{\small τ} descends to the exact Grothendieck group
and extends to differences by inversion of a series with constant term
one.

Define
\texttt{\small Φ\_\allowbreak{}n(\allowbreak{}ξ)\allowbreak{}=\allowbreak{}Σ\_\allowbreak{}(\allowbreak{}ν⊢n)\allowbreak{}m\_\allowbreak{}ν(\allowbreak{}ξ)\allowbreak{}ρ\_\allowbreak{}ν}.
The universal alphabet-splitting identity

\texttt{\small m\_\allowbreak{}ν(\allowbreak{}x+\allowbreak{}y)\allowbreak{}=\allowbreak{}Σ\_\allowbreak{}(\allowbreak{}α∪β=\allowbreak{}ν)\allowbreak{}m\_\allowbreak{}α(\allowbreak{}x)\allowbreak{}m\_\allowbreak{}β(\allowbreak{}y)\allowbreak{}}

and the actual induction identity
\texttt{\small ρ\_\allowbreak{}αρ\_\allowbreak{}β=\allowbreak{}ρ\_\allowbreak{}(\allowbreak{}α∪β)\allowbreak{}}
give
\texttt{\small Φ(\allowbreak{}ξ+\allowbreak{}η)\allowbreak{}=\allowbreak{}Φ(\allowbreak{}ξ)\allowbreak{}Φ(\allowbreak{}η)\allowbreak{}}.
The alphabet identity is an integral symmetric-function identity and
evaluates in the special λ-ring
\texttt{\small K₀(\allowbreak{}X)\allowbreak{}} of SGA 6 VI Theorem 3.3.
Thus \texttt{\small Φ} is the same kind of homomorphism as
\texttt{\small τ}. They agree on effective classes by the preceding
proof, hence on their entire group completion. This proves (A12.1) for
all \texttt{\small ξ}.

\subsection{Consequences for the
audit}\label{consequences-for-the-audit}

The polynomial
\texttt{\small Σ\_\allowbreak{}(\allowbreak{}α\_\allowbreak{}1,\allowbreak{}…,\allowbreak{}α\_\allowbreak{}n)\allowbreak{}}
on scan 028 is exactly \texttt{\small m\_\allowbreak{}ν}, with
\texttt{\small ν=\allowbreak{}(\allowbreak{}1\textasciicircum{}α\_\allowbreak{}1⋯n\textasciicircum{}α\_\allowbreak{}n)\allowbreak{}},
in every characteristic under the exact-K₀ convention now fixed. This is
a proved current-run reformulation using the known ABW integral
filtration. The page's direct-summand margin remains refuted by A-007's
rank-two example; the equality established here lives in the exact
Grothendieck group. The modular relation
\texttt{\small ρ\_\allowbreak{}1²=\allowbreak{}2ρ\_\allowbreak{}2} in
exact
\texttt{\small R\_\allowbreak{}(\allowbreak{}F\_\allowbreak{}2)\allowbreak{}(\allowbreak{}S\_\allowbreak{}2)\allowbreak{}}
also remains valid: no independence of the
\texttt{\small ρ\_\allowbreak{}ν} was used in the proof.

The unsigned-subset counterexample A-001 and the split quadratic
counterexample A-003 are unaffected. Their explicit comparison maps
still determine the exact targets of the identities.

\section{A-013. Exact checks of the supplied study's finite-set
additions}\label{a-013.-exact-checks-of-the-supplied-studys-finite-set-additions}

The assertions below are reported by
\texttt{\small SA17\_\allowbreak{}STUDY\_\allowbreak{}cote133\_\allowbreak{}lambda.md}
for cote133 scans20--27. Those scans were not part of the supplied
transcription set. All public original PDFpages20--27 have now been
retrieved and visually read. They confirm the injection/map-power
definitions, the deleted coarse induction formula, the
quotient-partition count, the representable/map-power definitions, the
single-map-power summand assertion, and the distinguished-functor and
eventual-equivalence definitions discussed below. These current
inspections are recorded separately from the earlier supplied study.

\subsection{Injection powers and map
powers}\label{injection-powers-and-map-powers}

For a finite \texttt{\small G}-set \texttt{\small E}, let
\texttt{\small τⁿ(\allowbreak{}E)\allowbreak{}=\allowbreak{}Mon(\allowbreak{}[\allowbreak{}n]\allowbreak{},\allowbreak{}E)\allowbreak{}},
with \texttt{\small G} acting on the target and
\texttt{\small S\_\allowbreak{}n} acting by inverse precomposition. For
\texttt{\small E⊔F}, the inverse image of \texttt{\small E} under an
injection is a subset \texttt{\small A⊂[\allowbreak{}n]\allowbreak{}}.
Fixing \texttt{\small |A|=\allowbreak{}i} and choosing the usual ordered
representative gives the equivariant bijection

\texttt{\small Mon(\allowbreak{}[\allowbreak{}n]\allowbreak{},\allowbreak{}E⊔F)\allowbreak{} ≅ ⊔\_\allowbreak{}(\allowbreak{}i+\allowbreak{}j=\allowbreak{}n)\allowbreak{} Ind\_\allowbreak{}(\allowbreak{}G×S\_\allowbreak{}i×S\_\allowbreak{}j)\allowbreak{}\textasciicircum{}(\allowbreak{}G×S\_\allowbreak{}n)\allowbreak{}(\allowbreak{}Mon(\allowbreak{}[\allowbreak{}i]\allowbreak{},\allowbreak{}E)\allowbreak{}×Mon(\allowbreak{}[\allowbreak{}j]\allowbreak{},\allowbreak{}F)\allowbreak{})\allowbreak{}}.

Explicitly, a coset representative \texttt{\small σ} and the two
injections produce the injection whose restrictions to
\texttt{\small σ(\allowbreak{}[\allowbreak{}i]\allowbreak{})\allowbreak{}}
and its complement are those given injections transported by
\texttt{\small σ⁻¹}. Changing \texttt{\small σ} within
\texttt{\small S\_\allowbreak{}i×S\_\allowbreak{}j} is precisely the
induction balancing relation. The inverse extracts the two restrictions.
This proves multiplicativity of the total injection series in the
induction product. Its degree-zero term is 1. The degree completion has
inverses for every series with that constant term, so the universal
group completion extends it to the Burnside group.

For
\texttt{\small Tⁿ(\allowbreak{}E)\allowbreak{}=\allowbreak{}Hom(\allowbreak{}[\allowbreak{}n]\allowbreak{},\allowbreak{}E)\allowbreak{}},
classify each map by the partition of
\texttt{\small [\allowbreak{}n]\allowbreak{}} into its nonempty fibers.
A type \texttt{\small α} has \texttt{\small α\_\allowbreak{}j} blocks of
size \texttt{\small j}, with
\texttt{\small Σjα\_\allowbreak{}j=\allowbreak{}n} and
\texttt{\small k=\allowbreak{}Σα\_\allowbreak{}j}. Fix one such
partition and a quotient map
\texttt{\small p\_\allowbreak{}α:\allowbreak{}[\allowbreak{}n]\allowbreak{}→\allowbreak{}[\allowbreak{}k]\allowbreak{}}.
Its stabilizer is

\texttt{\small W\_\allowbreak{}α=\allowbreak{}∏\_\allowbreak{}j(\allowbreak{}S\_\allowbreak{}j wr S\_\allowbreak{}(\allowbreak{}α\_\allowbreak{}j)\allowbreak{})\allowbreak{}},

and its action on blocks is the quotient
\texttt{\small W\_\allowbreak{}α→\allowbreak{}H\_\allowbreak{}α=\allowbreak{}∏\_\allowbreak{}j S\_\allowbreak{}(\allowbreak{}α\_\allowbreak{}j)\allowbreak{}⊂S\_\allowbreak{}k}.
The exact component of \texttt{\small Tⁿ(\allowbreak{}E)\allowbreak{}}
is

\texttt{\small Ind\_\allowbreak{}(\allowbreak{}W\_\allowbreak{}α)\allowbreak{}\textasciicircum{}(\allowbreak{}S\_\allowbreak{}n)\allowbreak{} Inf\_\allowbreak{}(\allowbreak{}H\_\allowbreak{}α)\allowbreak{}\textasciicircum{}(\allowbreak{}W\_\allowbreak{}α)\allowbreak{} Res\_\allowbreak{}(\allowbreak{}H\_\allowbreak{}α)\allowbreak{}\textasciicircum{}(\allowbreak{}S\_\allowbreak{}k)\allowbreak{} Mon(\allowbreak{}[\allowbreak{}k]\allowbreak{},\allowbreak{}E)\allowbreak{}}.

The actual bijection sends
\texttt{\small [\allowbreak{}σ,\allowbreak{}i]\allowbreak{}} to
\texttt{\small i∘p\_\allowbreak{}α∘σ⁻¹}; the stabilizer relation is
exactly what makes it well-defined, and recovering the kernel partition
and the injection on its quotient gives the inverse. The number of
source partitions of this type is the index
\texttt{\small n!/\allowbreak{}∏\_\allowbreak{}j(\allowbreak{}α\_\allowbreak{}j!(\allowbreak{}j!)\allowbreak{}\textasciicircum{}(\allowbreak{}α\_\allowbreak{}j)\allowbreak{})\allowbreak{}}.
This proves the study's corrected wreath-stabilizer formula. The deleted
induction-from-\texttt{\small S\_\allowbreak{}i} formula already fails
for \texttt{\small n=\allowbreak{}2}: diagonal maps give a trivial
\texttt{\small S\_\allowbreak{}2} action, whereas induction from
\texttt{\small S\_\allowbreak{}1} gives a free two-point action.

\subsection{Representables, limits and
composition}\label{representables-limits-and-composition}

Let \texttt{\small FI} be the category of finite sets and injections,
and let
\texttt{\small h\_\allowbreak{}I(\allowbreak{}E)\allowbreak{}=\allowbreak{}Mon(\allowbreak{}I,\allowbreak{}E)\allowbreak{}}.
The category of elements of \texttt{\small h\_\allowbreak{}I} has
initial object
\texttt{\small (\allowbreak{}I,\allowbreak{}id)\allowbreak{}}, so
\texttt{\small h\_\allowbreak{}I} is connected. Fiber partitions give a
natural disjoint union

\texttt{\small Hom(\allowbreak{}J,\allowbreak{}E)\allowbreak{}=\allowbreak{}⊔\_\allowbreak{}(\allowbreak{}π partition of J)\allowbreak{} Mon(\allowbreak{}J/\allowbreak{}π,\allowbreak{}E)\allowbreak{}}.

These are precisely its connected components. In particular,
\texttt{\small h\_\allowbreak{}I} is a direct summand of
\texttt{\small Hom(\allowbreak{}I,\allowbreak{}-\allowbreak{})\allowbreak{}}.

The statement equating \textbf{all finite sums of representables} with
direct summands of \textbf{a single}
\texttt{\small Hom(\allowbreak{}I,\allowbreak{}-\allowbreak{})\allowbreak{}}
is false. It appears in the original PDF page25/site24 itself: after the
finite-sum definition, the manuscript says that these are direct
summands of \texttt{\small H\_\allowbreak{}I} for \texttt{\small I} of
sufficiently large cardinality. The constant functor with two elements
is \texttt{\small h\_\allowbreak{}∅⊔h\_\allowbreak{}∅}. At a singleton
\texttt{\small E}, every
\texttt{\small Hom(\allowbreak{}I,\allowbreak{}E)\allowbreak{}} is a
singleton, so it cannot contain that two-element summand, for any
\texttt{\small I}. The corrected equality uses finite sums of the
\texttt{\small Hom(\allowbreak{}I,\allowbreak{}-\allowbreak{})\allowbreak{}}:
each representable is a summand of its own map-power functor, and every
summand of a finite sum of map powers is a union of their connected
representable components. This is a correction to the actual manuscript
assertion as well as to the study's summary.

Finite sums of representables are closed under finite limits. The
terminal object is \texttt{\small h\_\allowbreak{}∅}. The product
\texttt{\small h\_\allowbreak{}I×h\_\allowbreak{}J} is a finite disjoint
union of representables, indexed by the equivalence relations on
\texttt{\small I⊔J} whose restrictions to both summands are equality;
each relation is the overlap pattern of two injections into
\texttt{\small E}. Its quotient gives the corresponding representable.
To check equalizers, restrict two natural maps to a source component
\texttt{\small h\_\allowbreak{}I}. Yoneda identifies each with a choice
of target component \texttt{\small h\_\allowbreak{}J} and an injection
\texttt{\small J→\allowbreak{}I}. Different target components have empty
equalizer. Within the same component, two distinct injections remain
distinct after composition with every injection
\texttt{\small I→\allowbreak{}E}, so their equalizer is empty; identical
injections have the whole component as equalizer.

They are closed under composition as well. In
\texttt{\small h\_\allowbreak{}I(\allowbreak{}⊔\_\allowbreak{}a h\_\allowbreak{}(\allowbreak{}J\_\allowbreak{}a)\allowbreak{}(\allowbreak{}E)\allowbreak{})\allowbreak{}},
first choose for each element of \texttt{\small I} the component
\texttt{\small a} of its image. Then form the finite disjoint union of
the corresponding sets \texttt{\small J\_\allowbreak{}a}. The kernel
partition of its combined map to \texttt{\small E} records every
overlap. Admissible partitions restrict to equality on each
\texttt{\small J\_\allowbreak{}a} and make the resulting injections
distinct whenever two selected components coincide. Each admissible
partition gives exactly the representable of the quotient set. There are
finitely many component choices and partitions. Outer finite coproducts
preserve this decomposition.

The contracted-product description of the distinguished functors is
exact: decompose a finite \texttt{\small S\_\allowbreak{}n}-set
\texttt{\small X\_\allowbreak{}n} into its transitive orbits
\texttt{\small S\_\allowbreak{}n/\allowbreak{}H}, and

\texttt{\small Mon(\allowbreak{}[\allowbreak{}n]\allowbreak{},\allowbreak{}E)\allowbreak{}×\_\allowbreak{}(\allowbreak{}S\_\allowbreak{}n)\allowbreak{}(\allowbreak{}S\_\allowbreak{}n/\allowbreak{}H)\allowbreak{} ≅ Mon(\allowbreak{}[\allowbreak{}n]\allowbreak{},\allowbreak{}E)\allowbreak{}/\allowbreak{}H}.

A formulation by a finite sum requires finite support in
\texttt{\small n}. Arbitrary infinite sequences of finite
\texttt{\small X\_\allowbreak{}n} can still give finite values on each
finite \texttt{\small E}, but are not all finite sums; the powerset
functor is an immediate example.

\subsection{Eventual equivalences and injective transition
maps}\label{eventual-equivalences-and-injective-transition-maps}

For composable natural transformations that are bijective on all sets
beyond specified size thresholds, the set-theoretic two-out-of-three
property applies beyond the maximum threshold. Pullbacks of set
bijections are bijections, proving stability under base change.
Pointwise coproducts and products preserve injective transition maps.
Composition also preserves them: the inner functor sends each injection
to an injection of finite sets, and the outer functor sends that
injection to an injection.

All arguments above have been checked in full. Original PDFpages20--27
were visually inspected in the public-source follow-up. The numerical
quotient-partition count is on originalPDFpage23, and the
eventual-equivalence definition and its two properties are onPDFpage26;
the ledger now uses those precise original page locators.

\section{Public original-scan
follow-up}\label{public-original-scan-follow-up}

Run
\texttt{\small /\allowbreak{}root/\allowbreak{}archive\_\allowbreak{}audit},
GPT-6 Astra Pro, effort not exposed. This follow-up uses the public
University of Montpellier archive directly. No Library attachment was
read. The supplied inputs remain transcriptions; this file identifies
the additional original images actually inspected in this run.

\subsection{Access and exact page
mapping}\label{access-and-exact-page-mapping}

The official catalogue is
\url{https://grothendieck.umontpellier.fr/archives-grothendieck/} . Its
downloaded HTML,
\texttt{\small public\_\allowbreak{}scans/\allowbreak{}catalogue.html},
explicitly supplies viewer paths
\texttt{\small ../\allowbreak{}web/\allowbreak{}viewer.html?file=\allowbreak{}../\allowbreak{}4.pdf},
and likewise for 2,12,16,43,79,119,133. The direct PDF URLs are
therefore source-provided paths, not guessed endpoints. Web-search PDF
opens failed for 4,12,16,133, but ordinary HTTPS access and HTTP
byte-range requests succeeded. Headers identify the files as PDFs last
modified on 11 May 2017 and advertise byte ranges.

Complete cote4 was downloaded to
\texttt{\small public\_\allowbreak{}scans/\allowbreak{}4.pdf} (4,884,318
bytes;12 PDF pages). The catalogue's count11 refers to numbered archive
pages; the file includes an opening page. The supplied scan007 is PDF
page7, whose displayed archive/site number is6. This page mapping was
verified visually, not inferred only from arithmetic.

Targeted byte-range page extraction succeeded for
cotes2,12,16,79,119,133. Complete files4,12,16,43,79,89,119,133 have
also been downloaded.
\texttt{\small extract\_\allowbreak{}public\_\allowbreak{}ranges.py}
preserves the actual selected pages in standalone PDFs; no content is
regenerated. Download attempts are recorded in
\texttt{\small public\_\allowbreak{}scans/\allowbreak{}ACCESS\_\allowbreak{}ATTEMPTS.jsonl}.
An HTTP success is not itself an image-inspection claim. Every one of
the thirteen supplied target leaves has now been viewed as an original
public scan.

\subsection{Cote2, supplied scan017: the extension convention
verified}\label{cote2-supplied-scan017-the-extension-convention-verified}

Visually inspected originalPDFpages16--18, including
targetpage17/site16, rendered at150--160dpi. Page15 is an unrelated
printed verso and was recognized as such. Page16/site15 defines
\texttt{\small A(\allowbreak{}X,\allowbreak{}G)\allowbreak{}} as the
category of modules on the ringed topos with a group action, then
explicitly calls \texttt{\small u\_\allowbreak{}*} extension and
\texttt{\small u*} restriction. This confirms the induction/extension
convention used by A-009, rather than leaving its adjoint choice to an
inference from the symbol. The target's tensor and projection formulas,
external product, diagonal pullback and transitivity match the supplied
transcription. Page18 continues with the induction product along
\texttt{\small S\_\allowbreak{}p×S\_\allowbreak{}q→\allowbreak{}S\_\allowbreak{}(\allowbreak{}p+\allowbreak{}q)\allowbreak{}}.

\subsection{Cote43, supplied scan037: the Witt-series leaf
checked}\label{cote43-supplied-scan037-the-witt-series-leaf-checked}

Visually inspected originalPDFpage37/site36 at150dpi. The inserted
affine qualification, formal-series addition by multiplication, line
convolution, coefficient Frobenius, substitution Verschiebung,
\texttt{\small FV=\allowbreak{}VF=\allowbreak{}p}, twisted-line formula
and final line value of \texttt{\small ψ} match the supplied
transcription. The overwritten input letter of the first displayed
functor formula remains a literal ink issue; the surrounding sentence
explicitly varies the ring \texttt{\small A}. A-008 gives the actual
isomorphism to the standard big-Witt convention, and A-018 proves the
final additive affine-scheme uniqueness claim.

\subsection{Cote4, supplied scan007: context recovered, assumption
isolated}\label{cote4-supplied-scan007-context-recovered-assumption-isolated}

Visually inspected original PDF pages2--9, rendered with Poppler at110
dpi. Files are
\texttt{\small public\_\allowbreak{}scans/\allowbreak{}4-\allowbreak{}p002.png}
through \texttt{\small 4-\allowbreak{}p009.png}.

\begin{itemize}
\tightlist
\item
  PDFp3/sitep2, §8.1: \texttt{\small X} is smooth and proper over a
  perfect field \texttt{\small k} of characteristic \texttt{\small p>0};
  \texttt{\small W} is its Witt ring;
  \texttt{\small H(\allowbreak{}X)\allowbreak{}} is crystalline
  cohomology. The page displays the derived crystalline-to-de-Rham base
  change and the universal-coefficient sequence. It defines
  \texttt{\small E(\allowbreak{}Y)\allowbreak{}=\allowbreak{}coker(\allowbreak{}H\textasciicircum{}(\allowbreak{}n−1)\allowbreak{}(\allowbreak{}X)\allowbreak{}→\allowbreak{}H\textasciicircum{}(\allowbreak{}n−1)\allowbreak{}(\allowbreak{}Y)\allowbreak{})\allowbreak{}}
  for a hyperplane section of a connected, projective
  \texttt{\small n}-dimensional \texttt{\small X}.
\item
  PDFp4/sitep3, §8.3: the argument first invokes integral Lefschetz in
  the complex/ell-adic situation, then \textbf{expressly admits} the
  corresponding crystalline assertions. This is where torsion-freeness
  of \texttt{\small E(\allowbreak{}Y)\allowbreak{}} and the torsion
  restriction isomorphism enter. The underlying missing input is now
  precisely located, rather than described merely as missing notation.
\item
  PDFp5/sitep4, §8.4: \texttt{\small V\_\allowbreak{}ss} is the maximal
  Frobenius-stable subspace on which the semilinear Frobenius is
  bijective, equivalently the stable image of its iterates; Frobenius on
  the quotient is nilpotent. This is a Fitting stable-part construction.
  For example, an invertible unipotent Jordan block over
  \texttt{\small F\_\allowbreak{}p} has its whole space as this stable
  part, explicitly showing what the source terminology means.
\item
  PDFp5/sitep4, §8.4.2: the stable semilinear category is related to
  finite-dimensional \texttt{\small F\_\allowbreak{}p} Galois
  representations. This supplies the definition missing from the
  isolated target page.
\item
  PDFp6/sitep5, Lemma8.4.3: the augmentation from de-Rham cohomology to
  coherent cohomology induces an isomorphism on stable parts. The source
  gives a Hodge-to-de-Rham spectral-sequence argument. An elementary
  factorization proof is supplied below.
\item
  PDFp6/sitep5, the displayed ``Théorème''8.5: the text again explicitly
  admits the preceding crystalline ingredients. This remains an input
  gate for the geometric exact sequence; it is not promoted here to an
  unconditional theorem.
\item
  PDFp7/sitep6: the supplied transcription's p-torsion replacement in
  (8.5.1), the deleted cup-product passage, the marginal
  \texttt{\small Faux}, and the algebraicity qualification on
  \texttt{\small Λ} are all visible. The parenthetical comparison with
  ell-adic eigenvalues retains that stated qualification. The text
  inserts the inverse of the Galois Frobenius generator, consistent with
  the transcription.
\item
  PDFp8/sitep7, Remark8.5.7: injectivity of \texttt{\small v} is itself
  stated as an expectation, followed by a sufficient large-degree
  coherent-vanishing argument. This further confirms that no universal
  injectivity result should be extracted from scan007 alone.
\item
  PDFp9/sitep8: the following discussion turns to families of
  hypersurfaces, stable ranks and generic constancy. No explicit
  plethysm multiplicity or Reynolds-projector map appears in these
  inspected pages.
\end{itemize}

\subsubsection{A-014: the stable-part assertions that can be proved
directly}\label{a-014-the-stable-part-assertions-that-can-be-proved-directly}

Let \texttt{\small V} be a finite-dimensional space over a perfect field
of characteristic \texttt{\small p}, with a Frobenius-semilinear
endomorphism \texttt{\small F}. The descending images
\texttt{\small F\textasciicircum{}r(\allowbreak{}V)\allowbreak{}}
stabilize by finite dimension; write the stabilized image as
\texttt{\small V\_\allowbreak{}st}. Frobenius is bijective there and
nilpotent on \texttt{\small V/\allowbreak{}V\_\allowbreak{}st}. This
decomposition is functorial. It is exact on equivariant short exact
sequences: a sufficiently large common iterate gives surjectivity on
stable images; the intersection of the stable image with an invariant
subspace has bijective Frobenius by finite dimension and is exactly that
subspace's own stable image.

For the algebraic de-Rham complex, let
\texttt{\small π:\allowbreak{}Ω\_\allowbreak{}X\textasciicircum{}•→\allowbreak{}O\_\allowbreak{}X}
be the augmentation, and let
\texttt{\small ι\_\allowbreak{}F:\allowbreak{}O\_\allowbreak{}X→\allowbreak{}Ω\_\allowbreak{}X\textasciicircum{}•}
send a function to its pth power in degree zero. This is a semilinear
chain map because the derivative of a pth power is zero. Their
composites are the Frobenius maps on the two complexes. On
hypercohomology stable parts, \texttt{\small π} is therefore invertible,
with the explicitly typed inverse
\texttt{\small ι\_\allowbreak{}F∘F\_\allowbreak{}O\textasciicircum{}(\allowbreak{}-\allowbreak{}1)\allowbreak{}}
from coherent cohomology to de-Rham cohomology. This proves the source's
Lemma8.4.3 directly in the finite-dimensional situation without
requiring Hodge-to-de-Rham degeneration.

For a finite free
\texttt{\small W(\allowbreak{}F\_\allowbreak{}q)\allowbreak{}}-module
\texttt{\small M} with Frobenius-semilinear \texttt{\small F}, write
\texttt{\small q=\allowbreak{}p\textasciicircum{}f}. The iterate
\texttt{\small T=\allowbreak{}F\textasciicircum{}f} is linear. Its
characteristic polynomial is integral. Its reduction factors as a power
of the variable times a polynomial nonzero at zero; the latter degree
equals the dimension of the Fitting stable part of
\texttt{\small M/\allowbreak{}pM}. Hensel separation of these two
relatively prime factors shows that exactly that many characteristic
roots of \texttt{\small T} are p-adic units, counted with algebraic
multiplicity, and their reductions are the nonzero roots of the reduced
characteristic polynomial. This proves the precise lattice/Frobenius
linear algebra underlying the target page's count. Applying it to the
proposed geometric \texttt{\small E(\allowbreak{}Y)\allowbreak{}} still
requires the freeness and Frobenius-stability input identified in §8.3;
the ell-adic comparison additionally retains the source's algebraicity
input.

\subsection{Cote12, supplied scan041: definitions recovered, projected
theory
retained}\label{cote12-supplied-scan041-definitions-recovered-projected-theory-retained}

Visually inspected original PDFpages31--43, with a260dpi detail of
targetpage41. The foliation is again one less than the PDFpage number.
These are handwritten pages of the motive programme, not a previously
missing independent proof of the target's promised identity.

\begin{itemize}
\tightlist
\item
  PDFp31/site30, §9, explicitly introduces the filtration by weights.
  PDFpp32--34 display tensor-product compatibility and Tate twists
  shifting weight by two.
\item
  PDFpp36--37/site35--36, §10, introduce the effective and Tate-inverted
  categories and their Grothendieck rings
  \texttt{\small M⁺(\allowbreak{}X)\allowbreak{}} and
  \texttt{\small M(\allowbreak{}X)\allowbreak{}}. PDFp37 explicitly
  defines \texttt{\small ξ} as the class of the Tate object written
  \texttt{\small Q\_\allowbreak{}X(\allowbreak{}−1)\allowbreak{}} and
  gives
  \texttt{\small M(\allowbreak{}X)\allowbreak{}=\allowbreak{}M⁺(\allowbreak{}X)\allowbreak{}[\allowbreak{}1/\allowbreak{}ξ]\allowbreak{}}.
  The page states the tensor-category and finite-length/simple-object
  framework being used.
\item
  PDFpp38--40/site37--39 define the simple-object label sets
  \texttt{\small Σ⁺(\allowbreak{}X)\allowbreak{}}, their
  filtration/weight pieces, and the primitive labels
  \texttt{\small σ\_\allowbreak{}i(\allowbreak{}X)\allowbreak{}}
  obtained by removing positive Tate shifts of lower-weight labels. They
  express the additive Grothendieck groups as free abelian groups on
  these labels. Thus the \texttt{\small σ\_\allowbreak{}i} in the target
  are now identified; they are not left as undefined symbols.
\item
  PDFp41/site40 preserves the progression from a twisted-duality
  involution to the claim that it will be the identity. The sentence
  expressly says that the identity will result from the projected
  theory. The words ``théorie projetée'' are readable in the detail
  crop. The coefficient symbol before
  \texttt{\small (\allowbreak{}-\allowbreak{}i)\allowbreak{}} has been
  overwritten and is not clean enough to settle by letter shape alone;
  its Tate-object meaning is supplied by the neighboring definitions.
  The page's final finite-group sentence is a proposed method for the
  lambda operations.
\item
  PDFp42/site41 begins §11 on dimensional type/coniveau with the
  coniveau spectral sequence. PDFp43/site42 includes a crossed-out
  discussion of geometric interpretation and the Tate conjecture, then
  begins §12 on the initial ring of varieties. These following pages do
  not supply the promised proof of the twisted-duality identity.
\end{itemize}

The remaining input is therefore precise: a construction of the stated
motivic category together with a morphism realizing the claimed
self-duality of every primitive simple class, compatibly with the Tate
twist and the chosen equivalence on correspondences. The inspected pages
announce this as part of the projected theory. Merely manipulating the
involution does not establish that it fixes every simple class. The
actual symmetric-group/Schur maps needed for the lambda sentence are
supplied independently in A-007, A-009 and A-012.

\subsection{Cote16, supplied scan038: primitive projectors and concrete
failures}\label{cote16-supplied-scan038-primitive-projectors-and-concrete-failures}

Original PDFpages33,34,36,38 were visually read, including240dpi top and
bottom details of the target. The original full PDF was downloaded.
Pages35 and37 are reused printed versos, recognized as such rather than
mistaken for continuation of the handwritten argument.

PDFp33/site32 supplies the primitive-projector convention and the
coefficient-free Lefschetz lowering operator; the
\texttt{\small p\_\allowbreak{}i} are the primitive bottom/top pieces,
rather than the projectors onto whole cohomological degrees.
PDFp34/site33 supplies the hyperplane restriction strings and the extra
middle piece \texttt{\small E}. PDFp36/site35 supplies
\texttt{\small (\allowbreak{}α)\allowbreak{}} and the pushforward action
on the strings. Thus the earlier missing-definition gate has been
attacked and closed at the linear-algebra level.

\texttt{\small LEFSCHETZ\_\allowbreak{}FRAGMENT.md} (A-016) gives the
full typed model, its restriction and pushforward maps, and verifies the
repaired identities summand by summand. The manuscript target still
contains type/index problems; these are not silently edited into correct
statements. In particular, the last projected equality fails even after
replacing the misplaced pullback by pushforward: for
\texttt{\small X=\allowbreak{}P¹}, \texttt{\small Y=\allowbreak{}point},
one side sends \texttt{\small 1} to \texttt{\small 1} and the other
sends it to \texttt{\small 0}. The transcribed kernel bounds also admit
the explicit \texttt{\small P¹/\allowbreak{}point} and
\texttt{\small P²/\allowbreak{}P¹} counterexamples in A-016.
Algebraicity of the resulting cohomological operators is a further
geometric issue and is not inferred from the string calculation.

\subsection{Cote79, supplied scan051: determinant model and the
remaining reconstruction
question}\label{cote79-supplied-scan051-determinant-model-and-the-remaining-reconstruction-question}

Visually inspected original PDFpages46--51. PDFp47/site46,
formulas(114)--(117), defines a rank-two free module over
\texttt{\small Z/\allowbreak{}n}, its determinant-generator torsor
modulo sign, the scalar group
\texttt{\small (\allowbreak{}Z/\allowbreak{}n)\allowbreak{}×/\allowbreak{}\{±1\}},
and the unimodular-vector set modulo sign. PDFp48/site47,
formulas(118)--(123), defines the wedge-type map and the edge/face sets.
PDFpp49--50/site48--49 supplies \texttt{\small A}, \texttt{\small τ},
and the missing relation(125′).

These definitions permit a complete direct proof of targetformula(127).
\texttt{\small SCALAR\_\allowbreak{}TRIANGULATIONS.md} (A-017) defines
all maps, proves
\texttt{\small τ(\allowbreak{}T\_\allowbreak{}λs,\allowbreak{}T\_\allowbreak{}λt)\allowbreak{}=\allowbreak{}λ²τ(\allowbreak{}s,\allowbreak{}t)\allowbreak{}},
gives the resulting automorphism
\texttt{\small (\allowbreak{}T\_\allowbreak{}λ,\allowbreak{}λ²)\allowbreak{}}
of the family, and proves recovery of the scalar action from the
one-argument relation. The square is explicitly induced by the
determinant map
\texttt{\small Λ²(\allowbreak{}a·id)\allowbreak{}=\allowbreak{}a²·id};
no interpretation of that scalar square as an unspecified lambda
operation is needed.

The final question of scan051 remains exactly the reconstruction from
one abstract, unlabeled triangulation \texttt{\small P₀}. The concrete
module construction and its compatibility proof do not themselves
recover that module and the whole family from one member. The
vertex-action embedding of
\texttt{\small Aut(\allowbreak{}P₀)\allowbreak{}} is directly verified
in A-017.

\subsection{Cote89, supplied scan018: the omitted equation supplied next
door}\label{cote89-supplied-scan018-the-omitted-equation-supplied-next-door}

Visually inspected original PDFpages16--18. PDFp16/site15 defines
\texttt{\small V} as a free rank-two module over a commutative ring
\texttt{\small Λ}, its unimodular vectors, free direct-summand lines,
their generator torsors, and the determinant-generator torsor.
Disjointness is the direct-sum condition. PDFp17/site16,
formulas(7)--(8), explicitly states the wedge equality and its
equivalent equation
\texttt{\small u+\allowbreak{}v+\allowbreak{}w=\allowbreak{}0}. This
supplies the missing right-hand side on targetp18 from the manuscript
itself. The target's concluding reconstruction is expressly discussed
over a field.

The basis computation of A-011 proves the wedge equivalence over the
ring as well: \texttt{\small u,\allowbreak{}v} form a basis by the
disjointness condition, and equality of multiples of their determinant
generator forces the coefficients of \texttt{\small w} to be
\texttt{\small −1,\allowbreak{}−1}. The reconstruction of all addition
in the field case is already proved in A-011. No arbitrary-ring
reconstruction has been inferred from the shorter field argument.

\subsection{Cote119, supplied scan217: document and wording
identified}\label{cote119-supplied-scan217-document-and-wording-identified}

Visually inspected original PDFpages215--217. The title page215/site214
identifies the text as the1972 \textbf{Esquisse thématique des
principaux travaux mathématiques de A. Grothendieck}. Page216 is typed
reportpage2, on homological algebra; targetpage217/site216 is typed
reportpage3, beginning the topology section. Its words about lambda
operations and their links with symmetric-group operations match the
supplied transcription. The paragraph on topoi and ell-adic cohomology
is historical testimony in this report, with its sentence continuing
beyond the leaf. It does not state another precise coefficient theorem.
The mathematical maps supporting the lambda/symmetric-group link are the
explicit maps of A-007, A-009 and A-012.

\subsection{Cote133: literal checks and the supplied study
discrepancy}\label{cote133-literal-checks-and-the-supplied-study-discrepancy}

Visually inspected original PDFpages4,6,8,9 and every page19--29; pages6
and8 also at260dpi details and page19 at300dpi. Newly read
pages4,20--23,26--27,29 were rendered at140dpi.

\begin{itemize}
\tightlist
\item
  Page6/site5 confirms the literal group completion of the
  orthogonal-sum monoid, its pre-lambda assertion, the even-rank
  observation, and the two special-lambda questions. The displayed
  product polynomial ends at the handwritten \texttt{\small i−1}
  exponent. That literal truncation is mathematically false, as
  separately recorded in A-Q063. A-015 now independently refutes the
  full iterate identity for the literal split ring, as A-003 already
  refuted the product identity.
\item
  Page8/site7 supports the handwritten name \textbf{Lazard}, rather than
  the alternative Laborde in the supplied transcription. The hyperbolic
  relation and its warning about a classifying topos are visible. A-004b
  gives the precise characteristic-two counterexample.
\item
  Page9/site8 confirms that the hyperbolic-class vanishing and
  Witt-factorization passage is crossed out. A-006 treats the surviving
  characteristic-class claims and gives the counterexample to that
  deleted passage.
\item
  Page19/site18 visibly gives unsigned subsets and symmetric powers at
  the top, the linearization arrow, the two degree-two formulas, and the
  boxed
  \texttt{\small 2λ²(\allowbreak{}E)\allowbreak{}+\allowbreak{}E=\allowbreak{}E²}.
  This confirms that the supplied study's claimed equivariant check
  cannot be repaired by moving the formula to a different manuscript
  page. The free two-point \texttt{\small C₂}-set of A-001 refutes the
  boxed equivariant equality. Its cardinality image is the valid integer
  equality written underneath. The ambiguous overwritten sign in the
  separate geometric-series denominator is retained as a transcription
  issue; the mathematical inverse-series formula has the minus sign.
\item
  Page24/site23 defines the category of finite sets with injections,
  \texttt{\small h\_\allowbreak{}I=\allowbreak{}Mon(\allowbreak{}I,\allowbreak{}−)\allowbreak{}}
  and
  \texttt{\small H\_\allowbreak{}I=\allowbreak{}Hom(\allowbreak{}I,\allowbreak{}−)\allowbreak{}}.
  Page25/site24 itself asserts the equivalence between finite sums of
  \texttt{\small h\_\allowbreak{}I} and summands of one
  \texttt{\small H\_\allowbreak{}I} for sufficiently large
  \texttt{\small I}. A-013's constant-two counterexample therefore
  refutes a manuscript assertion, as well as its repetition in the
  study. At a singleton input every \texttt{\small H\_\allowbreak{}I}
  has one element.
\item
  Page28/site27 confirms the tensor-power/Young-induction expansion and
  its monomial-symmetric coefficient definition. A-012 proves the
  identity integrally in exact K₀ for arbitrary commutatively ringed
  topoi, using the universal Cauchy filtration and its actual
  multidegree maps. The direct-summand positivity reading is separately
  refuted in A-007; the global positive-count constructions G-009--G-012
  supply actual invariant-algebra/module presentations and terminating
  finite counts.
\item
  Page4/site3 contains the actual \texttt{\small c₁ ?} question, the
  hyperbolic line pair and its first Stiefel--Whitney class. A-022
  proves the line-bundle answer by an explicit nonisotropic-line
  pullback and the Kummer/Bockstein identity. The higher-rank hyperbolic
  formula is located and checked in EKVProposition5.5, pp183--184, with
  the source's regularity assumption recorded. A-022 then proves the
  full formula over every
  \texttt{\small Z[\allowbreak{}1/\allowbreak{}2]\allowbreak{}}-scheme
  by iterating the actual projector
  \texttt{\small z↦v f(\allowbreak{}z)\allowbreak{}/\allowbreak{}f(\allowbreak{}v)\allowbreak{}}
  on nonisotropic-line covers and descending through their injective
  cohomology pullbacks.
\item
  Pages20--21/site19--20 define injection powers and prove their
  disjoint-union induction product. Page22/site21 defines map powers and
  crosses out the induction-from-\texttt{\small S\_\allowbreak{}i}
  formula. Page23/site22 writes the quotient-partition count with its
  block factorials. A-013 gives the actual wreath-stabilizer maps and a
  two-element counterexample to the deleted formula. Page26/site25 gives
  distinguished functors and eventual equivalences; page27/site26
  defines the category with injective transition maps and its sum,
  product and composition.
\item
  Page29/site28 states the coefficient corollary and asks for all rings
  giving the polynomial induction ring, explicitly including local rings
  or rings with \texttt{\small K(\allowbreak{}k)\allowbreak{}≅Z}. A-020
  proves the study's modular rank count from the checked James theorem.
  A-021 resolves the \texttt{\small Z\_\allowbreak{}p}-lattice
  subquestion and extends it to all characteristic-zero DVRs by the
  explicit rationalization map. A-023 gives a characteristic-zero local
  counterexample with a determinant obstruction. These two results
  sharpen, rather than erase, the broader coefficient-ring
  classification question.
\end{itemize}

\subsection{Current assertion
register}\label{current-assertion-register}

\texttt{\small ARCHIVE\_\allowbreak{}QUESTIONS.tsv} has123 separate
source-assertion rows covering all thirteen supplied archive leaves,
plus the relevant lambda/finite-set additions in the supplied
study:80proved,25refuted,18precise gates. Each row records source kind,
cote/page, status, exact reason, and the current proof ID. A-012 closes
the arbitrary-characteristic exact-K₀ tensor-power identity. A-013
records the finite-set checks and the single-map-power summand
counterexample, now verified against the manuscript. A-014--A-017 close
the specified Frobenius, Lefschetz and scalar-action calculations, with
exact remaining geometric/reconstruction gates. A-018 closes uniqueness
of the additive scheme morphism on cote43. A-019 closes field hyperbolic
cancellation in every characteristic; A-020--A-023 supply the modular
rank, DVR polynomial ring, characteristic-class formula and
characteristic-zero local counterexample. Image-inspection fields are
updated only after a page is actually viewed.

\section{A-015. The separate exterior-iterate question on scan133/006 is
refuted}\label{a-015.-the-separate-exterior-iterate-question-on-scan133006-is-refuted}

This supplements A-003's product counterexample. Use exactly its field,
group, quadratic-form convention and split Green ring:

\texttt{\small k=\allowbreak{}algebraic closure of F₃},
\texttt{\small G=\allowbreak{}C₃},
\texttt{\small x=\allowbreak{}[\allowbreak{}J₃]\allowbreak{}},
\texttt{\small y=\allowbreak{}[\allowbreak{}J₂]\allowbreak{}},

\texttt{\small x²=\allowbreak{}3x}, \texttt{\small xy=\allowbreak{}2x},
\texttt{\small y²=\allowbreak{}x+\allowbreak{}1},

\texttt{\small λ\_\allowbreak{}t(\allowbreak{}x)\allowbreak{}=\allowbreak{}1+\allowbreak{}xt+\allowbreak{}xt²+\allowbreak{}t³},
\texttt{\small λ\_\allowbreak{}t(\allowbreak{}y)\allowbreak{}=\allowbreak{}1+\allowbreak{}yt+\allowbreak{}t²}.

The symmetric object \texttt{\small X} is the permutation module
\texttt{\small J₃} with its invariant dot-product form;
\texttt{\small H(\allowbreak{}J₂)\allowbreak{}} is hyperbolic. Take the
actual symmetric object

\texttt{\small A=\allowbreak{}X⊥H(\allowbreak{}J₂)\allowbreak{}}.

Its underlying class is
\texttt{\small a=\allowbreak{}x+\allowbreak{}2y}, of rank seven.
Forgetting symmetric forms commutes with every exterior operation.

\subsection{The exact universal
polynomial}\label{the-exact-universal-polynomial}

The special-lambda iterate polynomial for
\texttt{\small λ³(\allowbreak{}λ² a)\allowbreak{}} is

\texttt{\small Q\_\allowbreak{}(\allowbreak{}3,\allowbreak{}2)\allowbreak{}=\allowbreak{}e₁²e₄−2e₂e₄−e₁e₅+\allowbreak{}e₆+\allowbreak{}e₃²},

where
\texttt{\small e\_\allowbreak{}i=\allowbreak{}λⁱ(\allowbreak{}a)\allowbreak{}}.
Here is a finite universal verification independent of any Green-ring
example. In six indeterminates,
\texttt{\small e₃[\allowbreak{}e₂]\allowbreak{}} is the sum of products
attached to three distinct edges of the complete graph on those six
indices. The possible degree patterns and coefficients of their monomial
symmetric functions are

\begin{longtable}[]{@{}lr@{}}
\toprule\noalign{}
Pattern & Coefficient \\
\midrule\noalign{}
\endhead
\bottomrule\noalign{}
\endlastfoot
\texttt{\small (\allowbreak{}3,\allowbreak{}1,\allowbreak{}1,\allowbreak{}1)\allowbreak{}}
& 1 \\
\texttt{\small (\allowbreak{}2,\allowbreak{}2,\allowbreak{}1,\allowbreak{}1)\allowbreak{}}
& 2 \\
\texttt{\small (\allowbreak{}2,\allowbreak{}2,\allowbreak{}2)\allowbreak{}}
& 1 \\
\texttt{\small (\allowbreak{}2,\allowbreak{}1,\allowbreak{}1,\allowbreak{}1,\allowbreak{}1)\allowbreak{}}
& 6 \\
\texttt{\small (\allowbreak{}1,\allowbreak{}1,\allowbreak{}1,\allowbreak{}1,\allowbreak{}1,\allowbreak{}1)\allowbreak{}}
& 15 \\
\end{longtable}

They correspond respectively to a star, a path, a triangle, an adjacent
pair plus a disjoint edge, and a matching. Expanding the displayed
elementary polynomial gives these same five coefficients.
\texttt{\small check\_\allowbreak{}iterate\_\allowbreak{}defect.py}
verifies the equality of every monomial coefficient using integer
arithmetic. Specialization to six variables is injective in
symmetric-function degree six, since every partition of six has at most
six parts. Thus this is the universal polynomial identity, not only a
test at one numerical alphabet.

\subsection{Evaluation in the actual split Green
ring}\label{evaluation-in-the-actual-split-green-ring}

Multiplication of the exterior series gives

\texttt{\small λ\_\allowbreak{}t(\allowbreak{}a)\allowbreak{}=\allowbreak{}(\allowbreak{}1+\allowbreak{}xt+\allowbreak{}xt²+\allowbreak{}t³)\allowbreak{}(\allowbreak{}1+\allowbreak{}yt+\allowbreak{}t²)\allowbreak{}²},

whose coefficients from degree zero through seven are

\texttt{\small 1,\allowbreak{} x+\allowbreak{}2y,\allowbreak{} 6x+\allowbreak{}3,\allowbreak{} 10x+\allowbreak{}2y+\allowbreak{}1,\allowbreak{} 10x+\allowbreak{}2y+\allowbreak{}1,\allowbreak{} 6x+\allowbreak{}3,\allowbreak{} x+\allowbreak{}2y,\allowbreak{} 1}.

In particular,
\texttt{\small λ²(\allowbreak{}a)\allowbreak{}=\allowbreak{}6x+\allowbreak{}3}.
The degree-three coefficient of

\texttt{\small (\allowbreak{}1+\allowbreak{}xt+\allowbreak{}xt²+\allowbreak{}t³)\allowbreak{}⁶(\allowbreak{}1+\allowbreak{}t)\allowbreak{}³}

is \texttt{\small 441x+\allowbreak{}7}: the degree-three coefficient of
the first factor is \texttt{\small 270x+\allowbreak{}6}, its degree-two
coefficient is \texttt{\small 51x}, and its degree-one coefficient is
\texttt{\small 6x}, giving
\texttt{\small (\allowbreak{}270x+\allowbreak{}6)\allowbreak{}+\allowbreak{}3(\allowbreak{}51x)\allowbreak{}+\allowbreak{}3(\allowbreak{}6x)\allowbreak{}+\allowbreak{}1}.

The universal polynomial above, evaluated at the seven displayed
exterior coefficients, is instead

\texttt{\small 445x−4y+\allowbreak{}3}.

Consequently

\texttt{\small λ³(\allowbreak{}λ²(\allowbreak{}a)\allowbreak{})\allowbreak{}−Q\_\allowbreak{}(\allowbreak{}3,\allowbreak{}2)\allowbreak{}(\allowbreak{}a)\allowbreak{}=\allowbreak{}4+\allowbreak{}4y−4x=\allowbreak{}−4(\allowbreak{}x−y−1)\allowbreak{}}.

The split Green group is free abelian on
\texttt{\small 1,\allowbreak{}y,\allowbreak{}x}, so this difference is
nonzero. Because the forgetful map commutes with every exterior power,
the corresponding iterate identity also fails on the actual symmetric
object \texttt{\small A} in the literal split quadratic Grothendieck
ring of \texttt{\small BC₃}. The coefficients are integers in that
group, not scalars reduced modulo three. The exact quotient sends
\texttt{\small x→\allowbreak{}3}, \texttt{\small y→\allowbreak{}2} and
kills this difference, just as it kills the product defect of A-003.

Thus both universal questions on scan133/006 have separate explicit
negative answers for its stated split ring: A-003 for multiplication,
and A-015 for iterated exterior powers. The modern theorem for the
exact/metabolic quotient remains compatible with both counterexamples
through the explicitly given quotient map.

\section{A-016. Exact operators and a counterexample on cote16,
scan038}\label{a-016.-exact-operators-and-a-counterexample-on-cote16-scan038}

The original Montpellier PDF is
\texttt{\small public\_\allowbreak{}scans/\allowbreak{}16.pdf}
(21,774,525 bytes). The target is PDFpage38/sitepage37. The current run
visually inspected PDFpages33,34,36,38, including a240-dpi crop of the
target's final formulas. PDFpage33/site32 supplies the previously
missing definitions; PDFpage36/site35 supplies the identification
\texttt{\small (\allowbreak{}α)\allowbreak{}} and the pushforward
description. Unrelated printed versos were recognized as such and are
not used in the argument.

\subsection{Definitions recovered from the
original}\label{definitions-recovered-from-the-original}

PDFp33 calls \texttt{\small X} a smooth connected projective variety of
dimension \texttt{\small n}, polarized by a class in
\texttt{\small H²(\allowbreak{}X)\allowbreak{}}, and \texttt{\small Y} a
smooth hyperplane section. The map
\texttt{\small φ:\allowbreak{}Y→\allowbreak{}X} is the immersion. Its
\texttt{\small p\_\allowbreak{}i\textasciicircum{}X} are the projections
onto the primitive \textbf{bottom and top pieces}, indexed from0
through2n. They are not the projectors onto whole cohomological degrees.
The actual operator convention is fixed by the source equations

\texttt{\small Λ\_\allowbreak{}X L\_\allowbreak{}X=\allowbreak{}1−Σ\_\allowbreak{}(\allowbreak{}n≤\allowbreak{}i≤\allowbreak{}2n)\allowbreak{}p\_\allowbreak{}i\textasciicircum{}X},

\texttt{\small L\_\allowbreak{}X Λ\_\allowbreak{}X=\allowbreak{}1−Σ\_\allowbreak{}(\allowbreak{}0≤\allowbreak{}i≤\allowbreak{}n)\allowbreak{}p\_\allowbreak{}i\textasciicircum{}X}.

Thus \texttt{\small Λ} removes one \texttt{\small L} from a primitive
string with coefficient one. PDFp34/site33 identifies restriction with
truncation of each string and defines the extra middle piece
\texttt{\small E\textasciicircum{}(\allowbreak{}n−1)\allowbreak{}}.
PDFp36/site35 writes \texttt{\small (\allowbreak{}α)\allowbreak{}} and
identifies \texttt{\small φ\_\allowbreak{}*} with multiplication by
\texttt{\small L\_\allowbreak{}X} on the shared string pieces, zero on
\texttt{\small E}.

The following explicit construction verifies the resulting operator
statements and makes their maps unambiguous. It also supplies a precise
reading of the source's decompositions. It does not assert algebraicity
of any of these cohomological projectors or of \texttt{\small Λ}.

\subsection{The full explicit string
model}\label{the-full-explicit-string-model}

For a positive integer \texttt{\small n}, choose finite-dimensional
rational vector spaces
\texttt{\small P\_\allowbreak{}0,\allowbreak{}…,\allowbreak{}P\_\allowbreak{}n,\allowbreak{}E}.
Form

\texttt{\small V\_\allowbreak{}X=\allowbreak{}⊕\allowbreak{}\_\allowbreak{}(\allowbreak{}0≤\allowbreak{}i≤\allowbreak{}n)\allowbreak{} ⊕\allowbreak{}\_\allowbreak{}(\allowbreak{}0≤\allowbreak{}j≤\allowbreak{}n−i)\allowbreak{} P\_\allowbreak{}i[\allowbreak{}j]\allowbreak{}},

\texttt{\small V\_\allowbreak{}Y=\allowbreak{}(\allowbreak{}⊕\allowbreak{}\_\allowbreak{}(\allowbreak{}0≤\allowbreak{}i≤\allowbreak{}n−1)\allowbreak{} ⊕\allowbreak{}\_\allowbreak{}(\allowbreak{}0≤\allowbreak{}j≤\allowbreak{}n−1−i)\allowbreak{} P\_\allowbreak{}i[\allowbreak{}j]\allowbreak{})\allowbreak{} ⊕\allowbreak{} E}.

Place \texttt{\small P\_\allowbreak{}i[\allowbreak{}j]\allowbreak{}} in
degree \texttt{\small i+\allowbreak{}2j}, and \texttt{\small E} in
degree \texttt{\small n−1}. On each string define \texttt{\small L} to
move \texttt{\small j} to \texttt{\small j+\allowbreak{}1}, and
\texttt{\small Λ} to move \texttt{\small j} to \texttt{\small j−1}, both
zero at the relevant endpoint; both operators kill \texttt{\small E}.
Define
\texttt{\small r:\allowbreak{}V\_\allowbreak{}X→\allowbreak{}V\_\allowbreak{}Y}
to be the identity on the shared pieces and zero on the last term
\texttt{\small P\_\allowbreak{}i[\allowbreak{}n−i]\allowbreak{}} of
every \texttt{\small X} string. Define
\texttt{\small g:\allowbreak{}V\_\allowbreak{}Y→\allowbreak{}V\_\allowbreak{}X}
by
\texttt{\small g(\allowbreak{}P\_\allowbreak{}i[\allowbreak{}j]\allowbreak{})\allowbreak{}=\allowbreak{}P\_\allowbreak{}i[\allowbreak{}j+\allowbreak{}1]\allowbreak{}},
and \texttt{\small g(\allowbreak{}E)\allowbreak{}=\allowbreak{}0}.

The bottom projector
\texttt{\small p\_\allowbreak{}X\textasciicircum{}i},
\texttt{\small 0≤\allowbreak{}i≤\allowbreak{}n}, projects onto
\texttt{\small P\_\allowbreak{}i[\allowbreak{}0]\allowbreak{}}. For
\texttt{\small n<i≤\allowbreak{}2n}, the top projector
\texttt{\small p\_\allowbreak{}X\textasciicircum{}i} projects onto
\texttt{\small P\_\allowbreak{}(\allowbreak{}2n−i)\allowbreak{}[\allowbreak{}i−n]\allowbreak{}}.
The analogous definitions for \texttt{\small Y} use dimension
\texttt{\small n−1}; its middle projector also contains
\texttt{\small E}. These are explicit direct-sum projectors. In this
model \texttt{\small g r=\allowbreak{}L\_\allowbreak{}X} and
\texttt{\small r g=\allowbreak{}L\_\allowbreak{}Y}, checked on every
displayed summand, exactly as on PDFp36.

The kernel descriptions are

\texttt{\small ker L\_\allowbreak{}Y=\allowbreak{}(\allowbreak{}⊕\allowbreak{}\_\allowbreak{}(\allowbreak{}i=\allowbreak{}0)\allowbreak{}\textasciicircum{}(\allowbreak{}n−1)\allowbreak{} P\_\allowbreak{}i[\allowbreak{}n−1−i]\allowbreak{})\allowbreak{}⊕\allowbreak{}E},

\texttt{\small ker Λ\_\allowbreak{}Y=\allowbreak{}(\allowbreak{}⊕\allowbreak{}\_\allowbreak{}(\allowbreak{}i=\allowbreak{}0)\allowbreak{}\textasciicircum{}(\allowbreak{}n−1)\allowbreak{} P\_\allowbreak{}i[\allowbreak{}0]\allowbreak{})\allowbreak{}⊕\allowbreak{}E}.

The operator \texttt{\small Λ\_\allowbreak{}Y} is injective on all
remaining shared string pieces. The first kernel is related to the high
primitive pieces of \texttt{\small X} by the actual isomorphism

\texttt{\small rΛ\_\allowbreak{}X :\allowbreak{} ⊕\allowbreak{}\_\allowbreak{}(\allowbreak{}j=\allowbreak{}n+\allowbreak{}1)\allowbreak{}\textasciicircum{}(\allowbreak{}2n)\allowbreak{} im(\allowbreak{}p\_\allowbreak{}X\textasciicircum{}j)\allowbreak{} →\allowbreak{} ⊕\allowbreak{}\_\allowbreak{}(\allowbreak{}i=\allowbreak{}0)\allowbreak{}\textasciicircum{}(\allowbreak{}n−1)\allowbreak{} P\_\allowbreak{}i[\allowbreak{}n−1−i]\allowbreak{}}.

Consequently a display identifying these two spaces must use this
degree-minus-two map, rather than a literal equality of subspaces of one
graded object.

\subsection{Corrected operator identities and their complete
proof}\label{corrected-operator-identities-and-their-complete-proof}

Put
\texttt{\small Q\_\allowbreak{}X\textasciicircum{}top=\allowbreak{}Σ\_\allowbreak{}(\allowbreak{}j=\allowbreak{}n+\allowbreak{}1)\allowbreak{}\textasciicircum{}(\allowbreak{}2n)\allowbreak{}p\_\allowbreak{}X\textasciicircum{}j},
\texttt{\small Q\_\allowbreak{}Y\textasciicircum{}top=\allowbreak{}Σ\_\allowbreak{}(\allowbreak{}j=\allowbreak{}n−1)\allowbreak{}\textasciicircum{}(\allowbreak{}2n−2)\allowbreak{}p\_\allowbreak{}Y\textasciicircum{}j},
\texttt{\small Q\_\allowbreak{}X\textasciicircum{}bottom=\allowbreak{}Σ\_\allowbreak{}(\allowbreak{}i=\allowbreak{}0)\allowbreak{}\textasciicircum{}(\allowbreak{}n−1)\allowbreak{}p\_\allowbreak{}X\textasciicircum{}i},
and
\texttt{\small Q\_\allowbreak{}Y\textasciicircum{}bottom=\allowbreak{}Σ\_\allowbreak{}(\allowbreak{}i=\allowbreak{}0)\allowbreak{}\textasciicircum{}(\allowbreak{}n−1)\allowbreak{}p\_\allowbreak{}Y\textasciicircum{}i}.
Then the concrete maps satisfy

\texttt{\small Λ\_\allowbreak{}Y r=\allowbreak{}rΛ\_\allowbreak{}X(\allowbreak{}1−Q\_\allowbreak{}X\textasciicircum{}top)\allowbreak{}=\allowbreak{}(\allowbreak{}1−Q\_\allowbreak{}Y\textasciicircum{}top)\allowbreak{}rΛ\_\allowbreak{}X},

\texttt{\small gΛ\_\allowbreak{}Y=\allowbreak{}(\allowbreak{}1−Q\_\allowbreak{}X\textasciicircum{}bottom)\allowbreak{}Λ\_\allowbreak{}Xg=\allowbreak{}Λ\_\allowbreak{}Xg(\allowbreak{}1−Q\_\allowbreak{}Y\textasciicircum{}bottom)\allowbreak{}}.

For the first line, a nonfinal shared \texttt{\small X} term is lowered
identically on both sides. A final \texttt{\small X} term is killed by
\texttt{\small r}, while \texttt{\small rΛ\_\allowbreak{}X} would send
it to the final \texttt{\small Y} term; either displayed top projector
removes exactly this contribution. The isolated middle
\texttt{\small P\_\allowbreak{}n[\allowbreak{}0]\allowbreak{}} is
already killed by both lowering and restriction. This checks every
summand of \texttt{\small V\_\allowbreak{}X}.

For the second line, on a \texttt{\small Y} term with
\texttt{\small j≥\allowbreak{}1}, both composites land in
\texttt{\small P\_\allowbreak{}i[\allowbreak{}j]\allowbreak{}⊂V\_\allowbreak{}X}.
On a \texttt{\small Y} bottom term, \texttt{\small gΛ\_\allowbreak{}Y}
is zero but \texttt{\small Λ\_\allowbreak{}Xg} returns the corresponding
\texttt{\small X} bottom; either bottom projector removes exactly this
contribution. Every composite kills \texttt{\small E}. This checks every
summand of \texttt{\small V\_\allowbreak{}Y}.

The same basis check gives, for every nonnegative integer
\texttt{\small m},

\texttt{\small Λ\_\allowbreak{}Y\textasciicircum{}m r=\allowbreak{}rΛ\_\allowbreak{}X\textasciicircum{}m(\allowbreak{}1−Q\_\allowbreak{}X\textasciicircum{}top)\allowbreak{}}.

It also gives the correctly typed projector identities

\texttt{\small rΛ\_\allowbreak{}X p\_\allowbreak{}X\textasciicircum{}j=\allowbreak{}p\_\allowbreak{}Y\textasciicircum{}(\allowbreak{}j−2)\allowbreak{} rΛ\_\allowbreak{}X}
for \texttt{\small n+\allowbreak{}1≤\allowbreak{}j≤\allowbreak{}2n},

\texttt{\small p\_\allowbreak{}X\textasciicircum{}i Λ\_\allowbreak{}Xg=\allowbreak{}Λ\_\allowbreak{}Xg p\_\allowbreak{}Y\textasciicircum{}i}
for \texttt{\small 0≤\allowbreak{}i≤\allowbreak{}n−1}.

The first carries the final term of an \texttt{\small X} string onto the
final term of its truncated \texttt{\small Y} string. The second selects
exactly its bottom term. These descriptions prove the identities on
every summand and give the exact relation between the source's indices.

\subsection{The target's final chain fails on an actual hyperplane
section}\label{the-targets-final-chain-fails-on-an-actual-hyperplane-section}

The supplied target's final chain uses \texttt{\small φ*} in positions
whose source must be \texttt{\small H(\allowbreak{}Y)\allowbreak{}}; the
original crop retains the problematic star placement. Even after the
natural type repair replacing those two occurrences by
\texttt{\small g=\allowbreak{}φ\_\allowbreak{}*}, its additional
equality

\texttt{\small p\_\allowbreak{}X\textasciicircum{}i Λ\_\allowbreak{}X φ\_\allowbreak{}* =\allowbreak{} p\_\allowbreak{}X\textasciicircum{}i φ\_\allowbreak{}* Λ\_\allowbreak{}Y}

is false. Take \texttt{\small X=\allowbreak{}P¹\_\allowbreak{}C},
\texttt{\small Y} one point, \texttt{\small n=\allowbreak{}1},
\texttt{\small i=\allowbreak{}0}. Write \texttt{\small h} for the
hyperplane class, so
\texttt{\small H*(\allowbreak{}X,\allowbreak{}Q)\allowbreak{}=\allowbreak{}Q[\allowbreak{}h]\allowbreak{}/\allowbreak{}(\allowbreak{}h²)\allowbreak{}},
\texttt{\small φ\_\allowbreak{}*(\allowbreak{}1)\allowbreak{}=\allowbreak{}h},
\texttt{\small Λ\_\allowbreak{}X(\allowbreak{}h)\allowbreak{}=\allowbreak{}1},
and
\texttt{\small p\_\allowbreak{}X\textasciicircum{}0(\allowbreak{}1)\allowbreak{}=\allowbreak{}1}.
Since \texttt{\small Y} has only degree-zero cohomology,
\texttt{\small Λ\_\allowbreak{}Y=\allowbreak{}0}. Thus the two displayed
maps send \texttt{\small 1} to \texttt{\small 1} and \texttt{\small 0}
respectively. In the full string model the right side is always zero,
while the left side identifies the selected bottom piece. This is a
complete concrete-map counterexample, not only a typing objection.

Two other bounds in the supplied transcription can be settled in the
same elementary models. The expression for
\texttt{\small ker L\_\allowbreak{}Y} must include the middle projector
\texttt{\small p\_\allowbreak{}Y\textasciicircum{}(\allowbreak{}n−1)\allowbreak{}};
the transcribed range starting at \texttt{\small n} fails already for
the same \texttt{\small P¹/\allowbreak{}point} pair. The transcribed
bound \texttt{\small 0≤\allowbreak{}i≤\allowbreak{}n+\allowbreak{}1} for
\texttt{\small ker Λ\_\allowbreak{}Y} includes unwanted top pieces: for
\texttt{\small X=\allowbreak{}P²}, \texttt{\small Y=\allowbreak{}P¹}, it
includes \texttt{\small h∈H²(\allowbreak{}Y)\allowbreak{}}, whereas
\texttt{\small Λ\_\allowbreak{}Y(\allowbreak{}h)\allowbreak{}=\allowbreak{}1}.
The exact correct range is
\texttt{\small 0≤\allowbreak{}i≤\allowbreak{}n−1}. The corresponding
image readings remain distinguished from these mathematical repairs; no
uncertain handwritten sign has been silently normalized.

The source's first corrected pullback line also prints
\texttt{\small p\_\allowbreak{}X} next to
\texttt{\small id\_\allowbreak{}Y}. The explicit model above proves its
well-typed version with \texttt{\small p\_\allowbreak{}Y}. These
corrections and counterexamples close the operator-level questions in
the supplied fragment. They do not supply the algebraicity assertions
about Lefschetz projectors that the neighboring archive discussions
treat separately.

\section{A-017. Cote 79: the exact scalar and determinant
maps}\label{a-017.-cote-79-the-exact-scalar-and-determinant-maps}

The original public PDF pages47--51, particularly (114)--(125′), supply
the definitions missing from supplied scan051. Write
\texttt{\small R=\allowbreak{}Z/\allowbreak{}nZ},
\texttt{\small n≥\allowbreak{}2}, and let \texttt{\small M} be free of
rank two over \texttt{\small R}. Let
\texttt{\small M\_\allowbreak{}prim} be its unimodular vectors,
\texttt{\small D=\allowbreak{}Λ²\_\allowbreak{}R M}, and
\texttt{\small D\_\allowbreak{}gen} the generators of this rank-one
module. The source writes stars for these generator loci. Its objects
are

\texttt{\small U=\allowbreak{}R×/\allowbreak{}\{±1\}},
\texttt{\small S=\allowbreak{}M\_\allowbreak{}prim/\allowbreak{}\{±1\}},
\texttt{\small Π=\allowbreak{}D\_\allowbreak{}gen/\allowbreak{}\{±1\}},
\texttt{\small L′=\allowbreak{}D/\allowbreak{}\{±1\}}.

Here \texttt{\small Π} is a \texttt{\small U}-torsor. The map

\texttt{\small β:\allowbreak{}S×S→\allowbreak{}L′},
\texttt{\small (\allowbreak{}[\allowbreak{}u]\allowbreak{},\allowbreak{}[\allowbreak{}v]\allowbreak{})\allowbreak{}↦[\allowbreak{}u∧v]\allowbreak{}}

is well-defined. Changing either lift changes the wedge by a sign, and
interchanging the lifts also changes it by a sign. Thus β is symmetric
on the quotient, despite being induced by the alternating bilinear wedge
map before quotienting. Define
\texttt{\small A⊂P₂(\allowbreak{}S)\allowbreak{}} by
\texttt{\small β(\allowbreak{}s,\allowbreak{}t)\allowbreak{}∈Π}, and
define \texttt{\small τ:\allowbreak{}A→\allowbreak{}Π} by this value. A
unit wedge already ensures \texttt{\small s≠t}.

For \texttt{\small i∈Π}, set
\texttt{\small A\_\allowbreak{}i=\allowbreak{}τ⁻¹(\allowbreak{}i)\allowbreak{}}
and let \texttt{\small F\_\allowbreak{}i} be the three-element subsets
all of whose two-element subsets belong to
\texttt{\small A\_\allowbreak{}i}. These are the precise edge and face
constructions in (120)--(121), defining the source's family
\texttt{\small Σ(\allowbreak{}i)\allowbreak{}} on its vertex set. The
statements that each is a particular connected regular surface and that
one of them reconstructs all the rest are additional assertions; the
maps below do not assume a classification of those surfaces.

\subsection{The action and its square on determinant
types}\label{the-action-and-its-square-on-determinant-types}

For \texttt{\small λ=\allowbreak{}[\allowbreak{}a]\allowbreak{}∈U},
define
\texttt{\small T\_\allowbreak{}λ(\allowbreak{}[\allowbreak{}u]\allowbreak{})\allowbreak{}=\allowbreak{}[\allowbreak{}au]\allowbreak{}}.
This is independent of both representatives, has inverse
\texttt{\small T\_\allowbreak{}(\allowbreak{}λ⁻¹)\allowbreak{}}, and
satisfies
\texttt{\small T\_\allowbreak{}λT\_\allowbreak{}μ=\allowbreak{}T\_\allowbreak{}(\allowbreak{}λμ)\allowbreak{}}
and \texttt{\small T\_\allowbreak{}1=\allowbreak{}id}. Bilinearity of
the actual wedge map gives

\texttt{\small β(\allowbreak{}s,\allowbreak{}T\_\allowbreak{}λt)\allowbreak{}=\allowbreak{}λβ(\allowbreak{}s,\allowbreak{}t)\allowbreak{}},

and hence

\texttt{\small β(\allowbreak{}T\_\allowbreak{}λs,\allowbreak{}T\_\allowbreak{}λt)\allowbreak{}=\allowbreak{}λ²β(\allowbreak{}s,\allowbreak{}t)\allowbreak{}}.

Multiplication by units preserves generators of \texttt{\small D}, so
\texttt{\small T\_\allowbreak{}λ} preserves \texttt{\small A} and

\texttt{\small τ(\allowbreak{}\{T\_\allowbreak{}λs,\allowbreak{}T\_\allowbreak{}λt\})\allowbreak{}=\allowbreak{}λ²τ(\allowbreak{}\{s,\allowbreak{}t\})\allowbreak{}}.

Consequently the pair of maps
\texttt{\small (\allowbreak{}T\_\allowbreak{}λ,\allowbreak{} i↦λ²i)\allowbreak{}}
is an automorphism of
\texttt{\small (\allowbreak{}S,\allowbreak{}Π,\allowbreak{}A,\allowbreak{}τ)\allowbreak{}}
and transports the actual simplicial structure
\texttt{\small Σ(\allowbreak{}i)\allowbreak{}} to
\texttt{\small Σ(\allowbreak{}λ²i)\allowbreak{}}. Faces are carried to
faces because their three edge types are carried to the common new type.
This proves the target page's formula (127), with every source and
target specified. It also identifies the square in that formula by the
explicit determinant map:
\texttt{\small Λ²(\allowbreak{}a·id\_\allowbreak{}M)\allowbreak{}=\allowbreak{}a²·id\_\allowbreak{}D}.

\subsection{Recovery of the scalar action from the one-argument
relation}\label{recovery-of-the-scalar-action-from-the-one-argument-relation}

The full wedge relation (125) determines each
\texttt{\small T\_\allowbreak{}λ} uniquely. Indeed, primitive
\texttt{\small u} extends to a basis
\texttt{\small (\allowbreak{}u,\allowbreak{}w)\allowbreak{}}. For
\texttt{\small v=\allowbreak{}bu+\allowbreak{}cw}, equality of the
classes of \texttt{\small s∧u} and \texttt{\small s∧v} for all primitive
\texttt{\small s} first gives \texttt{\small c=\allowbreak{}0} by taking
\texttt{\small s=\allowbreak{}u}, then \texttt{\small b=\allowbreak{}±1}
by taking \texttt{\small s=\allowbreak{}w}. Therefore the map
\texttt{\small S→\allowbreak{}(\allowbreak{}L′)\allowbreak{}\textasciicircum{}S}
given by all wedge values is injective.

The restricted relation (125′) using only pairs whose wedge is a unit
also determines each such permutation for every \texttt{\small n≠4}. To
prove this, rescale and change the sign of \texttt{\small v} so that
comparison with \texttt{\small s=\allowbreak{}w} gives
\texttt{\small v=\allowbreak{}u+\allowbreak{}cw}. For every
\texttt{\small t∈R}, the vector \texttt{\small w+\allowbreak{}tu} is
primitive and has unit wedge with \texttt{\small u}; comparison gives
\texttt{\small 1−tc∈\{1,\allowbreak{}−1\}}. The additive subgroup
\texttt{\small cR} is therefore contained in
\texttt{\small \{0,\allowbreak{}2\}}. A nonzero such subgroup forces
\texttt{\small 2} to have additive order two, hence
\texttt{\small n=\allowbreak{}4}. Thus \texttt{\small c=\allowbreak{}0}
for \texttt{\small n≠4}. For \texttt{\small n=\allowbreak{}4}, the
acting group \texttt{\small U} is trivial, so its group action is
uniquely the identity. This proves the unique compatible
\texttt{\small U} action in the actual rank-two module model for all
\texttt{\small n≥\allowbreak{}2}.

The exception also explains why a bare point-separation assertion needs
care. In \texttt{\small R=\allowbreak{}Z/\allowbreak{}4}, the distinct
classes of \texttt{\small (\allowbreak{}1,\allowbreak{}0)\allowbreak{}}
and \texttt{\small (\allowbreak{}1,\allowbreak{}2)\allowbreak{}} have
exactly the same unit-wedge tests and unit-wedge types: the sole type is
the class of \texttt{\small ±1}, and both tests say that the second
coordinate of the testing vector is odd. The full wedge map separates
them by testing
\texttt{\small (\allowbreak{}1,\allowbreak{}0)\allowbreak{}}, whose two
wedge values are \texttt{\small 0} and \texttt{\small 2}.

Conversely, every \texttt{\small i∈Π} is the wedge type of a basis pair.
The expression
\texttt{\small τ(\allowbreak{}\{s,\allowbreak{}T\_\allowbreak{}λt\})\allowbreak{}}
for such a pair of type \texttt{\small i} equals \texttt{\small λi},
independently of the chosen pair. It therefore recovers the torsor
scalar action from the constructed action on vertices. This supplies
actual inverse recovery maps in this module model, rather than an
unproved comparison of structures.

\subsection{What remains a source
question}\label{what-remains-a-source-question}

The last paragraph of scan051 asks how a single abstract triangulation
\texttt{\small P₀=\allowbreak{}(\allowbreak{}S,\allowbreak{}A₀,\allowbreak{}F₀)\allowbreak{}}
determines every other \texttt{\small Σ(\allowbreak{}i)\allowbreak{}}
and the \texttt{\small U} action. The preceding concrete construction
starts with the rank-two module, its determinant line, and its scalar
action; it does not reconstruct them from one unlabeled complex. The
inspected pages do not provide that reconstruction. The displayed
embedding
\texttt{\small Aut(\allowbreak{}P₀)\allowbreak{}→\allowbreak{}Sym(\allowbreak{}S)\allowbreak{}}
is injective directly: an automorphism is the vertex permutation
preserving the specified edge and face subsets, so its action on
vertices determines it.

The entire proof above was checked by expanding the wedge map on a
basis, including the exceptional \texttt{\small n=\allowbreak{}4}
separation calculation. No broader reconstruction theorem is asserted.

\section{A-018. The additive Witt-series morphism is determined by line
values}\label{a-018.-the-additive-witt-series-morphism-is-determined-by-line-values}

This closes the remaining uniqueness assertion on cote43/scan037, in the
precise category written there: morphisms of affine group schemes over
\texttt{\small Z}. Let
\texttt{\small S∞=\allowbreak{}Spec Z[\allowbreak{}a₁,\allowbreak{}a₂,\allowbreak{}…]\allowbreak{}}
represent the group
\texttt{\small 1+\allowbreak{}tA[\allowbreak{}[\allowbreak{}t]\allowbreak{}]\allowbreak{}}
whose group operation is ordinary multiplication of series. Let
\texttt{\small G=\allowbreak{}Spec B} be an affine group scheme, and let
\texttt{\small f,\allowbreak{}g:\allowbreak{}S∞→\allowbreak{}G} be
additive group-scheme morphisms agreeing on the map

\texttt{\small A¹→\allowbreak{}S∞}, \texttt{\small x↦1+\allowbreak{}xt}.

We prove \texttt{\small f=\allowbreak{}g} by their actual coordinate
maps. For each \texttt{\small b∈B}, the difference
\texttt{\small f*(\allowbreak{}b)\allowbreak{}−g*(\allowbreak{}b)\allowbreak{}}
is a polynomial
\texttt{\small P(\allowbreak{}a₁,\allowbreak{}…,\allowbreak{}a\_\allowbreak{}N)\allowbreak{}}
in finitely many coordinates. Over the polynomial ring
\texttt{\small Z[\allowbreak{}x₁,\allowbreak{}…,\allowbreak{}x\_\allowbreak{}N]\allowbreak{}},
form

\texttt{\small u(\allowbreak{}t)\allowbreak{}=\allowbreak{}∏\_\allowbreak{}(\allowbreak{}j=\allowbreak{}1)\allowbreak{}\textasciicircum{}N(\allowbreak{}1+\allowbreak{}x\_\allowbreak{}jt)\allowbreak{}}.

Additivity and agreement on lines give
\texttt{\small f(\allowbreak{}u)\allowbreak{}=\allowbreak{}g(\allowbreak{}u)\allowbreak{}}.
Thus

\texttt{\small P(\allowbreak{}e₁(\allowbreak{}x)\allowbreak{},\allowbreak{}…,\allowbreak{}e\_\allowbreak{}N(\allowbreak{}x)\allowbreak{})\allowbreak{}=\allowbreak{}0}.

The map
\texttt{\small Z[\allowbreak{}a₁,\allowbreak{}…,\allowbreak{}a\_\allowbreak{}N]\allowbreak{}→\allowbreak{}Z[\allowbreak{}x₁,\allowbreak{}…,\allowbreak{}x\_\allowbreak{}N]\allowbreak{}},
\texttt{\small a\_\allowbreak{}j↦e\_\allowbreak{}j(\allowbreak{}x)\allowbreak{}},
is injective. One elementary proof uses lexicographic leading terms: the
leading monomial of
\texttt{\small e₁\textasciicircum{}(\allowbreak{}r₁)\allowbreak{}…e\_\allowbreak{}N\textasciicircum{}(\allowbreak{}r\_\allowbreak{}N)\allowbreak{}}
has exponent vector
\texttt{\small (\allowbreak{}r₁+\allowbreak{}…+\allowbreak{}r\_\allowbreak{}N,\allowbreak{}r₂+\allowbreak{}…+\allowbreak{}r\_\allowbreak{}N,\allowbreak{}…,\allowbreak{}r\_\allowbreak{}N)\allowbreak{}}.
These vectors are distinct for distinct nonnegative tuples
\texttt{\small (\allowbreak{}r₁,\allowbreak{}…,\allowbreak{}r\_\allowbreak{}N)\allowbreak{}},
and each leading coefficient is one, so a nonzero integral polynomial
cannot map to zero. Hence \texttt{\small P=\allowbreak{}0}. This holds
for every \texttt{\small b}, proving equality of the scheme morphisms.

Apply this to the affine p-typical Witt group scheme
\texttt{\small G=\allowbreak{}W\_\allowbreak{}p}. The established
big-to-p-typical ring morphism, in the source convention
\texttt{\small α(\allowbreak{}f)\allowbreak{}=\allowbreak{}1/\allowbreak{}f(\allowbreak{}−t)\allowbreak{}},
sends \texttt{\small 1+\allowbreak{}at} to the Teichmüller vector
\texttt{\small (\allowbreak{}a,\allowbreak{}0,\allowbreak{}0,\allowbreak{}…)\allowbreak{}}.
The preceding argument proves that these line values uniquely determine
its additive scheme morphism, over every base characteristic. It uses
finite dependence of a coordinate polynomial; it introduces no
unsupported continuity assumption and does not identify arbitrary
abstract homomorphisms of groups of series with scheme morphisms.

The full proof, including the coordinate variance, finite-variable
reduction and leading-monomial injection, was checked in this run.

\section{A-019. The field cancellation input on
cote133/008}\label{a-019.-the-field-cancellation-input-on-cote133008}

The original target's field cancellation step is valid for the specified
strictly nondegenerate quadratic forms, including characteristic two.
The exact primary locator is Elman--Karpenko--Merkurjev, \emph{The
Algebraic and Geometric Theory of Quadratic Forms}, §8, Theorem8.4, p48;
its proof uses Theorem8.3 and Lemma8.2, pp47--48. The author-hosted book
is \url{https://www.math.ucla.edu/~merkurev/Book/Kniga-final/Kniga.pdf}
.

In the source's notation, put
\texttt{\small V=\allowbreak{}q⊥(\allowbreak{}−q′)\allowbreak{}}. An
isometry

\texttt{\small V⊥H(\allowbreak{}F)\allowbreak{} ≅ H(\allowbreak{}E)\allowbreak{}}

over a field has \texttt{\small dim E≥\allowbreak{}dim F}. Choose a
vector space \texttt{\small G} of dimension \texttt{\small dim E−dim F},
and identify
\texttt{\small H(\allowbreak{}E)\allowbreak{}≅H(\allowbreak{}G)\allowbreak{}⊥H(\allowbreak{}F)\allowbreak{}}.
The polar form of \texttt{\small H(\allowbreak{}F)\allowbreak{}} is
nondegenerate in every characteristic. The cited cancellation theorem
therefore gives the actual isometry
\texttt{\small V≅H(\allowbreak{}G)\allowbreak{}}, removing the
stabilizer.

The complete proof of Lemma8.2 and Theorems8.3--8.4 was read and checked
in this run, including the reflection calculations, the finite-field
exception \texttt{\small F₂}, and the restriction of an extended
isometry to orthogonal complements. Its characteristic-two hypothesis is
nondegeneracy of the polar form of the cancelled summand, which
\texttt{\small H(\allowbreak{}F)\allowbreak{}} satisfies.

\section{A-020--A-021. The modular rank and the mixed-characteristic
lattice
ring}\label{a-020a-021.-the-modular-rank-and-the-mixed-characteristic-lattice-ring}

These results answer two separate questions in the supplied study of
cote133, scans28--29. Every Grothendieck group below uses all short
exact sequences in the specified category. The maps are explicit; no
assertion about the split group of isomorphism classes is substituted
for this exact Grothendieck group.

\subsection{A-020. The rank over the prime
field}\label{a-020.-the-rank-over-the-prime-field}

For a prime \texttt{\small p}, let
\texttt{\small G₀(\allowbreak{}F\_\allowbreak{}p[\allowbreak{}S\_\allowbreak{}n]\allowbreak{})\allowbreak{}}
be the Grothendieck group of finite-dimensional modules. The
composition-factor map identifies this group with the free abelian group
on the simple modules: Jordan--Hölder makes the map well defined, and
sending a simple label to its module class is its inverse.

James, \emph{The Representation Theory of the Symmetric Groups}, Lecture
Notes in Mathematics682 (1978), Theorem11.5, printedp40, classifies
these simples by the p-regular partitions of \texttt{\small n}. Here
p-regular means that no part occurs p or more times. For
\texttt{\small n=\allowbreak{}p}, the only excluded partition is
\texttt{\small (\allowbreak{}1\textasciicircum{}p)\allowbreak{}}. Hence

\texttt{\small rank\_\allowbreak{}Z G₀(\allowbreak{}F\_\allowbreak{}p[\allowbreak{}S\_\allowbreak{}p]\allowbreak{})\allowbreak{} =\allowbreak{} number\_\allowbreak{}of\_\allowbreak{}partitions(\allowbreak{}p)\allowbreak{} − 1}.

The source's alternative count by p-regular conjugacy classes gives the
same answer: among cycle partitions of \texttt{\small p}, only
\texttt{\small (\allowbreak{}p)\allowbreak{}} is excluded. James
Lemma10.2, p36, proves the equality of these two counts by expanding
\texttt{\small ∏(\allowbreak{}1−t\textasciicircum{}(\allowbreak{}mp)\allowbreak{})\allowbreak{}/\allowbreak{}(\allowbreak{}1−t\textasciicircum{}m)\allowbreak{}}
in the two ways. These are different label sets, linked by that
coefficient identity; their equal cardinalities have not been inferred
by identifying their individual partitions.

\subsection{A-021. Rationalization for lattices over a
characteristic-zero
DVR}\label{a-021.-rationalization-for-lattices-over-a-characteristic-zero-dvr}

Let \texttt{\small O} be a discrete valuation ring of characteristic
zero, with uniformizer \texttt{\small π}, fraction field
\texttt{\small K}, residue field \texttt{\small k}, and
\texttt{\small G=\allowbreak{}S\_\allowbreak{}n}. Write
\texttt{\small L\_\allowbreak{}O(\allowbreak{}G)\allowbreak{}} for the
exact category of finitely generated
\texttt{\small O[\allowbreak{}G]\allowbreak{}}-modules whose underlying
\texttt{\small O}-module is free. Its admissible sequences are the short
exact \texttt{\small O[\allowbreak{}G]\allowbreak{}}-sequences with all
three terms in this category. Rationalization is the exact functor

\texttt{\small L\_\allowbreak{}O(\allowbreak{}G)\allowbreak{} →\allowbreak{} Rep\_\allowbreak{}K(\allowbreak{}G)\allowbreak{},\allowbreak{}    M ↦ K⊗\allowbreak{}\_\allowbreak{}O M}.

It induces an isomorphism

\texttt{\small r\_\allowbreak{}n :\allowbreak{} K₀(\allowbreak{}L\_\allowbreak{}O(\allowbreak{}S\_\allowbreak{}n)\allowbreak{})\allowbreak{} →\allowbreak{} R\_\allowbreak{}K(\allowbreak{}S\_\allowbreak{}n)\allowbreak{}}.

Here is a direct proof, including the boundary construction that
supplies injectivity.

\subsubsection{The torsion boundary}\label{the-torsion-boundary}

For a finite \texttt{\small O}-torsion
\texttt{\small O[\allowbreak{}G]\allowbreak{}}-module \texttt{\small T},
choose a surjection \texttt{\small P→\allowbreak{}T} with
\texttt{\small P} finite free over
\texttt{\small O[\allowbreak{}G]\allowbreak{}}. Its kernel
\texttt{\small L} is finite free over the discrete valuation ring
\texttt{\small O}. Define

\texttt{\small δ(\allowbreak{}T)\allowbreak{}=\allowbreak{}[\allowbreak{}P]\allowbreak{}−[\allowbreak{}L]\allowbreak{} ∈ K₀(\allowbreak{}L\_\allowbreak{}O(\allowbreak{}G)\allowbreak{})\allowbreak{}}.

This class is independent of the chosen presentation. Indeed, for
presentations \texttt{\small P→\allowbreak{}T} and
\texttt{\small P′→\allowbreak{}T}, their fibre product
\texttt{\small Q=\allowbreak{}P×\_\allowbreak{}T P′} lies in the lattice
category. The exact sequences
\texttt{\small 0→\allowbreak{}L→\allowbreak{}Q→\allowbreak{}P′→\allowbreak{}0}
and
\texttt{\small 0→\allowbreak{}L′→\allowbreak{}Q→\allowbreak{}P→\allowbreak{}0}
give
\texttt{\small [\allowbreak{}P]\allowbreak{}−[\allowbreak{}L]\allowbreak{}=\allowbreak{}[\allowbreak{}P′]\allowbreak{}−[\allowbreak{}L′]\allowbreak{}}.
Both sequences split over \texttt{\small O[\allowbreak{}G]\allowbreak{}}
because their quotients are free over
\texttt{\small O[\allowbreak{}G]\allowbreak{}}, although their ordinary
exactness already suffices for the displayed equality.

The assignment is additive on a torsion exact sequence
\texttt{\small 0→\allowbreak{}T₁→\allowbreak{}T₂→\allowbreak{}T₃→\allowbreak{}0}.
Choose free presentations
\texttt{\small P\_\allowbreak{}i→\allowbreak{}T\_\allowbreak{}i} for
\texttt{\small i=\allowbreak{}1,\allowbreak{}3}, lift
\texttt{\small P₃→\allowbreak{}T₃} to \texttt{\small T₂}, and combine
this lift with \texttt{\small P₁→\allowbreak{}T₁→\allowbreak{}T₂}. The
resulting surjection \texttt{\small P₁⊕\allowbreak{}P₃→\allowbreak{}T₂}
has a kernel \texttt{\small L₂} in an exact sequence
\texttt{\small 0→\allowbreak{}L₁→\allowbreak{}L₂→\allowbreak{}L₃→\allowbreak{}0}.
This gives
\texttt{\small δ(\allowbreak{}T₂)\allowbreak{}=\allowbreak{}δ(\allowbreak{}T₁)\allowbreak{}+\allowbreak{}δ(\allowbreak{}T₃)\allowbreak{}}.

More generally, any exact sequence
\texttt{\small 0→\allowbreak{}A→\allowbreak{}B→\allowbreak{}T→\allowbreak{}0}
with \texttt{\small A,\allowbreak{}B} lattices satisfies

\texttt{\small δ(\allowbreak{}T)\allowbreak{}=\allowbreak{}[\allowbreak{}B]\allowbreak{}−[\allowbreak{}A]\allowbreak{}}.

To verify this without choosing a new invariant, form
\texttt{\small Q=\allowbreak{}P×\_\allowbreak{}T B}. Its two lattice
exact sequences are
\texttt{\small 0→\allowbreak{}A→\allowbreak{}Q→\allowbreak{}P→\allowbreak{}0}
and
\texttt{\small 0→\allowbreak{}L→\allowbreak{}Q→\allowbreak{}B→\allowbreak{}0}.
Thus
\texttt{\small [\allowbreak{}A]\allowbreak{}+\allowbreak{}[\allowbreak{}P]\allowbreak{}=\allowbreak{}[\allowbreak{}L]\allowbreak{}+\allowbreak{}[\allowbreak{}B]\allowbreak{}},
which is precisely the formula.

For every lattice \texttt{\small M}, multiplication by \texttt{\small π}
identifies \texttt{\small M} with the lattice \texttt{\small πM}. The
last formula applied to
\texttt{\small 0→\allowbreak{}πM→\allowbreak{}M→\allowbreak{}M/\allowbreak{}πM→\allowbreak{}0}
gives

\texttt{\small δ(\allowbreak{}M/\allowbreak{}πM)\allowbreak{}=\allowbreak{}0}.

\subsubsection{Specht reductions make the boundary
zero}\label{specht-reductions-make-the-boundary-zero}

Let \texttt{\small S\_\allowbreak{}O\textasciicircum{}λ} be the Specht
lattice with its standard integral polytabloid basis, extended from
\texttt{\small Z} to \texttt{\small O}. James Theorem8.4, pp29--30, and
Corollaries8.6--8.11, pp30--31, give its freeness, integer action
matrices and reduction
\texttt{\small k⊗\allowbreak{}S\_\allowbreak{}O\textasciicircum{}λ=\allowbreak{}S\_\allowbreak{}k\textasciicircum{}λ}.

In residue characteristic \texttt{\small p>0}, James Corollary12.2, p42,
shows that for a p-regular \texttt{\small λ}, the class
\texttt{\small [\allowbreak{}S\_\allowbreak{}k\textasciicircum{}λ]\allowbreak{}}
equals
\texttt{\small [\allowbreak{}D\_\allowbreak{}k\textasciicircum{}λ]\allowbreak{}}
plus a sum of classes indexed by strictly dominating p-regular
partitions. Equivalently, Corollary12.3, p43, exhibits the unitriangular
p-regular submatrix of the decomposition matrix. This works over the
present arbitrary residue field: Theorem11.5, p40, explicitly states
that every \texttt{\small D\textasciicircum{}λ} is absolutely
irreducible and every field is a splitting field for
\texttt{\small S\_\allowbreak{}n}. Tensoring a composition series over
\texttt{\small F\_\allowbreak{}p} with \texttt{\small k} is exact and
its simple factors remain simple. The standard polytabloid matrices
identify that tensor product with the Specht module over
\texttt{\small k}. Since the poset is finite, backward substitution over
\texttt{\small Z} expresses every simple class
\texttt{\small [\allowbreak{}D\_\allowbreak{}k\textasciicircum{}λ]\allowbreak{}}
as an integral linear combination of the classes
\texttt{\small [\allowbreak{}S\_\allowbreak{}k\textasciicircum{}μ]\allowbreak{}}.

In residue characteristic zero, the Specht modules themselves form the
complete set of simple
\texttt{\small k[\allowbreak{}S\_\allowbreak{}n]\allowbreak{}}-modules
by the characteristic-zero case of the same theorem. Their classes
already form a basis. This supplies the same generation conclusion
without a modular reduction argument.

The boundary vanishes on every
\texttt{\small [\allowbreak{}S\_\allowbreak{}k\textasciicircum{}μ]\allowbreak{}},
because it is the reduction of
\texttt{\small S\_\allowbreak{}O\textasciicircum{}μ}. Its additivity
therefore makes it vanish on every finite-dimensional
\texttt{\small k[\allowbreak{}G]\allowbreak{}}-module. A finite
\texttt{\small O}-torsion module has the finite filtration by powers of
\texttt{\small π}, whose successive quotients are finite-dimensional
\texttt{\small k[\allowbreak{}G]\allowbreak{}}-modules. Consequently

\texttt{\small δ(\allowbreak{}T)\allowbreak{}=\allowbreak{}0} for every
finite \texttt{\small O}-torsion
\texttt{\small O[\allowbreak{}G]\allowbreak{}}-module \texttt{\small T}.

\subsubsection{Commensurability and the inverse on Grothendieck
groups}\label{commensurability-and-the-inverse-on-grothendieck-groups}

Suppose lattices \texttt{\small M,\allowbreak{}N} have isomorphic
rationalizations. Use an actual
\texttt{\small K[\allowbreak{}G]\allowbreak{}}-isomorphism to place them
inside one finite-dimensional
\texttt{\small K[\allowbreak{}G]\allowbreak{}}-module. The intersection
\texttt{\small I=\allowbreak{}M∩N} is a \texttt{\small G}-stable
lattice, and both \texttt{\small M/\allowbreak{}I} and
\texttt{\small N/\allowbreak{}I} are finite \texttt{\small O}-torsion.
Applying the boundary formula gives
\texttt{\small [\allowbreak{}M]\allowbreak{}=\allowbreak{}[\allowbreak{}I]\allowbreak{}=\allowbreak{}[\allowbreak{}N]\allowbreak{}}
in the exact lattice group.

The algebra
\texttt{\small K[\allowbreak{}S\_\allowbreak{}n]\allowbreak{}} is split
semisimple. Its irreducibles are the rationalizations of the Specht
lattices. These facts prove surjectivity of
\texttt{\small r\_\allowbreak{}n}. An element of its kernel has the form
\texttt{\small [\allowbreak{}M]\allowbreak{}−[\allowbreak{}N]\allowbreak{}};
equality of the rationalized classes in a semisimple category gives an
actual rational isomorphism. The intersection argument proves
\texttt{\small [\allowbreak{}M]\allowbreak{}=\allowbreak{}[\allowbreak{}N]\allowbreak{}},
and hence injectivity. More explicitly, the inverse sends the rational
Specht class
\texttt{\small [\allowbreak{}S\_\allowbreak{}K\textasciicircum{}λ]\allowbreak{}}
to
\texttt{\small [\allowbreak{}S\_\allowbreak{}O\textasciicircum{}λ]\allowbreak{}};
the preceding argument verifies that this inverse also recovers the
class of every lattice.

\subsubsection{Induction multiplication}\label{induction-multiplication}

Form the graded group
\texttt{\small R\_\allowbreak{}O=\allowbreak{}⊕\allowbreak{}\_\allowbreak{}\{n≥\allowbreak{}0\}K₀(\allowbreak{}L\_\allowbreak{}O(\allowbreak{}S\_\allowbreak{}n)\allowbreak{})\allowbreak{}},
with product

\texttt{\small [\allowbreak{}M]\allowbreak{}·[\allowbreak{}N]\allowbreak{}=\allowbreak{}[\allowbreak{}Ind\_\allowbreak{}\{S\_\allowbreak{}a×S\_\allowbreak{}b\}\textasciicircum{}\{S\_\allowbreak{}\{a+\allowbreak{}b\}\}(\allowbreak{}M⊗\allowbreak{}\_\allowbreak{}O N)\allowbreak{}]\allowbreak{}}.

Tensor product is exact in each variable because the underlying modules
are free. The group algebra of
\texttt{\small S\_\allowbreak{}\{a+\allowbreak{}b\}} is free over that
of \texttt{\small S\_\allowbreak{}a×S\_\allowbreak{}b} on a coset basis,
so induction is exact and preserves lattices. Rationalization commutes
with these functors through the canonical tensor and induction
isomorphisms. Thus \texttt{\small ⊕\allowbreak{}r\_\allowbreak{}n} is an
isomorphism of graded rings.

The rational Frobenius characteristic sends the trivial representation
\texttt{\small ρ\_\allowbreak{}n} to the complete symmetric function
\texttt{\small h\_\allowbreak{}n}. Since
\texttt{\small Λ=\allowbreak{}Z[\allowbreak{}h₁,\allowbreak{}h₂,\allowbreak{}…]\allowbreak{}},
this yields the concrete generator-preserving isomorphism

\texttt{\small Z[\allowbreak{}ρ₁,\allowbreak{}ρ₂,\allowbreak{}…]\allowbreak{} ≅ R\_\allowbreak{}O}.

This includes \texttt{\small O=\allowbreak{}Z\_\allowbreak{}p} and
\texttt{\small O=\allowbreak{}Z\_\allowbreak{}(\allowbreak{}p)\allowbreak{}}
and resolves the study's expressly open
\texttt{\small Z\_\allowbreak{}p}-lattice case. No completeness or
unramifiedness assumption is needed. The proof uses that submodules of
finite free \texttt{\small O}-modules are free and that finite torsion
modules admit the displayed residue-field filtration. The exact boundary
of these steps is substantive: A-023 in
\texttt{\small LOCAL\_\allowbreak{}RING\_\allowbreak{}COUNTEREXAMPLE.md}
gives a characteristic-zero local ring with a rank-one representation
outside the subgroup generated by the Young permutation classes. The
broader classification of coefficient rings remains a question.

\subsection{Primary source and inspection
record}\label{primary-source-and-inspection-record}

James's original book: DOI
\texttt{\small 10.1007/\allowbreak{}BFb0067708}; public
university-hosted scan:
\url{https://www-users.cse.umn.edu/~webb/oldteaching/Year2010-11/the-representation-theory-of-the-symmetric-groups-SLN.pdf}
. The cited pages29--31,36,39--40,42--43 were read. The standard-basis
proof, the stated simple-module classification and its proof
dependencies, and the unitriangularity argument were checked. The
classification theorem's count invokes general modular representation
theory explicitly cited by James; this report does not claim a new proof
of that theory. The torsion boundary, pullback comparisons,
commensurability argument and graded-ring compatibility above were
independently checked in full in this run.

\section{A-022. The hyperbolic Stiefel--Whitney
formula}\label{a-022.-the-hyperbolic-stiefelwhitney-formula}

Write
\texttt{\small α∈H¹\_\allowbreak{}et(\allowbreak{}X,\allowbreak{}Z/\allowbreak{}2)\allowbreak{}}
for the Kummer class of \texttt{\small −1}, and let
\texttt{\small c̄\_\allowbreak{}j} denote the mod2 étale Chern class. For
a rank-r vector bundle \texttt{\small E}, the formula appearing in the
supplied study is

\texttt{\small w\_\allowbreak{}t(\allowbreak{}H(\allowbreak{}E)\allowbreak{})\allowbreak{} =\allowbreak{} Σ\_\allowbreak{}\{j=\allowbreak{}0\}\textasciicircum{}r (\allowbreak{}1+\allowbreak{}αt)\allowbreak{}\textasciicircum{}(\allowbreak{}r−j)\allowbreak{} c̄\_\allowbreak{}j(\allowbreak{}E)\allowbreak{}t\textasciicircum{}(\allowbreak{}2j)\allowbreak{}}.

The displayed hyperbolic formula holds for every scheme over
\texttt{\small Z[\allowbreak{}1/\allowbreak{}2]\allowbreak{}}, by the
explicit splitting morphisms proved below. The exact primary locator for
the formula is Esnault--Kahn--Viehweg, \emph{Coverings with odd
ramification and Stiefel--Whitney classes}, Journal für die reine und
angewandte Mathematik441 (1993), Proposition5.5, pp183--184. The actual
author-hosted scan is
\url{https://page.mi.fu-berlin.de/esnault/preprints/helene/29-es_ka_vi.pdf}
. Section5.1, p179, starts with a scheme over
\texttt{\small Z[\allowbreak{}1/\allowbreak{}2]\allowbreak{}};
section5.4, p182, introduces a regularity assumption. Thus the source
directly establishes the displayed formula for regular schemes over
\texttt{\small Z[\allowbreak{}1/\allowbreak{}2]\allowbreak{}}. Its
result also covers metabolic bundles with the Chern classes of a
Lagrangian. The proof here establishes the stated hyperbolic formula
over arbitrary
\texttt{\small Z[\allowbreak{}1/\allowbreak{}2]\allowbreak{}}-schemes.
No characteristic-two version is claimed.

Both pages of the proposition and proof were visually read. The proof
first computes the rank-two case using the complement of the isotropic
quadric, then constructs an affine-bundle pullback splitting off a
metabolic rank-two summand. The remaining Lagrangian fits in a
vector-bundle exact sequence, and the Chern-class Whitney formula gives
the displayed sum by induction. The injective étale-cohomology pullback
descends the identity. Zibrowius, \emph{Witt groups of curves and
surfaces}, §1.4, equations(6)--(8), printedp8 (PDFpage9) in the author
PDF, reproduces the formula and its first two components with this exact
EKV locator. That exposition explicitly treats schemes over
\texttt{\small Z[\allowbreak{}1/\allowbreak{}2]\allowbreak{}}; the
regularity inherited from EKV§5.4 is recorded here rather than silently
attributed away.

In particular,

\texttt{\small w₁(\allowbreak{}H(\allowbreak{}E)\allowbreak{})\allowbreak{}=\allowbreak{}rα},

\texttt{\small w₂(\allowbreak{}H(\allowbreak{}E)\allowbreak{})\allowbreak{}=\allowbreak{}binom(\allowbreak{}r,\allowbreak{}2)\allowbreak{}α²+\allowbreak{}c̄₁(\allowbreak{}E)\allowbreak{}}.

For a line bundle \texttt{\small L}, the latter reduces to
\texttt{\small w₂(\allowbreak{}H(\allowbreak{}L)\allowbreak{})\allowbreak{}=\allowbreak{}c̄₁(\allowbreak{}L)\allowbreak{}}.
This line-bundle identity admits the following direct proof over every
scheme over
\texttt{\small Z[\allowbreak{}1/\allowbreak{}2]\allowbreak{}}, without
the regularity assumption just identified in the source.

\subsection{The explicit line-bundle
calculation}\label{the-explicit-line-bundle-calculation}

Let \texttt{\small p:\allowbreak{}U→\allowbreak{}X} be the open
complement of the isotropic quadric in the projective bundle of lines in
\texttt{\small H(\allowbreak{}L)\allowbreak{}=\allowbreak{}L⊕\allowbreak{}L∨}.
Its tautological line \texttt{\small M} is nondegenerate. The first
projection \texttt{\small M→\allowbreak{}p*L} is an isomorphism: at
every geometric point the quadratic value of a tautological vector
\texttt{\small (\allowbreak{}v,\allowbreak{}f)\allowbreak{}} is
\texttt{\small f(\allowbreak{}v)\allowbreak{}≠0}, so both components are
nonzero. A fibrewise nonzero map of line bundles is invertible.
Orthogonal projection gives the actual decomposition

\texttt{\small p*H(\allowbreak{}L)\allowbreak{}=\allowbreak{}M ⊥ M\textasciicircum{}⊥}.

Set \texttt{\small x=\allowbreak{}w₁(\allowbreak{}M)\allowbreak{}}. The
determinant form of \texttt{\small H(\allowbreak{}L)\allowbreak{}} has
value \texttt{\small −1} in its canonical determinant trivialization;
hence
\texttt{\small w₁(\allowbreak{}M\textasciicircum{}⊥)\allowbreak{}=\allowbreak{}α+\allowbreak{}x}.
The Whitney formula now gives

\texttt{\small p*w₂(\allowbreak{}H(\allowbreak{}L)\allowbreak{})\allowbreak{}=\allowbreak{}x(\allowbreak{}α+\allowbreak{}x)\allowbreak{}}.

For any symmetric line bundle with class
\texttt{\small x∈H¹\_\allowbreak{}et(\allowbreak{}U,\allowbreak{}Z/\allowbreak{}2)\allowbreak{}},
its ordinary first Chern class satisfies

\texttt{\small c̄₁(\allowbreak{}M)\allowbreak{}=\allowbreak{}x²+\allowbreak{}αx}.

This is Lemma5.3 of EKV, pp181--182, whose cocycle calculation uses only
that2 is invertible. One can see the formula directly by writing
transition signs as
\texttt{\small (\allowbreak{}-\allowbreak{}1)\allowbreak{}\textasciicircum{}(\allowbreak{}a\_\allowbreak{}ij)\allowbreak{}},
choosing local square roots of \texttt{\small −1}, and lifting these
signs to fourth roots of unity. The Kummer coboundary has the two mod2
summands \texttt{\small a\_\allowbreak{}ij a\_\allowbreak{}jk} and
\texttt{\small α\_\allowbreak{}ij a\_\allowbreak{}jk}, respectively the
square and the cup product. Therefore

\texttt{\small p*w₂(\allowbreak{}H(\allowbreak{}L)\allowbreak{})\allowbreak{}=\allowbreak{}c̄₁(\allowbreak{}M)\allowbreak{}=\allowbreak{}p*c̄₁(\allowbreak{}L)\allowbreak{}}.

The pullback
\texttt{\small p*:\allowbreak{}H*\_\allowbreak{}et(\allowbreak{}X,\allowbreak{}Z/\allowbreak{}2)\allowbreak{}→\allowbreak{}H*\_\allowbreak{}et(\allowbreak{}U,\allowbreak{}Z/\allowbreak{}2)\allowbreak{}}
is injective by the orthogonal splitting principle, Zibrowius, \emph{The
gamma-filtration on the Witt ring of a scheme}, Proposition3.3, author
version arXiv:1411.3198v2, pp19--20. Its proposition explicitly assumes
a scheme over
\texttt{\small Z[\allowbreak{}1/\allowbreak{}2]\allowbreak{}}, which is
exactly the present scope. Injectivity gives the asserted equality on
\texttt{\small X}.

\subsection{A direct splitting morphism for an arbitrary vector
bundle}\label{a-direct-splitting-morphism-for-an-arbitrary-vector-bundle}

The same nonisotropic-line construction supplies the full higher-rank
formula without an additional affine-bundle invariance theorem. Let
\texttt{\small E} have rank \texttt{\small r>0} and form

\texttt{\small p:\allowbreak{}U=\allowbreak{}P\_\allowbreak{}nd(\allowbreak{}H(\allowbreak{}E)\allowbreak{})\allowbreak{}→\allowbreak{}X}.

Let \texttt{\small M⊂p*E⊕\allowbreak{}p*E∨} be the tautological
nondegenerate line. Its two projections give maps

\texttt{\small a:\allowbreak{}M→\allowbreak{}p*E},

\texttt{\small b:\allowbreak{}p*E→\allowbreak{}M∨}.

The composite
\texttt{\small β=\allowbreak{}b∘a:\allowbreak{}M→\allowbreak{}M∨} is an
isomorphism: in a local generator
\texttt{\small (\allowbreak{}v,\allowbreak{}f)\allowbreak{}} of
\texttt{\small M}, it is multiplication by the invertible value
\texttt{\small f(\allowbreak{}v)\allowbreak{}}. The nondegeneracy of the
restricted hyperbolic symmetric form says that
\texttt{\small 2f(\allowbreak{}v)\allowbreak{}} is a unit; here2 is
invertible. Therefore

\texttt{\small e=\allowbreak{}a∘β\textasciicircum{}(\allowbreak{}−1)\allowbreak{}∘b :\allowbreak{} p*E→\allowbreak{}p*E}

is an actual idempotent, since
\texttt{\small e²=\allowbreak{}aβ\textasciicircum{}(\allowbreak{}−1)\allowbreak{}(\allowbreak{}ba)\allowbreak{}β\textasciicircum{}(\allowbreak{}−1)\allowbreak{}b=\allowbreak{}e}.
It exhibits

\texttt{\small p*E=\allowbreak{}a(\allowbreak{}M)\allowbreak{} ⊕\allowbreak{} ker(\allowbreak{}b)\allowbreak{}}.

Both summands are vector bundles, of ranks1 and \texttt{\small r−1}. In
local coordinates the projector is
\texttt{\small z↦v f(\allowbreak{}z)\allowbreak{}/\allowbreak{}f(\allowbreak{}v)\allowbreak{}}.
Replacing \texttt{\small (\allowbreak{}v,\allowbreak{}f)\allowbreak{}}
by \texttt{\small t(\allowbreak{}v,\allowbreak{}f)\allowbreak{}} leaves
this projector unchanged, which also verifies the global gluing
directly. The line \texttt{\small a(\allowbreak{}M)\allowbreak{}} is
isomorphic to \texttt{\small M} through \texttt{\small a}.

Zibrowius Proposition3.3 and its first-step proof show that this
particular \texttt{\small p} induces an injection on mod2 étale
cohomology. Repeat the construction on
\texttt{\small ker(\allowbreak{}b)\allowbreak{}}, reducing its rank by
one at each step. After finitely many steps there is a composite
morphism \texttt{\small q:\allowbreak{}Y→\allowbreak{}X} with injective
pullback on mod2 étale cohomology and an actual direct-sum decomposition

\texttt{\small q*E=\allowbreak{}L₁⊕\allowbreak{}⋯⊕\allowbreak{}L\_\allowbreak{}r}.

The direct-sum projections and their dual inclusions give the orthogonal
isomorphism

\texttt{\small q*H(\allowbreak{}E)\allowbreak{} ≅ H(\allowbreak{}L₁)\allowbreak{} ⊥ ⋯ ⊥ H(\allowbreak{}L\_\allowbreak{}r)\allowbreak{}}.

By the already proved line-bundle result and the Whitney identities,

\texttt{\small q*w\_\allowbreak{}t(\allowbreak{}H(\allowbreak{}E)\allowbreak{})\allowbreak{}=\allowbreak{}∏\_\allowbreak{}\{i=\allowbreak{}1\}\textasciicircum{}r(\allowbreak{}1+\allowbreak{}αt+\allowbreak{}c̄₁(\allowbreak{}L\_\allowbreak{}i)\allowbreak{}t²)\allowbreak{}}.

Expanding by the set of indices that contribute their Chern term gives

\texttt{\small ∏\_\allowbreak{}\{i=\allowbreak{}1\}\textasciicircum{}r(\allowbreak{}1+\allowbreak{}αt+\allowbreak{}c̄₁(\allowbreak{}L\_\allowbreak{}i)\allowbreak{}t²)\allowbreak{}=\allowbreak{}Σ\_\allowbreak{}\{j=\allowbreak{}0\}\textasciicircum{}r(\allowbreak{}1+\allowbreak{}αt)\allowbreak{}\textasciicircum{}(\allowbreak{}r−j)\allowbreak{}c̄\_\allowbreak{}j(\allowbreak{}q*E)\allowbreak{}t\textasciicircum{}(\allowbreak{}2j)\allowbreak{}}.

Naturality and the injectivity of \texttt{\small q*} descend this
equality to \texttt{\small X}. Rank zero is the identity
\texttt{\small 1=\allowbreak{}1}, and a locally constant rank is treated
on its open-and-closed rank pieces. This proves the full hyperbolic
formula over every scheme over
\texttt{\small Z[\allowbreak{}1/\allowbreak{}2]\allowbreak{}} through
explicitly constructed bundle maps.

\subsection{What this resolves in the
archive}\label{what-this-resolves-in-the-archive}

Original Montpellier PDFpage4 was visually inspected and does contain
the \texttt{\small c₁ ?} question. The study's proposed answer for
\texttt{\small w₂(\allowbreak{}H(\allowbreak{}L)\allowbreak{})\allowbreak{}}
and its full hyperbolic formula are correct over every scheme over
\texttt{\small Z[\allowbreak{}1/\allowbreak{}2]\allowbreak{}}, as just
proved. The arbitrary-ringed-topos question about an augmented completed
lambda-ring target, recorded separately as A-Q080, still requires its
multiplication and operations to be specified; these
characteristic-class identities do not supply that missing target
definition.

The printed rank-two proof on EKVp183 includes a sentence asserting an
isomorphism on the unmodified Picard groups. That assertion is
unnecessary for the direct argument above. For example, the complement
for \texttt{\small H(\allowbreak{}L)\allowbreak{}} is the nonzero total
space of a square of \texttt{\small L}, and the pullback kills
\texttt{\small 2[\allowbreak{}L]\allowbreak{}}; consequently an integral
Picard-group isomorphism cannot be used in that generality. The
independently verified argument uses the stated injective mod2
cohomology pullback and the actual line isomorphism instead.

\section{A-023. A characteristic-zero local ring where the rho classes
do not
generate}\label{a-023.-a-characteristic-zero-local-ring-where-the-rho-classes-do-not-generate}

Take the explicit commutative ring

\texttt{\small A=\allowbreak{}Z\_\allowbreak{}(\allowbreak{}2)\allowbreak{}[\allowbreak{}ε]\allowbreak{}/\allowbreak{}(\allowbreak{}2ε,\allowbreak{} ε²)\allowbreak{}}.

As an abelian group it is
\texttt{\small Z\_\allowbreak{}(\allowbreak{}2)\allowbreak{}⊕\allowbreak{}F₂ε},
with product

\texttt{\small (\allowbreak{}a+\allowbreak{}bε)\allowbreak{}(\allowbreak{}c+\allowbreak{}dε)\allowbreak{}=\allowbreak{}ac+\allowbreak{}(\allowbreak{}ād+\allowbreak{}c̄b)\allowbreak{}ε}.

This proves that \texttt{\small ε≠0} and that \texttt{\small A} has
characteristic zero. An element \texttt{\small a+\allowbreak{}bε} is a
unit precisely when \texttt{\small a} is a unit of
\texttt{\small Z\_\allowbreak{}(\allowbreak{}2)\allowbreak{}}: its
inverse is
\texttt{\small a\textasciicircum{}(\allowbreak{}−1)\allowbreak{}−a\textasciicircum{}(\allowbreak{}−2)\allowbreak{}bε}.
Hence \texttt{\small A} is local with maximal ideal
\texttt{\small (\allowbreak{}2,\allowbreak{}ε)\allowbreak{}} and residue
field \texttt{\small F₂}.

The coefficient ring also has
\texttt{\small K₀(\allowbreak{}A)\allowbreak{}≅Z}, with the isomorphism
given by rank. Indeed, for a finite projective module \texttt{\small P},
lift a basis of
\texttt{\small P/\allowbreak{}(\allowbreak{}2,\allowbreak{}ε)\allowbreak{}P}
to obtain a surjection
\texttt{\small A\textasciicircum{}r→\allowbreak{}P} by Nakayama's lemma.
It splits because \texttt{\small P} is projective. Its complementary
finite projective module has zero residue quotient and is therefore zero
by the same lemma. Thus every finite projective \texttt{\small A}-module
is free, and rank identifies its Grothendieck group with
\texttt{\small Z}. The example therefore also addresses the
originalpage29 parenthetical question about coefficient rings with
\texttt{\small K(\allowbreak{}k)\allowbreak{}≅Z}.

Let \texttt{\small G=\allowbreak{}S₂=\allowbreak{}C₂}, with generator
\texttt{\small g}. The free rank-one \texttt{\small A}-module
\texttt{\small L=\allowbreak{}Ae} with action

\texttt{\small g(\allowbreak{}e)\allowbreak{}=\allowbreak{}(\allowbreak{}1+\allowbreak{}ε)\allowbreak{}e}

is a representation, since
\texttt{\small (\allowbreak{}1+\allowbreak{}ε)\allowbreak{}²=\allowbreak{}1}.
Its determinant character sends \texttt{\small g} to
\texttt{\small 1+\allowbreak{}ε}.

For every \texttt{\small A[\allowbreak{}G]\allowbreak{}}-module
\texttt{\small M} whose underlying \texttt{\small A}-module is finite
free, the determinant of each group element defines a character
\texttt{\small det(\allowbreak{}M)\allowbreak{}:\allowbreak{}G→\allowbreak{}A×}.
This gives a homomorphism

\texttt{\small det:\allowbreak{}K₀(\allowbreak{}L\_\allowbreak{}A(\allowbreak{}G)\allowbreak{})\allowbreak{}→\allowbreak{}Hom(\allowbreak{}G,\allowbreak{}A×)\allowbreak{}}

from the exact Grothendieck group (written additively) to the character
group (written multiplicatively). Indeed, every admissible exact
sequence splits as a sequence of underlying \texttt{\small A}-modules.
In an \texttt{\small A}-basis adapted to its submodule, a group element
acts by a block upper triangular matrix. Its determinant is the product
of the determinants on submodule and quotient. Thus the determinant
respects precisely the exact relations defining this \texttt{\small K₀}.

In degree two, the monomials in the proposed generators are

\texttt{\small ρ₂=\allowbreak{}[\allowbreak{}A\_\allowbreak{}trivial]\allowbreak{}},

\texttt{\small ρ₁²=\allowbreak{}[\allowbreak{}Ind\_\allowbreak{}\{S₁×S₁\}\textasciicircum{}\{S₂\}A]\allowbreak{}=\allowbreak{}[\allowbreak{}A[\allowbreak{}S₂]\allowbreak{}]\allowbreak{}}.

The corresponding determinant characters take \texttt{\small g} to
\texttt{\small 1} and \texttt{\small −1}, respectively: on the standard
basis of the regular module, \texttt{\small g} interchanges two basis
vectors. Consequently the determinant of every integer linear
combination \texttt{\small aρ₂+\allowbreak{}bρ₁²} takes
\texttt{\small g} to
\texttt{\small (\allowbreak{}-\allowbreak{}1)\allowbreak{}\textasciicircum{}b}.

The character \texttt{\small 1+\allowbreak{}ε} equals neither
\texttt{\small 1} nor \texttt{\small −1}. The first inequality follows
from \texttt{\small ε≠0}. Reduction modulo \texttt{\small ε} maps the
second proposed equality to \texttt{\small 1=\allowbreak{}−1} in the
characteristic-zero ring
\texttt{\small Z\_\allowbreak{}(\allowbreak{}2)\allowbreak{}}, which is
false. Therefore

\texttt{\small [\allowbreak{}L]\allowbreak{} ∉ Zρ₂+\allowbreak{}Zρ₁²}.

The graded map
\texttt{\small Z[\allowbreak{}ρ₁,\allowbreak{}ρ₂,\allowbreak{}…]\allowbreak{}→\allowbreak{}⊕\allowbreak{}\_\allowbreak{}n K₀(\allowbreak{}L\_\allowbreak{}A(\allowbreak{}S\_\allowbreak{}n)\allowbreak{})\allowbreak{}}
fails surjectivity already in degree two. This is an explicit
counterexample even among characteristic-zero local rings. It is
compatible with A-021: the ring \texttt{\small A} has the nonzero
torsion element \texttt{\small ε} and is not a discrete valuation ring.
The exact determinant map above is the witness for the failure; no
distinction between two unnamed representation categories is being used.

The supplied study marks the classification of coefficient rings as open
and its later annotation points to a characteristic-zero local
counterexample in another record. That other record was not read here.
The ring, representation, determinant map and non-generation proof above
were constructed and checked directly in this run.

\section{The five-point Cox ring and the Segre cubic, with the exact
group
action}\label{the-five-point-cox-ring-and-the-segre-cubic-with-the-exact-group-action}

\subsection{A-028. Source identification and actual algebra
map}\label{a-028.-source-identification-and-actual-algebra-map}

Let \(E=\mathfrak h_5\), let
\(Z=\{[\ell_i]\}_{i=1}^5\subset\mathbf P(E)\), and let
\(X=\operatorname{Bl}_Z\mathbf P(E)\). Write \(H\) for the pullback of
the hyperplane class and \(E_i\) for its five exceptional divisors.
Choose the canonical exceptional sections \(e_i\), with their literal
permutation action. The following explicit map identifies the Cox ring
with the Plücker coordinate ring of \(\operatorname{Gr}(2,6)\): \[
\boxed{
p_{i6}\longmapsto e_i,\qquad
p_{ij}\longmapsto
\frac{q_i-q_j}{\prod_{k\notin\{i,j\}}e_k}
\quad(1\le i<j\le5).
}
\tag{Cox.1}
\] Here \(\sum_iq_i=0\), and \(p_{ji}=-p_{ij}\). Its Picard grading is
\[
\deg p_{i6}=E_i,\qquad
\deg p_{ij}=H-\sum_{k\notin\{i,j\}}E_k.
\tag{Cox.2}
\] The \(S_5\) action permutes the first five Plücker indices and fixes
the sixth. It is equivariant under (Cox.1), including the signs arising
when an ordered pair is reversed.

The primary general Cox-ring locator is A.-M. Castravet and J. Tevelev,
\emph{Hilbert's 14th problem and Cox rings}, \textbf{Compositio Math.
142} (2006), 1479--1498, \textbf{Remark 3.9, p.~1487}. I opened and read
the \href{https://people.math.umass.edu/~tevelev/14problem.pdf}{author's
published PDF}. Its \(s=1,n=3\) instance identifies this Cox ring with
the Plücker ring of \(\operatorname{Gr}(2,6)\). The exact map and action
above are verified below rather than inferred from an abstract ring
isomorphism.

For the Segre map and the representation convention, I opened and read
I. Dolgachev, B. Farb and E. Looijenga, \emph{Geometry of the
Wiman--Edge pencil, I: algebro-geometric aspects},
\href{https://www.math.uchicago.edu/~farb/papers/wiman.pdf}{author PDF},
\textbf{§2.3, Table 1, p.~10}, and \textbf{Remark 2.4, p.~13}. Table 1
gives the five-dimensional representation with transposition trace 1;
Remark 2.4 describes the equivariant quadratic map through the five
frame points and identifies the ambient \(S_6\) representation as
\([(3,3)]\). The printed remark contains a dimension typo
\(\mathbf P^4=\mathbf P(V)\): the \(V\) fixed in §2.3 has dimension four
and its projectivization is \(\mathbf P^3\). The quadratic map and
target dimension in the surrounding discussion specify the intended
spaces.

\subsection{1. Proof of the Cox map through one additive-group
action}\label{proof-of-the-cox-map-through-one-additive-group-action}

Take the polynomial ring \(P=\mathbf Q[x_1,\ldots,x_5,y_1,\ldots,y_5]\)
and the action of \(\mathbf G_a\) given by \(y_i\mapsto y_i+t x_i\),
with \(x_i\) fixed. The invariant polynomials include \[
x_i,\qquad p_{ij}=x_i y_j-x_jy_i.
\] The invariant algebra is the Plücker ring of
\(\operatorname{Gr}(2,6)\), with \(p_{i6}=x_i\). Here is a direct
comparison map. Regard \((x_i,y_i)^T\) as the first five columns of a
\(2\times6\) matrix and set the sixth column equal to \((0,1)^T\). Its
stabilizer in \(SL_2\) is exactly the displayed additive group. An
invariant function of the first five columns extends uniquely to an
\(SL_2\)-invariant regular function on the open set where the sixth
column is nonzero: choose locally an \(SL_2\) element sending that
column to \((0,1)^T\), and use additive invariance to make the value
independent of the choice. The two standard coordinate charts of the
nonzero sixth column give regular choices. The complement has
codimension two in affine matrix space, so the function extends across
it by the intersection property of polynomial-ring localizations.
Restriction and extension are inverse ring maps.

Finally, the \(SL_2\) invariant algebra of six column vectors is the
image algebra of their \(2\times2\) minors. This can be seen directly
from the Cauchy decomposition: its degree \(2d\) part is
\(S_{(d,d)}\mathbf Q^6\), and products of the minors contain a nonzero
highest vector in every such summand. Their span is \(GL_6\)-stable, so
it contains the whole summand. This is precisely the Plücker coordinate
ring. The construction commutes with permutation of the first five
columns.

To identify this invariant algebra with the Cox ring using the specified
grading, put \[
u=\prod_i x_i,\qquad r_i=y_i/x_i,\qquad
q_i=u\left(\frac15\sum_jr_j-r_i\right).
\tag{Cox.3}
\] These \(q_i\) are polynomial invariants, have sum zero, and transform
by the standard representation. On \(u\ne0\), the differences of the
\(r_i\) are coordinates for the additive quotient. Thus the localized
invariant ring is exactly
\(\mathbf Q[x_i^{\pm1},q_1,\ldots,q_5]/(\sum q_i)\). Sending \(x_i\) to
\(e_i\) gives \[
x_i y_j-x_j y_i
=\frac{q_i-q_j}{\prod_{k\notin\{i,j\}}x_k},
\] which proves the formula for each generator in (Cox.1).

Assign \(\deg x_i=E_i\) and \(\deg y_i=H-\sum_{j\ne i}E_j\). Then
\(q_i\) has degree \(H\), and every homogeneous element of the localized
invariant ring is a polynomial \(f(q)\) of some degree \(c\), multiplied
by \(\prod_i x_i^{a_i}\). Its degree is \(cH+\sum_i a_iE_i\).

It remains to check which such elements are polynomial before
localization. Along \(x_i=0\), the vector \(q\) in (Cox.3) becomes the
nonzero scalar \(-y_i\prod_{j\ne i}x_j\) times \(\ell_i\) on the generic
locus. The normalized transverse coordinates have leading terms
\(x_i(r_j-r_l)/y_i\), with \(j,l\ne i\); their three independent
differences vary freely. The same lowest-term calculation as G-015
therefore proves \[
\operatorname{ord}_{x_i}f(q(X,Y))
=\operatorname{mult}_{[\ell_i]}f.
\] The scalar trivializing the homogeneous degree-\(c\) factor is a unit
there. Since the \(x_i\) are distinct primes, the Laurent expression is
polynomial exactly when \(\operatorname{mult}_{[\ell_i]}f+a_i\ge0\) for
every \(i\).

These are exactly the section conditions for
\(\mathcal O_X(cH+\sum_i a_iE_i)\). They follow directly on the affine
blow-up charts from divisibility by the required order of the point
ideal. Positive exceptional coefficients permit their canonical sections
but no additional poles downstairs: the removed points have codimension
three, and a pole divisor cannot be supported on them. The criterion
therefore identifies each Picard-graded piece, and multiplication is the
same multiplication of Laurent polynomials. This proves the full ring
isomorphism (Cox.1) with its exact action.

\subsection{2. The quadratic subalgebra as an actual Cox
component}\label{the-quadratic-subalgebra-as-an-actual-cox-component}

Put \(D=2H-\sum_iE_i\). Under (Cox.1), the space
\(H^0(X,\mathcal O_X(D))\) is the space \(U\) of quadrics through the
five points. Its Plücker multidegree is \((1,1,1,1,1,1)\), and its
Plücker degree is three.

To verify the degree conversion explicitly, let \(b_{ij}\) count a
factor \(p_{ij}\) with \(i,j\le5\), and let \(a_i\) count \(p_{i6}\). If
its Cox degree is \(cH-k\sum E_i\), then \[
\sum_{i<j}b_{ij}=c,\qquad
a_i+\sum_{j\ne i}b_{ij}=c-k.
\] Hence its six column weights are \[
(c-k,c-k,c-k,c-k,c-k,3c-5k),
\] and its total Plücker degree is \(4c-5k\). Substituting \(c=2,k=1\)
gives the asserted degree three and six equal weights.

The multilinear \(SL_2\) invariants of six columns are the Specht module
\([(3,3)]\). This follows from the tensor Schur--Weyl decomposition,
with the \(SL_2\)-invariant determinant power \(S_{(3,3)}\mathbf Q^2\).
Restriction from the literal \(S_6\) column permutation action to the
stabilizer of column six removes the unique removable box from
\((3,3)\), giving \([(3,2)]\). Thus the extension of the quadratic
\(S_5\) action to \(S_6\) is fixed by this actual column model; no
choice of a projective lift determines it implicitly.

The same degree conversion sends every \(H^0(X,kD)\) to column weights
\((k^6)\). Permuting all six columns preserves this diagonal component
and therefore gives an \(S_6\) action on the whole \(D\)-section ring.
It preserves the original \(S_5\) action on the first five indices.
Permutations involving the sixth index act on this diagonal ring
regardless of their action on the other Picard degrees.

By A-027 the algebra generated by \(U\) has the exact presentation \[
\mathbf Q[A,B,C,D,E]/
\bigl(AD(E-B-C)-BC(E-A-D)\bigr).
\] By A-029 the ordinary and saturated ideal powers agree in degree
\(2k\) for every \(k\). Consequently multiplication of the quadrics maps
this algebra isomorphically onto \(\bigoplus_{k\ge0}H^0(X,kD)\). This
identifies the full section ring along the displayed Cox degree with the
cubic algebra, including every graded piece and its action.

The cubic has sign character under \(S_5\), as the direct transposition
calculation in A-027 shows. T-038 separately constructs its six linear
forms and the actual exceptional six-element permutation action; that
explicit intertwiner compares its usual six-coordinate Segre equation
with the quadratic coordinates above. The present Cox calculation fixes
the geometric \(S_6\) extension as \([(3,3)]\) in degree one.

\subsection{3. The exact maps to ordinary and sheaf
powers}\label{the-exact-maps-to-ordinary-and-sheaf-powers}

There is an injective bigraded map \[
\mathcal R(I)=\bigoplus_{k\ge0}I^k t^k
\longrightarrow\operatorname{Cox}(X),\qquad
f(q)t^k\longmapsto f(q)\left(\prod_i e_i\right)^{-k}.
\] Its degree-\((c,k)\) image is \((I^k)_c\) in the Cox component
\(H^0(X,cH-k\sum E_i)\), which is the corresponding saturated fat-point
section space. The component map is the usual inclusion of ordinary
ideal powers into their sheaf sections. A-029 proves it is onto for
\(c\ge2k\). Thus the quadratic cubic algebra, the ordinary Rees algebra,
and the Cox section spaces are related here by specified inclusions and
the stated equalities of graded components.

The full flag problem G-015 additionally tensors by its specified flag
line bundle. The proofs here determine the point-ideal algebra and its
action and provide the exact cohomology inputs recorded in A-029. They
make no assertion that an alternating resolution alone has already
evaluated every \(S_5\)-covariant flag-section multiplicity.

\section{The five-point ideal: exact Rees map, linear powers, and the
quotient-bundle
sections}\label{the-five-point-ideal-exact-rees-map-linear-powers-and-the-quotient-bundle-sections}

\subsection{A-029. Completed result}\label{a-029.-completed-result}

Let \(S=\mathbf Q[x,y,z,w]\), let \(Z\subset\mathbf P^3\) be the four
coordinate points and \([1:1:1:1]\), and put \[
I=(xz-xy,xw-xy,yz-xy,yw-xy,zw-xy).
\] The literal \(S_5\) action is the action permuting the five canonical
projective-frame points in \(\mathbf P(\mathfrak h_5)\), transported by
the specified basis. Write \(U=I_2=[(3,2)]\). The following statements
have been proved with the hypotheses and degree conventions checked:

\begin{enumerate}
\def\labelenumi{\arabic{enumi}.}
\tightlist
\item
  The Rees algebra is the quotient by the actual five linear syzygies
  and the one cubic \(F\) below.
\item
  The kernel of \(\operatorname{Sym}^k_S(I(2))\to I^k(2k)\) is the
  degree-zero \(S_5\)-module
  \(\operatorname{sgn}\otimes\operatorname{Sym}^{k-3}U\), annihilated by
  \((x,y,z,w)\); a negative symmetric power means zero.
\item
  For every \(k\ge2\), \(I^k\) has a \(2k\)-linear free resolution. In
  particular its saturation agrees with \(I^k\) in every degree
  \(c\ge2k\), and \[
  H^1(\mathbf P^3,\mathcal I_Z^k(2k-1))=0.
  \]
\item
  For the tautological rank-three quotient \(\mathcal Q\) on
  \(\mathbf P(\mathfrak h_5)\), put
  \(M_k=H^0(\mathcal Q\otimes\mathcal I_Z^k(2k))\). Then \[
  \dim M_k=\frac{(k+1)(k+2)(3k+4)}2,
  \qquad
  \sum_{k\ge0}\dim M_k t^k=\frac{4+5t}{(1-t)^4}.
  \] These are actual statements for this ideal; no unverified
  general-position or acyclicity hypothesis remains in their use below.
\end{enumerate}

\subsection{Primary source and exact
locators}\label{primary-source-and-exact-locators}

The primary source is A. R. Kustin, C. Polini, B. Ulrich, \emph{The
equations defining blowup algebras of height three Gorenstein ideals},
\textbf{Algebra \& Number Theory 11} (2017), 1489--1525, DOI
\href{https://doi.org/10.2140/ant.2017.11.1489}{10.2140/ant.2017.11.1489}.
The \href{https://msp.org/ant/2017/11-7/ant-v11-n7-p01-p.pdf}{publisher
PDF} was opened and read; its local copy is
\texttt{\small primary\_\allowbreak{}m5/\allowbreak{}kustin\_\allowbreak{}polini\_\allowbreak{}ulrich\_\allowbreak{}2017.pdf}.

The exact passages used are formula \textbf{(6.1.4), p.~1509};
\textbf{Lemma 6.3 and Corollary 6.4, p.~1510}, with the corollary's
proof on p.~1511; and \textbf{Theorem 9.1, pp.~1520--1521}. Proposition
7.5, pp.~1514--1515, independently states the symmetric-power regularity
bounds agreeing with the calculation below. The 2016 author PDF has
different pagination: formula (6.1.4) p.~18, Lemma 6.3/Corollary 6.4
p.~19, Theorem 9.1 p.~28. The journal version supplies the locators in
this report.

I read the statements and their surrounding proofs, including the local
generator condition and the exceptional degree in (6.1.4). The use of
those established complexes is explicitly attributed; their 1992
construction is not presented as a newly checked construction here.

\subsection{1. The ideal and all local
hypotheses}\label{the-ideal-and-all-local-hypotheses}

The complete first resolution is the one constructed in T-038. Its
exactness can be verified from the displayed \(M,N\) without any extra
geometric hypothesis: \(f^TM=0\), \(MN=0\), a \(4\times4\) minor of
\(M\) is nonzero, and \(\gcd(N_1,\ldots,N_5)=1\). Over the fraction
field, the kernel of \(M\) is spanned by \(N\). A rational scalar
\(a/b\) with coprime numerator and denominator whose product with every
\(N_i\) is polynomial has \(b\) dividing all \(N_i\), so \(b\) is a
unit. Thus the polynomial kernel is exactly \(SN\). The Hilbert series
from G-016 then proves exactness at the preceding term, giving \[
0\longrightarrow S(-5)\longrightarrow S(-3)^5
\longrightarrow S(-2)^5\longrightarrow I\longrightarrow0.
\] In particular the presentation entries are linear and the ideal has
height three.

For direct compatibility with the primary source, the same matrix can be
made alternating. If \(f=(A(x),B(x),C(x),D(x),E(x))^T\) is the vector of
the five quadrics, then the terminal column is \(N=Pf\), where \[
P=\begin{pmatrix}
0&0&0&-1&0\\
0&0&1&0&-1\\
0&1&0&0&0\\
0&0&-1&1&0\\
1&-1&0&0&0
\end{pmatrix},\qquad\det P=-1.
\] The actual matrix \(\varphi=MP\) is \[
\begin{pmatrix}
0&y&w&0&-y\\
-y&0&0&-z&y\\
-w&0&0&-x&x\\
0&z&x&0&-x\\
y&-y&-x&x&0
\end{pmatrix}.
\] Its signed maximal Pfaffians are \(f\): their vector spans the same
one-dimensional kernel and its first entry is \(xz-xy\). The resolution
is self-dual, with terminal map \(f\), so the canonical module of
\(S/I\) is a grading shift of \(S/I\). Together with its dimension-one
Cohen--Macaulay property this verifies that \(I\) is a perfect
Gorenstein ideal. One can check the Cohen--Macaulay property directly: a
linear form nonzero at all five points is a nonzerodivisor in the
product embedding of the five line rings proved in G-016.

Away from the common vertex, the five affine lines are disjoint and
smooth. The ideal localized at any point of that punctured cone is
generated by three independent normal linear parameters. Thus the
source's \(G_4\) condition is satisfied. Its weaker condition
\(\mu(I_{\mathfrak p})\le\dim S_{\mathfrak p}+1\) also holds at every
prime: the only remaining vertex has five generators in a
four-dimensional regular local ring. Hence Lemma 6.3 applies with
\(s=4\) at every prime, not merely on the punctured spectrum. The
parameters in the source are \(d=4,n=5,\delta=2,D=1\).

\subsection{2. The exact Rees presentation and torsion
module}\label{the-exact-rees-presentation-and-torsion-module}

Let \(A,B,C,D,E\) now denote five independent variables mapping to the
five quadratic generators. The exact subalgebra kernel checked in A-027
is generated by \[
F=AD(E-B-C)-BC(E-A-D).
\] It has sign character. Put \(T=(A,B,C,D,E)\) and
\(L=(\text{entries of }T\varphi)\subset S[T]\). Then \(S[T]/L\) is the
symmetric algebra of the displayed presentation of \(I(2)\). The map to
the Rees ring sends each \(T_i\) to the corresponding \(f_i t\); it is
equivariant with the literal action on \(U\).

Corollary 6.4 gives fiber type for this ideal. Since its special-fiber
equation has just been proved to be exactly \(F\), it gives the concrete
presentation \[
\boxed{\mathcal R(I)=S[T]/(L,F).}
\] The same corollary says the symmetric-algebra torsion is annihilated
by \(\mathfrak m^{(n-d+1)/2}=\mathfrak m\). The linear syzygy ideal
\(L\) has positive \(x,y,z,w\)-degree and therefore has zero
intersection with \(\mathbf Q[T]\). It follows that the \(k\)-th torsion
piece is precisely \[
F\operatorname{Sym}^{k-3}U
\simeq\operatorname{sgn}\otimes\operatorname{Sym}^{k-3}U,
\] in polynomial degree zero after the shift by \(2k\). The
multiplication map giving this isomorphism is injective because
\(\mathbf Q[T]\) is a domain.

There is also an independent exact certificate for the annihilation
assertion. Write \(T\varphi=(x,y,z,w)B(T)\), where \[
B(T)=\begin{pmatrix}
0&0&D-E&E-C&C-D\\
E-B&A-E&0&0&B-A\\
0&D&0&-B&0\\
-C&0&A&0&0
\end{pmatrix}.
\] For the polynomial matrix \[
Q(T)=\begin{pmatrix}
AB&AD&-A(A+C-E)&-B(A+D-E)\\
AB&BC&-A(B+C-E)&-B(B+D-E)\\
BC&CD&-C(A+C-E)&-D(B+C-E)\\
AD&CD&-C(A+D-E)&-D(B+D-E)\\
AB&CD&-AC&-BD
\end{pmatrix},
\] direct multiplication gives \(B(T)Q(T)=F(T)I_4\). Consequently
\((x,y,z,w)F=(T\varphi)Q\), which explicitly places all four products in
\(L\). This multiplication was executed exactly; the receipt is
\texttt{\small FIVE\_\allowbreak{}POINT\_\allowbreak{}TORSION\_\allowbreak{}CERTIFICATE.txt}.

\subsection{\texorpdfstring{3. Linearity of every power from \(k=2\)
onward}{3. Linearity of every power from k=2 onward}}\label{linearity-of-every-power-from-k2-onward}

By the verified global hypothesis of Lemma 6.3, the complex
\(D^k_\bullet(\varphi)\) is a free resolution of
\(\operatorname{Sym}^k_S(I(2))\). Inspect formula (6.1.4) with
\(n=5,D=1\):

\begin{itemize}
\tightlist
\item
  \(k=2\): all nonzero terms in homological degrees \(0,1,2\) have shift
  equal to the homological degree.
\item
  \(k=3\): degrees \(0,1,2,3\) have that same shift. The exceptional
  \(Q_3\) term in degree four has shift \(3+(5-4+1)/2=4\), so it is
  linear as well.
\item
  \(k\ge4\): every nonzero term lies in homological degree \(r\le4\) and
  has shift \(r\).
\end{itemize}

Thus \(\operatorname{Sym}^k_S(I(2))\) has a linear resolution for all
\(k\ge2\). The excluded case \(k=1\) has the known shift three in
homological degree two, agreeing with its original resolution.

The torsion kernel in §2 is a vector space in degree zero killed by
\(\mathfrak m\), so its Koszul resolution is linear. Apply the graded
Tor exact sequence to \[
0\to\operatorname{sgn}\otimes\operatorname{Sym}^{k-3}U
\to\operatorname{Sym}^k_S(I(2))\to I^k(2k)\to0.
\] In homological degree \(i\), the possible degrees for the last Tor
group are at most \(i\). A module generated in degree zero has no
minimal \(i\)-th syzygies in degrees below \(i\). Therefore its Tor is
concentrated in degree \(i\): \(I^k(2k)\) has a linear resolution. This
proves \(\operatorname{reg}(I^k)=2k\), with the equality also forced by
its nonzero minimal generators in degree \(2k\).

Finite group equivariance is retained: the multiplication, its kernel,
and the Tor sequence are equivariant, and degreewise averaging supplies
equivariant choices of minimal generators whenever a free resolution is
chosen. This statement supplies actual resolutions and exact modules; it
does not claim an unevaluated representation-theoretic splitting to be a
positive coefficient formula.

\subsection{4. Saturation and the cohomology required in
G-017}\label{saturation-and-the-cohomology-required-in-g-017}

The saturation of \(I^k\) is the homogeneous ideal of the five
order-\(k\) fat points, since the ordinary power restricts to that ideal
away from the affine vertex. Its quotient by \(I^k\) is
\(H^0_{\mathfrak m}(S/I^k)\). The linear resolution bounds this module's
highest degree by \(2k-1\): in a free resolution of \(S/I^k\), the term
in homological degree four has shift at most \(2k+3\), and the top Čech
cohomology of \(S(-a)\) has highest degree \(a-4\). Thus \[
(I^k)_c=H^0(\mathbf P^3,\mathcal I_Z^k(c))
\quad(c\ge2k,k\ge2).
\] For \(k=1\), equality in all degrees follows from G-016's saturated
ideal calculation. In particular the degree-\(2k\) section space is the
image of \(\operatorname{Sym}^kU\), whose only relations are the
degree-\(k\) multiples of \(F\). Write \[
C_k=\binom{k+4}{4}-\binom{k+1}{4}
=\frac{(k+1)(k^2+2k+2)}2.
\] Here the second binomial is zero for \(k<3\).

The same resolution gives \(H^1(\mathcal I_Z^k(2k-1))=0\) for \(k\ge2\).
Equivalently, the highest degree of \(H^2_{\mathfrak m}(I^k)\) is at
most \(2k-2\). The fat-point scheme has length \(5\binom{k+2}{3}\),
because a point in a smooth threefold has \(\binom{k+2}{3}\) monomials
of degree below \(k\) in its three normal variables. Hence \[
h^0(\mathcal I_Z^k(2k-1))
=\binom{2k+2}{3}-5\binom{k+2}{3}
=\frac{k(k+1)(k-2)}2
\quad(k\ge2).
\] These are sheaf sections below the generator degree of the ordinary
power; the exact inclusion and saturation calculation above specifies
their relationship.

Tensor the tautological quotient sequence \[
0\to\mathcal O(-1)\to\mathfrak h_5\otimes\mathcal O
\to\mathcal Q\to0
\] with \(\mathcal I_Z^k(2k)\). It remains exact because the quotient
bundle is locally free. The proved first-cohomology vanishing makes its
global sections right exact for \(k\ge2\). Therefore \[
\dim M_k=4C_k-\frac{k(k+1)(k-2)}2
=\frac{(k+1)(k+2)(3k+4)}2.
\] For \(k=1\), evaluation in degree one is
\(\mathfrak h_5\hookrightarrow\operatorname{Perm}_5\), with trivial
one-dimensional cokernel. Thus \(H^1(\mathcal I_Z(1))=\mathbf Q\), while
\(H^0(\mathcal I_Z(1))=0\), and the same exact sequence gives
\(\dim M_1=4\cdot5+1=21\). For \(k=0\),
\(H^0(\mathcal Q)=\mathfrak h_5\) has dimension four. These agree with
the same polynomial formula.

Finally \[
\frac{(k+1)(k+2)(3k+4)}2
=4\binom{k+3}{3}+5\binom{k+2}{3},
\] so summing gives the stated series \((4+5t)/(1-t)^4\). This checks
precisely the cohomology and dimension step requested for G-017. The
later presentation of \(M\) by the root's matrix is a separate
proof-review event.

\part{Algebraic-geometric restatements and proof dictionaries}

\section{Separate algebraic-geometric
restatements}\label{separate-algebraic-geometric-restatements}

\subsection{The common exact
transport}\label{the-common-exact-transport}

Let \(E\) be a locally free sheaf of rank \(r\). Its frame torsor
\(\operatorname{Isom}(\mathcal O^r,E)\) transports every equivariant
polynomial map between the specified \(GL_r\)-modules to a natural map
between associated locally free sheaves. The representation ring is the
exact Grothendieck ring of the category of finite locally free
representations; the class of an exterior power is the corresponding
lambda operation. These are the constructions of SGA 3, Exposé I,
Section 4.5 and Propositions 4.7.2--4.7.3, and SGA 6, Exposé VI, Theorem
3.3. The conjugation \(S^{\mathrm{int}}_\lambda=L_{\lambda'}\) used in
A-012 fixes the integral Schur convention explicitly.

The evaluation homomorphism from the free special lambda ring sends
\(e_i\) to \([\Lambda^iE]\). In rank \(r\), it factors through
\(e_{r+1},e_{r+2},\ldots=0\); determinant inversion gives the Laurent
representation ring of \(GL_r\). The precise universal locators are SGA
6 V Proposition 4.2, Definition 4.4, Proposition 4.8, and Corollary 4.9.
Composition of the resulting operations is the Schur-functor composition
proved by the actual functorial maps in the preceding parts.

Characteristic-zero averaging uses the idempotent
\(|H|^{-1}\sum_{h\in H}h\). Its integral model is over
\(\mathbb Z[1/|H|]\), a scheme over \(\mathbb Z\). Each additional
rational presentation has finitely many coefficients and a finite
Gröbner basis; invert the product of their denominators and leading
coefficients. The resulting monic relations, module maps, and flat
Gröbner family define the exact same generic-fiber construction over
that localization. This construction specifies its base rather than
silently extending an averaging argument across primes dividing its
denominator.

\subsection{Dictionary for every theta and positive-formula
finding}\label{dictionary-for-every-theta-and-positive-formula-finding}

\begin{longtable}[]{@{}
  >{\raggedright\arraybackslash}p{(\columnwidth - 4\tabcolsep) * \real{0.3333}}
  >{\raggedright\arraybackslash}p{(\columnwidth - 4\tabcolsep) * \real{0.3333}}
  >{\raggedright\arraybackslash}p{(\columnwidth - 4\tabcolsep) * \real{0.3333}}@{}}
\toprule\noalign{}
\begin{minipage}[b]{\linewidth}\raggedright
Finding
\end{minipage} & \begin{minipage}[b]{\linewidth}\raggedright
Algebraic-geometric restatement and the map carrying the proof
\end{minipage} & \begin{minipage}[b]{\linewidth}\raggedright
Exact framework
\end{minipage} \\
\midrule\noalign{}
\endhead
\bottomrule\noalign{}
\endlastfoot
T-001 & The complexification of the compact pair gives the polynomial
algebra on \(E\otimes F\). Its two commuting group-scheme actions yield
the Cauchy summands; the stated determinant square root records the
oscillator pullback before its explicit removal. & SGA 3 I 4.5, 4.7.2;
SGA 6 VI 3.3. \\
T-002 & Extract the \((n^m)\)-graded summand of the diagonal-torus
comodule and apply the constant \(S_m\) idempotent. This associated
summand is exactly the multiplicity module of \(S_\lambda E\). & SGA 3 I
4.7.3; SGA 6 V 4.2. \\
T-003 & Pull back the \(GL(S_\nu E)\) representation along its specified
Schur-functor representation morphism; the wreath-group Hom is the
equivalent tensor idempotent summand. & SGA 3 I 4.7.2; SGA 6 VI 3.3. \\
T-004 & Interlacing patterns index the torus-graded Schur sheaf before
the finite quotient; its Hilbert function is the stated weight
multiplicity. & SGA 3 I 4.7.3. \\
T-005 & Tensoring by \(\det^t F\) shifts the torus grading and restricts
to \(\operatorname{sgn}^t\) on the permutation subgroup. This
line-bundle isomorphism transports the precise projector. & SGA 3 I 4.5,
4.7.2; SGA 6 V 4.8--4.9. \\
T-006 & The rank-two weight line is \(\det^j F\) times a balanced
symmetric monomial; its swap scalar computes the invariant direct
summand. & SGA 3 I 4.7.3; SGA 6 VI 3.3. \\
T-007 & The even-row decomposition becomes the indicated equality of
associated Schur-sheaf classes after the frame-torsor pullback. & SGA 6
VI 3.3. \\
T-008 & Products of highest-weight sections in the stated
integral-domain coordinate algebras give nonzero sections of the sum
weight. This is the actual multiplication map proving support addition.
& EGA I 9.5.1--9.5.3; SGA 3 I 4.7.2. \\
T-009 & The specified stretching Hilbert function has the recorded
period-six values; Ehrhart reciprocity contradicts the proposed
homogeneous-polytope Hilbert function. It leaves the proved cellwise
presentations intact. & EGA I 9.5.3; the exact Hilbert grading is
specified in T-009. \\
T-010 & The stable equality is an identity of Schur coefficients in the
completed lambda ring; its transpose is the characteristic of the
explicitly graded symmetric/exterior algebra \(\mathcal M_k\). & SGA 6 V
4.2, VI 3.3. \\
T-011 & The integer arrays index Littlewood--Richardson multiplicity
bases. Applying the associated-bundle construction identifies the same
counts in each tensor product of Schur sheaves. & SGA 6 VI 3.3. \\
T-012 & The width-five tensor factors, their finite exterior-cube table,
and their glued LR arrays give the Hilbert components of
\(\mathcal M_5\). & SGA 3 I 4.7.3; SGA 6 V 4.2. \\
T-013 & The independently certified 156-character table and the
displayed tensor-factor maps give the identical statement for
\(\mathcal M_6\). & SGA 3 I 4.7.3; SGA 6 VI 3.3. \\
T-014 & The specified products of highest-weight sections of the
symmetric factors and degree-zero/one sections of the exterior factors
give the claimed nonzero Schur-sheaf summand. & SGA 6 VI 3.3. \\
T-015 & The historical endpoint check concerns the same graded Chow
pullback as G-003; its parameter range is now contained in the cited Sra
bound. & EGA I 9.5.2; M-001 source audit. \\
T-016 & The cubic metric is a morphism of associated bundles; division
of its adjugate by the square of the discriminant is regular by the
proved divisor valuation. & EGA I 9.5.1--9.5.3; SGA 3 I 4.7.2. \\
T-017 & Matrix coefficients identify the specified homogeneous-bundle
sections on \(GL_7/(G_2\times\mu_3)\) with stabilizer invariants. The
harmonic and Pieri proof computes their rank. & SGA 3 I 4.7.2. \\
T-018 & Inclusion into the discriminant localization identifies
extension with a nonnegative divisor valuation. The three polynomial
highest sections generate the computed sector; their degrees are the
polytope equations. & EGA I 9.5.3; SGA 3 I 4.7.3. \\
T-019 & Differentiation of the degree-seven section is the stated
equivariant bundle map. Multiplication with the metric, cometric,
determinant and remaining factors gives the associated stable Schur
components. & SGA 3 I 4.7.2; SGA 6 VI 3.3. \\
T-020 & The filtration by order along the boundary divisor is the
explicit regular-extension subspace of the homogeneous-bundle section
space. & EGA I 9.5.1--9.5.3. \\
\end{longtable}

\subsection{Dictionary for the root
constructions}\label{dictionary-for-the-root-constructions}

\begin{longtable}[]{@{}
  >{\raggedright\arraybackslash}p{(\columnwidth - 4\tabcolsep) * \real{0.3333}}
  >{\raggedright\arraybackslash}p{(\columnwidth - 4\tabcolsep) * \real{0.3333}}
  >{\raggedright\arraybackslash}p{(\columnwidth - 4\tabcolsep) * \real{0.3333}}@{}}
\toprule\noalign{}
\begin{minipage}[b]{\linewidth}\raggedright
Finding
\end{minipage} & \begin{minipage}[b]{\linewidth}\raggedright
Algebraic-geometric statement and proof transport
\end{minipage} & \begin{minipage}[b]{\linewidth}\raggedright
Exact framework
\end{minipage} \\
\midrule\noalign{}
\endhead
\bottomrule\noalign{}
\endlastfoot
G-001 & The exact Pieri multiplication sequence restricts on its
all-ones weight to the augmentation kernel of the permutation module;
determinant tensoring gives the sign twist. Its constant \(A_5\)
restriction proves the stated monomial-basis obstruction. & SGA 3 I
4.7.2--4.7.3. \\
G-002 & The matrices are maps of free modules obtained from the two
constant-group correspondences. Their determinantal ideals describe
their rank strata; the Cauchy--Binet sum computes the same rank over the
real generic fiber. & Stacks Project 07Z6, Lemmas 15.8.2 and 15.8.4;
05P8, Lemma 31.9.6. No unverified EGA Fitting locator is supplied. \\
G-003 & The divided-power coordinate pullback defines the integral
scheme image and its torsion-free, graded-free coordinate algebra. Its
rational grading is the canonical-map image. Hilbert and Quot schemes
represent the stated flat projective families. & EGA I Prop.9.5.1,
Cor.9.5.2, Def.9.5.3, Prop.9.5.4; FGA Exp.221 Thms.3.1--3.2, after
Prop.3.8. \\
G-004 & The two rational Artin correspondences on a degree-two finite
etale scheme have Frobenius eigenvalues \(1,-1\). The actual projector
ranks are one and zero. & Weil II Defs.1.2.1--1.2.2, Variants
1.2.3--1.2.4, stability 1.2.5(i). \\
G-005 & The tree marks and polytope coordinates are identical; the
rank-two determinant line and frame-torsor map transport the known
branching count to Schur-sheaf operations. & SGA 6 VI 3.3; SGA 3 I
4.7.3. \\
G-006 & Binary factorization makes the Chow morphism dominant; its
degreewise coordinate pullback is the specified canonical isomorphism.
Natural Schur-functor decomposition carries every retained multiplicity
map to any larger frame rank. & EGA I 9.5.2; SGA 6 VI 3.3. \\
G-007 & Restriction to the constant \(S_3\) subgroup gives three trace
values on the balanced torus summand. The three character idempotents
give the displayed ranks in the representation Grothendieck ring. & SGA
3 I 4.7.2--4.7.3; SGA 6 V 4.2. \\
G-008 & Exact quartic character certificates produce invariant sections;
multiplication in their domain and the conjugation involution yield the
displayed repeated-part Schur coefficients. & SGA 6 V 4.2, VI 3.3. \\
G-009 & The successive affine flag, torus, and finite-group quotients
have the explicitly presented invariant ring. A flat monomial
degeneration preserves its Hilbert function; disjoint standard-monomial
cells yield the bounded lattice sets. & SGA 3 I 4.7.2--4.7.3; EGA I
9.5.2--9.5.3. \\
G-010 & Adjoining the symmetric algebra of the dual wreath
representation gives its covariant module in auxiliary degree one. Graph
elimination produces its exact graded presentation. & SGA 3 I 4.7.2; SGA
6 VI 3.3; EGA I 9.5.2. \\
G-011 & The two finite-group correspondences act linearly over their
common invariant ring; their image is a finite coherent module. The
syzygy/elimination presentation and module degeneration compute each
actual image rank. & EGA I 9.5.2; SGA 3 I 4.7.2. \\
G-012 & The explicit flag quotient has two monic quadratic equations.
Its four-element free module over the polynomial subring is the exact
geometric source of the four polytopes. & SGA 3 I 4.7.3; EGA I 9.5.3. \\
\end{longtable}

\subsection{Dictionary for rectangular
findings}\label{dictionary-for-rectangular-findings}

R-001 is the determinant line in the balanced Schur sheaf; its \(S_a\)
restriction is the sign character to the stated power. R-002 identifies
the contraction of the two normalized tensor projectors with the signed
Latin-square count and with the Haar period, with all factorials
retained. The finite scheme used in R-006 has one component per Latin
square and the indicated sign function; its pushforward trace gives that
same integer. These are constant finite-group schemes and their
associated locally free sheaves, in SGA 3 I 4.7.2--4.7.3 and SGA 6 VI
3.3.

R-003 restricts the specified wreath permutation representation and
computes its exterior powers; the explicit comparison with the
rectangular Kronecker summand is via those two functor multiplicities.
R-004's indicator variables give a bijection between the cited
rectangular-hook basis and a zero-one polytope. R-005 is the natural
inner-degree multiplication map of sheaves, with its nonzero determinant
pairing and highest-weight dimension. The coordinate-image transfer uses
EGA I 9.5.2 and the exact BIP morphism cited in its proof.

R-007 transports the two prime-adjacent integer congruences through the
exact scalar and Haar identities of R-002, yielding the stated
valuations in \(\mathbb Z_{(p)}\). R-008 works over that local base for
its normalized tensors, then applies the rational central Schur
projector to produce an actual preimage. R-009 uses the graded algebra
morphism induced by the Chow multiplication map; its product equality
persists under the frame-torsor transport. R-010 is an explicit
differential-operator map on the polynomial coordinate sheaf, with its
written kernel vectors. These are the same maps of SGA 3 representations
and EGA I scheme images, rather than dimension comparisons used in place
of maps.

\subsection{Dictionary for every archive
finding}\label{dictionary-for-every-archive-finding}

\begin{longtable}[]{@{}
  >{\raggedright\arraybackslash}p{(\columnwidth - 4\tabcolsep) * \real{0.3333}}
  >{\raggedright\arraybackslash}p{(\columnwidth - 4\tabcolsep) * \real{0.3333}}
  >{\raggedright\arraybackslash}p{(\columnwidth - 4\tabcolsep) * \real{0.3333}}@{}}
\toprule\noalign{}
\begin{minipage}[b]{\linewidth}\raggedright
Finding and archive locus
\end{minipage} & \begin{minipage}[b]{\linewidth}\raggedright
Exact transported statement
\end{minipage} & \begin{minipage}[b]{\linewidth}\raggedright
Locator
\end{minipage} \\
\midrule\noalign{}
\endhead
\bottomrule\noalign{}
\endlastfoot
A-001, 133/019 & The Burnside ring marks, rank map, and signed
permutation linearization compute the displayed discrepancy;
exterior/symmetric operations remain the exact-K0 operations. & SGA 6 V
2.1, 2.2.1, VI 3.3. \\
A-002, 133/006 & The rank augmentation on strict quadratic spaces in
characteristic two makes the proposed exterior-square rank impossible. &
Archive definition; SGA 6 V 2.1 fixes the pre-lambda identities
tested. \\
A-003, 133/006 & The displayed quotient from split orthogonal group
completion to the metabolic Grothendieck--Witt ring kills the computed
nonzero Adams defect. The forgetful maps specify both K0 targets. &
Zibrowius 2015 Def.1.2 and main theorem; SGA 6 V 2.4.1 for the identity
tested. \\
A-004, 133/008 & The written bilinear isometry identifies \(q\perp(-q)\)
with a hyperbolic bundle after inverting two. The characteristic-two
group-scheme representation gives the precise failure of the proposed
relation. & SGA 3 I 4.7.2; archive 133/008. \\
A-005, 133/008 & The natural tensor/hyperbolic isometry proves the ideal
statement; the computed exterior-square image proves failure of descent
through the indicated Witt quotient. & SGA 6 V 2.1; archive quotient
map. \\
A-006, 133/009 & The cited injective etale-cohomology pullback to a
splitting space transports the root formulas. The real Kummer class is
the explicit obstruction to the struck-out descent. & Zibrowius gamma
paper Prop.3.3; SGA 6 V 3.2--3.4, 3.10. \\
A-007, 133/028 & The tensor-power operation and Young induction product
give the exact virtual equality; the rank-two Adams example computes the
failure of an effective-summand reading. & SGA 6 V 4.2 and VI 3.3;
integral completion A-012. \\
A-008, 43/037 & The series involution \(f(t)\mapsto f(-t)^{-1}\)
transports the Witt conventions, with the stated Frobenius/Verschiebung
identities and ghost maps. & Archive 43/037; SGA 6 V 3.2--3.4 for the
related lambda-series substitution. \\
A-009, 2/017 & Restriction is pullback of representations, induction is
the explicit coset tensor construction, and the written map proves
projection. The coinvariant example computes loss of local freeness for
an arbitrary group map. & SGA 3 I 4.7.2; SGA 6 VI 3.3. \\
A-010, 133 and SGA 6 & The free-lambda evaluation, rank-r quotient and
determinant localization are the actual universal homomorphisms; the
gamma substitution computes the rank-kernel filtration. & SGA 6 V 4.2,
4.4, 4.8--4.9; V 3.2--3.4, 3.10; VI 3.3. \\
A-011, 89/018 & The exterior pairing is a morphism of rank-two locally
free sheaves; the displayed reconstruction is equivalent to the zero-sum
relation under its explicitly proved nondegeneracy. & SGA 3 I 4.5;
archive 89/018. \\
A-012, 133/028 & The integral Cauchy filtration is equivariant and
compatible with base change. Three exact torus weight extractions give
the tensor identity in the declared exact equivariant K0 on every
commutatively ringed topos. & ABW II.2.16 and III.1.4; SGA 3 I 4.7.3;
SGA 6 VI 3.3. \\
\end{longtable}

The motivic and Lefschetz expressions on the remaining archive fragments
require their own definitions before a scheme or realization map can be
written. They are recorded in the unresolved part; the proved Artin
example G-004 supplies an actual motive and projector in its stated
scope.

\subsection{Additional explicit dictionaries for the small presentations
and the uniform end-row
family}\label{additional-explicit-dictionaries-for-the-small-presentations-and-the-uniform-end-row-family}

\begin{longtable}[]{@{}
  >{\raggedright\arraybackslash}p{(\columnwidth - 4\tabcolsep) * \real{0.3333}}
  >{\raggedright\arraybackslash}p{(\columnwidth - 4\tabcolsep) * \real{0.3333}}
  >{\raggedright\arraybackslash}p{(\columnwidth - 4\tabcolsep) * \real{0.3333}}@{}}
\toprule\noalign{}
\begin{minipage}[b]{\linewidth}\raggedright
Finding
\end{minipage} & \begin{minipage}[b]{\linewidth}\raggedright
Restatement in schemes, comodules, or Schur operations, with the
proof-carrying map
\end{minipage} & \begin{minipage}[b]{\linewidth}\raggedright
Framework locator
\end{minipage} \\
\midrule\noalign{}
\endhead
\bottomrule\noalign{}
\endlastfoot
T-021 & Pullback of the coefficients of a product of four binary linear
forms sends the quartic coefficient \(x_i\) to
\(\binom4i^{-1}\sum_{|I|=i}\prod_{j\in I}b_j\prod_{j\notin I}a_j\). The
resulting highest-weight coordinate algebra is the stated kernel of the
locally nilpotent derivation, presented by \(T^2+4H^3-f^2HI+f^3J=0\).
The localization and modulo-\(f\) injection in the proof establish the
actual scheme-coordinate isomorphism. & EGA I 9.5.1--9.5.3; SGA 3 I
4.7.2. \\
T-022 & The monic equation gives a free module with basis \(1,T\) over
\(\mathbf Q[f,H,I,J]\). Its two degree equations are exactly the two
polytope systems, so the Schur coefficient is the corresponding graded
rank. & SGA 3 I 4.7.3; SGA 6 VI 3.3. \\
T-023 & The specified weighted initial monomials define a subalgebra
degeneration with the proved semigroup basis; the cusp equation is its
initial relation. The exact basis correspondence proves flatness over
the degeneration parameter and preserves each multigraded rank. & EGA I
9.5.3; the Rees construction and monomial independence are proved in
T-023. \\
T-024 & The flag-face map is the affine incidence morphism with
coordinates \(x_i,y_i,z\), relation \(\sum x_iy_i=0\), and
determinant-one torus quotient
\(\mathbf Q[u,v,z,t_i]/(\sum t_i,uv-\prod t_i)\). Literal row
permutations give the recorded signed action. The dense incidence chart
proves the kernel, rather than inferring an algebra from equal
dimensions. & SGA 3 I 4.7.2--4.7.3; EGA I 9.5.2--9.5.3. \\
T-025 & The alternating-group quotient is a quadratic discriminant
algebra with basis \(1,D\). Taking the remaining involution's
eigenspaces selects the written parity of the standard monomials. Their
weight equations reduce bijectively to \(\sum_{i=2}^m ic_i=R\), giving
the indicated Schur-sheaf multiplicity and its exact support. & SGA 3 I
4.7.3; SGA 6 V 4.2 and VI 3.3. \\
T-026 & The weighted discriminant relation gives a monic quadratic Rees
family, free with basis \(1,D\). The explicit rational rescaling turns
the special fiber into the displayed binomial semigroup algebra. Taking
the diagonal order-two invariants preserves the written parity basis and
flatness over the localized integral base. & EGA I 9.5.3; SGA 3 I
4.7.2. \\
T-027 & Multiplication by \(u^{(n-\beta)_+}v^{(\beta-n)_+}z^\gamma\)
identifies the balanced comodule with
\(\operatorname{sgn}^{\beta}\otimes\operatorname{Sym}^{S}\mathfrak h_m\),
where \(\mathfrak h_m\) is the quotient of the permutation bundle by its
constant line. Taking the \([\mu]\)-covariant summand and computing its
cycle traces gives the tableau generating function. & SGA 3 I
4.7.2--4.7.3; SGA 6 VI 3.3. \\
T-028 & The assignment matrix, descent variables, and partition
variables have a proved one-to-one correspondence with the tableaux and
partitions in that character formula. Hence the same covariant module
has rank equal to the one explicitly described lattice count. The
minimum and next graded terms yield the exact support criterion. & SGA 3
I 4.7.3; SGA 6 V 4.2. \\
T-029 & Each partial-flag face is an affine graded subalgebra. Inverting
its invariant product of selected minors produces an explicit
Laurent-polynomial algebra over its torus quotient. Its dimension, and
the restriction of a finite homogeneous normalization to every face,
force the negative integer \(c_{\{1,2,m-1\}}=2-m\). Thus the specified
finite polynomial-subalgebra morphism cannot exist for the full
highest-weight grading. & EGA I 9.5.3 for the affine algebra
correspondence; SGA 3 I 4.7.3 for torus comodules; all dimension and
face arguments are proved in T-029. \\
T-030 & The \(\omega_2\) face is the Plücker quadric. Its torus quotient
is the trace-free matching plane with a determinant coordinate; the
quotient homomorphism \(S_4\to S_3\) and the displayed alternating-form
intertwiner identify every outer character. Its invariant affine scheme
has equation \(27W^2+Z(D^2+4e^3)=0\). & SGA 3 I 4.7.2--4.7.3; SGA 6 VI
3.3; EGA I 9.5.3. \\
T-031 & Complementary-column exterior multiplication gives the perfect
pairing into \(\det^H\); it transports the finite permutation projector
with the exact sign twist. Multiplication by \(\prod_{i<j}[ij]\) gives
the binary exterior-to-symmetric sheaf isomorphism, including
determinant power \(\binom m2\). These two morphisms carry the
proper-face formulas to the quartic presentation. & SGA 3 I 4.5 and
4.7.2; SGA 6 VI 3.3. \\
\end{longtable}

\subsection{Additional dictionaries for canonical images and rectangular
comparisons}\label{additional-dictionaries-for-canonical-images-and-rectangular-comparisons}

\textbf{G-013.} The affine diagram \(j\Pi_a=\Pi_b\iota\) is defined by
multiplication of forms and appending fixed basis forms. Its coordinate
pullback is precisely
\(h_{b,a}\operatorname{Sym}^b(j^*)=\iota^*h_{b,b}\). Block determinant
evaluation takes the normalized square section to \(I_{a,b}\).
Multiplication in the image coordinate ring then realizes every printed
seed expression. The Apéry paths specify the exponents of actual
preimage sections. The rational denominators specify a localization of
\(\mathbf Z\) for each section. This is the scheme-image construction of
EGA I 9.5.1--9.5.4 and the associated-representation construction of SGA
3 I 4.7.2; its Schur coefficient lies in the representation \(K_0\) of
SGA 6 VI 3.3.

\textbf{G-014.} On the principal open where the pure-power Chow
coefficient is invertible, the ordered-factor coordinates give the
multisymmetric chart. The polynomial identity
\(\mu_a(C_\beta)=uE_\beta\) and the degree-tracked homogenization prove
equality of the image and normalization on each component with
\(an-\lambda_1\le n\). This is a precise graded coordinate-ring
statement; applying the Schur coefficient functional gives the
canonical-map rank. The adjoint tensor projections give the stated
isomorphism on the common first-row range. The relevant affine and
representation locators are EGA I 9.5.2--9.5.3, SGA 3 I 4.7.2--4.7.3,
and SGA 6 VI 3.3.

\textbf{R-011, R-012, and R-015.} Each common-polarization operator is a
natural morphism of polynomial representation bundles. Its integer
highest-weight presentation is specified by flag-minor Reynolds vectors.
A nonzero minor of its evaluated image matrix modulo the displayed prime
is a nonzero integral minor and therefore proves the asserted
generic-fiber injectivity. The certificate lists every source
multiplicity, so the argument covers the entire source functor. In the
language of determinantal ideals, these computations exhibit the
required nonzero minors of each finite multiplicity map; the exact
general locators are Stacks 07Z6, Lemmas 15.8.2 and 15.8.4, and 05P8,
Lemma 31.9.6. The bundle transport is SGA 3 I 4.7.2.

\textbf{R-013.} The morphism \(\nu_{L,b}:(\ell,q)\mapsto\ell^Lq\) is an
integral affine-scheme morphism. Its degree-\(r\) coordinate pullback
equals \(\binom bL^{-r}T_{r,L,b}\) over \(\mathbf Q\), with the scalar
proved on pure products. The kernel is the degree-\(r\) ideal of the
scheme image by EGA I 9.5.2. Multiplication in the ambient polynomial
domain transports the certified vanishing to every lower degree. The
cited Kadish--Landsberg Theorem 1.7 and its section 3 proof provide the
known geometric map; the stated finite boundary kernels and scalar
convention are evaluated here.

\textbf{R-014.} Padding and determinant realization give the actual
inclusion \(\Gamma_n^M\hookrightarrow\mathrm{Det}_M\), hence the
equivariant surjection of coordinate rings. The cited highest-weight
padding injection and the stable rectangular-Kronecker identification
transport this through the finite chain written in R-014. The end-row
lattice count then gives an explicit effective lower bound in the
representation Grothendieck group. The geometric framework is EGA I
9.5.2 and SGA 3 I 4.7.2, with the exact Ikenmeyer--Panova locators
Proposition 2.6(b), Lemma 2.7, Proposition 2.8, and Theorem 2.1.

\subsection{Additional dictionaries for the recovered archive
statements}\label{additional-dictionaries-for-the-recovered-archive-statements}

\begin{longtable}[]{@{}
  >{\raggedright\arraybackslash}p{(\columnwidth - 4\tabcolsep) * \real{0.3333}}
  >{\raggedright\arraybackslash}p{(\columnwidth - 4\tabcolsep) * \real{0.3333}}
  >{\raggedright\arraybackslash}p{(\columnwidth - 4\tabcolsep) * \real{0.3333}}@{}}
\toprule\noalign{}
\begin{minipage}[b]{\linewidth}\raggedright
Finding and archive locus
\end{minipage} & \begin{minipage}[b]{\linewidth}\raggedright
Geometric or exact-K0 restatement and the actual comparison
\end{minipage} & \begin{minipage}[b]{\linewidth}\raggedright
Exact locator
\end{minipage} \\
\midrule\noalign{}
\endhead
\bottomrule\noalign{}
\endlastfoot
A-013, 133/024--029 and supplied study & Injection powers and kernel
partitions give explicit isomorphisms of finite-set-valued functors;
composition and finite limits are computed by the written maps. The
two-copy constant functor evaluated at a singleton disproves the
single-power summand assertion. Applying the constant finite-scheme
functor preserves the specified maps and products. & Archive
133/024--029; SGA 3 I 4.5 for the associated constant actions. \\
A-014, 4/007 & The two chain maps between the de Rham complex and
\(\mathcal O\) have Frobenius composites. Passing to the Fitting-stable
images gives their stated inverse. On a finite free Witt-vector lattice,
Hensel separation of unit roots computes the stable reduction. The
crystalline freeness and algebraicity inputs explicitly admitted in the
original pages remain the exact application gates. & Original cote4, PDF
pp3--7, §§8.3--8.5; Weil II 1.2.1--1.2.5 gives the weight terminology,
not those missing crystalline inputs. \\
A-015, 133/006 & The integral universal iterate polynomial is evaluated
on the concrete modular orthogonal representation. Its nonzero defect in
the split group completion maps to zero under the specified metabolic
quotient. & SGA 6 V 2.4.1, 2.4.2; archive category and quotient in
A-003. \\
A-016, 16/038 & On the explicitly defined graded primitive-string
modules, restriction and pushforward act by truncation and shift. The
corrected projector identities hold on every displayed basis vector. For
\(\mathbf P^1\) and a point, the actual cohomological maps refute the
source's final projected equality. No algebraic-motivic realization of
an unstated lowering operator is inferred. & Original cote16 PDF
pp33--38; SGA 5 III 3.1--3.3 for the correspondence-to-cohomology
construction. \\
A-017, 79/051 & Exterior multiplication of a rank-two
\(\mathbf Z/n\)-module descends through the sign quotient. The identity
\(\Lambda^2(a\,\mathrm{id})=a^2\mathrm{id}\) transports edge and face
types. The wedge tests and the \(n=4\) calculation prove the stated
recovery maps. These finite sets and maps may be represented as constant
finite schemes; the reconstruction from one unlabeled complex remains
separately recorded. & Original cote79 PDF pp47--51, (114)--(127); SGA 3
I 4.5 and 4.7.2. \\
A-018, 43/037 & A morphism from
\(\operatorname{Spec}\mathbf Z[a_1,a_2,\ldots]\) depends on finitely
many coordinates of each target function. Pullback along products of the
universal line series sends \(a_i\) to the elementary symmetric
polynomial \(e_i\); its proved injectivity forces equality of the two
affine group-scheme morphisms. & Archive 43/037; EGA I 9.5.3 for the
affine coordinate dictionary; SGA 6 V 4.2 for the universal
symmetric-operation algebra. \\
A-019, 133/008 & The hyperbolic summand has a nonsingular polar form in
every characteristic. The cited field cancellation isometry removes that
actual summand and identifies the remaining form with the hyperbolic
form of the dimension difference. & Elman--Karpenko--Merkurjev §8, Lemma
8.2 and Theorems 8.3--8.4, pp47--48; archive 133/008. \\
A-020, 133/028--029 study & The composition-factor homomorphism
identifies exact \(G_0\) over \(\mathbf F_p\) with the free group on the
absolutely simple modules. The decomposition matrix relates the Specht
reductions to the simple labels; its rank at degree \(p\) is the proved
modular obstruction to a polynomial induction ring on all \(\rho_i\). &
James Theorem 11.5, p40; Corollaries 12.2--12.3, pp42--43; SGA 6 VI 3.3
for exact representation K0. \\
A-021, 133/028--029 study & Rationalization of finite locally free
representations over a characteristic-zero DVR induces an isomorphism of
exact representation groups, proved through the torsion boundary and
commensurable lattices. It commutes with the actual tensor and induction
functors, yielding the stated polynomial graded ring. & James 8.4,
8.6--8.11, 11.5, 12.2--12.3; SGA 3 I 4.7.2; SGA 6 VI 3.3. \\
A-022, 133/009 and study & The nonisotropic-line pullback is injective
on mod2 étale cohomology. Its actual symmetric-line summand and Kummer
connecting map give \(w_2(H(L))=\bar c_1(L)\). The iterated splitting
and Whitney product yield the displayed general hyperbolic formula in
the proved scheme scope. & EKV Lemma 5.3 and Proposition 5.5,
pp181--184; Zibrowius gamma paper Proposition 3.3, pp19--20; SGA 6 V §3
for the gamma-filtration comparison. \\
A-023, 133/028--029 study & Over the explicit local scheme
\(\operatorname{Spec}(\mathbf Z_{(2)}[\varepsilon]/(2\varepsilon,\varepsilon^2))\),
the rank-one constant-\(C_2\) comodule has character \(1+\varepsilon\).
The determinant functor on exact sequences detects its class outside the
subgroup generated by the two degree-two \(\rho\)-monomials. Rational
base change kills \(\varepsilon\), so the exhibited class difference
also lies in the kernel of rationalization. & SGA 3 I 4.7.2 for
constant-group comodules; SGA 6 VI 3.3 for the exact representation
group; the complete ring and determinant calculation is A-023. \\
\end{longtable}

\subsection{Étale and correspondence locators used in the
transport}\label{uxe9tale-and-correspondence-locators-used-in-the-transport}

For A-009, SGA 4, Exposé IV, \textbf{4.5.1} gives the restriction and
induction morphisms of classifying topoi induced by a group
homomorphism; \textbf{4.5.2} records the profinite version. The explicit
coset tensor map in A-009 is the corresponding locally free
representation calculation. These passages were opened and read, rather
than inferred from the title of the exposé.

For G-004 and the realized cohomological examples of A-016, SGA 5,
Exposé III, \textbf{3.1}, equations \textbf{(3.1.1)--(3.1.3)}, defines a
cohomological correspondence; \textbf{3.2(b)(i)--(ii)} gives the
diagonal and Frobenius-graph examples; \textbf{3.3}, equation
\textbf{(3.3.3)}, gives its action on cohomology. The Artin projectors
of G-004 are \((1\pm F)/2\) after passing to rational \(\ell\)-adic
coefficients. Their finite étale realization and weight-zero eigenvalues
are explicit. The projectors' ranks follow from the action on that
two-dimensional space; no mod2 projector with denominator two is
asserted.

\subsection{Exact dictionaries for the general chart, chamber sheaves,
and
truncation}\label{exact-dictionaries-for-the-general-chart-chamber-sheaves-and-truncation}

\textbf{T-032.} The balanced diagonal-torus comodule algebra is
localized at the explicit invariant section \(u=\prod x_i\). Its
coordinate homomorphism sends each flag minor to
\(u\mathbf1\wedge\bar w_2\wedge\cdots\), and the prescribed
centered-matrix construction proves the inverse on the principal open.
The extension subspace is the ideal-divisibility condition
\(u^k\mid F(v(X))\) in the original affine coordinate ring. This is an
explicit isomorphism of affine schemes with constant symmetric-group
action (EGA I 9.5.3; SGA 3 I 4.7.2--4.7.3). Its Schur-coefficient
functor is the corresponding exact representation-group functional in
SGA 6 VI 3.3.

\textbf{T-033--T-034.} The homomorphism \(S_4\to SO(W)\), with
\(W=\mathfrak h_4\otimes\operatorname{sgn}\), is specified by the
invariant form and determinant-one row-permutation action. Restriction
of the Schur comodule, the harmonic quotient by that form, and the
explicit five character projectors give the finite polynomial counts.
The complementary-column pairing transports the upper chamber with its
computed determinant character. These are restriction and direct-summand
morphisms of the SGA 3 I 4.7.2 representation category; their graded
classes and exterior determinant operations are those of SGA 6 VI 3.3.
The Elliott character calculation is proved in full in T-033, including
the required intertwining action.

\textbf{T-035--T-036.} The affine incidence presentation yields the flag
scheme \(\mathcal F\subset\mathbf P(W)\times\mathbf P(W)\) and its
projective-line-bundle morphism \(\pi\). The exact coordinate equations
identify each original divisor valuation with the order in the
pulled-back ideal of one of four specified points. Thus the balanced
comodule is the global-section representation of
\(\pi^*\mathcal I_Z^k\otimes\mathcal O(c,b)\), with the printed
determinant character. The two-quadric ideal resolution and the
tautological quotient sequence calculate those sections, including the
nonzero first-cohomology summand in T-036. The group-equivariant
splitting is obtained by averaging an actual linear projection. These
affine, ideal-sheaf, and locally free sheaf maps provide the full
geometric restatement; the associated representation and K0 maps are SGA
3 I 4.7.2 and SGA 6 VI 3.3.

\textbf{R-016.} The normalized pullback of R-013 identifies each of the
three computed kernels with the degree component of the ideal sheaf of
the padded-form scheme image. Every source summand is listed; the
selected nonzero minors prove the ranks and the displayed determinant
polynomial supplies the entire missing irreducible summand. The
scheme-image interpretation is EGA I 9.5.2, with the same
determinantal-ideal maps already specified for R-015.

\textbf{R-017--R-018.} Separate homogenization and first-coordinate
restriction are inverse maps between the indicated degree-bounded
polynomial spaces. Coefficient extraction becomes the explicitly
diagonal truncation operator with weights \(\prod_i\binom{b-d_i}{L}\).
The flag-minor generators bound each block degree by the second row.
Inverting those actual nonzero weights proves injectivity on the
specified entire Schur-isotypic components, and applying each
finite-group projector preserves it. Via R-013 this is the vanishing of
precisely the indicated graded ideal components of the padded-form
scheme image. These are equivariant coordinate-sheaf maps under EGA I
9.5.2--9.5.3 and SGA 3 I 4.7.2--4.7.3; their exact Schur operations lie
in SGA 6 VI 3.3.

\subsection{Five-point modules and the complete four-point
formula}\label{five-point-modules-and-the-complete-four-point-formula}

\textbf{T-037.} The exact four-point incidence quotient is free over the
polynomial algebra of its two quadratic generators through the averaged
multiplication section. The two-quadric power resolution and
tautological quotient sequence identify ordinary powers with the
required flag sections in the proved degree range; the boundary
correction is the displayed connecting cohomology module. The explicit
complementary pairing and sign twist transport the resulting positive
modules to every balanced Schur weight. These are ideal-sheaf,
associated-bundle and exact-representation maps in EGA I 9.5.1--9.5.3,
SGA 3 I 4.7.2--4.7.3 and SGA 6 VI 3.3.

\textbf{G-015--G-016.} The general balanced chart identifies extension
across each original coordinate divisor with the ordinary power of a
smooth flag fiber's ideal. Its leading normal ratios are independent
over the fiber function field, including every higher Plücker frame. The
point-ideal generators and canonical affine evaluation maps identify the
saturated projective-frame scheme and its actual permutation quotient.
The schemes and coordinate maps are the EGA I 9.5.1--9.5.3
constructions; the linearized flag line bundle is transported by the SGA
3 associated-bundle map, and its Schur class is the SGA 6 VI 3.3
operation. The exact Schur/section locator used is Magyar, Proposition1
pp5--6 and Proposition4(a)--(c) p10.

\textbf{T-038, A-028--A-029.} The five-point ideal has the written
equivariant alternating resolution. Its actual quadratic image is the
displayed cubic hypersurface, with the precise six-letter intertwiner.
The Cox comparison sends exceptional sections and each Plücker bracket
by (Cox.1), and identifies every Picard-graded section by the
normal-order calculation. The ordinary Rees map is the specified
inclusion into that Cox ring. Applied Kustin--Polini--Ulrich
formula6.1.4 and Lemma6.3/Corollary6.4 give the linear-power complexes
after every local hypothesis is checked. Saturation and cohomology
follow from those complexes. These are concrete graded scheme and
coherent-sheaf morphisms under EGA I 9.5.1--9.5.3; the flat-family
Hilbert/Quot interpretation is FGA Exposé221 Theorems3.1--3.2. No new
generic-position assertion is used.

\textbf{G-017.} The polynomial matrix (Q.9) presents the actual
quotient-bundle section module. Its determinant, adjugate and generic
evaluation prove an injection; the established ideal-power cohomology
proves surjectivity degree by degree. Four explicit invariant parameters
are regular by the determinant argument, and averaging a graded section
gives the actual free multiplication isomorphism (Q.21). The
corresponding coefficient is the constant-group covariant summand of the
geometric module, carried by G-015 and T-003. The exact K0 operation is
SGA6 VI3.3, and all determinant dualities are specified exterior
pairings of locally free sheaves.

\textbf{G-018.} The dual tautological inclusion and linear
multiplication identify the exterior quotient module with the concrete
kernel (W.6). Stabilizer-invariant affine charts give the entire
fat-point jet action, including its line-bundle trivialization. The
finite Newton count is the exact rational character identity of that
kernel. The EGA coherent sheaf, SGA3 comodule and SGA6
exterior-operation maps are precisely those just specified.

\textbf{G-019 and T-039.} The symmetric quotient sequence (SQ.4) is
locally split before tensoring the ideal. Its global cohomology gives
the exact representation-K0 formula, including the one nonzero
connecting term at point order1. Relative exterior pairings give the
determinant character in (SQ.3). Its fixed bivariate positive formula
therefore computes Schur-operation multiplicities of actual
ideal-twisted flag sections. The transport uses SGA3 I4.7.2--4.7.3 and
SGA6 V4.2,4.4,4.8--4.9,VI3.3, with the projective-family interpretation
of FGA221 Theorems3.1--3.2.

\textbf{T-040.} The degree-one Rees component has the explicit B-matrix
resolution. The tautological inclusion into its tensor product with E
gives the displayed quotient, whose full free resolution is verified by
the rank-three reduction modulo the cubic. This supplies the actual
module behind the positive residue formula. Its four-parameter
obstruction is proved by the exact multiplication-sequence Hilbert
series, which would have a negative finite quotient coefficient. The
ordinary-power, saturation and quotient maps specify the geometric
comparison.

\textbf{T-041.} The hook Schur functor is the kernel of the split
symmetric multiplication map. The explicit two-step tautological hook
resolution and its top-Cech identity prove all needed cohomology groups.
Multiplication by the five quadrics is the actual section isomorphism
(5H.12). These kernel, exterior-complement and symmetric operations are
the locally free Schur operations transported in SGA6 VI3.3.

\subsection{Further actual canonical-image
dictionaries}\label{further-actual-canonical-image-dictionaries}

\textbf{R-019--R-020.} Complementary column pairings identify the entire
first-row source band with a multilinear symmetric-group component. The
explicit cofactor highest vector and nonzero coefficient prove the
actual polarization rank. Through the padded-form scheme-image map
R-013, this removes precisely those graded ideal components. The
transport is the EGA I9.5.2 scheme-image coordinate map and the
SGA3/SGA6 associated Schur-sheaf map.

\textbf{R-021.} The block specialization and multiplication in the
graded Chow coordinate algebra carry prime-adjacent square determinant
preimages to every stated rectangular and row-extended preimage.
Nagura's exact prime interval determines the nonnegative exponents. Thus
the nonzero coordinate section is in the actual canonical image; its
determinant and highest-weight lines are the specified Schur operations
in SGA6 VI3.3.

\textbf{R-022--R-024.} Normalized last-slot polarization sends the Chow
image annihilator into the exact common-polarization kernel or into the
lower-degree annihilator via (P.11). Slot transvections prove the
rectangular prolongation equality, and the actual determinant preimage
removes that summand. The two Pieri inclusions then prove the stated
whole-isotypic vanishing of the graded ideal, including every
multiplicity. These are maps of coordinate modules of the scheme image
in EGA I9.5.2, transported to associated Schur sheaves by SGA3 I4.7.2
and SGA6 VI3.3. The degreewise minors give their rank strata as in the
G-002 determinantal dictionary.

\part{Source register and verification}

\section{Verified source locators for the archive and geometric
dictionaries}\label{verified-source-locators-for-the-archive-and-geometric-dictionaries}

Run
\texttt{\small /\allowbreak{}root/\allowbreak{}archive\_\allowbreak{}audit};
model GPT-6 Astra Pro; effort not exposed. This document records source
inspection, not an independent reproof of every cited theorem. Archive
evidence consists of the supplied TeX transcriptions. No manuscript scan
images or Lean files were supplied.

\subsection{Archive and supplied secondary
material}\label{archive-and-supplied-secondary-material}

All thirteen \texttt{\small ARCHIVE\_\allowbreak{}*.tex} and all of
\texttt{\small SA17\_\allowbreak{}STUDY\_\allowbreak{}cote133\_\allowbreak{}lambda.md}
were read. The report preserves the difference between transcribed
statements, transcription uncertainty, supplied study additions, and
current-run arguments. In particular, the study's alleged equivariant
finite-set check on 133/019 is refuted in A-001. No claim about scans
merely mentioned in the study is treated as a new scan inspection.

\subsection{SGA 6: exact K theory, exterior operations, and the free
λ-ring}\label{sga-6-exact-k-theory-exterior-operations-and-the-free-ux3bb-ring}

Bibliography: P. Berthelot, A. Grothendieck, L. Illusie, with
collaborators, \emph{Théorie des intersections et théorème de
Riemann--Roch}, SGA 6, Lecture Notes in Mathematics 225, Springer, 1971.

\begin{longtable}[]{@{}
  >{\raggedright\arraybackslash}p{(\columnwidth - 4\tabcolsep) * \real{0.3333}}
  >{\raggedright\arraybackslash}p{(\columnwidth - 4\tabcolsep) * \real{0.3333}}
  >{\raggedright\arraybackslash}p{(\columnwidth - 4\tabcolsep) * \real{0.3333}}@{}}
\toprule\noalign{}
\begin{minipage}[b]{\linewidth}\raggedright
Locator
\end{minipage} & \begin{minipage}[b]{\linewidth}\raggedright
Statement used
\end{minipage} & \begin{minipage}[b]{\linewidth}\raggedright
Inspection
\end{minipage} \\
\midrule\noalign{}
\endhead
\bottomrule\noalign{}
\endlastfoot
Exposé V, Definition 2.1 & Pre-λ definition & French mathematical
transcription inspected \\
V, Lemma 2.2.1 & Exterior-power filtration for short exact sequences &
Exact citation confirmed in Zibrowius 2016, Lemma 2.2 \\
V, 2.4.1 & Universal product and iterate identities & Exact citation
confirmed in Zibrowius 2015, §3.1 \\
V, 3.2--3.4 & γ substitution and equivalent formulas & French
mathematical transcription inspected \\
V, 3.10 & γ filtration & French mathematical transcription inspected \\
V, Proposition 4.2; Definition 4.4 & Unique evaluation of the free
λ-ring & Statement and proof text inspected \\
V, Proposition 4.8; Corollary 4.9 & Finite-rank quotient and its λ-ideal
& Statement and proof text inspected \\
V, §6 & Chern-ring universal construction & Section located; entire
proof not independently rederived \\
VI, Theorem 3.3 & λ structure on exact K₀ of a locally ringed topos;
transfer via frame torsors & Theorem and proof transcription inspected,
exact locator corroborated in Zibrowius 2015 p.1057 and 2016 §2.2 \\
\end{longtable}

Inspected transcriptions:
\url{https://grothendiecksga.com/read/sga6/fr/V-3.html} ;
\url{https://grothendiecksga.com/read/sga6/fr/V-4.html} ;
\url{https://grothendiecksga.com/read/sga6/fr/VI-3.html} . Local
HTML/text copies are in this directory. The attempted original scan
download from \texttt{\small iut-\allowbreak{}theory.org} returned HTTP
503. The paper must not say the original SGA 6 scan was visually
inspected.

\subsection{Zibrowius: precise quadratic-form quotient and
splitting}\label{zibrowius-precise-quadratic-form-quotient-and-splitting}

Marcus Zibrowius, \emph{Symmetric representation rings are λ-rings}, New
York J. Math. 21 (2015), 1055--1092. Primary PDF:
\url{https://nyjm.albany.edu/j/2015/21-49v.pdf} .

\begin{itemize}
\tightlist
\item
  Main theorem, p.1056: special λ structure on the Grothendieck--Witt
  representation ring defined with metabolic relations.
\item
  Lemma 1.1, p.1059: the natural hyperbolic isometry for a symmetric
  object and its negative, with 2 invertible.
\item
  Definition 1.2, p.1060: the exact K₀ relations and, explicitly, the
  additional relation
  \texttt{\small [\allowbreak{}M]\allowbreak{}=\allowbreak{}[\allowbreak{}H(\allowbreak{}L)\allowbreak{}]\allowbreak{}}
  for metabolic \texttt{\small M} with Lagrangian \texttt{\small L}.
\item
  §3.1, pp.1068--1069: universal special-λ identities and their SGA 6
  locator.
\end{itemize}

Marcus Zibrowius, \emph{The γ-filtration on the Witt ring of a scheme},
arXiv:1411.3198v2, submitted 11 August 2016, title-page date 12 August
2016. Primary PDF: \url{https://arxiv.org/pdf/1411.3198} .

\begin{itemize}
\tightlist
\item
  Theorem 2, displayed p.3, and Theorem 2.5, displayed p.11: schemes
  over a field of characteristic different from two.
\item
  Proposition 2.15, displayed pp.17--18: orthogonal torsors carrying
  pairs of symmetric bundles.
\item
  \textbf{Proposition 3.3, displayed pp.19--20: the preceding sentence
  explicitly says ``any scheme over Z{[}1/2{]}''.} This proposition
  supplies a morphism that splits the given symmetric bundle into
  symmetric lines and induces an injection on étale cohomology with Z/2
  coefficients. Its scope is broader than Theorem 2. The report's
  scheme-over-Z{[}1/2{]} scope is verified.
\item
  Proposition 4.8, displayed p.23: again explicitly over Z{[}1/2{]},
  consistent with that broader splitting scope; not needed for our main
  argument.
\end{itemize}

\subsection{Macdonald: symmetric functions and
Schur--Weyl}\label{macdonald-symmetric-functions-and-schurweyl}

I. G. Macdonald, \emph{Symmetric Functions and Hall Polynomials}, second
edition, Oxford, 1995. Primary book PDF inspected:
\url{https://math.berkeley.edu/~corteel/MATH249/macdonald.pdf} . Local
copy \texttt{\small Macdonald.pdf} and extracted text
\texttt{\small Macdonald.txt}.

\begin{longtable}[]{@{}
  >{\raggedright\arraybackslash}p{(\columnwidth - 2\tabcolsep) * \real{0.5000}}
  >{\raggedright\arraybackslash}p{(\columnwidth - 2\tabcolsep) * \real{0.5000}}@{}}
\toprule\noalign{}
\begin{minipage}[b]{\linewidth}\raggedright
Locator and printed page
\end{minipage} & \begin{minipage}[b]{\linewidth}\raggedright
Exact role
\end{minipage} \\
\midrule\noalign{}
\endhead
\bottomrule\noalign{}
\endlastfoot
I (4.2), p.62 & Cauchy kernel expanded as the sum of monomial times
complete symmetric functions \\
I (4.3), p.63 & Cauchy kernel expanded as the sum of Schur times Schur
functions \\
I (7.3), pp.113--114 & Frobenius characteristic is an isometric
graded-ring isomorphism; its proof sends induced trivial Young modules
to \texttt{\small h\_\allowbreak{}ν} \\
I (7.5)--(7.6), p.114 & Irreducible character basis maps to the Schur
basis \\
Appendix A (5.3)--(5.4), p.155 & Homogeneous polynomial functors from
symmetric-group modules and the characteristic-zero equivalence \\
Appendix A (5.6), p.156 & Functorial tensor-power decomposition
\texttt{\small Tⁿ(\allowbreak{}X)\allowbreak{}=\allowbreak{}⊕\allowbreak{}M\_\allowbreak{}λ⊗\allowbreak{}F\_\allowbreak{}λ(\allowbreak{}X)\allowbreak{}} \\
Appendix A (7.2), p.160 & Character of the irreducible polynomial
functor is \texttt{\small s\_\allowbreak{}λ} \\
Appendix A (7.4), p.160 & Nonnegative integral Schur coefficients for
Schur plethysm through functor composition \\
\end{longtable}

These numbered statements were inspected in the primary PDF, including
the previously pending Appendix A (5.6). The PDF's page index is printed
page plus 10.

\subsection{Akin--Buchsbaum--Weyman: the integral Cauchy
filtration}\label{akinbuchsbaumweyman-the-integral-cauchy-filtration}

K. Akin, D. A. Buchsbaum, J. Weyman, \emph{Schur functors and Schur
complexes}, Advances in Mathematics 44 (1982), 207--278, DOI
10.1016/0001-8708(82)90039-1. The author's current research page is
\url{https://jerzyweyman.webflow.io/research/schur} . Its linked
original paper was downloaded completely from
\url{https://cdn.prod.website-files.com/69c19dc391eb37d617c813ce/69cc3be3405f5e9e171b7db4_JW\%231_Schur.pdf}
and retained as \texttt{\small ABW1982.pdf}, with extracted text
\texttt{\small ABW1982.txt}.

\begin{itemize}
\tightlist
\item
  Theorem II.2.16, printed p.232: the universal standard-tableau basis
  for a Schur module, proving freeness and compatibility with arbitrary
  scalar extension.
\item
  Theorem III.1.4, printed pp.245--246: the natural integral filtration
  of a symmetric power of a tensor product, with tensor products of
  Schur modules as its associated graded.
\end{itemize}

The paper's indexing uses
\texttt{\small L\_\allowbreak{}(\allowbreak{}a)\allowbreak{}=\allowbreak{}Λᵃ}.
A-012 explicitly defines
\texttt{\small S\_\allowbreak{}λ\textasciicircum{}int=\allowbreak{}L\_\allowbreak{}(\allowbreak{}λ′)\allowbreak{}}
to translate its convention to the archive's usual Schur partitions.
Three exact torus-weight extractions from this one filtration prove the
arbitrary-characteristic tensor-power formula, including its virtual
exact-K₀ extension. The source theorem is carried as verified
background; A-012's deduction is independently checked in full.

\subsection{Witt vectors and Borger}\label{witt-vectors-and-borger}

Michiel Hazewinkel, \emph{Witt vectors. Part 1}, primary author PDF:
\url{https://www.sas.rochester.edu/mth/sites/doug-ravenel/otherpapers/Hazewinkel-Witt.pdf}
.

\begin{itemize}
\tightlist
\item
  §§14.1--14.4, pp.71--72: big Witt operation polynomials in coordinate
  \texttt{\small n} involve only divisor coordinates. The resulting
  \texttt{\small p}-power coordinate polynomials define the
  \texttt{\small p}-typical Witt ring functor and the natural quotient
  from the big Witt functor.
\item
  Definition 14.5 and §14.8, pp.72--73: divisor-stable index sets
  (``nests'') and their Witt functors.
\item
  §14.15, equation (14.16): canonical big-to-typical projection. No
  \texttt{\small p}-local hypothesis is needed for the quotient. The
  stronger decomposition involving sections later in §14.21 uses
  \texttt{\small p}-local rings and is not invoked here.
\end{itemize}

James Borger, supplied \texttt{\small lrfoe13.tex}, §§1.1--1.2 and §2.2,
inspected in the input file. The first locators give the free λ-ring,
arithmetic jet-space, and Witt-space adjunction framework; §2.2 gives
monoidal Frobenius lifts. The report verifies the explicit
character-ring map
\texttt{\small ψᵖ f(\allowbreak{}x)\allowbreak{}=\allowbreak{}f(\allowbreak{}xᵖ)\allowbreak{}}
and its congruence modulo \texttt{\small p} directly.

\subsection{EGA I: schematic image}\label{ega-i-schematic-image}

A. Grothendieck and J. Dieudonné, \emph{Éléments de géométrie
algébrique. I. Le langage des schémas}, Publ. Math. IHÉS 4 (1960),
5--228. Primary original PDF downloaded completely and its text
inspected: \url{https://www.numdam.org/item/PMIHES_1960__4__5_0.pdf} .
Local copies \texttt{\small EGAI.pdf}, \texttt{\small EGAI.txt}.

\begin{longtable}[]{@{}
  >{\raggedright\arraybackslash}p{(\columnwidth - 2\tabcolsep) * \real{0.5000}}
  >{\raggedright\arraybackslash}p{(\columnwidth - 2\tabcolsep) * \real{0.5000}}@{}}
\toprule\noalign{}
\begin{minipage}[b]{\linewidth}\raggedright
Locator and printed page
\end{minipage} & \begin{minipage}[b]{\linewidth}\raggedright
Exact role
\end{minipage} \\
\midrule\noalign{}
\endhead
\bottomrule\noalign{}
\endlastfoot
Proposition 9.5.1, p.176 & Existence of a smallest closed subprescheme
through which the morphism factors when direct image of the structure
sheaf is quasi-coherent \\
Corollary 9.5.2, p.176 & The defining ideal is
\texttt{\small ker(\allowbreak{}O\_\allowbreak{}Y→\allowbreak{}f\_\allowbreak{}*O\_\allowbreak{}X)\allowbreak{}} \\
Definition 9.5.3, p.177 & Names that closed image \\
Proposition 9.5.4, p.177 & Its underlying space is the closure of the
image \\
Proposition 9.5.8, p.178 & Formation commutes with restriction to an
open of the target; the following paragraph points forward to flat base
change \\
Proposition 9.5.9, p.178 & Reducedness and irreducibility/integrality
pass to the closed image in the stated range \\
\end{longtable}

For the parent's affine morphism
\texttt{\small μ:\allowbreak{}Spec B→\allowbreak{}Spec C}, the actual
algebra map \texttt{\small C→\allowbreak{}B} identifies the closed image
with
\texttt{\small Spec(\allowbreak{}C/\allowbreak{}ker μ*)\allowbreak{}=\allowbreak{}Spec(\allowbreak{}im μ*)\allowbreak{}}.
Affineness gives the required quasi-coherence directly. The elementary
tensor-exactness proof gives compatibility with a flat extension of
scalars in this affine case; no unverified global EGA base-change number
is required.

\subsection{FGA: Quot and Hilbert}\label{fga-quot-and-hilbert}

A. Grothendieck, \emph{Techniques de construction et théorèmes
d'existence en géométrie algébrique. IV. Les schémas de Hilbert},
Séminaire Bourbaki, exposé 221 (1960--1961), pp.249--276 in the Numdam
collected pagination. Primary PDF inspected:
\url{https://www.numdam.org/item/SB_1960-1961__6__249_0.pdf} .

\begin{itemize}
\tightlist
\item
  Theorem 3.1, printed p.259 (PDF index 11): Quot represented by a
  disjoint sum of projective schemes in the setting stated there.
\item
  Theorem 3.2, printed p.260 (PDF index 12): the
  fixed-Hilbert-polynomial Quot piece is projective.
\item
  Following Proposition 3.8, printed p.265 (PDF index 17):
  \texttt{\small Quot(\allowbreak{}O\_\allowbreak{}X)\allowbreak{}} is
  the Hilbert functor representing flat families of closed subschemes.
\item
  Remarks 3.9, printed p.266: reduction used for projective schemes over
  Z.
\end{itemize}

\subsection{Fitting ideals: verified alternative locator, no invented
EGA
citation}\label{fitting-ideals-verified-alternative-locator-no-invented-ega-citation}

The requested exact EGA locator for the determinantal/Fitting
formulation was not established in this run. The precise primary
reference inspected is the Stacks Project, Section 15.8, tag
\texttt{\small 07Z6}: Lemma 15.8.2 defines the presentation-independent
ideal of minors; Definition 15.8.3 names it; Lemma 15.8.4(3) proves
arbitrary scalar-extension compatibility.
\url{https://stacks.math.columbia.edu/tag/07Z6} . For locally closed
constant-rank strata, Lemma 31.9.6, tag \texttt{\small 05P8}, supplies
the sheaf statement: \url{https://stacks.math.columbia.edu/tag/05P8} .

In the parent's explicit matrix model, all claims about a rank bound are
proved directly by the vanishing/nonvanishing of minors. Those proofs do
not require an inaccurately numbered EGA theorem.

\subsection{SGA 3: the group scheme and comodule
dictionary}\label{sga-3-the-group-scheme-and-comodule-dictionary}

M. Demazure and A. Grothendieck, \emph{Schémas en groupes}, SGA 3,
Exposé I, by M. Demazure. The primary editor-maintained PDFs were
located and inspected:
\url{https://webusers.imj-prg.fr/~patrick.polo/SGA3/origExp1.pdf}
(original text in the editor's typesetting, carrying the original
marginal pagination) and
\url{https://webusers.imj-prg.fr/~patrick.polo/SGA3/Exp1-13oct24.pdf}
(corrected edition of 13 October 2024).

\begin{itemize}
\tightlist
\item
  I §4.5, original marginal p.29: constructs the classical general
  linear group by the determinant-open coordinate algebra. An earlier
  tentative §4.6 locator is corrected to \textbf{§4.5}.
\item
  I Definition 4.6.1 and 4.6.3: the
  \texttt{\small V(\allowbreak{}F)\allowbreak{}} and
  \texttt{\small W(\allowbreak{}F)\allowbreak{}} module functors, with
  the former represented by \texttt{\small Spec Sym F}; EGA II 1.7.8 is
  explicitly cited for this representing scheme.
\item
  I Proposition 4.7.2, original marginal pp.34--35; corrected edition
  displayed pp.29--30: a group action is a coaction
  \texttt{\small F→\allowbreak{}F⊗\allowbreak{}A(\allowbreak{}G)\allowbreak{}}
  satisfying the counit and coassociativity diagrams. This identifies
  representations with comodules and supplies the precise dictionary for
  integral polynomial representations.
\item
  I Proposition 4.7.3, original marginal pp.35--36: diagonalizable-group
  actions correspond to weight gradings, by the orthogonal projectors
  given by coefficient extraction in the group algebra.
\end{itemize}

No manuscript scan inspection is implied by these typeset primary
editions.

\subsection{Weil II: purity for the finite étale sign
example}\label{weil-ii-purity-for-the-finite-uxe9tale-sign-example}

P. Deligne, \emph{La conjecture de Weil. II}, Publ. Math. IHÉS
\textbf{52} (1980), 137--252; DOI 10.1007/BF02684780. Primary original
scan inspected at
\url{https://publications.ias.edu/sites/default/files/Number40.pdf} .
The IAS landing-page metadata says volume 42; the correct journal volume
is 52, as recorded by Numdam and the DOI.

\begin{itemize}
\tightlist
\item
  Definition 1.2.1 and Definition 1.2.2, printed p.153 (PDF index 16):
  numerical weights and pointwise purity of a sheaf.
\item
  Variants 1.2.3--1.2.4, p.154: terminology for Galois representations
  over a finite field and vector spaces with Frobenius.
\item
  Stability 1.2.5(i), p.154: pointwise pure sheaves of a fixed weight
  are stable under subquotients and finite direct image.
\item
  §1.1.7, pp.150--151: geometric Frobenius is inverse to the field's
  power Frobenius, and the finite-field
  sheaf/continuous-Galois-representation equivalence is spelled out.
\end{itemize}

The sign summand of the degree-two finite étale permutation
representation has geometric Frobenius \texttt{\small −1}; every complex
conjugate has absolute value one. This directly verifies weight zero
from the cited definitions. It also follows from the finite-direct-image
and subobject assertions in 1.2.5(i). No general cohomological purity
theorem is needed for that example.

\subsection{Remaining historical-source
gates}\label{remaining-historical-source-gates}

\begin{itemize}
\tightlist
\item
  Original SGA 6 scan visual verification was not completed; exact
  source numbering is corroborated as detailed above.
\item
  An exact EGA Fitting-ideal locator remains unverified; the Stacks
  primary locators above are verified.
\end{itemize}

\subsection{Field Witt cancellation, exact scope
checked}\label{field-witt-cancellation-exact-scope-checked}

Elman--Karpenko--Merkurjev, \emph{The Algebraic and Geometric Theory of
Quadratic Forms}, AMS Colloquium Publications56 (2008), §8: Lemma8.2 and
Theorem8.3 pp47--48, Theorem8.4 p48. Author-hosted PDF:
\url{https://www.math.ucla.edu/~merkurev/Book/Kniga-final/Kniga.pdf} .
Full proof of these three statements read; characteristic-free
cancellation requires zero polar radical for the cancelled summand. This
covers the hyperbolic stabilizer in original133/008. Local copies
EKM2008.pdf/.txt; consequence A019 in FIELD\_CANCELLATION.md.

\subsection{James: modular rank and the integral Specht
lattices}\label{james-modular-rank-and-the-integral-specht-lattices}

G.D. James, \emph{The Representation Theory of the Symmetric Groups},
Lecture Notes in Mathematics682, Springer,1978; DOI10.1007/BFb0067708.
Public university-hosted original book scan:
\url{https://www-users.cse.umn.edu/~webb/oldteaching/Year2010-11/the-representation-theory-of-the-symmetric-groups-SLN.pdf}
. Local copies \texttt{\small James1978.pdf/\allowbreak{}.txt}.

\begin{itemize}
\tightlist
\item
  Theorem8.4, pp29--30: standard polytabloid basis.
  Corollaries8.6--8.11, pp30--31: integral coefficients/action matrices
  and reduction of the integral lattice to the modular Specht module.
\item
  Definition10.1 and Lemma10.2, p36: p-regular partition labels and
  equality of their count with the count of p-regular conjugacy classes,
  by the displayed Euler-product identity.
\item
  Theorem11.5, p40: complete simple-module list over every field,
  absolute irreducibility and splitting over every field. Its proof and
  the explicitly cited general modular-theory dependencies were read.
\item
  Corollary12.2, p42, and Corollary12.3, p43: the p-regular Specht rows
  give a unitriangular decomposition submatrix. The preceding
  Theorem12.1 proof was read as well.
\end{itemize}

These statements supply the exact basis/absolute-base-change input for
A-020 and A-021. The torsion boundary, commensurability and
generator-preserving rationalization proof in
\texttt{\small MIXED\_\allowbreak{}CHAR\_\allowbreak{}LATTICES.md} were
checked independently. The original standard-basis page locator
is29--31, rather than32--33.

\subsection{EKV: the hyperbolic characteristic-class
formula}\label{ekv-the-hyperbolic-characteristic-class-formula}

H. Esnault, B. Kahn, E. Viehweg, \emph{Coverings with odd ramification
and Stiefel--Whitney classes}, Journal für die reine und angewandte
Mathematik441 (1993),145--188. Author-hosted original scan:
\url{https://page.mi.fu-berlin.de/esnault/preprints/helene/29-es_ka_vi.pdf}
. Local copy \texttt{\small EKV1993.pdf}.

\begin{itemize}
\tightlist
\item
  §5.1, p179: scheme over
  \texttt{\small Z[\allowbreak{}1/\allowbreak{}2]\allowbreak{}} and
  unimodular quadratic bundles.
\item
  §5.2, pp179--181: the nonisotropic-line pullback and its
  étale-cohomology decomposition.
\item
  Lemma5.3, pp181--182: the line-class identity
  \texttt{\small c̄₁(\allowbreak{}M)\allowbreak{}=\allowbreak{}w₁(\allowbreak{}M)\allowbreak{}²+\allowbreak{}(\allowbreak{}-\allowbreak{}1)\allowbreak{}w₁(\allowbreak{}M)\allowbreak{}},
  with the Kummer/twisted-Bockstein cocycle proof.
\item
  §5.4, p182: regularity is introduced. This hypothesis is retained when
  citing the following proposition.
\item
  Proposition5.5, pp183--184: the full metabolic/hyperbolic class
  formula in the Chern classes of the Lagrangian. Both pages and the
  complete proof were visually read. The rank-two proof's integral
  Picard-isomorphism sentence needs qualification; A-022 gives a direct
  argument using injective mod2 cohomology instead.
\end{itemize}

Marcus Zibrowius, \emph{Witt groups of curves and surfaces}, author PDF
\url{https://www.math.uni-duesseldorf.de/~zibrowius/Research/Zibrowius_WCS.pdf}
, §1.4, printedp8/PDFpage9, equations(6)--(8), reproduces the formula
with exact reference to EKVProposition5.5. Its preceding statement
explicitly discusses schemes over
\texttt{\small Z[\allowbreak{}1/\allowbreak{}2]\allowbreak{}}.
Zibrowius's gamma-filtration paper Proposition3.3, pp19--20,
independently supplies the injectivity used in A-022's
arbitrary-\texttt{\small Z[\allowbreak{}1/\allowbreak{}2]\allowbreak{}}-scheme
line-bundle proof.

A-022 now proves the full hyperbolic formula over every
\texttt{\small Z[\allowbreak{}1/\allowbreak{}2]\allowbreak{}}-scheme
using that same exact injectivity result. The construction applies
\texttt{\small P\_\allowbreak{}nd} to
\texttt{\small H(\allowbreak{}E)\allowbreak{}}; its tautological pair
\texttt{\small (\allowbreak{}v,\allowbreak{}f)\allowbreak{}} gives the
globally defined projector
\texttt{\small z↦v f(\allowbreak{}z)\allowbreak{}/\allowbreak{}f(\allowbreak{}v)\allowbreak{}},
so \texttt{\small E} splits a line after this cohomologically injective
pullback. Iteration and the actual orthogonal hyperbolic-sum isomorphism
supply the formula. Thus no affine-bundle invariance theorem is needed
for the extension.

For reference, the exact affine-space invariance locator requested
during this check is \textbf{SGA4 XV, Corollary2.2}, with the meaning of
acyclicity in \textbf{XV, Proposition1.6 and Definition1.7}. The French
transcription of these statements was read at
\url{https://grothendiecksga.com/read/sga4/fr/XV-2.html} and
\url{https://grothendiecksga.com/read/sga4/fr/XV-1.html} . Corollary2.2
treats affine space over an arbitrary scheme, with torsion primes
invertible on the base. No original SGA4 scan inspection is claimed for
this locator check.

\section{Additional exact source
register}\label{additional-exact-source-register}

The attached primary sources are inventoried and hashed in
\texttt{\small work/\allowbreak{}source\_\allowbreak{}inventory.json}.
The following supplementary results were used at their stated scope.
Each event in \texttt{\small CHECK\_\allowbreak{}EVENTS.tsv} records
whether the full proof, a statement, or a computation was actually
checked.

\begin{longtable}[]{@{}
  >{\raggedright\arraybackslash}p{(\columnwidth - 2\tabcolsep) * \real{0.5000}}
  >{\raggedright\arraybackslash}p{(\columnwidth - 2\tabcolsep) * \real{0.5000}}@{}}
\toprule\noalign{}
\begin{minipage}[b]{\linewidth}\raggedright
Source
\end{minipage} & \begin{minipage}[b]{\linewidth}\raggedright
Exact locator and use
\end{minipage} \\
\midrule\noalign{}
\endhead
\bottomrule\noalign{}
\endlastfoot
Sturmfels, \emph{Algorithms in Invariant Theory}, second edition &
Theorem 2.1.4, printed p.27: finite-group Noether degree bound. Theorems
1.2.4 and 1.2.6: Gröbner bases and normal forms. Algorithm 1.4.5,
printed pp.22--23: Hilbert-basis calculation. These are the finite
algorithmic inputs of G-009--G-011. \\
Bruns--Gubeladze--Trung, \emph{Problems and algorithms for affine
semigroups}, arXiv:math/0101116v4 & Proposition 2.1.1(a): the
affine-semigroup finiteness input used for the balanced monoid. \\
Pak--Panova--Swanson, \emph{A combinatorial interpretation for certain
plethysm and Kronecker coefficients}, 4 November 2025 & Lemma 2.1;
Theorems 3.1 and 3.6; Definitions 3.2, 3.5, 3.8, 4.4 and 5.1;
Propositions 3.4 and 3.9; Theorems 4.5 and 5.2. G-005 proves the exact
mark-to-polytope change. Two local printed proof errors are repaired in
the independent review. \\
Maas--Gariépy--Tétreault, \emph{3-Plethysms of homogeneous and
elementary symmetric functions}, arXiv:2209.14942v1 & Section 4.2,
printed pp.8--9, for the classical Thrall--Chen coefficient formulas;
Theorem 1, Section 7.3, printed p.20, for the tableau formula. G-007
proves the normalizer-trace version. \\
Sarkar, \emph{Stability of plethysm coefficients and modified polynomial
induction}, arXiv:2610.07929v2, 14 October 2026 & Theorem D = 6.4,
Proposition 7.1, Theorem E = 7.3. The exact hypotheses and character
transpose are written in T-010--T-013. \\
Anderson, \emph{Chern class formulas for G2 Schubert loci} & Appendix A,
Proposition A.1 and Corollary A.2, printed p.29: stabilizer and open
orbit. \\
Manivel, \emph{Prehomogeneous spaces and projective geometry} &
Proposition 3.1.4 p.44; Examples 4.3.3 pp.59--60 and 4.4.2 pp.61--62;
Section 5.1.2 p.68: the width-seven orbit, degree-seven boundary and
metric ranks. \\
Brion, \emph{Invariants d'un sous-groupe unipotent maximal d'un groupe
semi-simple}, Ann. Inst. Fourier 33(1), 1983 & Theorem 3 pp.12--13, case
\((A_6,\omega_3)\) in III.1 p.14, and Corollary 2 to Theorem 2: the
precise nonpolynomial invariant-ring calculation used in T-020. \\
Akin--Buchsbaum--Weyman, \emph{Schur functors and Schur complexes}, 1982
& Theorem II.2.16 p.232 and Theorem III.1.4 pp.245--246; exact integral
use and conjugated notation are proved in A-012. \\
Mulmuley, \emph{Geometric Complexity Theory VII: Nonstandard quantum
group for the plethysm problem}, supplied source & Problem 1.1 asks for
an explicit positive \#P formula. Theorem 1.2 states the nonstandard
quantization and complete reducibility/Peter--Weyl assertions. The
subsequent duality, reciprocity and canonical-basis requirements are
stated as unresolved in that source. The present proofs use no such
unproved requirement. \\
\end{longtable}

Primary URLs:
\url{https://www.math.ucla.edu/~pak/papers/plethysm-Kronecker-trees.pdf};
\url{https://arxiv.org/abs/2209.14942};
\url{https://arxiv.org/abs/2610.07929};
\url{https://arxiv.org/abs/math/0101116};
\url{https://people.math.osu.edu/anderson.2804/papers/g2-chern.pdf};
\url{https://manivel.perso.math.cnrs.fr/trento2012.pdf};
\url{https://www.numdam.org/item/10.5802/aif.902.pdf}.

The source-copy manifest records exact local bytes used for
consultation. Full third-party articles downloaded during the run are
not redistributed in the ZIP; their mathematical uses, URLs, hashes, and
the complete current proofs and certificates are included. The
user-supplied transcriptions and literature sources remain included as
supplied input evidence.

\section{Independent full review of the outer-four fiber and polytope
formula}\label{independent-full-review-of-the-outer-four-fiber-and-polytope-formula}

\subsection{Review outcome and scope}\label{review-outcome-and-scope}

\textbf{PASS.} I read and checked the complete proofs in T-032
\(`general_chart_report.md`\), T-035--T-036
\(`outer4/four_fibers_report.md`\), and T-037
\(`outer4/full_outer4_report.md`, including A.29–A.31\). I also checked
the exact complement pairing F.12--F.15, the Elliott character
calculation H.5--H.7, the harmonic character table H.8--H.12, and the
three \(S_3\) series P.10 used by T-037. The checks below establish the
actual module maps, all character twists, saturation in the full
oriented range, and the fixed positive lattice counts.

T-037 gives all coefficients \(p_{\mu,(n)}^\lambda\) with
\(\mu\vdash4\), including all central cases. It supplies at most 108
tagged bounded rational polytopes in at most five variables; the bounds
are 36 for the symmetric and exterior outer labels and 72 for \((2,2)\).
The exact fixed-table identities and arithmetic checks were
independently executed in
\texttt{\small check\_\allowbreak{}outer4\_\allowbreak{}fixed\_\allowbreak{}tables.py};
its receipt is
\texttt{\small OUTER4\_\allowbreak{}FIXED\_\allowbreak{}TABLE\_\allowbreak{}RECEIPT.json}.
I did not rerun the other agent's 1,980-coefficient Murnaghan--Nakayama
sweep and do not count that sweep as my own event.

The only wording correction requested during this review was an explicit
redefinition of (a,b) after the complement operation. The author has
made that correction in A.7. The original and retwisted \(q\)-actions
were checked separately; T-037 uses the correct cancellation.

\subsection{A-024. Actual free decomposition of the incidence
algebra}\label{a-024.-actual-free-decomposition-of-the-incidence-algebra}

Write \(V=\mathfrak h_4\), with its invariant dot product,
\(W=\Lambda^2V=V\otimes\operatorname{sgn}\), and \[
R=\mathbf Q[q\in V,t\in W]/(f),\qquad f=q\cdot t.
\] The notation identifies each self-dual representation with the
corresponding polynomial generator space through the displayed invariant
form. The generator \(f\) has bidegree ((1,1)) and character
\(\operatorname{sgn}\). In particular \(\det V=\operatorname{sgn}\) and
\(\det W=1\).

Let \(\ell_i=e_i-\frac14\mathbf1\). The two quadrics \(U\) vanishing on
their four projective lines give, after the specified linear change of
coordinates, \[
A=\mathbf Q[q]/(U)
=\mathbf Q[X,Y,Z]/(X^2-Z^2,Y^2-Z^2).
\] The monic reductions give an actual free \(\mathbf Q[Z]\)-basis
(1,X,Y,XY). After inverting \(Z\), the ring maps isomorphically to the
product of the four Laurent line rings indexed by
\(X/Z,Y/Z\in\{1,-1\}\). Freeness over \(\mathbf Q[Z]\) makes the map to
the product of the four polynomial line rings injective before
localization. Thus no component or embedded contribution at the origin
has been removed. The two displayed quadrics form a regular sequence in
the polynomial ring: after quotienting by the first, the second is monic
in \(Y\) and hence acts injectively.

On the \(i\)-th component, \(q=s\ell_i\) and \(f=s(\ell_i\cdot t)\).
Both factors are nonzero in that component's polynomial domain. The
product embedding therefore proves that multiplication by \(f\) on
(A{[}t{]}) is injective. With \(\mathcal H=R/(U)\), the resulting exact
Hilbert-series identity is \[
\operatorname{Hilb}\mathcal H=(1-s^2)^2\operatorname{Hilb}R.
\]

Choose a rational section of \(R_{C,b}\to\mathcal H_{C,b}\) in each
finite bidegree, and average it over the 24 elements of \(S_4\). This
gives an actual equivariant section \(\sigma\). The multiplication map
\[
M:\operatorname{Sym}(U)\otimes\mathcal H\longrightarrow R,
\qquad p\otimes h\longmapsto p(U)\sigma(h)
\] is surjective by repeatedly reducing modulo (\(U\)); every reduction
lowers the \(q\)-degree by two. Its source and target have equal
dimensions in every bidegree by the preceding series identity. Hence
\(M\) is an equivariant \(\operatorname{Sym}(U)\)-linear isomorphism.
This also proves algebraic independence of the two quadrics, and that
they form a regular sequence in \(R\). It directly identifies \[
(U)^kR_{c,b}
\simeq\bigoplus_{j=k}^{\lfloor c/2\rfloor}
\operatorname{Sym}^jU\otimes\mathcal H_{c-2j,b}.
\] There is no unproved choice of a multiplicative basis in this
argument. The section is a linear section; its explicit averaged
multiplication map supplies exactly the required free module structure.

\subsubsection{Evaluation and every quotient
piece}\label{evaluation-and-every-quotient-piece}

Evaluation at the actual affine vectors \(\ell_i\) is
\(S_4\)-equivariant with no parity character. It gives \(A_0=\mathbf1\),
\(A_1=V\subset\operatorname{Perm}_4\), and \(A_C=\operatorname{Perm}_4\)
for all \(C\ge2\). An explicit surjectivity check is available without
any projective scaling convention: put \[
E_i(q)=2q_i^2-\tfrac16\sum_lq_l^2.
\] Then \(E_i(\ell_j)=\delta_{ij}\); for \(C\ge2\), the homogeneous
polynomial \(E_i(q)(4q_i/3)^{C-2}\) has the same evaluations. The
Hilbert series of \(A\) gives dimension four in all those degrees. In
degree two the kernel of the evaluation quotient
\(\operatorname{Sym}^2V\to\operatorname{Perm}_4\) is exactly
\(U=[(2,2)]\). Its matching-plane realization gives
\(\det U=\operatorname{sgn}\).

The nonzerodivisor just proved yields the exact equivariant sequence \[
0\longrightarrow\operatorname{sgn}\otimes A_{C-1}\otimes\operatorname{Sym}^{b-1}W
\xrightarrow{f}A_C\otimes\operatorname{Sym}^bW
\longrightarrow\mathcal H_{C,b}\longrightarrow0.
\] For \(C\ge3\), evaluation identifies the first map point by point
with multiplication by \(\ell_i\cdot t\). The stabilizer \(S_3\) of
\(\ell_i\) acts by sign on this line in \(W\), and its quotient is
\(\mathfrak h_3\otimes\operatorname{sgn}\). The explicit rational
area-form intertwiner identifies this quotient with \(\mathfrak h_3\).
Thus the cokernel is
\(D_b=\operatorname{Ind}_{S_3}^{S_4}\operatorname{Sym}^b\mathfrak h_3\).

For \(C=2\), extend the source \(A_1=V\) to \(\operatorname{Perm}_4\).
The extended map remains injective point by point, and has the preceding
cokernel \(D_b\). The quotient of its image by the original image is
\(\operatorname{sgn}\otimes\operatorname{Sym}^{b-1}W\), because
\(\operatorname{Perm}_4/V=\mathbf1\). The resulting short exact sequence
splits by an averaged projection.

For \(C=0\), the quotient is \(\operatorname{Sym}^bW\). For \(C=1\), the
incidence multiplication removes the \(S_{(b,b,1)}V\) term from
\(V\otimes\operatorname{Sym}^b(\Lambda^2V)\); the quotient is the Cartan
term \(S_{(b+1,b)}V\). Twisting \(V=W\otimes\operatorname{sgn}\)
contributes sign because (2b+1) is odd. The Elliott rule with Dynkin
labels ((1,b)) gives \(\bigoplus_{L=1}^{b+1}H_L\), including \(b=0\).
Consequently every row of A.5 follows by an actual exact sequence or
intertwiner, with the sign in the \(C=1\) row required.

\subsection{A-025. Chart, fibers, saturation, and
complement}\label{a-025.-chart-fibers-saturation-and-complement}

\subsubsection{The general chart and the fiber
valuations}\label{the-general-chart-and-the-fiber-valuations}

In T-032 the centered-ratio chart is equivariant because row permutation
fixes \(\mathbf1\), fixes \(u=\prod_i x_i\), and commutes with
centering. Its flag-coordinate map is multiplication by the explicitly
displayed products of first-column variables. A balanced monomial of
weight \((n^m)\) and highest weight \(\lambda\) becomes
\(u^{n-\lambda_2}\) times a tail-flag monomial: the number of columns of
height at least two is exactly \(\lambda_2\), and the remaining exponent
of every \(x_i\) is \(n\). This proves the claimed surjectivity. The
inverse chart prescribing nonzero \(u\), then choosing the \(x_i\) with
that product and the remaining centered columns, proves injectivity on a
dense open set. Polynomiality is precisely divisibility by the stated
power of \(u\), and equality there extends as a polynomial identity.
These are actual algebra maps.

For \(m=4\), the map gives \(q\in V\), \(t\in W\), and the determinant
variable \(z\in\operatorname{sgn}\). At \(x_4=\varepsilon\), use the
affine chart with \(q_4-q_1\ne0\). The two normal coordinates at
\(\ell_4\) are \[
\frac{q_j-q_1}{q_4-q_1}
=\frac{\varepsilon(r_j-r_1)}{S-\varepsilon r_1},\qquad j=2,3.
\] Their leading coefficients vary independently. Setting the last entry
of the third column to zero makes the limiting \(t_1,t_2,t_3\) depend on
the first and third columns and \(S\), independently of those two normal
coefficients; they range over the fiber \(\sum_{i=1}^3t_i=0\). The
fourth column makes \(z\ne0\) on this same generic locus. Therefore the
lowest nonzero transverse term cannot disappear upon substitution, and
the \(x_4\)-valuation equals the ideal order at the generic point of the
entire fiber. Each graded normal quotient is a vector bundle on the
integral fiber, so a section vanishing at its generic point vanishes
everywhere. This upgrades the generic valuation statement to membership
in the full ideal power. Permutation equivariance supplies the other
three fibers.

The \(x_i\) are pairwise nonassociate primes, so the individual
valuations characterize divisibility by \(u^k\). The four fibers are
disjoint, and the incidence projection is a projective-line bundle,
hence flat. Their simultaneous order-\(k\) condition is exactly
\(\pi^*\mathcal I_Z^k\). These facts establish the stated section-space
model by the actual chart substitution.

\subsubsection{Character convention and
orientation}\label{character-convention-and-orientation}

T-035 initially uses \(q\in W\), so it has the outside twist
\(\operatorname{sgn}^{c+d}\). Replacing that \(q\)-action by the
original \(q\in V\) changes the degree-\(c\) section module by
\(\operatorname{sgn}^c\). The two factors cancel and leave precisely
\(\operatorname{sgn}^d\) in A.16 and A.9.

As an independent check that detects this sign, take \(n=4\),
\(\lambda=(11,5,0,0)\), so \(c=5,k=1,b=d=0\). The original-action ideal
piece is
\(U\otimes\operatorname{Perm}_4\oplus\operatorname{Sym}^2U\otimes V\).
It has dimension 17 and transposition trace 1. Independently, the
balanced trace in \(s_{(11,5)}=h_{11}h_5-h_{12}h_4\), at cycle type
((2,1,1)), is the number of triples \(0\le x,y,z\le4\) with
\(2x+y+z=11\), minus the corresponding number for 12. It is \(10-9=1\).
An uncancelled extra sign would give (-1).

The complement map uses the nonzero Schur pairing
\(S_\lambda V_4\otimes S_{\widehat\lambda}V_4\to(\det V_4)^H\), obtained
by complementary exterior columns. It is perfect by irreducibility and
carries the balanced weight to the balanced weight for
\(\widehat n=H-n\). Literal permutation matrices act on the target by
\(\operatorname{sgn}^H\); self-duality of rational Specht modules gives
exactly the conjugation rule in A.6. The operation exchanges the
opposite gaps (a,b) and preserves \(c\). After redefining the gaps for
the resulting partition, \(a\ge b\), and \(c-2k=(a-b)/2\ge0\) whenever
\(k>0\). Thus every input reaches the proved saturation range in at most
one complement operation.

\subsubsection{Saturation above the
boundary}\label{saturation-above-the-boundary}

The two-quadric power resolution pulls back along the flat projection
and becomes \[
0\to\operatorname{Sym}^{k-1}U\otimes\det U\otimes
\mathcal O_{\mathcal F}(c-2k-2,b)
\to\operatorname{Sym}^kU\otimes\mathcal O_{\mathcal F}(c-2k,b)
\to\pi^*\mathcal I_Z^k\otimes\mathcal O_{\mathcal F}(c,b)\to0.
\] For \(c\ge2k+1\), the left degree is at least (-1). The incidence
hypersurface sequence on \(\mathbf P^2\times\mathbf P^2\) proves
\(H^1(\mathcal O_{\mathcal F}(a,b))=0\) for \(a\ge-1,b\ge0\): both
ambient \(H^1\) and the preceding ambient \(H^2\) vanish, since the
latter twists are at least (-2,-1). The same sequence identifies global
sections with \(R_{a,b}\) for \(a\ge0\), and gives zero for \(a=-1\).

Global sections are therefore the cokernel of exactly the multiplication
map obtained by extending the two-variable power-ideal resolution along
\(\operatorname{Sym}(U)\to R\). A-024 proves that \(R\) is free over
this polynomial algebra, so that extension is exact. Its cokernel is
\((J_R^k)_{c,b}\). This proves equality with the sheaf section space by
the same map, including the first endpoint \(c=2k+1\).

At \(c=2k\), the cohomology calculation in T-036 gives the extra summand
\[
\operatorname{sgn}\otimes\operatorname{Sym}^{k-1}U
\otimes\operatorname{Sym}^{b-1}W.
\] The character can also be checked entirely in the original
convention. Here
\(H^2(\mathbf P(V),\mathcal O(-3))=\det V=\operatorname{sgn}\), while
\(\det U=\operatorname{sgn}\); their product makes
\(H^1(\mathcal I_Z^k(2k-1))=\operatorname{Sym}^{k-1}U\). The incidence
injection itself has sign, supplying the displayed remaining character.
An averaged projection splits the resulting exact sequence. Negative
symmetric powers are zero; \(k=0\) is the ordinary whole-ring case and
gives no correction. This proves A.9 for every oriented input.

\subsection{A-026. Exact fixed tags, lattice counts, and
arithmetic}\label{a-026.-exact-fixed-tags-lattice-counts-and-arithmetic}

The harmonic trace series on the five classes are \[
\frac{1+q}{(1-q)^2},\quad\frac1{1+q},\quad\frac1{1+q},\quad
\frac{1+q}{1+q+q^2},\quad\frac{1+q}{1+q^2}.
\] They follow by dividing \(1-q^2\) by the characteristic polynomial on
\(W\). Using the explicit five-row character table and class sizes
(1,6,3,8,6), I checked every rational identity giving the stated
\(E_\tau\) numerator over \((1-q^4)(1-q^6)\). I also checked all three
\(S_3\) averages giving \(F\), all 75 relevant tensor entries by
character inner products, and the induction branching giving \(G\). All
entries of the tensor table are zero or one.

The Elliott proof H.5--H.7 is a full character calculation: its finite
geometric sum distinguishes both parities, the extra \(\delta/(1+z)\)
term cancels from \(P(z)-z^{-1}P(z^{-1})\), and the remaining six
monomials match the specialized Weyl numerator. Thus the use of its
((1,b)) specialization in A.5 has no unproved branching step.

In the positive module sum, put \(C=c-2j\). The terms with \(C\ge3\) are
exactly A.25, with \(j=f+2x_2+3x_3\) recoverable uniquely from each
lattice point. Odd \(c\) leaves \(C=1\), producing A.27 and its required
label conjugation. Even \(c\) leaves \(C=0\), and also \(C=2\) precisely
when \(c\ge2k+2\); these give the displayed A.26 term and the two A.5
terms. At \(c=2k\), only the boundary correction remains in addition to
\(C=0\). Consequently A.28 neither omits nor duplicates any summand.

The exact numbers of fixed tags are:

\begin{longtable}[]{@{}lrrr@{}}
\toprule\noalign{}
Target label & A.25 & A.26 or A.27 & Full even bound \\
\midrule\noalign{}
\endhead
\bottomrule\noalign{}
\endlastfoot
((4)), \((1^4)\) & 6 & 12 & 36 \\
((2,2)) & 12 & 24 & 72 \\
((3,1)), ((2,1,1)) & 18 & 36 & 108 \\
\end{longtable}

Every variable has a positive coefficient in an equality or upper bound.
Thus each listed region is a bounded rational polytope. Its constants
are (O\(n\)), including after complement, and the matrices are fixed.

The elementary scalar formulas A.29--A.31 are also correct. Splitting
\(y=2t+e\) proves the arithmetic-progression formula for
\(P(N)=\#\{2x+3y\le N\}\); the \(T_e=-1\) endpoint contributes zero. The
weight-2 coordinate in A.26 is the uniquely determined slack variable
after division by two. For A.27 the lower exclusion is exactly
\(P(\lfloor-e/2\rfloor)\), including \(e=0\), where it removes \(L=0\).
I directly enumerated the scalar formulas for \(-12\le N\le120\) and all
310 intervals \(0\le b\le30,0\le e\le9\), in addition to the algebraic
proofs.

Each fixed template count is \(O((n+1)^3)\); its free count has
dimension at most three. The finite number of arithmetic operations uses
integers of \(O(\log(n+2))\) bits. Thus the stated schoolbook
\(O((\log(n+2))^2)\) arithmetic bound follows for the coefficient
evaluator. The witness enumerator is separately output-dependent, as its
implementation states.

\subsubsection{Endpoint checks}\label{endpoint-checks}

\begin{itemize}
\tightlist
\item
  \(n=0\), \(\lambda=0\): \(c=b=d=k=0\). Only \(\mathsf W_\mu(0,0)\)
  survives, giving 1 for \(\mu=(4)\) and zero for the other four labels,
  exactly the Schur functors of the one-dimensional
  \(\operatorname{Sym}^0\).
\item
  \(b=0\): every term with (b-1) vanishes; the \(C=1\) row is the
  correct standard \(V\), since \(\operatorname{sgn}\otimes H_1=V\).
\item
  \(c=0\): \(k=0\), and only the degree-zero quotient piece remains.
\item
  \(c=1\): \(k=0\), and the only term is \(\mathsf L_{\nu'}(0,b)\).
\item
  \(c=2k+1\): the bulk interval is empty, so only the \(C=1\) term
  remains.
\item
  \(c=2k\): the bulk and \(C=2\) terms are absent; the two \(\mathsf W\)
  terms give exactly the boundary sequence.
\item
  Valid partitions with more than four parts have zero coefficient: four
  successive horizontal-strip additions produce at most four rows in the
  ambient tensor power, and the Schur symmetrizer maps the requested
  outer functor into that power.
\end{itemize}

All these checks use the exact maps or formulas above. No unresolved
module, saturation, parity, or positive-count gate remains for outer
degree four in the stated family. This review makes no assertion of
historical novelty.

\section{Independent review of G-015 and
G-016}\label{independent-review-of-g-015-and-g-016}

\subsection{Status}\label{status}

\textbf{PASS, complete proof read.} The general projective-frame section
model and the full quadratic ideal calculation are correct in their
stated characteristic-zero scope. The original \(S_m\) action and
determinant line are retained throughout. The proof gives the actual
comparison map and the ordinary sheaf ideal powers; it proves the needed
equality of orders without assuming a higher-cohomology vanishing
statement.

I separately opened the original primary source
\href{https://arxiv.org/pdf/alg-geom/9411014}{P. Magyar, \emph{A
Borel--Weil Construction for Schur Modules}, arXiv:alg-geom/9411014v1}
and read Proposition 1, printed pp.~5--6, the Plücker determinant bundle
convention on p.~7, and Proposition 4\(a\)--\(c\), printed p.~10. These
give the dual Schur convention and the precise full-flag pullback used
in GF.2 and GF.7. All maps in this specialization are defined over
\(\mathbf Q\). The source's algebraically closed field convention in §2
causes no rational-field gap: global sections commute with field
extension, as follows from an affine-cover equalizer and flatness, and a
map becoming an isomorphism after the faithfully flat extension
\(\mathbf Q\subset\overline{\mathbf Q}\) was already an isomorphism.

\subsection{1. The exact general chart and its
linearization}\label{the-exact-general-chart-and-its-linearization}

The chart T-032 takes a smaller flag polynomial \(F\) to \[
u^{n-\lambda_2}F(v(X)),\qquad
v_{j,I}=\left(\prod_{i\notin I}x_i\right)p_{j,I},
\quad u=\prod_i x_i.
\] The map \(w\mapsto\mathbf1\wedge w\) identifies the smaller exterior
coordinates with the displayed minors. Row permutations fix
\(\mathbf1,u\), commute with centering, and act on every exterior power
by its literal permutation matrix. Thus the chart preserves the exact
\(S_m\) action. The full-height smaller flag factor is \(\det E\); its
sign character is the last factor of \(\mathcal L_\tau\), not a separate
omitted twist.

The highest-weight and row-weight calculation in C.12--C.13 gives
precisely \(\lambda_2\) factors of height at least two and the common
residual exponent \(n\) for every \(x_i\). The inverse dense chart C.14
recovers any prescribed smaller matrix and any nonzero \(u\).
Consequently polynomial extension is exactly divisibility by \(u^k\), as
asserted in GF.8. The invariant form on \(E\) is nondegenerate and
equivariant over \(\mathbf Q\), so its Schur functor gives the specified
map from the section convention \(S_\tau E^*\) to \(S_\tau E\).

\subsection{2. Every higher flag frame is independent of the leading
normal
directions}\label{every-higher-flag-frame-is-independent-of-the-leading-normal-directions}

At \(x_m=\varepsilon\), with \(s_m=S\ne0\), the first-line coordinates
are exactly \[
a_j=\varepsilon\frac{r_j-r_1}{S-\varepsilon r_1},
\qquad 2\le j\le m-1.
\] The (m-2) coordinates are a regular parameter system at \([\ell_m]\)
in \(\mathbf P(E)\). Pullback along the flag bundle gives a regular
parameter system transverse to its whole fiber.

The subfamily \(X_{mj}=0\) for \(j\ge3\) is sufficient to test the
possible leading-coefficient cancellation. At \(\varepsilon=0\), a
nonzero minor must contain row \(m\); expansion along that row uses only
the entry \(S\) in column two. After factoring out the \(x_i\), every
remaining minor is an exterior coordinate in the columns \[
\mathbf1,\quad (X_{i3}/x_i)_{i<m},\quad\ldots,
(X_{im}/x_i)_{i<m}.
\] No such coordinate depends on the \(r_j\). The quotient
\(E/\mathbf Q\ell_m\) is exactly \(\mathbf Q^{m-1}/\mathbf Q\mathbf1\):
restricting the first (m-1) coordinates is surjective onto that
quotient, with kernel \(\mathbf Q\ell_m\). The remaining ratio columns
are arbitrary, so their classes realize every generic flag and every
choice of nonzero successive frame scalings in the quotient. Varying
\(S\prod_{i<m}x_i\) supplies the first-line frame; successive column
scalings supply all higher exterior frames. The last one is nonzero
whenever those quotient columns form a basis, so the determinant factor
also has valuation zero.

Meanwhile the (m-2) differences \(r_j-r_1\) vary freely with all those
image flag and frame coordinates fixed. This proves the required
algebraic independence over the image fiber's function field. It is
stronger than independence of only the first line: it includes every
higher coordinate which can occur in \(F\).

\subsection{3. Equality of orders and the ordinary ideal
powers}\label{equality-of-orders-and-the-ordinary-ideal-powers}

Choose local nonzero Plücker coordinates at the generic fiber point. On
the exhibited family each has valuation zero, so the local
trivialization of \(\mathcal L_\tau\) introduces no extra factor of
\(\varepsilon\). A section with lowest nonzero transverse term of degree
\(h\) has leading coefficient equal to that nonzero polynomial evaluated
at the independent \((r_j-r_1)/S\). It cannot vanish. The opposite
inequality follows because every transverse parameter is divisible by
\(\varepsilon\). This proves exact equality of the original
coordinate-divisor valuation and the fiber-ideal order.

The special subfamily used here need not dominate the entire original
divisor: the lower-order divisibility is valid on its full open chart,
and nonvanishing on this subfamily proves nonvanishing of the universal
leading coefficient. Its independence statement is exactly what is
needed for the reverse inequality.

The generic order detects the ordinary sheaf power everywhere on the
fiber. Indeed the fiber is a regular embedding and its graded normal
quotients are locally free over an integral flag variety. A section of
such a quotient which is zero at the generic point is zero globally.
Induction through the first \(k\) graded pieces therefore gives
membership in \(\mathcal I_i^k\) at every point, not just symbolic
membership at one generic point.

The original \(x_i\) are pairwise nonassociate primes. Hence the \(m\)
valuations are equivalent to divisibility by \(\prod_i x_i^k\). The
fibers are disjoint, so their powers are pairwise comaximal and
intersection equals product. Flatness of the flag projection identifies
that product with \(\pi^*\mathcal I_{Z_m}^k\). Applying the actual chart
map proves GF.4, and the specified Specht covariant functor proves GF.5.

The cases \(k=0\), \(n=0\), and \(m=3\) cause no change in the argument.
For \(m=3\), the quotient flag fiber is a point, the transverse system
has one parameter, and the same independence test still applies.

\subsection{4. The full quadratic ideal and its
action}\label{the-full-quadratic-ideal-and-its-action}

Set \(r=m-1\ge3\) and use the \(r\) coordinate points plus
\([1:\cdots:1]\). In degree two, vanishing at the coordinate points
removes every square, while vanishing at the last point imposes exactly
one sum-zero relation on the mixed coefficients. This gives the listed
differences \(z_i z_j-z_1z_2\).

In their quotient, write \(z=z_1z_2\). For distinct (i,j), choose
\(a\notin\{i,j\}\); then \[
z_i z=z_i(z_jz_a)=z_j(z_i z_a)=z_jz.
\] Every mixed degree-\(C\) monomial contains a cross-product, and this
identity replaces every remaining factor by \(z_1\). Thus the pure
powers \(z_i^C\) and the single mixed class \(zz_1^{C-2}\) span.
Evaluation at the coordinate points isolates the coefficients of the
pure powers, and evaluation at the last point then isolates the mixed
coefficient. They are independent. This proves the asserted basis in
every degree, including the low degrees separately.

The resulting graded map to the product of line coordinate rings is
injective. Its kernel before quotienting is exactly the intersection of
the homogeneous linear prime ideals of those lines. Hence the generated
ideal is radical and saturated as well as the full projective ideal. The
dimension of its quadratic piece is \(m(m-3)/2\).

Returning to the exact affine representatives \(\ell_i\), evaluation is
equivariant without a scaling character. The two eigenvalues for the
matrix of the values of \(q_i^2\) are \((m-2)/m\) and \((m-1)/m\); for
\(q_i^3\) they are \(1-3/m+3/m^2\) and \((m-1)(m-2)/m^2\). These are
nonzero for the stated \(m\). The invariant quadratic has the same
nonzero value at each point and propagates surjectivity by two degrees.
The established dimensions make evaluation an isomorphism onto
\(\operatorname{Perm}_m\) for every \(C\ge2\). Degrees zero and one give
its constant and standard submodules. All the actions in GF.15 are
therefore literal evaluation maps.

For \(m=3\), the three differences \(q_i-q_j\) are distinct line
equations on \(E\), each vanishing at one frame point. Their product
gives exactly the reduced degree-three divisor ideal and changes by sign
under permutations. This is the stated separate cubic formula.

\subsection{\texorpdfstring{5. Independent check supplied for the next
\(m=5\)
computation}{5. Independent check supplied for the next m=5 computation}}\label{independent-check-supplied-for-the-next-m5-computation}

The subsequent cubic and full resolution are being written by the theta
agent in T-038. I supplied an independent verification of its
quadratic-subalgebra kernel; the full resolution was not audited in this
event.

For the five quadrics \[
A=z_1z_3-z_1z_2,\ B=z_1z_4-z_1z_2,\
C=z_2z_3-z_1z_2,\ D=z_2z_4-z_1z_2,\ E=z_3z_4-z_1z_2,
\] the nonzero cubic \[
F=AD(E-B-C)-BC(E-A-D)
\] maps to zero; exact symbolic expansion confirms it. Put \(x=z_1z_2\).
In the image fraction field, \[
x=\frac{AD}{E-A-D},\qquad
z_1^2=\frac{x(x+A)}{x+C}.
\] All six pair products follow from \(x,A,B,C,D,E\), and all \(z_i\)
follow from them and \(z_1\). Thus \(\mathbf Q(z_1,z_2,z_3,z_4)\) is
algebraic of degree two over the image field, so that field has
transcendence degree four. The cubic is primitive linear in \(E\): its
coefficient \(AD-BC\) is irreducible, and does not divide its constant
term, as the substitution \((A,B,C,D)=(1,2,3,6)\) gives constant 12.
Gauss's lemma proves irreducibility. The kernel is consequently the
height-one prime \((F)\).

Under the actual coordinate transposition \(z_1\leftrightarrow z_2\),
the quadrics transform by \(A\leftrightarrow C\),
\(B\leftrightarrow D\), and \(E\mapsto E\), so \(F\mapsto-F\). The
degree-three kernel is one-dimensional and stable under \(S_5\), and a
one-dimensional character of \(S_5\) is determined by a transposition.
Hence this relation has sign character, agreeing with theta's
independent calculation. This is an explicit algebra map and relation,
not a dimension-only comparison.

\section{Independent complete reviews of the five-point
formulas}\label{independent-complete-reviews-of-the-five-point-formulas}

Reviewer: GPT-6 Astra Pro, run /root/archive\_audit. The reasoning
setting is not exposed. The complete reports T-038, T-039 and G-017 were
read, their arguments checked, and the finite identities below
independently executed. The events are recorded in CHECK\_EVENTS.tsv.
This review concerns the precise families and maps stated in those
reports.

\subsection{A-030. Complete review of
T-038}\label{a-030.-complete-review-of-t-038}

\textbf{Status: PASS.} The full source is
work/theta/outer5/five\_point\_resolution\_report.md, equations
(5R.1)--(5R.20).

\subsubsection{The reduced scheme and the complete
resolution}\label{the-reduced-scheme-and-the-complete-resolution}

Let \(S=\mathbf Q[x,y,z,w]\), and let \(I\) be generated by the five
displayed differences of cross-products. In \(S/I\), all six
cross-products have the same image \(p\). With a third distinct variable
available, \(z_i p=z_jp\). Thus every mixed monomial of degree \(C\ge2\)
equals \(px^{C-2}\). The four pure powers and this mixed monomial have
independent evaluations at the four coordinate points and at
\((1,1,1,1)\). This proves the stated homogeneous basis and the
injection into the product of the five line-coordinate rings. The kernel
is consequently the intersection of the five homogeneous line ideals.
This proves reducedness and saturation directly, as well as the complete
Hilbert series used later.

I independently multiplied the displayed matrices and verified
\(f^tM=0\) and \(MN=0\). The specified four-by-four minor is nonzero.
The entries \(y(x-w)\) and \(z(y-w)\) already have greatest common
divisor one. Over the fraction field, the kernel of \(M\) is therefore
the line generated by \(N\); clearing a reduced denominator and using
primitivity shows that its kernel over \(S\) is precisely \(SN\). The
Hilbert-series comparison with \(\ker f^t\) then proves exactness in
every degree. The injectivity at the left and positivity of every
entry's degree also prove minimality.

The change-of-basis matrix has determinant \(-1\), satisfies \(N=Pf\),
and turns \(M\) into the alternating matrix in (5R.14). All five signed
maximal Pfaffians independently agree with \(f\). These calculations
give the actual height-three alternating presentation needed in A-029.

\subsubsection{Group actions and the cubic
map}\label{group-actions-and-the-cubic-map}

The affine representatives \(\ell_i=e_i-\frac15\mathbf1\) are permuted
exactly. Their use in evaluation proves
\((S/I)_C\simeq\operatorname{Perm}_5\) for every \(C\ge2\), including
odd \(C\). The fifth representative has coordinates \((-1,-1,-1,-1)\) in
the chosen basis; the proof of independence remains valid with that
representative. There is therefore no unrecorded degree-dependent
character in (5R.9).

For each of the seven conjugacy classes, I checked the rational
character identity \[
\frac{1+\chi_E(g)t+\chi_{\operatorname{Perm}_5}(g)t^2/(1-t)}
     {\operatorname{tr}(g\mid\operatorname{Sym}E)}
=1-\chi_U(g)t^2+\chi_{\epsilon U}(g)t^3-\chi_\epsilon(g)t^5.
\] It verifies the generator, first-syzygy and last-syzygy labels
simultaneously. The report's transposition calculation also identifies
the final one-dimensional representation without requiring this larger
check.

For the subalgebra map \(\mathbf Q[A,B,C,D,E_0]\to S\), the cubic \(F\)
vanishes by expansion. Its image field recovers \(xy\), all other
cross-products, and \(x^2\) by the rational expressions in (5R.16).
Adjoining \(x\) recovers \(y,z,w\), so the image has transcendence
degree four. Viewed as a polynomial in \(E_0\), \(F\) has irreducible
leading coefficient \(AD-BC\); the test \((A,B,C,D)=(1,2,3,6)\) proves
that coefficient does not divide the constant term. Gauss's lemma proves
irreducibility, and the height-one prime kernel is exactly \((F)\). This
verifies the algebra isomorphism, rather than merely a relation in the
image.

The six forms have rank five, sum zero and sum of cubes \(-24F\). I
separately constructed the six subgroups of order five in the specified
lexicographic order. For each of the four adjacent transpositions,
direct substitution in the quadratic forms gives \[
g(t_i)=-t_{g(i)}.
\] The target indices are exactly those obtained by conjugating the
subgroups. These four generators prove the full \(S_5\)-intertwining
identity. In particular the signed six-letter action used subsequently
in G-017 is verified with its actual indexing.

The primary representation conventions and geometric comparison were
read in Dolgachev--Farb--Looijenga, author PDF, §2.3 Table 1, p.~10, and
Remark 2.4, p.~13, as recorded in A-028. The report's explicit cubic
morphism proves the needed identification independently of the
historical attribution.

\subsection{A-031. Complete review of
T-039}\label{a-031.-complete-review-of-t-039}

\textbf{Status: PASS.} The full source is
work/theta/outer5/five\_point\_family\_report.md, including its endpoint
and support arguments.

\subsubsection{The bundle map and its
linearization}\label{the-bundle-map-and-its-linearization}

For the stated partition, \[
\lambda_2-n=1,\qquad
(\lambda_2,\lambda_3,\lambda_4,\lambda_5)
=(d+a+2,d+a,d+a,d).
\] The actual G-015 chart therefore gives the flag section module with
ideal factor \(\mathcal I_Z\). The relative Plücker map on the quotient
flag gives \[
(\det E^*)^d\otimes\mathcal O(a+2)\otimes S_{(a,a)}Q^*.
\] The exterior pairing \(\Lambda^2Q^*\simeq Q\otimes\det Q^*\), its
symmetric power, and the determinant of the tautological sequence give
\[
(\det E^*)^{d+a}\otimes\mathcal O(2)\otimes\operatorname{Sym}^aQ.
\] For rank three, the asserted identification of \(S_{(a,a)}Q^*\) with
the symmetric power of \(\Lambda^2Q^*\) has the specified highest-vector
map and equal dimensions \(\binom{a+2}{2}\). It is valid also for
\(a=0\).

The local product description of the flag bundle verifies projection
with the possibly nonflat ideal. Over a trivializing open set, the
relevant Čech complex is a complex of rational vector spaces tensored
with the base ring. Its kernel and image sequences split over
\(\mathbf Q\); tensoring by the ideal preserves this calculation of
degree-zero cohomology. The local maps agree under changes of quotient
frame. This proves the stated projection map with the ideal factor.

Since \(\det E=\epsilon\), the resulting sign is precisely
\(\epsilon^{d+a}\). No determinant factor is lost in either exterior
pairing.

\subsubsection{Cohomology and the positive exact
sequence}\label{cohomology-and-the-positive-exact-sequence}

Multiplication by the tautological line gives the exact symmetric-power
sequence (5F.11). Its fiberwise kernel and injectivity are correct. On
\(\mathbf P(E)\), parametrizing lines, the canonical bundle is
\(\mathcal O(-4)\otimes\det E^*\); thus
\(H^3(\mathcal O(-4))=\det E=\epsilon\). This agrees with the ordered
Čech computation.

The three cohomology rows in (5F.12) follow from the full long exact
sequences: \(B_a\) has only \(H^0=\operatorname{Sym}^aE\); \(B_a(-1)\)
is acyclic; and \(B_a(-3)\) has only
\(H^2=\epsilon\operatorname{Sym}^{a-1}E\). The last group is zero at
\(a=0\).

Sheafifying the exact T-038 resolution and tensoring by \(B_a(2)\)
produces the two short exact sequences in the report. They give
\(H^0(K)=0\) and \(H^1(K)=\operatorname{Sym}^{a-1}E\). The last-syzygy
sign and projective-cohomology sign multiply to the trivial character.
Consequently the actual sequence is \[
0\longrightarrow U\otimes\operatorname{Sym}^aE
\longrightarrow H^0(\mathcal I_Z(2)\otimes B_a)
\longrightarrow\operatorname{Sym}^{a-1}E\longrightarrow0.
\] The first arrow is quadratic multiplication followed by the
quotient-bundle map; the last arrow is the connecting map from the
stated resolution. Averaging a section over \(S_5\) gives the direct-sum
isomorphism. This proves (5F.2) in every \(a,d\ge0\).

\subsubsection{All fixed tags and
endpoints}\label{all-fixed-tags-and-endpoints}

I independently enumerated the 26 standard tableaux of the seven shapes
of size five and computed their stated comajor indices. Every entry,
including the repeated value \(5\) for shape \((3,1,1)\), agrees. I then
verified all seven rational identities for the multiplicities in
\(\operatorname{Sym}^aE\): \[
\sum_{g\text{ classes}}\frac{|g^{S_5}|\chi_\rho(g)}{120}
\frac{1-t}{\prod_{\ell\in\operatorname{cycles}(g)}(1-t^\ell)}
=\frac{\sum_{r\in D_\rho}t^r}
{(1-t^2)(1-t^3)(1-t^4)(1-t^5)}.
\] The 49 entries of the tensor table \(U\otimes[\rho]\) were
independently computed by exact character inner products. They all
agree. The resulting numbers of tagged polytopes, in the report's row
order, are \[
6,\ 24,\ 30,\ 36,\ 30,\ 24,\ 6.
\] These totals equal \(6\dim[\nu]\) exactly.

The twist replaces \(\mu\) by \(\mu^t\) exactly when \(d+a\) is odd. In
the symmetric outer case, the even list is \(1,\ldots,6\), and the odd
list is \(4,\ldots,8,11\). The semigroup generated by \(2,3,4,5\) is
\(\{0\}\cup\mathbf Z_{\ge2}\). This proves the stated support thresholds
\(a\ge1\) and \(a\ge4\), including the consecutive first two degrees
that handle residual zero and the missing residual one. At \(a=0\), all
counting terms have the correct value for the single twisted module
\(U\). Every polytope is bounded because its nonnegative coordinates
have positive weights.

\subsection{A-032. Complete review of
G-017}\label{a-032.-complete-review-of-g-017}

\textbf{Status: PASS.} The full source is
work/root/five\_point\_q\_module\_report.md, equations (Q.1)--(Q.26).

\subsubsection{The actual section module and its
generators}\label{the-actual-section-module-and-its-generators}

The partition has \(\lambda_2-n=k\). Its relative quotient-flag weight
is \((1,1,0)\), so the same exterior pairing as above gives the exact
module \[
(S_\lambda\mathbf Q^5)_{(n^5)}
\simeq\epsilon^{d+1}M_k,\qquad
M_k=H^0(\mathcal I_Z^k(2k)\otimes Q).
\] The specified inverse chart and G-015 prove its extension criterion.
The parity convention for \(\nu\) in (Q.3) follows: it equals \(\mu\)
for odd \(d\) and \(\mu^t\) for even \(d\).

The quadratic vector \((z_i^2)_i\) vanishes in the quotient fibers at
all five points. Its invariant description in centered coordinates
differs from this vector by \(-\frac25(\sum z_i)z\), a tautological-line
vector. This proves that \(\theta\) is an invariant section of the
stated sheaf.

I independently formed \(P=\operatorname{adj}K/F\) from the printed
\(K\). Every entry is polynomial. The determinant identities, both
matrix products, the last row and all twenty line-corrected substitution
identities (Q.12) hold exactly. The nonzero two-by-two minor agrees with
the printed expression.

The relation-space linearization is justified by the report's
injectivity argument. A degree-three polynomial double at the four
coordinate points is a combination of the four squarefree cubics. The
four derivatives at the fifth point force every coefficient to vanish.
Thus \(H^0(\mathcal I_Z^2(3))=0\). Together with the absence of a
degree-two cubic-ring relation, this proves the required injectivity of
\(E\otimes\operatorname{Sym}^2U\to M_2\). It makes the coefficient
vector attached to \(f\theta\) unique; applying the group action proves
that the five relations form \(U\).

\subsubsection{Injection and surjection in every
degree}\label{injection-and-surjection-in-every-degree}

The adjugate argument establishing \(C\)-torsion-freeness of the
cokernel is valid. In \(gv=Pw\), where \(F\nmid g\), each component of
\(\operatorname{adj}P\,v\) is divisible by \(F^3\), because
\(\det P=81F^3\) and \(S\) is a unique-factorization domain. The
resulting polynomial vector \(u\) satisfies \(Pu=v\).

Over the discrete valuation ring \(S_{(F)}\), the elementary divisors of
\(P\) have exponents at most one by \(PK=3FI\). Their sum is three by
the determinant. The cokernel therefore has rank three over \(C\). Upon
extending its fraction field to \(\mathbf Q(z_1,z_2,z_3,z_4)\), the four
constant generators span the three-dimensional quotient by the
tautological line. The resulting generic map is an isomorphism; its
kernel before localization is \(C\)-torsion and hence zero. This proves
the actual injection into the section module.

A-029 proves the needed \(2k\)-linear resolutions for every \(k\ge2\),
saturation in degree \(2k\), and \(H^1(\mathcal I_Z^k(2k-1))=0\). The
latter degree is a sheaf-section calculation: these sections may be
nonzero even though the ordinary polynomial ideal power has minimum
degree \(2k\). Formula (Q.15) correctly uses the fat-point length and
sheaf sections.

The tautological sequence is locally split, so it remains exact after
tensoring with the ideal. Its cohomology gives \[
\dim M_k=4\left(\binom{k+4}{4}-\binom{k+1}{4}\right)
-\left(\binom{2k+2}{3}-5\binom{k+2}{3}\right)
=\frac{(k+1)(k+2)(3k+4)}2.
\] At \(k=1\), the additional \(H^1(\mathcal I_Z(1))=\mathbf1\) gives
\(21\); at \(k=0\), the space is \(E\) of dimension four. Both values
agree with the polynomial. Thus the injective presentation and \(M\)
have equal dimensions in every degree, proving surjectivity and all the
stated relations.

\subsubsection{The invariant parameters and positive
formula}\label{the-invariant-parameters-and-positive-formula}

The successive-root construction proves that the polynomial ring in six
variables is free over the elementary symmetric coefficient ring. Each
adjoining step is a monic quotient of the indicated degree; after the
final root, coefficient elimination identifies the result with the
original polynomial ring. Quotienting by \(e_1\), replacing \(e_5\) by
\(e_5^2\), and substituting \(e_3=-8F\) therefore prove that \(S\) is
free of rank \(1440\) over \[
\mathbf Q[e_2,e_4,e_6,e_5^2,F].
\] The signed six-letter action makes the first four displayed
generators invariant.

After any preceding parameter quotient, both the next parameter and
\(F\) remain nonzerodivisors. The determinant \(81F^3\) then proves that
\(P\) stays injective after quotienting by the next parameter. Reducing
\(pv=Pw\) modulo that parameter gives \(w=pw'\), and cancellation gives
\(v=Pw'\). This proves regularity on the cokernel at every step.

The finite quotient \(H=M/B_+M\) consequently has the stated Hilbert
series, degree 19 and dimension 4320. A graded averaged section followed
by multiplication gives \(B\otimes H\to M\); degree induction proves
surjectivity and the regular-sequence Hilbert identity proves
injectivity. This is the required actual equivariant free decomposition.

I independently derived all seven rational characters of \(M\) by Newton
identities from the cycle-type character of \(U\). All 140
multiplicities in the displayed \(H\)-table agree. Its row totals are \[
36,\ 144,\ 180,\ 216,\ 180,\ 144,\ 36,
\] and its total dimension is 4320. Every nonidentity character vanishes
at \(t=1\), which verifies the regular-character check.

All support criteria follow from the first two consecutive nonzero
entries of each row and the weight-two parameter. This includes
\(k=0,d=0\), where the balanced module is \(\epsilon E=[2,1,1,1]\). The
residue decomposition for weights \(1,2,3,5\) gives exactly (Q.26), with
residue polynomial degree \(29+28+27+25=109\) and at most four summands.
An independent dynamic-programming count verified it for arguments
\(0,\ldots,240\); the proof by unique residue decomposition establishes
it for every nonnegative argument.

\subsection{Check files and precise
scope}\label{check-files-and-precise-scope}

The independently written check script is
work/archive/check\_outer5\_fixed\_data.py. Its executed receipt is
work/archive/OUTER5\_FIXED\_DATA\_RECEIPT.json. It verifies the exact
fixed identities described above. One initial script assertion compared
expanded and factored polynomial expressions structurally; changing that
comparison to expanded polynomial equality made the intended check
exact. No mathematical statement was changed.

The whole arguments, including all zero and parity cases, have been
reviewed separately from these calculations. This event does not claim
to rerun the authors' larger variable coefficient sweeps or their
witness enumeration. Those checks retain their own authorship and
receipts.

\section{Five-point jet characters and independent positive-formula
checks}\label{five-point-jet-characters-and-independent-positive-formula-checks}

Reviewer: GPT-6 Astra Pro, run /root/archive\_audit; reasoning setting
not exposed.

\subsection{A-033. Exact jet maps and complete independent review of
G-019}\label{a-033.-exact-jet-maps-and-complete-independent-review-of-g-019}

\textbf{Status: PASS.} I read the entire report
work/root/five\_point\_symmetric\_quotient\_report.md, including
(SQ.1)--(SQ.16), every endpoint, the full zero set and the stated
framework transport. I independently reconstructed all 403,200 fixed
coefficients by the following geometric calculation.

\subsubsection{1. A uniform equivariant jet
description}\label{a-uniform-equivariant-jet-description}

Let \(E=\mathfrak h_5\), let \(\ell_i=e_i-\frac15\mathbf1\), and let
\(Z\) be the five corresponding projective points. Use the invariant
inner product to write \[
E=\mathbf Q\ell_i\oplus W_i,\qquad
\alpha_i(v)=\frac{\langle v,\ell_i\rangle}{\langle\ell_i,\ell_i\rangle}.
\] The affine chart \(\alpha_i\ne0\) is identified with \(W_i\) by \[
[v]\longmapsto v/\alpha_i(v)-\ell_i.
\] This map and its inverse \(w\mapsto[\ell_i+w]\) commute with the
stabilizer \(S_4\) of the point. The degree-\(c\) line bundle is
trivialized on this chart by \(\alpha_i^c\). These choices are
transported exactly when the five points are permuted.

Consequently the order-\(k\) jets of \(\mathcal O(c)\) at this point are
the actual module \[
\bigoplus_{j=0}^{k-1}\operatorname{Sym}^j(W_i^*).
\] The quotient by the \(k\)-th power of the coordinate maximal ideal
gives this decomposition directly. The invariant form identifies
\(W_i^*\) with \(W_i\), and \(E=\mathbf1\oplus W_i\) under the
stabilizer. The complete five-point jet module is induced from this
stabilizer module. This specifies the action for every integer \(c\);
there is no scalar rescaling of the affine representatives.

For \(g\in S_5\), let \(m_1(g)\) be its number of fixed letters, and put
\[
h_r(g)=\operatorname{tr}(g\mid\operatorname{Sym}^rE),\qquad h_r=0\quad(r<0).
\] Only fixed points contribute to the trace on the five-point direct
sum. At any such point, \[
\sum_{j=0}^{k-1}
\operatorname{tr}(g\mid\operatorname{Sym}^jW_i)
=h_{k-1}(g).
\] This identity is the decomposition of
\(\operatorname{Sym}^{k-1}(\mathbf1\oplus W_i)\). Therefore the exact
jet character is \[
\boxed{\operatorname{tr}\bigl(g\mid
H^0(\mathcal O/\mathcal I_Z^k(c))\bigr)
=m_1(g)h_{k-1}(g).}
\tag{J.1}
\] For \(k=0\), both sides are zero.

\subsubsection{2. The two scalar Euler terms and the symmetric quotient
map}\label{the-two-scalar-euler-terms-and-the-symmetric-quotient-map}

The restriction sequence \[
0\longrightarrow\mathcal I_Z^k(c)
\longrightarrow\mathcal O(c)
\longrightarrow\mathcal O/\mathcal I_Z^k(c)
\longrightarrow0
\] is equivariant. Its cohomology, with the complete vanishings in
A-029, gives \[
\operatorname{tr}(g\mid C_k)
=h_{2k}(g)-m_1(g)h_{k-1}(g),
\tag{J.2}
\] and \[
\operatorname{tr}\bigl(g\mid
H^0(\mathcal I_Z^k(2k-1))-H^1(\mathcal I_Z^k(2k-1))\bigr)
=h_{2k-1}(g)-m_1(g)h_{k-1}(g).
\tag{J.3}
\] The notation in (J.3) denotes the exact representation-group class of
the displayed cohomology sequence.

These equations include the endpoints. At \(k=0\), (J.3) is zero. At
\(k=1\), it is \(-\mathbf1\), since \(h_1(g)=m_1(g)-1\); the actual
groups are \(H^0(\mathcal I_Z(1))=0\) and
\(H^1(\mathcal I_Z(1))=\mathbf1\). At \(k\ge2\), the latter cohomology
vanishes by A-029. In degree \(2k\), \(H^1\) vanishes for every \(k\),
including the direct degree-two evaluation check at \(k=1\).

Multiplication by the tautological line gives the exact, locally split
symmetric-quotient sequence used in (SQ.4). Tensoring with the ideal
preserves it. Its cohomology, including the \(k=1\) connecting term,
therefore gives \[
\boxed{\begin{aligned}
\operatorname{tr}(g\mid M_{a,k})
={}&h_a(g)\bigl(h_{2k}(g)-m_1(g)h_{k-1}(g)\bigr)\\
&+h_{a-1}(g)\bigl(m_1(g)h_{k-1}(g)-h_{2k-1}(g)\bigr).
\end{aligned}}
\tag{J.4}
\] This independent formula uses the standard representation and the
actual Taylor maps. It does not require the cubic presentation or
G-017's fixed character table.

The flag comparison in G-019 has point order \(k\), first-line gap
\(2k\), and relative weight \((a,a,0)\). Its first-line weight after the
full determinant is \(a+2k\). The exterior pairing contributes
\((\det Q^*)^a\), giving \(\mathcal O(2k)\) and sign \(\epsilon^{d+a}\).
Thus (SQ.3) and its parity convention agree with the actual local frame
maps. At \(a=0\), the second term in (J.4) is zero. At \(a=k=0\), the
module is trivial before the determinant twist, including the
empty-partition endpoint.

\subsubsection{3. Independent reconstruction of every fixed
coefficient}\label{independent-reconstruction-of-every-fixed-coefficient}

For each of the seven cycle types, I formed \[
\det(1-tg\mid E)=
\frac{\prod_{\ell\in\operatorname{cycles}(g)}(1-t^\ell)}{1-t}
=1+d_1t+\cdots+d_4t^4.
\] The exact recurrence \[
h_0=1,\qquad
h_r=-\sum_{j=1}^{\min(4,r)}d_jh_{r-j}
\] computes all required characters through degree 478. Equations (J.2)
and (J.3) give the two scalar arrays, and (J.4) gives their products
with \(h_a,h_{a-1}\).

I computed the order-zero through order-three step-60 differences of
each separate array. Multiplying those differences according to (J.4),
followed by the seven fixed character inner products, gives the mixed
differences directly. This differs from computing a two-dimensional
multiplicity grid and then taking differences of that grid.

All 403,200 resulting integers agree exactly with the complete root
table. Every entry is nonnegative. The \((3,3)\) entries are zero. The
other fourteen positive-order positions have strictly positive minima
for every label, as stated in the report.

The rational-function verification was independent as well. With
\(H(x)=1/\det(1-xg\mid E)\), the even and odd subsequences are \[
\sum_{k\ge0}h_{2k}t^k
=\left.\frac{H(x)+H(-x)}2\right|_{x^2=t},
\quad
\sum_{k\ge0}h_{2k+1}t^k
=\left.\frac{H(x)-H(-x)}{2x}\right|_{x^2=t}.
\] Using these expressions in (J.2)--(J.3), I checked all 28
one-variable factors against the common denominator \((1-t^{60})^4\).
Their numerators are polynomials of degree at most 239. This proves the
residue-polynomial argument from the first nonnegative quotient value,
including its initial zero binomial terms.

The two-variable Newton identity consequently holds for all \(a,k\),
with degree at most three in each quotient variable. Its bounded
polytope conversion uses the two fixed residue equations, one bounded
multiplicity coordinate, and two increasing subset chains. The claimed
upper bounds of 57,600 tags and nine variables are valid. The constant
number of arithmetic operations and the quadratic schoolbook bit-cost
bound follow because all binomial degrees and table dimensions are
fixed.

I also verified the dimension identity (SQ.15). The common factor \[
C_k+\bigl(M_k-4C_k\bigr)=(k+1)(2k+1)
\] at the level of dimensions gives exactly the displayed expression
after applying Pascal's identity in \(a\). The exceptional \(k=1\) term
is retained.

The exact zero list has 44 triples \((\nu,a,k)\), all in the first
60-by-60 square, and agrees entry by entry with the report. Strict
positivity of the \((1,0)\) and \((0,1)\) coefficients on every residue
excludes any further zeros. This proves the stated complete support,
including both axes and all determinant parities.

The independent script is
work/archive/check\_symmetric\_quotient\_by\_jets.py. The full receipt,
including target hash, separate coefficient minima and maxima, all
rational bounds and all 44 zeros, is
work/archive/SYMMETRIC\_QUOTIENT\_JET\_RECEIPT.json.

\subsection{A-034. Complete review of T-040 and a second jet
reconstruction}\label{a-034.-complete-review-of-t-040-and-a-second-jet-reconstruction}

\textbf{Status: PASS.} I read all of
work/theta/outer5/q\_plus\_report.md, equations (QP.1)--(QP.23),
including the matrix presentation, its equivariant splitting, every
Newton coefficient formula and the complete support argument.

The ordinary/saturated comparisons in (QP.5) use precisely the proved
degree ranges of A-029. The degree-one part of the Rees presentation is
the cokernel of \(\mathsf B\); the identity \(\mathsf B\mathsf Q=FI_4\)
shows that its additional \(F\)-relations already lie in that image. All
five maximal minors are the stated \(-FT_i\), and the primitive column
\(T\) generates the entire polynomial kernel. The sign labels follow
from the actual change of basis in the T-038 resolution.

The tautological sequence gives the exact injection \[
\mathcal A_k\longrightarrow E\otimes\mathcal C_k,\qquad
s\longmapsto\sum e_i\otimes x_i s.
\] The quotient is \(\mathcal N_k\) because
\(H^1(\mathcal I_Z^k(2k))=0\), including \(k=0,1\). This identifies the
cokernel in (QP.13) with the section module.

In the kernel proof, the three-by-three minor \(AD(E_0-B)\) is not
divisible by the irreducible cubic \(F\). Thus \(\mathsf B\) has rank
three modulo \(F\). From \(Y\mathsf B^t=hI_4\), passage to the fraction
field of \(R/(F)\) proves \(F\mid h\). Subtracting \((h/F)\mathsf Q^t\)
leaves rows in \(RT\). Multiplication by \(\mathsf B^t\) then proves
injectivity of the first map. The report's averaging of the quadratic
lift gives the asserted equivariant free decomposition of this kernel;
changing that lift only adds a vector in the preceding \(rT^t\) summand.
All maps and characters therefore agree.

The negative highest coefficient after multiplying the dimension series
by the four parameter factors gives the stated concrete obstruction to
regularity of that particular parameter sequence. It does not interfere
with the subsequent exact resolution and positive Newton count.

For an independent character calculation, the same jet map (J.1) gives
\[
\boxed{\operatorname{tr}(g\mid\mathcal N_k)
=e_g h_{2k+1}(g)-h_{2k}(g)
-m_1(g)(e_g-1)h_{k-1}(g),}
\tag{J.5}
\] where \(e_g=m_1(g)-1\). This follows from the actual quotient-bundle
exact sequence, with the scalar jet terms substituted. It includes
\(\mathcal N_0=\operatorname{End}(E)/\mathbf1\).

I derived its seven rational series by even/odd extraction from
\(H(x)\). They agree exactly with (QP.16), and every numerator over
\((1-t^{60})^4\) has degree 238. The recurrence calculation from (J.5)
reproduces all 1,680 Newton constants, their three closed positive
difference formulas, their maximum 197280 and all seven support sets.
The coefficient bounds used in the positivity proof are valid because
the absolute character values are at most the representation dimensions.

The fixed-residue interpolation applies from \(q=0\), and its
degree-three Newton basis gives the displayed positive count. The
bounded multiplicity coordinate contributes exactly the fixed integer
coefficient; the increasing chain counts the appropriate binomial
coefficient. The limits of 240 tags, five variables and four active tags
are correct, and the determinant parity gives exactly the two
symmetric-outer thresholds.

The independent check script is
work/archive/check\_q\_plus\_by\_jets.py. Its executed receipt is
work/archive/Q\_PLUS\_JET\_RECEIPT.json. The five signed maximal minors
and the specified rank-three minor were also independently expanded. The
earlier A-029 certificate already independently verifies
\(\mathsf B\mathsf Q=FI_4\).

\part{Unresolved questions and exact remaining gates}

\section{Unresolved items and exact next
operations}\label{unresolved-items-and-exact-next-operations}

G-009--G-011 prove terminating positive formulas for every coefficient
and every actual canonical-map image multiplicity. Uniform polynomial
bounds on their preprocessing and witness data remain open. The exact
missing bounds concern the Hilbert basis, invariant list, and
elimination presentation.

The complete direct outer-degree-four formula is proved in T-037,
including every central unequal-gap shape. The next direct-evaluation
gate is outer degree five and higher. G-015 identifies every coefficient
with the covariants of an explicit projective-frame ideal power on a
flag variety. T-038 and A-029 evaluate the five-point ideal algebra and
its ordinary powers. G-019 evaluates every symmetric quotient power on
the quadratic ray, with a fixed positive count and complete support.
G-018, T-040 and T-041 evaluate additional exterior, shifted and hook
families. Remaining direct objects are the other first-line twists and
arbitrary two-row Schur powers of the rank-three quotient, with their
exact flag comparison already proved.

For the canonical map, G-013 gives prime-adjacent semigroup preimages,
G-014 the complete long-first-row range, and R-016 determines three
finite boundary kernels. R-017 identifies common polarization with exact
block truncation; R-018--R-020 eliminate its complete second-row range
and first-row boundary band. R-021--R-023 construct an injection of the
critical canonical annihilator into the remaining common kernel tensored
with one row, after removing the rectangular prolongation by an actual
preimage. Its remaining first-factor labels have first row at least r²+1
and second row at least r(r−1). R-024 gives explicit recursive
second-row ranges from the full degree-five base and every established
square base. The complete varying critical-boundary kernel remains the
next rank gate; all component statements retain their actual maps.

For width seven, T-018 evaluates the end-node sector and G-010 supplies
a general finite count. The smaller direct-formula gate is the remaining
discriminant-valuation filtration on the stated G2 invariant spaces.

The archive has 123 assertion rows, of which 18 retain precise gates
after inspection of all thirteen original target leaves and relevant
neighboring pages. These include admitted crystalline Lefschetz/freeness
inputs, the projected motivic self-duality map, reconstruction from one
unlabeled triangulation, and the unspecified completed
characteristic-class target. Literal type errors are refuted with their
corrected maps proved. The full classification of coefficient rings
remains open after the DVR theorem and explicit local counterexample.

The Alpöge theorem is accepted as requested. Its transfer to the
projectors and canonical-map image in this paper still needs the
explicit equivariant representation/period morphism. The named Lean
source files and scan report are absent. The requested Masullo work has
no supplied statement or identifier. No locator or missing comparison
has been invented.

No formal progress review is due before three elapsed hours, at 03:01:08
Europe/Berlin. This checkpoint retains active work and makes no
assertion that the stopping rule has been reached.

\end{document}
